\documentclass{report}
\usepackage[T1]{fontenc}
\usepackage[french,english]{babel}

\usepackage[colorinlistoftodos,draft]{todonotes}

\usepackage{marginnote}
\setlength{\marginparwidth}{2cm}
\usepackage[dvipsnames,x11names]{xcolor}
\usepackage{pgfplots,diagbox,colortbl,tikz,graphicx,algorithm2e,cancel,verbatim,
listings,float,amsmath,amssymb,enumitem,array,subfiles,bussproofs,
rotating,MnSymbol,mathtools,subcaption,caption}
\pgfplotsset{compat=1.18}
\usetikzlibrary{automata, positioning,graphs,shapes, arrows, calc}
% CORRECT LOCATION (preamble):
\usepackage{csquotes}

\hyphenation{ma-chine-check-a-ble}
\hyphenation{dou-ble-push-out}
%%%%%%% careful : order is important begin
\usepackage{amsthm}
\usepackage{thmtools,hyperref}
\newtheorem{theorem}{Theorem}[chapter]
\newtheorem{corollary}[theorem]{Corollary}
\newtheorem{lemma}[theorem]{Lemma}
\newtheorem{idea}[theorem]{Idea}
\newtheorem{definition}[theorem]{Definition}
\newtheorem{proposition}[theorem]{Proposition}
\newtheorem{remark}[theorem]{Remark}
\newtheorem{example}[theorem]{Example}
\newtheorem{Plan}[theorem]{Plan}
\newtheorem{claim}[theorem]{Claim}
\newtheorem{notation}[theorem]{Notation}
\renewcommand{\chapterautorefname}{Chapter}
\usepackage[section]{placeins}
\usepackage{minted}



\newcommand{\set}[1]{\{#1\}}
\newcommand{\graphbox}[8]{
  \begin{scope}[xshift=#2,yshift=#3]
    \draw [rounded corners=2mm] (0,0) rectangle (#4,-#5);
    \node at (0,0mm) [anchor=north west,inner sep=1mm] {#1};
    \begin{scope}[xshift=#4/2+#6,yshift=#7] 
    #8
    \end{scope}
  \end{scope}
}
\newcommand{\card}[1]{|#1|}
\newcommand{\homset}[2]{\operatorname{Hom}(#1,#2)} 
\newcommand{\isdef}{\overset{\operatorname{def}}{=}}
\newcommand{\cat}[1]{\mathcal{#1}}
\newcommand{\f}[1]{\mathfrak{#1}}
\newcommand{\rs}[1]{\mathcal{#1}}
\newcommand{\opn}[1]{\operatorname{#1}}
%arrow
\newcommand{\itrs}{\overset{i}{\to}} 

\usepackage{tocloft}

% % Customizing the table of contents
% % \setlength{\cftpartindent}{0pt} %  space before part titles
\setlength{\cftchapindent}{15pt} %  space before chapter titles
\setlength{\cftsecindent}{30pt} %  space before sec titles
\setlength{\cftsubsecindent}{40pt} %  space before subsec titles
\setlength{\cftsubsubsecindent}{50pt} %  space before subsubsec titles


\newcommand{\textdef}{\text{$\operatorname{Def.}$}}
\newcommand{\pbpop}{$\operatorname{PBPO}^+$}
\addtolength{\cftbeforechapskip}{-10pt} % Reduce spacing by 10pt
\addtolength{\cftbeforepartskip}{-20pt} % Reduce spacing by 10pt
\newif\iflongversion
\longversiontrue
\newif\iftrackChange
\trackChangetrue
\newcommand{\trackedtext}[1]{%
    \iftrackChange
        \textcolor{red}{#1}% Apply red color if trackChange is true
    \else
        #1% No color change if trackChange is false
    \fi
} 


% \usepackage[style=alphabetic, maxbibnames=999]{biblatex}
\usepackage[maxbibnames=999]{biblatex}
\addbibresource{bib/these.bib} % Add your bibliography file here

\usepackage{titlesec} 
 
\usepackage{tikz-cd}
\usepackage{pifont}      
\usepackage{makecell}    
\usepackage{nicematrix}  
\usepackage{wrapfig}     
\usepackage{lipsum}      

% 

% \usepackage{diagbox,colortbl,tikz,graphicx,algorithm2e,cancel,verbatim,
% listings,float,amsthm,amsmath,amssymb,enumitem,array,subfiles,bussproofs,
% rotating,MnSymbol,hyperref,mathtools,subcaption,caption}
\begin{document}  
   

\begin{titlepage}

    \unitlength 1cm
    \begin{center}
    
    \vspace*{-2.5cm}
    \begin{figure}[h]
        \centering
        \includegraphics[width=0.8\textwidth]{LogoLyon1Off_CoulCmjn300dpi.jpg}
    \end{figure}
    
               {\large\bf THÈSE de DOCTORAT DE L'UNIVERSITÉ DE LYON\\}
   
               {\large \bf l'Université Claude Bernard Lyon 1\\}
    \vspace{12pt}
               {\large \textbf{Ecole Doctorale} 512 \\InfoMaths}

    \vspace{12pt}
    
               {\large \textbf{Spécialité de doctorat~:} Informatique \\}
            %    {\large \textbf{Discipline~:} (éventuellement)}
    
    \vspace{0.8cm}
    
               {Soutenue publiquement le jj/mm/aaaa, par~:\\}
               {\Large\bf {Qi QIU}}
    \vspace{0.5cm}           
    
    
    \rule{5cm}{1pt}
    
    \vspace{12pt}
               
               {\LARGE \bf Automated proofs of termination: contributions to graph rewriting using extended weighted type graph and morphism counting}\\
               \vspace{0.3cm}
               {\Large \bf }
               
    \vspace{12pt}
    \rule{5cm}{1pt}
    
    \vspace{0.5cm}
    
    \end{center}
    
    Devant le jury composé de~:\\
    
    \small {
    \begin{tabular}{ll}
    % \textbf{NAME}           &\textbf{Président.e}\\
    %   Grade/qualité, Établissement/entreprise\\
    % \textbf{Clara Bertolissi}           &\textbf{Examinatrice}\\
    %   Grade/qualité, Établissement/entreprise\\
    % \textbf{Évelyne Contejean}           &\textbf{Examinatrice}\\
    %   Grade/qualité, Établissement/entreprise\\
    \textbf{Jean Krivine}           &\textbf{Rapporteur}\\
      Grade/qualité, Établissement/entreprise\\
    \textbf{Pascale Le Gall}           &\textbf{Rapporteure}\\
      Grade/qualité, Établissement/entreprise\\
    % \textbf{Philippe Malbos}           &\textbf{Examinatrice}\\
    %   Grade/qualité, Établissement/entreprise\\
    % \textbf{Frédéric Prost}           &\textbf{Examinatrice}\\
    %   Grade/qualité, Établissement/entreprise\\
    \textbf{Xavier URBAIN}           &\textbf{Directeur de thèse}\\
      Grade/qualité, Établissement/entreprise\\
    % \textbf{NAME}     &\textbf{Invité.e} (le cas échant)\\
    %   Grade/qualité, Établissement/entreprise\\
    \end{tabular}
    }
    
    %\vfill \clearpage{\pagestyle{empty}\cleardoublepage}
    
    %-------------------------------- Remerciement  -------------------
    \end{titlepage}  
 
\newpage
\section*{Abstract}
Graph rewriting systems are formal models that describe how graphs can be transformed by applying a set of rules. The field of automated termination analysis of graph rewriting systems studies techniques for proving that a graph rewriting system terminates, meaning that, for any graph, the application of its transformation rules cannot last indefinitely. 

The contribution of this thesis is threefold.
Firstly, we extend an existing technique, called the type graph method, for graph rewriting systems based on the double-pushout (DPO) approach. The type graph method assigns weights to morphisms targeting a weighted type graph over a semiring, and the weight of a graph is defined as the sum of the weights of all morphisms from that graph to the type graph. The termination property is then proved by showing that the weight of the graph strictly decreases with each application of a transformation rule if the set of weights is well-founded. However, requiring the weights to be well-founded can have practical limitations because the complexity of searching for a suitable type graph is extremely high.
Our extension reduces the complexity of searching for a suitable type graph by allowing the use of weights from non-well-founded semirings. This enables more efficient implementation of the method in practice in some cases.

Secondly, we develop a new machine-checkable sufficient condition, called subgraph counting, for injective DPO graph rewriting systems on edge-labeled directed multigraphs. It is based on the idea that the transformation cannot last indefinitely if the number of a specific subgraph in the graph strictly decreases every time a transformation is performed.
Thirdly, we implement an automated termination tool, called LyonParallel, for edge-labeled directed multigraph transformation, including both the existing
type graph method and our extension, as well as our subgraph counting technique. 







% We study automated techniques for proving the transformation of any graph according to a set of graph transformation rules cannot last indefinitely.
%  The contribution of this thesis is threefold. 
%  Firstly, we extend an existing technique, called the type graph method.
%    Secondly, we develop a new machine-checkable sufficient condition 
%    for injective DPO graph rewriting systems on edge-labeled directed multigraphs based on the idea that if the number of a specific subgraph in the graph strictly decreases every time a transformation is performed, then the transformation cannot last indefinitely.
%     Thirdly, we implement a termination tool, called LyonParallel, including the existing 
%     type graph method and our extension for edge-labeled directed multigraph transformation, and our techniques.
 
\newpage  
\section*{Resumé}
La réécriture de graphe est un modèle formel qui décrit comment des graphes peuvent être transformés en appliquant un ensemble de règles appelé système de réécriture de graphes. Nous nous intéressons à la preuve automatique de la terminaison de ces systèmes, c'est-à-dire aux techniques permettant de montrer que, pour tout graphe, l'application des règles de transformation d'un système donné ne peut se répéter indéfiniment.

Nous contribuons au développement de techniques automatiques pour l'analyse de la terminaison des systèmes de réécriture de graphes avec double pushout (DPO), en proposant quatre apports.

Premièrement, nous étendons une technique existante,
reposant sur des graphes de type. Cette méthode attribue des poids à des morphismes ciblant un graphe de type; le poids d'un graphe sujet à transformation est défini comme la somme des poids de tous les morphismes depuis ce graphe vers le graphe de type. La terminaison est alors établie en montrant que le poids du graphe diminue strictement à chaque application d'une règle de transformation, à condition que l'ensemble des poids soit bien fondé. Exiger que les poids soient bien fondés peut, cependant, se révéler limitant en pratique, car la complexité de la recherche d'un graphe de type adapté est extrêmement élevée. Notre extension réduit cette complexité en autorisant des poids issus d'ensembles non bien fondés. \todo{todo to do : il faut une analyse du gain en complexite dans le document}

Deuxièmement, nous développons une nouvelle condition suffisante, appelée comptage de morphismes. Cette méthode s'applique aux systèmes de réécriture de graphes avec DPO à règles injectives sur des multigraphes orientés dont les arêtes sont étiquetées. Elle repose sur l'idée que la réécriture ne peut pas durer indéfiniment si le nombre de morphismes ayant un sous-graphe donné comme domaine diminue strictement à chaque transformation.

Troisièmement, nous étendons notre méthode de comptage de morphismes pour une analyse plus fine. Plus particulièrement, nous considérons les morphismes ayant un sous-graphe spécifique comme domaine et dont l'image n'est pas incluse dans un contexte interdit.

Nous mettons enfin en œuvre ces techniques en les implantant au sein d'un nouvel outil logiciel : LyonParallel. Dédié à la transformation avec DPO des multigraphes orientés étiquetés sur les arêtes, cet outil permet d'exhiber automatiquement la terminaison d'exemples que les approches antérieures ne pouvaient pas traiter.
  
 
\newpage  
\section*{Acknowledgements} 
First of all, I thank my supervisor, Xavier Urbain, for providing me with the opportunity to work on this subject and for his guidance not only in my research but, more importantly, in my personal development: his integrity is a great example for me to follow. Midway through my PhD, we had different opinions regarding the research direction, and he generously allowed me to pursue the path which inspired me the most. This freedom has been invaluable to me, and I am grateful for his kindness. I am also grateful for his careful, tireless reading of this dissertation; when I initially disagreed with his suggestions he was patient and later restated them in a way I found considerate. His help significantly improved the quality of my writing.

I thank Evelyne Contejean and Victor Ostromoukhov for serving as members of the \enquote{Comité de Suivi de Thèse}; Pierre Courtieu for his help in the early stages of my PhD and for his extraordinary kindness; Rachid Echahed and Nicolas Peltier for hosting me as temporary teaching and research associate (ATER) at the University of Grenoble Alpes;
Tom Hirschowitz for his invitation to give a talk at the University of Chambéry, which meant a lot to me as a follower of Jean-Jacques Rousseau; 
Zhilin Wu for his constructive questions on the first draft;
the anonymous reviewers of my work for their time and valuable feedback; and, last but not least, Pascale Legall and Jean Krivine for agreeing to serve as the examiners of this dissertation.

I extend my sincere thanks to all the faculty members who contributed to my academic development during my studies in France, for their generosity in sharing knowledge, and for their kindness and patience towards their students. I am especially grateful to Benoît Mariou, who taught my first course in logic; it was in that class that I realized logic was what I had been looking for.

I thank my parents, Shuangluan Ye and Guowen Qiu, for their unconditional love and sacrifices that have allowed me to pursue my dreams. I also thank my wife, Lu Wang, for her patience and understanding during the challenging times of my PhD.

I thank my friend Pierre Bourquat, the most authentic person I have ever met. I still remember one day at the very beginning of my PhD, when I was working in my cubicle alone, he gently tapped on the glass side of my cubicle and asked me if I was new ... then introduced me to other friends Thomas Ranvier, Auday Berro, Mehdi Hennequin and Yacine Gaci. I also thank Zhongyun Yang, Yaqing Fang, Jiayi Wei, Chukun Zhang, Xin Xie and all those who have come into my life, as each person has played a role in shaping who I am today.

The author acknowledges the use of large language models, but solely for the purpose of proofreading and wording edits.

Finally, I would like to express my gratitude to the French National Research Agency (ANR) for funding the project Safe, Adaptive, and Provable Protocols for Oblivious Robots Operation (SAPPORO) under Grant 2019-CE25-0005, which has provided essential resources and support for my research endeavors.  

\newpage       
 
  
\tableofcontents   
\newpage      

%%%%%%%%%%%%%%%%%%%%%%%%%%%%%%%%%%%%%%%%%%%%%%%%%%%
%%%%   Chapter : Non-Well-Founded Semirings   %%%%%
%%%%   what did I do ?
%%%%%%%%%%%%%%%%%%%%%%%%%%%%%%%%%%%%%%%%%%%%%%%%%%%
\chapter{Introduction}  

% \paragraph{Ensuring the correctness of distributed algorithms is crucial} Distributed algorithms play an important role in our modern society. For example, they control machines that we use in daily basis such as smart phone but also critical systems such as medical devises, trains, airplanes, spaceship. Errors in distributed algorithms can have severe consequences, such as the loss of critical data or even the loss of human lives. 


% \paragraph{Ensuring the correctness of distributed algorithms is challenging}
% It is challenging to ensure the correctness of distributed algorithms.

% \paragraph{Interactive theorem proving can be used to ensure the correctness of distributed algorithms, but they are too demanding for most users}
% Formal methods can address this challenge by rigorous mathematical methods. Among these methods, mechanically verified proofs using proof assistants can provide a high level of assurance in the correctness of distributed algorithms.
% However, this approach is too demanding for most users, as they require a high level of expertise and are often time consuming.  

% \paragraph{Automated tools that check specific properties of distributed algorithms can alleviate the problem}
% Automated tools that check specific properties of distributed algorithms and can generate certificates that can be checked by proof assistants alleviate this problem: they can be used by users who are not experts in formal methods, and they can also speed up the verification process for experts. 


%   type graph method for edge-labeled directed multigraph transforms and the subgraph counting technique in a unified tool, called, together with the existing type graph method by Endrullis and Overbeek~\cite{overbeek2024termination_lmcs} for edge-labeled directed multigraph transforms.

\paragraph{Distributed algorithms and their correctness}
Distributed systems play a crucial role in improving performance, 
scalability, and efficiency in modern computing.
For example, they control machines that we use in daily basis such as smart phone but also critical systems such as medical devises, trains, airplanes, spaceship. Errors in distributed algorithms can have severe consequences, such as the loss of critical data or even the loss of human lives. 
Ensuring their correctness is vital but challenging due to the inherent complexity of these systems~\cite{heiser2010theroad, lamport2019thebyzantine}.

\paragraph{Automated methods for verifying distributed algorithms}
While interactive theorem proving~\cite{harrison2014history} offers a promising approach to formally verifying properties of distributed systems~\cite{plump2024formalisingDPO,potop2019formal,courtieu2016certified},  
it requires proficiency in proof assistants and can be
time-consuming. Automated methods that check specific properties of distributed algorithms and can generate certificates that can be checked by proof assistants provide a practical alternative~\cite{contejean2011automated,giesl2014proving}: while they can only be applied to specific properties for some classes of distributed algorithms, they can be used by users who are not experts in formal methods, and they can also speed up the verification process for experts. To develop such automated methods, distributed algorithms need to be modeled in a formal way enabling rigorous analysis and automated verification of their properties.


 \paragraph{Graph transformation}
 Graph transformation provide an intuitive yet formal way to model distributed systems: computation units are represented by nodes; communication canals are represented by edges; states of the system are modeled by graphs whose edges have labels representing information encoded in computation units and states of communication canals; state changes of the system is modeled by replacing subgraphs in the graph by other subgraphs according to transformation rules.

 A predominant school of thought in this area is called algebraic graph rewriting which uses concepts from category theory to define the transformation of graphs. One advantage of this approach is that the language of category theory abstract away from the details of the different notions of graphs, provided that the notion of graph satisfies some properties. In fact, different notions of graphs have been proposed in the literature for different applications, such as edge-labeled directed multigraphs (todo ref), single graphs (todo ref), hypergraphs (todo ref), attributed graphs (todo ref), etc. Another advantage is that transformation is defined up to isomorphism, which simplifies reasoning about graph transformation.

  To transform a host graph $G$ using a rule which replaces an occurrence of a graph $L$ in $G$ by the occurrence of another graph $R$, one firstly needs to decompose the host graph $G$ into a context $C$ and an occurrence of $L$, then remove some nodes and edges of the occurrence of $L$ and finally adds some fresh nodes and edges to form the result graph. For example, if we consider a graph transformation rule which identifies an occurrence of the following graph, then removes the middle node and edges of the occurrence

    \begin{tikzpicture}
           \graphbox{\( \)}{0mm}{-3mm}{34mm}{12mm}{2mm}{2mm}{
                \coordinate (o) at (0mm,-8mm); 
                \node[draw,circle] (l1) at ($(o)+(-10mm,0mm)$) {1};
                \node[draw,circle] (l2) at ($(l1)+(2,0)$) {2};
                \node[draw,circle] (l3) at ($(l1) + (1,0)$) {3};
                \draw[->] (l1) -- (l3) node[midway,above] {a};
                \draw[->] (l3) -- (l2) node[midway,above] {a};
            } 
    \end{tikzpicture}

  The graph $G$, illustrated below, can be transformed into the graph $H$, illustrated below, by applying the rule.

  \begin{tikzpicture}
     \graphbox{\( G \)}{0mm}{-22mm}{34mm}{22mm}{2mm}{-3mm}{
              \coordinate (o) at (0mm,-3mm); 
              \node[draw,circle] (l1) at ($(o)+(-10mm,0mm)$) {1};
              \node[draw,circle] (l2) at ($(l1)+(2,0)$) {2};
              \node[draw,circle] (l3) at ($(l1) + (1,0)$) {3};
              \node[draw,circle] (l4) at ($(l2) + (0,-1)$) {6};
              \draw[] (l1) -- (l3) node[midway,above] {a};
              \draw[] (l3) -- (l2) node[midway,above] {a};
              \draw[ ] (l2) -- (l4) node[midway,right] {b};
              \node[draw,circle] (l6) at ($(l1) + (0,-1)$) {7};
              \draw[] (l1) -- (l6) node[midway,left] {b};
          }    
          \graphbox{\( H \)}{40mm}{-22mm}{34mm}{22mm}{2mm}{-3mm}{
              \coordinate (o) at (0mm,-3mm); 
              \node[draw,circle] (l1) at ($(o)+(-10mm,0mm)$) {1};
              \node[draw,circle] (l2) at ($(l1)+(2,0)$) {2};
              \node[draw,circle] (l4) at ($(l2) + (0,-1)$) {6};
              \draw[ ] (l2) -- (l4) node[midway,right] {b};
              \node[ draw,circle] (l6) at ($(l1) + (0,-1)$) {7};
              \draw[ ] (l1) -- (l6) node[midway,left] {b};
          }    
  \end{tikzpicture}

  There are different algebraic graph transformation systems according to whether the rule mentioned above can be applied to the following graph $G'$:

  \begin{tikzpicture}
    \graphbox{\( G' \)}{0mm}{-22mm}{34mm}{22mm}{2mm}{-3mm}{
                \coordinate (o) at (0mm,-3mm); 
                \node[draw,circle] (l1) at ($(o)+(-10mm,0mm)$) {1};
                \node[draw,circle] (l2) at ($(l1)+(2,0)$) {2};
                \node[draw,circle] (l3) at ($(l1) + (1,0)$) {3};
                \node[draw,circle] (l4) at ($(l2) + (0,-1)$) {6};
                \draw[->,red] (l3) -- (l4) node[midway,above] {a};
                \draw[->] (l1) -- (l3) node[midway,above] {a};
                \draw[->] (l3) -- (l2) node[midway,above] {a};
                \draw[->] (l2) -- (l4) node[midway,right] {b};
                \node[draw,circle] (l6) at ($(l1) + (0,-1)$) {7};
                \draw[<-] (l1) -- (l6) node[midway,left] {b};
                \draw[->] (l2) edge[out=-135,in=-45]node[midway,below] {a} (l1) ;
            }   
  \end{tikzpicture}

   The most studied algebraic graph transformation systems are double-pushout (DPO) graph rewriting systems~\cite{corradini1997algebraic}. It does not allow this kind of transformation because when the node 3 is removed, the edge from node 3 to node 6 becomes dangling. Single pushout (SPO) graph rewriting system~\cite{corradini1997algebraic} allows the application of the rule to transform the graph $G'$ and remove all dangling edges. PBPO and PBPO+ graph rewriting systems avoid implicit edge deletion by providing a mechanism for each rule to impose restrictions on the context $C$. (to do SqPO and Agree)

   We focus on DPO graph rewriting systems in this thesis, because it is the most studied and the easier in the sense that its definition uses only the concept of pushout as one of the algebraic graph rewriting's problem is that it is not easy to understand the definition of a graph rewriting system. DPO graph rewriting systems are also powerful enough to model many distributed algorithms and provide a simple formalism easier to develop termination techniques because it does not allow implicit edge deletion and creation.



\paragraph{DPO graph rewriting systems on edge-labeled directed multigraphs}
 While termination techniques can be developed for DPO graph rewriting systems using different notions of graphs that satisfy certain properties, we focus on developing termination technique for DPO graph rewriting systems on edge-labeled directed multigraphs.
 There are two main reasons for this choice.
 Firstly, we want our technique to be easy to check by software engineers who are not experts in category theory and graph rewriting.  
 By this choice, we hope that our automated techniques can be adopt with the real confidence.
 For example, consider the new automated technique presented in \autoref{chap:subgraph_counting}, its proofs using the language of category theory are long and detailed but can be understood with very basic knowledge in graph theory and very basic knowledge in category theory.
Secondly, edge-labeled directed multigraphs are a powerful and elegant notion of graphs that can be used to model many distributed algorithms.


\paragraph{Relative Termination}
We focus on relative termination property. Given two rule sets \( \mathcal{A} \) and \( \mathcal{B} \), relative termination ensures that, for any graph finite graph $G$,
rules in $\mathcal{A}$ can only be applied finitely many times in any sequence of rule applications from $\mathcal{A} \cup \mathcal{B}$. 


\paragraph{Existing techniques}
There are not many techniques for proving relative termination of DPO graph rewriting systems. Below, we list some of the existing techniques.
 
\begin{itemize}
    \item Overbeek and Endrullis have developed a termination technique for PBPO+ graph rewriting systems~\cite{overbeek2024termination_lmcs}. This technique counts counts the number of occurrences of certain subgraphs before and after a rewriting steps. This technique can be applied to PBPO+ graph rewriting systems on many categories including edge-labeled directed multigraphs. It can also be applied to left-injective DPO graph rewriting systems.
    \item The type graph method, which weighs an object by summing the weights of morphisms from the object to a type graph, was initially introduced by Zantema, K{\"o}nig and Bruggink \cite{zantema2014termination} for cycle-rewriting systems. 
    This method has since been generalized to edge-labeled multigraphs by Bruggink et al. \cite{bruggink2014termination} for DPO rewriting with monic matches and injective rules, later extended to DPO rewriting in general by Bruggink et al. \cite{bruggink2015proving}, and further adapted to more categories and different DPO variants by Endrullis et al. \cite{endrullis2024generalized_arxiv_v2}. 
    \item Plump \cite{plump1995ontermination} introduced a necessary and sufficient termination condition for left-injective DPO hypergraph rewriting via forward closure, though verifying this condition is undecidable. 
    \item Plump \cite{plump2018modular} later proposed a modular critical pair-based strategy for left-injective DPO hypergraph rewriting with monic matches. 
    \item Levendovszky et al. \cite{levendovszky2007termination} propose a termination criterion for DPO rewriting (monic matches, injective rules, negative application condition), though automated verification is hard as explained in \cite[\textsection 6]{levendovszky2007termination}. 
    \item Bottoni et al. \cite{bottoni2005termination} propose a termination criterion for DPO/SPO rewriting on high-level replacement units. Their method imposes a strongly constrained measuring function and the only concretes measuring function proposed are node-counting and edge-counting.
    \item Bottoni et al. \cite{bottoni2010atermination} presents a criterion for termination of DPO rewriting with monic matches, injective rules and negative application conditions, based on the construction of a labeled transition system. 
\end{itemize}

\textcolor{red}{todo :adapt section related work in \autoref{chap:nwf_type_graph} and \autoref{chap:subgraph_counting}}

% \paragraph{Position of our work}
% \begin{description}
%     \item[general introduction] We extend the type graph method to non-well-founded semirings and develop a termination technique for injective DPO graph rewriting systems on edge-labeled directed multigraphs using subgraph counting.
%     \item[our general philosphy] Different from the philosphy of Overbeek and Endrullis, we focus on 
% developping specific termination techniques for DPO graph rewriting systems on
% edge-labeled directed multigraphs, which are graphs with directed edges that can have multiple edges between two nodes and edges can have labels. These techniques cannot be applied to other categories in general, for some of them, they can be applied with some minor modifications. This choice is motivated by the fact that developping termination techniques is a difficult task, we want more arguments specific to DPO graph rewriting systems on edge-labeled directed multigraphs be available in the process of developing termination techniques. Another reason is that we want our techniques easy to be understood: to understand our techniques, only undergraduate knowledge in graph theory and some basic knowledge in category theory are required. The second reason is important because in our opinion, automated tools are not only for experts, and non experts should be able to easily understand the techniques and are not supposed to be experts in category theory. 
% \end{description}

\paragraph{Contribution 1 : Termination of Graph Rewriting using Weighted Type Graphs over Non-well-founded Semirings} 
While the type graph method is a powerful termination technique for DPO graph rewriting systems, previous approaches use weighted type graphs over well-founded semirings, which have practical limitations when applied to edge-labeled directed multigraph rewriting. In \autoref{chap:nwf_type_graph}, we investigate the use of non-well-founded semirings to overcome these limitations. This is a step towards a more practical termination tool.
We have implemented both our method and the prior approach by Endrullis and Overbeek~\cite{overbeek2024termination_lmcs} for edge-labeled graph rewriting within a unified tool.

While the weighted type graph method is a powerful technique for proving termination of double-pushout (DPO) graph rewriting systems. 
Existing approaches use weighted type graphs over well-founded semirings, but these have practical limitations when applied to edge-labeled directed multigraph rewriting.
We investigate the use of non-well-founded semirings to overcome these limitations.
We have implemented both our method and a prior approach for edge-labeled graph rewriting within a unified tool.

\autoref{chap:nwf_type_graph} is based on the following paper:
\begin{itemize}
    \item Qi Qiu. Termination of Graph Rewriting using Weighted Type Graphs over Non-well-founded Semirings. 16th International Workshop on Graph Computation Models, Jun 2025, Koblenz, Germany. 2025. ⟨hal-04954960v3⟩
\end{itemize}
 
\paragraph{Contribution 2 : Termination of Injective DPO Graph Rewriting Systems using Subgraph Counting} 

\autoref{chap:subgraph_counting} presents a machine-checkable sufficient condition for relative termination of double-pushout (DPO) graph rewriting systems with injective rules on edge-labeled multigraphs. 
the method defines a graph's weight as the sum of weights of occurrences of a set of graphs within it. By ensuring 
(1) every rewriting step using rules in a set A strictly decreases the host graph's weight, and 
(2) every rewriting step using rules in a set B never increases it, we guarantee that rules in the set A can be applied only finitely many times in any rewriting sequence with rules in the union of A and B.  
The method resolves termination cases that prior interpretation-based methods cannot. 
We also propose an implementation of our technique.

\autoref{chap:subgraph_counting} is based on the following paper:
\begin{itemize}
    \item Qiu, Q. (2025). Termination of Injective DPO Graph Rewriting Systems Using Subgraph Counting. In: Endrullis, J., Tichy, M. (eds) Graph Transformation. ICGT 2025. Lecture Notes in Computer Science, vol 15720. Springer, Cham. \url{https://doi.org/10.1007/978-3-031-94706-3_1}
\end{itemize}


%%%%%%%%%%%%%%%%%%%%%%%%%%%%%%%%%%%%%%%%%%%%%%%%%%
%%%   Chapter : Preliminaries
%%%   What basic concepts do you need to understand my work in detail?
%%%%%%%%%%%%%%%%%%%%%%%%%%%%%%%%%%%%%%%%%%%%%%%%%%
\chapter{Preliminaries}
\label{chap:preliminaries}
In this chapter we recall definitions related to the double-pushout approach to edge-labeled multigraph rewriting and to relative termination. 
% The concepts and notation in~\autoref{Preliminaries:Graphs} and~\autoref{Preliminaries:Pushout} follow the treatments in~\cite{pierce1991basic,barr1990category}; the presentation in~\autoref{sec:dpo} is based on~\cite{endrullis2024generalized_icgt}.
\section{Directed Edge-labeled Multigraph} 
\label{Preliminaries:Graphs}
 \begin{definition}[\cite{barr1990category}]
    \label{def:graph:unlabeled}
    An \textbf{unlabeled graph} \( G \) consists of a collection of \textbf{nodes} (also called \textbf{objects}) and a collection of \textbf{edges}, each equipped with a \textbf{source} (or \textbf{domain}) node and a \textbf{target} (or \textbf{codomain}) node. 
    For an unlabeled graph \( G \), we denote by \( G_0 \) its collection of nodes, \( G_1 \) its collection of edges, \( \operatorname{dom}:G_1{\to}G_0 \) the domain function, and \( \operatorname{cod}:G_1{\to}G_0 \) the codomain function. An unlabeled graph is \textbf{finite} if \( G_0 \) and \( G_1 \) are both finite sets.
    We write \( a: s \to t \) to indicate that \( a \) is a directed edge from \( s \) to \( t \).
\end{definition}   

\begin{example}
    In an unlabeled graph, edges are directed and will be visualized as arrows; between any two nodes, there can be multiple edges, and loops are allowed. The following is an example of an unlabeled graph:
     \begin{center}
        \begin{tikzpicture}
            \graphbox{}{0mm}{0mm}{32mm}{28mm}{-10mm}{-14mm}{
                \node[draw,circle] (1) at (0,0) {1};
                \node[draw,circle] (2) at (2,0) {2};
                \draw[->] (1) edge[loop above] (1) ;
                \draw[->] (1) edge[loop below] (1) ;
                \draw[->] (1) edge[bend left] (2)  ;
                \draw[->] (2) edge[bend left] (1)   ;
            }
        \end{tikzpicture}
    \end{center}
\end{example}
% A homomorphism of unlabeled graphs is a mapping between the nodes and edges of two graphs that preserves the graph structure.
\begin{definition}
    \label{def:unlabeled_graph:homomorphism}
    Let \( G \) and \( H \) be unlabeled graphs. A \textbf{homomorphism of unlabeled graphs} $h: G \to H$ is a pair of functions $h_V: V(G) \to V(H) $ and $h_E: E(G) \to E(H)$ such that for every edge \( a: s \to t \) in \( G \), the edge \( h_E(a) \) is from \( h_V(s) \) to \( h_V(t) \).
\end{definition}


\begin{definition}
    \label{def:graph:composition}
    Let $f: (G,\lambda) \to (G',\lambda')$ and $g: (G',\lambda') \to (G'',\lambda'')$ be two homomorphisms of unlabeled graphs. Their composition, denoted $g \circ f$, is defined as the pair of functions $g_V \circ f_V$ and $g_E \circ f_E$.
\end{definition}
\begin{definition} 
    \label{def:graph}
     Let \(\Sigma\) be a finite set of labels. A \textbf{directed edge-labeled multigraph} is a structure \((V, E, \opn{dom}, \opn{cod},\lambda)\) where
    \begin{itemize}
        \item $V$ is a collection of distinct elements called \textbf{nodes} (or \textbf{objects}),
        \item $E$ is a collection of distinct elements, disjoint from $V$, called \textbf{edges} (or \textbf{arrows}),
        \item $\opn{dom}: E \to V$ is the \textbf{domain} (or \textbf{source}) function assigning to each edge its source node,
        \item $\opn{cod}: E \to V$ is the \textbf{codomain} (or \textbf{target}) function assigning to each edge its target node,
        \item $\lambda: E \to \Sigma$ is the \textbf{labeling} function assigning to each edge a label from $\Sigma$.
    \end{itemize}
    A directed edge-labeled multigraph is said to be \textbf{finite} if $V$ and $E$ are finite, and \textbf{discrete} if its edge set \(E\) is empty.
    For a directed edge-labeled multigraph \( G \), we write \( V(G) \) for its set of nodes, \( E(G) \) for its set of edges, \( \opn{dom}_G \) for its domain function, \( \opn{cod}_G \) for its codomain function, and \( \lambda_G \) for its labeling function. $a : s\overset{l}{\rightarrow} t$ denotes an arrow $a$ labeled by $l$ from $s$ to $t$.
\end{definition}
A directed edge-labeled multigraph may have multiple edges from one node to another, as well as loops (edges from a node to itself), because different edges $u,v$ can have the same source and target nodes, i.e. $\opn{dom}(u) = \opn{dom}(v)$ and $\opn{cod}(u) = \opn{cod}(v)$.

    Throughout, a directed edge-labeled multigraph will be simply referred to as a \enquote{labeled graph} or \enquote{graph} when the context makes it clear, and we fix a set $\Sigma$ of labels.

% \begin{definition} 
%     \label{def:graph}
%     Let \(\Sigma\) be a finite set of labels. A \textbf{directed edge-labeled multigraph} is an ordered pair \((G,\lambda)\) where \( G \) is an unlabeled graph and \( \lambda : E(G) \rightarrow \Sigma\) is an edge-labeling function. 
%     It is called \textbf{finite} if its underlying unlabeled graph is finite.
%     Throughout, a directed edge-labeled multigraph will be simply referred to as a \enquote{labeled graph} when the context makes it clear.
%     An arrow $a$ labeled by $l$ from $s$ to $t$ is denoted by $a : s\overset{l}{\rightarrow} t$.
% \end{definition}
%  The definition of labeled graphs extends the definition of unlabeled graphs, as unlabeled graphs can be seen as labeled graphs with a single label for all edges, i.e., the label function is a constant function.


\begin{example}
    Consider the graph shown in Figure~\ref{fig:preliminaries:labeled_graph}.
    It has two nodes which are marked with numbers to facilitate discussion, but in general, nodes are not labeled.
    There are five edges: two loops on node $1$ labeled \(a\) and \(b\), two edges from node $1$ to node $2$ both labeled \(a\), and one edge from node $2$ to node $1$ labeled \(a\). Note that the two edges from node $1$ to node $2$ are distinct edges, even though they share the same source, target, and label.

    \begin{figure}[H]
       \centering
       \resizebox{0.4\textwidth}{!}{
        \begin{tikzpicture}
            \graphbox{}{0mm}{0mm}{32mm}{28mm}{-10mm}{-14mm}{
                \node[draw,circle] (1) at (0,0) {1};
                \node[draw,circle] (2) at (2,0) {2};
                \draw[->] (1) edge[loop above] node[midway, above] {$a$} (1) ;
                \draw[->] (1) edge[loop below] node[midway, below] {$b$} (1) ;
                \draw[->] (1) edge[bend left] node[midway, above] {$a$}  (2)  ;
                \draw[->] (1) edge node[midway, above] {$a$}  (2)  ;
                \draw[->] (2) edge[bend left] node[midway, below] {$a$} (1)   ;
            }
        \end{tikzpicture}
       }
        \caption{Labeled graph.}
        \label{fig:preliminaries:labeled_graph}
    \end{figure}

\end{example}
  A homomorphism of labeled graphs is a homomorphism of unlabeled graphs that preserves the labels assigned to the edges. 
% \begin{definition}
%     \label{def:graph:homomorphism}
%     Let \( (G,\lambda) \) and \( (H,\lambda') \) be labeled graphs. A \textbf{homomorphism of labeled graphs} $h:(G,\lambda) \mathop{\rightarrow} (H,\lambda')$ is a homomorphism of unlabeled graphs such that for each edge \( a \) in \( G \), we have \( \lambda (a) \mathop{=} \lambda' (h_E (a)) \).
% \end{definition}

% \begin{definition}
%     \label{def:unlabeled_graph:homomorphism}
%     Let \( G \) and \( H \) be unlabeled graphs. A \textbf{homomorphism of unlabeled graphs} $h: G \mathop{\to} H$ 
% \end{definition}
\begin{definition}
    \label{def:graph:homomorphism}
    Let \( G \) and \( H \) be labeled graphs. A \textbf{homomorphism of labeled graphs}\index{Homomorphism of labeled graphs} $h:G \mathop{\rightarrow} H$ is a pair of functions $(h_V: V(G) \mathop{\to} V(H), h_E: E(G) \mathop{\to} E(H))$ such that for every edge \( a: s \mathop{\to} t \) in \( G \), the edge \( h_E(a) \) in $H$ is from node \( h_V(s) \) to node \( h_V(t) \), and \( \lambda_H(h_E(a)) \mathop{=} \lambda_G(a) \).
\end{definition}

\begin{notation}
    \label{notation:graph_homomorphism}
    We use the notation from Overbeek and Endrullis~\cite[Notation 1]{overbeek2023apbpotutorial} to visualize edge-labeled graph homomorphisms. Labeled graphs are enclosed in boxes with their names displayed in the top-left corner. Nodes and edges are assigned subsets of \(\mathbb{N}\) as identifiers, and these identifiers are chosen such that: (i) Each node \( y \) (resp. edge \( y \)) in the codomain graph is assigned the union of the identifiers of all nodes (resp. edges) in the domain graph that are mapped to \( y \); (ii) The graph homomorphism is uniquely determined by this assignment. To improve readability, we represent sets by listing their elements. Additionally, we omit identifiers when doing so does not cause confusion. An example of a homomorphism of labeled graphs is shown below.
    % in Figure~\ref{fig:preliminaries:graph_homomorphism}.
    In this example, the sets \(\{1\}\), \(\{2\}\), \(\{3\}\), \(\{4\}\), and \(\{3,4\}\) are represented as \(1\), \(2\), \(3\), \(4\), and \(3\ 4\), respectively. Edge identifiers are omitted.
    % \begin{figure}[H]
    %     \centering
    \begin{center}
        \resizebox{0.8\textwidth}{!}{
        \begin{tikzpicture}
            \graphbox{\( G \)}{00mm}{-20mm}{45mm}{24mm}{2mm}{-8mm}{
                \coordinate (o) at (-5mm,-8mm); 
                \node[draw,circle] (l1) at ($(o)+(-10mm,0mm)$) {1};
                \node[draw,circle] (l2) at ($(l1)+(3,0)$) {2};
                \node[draw,circle] (l3) at ($(l1)+(1,0)$) {3};
                \node[draw,circle] (l4) at ($(l1)+(2,0)$) {4};
                \draw[->] (l1) -- (l3) node[midway,above] {$a$};
                \draw[->] (l3) -- (l4) node[midway,above] {$b$};
                \draw[->] (l4) -- (l2) node[midway,above] {$a$};
            }  
            \graphbox{\( H \)}{50mm}{-20mm}{34mm}{24mm}{2mm}{-8mm}{
                \coordinate (o) at (0mm,-8mm); 
                \node[draw,circle] (l1) at ($(o)+(-10mm,0mm)$) {1};
                \node[draw,circle] (l2) at ($(l1)+(2,0)$) {2};
                \node[draw,circle] (l3) at ($(l1)+(1,0)$) {3\ 4};
                \draw[->] (l1) -- (l3) node[midway,above] {$a$};
                \draw[->] (l3) edge[loop above] node[midway,above] {$b$} (l3) ;
                \draw[->] (l3) -- (l2) node[midway,above] {$a$};
            }      
            \node () at (48mm,-29mm) {$\rightarrow$};
        \end{tikzpicture}
    }
    \end{center}
    % \caption{}
    % \label{fig:preliminaries:graph_homomorphism}
    % \end{figure}
\end{notation} 


\begin{definition}
    \label{def:graph:composition}
    Let $G, G'$ and $G''$ be labeled graphs, and $f: G \mathop{\to} G'$ and $g: G' \mathop{\to} G''$ homomorphisms of labeled graphs. Their composition, denoted by $g \circ f$, is defined as the component-wise composition of their node and edge mapping functions $(g_V \circ f_V, g_E \circ f_E)$.
\end{definition}

\begin{definition}
    \label{def:graph:unlabeled}
    An \textbf{unlabeled multigraph}\index{Unlabeled multigraph} \( G \) is a labeled graph such that the labeling function is constant, i.e., there is a label \( l \mathop{\in} \Sigma \) such that \( \lambda_G(e) \mathop{=} l \) for all edges \( e \) in $G$.
\end{definition}
    Hereafter, we omit the labeling function when referring to an unlabeled graph, and \( a: s \mathop{\to} t \) denotes an arrow \( a \) from \( s \) to \( t \). 
\section{Pushout and pullback}  
\label{Preliminaries:Pushout}
 The concepts and notation in this section follow the treatments in~\cite{pierce1991basic,barr1990category}.

\begin{definition}
    \label{def:cat}
    A \textbf{category} is an unlabeled graph \( C \) together with a total function \( u : V(C) \to E(C) \) and a partial function \( \star: E(C) \times E(C) \to E(C) \) such that 
        \begin{itemize}
            \item for all edges \( f:X \to Y \) and \( g:Y \to Z \), the edge \( f \star g :X \to Z \) is defined; 
            \item  for every node \( X \), \( u(X) \) is an edge from \( X \) to \( X \);
            \item for every \( f:X \to Y \), we have \(u(X) \star f = f = f \star u(Y)\);
            \item for all edges \( f \), \( g \) and \(h\), we have \( (f \star g) \star h = f \star (g \star h) \) whenever either side is defined.
        \end{itemize}
    Edges are called \textbf{morphisms}. The function $\star$ is called \textbf{composition}. For all \( X \in V(C) \), the edge \( u(X) \) is denoted by \( \operatorname{id}_X \) and is called the \textbf{identity} of the object \( X \).
    % \( C \) is called the \textbf{underlying graph} of the category \( \mathcal{C} \).
\end{definition} 
\begin{definition}
    A category \(\mathcal{C}\) is said to be \textbf{locally small} if for all objects \(X,Y\) in \(\mathcal{C}\), the collection $\opn{Hom}(X,Y)$ of morphisms from \(X\) to \(Y\) is a set (called a \textbf{hom-set}). For a locally small category, $\opn{Mono}(X,Y)$ denotes the set of all monomorphisms from $X$ to $Y$.
\end{definition}
 Throughout this section, we fix a locally small category \( \mathcal{C} \).
\begin{example}
    Consider the unlabeled graph shown in~\autoref{fig:preliminaries:category}. 
        \begin{figure}[H]
        \centering
        \begin{tikzpicture}
            \graphbox{}{0mm}{0mm}{32mm}{16mm}{-10mm}{-8mm}{
                \node[draw,circle] (1) at (0,0) {};
                \node[draw,circle] (2) at (2,0) {};
                \draw[->] (1) edge[loop above] node[midway, above] { } (1) ;
                \draw[->] (2) edge[loop above] node[midway, above] { } (2) ;
                \draw[->] (1) edge[bend left] node[midway, above] {}  (2)  ;
                \draw[->] (2) edge[bend left] node[midway, below] {} (1)   ;
            }
        \end{tikzpicture}
        \caption{}
        \label{fig:preliminaries:category}
    \end{figure}
    It can be considered as a category where the objects are the nodes and the morphisms are the paths between nodes; composition is path concatenation. The identity of a node is the self-loop of the node. There are at least three morphisms from the left node to itself: the identity morphism (the self-loop), the path that traverses the self-loop twice, and the path that goes to the right node and back.
\end{example}

% cat notation * 
\begin{notation}
    The composition of morphisms \( f : X \to Y \) and \( g : Y \to Z \) is written in diagrammatic order as \( f \star g \), rather than in functional order \( g \circ f \) as is common in litterature. The advantage is that, when reading from left to right, the morphisms appear in the same order as in the corresponding diagram, making the reasoning accompanying diagrams more intuitive. 
\end{notation}  

\begin{definition} 
    \label{def:cat:homo}
    A morphism \( f : X \to Y \) is said to be a \textbf{monomorphism} (is \textbf{monic}) if given any morphisms \( g,h: Z \to X  \), \( g \star f = h \star f \) implies \( g = h \). 
    In this case, we write $f : X \rightarrowtail Y$ to indicate that $f$ is a monomorphism.
\end{definition} 

\begin{example}
The category \(\mathbf{Set}\) has sets as objects and total functions between them as morphisms. For \(f\colon A\to B\) and \(g\colon B\to C\), composition is given by \(g\circ f\), and the identity morphism on a set \(A\) is the identity function \(\mathrm{id}_A\).
\end{example}

\begin{example} 
    Finite labeled graphs and their homomorphisms form a category, hereafter denoted by \textbf{Graph}. Its objects are labeled graphs, its morphisms are graph homomorphisms, and the monomorphisms are homomorphisms. 
    $\textbf{Graph}$ is locally small. 
\end{example}

% \begin{definition}[Span \cite{lowe2010graph}]
%     A pair \( (\alpha : A \to B,~\beta : A \to C) \) of morphisms with a common domain is called a \textbf{span}, denoted by \( B \overset{\alpha}{\leftarrow} A \overset{\beta}{\rightarrow} C \).
% \end{definition}
 
% \begin{definition}[Cospan]
%     A pair \( (\beta' : B \to D,~\alpha' : C \to D) \) of morphisms with a common codomain is called a \textbf{cospan}, denoted by \( B \overset{\beta'}{\rightarrow} D \overset{\alpha'}{\leftarrow} C \). 
% \end{definition} 
A span (resp. cospan) is a couple of morphisms with a common domain (resp. codomain).
\begin{definition}
An ordered pair \((\alpha : A \to B,\, \beta : A \to C)\) of morphisms with a common domain is called a \textbf{span} \cite{lowe2010graph}, denoted by
\(
B \overset{\alpha}{\leftarrow} A \overset{\beta}{\rightarrow} C
\). 
% An example of a span $(\alpha, \beta)$ is shown below.
Likewise, an ordered pair \((\beta' : B \to D,\, \alpha' : C \to D)\) of morphisms with a common codomain is called a \textbf{cospan}, denoted by
\(
B \overset{\beta'}{\rightarrow} D \overset{\alpha'}{\leftarrow} C
\). 
\end{definition}
\begin{example}
Consider the diagram in the category \textbf{Graph} shown in~\autoref{fig:preliminaries:span_cospan}. $(\alpha, \beta)$ is a span, and $(\beta', \alpha')$ is a cospan.
\end{example}


\begin{figure}[H]
   \centering
        \resizebox{0.7\textwidth}{!}{
        \begin{tikzpicture} 
            \graphbox{\( L \)}{40mm}{20mm}{34mm}{12mm}{2mm}{2mm}{
                \coordinate (o) at (0mm,-8mm); 
                \node[draw,circle] (l1) at ($(o)+(-10mm,0mm)$) {1};
                \node[draw,circle] (l2) at ($(l1)+(2,0)$) {2};
                \node[draw,circle,red] (l3) at ($(l1) + (1,0)$) {3};
                \draw[->,red] (l1) -- (l3) node[midway,above] {$a$};
                \draw[->,red] (l3) -- (l2) node[midway,above] {$a$};
            } 
    
            \graphbox{\( K \)}{0mm}{0mm}{34mm}{12mm}{2mm}{2mm}{
                \coordinate (o) at (0mm,-8mm); 
                \node[draw,circle] (l1) at ($(o)+(-10mm,0mm)$) {1};
                \node[draw,circle] (l2) at ($(l1)+(2,0)$) {2};
            }  
            \graphbox{\(G = L \cup C \)}{90mm}{5mm}{34mm}{27mm}{2mm}{-3mm}{
                \coordinate (o) at (0mm,-8mm); 
                \node[draw,circle] (l1) at ($(o)+(-10mm,0mm)$) {1};
                \node[draw,circle] (l2) at ($(l1)+(2,0)$) {2};
                \node[draw,circle,red] (l3) at ($(l1) + (1,0)$) {3};
                \node[draw,circle,blue] (l4) at ($(l2) + (0,-1)$) {6};
                \draw[->,red] (l1) -- (l3) node[midway,above] {$a$};
                \draw[->,red] (l3) -- (l2) node[midway,above] {$a$};
                \draw[->,blue] (l2) -- (l4) node[midway,right] {$a$};
                \node[draw,circle,blue] (l6) at ($(l1) + (0,-1)$) {7};
                \draw[<-,blue] (l1) -- (l6) node[midway,left] {$a$};
                \draw[->,blue] (l2) edge[out=-135,in=-45]node[midway,below] {$a$} (l1) ;
            }   
     
            \graphbox{\( C \)}{40mm}{-20mm}{34mm}{27mm}{2mm}{-3mm}{
                \coordinate (o) at (0mm,-8mm); 
                \node[draw,circle] (l1) at ($(o)+(-10mm,0mm)$) {1};
                \node[draw,circle] (l2) at ($(l1)+(2,0)$) {2};
                \node[draw,circle,blue] (l4) at ($(l2) + (0,-1)$) {6};
                \draw[->,blue] (l2) -- (l4) node[midway,right] {$a$};
                \draw[->,blue] (l2) edge[out=-135,in=-45]node[midway,below] {$a$} (l1) ;
                \node[ draw,circle,blue] (l6) at ($(l1) + (0,-1)$) {7};
                \draw[<-,blue] (l1) -- (l6) node[midway,left] {$a$};
            }      
            % K to L
            \draw[>->] (17mm,5mm) -- node[above] {$\alpha$} (37mm,15mm);
            % C to G
            \draw[>->] (76mm,-28mm)-- node[below] {$\alpha'$} (104mm,-24mm) ;
            % K to C
            \draw[>->] (17mm,-17mm) -- node[below] {$\beta$} (37mm,-28mm);
            % L to G
            \draw[>->] (76mm,16mm) -- node[above] {$\beta'$} (104mm,7mm);
            % \node () at (57mm,-6mm) {$PO$};
        \end{tikzpicture}
        }
        \caption{}
        \label{fig:preliminaries:span_cospan}
\end{figure}

\begin{definition}[\cite{barr1990category}]
    \label{def:cat:diagram}
    Let $\mathcal{C}$ be a category, and \( G \) an unlabeled graph. A \textbf{diagram} (of shape \( G \)) is a homomorphism of unlabeled graphs \( h : G \to \mathcal{C} \) where \( \mathcal{C} \) is considered as an unlabeled graph. A diagram is said to be \textbf{commutative} if, for all nodes \( u \), \( v \), and any two paths from \( u \) to \( v \) in the unlabeled graph \( G \):

    \begin{center}
    \resizebox{12cm}{!}{
        \begin{tikzpicture}
        \node (u) at (0,0) {\( u \)};
        \node (k1) at (2,0.5) {\( k_1 \)};
        \node (k2) at (4,0.5) {\( k_2 \)};
        \node (ketc) at (6,0.5) {\( \dots \)};
        \node (knm2) at (8,0.5) {\( k_{n-2} \)};
        \node (knm1) at (10,0.5) {\( k_{n-1} \)};
        \node (v) at (12,0) {\( v \)};
        \node (l1) at (2,-0.5) {\( l_1 \)};
        \node (l2) at (4,-0.5) {\( l_2 \)};
        \node (letc) at (6,-0.5) {\( \dots \)};
        \node (lnm2) at (8,-0.5) {\( l_{m-2} \)};
        \node (lnm1) at (10,-0.5) {\( l_{m-1} \)};
        \draw[->] (u) -- (k1) node [midway,above] {\( s_1 \)};
        \draw[->] (k1) -- (k2) node [midway,above] {\( s_2 \)};
        \draw[->] (k2) -- (ketc);
        \draw[->] (ketc) -- (knm2); 
        \draw[->] (knm2) -- (knm1) node[midway,above] {\( s_{n-2} \)}; 
        \draw[->] (knm1) -- (v) node[midway,above] {\( s_{n-1} \)}; 
        \draw[->] (u) -- (l1) node[midway,below] {\( t_1 \)};
        \draw[->] (l1) -- (l2) node[midway,below] {\( t_2 \)};
        \draw[->] (l2) -- (letc);
        \draw[->] (letc) -- (lnm2); 
        \draw[->] (lnm2) -- (lnm1) node[midway,below] {\( t_{m-2} \)}; 
        \draw[->] (lnm1) -- (v) node[midway,below] {\( t_{m-1} \)}; 
        \end{tikzpicture}
    }
    \end{center}
    \noindent
    the equality \( h(s_1) \star h(s_2) \star \dots  \star h(s_{n-1}) = h(t_1) \star h(t_2) \star \dots  \star h(t_{m-1}) \) holds.
\end{definition}

\begin{example}
    A commutative diagram in the category \textbf{Graph} of finite directed edge-labeled multigraphs is illustrated in~\autoref{fig:preliminaries:commutative_diagram}. The symbol $\mathop{=}$ in the center of the diagram
    is used to indicate that the two paths from node 1 to node 2 are equal.
 
    \begin{figure}[H]
        \centering
        \resizebox{0.7\textwidth}{!}{
        \begin{tikzpicture} 
            \graphbox{\( L \)}{40mm}{20mm}{34mm}{12mm}{2mm}{2mm}{
                \coordinate (o) at (0mm,-8mm); 
                \node[draw,circle] (l1) at ($(o)+(-10mm,0mm)$) {1};
                \node[draw,circle] (l2) at ($(l1)+(2,0)$) {2};
                \node[draw,circle,red] (l3) at ($(l1) + (1,0)$) {3};
                \draw[->,red] (l1) -- (l3) node[midway,above] {$a$};
                \draw[->,red] (l3) -- (l2) node[midway,above] {$a$};
            } 
    
            \graphbox{\( K \)}{0mm}{0mm}{34mm}{12mm}{2mm}{2mm}{
                \coordinate (o) at (0mm,-8mm); 
                \node[draw,circle] (l1) at ($(o)+(-10mm,0mm)$) {1};
                \node[draw,circle] (l2) at ($(l1)+(2,0)$) {2};
            }  
            \graphbox{\(G = L \cup C \)}{90mm}{5mm}{34mm}{27mm}{2mm}{-3mm}{
                \coordinate (o) at (0mm,-8mm); 
                \node[draw,circle] (l1) at ($(o)+(-10mm,0mm)$) {1};
                \node[draw,circle] (l2) at ($(l1)+(2,0)$) {2};
                \node[draw,circle,red] (l3) at ($(l1) + (1,0)$) {3};
                \node[draw,circle,blue] (l4) at ($(l2) + (0,-1)$) {6};
                \draw[->,red] (l1) -- (l3) node[midway,above] {$a$};
                \draw[->,red] (l3) -- (l2) node[midway,above] {$a$};
                \draw[->,blue] (l2) -- (l4) node[midway,right] {$a$};
                \node[draw,circle,blue] (l6) at ($(l1) + (0,-1)$) {7};
                \draw[<-,blue] (l1) -- (l6) node[midway,left] {$a$};
                \draw[->,blue] (l2) edge[out=-135,in=-45]node[midway,below] {$a$} (l1) ;
            }   
     
            \graphbox{\( C \)}{40mm}{-20mm}{34mm}{27mm}{2mm}{-3mm}{
                \coordinate (o) at (0mm,-8mm); 
                \node[draw,circle] (l1) at ($(o)+(-10mm,0mm)$) {1};
                \node[draw,circle] (l2) at ($(l1)+(2,0)$) {2};
                \node[draw,circle,blue] (l4) at ($(l2) + (0,-1)$) {6};
                \draw[->,blue] (l2) -- (l4) node[midway,right] {$a$};
                \draw[->,blue] (l2) edge[out=-135,in=-45]node[midway,below] {$a$} (l1) ;
                \node[ draw,circle,blue] (l6) at ($(l1) + (0,-1)$) {7};
                \draw[<-,blue] (l1) -- (l6) node[midway,left] {$a$};
            }      
            % K to L
            \draw[>->] (17mm,5mm) -- node[above] {$\alpha$} (37mm,15mm);
            % C to G
            \draw[>->] (76mm,-28mm)-- node[below] {$\alpha'$} (104mm,-24mm) ;
            % K to C
            \draw[>->] (17mm,-17mm) -- node[below] {$\beta$} (37mm,-28mm);
            % L to G
            \draw[>->] (76mm,16mm) -- node[above] {$\beta'$} (104mm,7mm);
            \node () at (57mm,-6mm) {$\mathop{=}$};
        \end{tikzpicture}
        }
        \caption{}
        \label{fig:preliminaries:commutative_diagram}
    \end{figure}
\end{example}

\begin{notation}
    When the context makes it clear, \( h_{AB} \) denotes a morphism \( h : A \to B \), and we refer to diagrams by listing their nodes, as is standard in geometry. 
    For example, we denote the diagram shown in \autoref{def:cat:po} by \( ACDB \) or \( ABDC \).
\end{notation}   

The pushout is a construction in category theory that can often be thought of as the construction of a new structure from two given structures by gluing them along a common interface structure.
\begin{definition}
    \label{def:cat:po} 
    A \textbf{pushout} of a span \( B \overset{\alpha}{\leftarrow} A \overset{\beta}{\rightarrow} C \), as shown in~\autoref{fig:preliminaries:pushout_sdfkjasdlgjfl}
    is defined as a cospan \( B \overset{\beta'}{\rightarrow} D \overset{\alpha'}{\leftarrow} C \) such that \( \alpha \star \beta' = \beta \star \alpha' \), and for every cospan \( B \overset{\gamma'}{\rightarrow} E \overset{\gamma}{\leftarrow} C \), if \( \alpha \star \gamma' = \beta \star \gamma \) holds, then there exists a unique morphism \(\delta : D \to E\) such that \( \gamma' = \beta' \star \delta \) and \( \gamma = \alpha' \star \delta \).
    \begin{figure}[H]
        \centering
        \resizebox{0.4\textwidth}{!}{
            \begin{tikzpicture}
                    \node (i) at (0,0) {A};
                    \node (r) at (1,1) {B};
                    \node (c) at (1,-1) {C};
                    \node (h) at (2,0) {D};
                    % \node () at (1,-1) {\( \Delta \)};
                    \draw[->]  (i) -- (r) node [midway,left] {$ \alpha $};
                    \draw[->] (c) -- (h) node [midway,left] {$ \alpha' $};
                    \draw[->] (r) -- (h) node[midway, left] {$ \beta' $};
                    \draw[->] (i) -- (c) node[midway, left] {$ \beta $};
                    \node (d') at (4,0) {E};
                    \draw[->] (c) -- (d') node [midway,below]{$ \gamma $};
                    \draw[->] (r) -- (d') node [midway,above]{$ \gamma' $};
                    \draw[->,dashed] (h) -- (d') node [midway]{$ \delta $};
                \end{tikzpicture}
        }
        \caption{}
        \label{fig:preliminaries:pushout_sdfkjasdlgjfl}
    \end{figure}
The diagram involving \( (\alpha, \beta, \alpha', \beta') \) is called the \textbf{pushout square}, with \(D\) as the \textbf{pushout object}. The existence of a unique morphism is known as the \textbf{universal mapping property of the pushout}.
\end{definition} 

\begin{example}
    \label{ex:cat:posfjsdlkgja}
     Pushouts of two morphisms always exist in \(\mathbf{Set}\), and (up to isomorphism) can be described as follows. Let
    \( B \overset{\alpha}{\leftarrow} A \overset{\beta}{\rightarrow} C \) be a span. Its pushout is the cospan \( B \overset{\beta'}{\rightarrow} D \overset{\alpha'}{\leftarrow} C \) where the pushout object of \((\alpha,\beta)\) is the set
    \[
    D \;=\; (B + C)/{\sim}
    \]
    where \(B + C\) denotes the disjoint union of $B$ and $C$ and \(\sim\) is the smallest equivalence relation that includes \(\set{(\alpha(a),\beta(b))\mid a \in A }\). The maps
    \(\beta' \colon B\to D\) and \(\alpha' \colon C\to D\) send each element to its equivalence class.

    %  Consider the functions \(b\) and \(c\) shown in~\autoref{fig:preliminaries:a_rewriting_step_dfjalsdkjflg} in the category \(\mathbf{Set}\). Each set is drawn as a box; elements of a set are represented by circles. The numbers inside circles indicate how the functions map those elements. 
    %  The element labeled by $n$ in set $X$ is denoted by $n_X$ to avoid ambiguity.
    %  The equivalence relation \(\sim\) is the set \(\set{(1_B,1_C),(2_B,2_C)}\). The disjoint union of \(B\) and \(C\) is $D = \set{1_B,2_B,4_B,1_C,2_C,3_C}$. The quotient set $D/{\sim}$ is $\set{[1_B],[2_B],[3_C],[4_B]}$, where $[x]$ denotes the equivalence class of $x$. Note that $[1_C] = [1_B]$ and $[2_C] = [2_B]$, and $\set{[1_B],[2_B],[3_C],[4_B]}$ is isomorphic to $\set{1_D,2_D,3_D,4_D}$.
    For example, consider the functions \(\alpha\) and \(\beta\) shown in Figure~\ref{fig:preliminaries:a_rewriting_step_dfjalsdkjflg} in the category \(\mathbf{Set}\) where each set is drawn as a box and its elements are represented by circles. The numbers inside circles indicate how the functions map those elements.
    \begin{figure}[H]
      \centering 
      \resizebox{0.7\textwidth}{!}{
      \begin{tikzpicture}
          \graphbox{\( A\)}{40mm}{-3mm}{34mm}{12mm}{2mm}{2mm}{
              \coordinate (o) at (0mm,-8mm); 
              \node[draw,circle] (l1) at ($(o)+(-10mm,0mm)$) {1};
              \node[draw,circle] (l2) at ($(l1)+(2,0)$) {2};
          }  
          \graphbox{\( B \)}{80mm}{-3mm}{45mm}{12mm}{2mm}{2mm}{
              \coordinate (o) at (-5mm,-8mm); 
              \node[draw,circle] (l1) at ($(o)+(-10mm,0mm)$) {1};
              \node[draw,circle] (l2) at ($(l1)+(3,0)$) {2};
              \node[draw,circle] (l3) at ($(l1) + (1,0)$) {4};
            %   \node[draw,circle] (l4) at ($(l1) + (2,0)$) {5};
            %   \draw[ ] (l1) -- (l3) node[midway,above] {$a$};
            %   \draw[ ] (l3) -- (l4) node[midway,above] {$b$};
            %   \draw[ ] (l4) -- (l2) node[midway,above] {$a$};
          }     
          \graphbox{\( C  \)}{40mm}{-22mm}{34mm}{22mm}{2mm}{-3mm}{
              \coordinate (o) at (0mm,-3mm); 
              \node[draw,circle] (l1) at ($(o)+(-10mm,0mm)$) {1};
              \node[draw,circle] (l2) at ($(l1)+(2,0)$) {2};
            %   \node[draw,circle] (l4) at ($(l2) + (0,-1)$) {6};
              \node[ draw,circle] (l6) at ($(l1) + (0,-1)$) {3};
            %   \draw[ ] (l1) -- (l6) node[midway,left] {$a$};
            %   \draw[ ] (l2) -- (l4) node[midway,right] {$a$};
          }    
          \graphbox{\( D \)}{80mm}{-22mm}{45mm}{22mm}{2mm}{-3mm}{
              \coordinate (o) at (-5mm,-3mm); 
              \node[draw,circle] (l1) at ($(o)+(-10mm,0mm)$) {1};
              \node[draw,circle] (l2) at ($(l1)+(3,0)$) {2};
              \node[draw,circle] (l3) at ($(l1) + (1,0)$) {4};
            %   \node[draw,circle] (l4) at ($(l1) + (2,0)$) {5};
            %   \node[ draw,circle] (l5) at ($(l2) + (0,-1)$) {6};
              \node[ draw,circle] (l6) at ($(l1) + (0,-1)$) {3};
            %   \draw[ ] (l1) -- (l6) node[midway,left] {$a$};
            %   \draw[] (l1) -- (l3) node[midway,above] {$a$};
            %   \draw[] (l3) -- (l4) node[midway,above] {$b$};
            %   \draw[ ] (l4) -- (l2) node[midway,above] {$a$};
            %   \draw[ ] (l2) -- (l5) node[midway,right] {$a$};
          }    
          \node () at (77mm,-8mm) {\( \overset{\alpha}{\rightarrowtail} \)}; % K -> R
          \node () at (58mm,-18mm) {\( \beta\downarrowtail \)};
          \node () at (102mm,-18mm) {\( \beta'\downarrowtail \)};
          \node () at (77mm,-33mm) {\( \overset{\alpha'}{\rightarrowtail} \)}; % C -> H
      \end{tikzpicture}
      }
      \caption{}
      \label{fig:preliminaries:a_rewriting_step_dfjalsdkjflg}
  \end{figure}
    Throughout this example, an element labeled \(n\) in a set \(X\) is denoted by \(n_X\) to avoid ambiguity.        
    The binary relation \(\sim\) is the reflexive, symmetric and transitive closure of the binary relation $\{(1_B,1_C),(2_B,2_C)\}$.
 
    The disjoint union of \(B\) and \(C\) is
    \[ 
    D' = \{1_B,2_B,4_B,1_C,2_C,3_C\},
    \]
    and the quotient \(D'/\sim\) is
    \[
    D'/\sim = \{[1_B],[2_B],[3_C],[4_B]\},
    \]
    We have \([1_C]=[1_B]\) and \([2_C]=[2_B]\), and the maps
    \(\beta'' \colon B\to D'\) and \(\alpha'' \colon C\to D'\) send each element to its equivalence class. Note that \(\{[1_B],[2_B],[3_C],[4_B]\}\) is isomorphic to \(D\), shown in \autoref{fig:preliminaries:a_rewriting_step_dfjalsdkjflg}, which is expected because the pushout of a span is unique up to isomorphism.
\end{example}

\begin{proposition}{\cite{corradini1997algebraic}[p.188]}
    \label{prop:pushout_graph_always_exists}
    In category \textbf{Graph} the pushout of two arrows always exists: It can be computed componentwise (as a pushout in \textbf{Set}) for the nodes and for the edges, and the source, target, and labeling mappings are uniquely determined.
\end{proposition}

\begin{example}
    The diagram in~\autoref{fig:preliminaries:pushout_injective} is a pushout square in the category \textbf{Graph}. In this example, both $\alpha$ and $\beta$ are injective morphisms. Therefore, the pushout object $G$ can be constructed easily by taking the interface graph $K$ and adding elements from $L$ and $C$ which are not present in $K$.
    \begin{figure}[H]
        \centering
        \resizebox{0.7\textwidth}{!}{
        \begin{tikzpicture} 
            \graphbox{\( L \)}{40mm}{20mm}{34mm}{12mm}{2mm}{2mm}{
                \coordinate (o) at (0mm,-8mm); 
                \node[draw,circle] (l1) at ($(o)+(-10mm,0mm)$) {1};
                \node[draw,circle] (l2) at ($(l1)+(2,0)$) {2};
                \node[draw,circle,red] (l3) at ($(l1) + (1,0)$) {3};
                \draw[->,red] (l1) -- (l3) node[midway,above] {$a$};
                \draw[->,red] (l3) -- (l2) node[midway,above] {$a$};
            } 
    
            \graphbox{\( K \)}{0mm}{0mm}{34mm}{12mm}{2mm}{2mm}{
                \coordinate (o) at (0mm,-8mm); 
                \node[draw,circle] (l1) at ($(o)+(-10mm,0mm)$) {1};
                \node[draw,circle] (l2) at ($(l1)+(2,0)$) {2};
            }  
            \graphbox{\(G  \)}{90mm}{5mm}{34mm}{27mm}{2mm}{-3mm}{
                \coordinate (o) at (0mm,-8mm); 
                \node[draw,circle] (l1) at ($(o)+(-10mm,0mm)$) {1};
                \node[draw,circle] (l2) at ($(l1)+(2,0)$) {2};
                \node[draw,circle,red] (l3) at ($(l1) + (1,0)$) {3};
                \node[draw,circle,blue] (l4) at ($(l2) + (0,-1)$) {6};
                \draw[->,red] (l1) -- (l3) node[midway,above] {$a$};
                \draw[->,red] (l3) -- (l2) node[midway,above] {$a$};
                \draw[->,blue] (l2) -- (l4) node[midway,right] {$a$};
                \node[draw,circle,blue] (l6) at ($(l1) + (0,-1)$) {7};
                \draw[<-,blue] (l1) -- (l6) node[midway,left] {$a$};
                \draw[->,blue] (l2) edge[out=-135,in=-45]node[midway,below] {$a$} (l1) ;
            }   
     
            \graphbox{\( C \)}{40mm}{-20mm}{34mm}{27mm}{2mm}{-3mm}{
                \coordinate (o) at (0mm,-8mm); 
                \node[draw,circle] (l1) at ($(o)+(-10mm,0mm)$) {1};
                \node[draw,circle] (l2) at ($(l1)+(2,0)$) {2};
                \node[draw,circle,blue] (l4) at ($(l2) + (0,-1)$) {6};
                \draw[->,blue] (l2) -- (l4) node[midway,right] {$a$};
                \draw[->,blue] (l2) edge[out=-135,in=-45]node[midway,below] {$a$} (l1) ;
                \node[ draw,circle,blue] (l6) at ($(l1) + (0,-1)$) {7};
                \draw[<-,blue] (l1) -- (l6) node[midway,left] {$a$};
            }      
            % K to L
            \draw[>->] (17mm,5mm) -- node[above] {$\alpha$} (37mm,15mm);
            % C to G
            \draw[>->] (76mm,-28mm)-- node[below] {$\alpha'$} (104mm,-24mm) ;
            % K to C
            \draw[>->] (17mm,-17mm) -- node[below] {$\beta$} (37mm,-28mm);
            % L to G
            \draw[>->] (76mm,16mm) -- node[above] {$\beta'$} (104mm,7mm);
            \node () at (57mm,-6mm) {$PO$};
        \end{tikzpicture}
        }
        \caption{Pushout square with injective morphisms.}
        \label{fig:preliminaries:pushout_injective}
    \end{figure}
\end{example}

\begin{example}
    The diagram in~\autoref{fig:preliminaries:pushout_non_injective} is a pushout square in the category \textbf{Graph}. In this example, $\beta$ is not injective. Therefore, some elements are merged in the pushout object $G$.

    \begin{figure}[H]
        \centering 
        \resizebox{0.7\textwidth}{!}{
        \begin{tikzpicture} 
            \graphbox{\( L \)}{40mm}{20mm}{34mm}{20mm}{2mm}{2mm}{
                \coordinate (o) at (0mm,-8mm); 
                \node[draw,circle] (l1) at ($(o)+(-10mm,0mm)$) {1};
                \node[draw,circle] (l2) at ($(l1)+(2,0)$) {2};
                \draw[->,red] (l2) edge[out=-135,in=-45]node[midway,below] {$a$} (l1) ;
                \node[draw,circle,red] (l3) at ($(l1) + (1,0)$) {3};
                \draw[->,red] (l1) -- (l3) node[midway,above] {$a$};
                \draw[->,red] (l3) -- (l2) node[midway,above] {$a$};
            } 
    
            \graphbox{\( K \)}{0mm}{0mm}{34mm}{12mm}{2mm}{2mm}{
                \coordinate (o) at (0mm,-8mm); 
                \node[draw,circle] (l1) at ($(o)+(-10mm,0mm)$) {1};
                \node[draw,circle] (l2) at ($(l1)+(2,0)$) {2};
            }  
            \graphbox{\(G  \)}{110mm}{5mm}{34mm}{40mm}{2mm}{-3mm}{
                \coordinate (o) at (0mm,-20mm); 
                 \node[draw,circle] (l1) at ($(o)+(0,0)$) {1\ 2};
                % \node[draw,circle] (l1) at ($(o)+(-10mm,0mm)$) {1};
                % \node[draw,circle] (l2) at ($(l1)+(2,0)$) {2};
                \draw[->,red] (l1) edge[loop below] node[midway, below] {$a$} (l1) ;
                \node[draw,circle,red] (l3) at ($(l1) + (0,1.4)$) {3};
                \node[draw,circle,blue] (l4) at ($(l1) + (1,-1)$) {6};
                \draw[->,red] (l1) edge[bend left] node[midway,left] {$a$} (l3);
                \draw[->,red] (l3) edge[bend left] node[midway,right] {$a$} (l1);
                \draw[->,blue] (l1) edge  node[midway,right] {$a$} (l4);
                \node[draw,circle,blue] (l6) at ($(l1) + (-1,-1)$) {7};
                \draw[<-,blue] (l1) edge node[midway,left] {$a$} (l6) ;
                % \draw[->,blue] (l1) edge[out=-135,in=-45]node[midway,below] {$a$} (l1) ;
            }   
     
            \graphbox{\( C \)}{40mm}{-20mm}{34mm}{35mm}{2mm}{-3mm}{
                \coordinate (o) at (0mm,-14mm); 
                \node[draw,circle] (l1) at ($(o)+(0,0)$) {1\ 2};
                % \node[draw,circle] (l2) at ($(l1)+(2,0)$) {2};
                \node[draw,circle,blue] (l4) at ($(l1) + (1,-1)$) {6};
                \draw[->,blue] (l1) -- (l4) node[midway,right] {$a$};
                % \draw[->,blue] (l1) edge[loop above] node[midway, above] {$a$} (l1) ;
                \node[ draw,circle,blue] (l6) at ($(l1) + (-1,-1)$) {7};
                \draw[<-,blue] (l1) -- (l6) node[midway,left] {$a$};
            }
            % K to L  
            \draw[>->] (17mm,5mm) -- node[above] {$\alpha$} (37mm,15mm);
            % C to G
            \draw[>->] (76mm,-28mm)-- node[below] {$\alpha'$} (104mm,-24mm) ;
            % K to C
            \draw[->] (17mm,-17mm) -- node[below] {$\beta$} (37mm,-28mm);
            % L to G
            \draw[->] (76mm,16mm) -- node[above] {$\beta'$} (104mm,7mm);
            \node () at (57mm,-6mm) {$PO$};
        \end{tikzpicture}
        }
        \caption{Pushout square with a non-injective morphism.}
        \label{fig:preliminaries:pushout_non_injective}
    \end{figure} 
\end{example}
The \emph{pullback}, defined below, is the dual construction of the pushout in category theory, and it can be thought of as construction of the interface structure along which two structures are glued together. 
\begin{definition} 
    \label{def:cat:pb}
   A \textbf{pullback} of a cospan \(B \overset{\beta'}{\rightarrow} D \overset{\alpha'}{\leftarrow} C \), as shown in~\autoref{fig:preliminaries:pullback_ssdsfd},
is a span \( B \overset{\alpha}{\leftarrow} A \overset{\beta}{\rightarrow} C \) such that \( \alpha \star \beta' = \beta \star \alpha' \), and for every span \( B \overset{\gamma'}{\leftarrow} E \overset{\gamma}{\rightarrow} C \) if \(\gamma' \star \beta' = \gamma \star \alpha'\) holds, then there exists a unique morphism \(\delta: E \to A\) such that $\gamma' = \delta \star \alpha$ and $\gamma = \delta \star \beta$. 
    \begin{figure}[H]
        \centering  
        \resizebox{0.4\textwidth}{!}{
                    \begin{tikzpicture}
                        \node (i) at (0,0) {A};
                        \node (r) at (1,1) {B};
                        \node (c) at (1,-1) {C};
                        \node (h) at (2,0) {D}; 
                        % \node () at (1,-1) {\( \Delta \)};
                        \draw[->]  (i) -- (r) node [midway,left] {$\alpha$};
                        \draw[->] (c) -- (h) node [midway,left] {$\alpha'$};
                        \draw[->] (r) -- (h) node[midway, left] {$\beta'$};
                        \draw[->] (i) -- (c) node[midway, left] {$\beta$};
                        \node (d') at (-2,0) {E};
                        \draw[<-] (c) -- (d') node [midway,below]{$\gamma$};
                        \draw[<-] (r) -- (d') node [midway,above]{$\gamma'$};
                        \draw[->, dashed] (d') -- (i) node [midway]{$\delta$};
                    \end{tikzpicture}
        }
        \caption{}
        \label{fig:preliminaries:pullback_ssdsfd}
    \end{figure}
The diagram involving \( (\alpha, \beta, \alpha', \beta') \) is called the \textbf{pullback square}, with \(A\) as the \textbf{pullback object}. The existence of a unique morphism is known as the \textbf{universal mapping property of the pullback}.
\end{definition} 
In the category \textbf{Graph} pullbacks always exist and are computed pointwise: take the pullback in \textbf{Set} of the node sets and of the edge sets, and equip the resulting graph with source, target and labeling maps induced componentwise; the projection graph homomorphisms are the corresponding componentwise projections and satisfy the universal property.
\begin{example} 
    In the category \textbf{Set},
    the pullback 
    of a cospan \(B \overset{\beta'}{\rightarrow} D \overset{\alpha'}{\leftarrow} C \) is the span \( B \overset{\alpha}{\leftarrow} A \overset{\beta}{\rightarrow} C \) where
    the pullback object $A$ is the set $A = \set{(b,c) \mid \beta (c) = \alpha (b)} \subseteq B \times C$; $\alpha$ and $\beta$ are defined as the corresponding projections, e.g., $\alpha((b, c)) = b$ and $\beta((b, c)) = c$.

    Consider the cospan \(B \overset{\beta'}{\rightarrow} D \overset{\alpha'}{\leftarrow} C \) shown in~\autoref{fig:preliminaries:a_rewriting_step_dfjalsdkdfsdfjflg}. 
    The span \( B \overset{\alpha}{\leftarrow} A \overset{\beta}{\rightarrow} C \) shown in~\autoref{fig:preliminaries:a_rewriting_step_dfjalsdkdfsdfjflg} is the pullback of this cospan.
    Indeed, the pullback object can be taken as the set $A' = \set{(1_B,1_C),(2_B,2_C)} \subseteq B \times C$, and the morphisms $\alpha'' \colon A' \to B$ and $\beta'' \colon A' \to C$ are defined as the corresponding projections, e.g., $\alpha''((x,y)) = x$ and $\beta''((x,y)) = y$. Note that $A$ is isomorphic to $A'$, which is expected because the pullback of a cospan is unique up to isomorphism.
      \begin{figure}[H]
      \centering 
      \resizebox{0.7\textwidth}{!}{
      \begin{tikzpicture}
          \graphbox{\( A\)}{40mm}{-3mm}{34mm}{12mm}{2mm}{2mm}{
              \coordinate (o) at (0mm,-8mm); 
              \node[draw,circle] (l1) at ($(o)+(-10mm,0mm)$) {1};
              \node[draw,circle] (l2) at ($(l1)+(2,0)$) {2};
          }  
          \graphbox{\( B \)}{80mm}{-3mm}{45mm}{12mm}{2mm}{2mm}{
              \coordinate (o) at (-5mm,-8mm); 
              \node[draw,circle] (l1) at ($(o)+(-10mm,0mm)$) {1};
              \node[draw,circle] (l2) at ($(l1)+(3,0)$) {2};
              \node[draw,circle] (l3) at ($(l1) + (1,0)$) {4};
            %   \node[draw,circle] (l4) at ($(l1) + (2,0)$) {5};
            %   \draw[ ] (l1) -- (l3) node[midway,above] {$a$};
            %   \draw[ ] (l3) -- (l4) node[midway,above] {$b$};
            %   \draw[ ] (l4) -- (l2) node[midway,above] {$a$};
          }     
          \graphbox{\( C  \)}{40mm}{-22mm}{34mm}{22mm}{2mm}{-3mm}{
              \coordinate (o) at (0mm,-3mm); 
              \node[draw,circle] (l1) at ($(o)+(-10mm,0mm)$) {1};
              \node[draw,circle] (l2) at ($(l1)+(2,0)$) {2};
            %   \node[draw,circle] (l4) at ($(l2) + (0,-1)$) {6};
              \node[ draw,circle] (l6) at ($(l1) + (0,-1)$) {3};
            %   \draw[ ] (l1) -- (l6) node[midway,left] {$a$};
            %   \draw[ ] (l2) -- (l4) node[midway,right] {$a$};
          }    
          \graphbox{\( D \)}{80mm}{-22mm}{45mm}{22mm}{2mm}{-3mm}{
              \coordinate (o) at (-5mm,-3mm); 
              \node[draw,circle] (l1) at ($(o)+(-10mm,0mm)$) {1};
              \node[draw,circle] (l2) at ($(l1)+(3,0)$) {2};
              \node[draw,circle] (l3) at ($(l1) + (1,0)$) {4};
            %   \node[draw,circle] (l4) at ($(l1) + (2,0)$) {5};
            %   \node[ draw,circle] (l5) at ($(l2) + (0,-1)$) {6};
              \node[ draw,circle] (l6) at ($(l1) + (0,-1)$) {3};
            %   \draw[ ] (l1) -- (l6) node[midway,left] {$a$};
            %   \draw[] (l1) -- (l3) node[midway,above] {$a$};
            %   \draw[] (l3) -- (l4) node[midway,above] {$b$};
            %   \draw[ ] (l4) -- (l2) node[midway,above] {$a$};
            %   \draw[ ] (l2) -- (l5) node[midway,right] {$a$};
          }    
          \node () at (77mm,-8mm) {\( \overset{\alpha}{\rightarrowtail} \)};
          \node () at (58mm,-18mm) {\( \beta\downarrowtail \)};
          \node () at (102mm,-18mm) {\( \beta'\downarrowtail \)};
          \node () at (77mm,-33mm) {\( \overset{\alpha'}{\rightarrowtail} \)};
      \end{tikzpicture}
      }
      \caption{}
      \label{fig:preliminaries:a_rewriting_step_dfjalsdkdfsdfjflg}
  \end{figure}
\end{example}

\begin{example}
    Consider the diagram shown in~\autoref{fig:preliminaries:pullback} in the category \textbf{Graph}.
    The span $L \overset{\alpha}{\leftarrow} K \overset{\beta}{\rightarrow} C$ is a pullback of the cospan $L \overset{\beta'}{\rightarrow} G \overset{\alpha'}{\leftarrow} C$.
    \begin{figure}[H]
        \centering
        \resizebox{0.9\textwidth}{!}{
        \begin{tikzpicture} 
            \graphbox{\( L \)}{40mm}{20mm}{34mm}{20mm}{2mm}{2mm}{
                \coordinate (o) at (0mm,-8mm); 
                \node[draw,circle] (l1) at ($(o)+(-10mm,0mm)$) {1};
                \node[draw,circle] (l2) at ($(l1)+(2,0)$) {2};
                \draw[->,red] (l2) edge[out=-135,in=-45]node[midway,below] {$a$} (l1) ;
                \node[draw,circle,red] (l3) at ($(l1) + (1,0)$) {3};
                \draw[->,red] (l1) -- (l3) node[midway,above] {$a$};
                \draw[->,red] (l3) -- (l2) node[midway,above] {$a$};
            } 
    
            \graphbox{\( K \)}{0mm}{0mm}{34mm}{12mm}{2mm}{2mm}{
                \coordinate (o) at (0mm,-8mm); 
                \node[draw,circle] (l1) at ($(o)+(-10mm,0mm)$) {1};
                \node[draw,circle] (l2) at ($(l1)+(2,0)$) {2};
            }  
            \graphbox{\(G  \)}{110mm}{5mm}{34mm}{40mm}{2mm}{-3mm}{
                \coordinate (o) at (0mm,-20mm); 
                 \node[draw,circle] (l1) at ($(o)+(0,0)$) {1\ 2};
                % \node[draw,circle] (l1) at ($(o)+(-10mm,0mm)$) {1};
                % \node[draw,circle] (l2) at ($(l1)+(2,0)$) {2};
                \draw[->,red] (l1) edge[loop below] node[midway, below] {$a$} (l1) ;
                \node[draw,circle,red] (l3) at ($(l1) + (0,1.4)$) {3};
                \node[draw,circle,blue] (l4) at ($(l1) + (1,-1)$) {6};
                \draw[->,red] (l1) edge[bend left] node[midway,left] {$a$} (l3);
                \draw[->,red] (l3) edge[bend left] node[midway,right] {$a$} (l1);
                \draw[->,blue] (l1) edge  node[midway,right] {$a$} (l4);
                \node[draw,circle,blue] (l6) at ($(l1) + (-1,-1)$) {7};
                \draw[<-,blue] (l1) edge node[midway,left] {$a$} (l6) ;
                % \draw[->,blue] (l1) edge[out=-135,in=-45]node[midway,below] {$a$} (l1) ;
            }   
     
            \graphbox{\( C \)}{40mm}{-20mm}{34mm}{35mm}{2mm}{-3mm}{
                \coordinate (o) at (0mm,-14mm); 
                \node[draw,circle] (l1) at ($(o)+(0,0)$) {1\ 2};
                % \node[draw,circle] (l2) at ($(l1)+(2,0)$) {2};
                \node[draw,circle,blue] (l4) at ($(l1) + (1,-1)$) {6};
                \draw[->,blue] (l1) -- (l4) node[midway,right] {$a$};
                % \draw[->,blue] (l1) edge[loop above] node[midway, above] {$a$} (l1) ;
                \node[ draw,circle,blue] (l6) at ($(l1) + (-1,-1)$) {7};
                \draw[<-,blue] (l1) -- (l6) node[midway,left] {$a$};
            }
            % K to L  
            \draw[>->] (17mm,5mm) -- node[above] {$\alpha$} (37mm,15mm);
            % C to G
            \draw[>->] (76mm,-28mm)-- node[below] {$\alpha'$} (104mm,-24mm) ;
            % K to C
            \draw[->] (17mm,-17mm) -- node[below] {$\beta$} (37mm,-28mm);
            % L to G 
            \draw[->] (76mm,16mm) -- node[above] {$\beta'$} (104mm,7mm);
            \node () at (57mm,-6mm) {$PB$};
        \end{tikzpicture}
        }
        \caption{Pullback.}
        \label{fig:preliminaries:pullback}
    \end{figure} 

 
    % Specifically, for a possible span \( B \overset{\gamma'}{\leftarrow} E \overset{\gamma}{\rightarrow} C \) to satisfy \(\gamma' \star \beta' = \gamma \star \alpha'\), graph $E$ must not contain any edges as $\alpha'$ and $\beta'$ do not map any edges from $C$ and $L$ to a common edge in $G$.
    %  For all graph $E$ without edges and all morphisms $\gamma': E \rightarrow L $ and $\gamma : E \rightarrow C$, we have \(\gamma' \star \beta' = \gamma \star \alpha'\). Three examples of such spans are illustrated in~\autoref{fig:preliminaries:three_squares}.
    % \begin{figure}[H]
    %     \centering
    %         \resizebox{0.7\textwidth}{!}{
    %         \begin{tikzpicture} 
    %             \graphbox{\( E_1 \)}{-50mm}{0mm}{34mm}{12mm}{2mm}{2mm}{
    %                 \coordinate (o) at (0mm,-8mm); 
    %                 \node[draw,circle] (l1) at ($(o)+(-10mm,0mm)$) {1};
    %                 \node[draw,circle] (l2) at ($(l1)+(2,0)$) {2};
    %             }  
    %             \graphbox{\( L \)}{40mm}{20mm}{34mm}{20mm}{2mm}{2mm}{
    %                 \coordinate (o) at (0mm,-8mm); 
    %                 \node[draw,circle] (l1) at ($(o)+(-10mm,0mm)$) {1};
    %                 \node[draw,circle] (l2) at ($(l1)+(2,0)$) {2};
    %                 \draw[->,red] (l2) edge[out=-135,in=-45]node[midway,below] {$a$} (l1) ;
    %                 \node[draw,circle,red] (l3) at ($(l1) + (1,0)$) {3};
    %                 \draw[->,red] (l1) -- (l3) node[midway,above] {$a$};
    %                 \draw[->,red] (l3) -- (l2) node[midway,above] {$a$};
    %             } 
        
    %             \graphbox{\( K \)}{0mm}{0mm}{34mm}{12mm}{2mm}{2mm}{
    %                 \coordinate (o) at (0mm,-8mm); 
    %                 \node[draw,circle] (l1) at ($(o)+(-10mm,0mm)$) {1};
    %                 \node[draw,circle] (l2) at ($(l1)+(2,0)$) {2};
    %             }  

    %             \graphbox{\( C \)}{40mm}{-20mm}{34mm}{35mm}{2mm}{-3mm}{
    %                 \coordinate (o) at (0mm,-14mm); 
    %                 \node[draw,circle] (l1) at ($(o)+(0,0)$) {1\ 2};
    %                 % \node[draw,circle] (l2) at ($(l1)+(2,0)$) {2};
    %                 \node[draw,circle,blue] (l4) at ($(l1) + (1,-1)$) {6};
    %                 \draw[->,blue] (l1) -- (l4) node[midway,right] {$a$};
    %                 % \draw[->,blue] (l1) edge[loop above] node[midway, above] {$a$} (l1) ;
    %                 \node[ draw,circle,blue] (l6) at ($(l1) + (-1,-1)$) {7};
    %                 \draw[<-,blue] (l1) -- (l6) node[midway,left] {$a$};
    %             }

    %             \graphbox{\(G  \)}{110mm}{5mm}{34mm}{40mm}{2mm}{-3mm}{
    %                 \coordinate (o) at (0mm,-20mm); 
    %                 \node[draw,circle] (l1) at ($(o)+(0,0)$) {1\ 2};
    %                 % \node[draw,circle] (l1) at ($(o)+(-10mm,0mm)$) {1};
    %                 % \node[draw,circle] (l2) at ($(l1)+(2,0)$) {2};
    %                 \draw[->,red] (l1) edge[loop below] node[midway, below] {$a$} (l1) ;
    %                 \node[draw,circle,red] (l3) at ($(l1) + (0,1.4)$) {3};
    %                 \node[draw,circle,blue] (l4) at ($(l1) + (1,-1)$) {6};
    %                 \draw[->,red] (l1) edge[bend left] node[midway,left] {$a$} (l3);
    %                 \draw[->,red] (l3) edge[bend left] node[midway,right] {$a$} (l1);
    %                 \draw[->,blue] (l1) edge  node[midway,right] {$a$} (l4);
    %                 \node[draw,circle,blue] (l6) at ($(l1) + (-1,-1)$) {7};
    %                 \draw[<-,blue] (l1) edge node[midway,left] {$a$} (l6) ;
    %                 % \draw[->,blue] (l1) edge[out=-135,in=-45]node[midway,below] {$a$} (l1) ;
    %             }   
    %             % E to K
    %             \draw[>->] (-15mm,-6mm) -- node[above] {$\delta_1$} (-1mm,-6mm);
    %             % E to L
    %             \draw[>->] (-33mm,5mm) -- node[above] {$\gamma'_1$} (33mm,15mm);
    %             % K to L  
    %             \draw[>->] (17mm,5mm) -- node[above] {$\alpha$} (37mm,15mm);
    %             % C to G
    %             \draw[>->] (76mm,-28mm)-- node[below] {$\alpha'$} (104mm,-24mm) ;
    %             % K to C
    %             \draw[->] (17mm,-17mm) -- node[below] {$\beta$} (37mm,-28mm);
    %             % E to C
    %             \draw[->] (-33mm,-17mm) -- node[below] {$\gamma_1$} (33mm,-28mm);
    %             % L to G
    %             \draw[->] (76mm,16mm) -- node[above] {$\beta'$} (104mm,7mm);
    %             % \node () at (57mm,-6mm) {$PB$};
    %         \end{tikzpicture}
    %         }

    %         \vspace{5mm}
    %         \resizebox{0.7\textwidth}{!}{
    %         \begin{tikzpicture} 
    %             \graphbox{\( E_2 \)}{-50mm}{0mm}{34mm}{12mm}{2mm}{2mm}{
    %                 \coordinate (o) at (0mm,-8mm); 
    %                 \node[draw,circle] (l1) at ($(o)+(-10mm,0mm)$) {1};
    %                 % \node[draw,circle] (l2) at ($(l1)+(2,0)$) {2};
    %             }  
    %             \graphbox{\( L \)}{40mm}{20mm}{34mm}{20mm}{2mm}{2mm}{
    %                 \coordinate (o) at (0mm,-8mm); 
    %                 \node[draw,circle] (l1) at ($(o)+(-10mm,0mm)$) {1};
    %                 \node[draw,circle] (l2) at ($(l1)+(2,0)$) {2};
    %                 \draw[->,red] (l2) edge[out=-135,in=-45]node[midway,below] {$a$} (l1) ;
    %                 \node[draw,circle,red] (l3) at ($(l1) + (1,0)$) {3};
    %                 \draw[->,red] (l1) -- (l3) node[midway,above] {$a$};
    %                 \draw[->,red] (l3) -- (l2) node[midway,above] {$a$};
    %             } 
        
    %             \graphbox{\( K \)}{0mm}{0mm}{34mm}{12mm}{2mm}{2mm}{
    %                 \coordinate (o) at (0mm,-8mm); 
    %                 \node[draw,circle] (l1) at ($(o)+(-10mm,0mm)$) {1};
    %                 \node[draw,circle] (l2) at ($(l1)+(2,0)$) {2};
    %             }  

    %             \graphbox{\( C \)}{40mm}{-20mm}{34mm}{35mm}{2mm}{-3mm}{
    %                 \coordinate (o) at (0mm,-14mm); 
    %                 \node[draw,circle] (l1) at ($(o)+(0,0)$) {1\ 2};
    %                 % \node[draw,circle] (l2) at ($(l1)+(2,0)$) {2};
    %                 \node[draw,circle,blue] (l4) at ($(l1) + (1,-1)$) {6};
    %                 \draw[->,blue] (l1) -- (l4) node[midway,right] {$a$};
    %                 % \draw[->,blue] (l1) edge[loop above] node[midway, above] {$a$} (l1) ;
    %                 \node[ draw,circle,blue] (l6) at ($(l1) + (-1,-1)$) {7};
    %                 \draw[<-,blue] (l1) -- (l6) node[midway,left] {$a$};
    %             }

    %             \graphbox{\(G  \)}{110mm}{5mm}{34mm}{40mm}{2mm}{-3mm}{
    %                 \coordinate (o) at (0mm,-20mm); 
    %                 \node[draw,circle] (l1) at ($(o)+(0,0)$) {1\ 2};
    %                 % \node[draw,circle] (l1) at ($(o)+(-10mm,0mm)$) {1};
    %                 % \node[draw,circle] (l2) at ($(l1)+(2,0)$) {2};
    %                 \draw[->,red] (l1) edge[loop below] node[midway, below] {$a$} (l1) ;
    %                 \node[draw,circle,red] (l3) at ($(l1) + (0,1.4)$) {3};
    %                 \node[draw,circle,blue] (l4) at ($(l1) + (1,-1)$) {6};
    %                 \draw[->,red] (l1) edge[bend left] node[midway,left] {$a$} (l3);
    %                 \draw[->,red] (l3) edge[bend left] node[midway,right] {$a$} (l1);
    %                 \draw[->,blue] (l1) edge  node[midway,right] {$a$} (l4);
    %                 \node[draw,circle,blue] (l6) at ($(l1) + (-1,-1)$) {7};
    %                 \draw[<-,blue] (l1) edge node[midway,left] {$a$} (l6) ;
    %                 % \draw[->,blue] (l1) edge[out=-135,in=-45]node[midway,below] {$a$} (l1) ;
    %             }   
    %             % E to K
    %             \draw[>->] (-15mm,-6mm) -- node[above] {$\delta_2$} (-1mm,-6mm);
    %             % E to L
    %             \draw[>->] (-33mm,5mm) -- node[above] {$\gamma'_2$} (33mm,15mm);
    %             % K to L  
    %             \draw[>->] (17mm,5mm) -- node[above] {$\alpha$} (37mm,15mm);
    %             % C to G
    %             \draw[>->] (76mm,-28mm)-- node[below] {$\alpha'$} (104mm,-24mm) ;
    %             % K to C
    %             \draw[->] (17mm,-17mm) -- node[below] {$\beta$} (37mm,-28mm);
    %             % E to C
    %             \draw[->] (-33mm,-17mm) -- node[below] {$\gamma_2$} (33mm,-28mm);
    %             % L to G
    %             \draw[->] (76mm,16mm) -- node[above] {$\beta'$} (104mm,7mm);
    %             % \node () at (57mm,-6mm) {$PB$};
    %         \end{tikzpicture}
    %         }

    %         \vspace{5mm}
    %         \resizebox{0.7\textwidth}{!}{
    %         \begin{tikzpicture} 
    %             \graphbox{\( E_3 \)}{-50mm}{0mm}{34mm}{12mm}{2mm}{2mm}{
    %                 \coordinate (o) at (0mm,-8mm); 
    %                 % \node[draw,circle] (l1) at ($(o)+(-10mm,0mm)$) {1};
    %                 \node[draw,circle] (l2) at ($(o)+(10mm,0mm)$) {2};
    %             }  
    %             \graphbox{\( L \)}{40mm}{20mm}{34mm}{20mm}{2mm}{2mm}{
    %                 \coordinate (o) at (0mm,-8mm); 
    %                 \node[draw,circle] (l1) at ($(o)+(-10mm,0mm)$) {1};
    %                 \node[draw,circle] (l2) at ($(l1)+(2,0)$) {2};
    %                 \draw[->,red] (l2) edge[out=-135,in=-45]node[midway,below] {$a$} (l1) ;
    %                 \node[draw,circle,red] (l3) at ($(l1) + (1,0)$) {3};
    %                 \draw[->,red] (l1) -- (l3) node[midway,above] {$a$};
    %                 \draw[->,red] (l3) -- (l2) node[midway,above] {$a$};
    %             } 
        
    %             \graphbox{\( K \)}{0mm}{0mm}{34mm}{12mm}{2mm}{2mm}{
    %                 \coordinate (o) at (0mm,-8mm); 
    %                 \node[draw,circle] (l1) at ($(o)+(-10mm,0mm)$) {1};
    %                 \node[draw,circle] (l2) at ($(l1)+(2,0)$) {2};
    %             }  

    %             \graphbox{\( C \)}{40mm}{-20mm}{34mm}{35mm}{2mm}{-3mm}{
    %                 \coordinate (o) at (0mm,-14mm); 
    %                 \node[draw,circle] (l1) at ($(o)+(0,0)$) {1\ 2};
    %                 % \node[draw,circle] (l2) at ($(l1)+(2,0)$) {2};
    %                 \node[draw,circle,blue] (l4) at ($(l1) + (1,-1)$) {6};
    %                 \draw[->,blue] (l1) -- (l4) node[midway,right] {$a$};
    %                 % \draw[->,blue] (l1) edge[loop above] node[midway, above] {$a$} (l1) ;
    %                 \node[ draw,circle,blue] (l6) at ($(l1) + (-1,-1)$) {7};
    %                 \draw[<-,blue] (l1) -- (l6) node[midway,left] {$a$};
    %             }

    %             \graphbox{\(G  \)}{110mm}{5mm}{34mm}{40mm}{2mm}{-3mm}{
    %                 \coordinate (o) at (0mm,-20mm); 
    %                 \node[draw,circle] (l1) at ($(o)+(0,0)$) {1\ 2};
    %                 % \node[draw,circle] (l1) at ($(o)+(-10mm,0mm)$) {1};
    %                 % \node[draw,circle] (l2) at ($(l1)+(2,0)$) {2};
    %                 \draw[->,red] (l1) edge[loop below] node[midway, below] {$a$} (l1) ;
    %                 \node[draw,circle,red] (l3) at ($(l1) + (0,1.4)$) {3};
    %                 \node[draw,circle,blue] (l4) at ($(l1) + (1,-1)$) {6};
    %                 \draw[->,red] (l1) edge[bend left] node[midway,left] {$a$} (l3);
    %                 \draw[->,red] (l3) edge[bend left] node[midway,right] {$a$} (l1);
    %                 \draw[->,blue] (l1) edge  node[midway,right] {$a$} (l4);
    %                 \node[draw,circle,blue] (l6) at ($(l1) + (-1,-1)$) {7};
    %                 \draw[<-,blue] (l1) edge node[midway,left] {$a$} (l6) ;
    %                 % \draw[->,blue] (l1) edge[out=-135,in=-45]node[midway,below] {$a$} (l1) ;
    %             }   
    %             % E to K
    %             \draw[>->] (-15mm,-6mm) -- node[above] {$\delta_3$} (-1mm,-6mm);
    %             % E to L
    %             \draw[>->] (-33mm,5mm) -- node[above] {$\gamma'_3$} (33mm,15mm);
    %             % K to L  
    %             \draw[>->] (17mm,5mm) -- node[above] {$\alpha$} (37mm,15mm);
    %             % C to G
    %             \draw[>->] (76mm,-28mm)-- node[below] {$\alpha'$} (104mm,-24mm) ;
    %             % K to C
    %             \draw[->] (17mm,-17mm) -- node[below] {$\beta$} (37mm,-28mm);
    %             % E to C
    %             \draw[->] (-33mm,-17mm) -- node[below] {$\gamma_3$} (33mm,-28mm);
    %             % L to G
    %             \draw[->] (76mm,16mm) -- node[above] {$\beta'$} (104mm,7mm);
    %             % \node () at (57mm,-6mm) {$=$};
    %         \end{tikzpicture}
    %         }
    %     \caption{}
    %     \label{fig:preliminaries:three_squares} 
    % \end{figure}
\end{example} 
An important phenomenon about pushouts and pullbacks in the category \textbf{Graph} is stated in the following proposition. For examples, consider the diagrams in~\autoref{fig:preliminaries:pushout_non_injective} and~\autoref{fig:preliminaries:a_rewriting_step_dfjalsdkdfsdfjflg}
\begin{proposition}[\text{\cite[Lemma 13]{lack2004adhesive}}]
    \label{prop:pb_eq_po}
    In the category \textbf{Graph}, pushouts along monomorphisms are also pullbacks. 
\end{proposition}

\section{Double-Pushout Rewriting} 
\label{sec:dpo}
\begin{definition}[Rewriting rule and match~\cite{corradini1997algebraic}]
  \label{def:grs:dpo_rule}
A \textbf{DPO rewriting rule} $\rho$ is a span \( L \overset{l}{\leftarrow} K \overset{r}{\rightarrow} R \), where \( K \) is the \textbf{interface}, \( L \) is the \textbf{left-hand-side graph}, denoted \( \operatorname{lhs}(\rho) \), and \( R \) is the \textbf{right-hand-side graph}, denoted \( \operatorname{rhs}(\rho) \). We denote $\rho^{-1}$ as the rule $(R \overset{r}{\leftarrow} K \overset{l}{\rightarrow} L)$. The rule is \textbf{left monic} if \( l \) is monic, and \textbf{right monic} if \( r \) is monic. 
The rule is \textbf{monic} if $l$ and $r$ are both monic.
A \textbf{match} of the rule in a graph \( G \) is a morphism \( m: L \rightarrow G \).
\end{definition}
In the case of the category \textbf{Graph} of edge-labeled directed multigraphs, the rule is left monic if the morphism \( l \) is injective, and right monic if the morphism \( r \) is injective. The rule is monic if both \( l \) and \( r \) are injective.

\begin{example}
  \label{ex:grsaa}
  The follwing DPO rewriting rule on edge-labeled directed multigraphs is from \cite[Example 6]{bruggink2014termination}:
  \begin{center} 
      \resizebox{0.7\textwidth}{!}{
      \begin{tikzpicture}
          \graphbox{$L$}{0mm}{0mm}{34mm}{15mm}{2mm}{-5mm}{
              \coordinate (o) at (0mm,-3mm); 
              \node[draw,circle] (l1) at ($(o)+(-10mm,0mm)$) {1};
              \node[draw,circle] (l2) at ($(l1)+(2,0)$) {2};
              \node[draw,circle] (l3) at ($(l1) + (1,0)$) {3};
              \draw[->] (l1) -- (l3) node[midway,above] {a};
              \draw[->] (l3) -- (l2) node[midway,above] {a};
          }     
          \graphbox{$K$}{40mm}{0mm}{24mm}{15mm}{2mm}{-5mm}{
              \coordinate (o) at (5mm,-3mm); 
              \node[draw,circle] (l1) at ($(o)+(-10mm,0mm)$) {1};
              \node[draw,circle] (l2) at ($(l1)+(1,0)$) {2};
              % \node[draw,circle] (l3) at ($(l1) + (1,0)$) {$\ $};
              % \draw[->] (l1) -- (l3) node[midway,above] {a};
              % \draw[->] (l3) -- (l2) node[midway,above] {a};
          }    
          \graphbox{$R$}{70mm}{0mm}{45mm}{15mm}{2mm}{-5mm}{
              \coordinate (o) at (-5mm,-3mm); 
              \node[draw,circle] (l1) at ($(o)+(-10mm,0mm)$) {1};
              \node[draw,circle] (l2) at ($(l1)+(3,0)$) {2};
              \node[draw,circle] (l3) at ($(l1) + (1,0)$) {4};
              \node[draw,circle] (l4) at ($(l1) + (2,0)$) {5};
              \draw[->] (l1) -- (l3) node[midway,above] {a};
              \draw[->] (l3) -- (l4) node[midway,above] {b};
              \draw[->] (l4) -- (l2) node[midway,above] {a};
          }    
          \node () at (37mm,-8mm) {$\overset{l}{\leftarrowtail}$};
          \node () at (67mm,-8mm) {$\overset{r}{\rightarrowtail}$};
          % \draw[>->] (51mm,2mm) -- (52mm,3mm);
      \end{tikzpicture}
      }
  \end{center}
  This rule has a discrete interface \( K \). It is monic (or injective). It replaces a chain \enquote{aa} whose middle node has no other incident edges by a chain \enquote{aba} with the same endpoints and all other nodes and edges are fresh. 
\end{example}


Intuitively, $K$ is the common part of $L$ and $R$. To apply the rule to an object $G$, we need to find a match $m:L \to G$. The rule is then applied by constructing the diagram with double pushout shown in \autoref{def:rewriting_step} and once the pushout squares are constructed, one says that the object $G$ rewrites to $H$ using the rule $\rho$ and the match $m$. 

% \begin{center}
%      \resizebox{0.4\textwidth}{!}{
%           \begin{tikzpicture}
%             % [node distance=11mm]
%             \node (I) at (0,0) {$K$};
%             \node (L) at (-2,0) {$L$};
%             \node (R) at (2,0) {$R$};
%             \node (G) at (-2,-2) {$G$};
%             \node (C) at (0,-2) {$C$};
%             \node (H) at (2,-2) {$H$};
%             \draw [->] (I) to  node [midway,below] {$l$} (L);
%             \draw [->] (I) to  node [midway,below] {$r$} (R);
%             \draw [->] (L) to node [midway,right] {$m$} (G);
%             \draw [->] (I) to node [midway,right] {$u$} (C);
%             \draw [->] (R) to node [midway,left] {$m'$} (H);
%             \draw [->] (C) to node [midway,above] {$l'$} (G);
%             \draw [->] (C) to node [midway,above] {$r'$} (H);
%             \node [at=($(I)!.5!(G)$)] {\normalfont PO};
%             \node [at=($(I)!.5!(H)$)] {\normalfont PO};
%           \end{tikzpicture}
%         % \end{center}
%         }
% \end{center}

\begin{definition}[DPO Rewriting step~\cite{endrullis2024generalized_arxiv_v2}]
  \label{def:rewriting_step}
    \ \newline
    \noindent
    \begin{minipage}{0.72\textwidth}
      A DPO diagram $\delta$ is a diagram as shown on the right.
      This diagram $\delta$ is a witness for the \textbf{rewriting step} from \( G \) to \( H \) using 
      a rule \( \rho = (L \overset{l}{\leftarrow} K \overset{r}{\rightarrow} R) \) and a \textbf{match}, denoted \( G \Rightarrow_\rho^m H \) or \( G \Rightarrow_\rho^\delta H \). We denote $\operatorname{left}(\delta)$ and $\operatorname{right}(\delta)$ the pushout squares $KLGC$ and $KRHC$, respectively.
    \end{minipage}
    \hfill
    \begin{minipage}{0.28\textwidth}
          % \begin{center}
          \hfill
          \resizebox{0.85\textwidth}{!}{
          \begin{tikzpicture}
            % [node distance=11mm]
            \node (I) at (0,0) {$K$};
            \node (L) at (-2,0) {$L$};
            \node (R) at (2,0) {$R$};
            \node (G) at (-2,-2) {$G$};
            \node (C) at (0,-2) {$C$};
            \node (H) at (2,-2) {$H$};
            \draw [->] (I) to  node [midway,below] {$l$} (L);
            \draw [->] (I) to  node [midway,below] {$r$} (R);
            \draw [->] (L) to node [midway,right] {$m$} (G);
            \draw [->] (I) to node [midway,right] {$u$} (C);
            \draw [->] (R) to node [midway,left] {$m'$} (H);
            \draw [->] (C) to node [midway,above] {$l'$} (G);
            \draw [->] (C) to node [midway,above] {$r'$} (H);
            \node [at=($(I)!.5!(G)$)] {\normalfont PO};
            \node [at=($(I)!.5!(H)$)] {\normalfont PO};
          \end{tikzpicture}
        % \end{center}
        }
        \end{minipage}
  \end{definition}
\begin{example}
  \label{ex:rewriting_step_grs_aa}
  The DPO diagram below defines a rewriting step using the rule in \autoref{ex:grsaa}.
  \begin{center} 
      \resizebox{0.7\textwidth}{!}{
      \begin{tikzpicture}
          \graphbox{\( L \)}{0mm}{-3mm}{34mm}{12mm}{2mm}{2mm}{
              \coordinate (o) at (0mm,-8mm); 
              \node[draw,circle] (l1) at ($(o)+(-10mm,0mm)$) {1};
              \node[draw,circle] (l2) at ($(l1)+(2,0)$) {2};
              \node[draw,circle] (l3) at ($(l1) + (1,0)$) {3};
              \draw[] (l1) -- (l3) node[midway,above] {a};
              \draw[] (l3) -- (l2) node[midway,above] {a};
          } 
          \graphbox{\( K \)}{40mm}{-3mm}{34mm}{12mm}{2mm}{2mm}{
              \coordinate (o) at (0mm,-8mm); 
              \node[draw,circle] (l1) at ($(o)+(-10mm,0mm)$) {1};
              \node[draw,circle] (l2) at ($(l1)+(2,0)$) {2};
          }  
          \graphbox{\( R \)}{80mm}{-3mm}{45mm}{12mm}{2mm}{2mm}{
              \coordinate (o) at (-5mm,-8mm); 
              \node[draw,circle] (l1) at ($(o)+(-10mm,0mm)$) {1};
              \node[draw,circle] (l2) at ($(l1)+(3,0)$) {2};
              \node[draw,circle] (l3) at ($(l1) + (1,0)$) {4};
              \node[draw,circle] (l4) at ($(l1) + (2,0)$) {5};
              \draw[ ] (l1) -- (l3) node[midway,above] {a};
              \draw[ ] (l3) -- (l4) node[midway,above] {b};
              \draw[ ] (l4) -- (l2) node[midway,above] {a};
          }    
          \graphbox{\( G \)}{0mm}{-22mm}{34mm}{22mm}{2mm}{-3mm}{
              \coordinate (o) at (0mm,-3mm); 
              \node[draw,circle] (l1) at ($(o)+(-10mm,0mm)$) {1};
              \node[draw,circle] (l2) at ($(l1)+(2,0)$) {2};
              \node[draw,circle] (l3) at ($(l1) + (1,0)$) {3};
              \node[draw,circle] (l4) at ($(l2) + (0,-1)$) {6};
              \draw[] (l1) -- (l3) node[midway,above] {a};
              \draw[] (l3) -- (l2) node[midway,above] {a};
              \draw[ ] (l2) -- (l4) node[midway,right] {a};
              \node[draw,circle] (l6) at ($(l1) + (0,-1)$) {7};
              \draw[] (l1) -- (l6) node[midway,left] {a};
          }    
          \graphbox{\( C  \)}{40mm}{-22mm}{34mm}{22mm}{2mm}{-3mm}{
              \coordinate (o) at (0mm,-3mm); 
              \node[draw,circle] (l1) at ($(o)+(-10mm,0mm)$) {1};
              \node[draw,circle] (l2) at ($(l1)+(2,0)$) {2};
              \node[draw,circle] (l4) at ($(l2) + (0,-1)$) {6};
              \draw[ ] (l2) -- (l4) node[midway,right] {a};
              \node[ draw,circle] (l6) at ($(l1) + (0,-1)$) {7};
              \draw[ ] (l1) -- (l6) node[midway,left] {a};
          }    
          \graphbox{\( H \)}{80mm}{-22mm}{45mm}{22mm}{2mm}{-3mm}{
              \coordinate (o) at (-5mm,-3mm); 
              \node[draw,circle] (l1) at ($(o)+(-10mm,0mm)$) {1};
              \node[draw,circle] (l2) at ($(l1)+(3,0)$) {2};
              \node[draw,circle] (l3) at ($(l1) + (1,0)$) {4};
              \node[draw,circle] (l4) at ($(l1) + (2,0)$) {5};
              \node[ draw,circle] (l5) at ($(l2) + (0,-1)$) {6};
              \node[ draw,circle] (l6) at ($(l1) + (0,-1)$) {7};
              \draw[ ] (l1) -- (l6) node[midway,left] {a};
              \draw[] (l1) -- (l3) node[midway,above] {a};
              \draw[] (l3) -- (l4) node[midway,above] {b};
              \draw[ ] (l4) -- (l2) node[midway,above] {a};
              \draw[ ] (l2) -- (l5) node[midway,right] {a};
          }    
          \node () at (37mm,-8mm) {\( \leftarrowtail \)}; % K -> L
          \node () at (77mm,-8mm) {\( \rightarrowtail \)}; % K -> R
          \node () at (15mm,-18mm) {\( m\ \downarrowtail \)};
          \node () at (37mm,-33mm) {\( \leftarrowtail \)};
          \node () at (58mm,-18mm) {\( u\downarrowtail \)};
          \node () at (102mm,-18mm) {\( \downarrowtail \)};
          \node () at (77mm,-33mm) {\( \rightarrowtail \)}; % C -> H
      \end{tikzpicture}
      }
  \end{center}
\end{example}
In category \textbf{Graph}, the pushout of two arrows always exsits~\cite[p.188]{corradini1997algebraic}. Therefore, once the first pushout square is constructed, the second pushout square can always be constructed and is unique up to isomorphism because of the universal property. However, the first pushout square cannot always be constructed.

\begin{proposition}[Existence of pushout complements~\cite{corradini1997algebraic}]
    Let $A \overset{l}{\rightarrow} B \overset{m}{\rightarrow} D$ be two morphisms in the category \textbf{Graph}. There exist two morphisms $A \overset{u}{\rightarrow} C \overset{l'}{\rightarrow} D$ such that $A \overset{l}{\rightarrow} B \overset{m}{\rightarrow} D$ and $A \overset{u}{\rightarrow} C \overset{l'}{\rightarrow} D$ form a pushout square as shown below

    \begin{center}
        \resizebox{0.2\textwidth}{!}{
            \begin{tikzpicture}
                \node (I) at (0,0) {$K$};
                \node (L) at (-2,0) {$L$};
                \node (G) at (-2,-2) {$G$};
                \node (C) at (0,-2) {$C$};
                \draw [->] (I) to  node [midway,below] {$l$} (L);
                \draw [->] (L) to node [midway,right] {$m$} (G);
                \draw [->] (I) to node [midway,right] {$u$} (C);
                \draw [->] (C) to node [midway,above] {$l'$} (G);
                \node [at=($(I)!.5!(G)$)] {\normalfont PO};
            \end{tikzpicture}
        }
    \end{center}
    if and only if the following conditions are satisfied:
    \begin{itemize}
        \item{[Dangling edge condition]} No edge of $E(G) \setminus m(E(L))$ is incident to any node in $m(E(L) \setminus l(E(K)))$ are the edge-sets of $G,L,K$,
        \item{[Identification condition]} There is no $x,y \in V(L) \cup E(L)$ such that $x \not = y$, $m(x) = m(y)$ and $y \notin l(V(K) \cup E(K))$ where 
    \end{itemize} 
\end{proposition}
Intuitively, an edge $e$ in $G_1 \setminus m(E(L))$ is an edge which is not in the image of the match $m$. If such an edge $e$ is incident to a node in $m(V(L) \setminus l(V(K)))$ which are nodes that will be removed in the rewriting step and thus neither contained in $C$ nor in $H$, the edge $e$ will be dangling in $H$ after the rewriting step. A match which does not satisfy the dangling edge condition is shown below. In this example, when node 3 is removed, the edge from node 3 to node 6 will be dangling. (to double check)
    \begin{center}
        \resizebox{0.9\textwidth}{!}{
        \begin{tikzpicture}
            \graphbox{\( L \)}{0mm}{-3mm}{34mm}{12mm}{2mm}{2mm}{
                \coordinate (o) at (0mm,-8mm); 
                \node[draw,circle] (l1) at ($(o)+(-10mm,0mm)$) {1};
                \node[draw,circle] (l2) at ($(l1)+(2,0)$) {2};
                \node[draw,circle] (l3) at ($(l1) + (1,0)$) {3};
                \draw[->] (l1) -- (l3) node[midway,above] {a};
                \draw[->] (l3) -- (l2) node[midway,above] {a};
            } 
    
            \graphbox{\( K \)}{40mm}{-3mm}{34mm}{12mm}{2mm}{2mm}{
                \coordinate (o) at (0mm,-8mm); 
                \node[draw,circle] (l1) at ($(o)+(-10mm,0mm)$) {1};
                \node[draw,circle] (l2) at ($(l1)+(2,0)$) {2};
            }  
    
            \graphbox{\( R \)}{80mm}{-3mm}{45mm}{12mm}{2mm}{2mm}{
                \coordinate (o) at (-5mm,-8mm); 
                \node[draw,circle] (l1) at ($(o)+(-10mm,0mm)$) {1};
                \node[draw,circle] (l2) at ($(l1)+(3,0)$) {2};
                \node[draw,circle] (l3) at ($(l1) + (1,0)$) {4};
                \node[draw,circle] (l4) at ($(l1) + (2,0)$) {5};
                \draw[->] (l1) -- (l3) node[midway,above] {a};
                \draw[->] (l3) -- (l4) node[midway,above] {b};
                \draw[->] (l4) -- (l2) node[midway,above] {a};
            }    
    
            \graphbox{\( G' \)}{0mm}{-22mm}{34mm}{22mm}{2mm}{-3mm}{
                \coordinate (o) at (0mm,-3mm); 
                \node[draw,circle] (l1) at ($(o)+(-10mm,0mm)$) {1};
                \node[draw,circle] (l2) at ($(l1)+(2,0)$) {2};
                \node[draw,circle] (l3) at ($(l1) + (1,0)$) {3};
                \node[draw,circle] (l4) at ($(l2) + (0,-1)$) {6};
                \draw[->,red] (l3) -- (l4) node[midway,above] {a};
                \draw[->] (l1) -- (l3) node[midway,above] {a};
                \draw[->] (l3) -- (l2) node[midway,above] {a};
                \draw[->] (l2) -- (l4) node[midway,right] {a};
                \node[draw,circle] (l6) at ($(l1) + (0,-1)$) {7};
                \draw[<-] (l1) -- (l6) node[midway,left] {a};
                \draw[->] (l2) edge[out=-135,in=-45]node[midway,below] {a} (l1) ;
            }    
    
            \graphbox{\( C'  \)}{40mm}{-22mm}{34mm}{22mm}{2mm}{-3mm}{
                \coordinate (o) at (0mm,-3mm); 
                \node[draw,circle] (l1) at ($(o)+(-10mm,0mm)$) {1};
                \node[draw,circle] (l2) at ($(l1)+(2,0)$) {2};
                \node[draw,circle] (l4) at ($(l2) + (0,-1)$) {6};
                \node[draw,circle,dashed,red] (l3) at ($(l1) + (1,0)$) {3};
                % \draw[->,red] (l3) -- (l4) node[midway,above] {a};
                \draw[->,red] (l3) -- (l4) node[midway,above] {a};
                \draw[->] (l2) -- (l4) node[midway,right] {a};
                \draw[->] (l2) edge[out=-135,in=-45]node[midway,below] {a} (l1) ;
                \node[draw,circle] (l6) at ($(l1) + (0,-1)$) {7};
                \draw[<-] (l1) -- (l6) node[midway,left] {a};
            }    
     
            % \graphbox{\( H' \)}{80mm}{-22mm}{45mm}{22mm}{2mm}{-3mm}{
            %     \coordinate (o) at (-5mm,-3mm); 
            %     \node[draw,circle] (l1) at ($(o)+(-10mm,0mm)$) {1};
            %     \node[draw,circle] (l2) at ($(l1)+(3,0)$) {2};
            %     \node[draw,circle] (l3) at ($(l1) + (1,0)$) {4};
            %     \node[draw,circle] (l4) at ($(l1) + (2,0)$) {5};
            %     \node[draw,circle] (l5) at ($(l2) + (0,-1)$) {6};
            %     \node[draw,circle] (l6) at ($(l1) + (0,-1)$) {7};
            %     \draw[<-] (l1) -- (l6) node[midway,left] {a};
            %     \draw[->] (l1) -- (l3) node[midway,above] {a};
            %     \draw[->] (l3) -- (l4) node[midway,above] {b};
            %     \draw[->] (l4) -- (l2) node[midway,above] {a};
            %     \draw[->] (l2) -- (l5) node[midway,right] {a};
            %     \draw[->] (l2) edge[out=-135,in=-45]node[midway,below] {a} (l1) ;
            %     \node[draw,circle,dashed,red] (l3) at (-0mm,-7mm) {3};
            %     \draw[->] (l3) -- (l5) node[midway,above] {a};
            % }    
    
            \node () at (37mm,-8mm) {\( \leftarrowtail \)}; % K -> L
            \node () at (77mm,-8mm) {\( \rightarrowtail \)}; % K -> R
            \node () at (15mm,-18mm) {\(\downarrowtail \)};
            \node () at (37mm,-33mm) {\( \leftarrowtail \)};
            % \node () at (37mm,-18mm) {PO};
            \node () at (58mm,-18mm) {\( \downarrowtail \)}; 
            % \node () at (80mm,-18mm) {PO};
            % \node () at (102mm,-18mm) {\( \downarrowtail \)};
            % \node () at (77mm,-33mm) {\( \rightarrowtail \)}; % C -> H
        \end{tikzpicture}
        }         
    \end{center}
To give an intuition on the second condition, suppose that there is a graph C and morphisms $\beta$ and $\alpha'$ such that the following diagram is a pushout square.
\begin{center}
        \resizebox{0.7\textwidth}{!}{
        \begin{tikzpicture} 
            \graphbox{\( L \)}{40mm}{20mm}{34mm}{20mm}{2mm}{2mm}{
                \coordinate (o) at (0mm,-8mm); 
                \node[draw,circle] (l1) at ($(o)+(-10mm,0mm)$) {1};
                \node[draw,circle] (l2) at ($(l1)+(2,0)$) {2};
                % \draw[->,red] (l2) edge[out=-135,in=-45]node[midway,below] {a} (l1) ;
                % \node[draw,circle,red] (l3) at ($(l1) + (1,0)$) {3};
                % \draw[->,red] (l1) -- (l3) node[midway,above] {a};
                % \draw[->,red] (l3) -- (l2) node[midway,above] {a};
            } 
    
            \graphbox{\( K \)}{0mm}{0mm}{34mm}{12mm}{2mm}{2mm}{
                \coordinate (o) at (0mm,-8mm); 
                \node[draw,circle] (l1) at ($(o)+(-10mm,0mm)$) {1};
                % \node[draw,circle] (l2) at ($(l1)+(2,0)$) {2};
            }  
            \graphbox{\(G  \)}{110mm}{5mm}{34mm}{40mm}{2mm}{-3mm}{
                \coordinate (o) at (0mm,-20mm); 
                 \node[draw,circle] (l1) at ($(o)+(0,0)$) {1\ 2};
                % \node[draw,circle] (l1) at ($(o)+(-10mm,0mm)$) {1};
                % \node[draw,circle] (l2) at ($(l1)+(2,0)$) {2};
                % \draw[->,red] (l1) edge[loop below] node[midway, below] {$a$} (l1) ;
                % \node[draw,circle,red] (l3) at ($(l1) + (0,1.4)$) {3};
                % \node[draw,circle,blue] (l4) at ($(l1) + (1,-1)$) {6};
                % \draw[->,red] (l1) edge[bend left] node[midway,left] {a} (l3);
                % \draw[->,red] (l3) edge[bend left] node[midway,right] {a} (l1);
                % \draw[->,blue] (l1) edge  node[midway,right] {a} (l4);
                % \node[draw,circle,blue] (l6) at ($(l1) + (-1,-1)$) {7};
                % \draw[<-,blue] (l1) edge node[midway,left] {a} (l6) ;
                % \draw[->,blue] (l1) edge[out=-135,in=-45]node[midway,below] {a} (l1) ;
            }   
            \graphbox{\( C \)}{50mm}{-28mm}{20mm}{10mm}{2mm}{-3mm}{
            }
            % K to L  
            \draw[->] (17mm,5mm) -- node[above] {$\alpha$} (37mm,15mm);
            % C to G
            \draw[->] (76mm,-28mm)-- node[below] {$\alpha'$} (104mm,-24mm) ;
            % K to C
            \draw[->] (17mm,-17mm) -- node[below] {$\beta$} (37mm,-28mm);
            % L to G
            \draw[->] (76mm,16mm) -- node[above] {$\beta'$} (104mm,7mm);
            \node () at (57mm,-6mm) {$PO$};
        \end{tikzpicture}
        }
    \end{center} 
   We have the following commutative diagram where the morphism $\gamma$ maps all nodes in $C$ to the node 1 in $E$.
   \begin{center}
        \resizebox{0.7\textwidth}{!}{
        \begin{tikzpicture} 
            \graphbox{\( L \)}{40mm}{20mm}{34mm}{20mm}{2mm}{2mm}{
                \coordinate (o) at (0mm,-8mm); 
                \node[draw,circle] (l1) at ($(o)+(-10mm,0mm)$) {1};
                \node[draw,circle] (l2) at ($(l1)+(2,0)$) {2};
                % \draw[->,red] (l2) edge[out=-135,in=-45]node[midway,below] {a} (l1) ;
                % \node[draw,circle,red] (l3) at ($(l1) + (1,0)$) {3};
                % \draw[->,red] (l1) -- (l3) node[midway,above] {a};
                % \draw[->,red] (l3) -- (l2) node[midway,above] {a};
            } 
    
            \graphbox{\( K \)}{0mm}{0mm}{34mm}{12mm}{2mm}{2mm}{
                \coordinate (o) at (0mm,-8mm); 
                \node[draw,circle] (l1) at ($(o)+(-10mm,0mm)$) {1};
                \node[draw,circle] (l2) at ($(l1)+(2,0)$) {2};
            }  
            \graphbox{\(E \)}{110mm}{5mm}{45mm}{40mm}{2mm}{-3mm}{
                \coordinate (o) at (-10mm,-20mm); 
                 \node[draw,circle] (l1) at ($(o)+(0,0)$) {1};
                % \node[draw,circle] (l1) at ($(o)+(-10mm,0mm)$) {1};
                \node[draw,circle] (l2) at ($(l1)+(2,0)$) {2};
                % \draw[->,red] (l1) edge[loop below] node[midway, below] {$a$} (l1) ;
                % \node[draw,circle,red] (l3) at ($(l1) + (0,1.4)$) {3};
                % \node[draw,circle,blue] (l4) at ($(l1) + (1,-1)$) {6};
                % \draw[->,red] (l1) edge[bend left] node[midway,left] {a} (l3);
                % \draw[->,red] (l3) edge[bend left] node[midway,right] {a} (l1);
                % \draw[->,blue] (l1) edge  node[midway,right] {a} (l4);
                % \node[draw,circle,blue] (l6) at ($(l1) + (-1,-1)$) {7};
                % \draw[<-,blue] (l1) edge node[midway,left] {a} (l6) ;
                % \draw[->,blue] (l1) edge[out=-135,in=-45]node[midway,below] {a} (l1) ;
            }   
     
            % \node () at (50mm,-28mm) {C};
            \graphbox{\( C \)}{50mm}{-28mm}{20mm}{10mm}{2mm}{-3mm}{
            }
            % K to L  
            \draw[->] (17mm,5mm) -- node[above] {$\alpha$} (37mm,15mm);
            % C to G
            \draw[->] (76mm,-28mm)-- node[below] {$\gamma$} (104mm,-24mm) ;
            % K to C
            \draw[->] (17mm,-17mm) -- node[below] {$\beta$} (37mm,-28mm);
            % L to G
            \draw[->] (76mm,16mm) -- node[above] {$\gamma'$} (104mm,7mm);
            % \node () at (57mm,-6mm) {$PO$};
        \end{tikzpicture}
        }
    \end{center} 
    However, there does not exist a morphism \(\delta : G \to E\) such that \( \gamma' = \beta' \star \delta \) and \( \gamma = \alpha' \star \delta \), which contradicts the universal property of the pushout. Therefore, there do not exist a graph \( C \) and morphisms \( \beta \) and \( \alpha' \) such that the diagram in the assumption is a pushout square.

    \begin{definition}[Rewriting framework \cite{endrullis2024generalized_arxiv_v2}]
  \label{def:rewriting_framework}
    A \textbf{DPO rewriting framework} $\mathfrak{F}$ is a mapping of DPO rewriting rules to collections of DPO diagrams. Specifically, for every rule \( \rho = (L \overset{l}{\leftarrow} K \overset{r}{\rightarrow} R) \), the collection $\mathfrak{F}(\rho)$ consists of DPO diagrams of the form shown below.
 \begin{center}
      \resizebox{0.4\textwidth}{!}{
      \begin{tikzpicture}
            \node (I) at (0,0) {$K$};
            \node (L)  at (-2,0) {$L$};
            \node (R)  at (2,0) {$R$};
            \node (G)  at (-2,-2) {$G$};
            \node (C)  at (0,-2) {$C$};
            \node (H)  at (2,-2) {$H$};
            \draw [->] (I) to  node [midway,below] {$l$} (L);
            \draw [->] (I) to  node [midway,below] {$r$} (R);
            \draw [->] (L) to node [midway,right] {} (G);
            \draw [->] (I) to  node [midway,right] 
            % {$u$}
            {} (C); 
            \draw [->] (R) to  node [midway,right] 
            {}
            (H);
            \draw [->] (C) to node [midway,above] {} (G);
            \draw [->] (C) to node [midway,above] 
            {} 
            (H);
        \end{tikzpicture}
        }
    \end{center}
    The \textbf{DPO rewriting relation $\Rightarrow_{\rho,\mathfrak{F}}$ induced by a DPO rewriting rule $\rho$ in $\mathfrak{F}$} is defined as follows: $G \Rightarrow_{\rho,\mathfrak{F}} H$ iff $G \Rightarrow_\rho^\delta H$ for some $\delta \in \mathfrak{F}(\rho)$. 
     The \textbf{DPO rewriting relation $\Rightarrow_{\mathcal{R},\mathfrak{F}}$ induced by a set $\mathcal{R}$ of DPO rewriting rules in $\mathfrak{F}$} is given by: $G \Rightarrow_{\mathcal{R},\mathfrak{F}} H$ iff $G \Rightarrow_{\rho,\mathfrak{F}} H$ for some $\rho \in \mathcal{R}$. When $\mathfrak{F}$ is clear from the context, we 
    suppress $\mathfrak{F}$ and 
    write $\Rightarrow_{\rho}$ and $\Rightarrow_{\mathcal{R}}$.
  \end{definition}
Let \(\mathfrak{F}\) denote the DPO rewriting framework that associates each rule \( \rho = (L \overset{l}{\leftarrow} K \overset{r}{\rightarrow} R) \) with the collection of all DPO diagrams of the form shown in~\autoref{def:rewriting_framework}.
Let \(\mathfrak{M}\) denote the DPO rewriting framework that associates each rule with the collection of all DPO diagrams of the following form with monic match $m$.  
\begin{center}
           \resizebox{0.4\textwidth}{!}{
      \begin{tikzpicture}
            \node (I) at (0,0) {$K$};
            \node (L)  at (-2,0) {$L$};
            \node (R)  at (2,0) {$R$};
            \node (G)  at (-2,-2) {$G$};
            \node (C)  at (0,-2) {$C$};
            \node (H)  at (2,-2) {$H$};
            \draw [->] (I) to  node [midway,below] {$l$} (L);
            \draw [->] (I) to  node [midway,below] {$r$} (R);
            \draw [>->] (L) to node [midway,right] {$m$} (G);
            \draw [>->] (I) to  node [midway,right] 
            % {$u$}
            {} (C);
            \draw [>->] (R) to  node [midway,right] 
            {}
            (H);
            \draw [->] (C) to node [midway,above] {} (G);
            \draw [->] (C) to node [midway,above] 
            {} 
            (H);
        \end{tikzpicture}
        }
    \end{center} 
 
\section{Relative Termination} 
\label{sec:relative_termination}
\begin{definition}[Rewriting sequence]
    Let \(\mathcal{R}\) be a set of rewriting rules and let $\mathfrak{F}$ be a DPO rewriting framework.
    A \textbf{$(\mathcal{R},\mathfrak{F})$-rewriting sequence} is either a finite sequence \(s_0,s_1,\hdots, s_m\) of objects such that \(s_n \Rightarrow_{\mathcal{R},\mathfrak{F}} s_{n+1}\) for each \( 0 \leq n \leq m-1\), or an infinite sequence \(s_0,s_1,\hdots\) of objects such that \(s_n \Rightarrow_{\mathcal{R},\mathfrak{F}} s_{n+1}\) for each \(n \in \mathbb{N}\).
    A $(\mathcal{R},\mathfrak{F})$-rewriting sequence from \( s_0 \) will be denoted as:
    \(
    s_0 \Rightarrow_{\mathcal{R},\mathfrak{F}} s_1 \Rightarrow_{\mathcal{R},\mathfrak{F}} s_2 \Rightarrow_{\mathcal{R},\mathfrak{F}} \cdots 
    \)
\end{definition}

\begin{example}
    Consider the rewriting rules in~\autoref{fig:preliminaries:graph_transformation_rule_nonterminating}.
    \begin{figure}[hbtp]
        \centering
            \resizebox{0.85\textwidth}{!}{
                \begin{tikzpicture}[baseline=-3ex]
                    \graphbox{\( L \)}{0mm}{-3mm}{34mm}{15mm}{2mm}{2mm}{
                        \coordinate (o) at (0mm,-11mm); 
                        \node[draw,circle] (l1) at ($(o)+(-10mm,0mm)$) {1};
                        \node[draw,circle] (l2) at ($(l1)+(2,0)$) {2};
                        \node[draw,circle] (l3) at ($(l1) + (1,0)$) {3};
                        \draw[->] (l1) -- (l3) node[midway,above] {a};
                        \draw[->] (l3) -- (l2) node[midway,above] {b};
                    } 
            
                    \graphbox{\( K \)}{40mm}{-3mm}{34mm}{15mm}{2mm}{2mm}{
                        \coordinate (o) at (0mm,-11mm); 
                        \node[draw,circle] (l1) at ($(o)+(-10mm,0mm)$) {1};
                        \node[draw,circle] (l2) at ($(l1)+(2,0)$) {2};
                    }  
            
                    \graphbox{\( R \)}{80mm}{-3mm}{35mm}{15mm}{2mm}{2mm}{
                        \coordinate (o) at (-5mm,-11mm); 
                        \node[draw,circle] (l1) at ($(o)+(-10mm,0mm)$) {1};
                        % \node[draw,circle] (l2) at ($(l1)+(3,0)$) {2};
                        \node[draw,circle] (l3) at ($(l1) + (1,0)$) {4};
                        \node[draw,circle] (l4) at ($(l1) + (2,0)$) {2};
                        \draw[->] (l1) -- (l3) node[midway,above] {b};
                        \draw[->] (l3) -- (l4) node[midway,above] {a};
                        % \draw[->] (l4) -- (l2) node[midway,above] {a};
                    }    
                    \node () at (37mm,-10mm) {\( \leftarrowtail \)}; % K -> L
                    \node () at (77mm,-10mm) {\( \rightarrowtail \)}; % K -> R
                \end{tikzpicture}
                }
        \caption{}
        \label{fig:preliminaries:graph_transformation_rule_nonterminating}
    \end{figure} 
    A rewriting sequence using the rule with matches in red is shown in~\autoref{fig:preliminaries:sequence_of_transformation_infinite}.
       
        \begin{figure}[hbtp]
            \centering
          \resizebox{0.85\textwidth}{!}{
            \tikz
            [baseline=-0.5ex]
            { 
                \node (x) at (0,0) {$\bullet$};  
                \node (y) at (1,0) {$\bullet$};
                \node (z) at (0.5,0.86) {$\bullet$};
                \draw[->,red] (x) -- node[midway,below] {a} (y) ;
                \draw[->,red] (y) -- node[midway,right] {b} (z) ;
                \draw[->] (z) -- node[midway,left] {b} (x) ;
            } 
            $\Rightarrow$ 
            \tikz[baseline=-0.5ex]{ 
                \node (x) at (0,0) {$\bullet$};  
                \node (y) at (1,0) {$\bullet$};
                \node (z) at (0.5,0.86) {$\bullet$};
                \draw[->] (x) -- node[midway,below] {b} (y) ;
                \draw[->,red] (y) -- node[midway,right] {a} (z) ;
                \draw[->,red] (z) -- node[midway,left] {b} (x) ;
            }
            $\Rightarrow$ 
            \tikz[baseline=-0.5ex]{ 
                \node (x) at (0,0) {$\bullet$};  
                \node (y) at (1,0) {$\bullet$};
                \node (z) at (0.5,0.86) {$\bullet$};
                \draw[->,red] (x) -- node[midway,below] {b} (y) ;
                \draw[->] (y) -- node[midway,right] {b} (z) ;
                \draw[->,red] (z) -- node[midway,left] {a} (x) ;
            }
            $\Rightarrow$ 
            \tikz[baseline=-0.5ex]{ 
                \node (x) at (0,0) {$\bullet$};   
                \node (y) at (1,0) {$\bullet$};
                \node (z) at (0.5,0.86) {$\bullet$};
                \draw[->,red] (x) -- node[midway,below] {a} (y) ;
                \draw[->,red] (y) -- node[midway,right] {b} (z) ;
                \draw[->] (z) -- node[midway,left] {b} (x) ;
            }
          }
          \caption{}
          \label{fig:preliminaries:sequence_of_transformation_infinite}
        \end{figure}
\end{example}
We adapt the concept of relative termination~\cite{klop1987term,geser1990relative} to DPO rewriting.
\begin{definition}[Relative termination]
    \label{termination:def:relative_termination}
     Let $\mathcal{A}$ and $\mathcal{B}$ be sets of rewriting rules and let $\mathfrak{F}$ be a DPO rewriting framework. We say that $\Rightarrow_{\mathcal{A},\mathfrak{F}}$ is \textbf{terminating relative to} $\Rightarrow_{\mathcal{B}, \mathfrak{F}}$ if any infinite $(\mathcal{A} \cup \mathcal{B},\mathfrak{F})$-rewriting sequence has only a finite number of rewriting steps with rules in $\mathcal{A}$.
\end{definition}

\begin{example}
    Consider the rewriting system with two rules.
    
     
    \begin{figure}[hbtp]
     $\alpha$ = \resizebox{0.85\textwidth}{!}{ 
                \begin{tikzpicture}[baseline=-3ex]
                    \graphbox{\( L \)}{0mm}{-3mm}{34mm}{15mm}{2mm}{2mm}{
                        \coordinate (o) at (0mm,-11mm); 
                        \node[draw,circle] (l1) at ($(o)+(-10mm,0mm)$) {1};
                        \node[draw,circle] (l2) at ($(l1)+(2,0)$) {2};
                        % \node[draw,circle] (l3) at ($(l1) + (1,0)$) {3};
                        \draw[->] (l1) -- (l2) node[midway,above] {a};
                    } 
            
                    \graphbox{\( K \)}{40mm}{-3mm}{34mm}{15mm}{2mm}{2mm}{
                        \coordinate (o) at (0mm,-11mm); 
                        \node[draw,circle] (l1) at ($(o)+(-10mm,0mm)$) {1};
                        \node[draw,circle] (l2) at ($(l1)+(2,0)$) {2};
                    }  
            
                    \graphbox{\( R \)}{80mm}{-3mm}{35mm}{15mm}{2mm}{2mm}{
                        \coordinate (o) at (-5mm,-11mm); 
                        \node[draw,circle] (l1) at ($(o)+(-10mm,0mm)$) {1};
                        % \node[draw,circle] (l2) at ($(l1)+(3,0)$) {2};
                        % \node[draw,circle] (l3) at ($(l1) + (1,0)$) {4};
                        \node[draw,circle] (l4) at ($(l1) + (2,0)$) {2};
                        \draw[->] (l1) -- (l4) node[midway,above] {b};
                        % \draw[->] (l4) -- (l2) node[midway,above] {a};
                    }    
                    \node () at (37mm,-10mm) {\( \leftarrowtail \)}; % K -> L
                    \node () at (77mm,-10mm) {\( \rightarrowtail \)}; % K -> R
                \end{tikzpicture}
                }
 
     $\beta$ = \resizebox{0.85\textwidth}{!}{ 
                \begin{tikzpicture}[baseline=-3ex]
                    \graphbox{\( L \)}{0mm}{-3mm}{34mm}{15mm}{2mm}{2mm}{
                        \coordinate (o) at (0mm,-11mm); 
                        \node[draw,circle] (l1) at ($(o)+(-10mm,0mm)$) {1};
                        \node[draw,circle] (l2) at ($(l1)+(2,0)$) {2};
                        % \node[draw,circle] (l3) at ($(l1) + (1,0)$) {3};
                        \draw[->] (l1) -- (l2) node[midway,above] {a};
                    } 
            
                    \graphbox{\( K \)}{40mm}{-3mm}{34mm}{15mm}{2mm}{2mm}{
                        \coordinate (o) at (0mm,-11mm); 
                        \node[draw,circle] (l1) at ($(o)+(-10mm,0mm)$) {1};
                        \node[draw,circle] (l2) at ($(l1)+(2,0)$) {2};
                    }  
            
                    \graphbox{\( R \)}{80mm}{-3mm}{35mm}{15mm}{2mm}{2mm}{
                        \coordinate (o) at (-5mm,-11mm); 
                        \node[draw,circle] (l1) at ($(o)+(-10mm,0mm)$) {1};
                        % \node[draw,circle] (l2) at ($(l1)+(3,0)$) {2};
                        % \node[draw,circle] (l3) at ($(l1) + (1,0)$) {4};
                        \node[draw,circle] (l4) at ($(l1) + (2,0)$) {2};
                        \draw[->] (l1) -- (l4) node[midway,above] {a};
                        % \draw[->] (l4) -- (l2) node[midway,above] {a};
                    }    
                    \node () at (37mm,-10mm) {\( \leftarrowtail \)}; % K -> L
                    \node () at (77mm,-10mm) {\( \rightarrowtail \)}; % K -> R
                \end{tikzpicture}
                }
    \end{figure}
    The rewriting system is not terminating because of the rule $\beta$, but the rule $\alpha$ is terminating relative to $\beta$ because it can only be applied finitely many times in any rewriting sequence that uses $\alpha$ and $\beta$.
\end{example}
% The concept of relative termination generalizes the notion of uniform termination of DPO graph rewriting systems which is a fundamental property of algorithms because many other properties are based on it. Given a set of rules, the uniform termination means that any rewriting sequence using these rules has a finite length. Uniform termination is undecidable in general~\cite{plump1998terminationundecidable}. It is a fundamental property of rewriting systems because it is a prerequisite for reasoning about the correctness of algorithms based on these systems.
% For example, there is no point to talk about the correctness of the result of an algorithm if the algorithm does not terminate. In the context of graph rewriting, termination means that any graph can not be transformed indefinitely by applying rewriting rules. Plump has shown that termination of DPO graph rewriting systems is undecidable in general~\cite{plump1998terminationundecidable}.

\section{Dependency Graph}
In the appendices of Chapter~\ref{chap:subgraph_counting} and Chapter~\ref{chap:antipattern}, and~\textsection~\ref{antipattern:sec:termination_criterion}, we present some propositions, lemmas, and theorems.
The dependency structure among these results is clarified in~\autoref{fig:dependency_graph}.
\begin{figure}[H]
\resizebox{0.9\textwidth}{!}{

\begin{tikzpicture}

    \node[draw] (prop_pb_eq_po) at (0,12) {Prop.$\ref{prop:pb_eq_po}$};
    \node[draw] (kpcpck_pullback) at (3,12) {Lem.$\ref{kpcpck_pullback}$};
    \node[draw] (prop_graph_adhesive) at (6,12) {Prop.$\ref{prop:graph_adhesive}$};
  \node[draw] (lem_xinXcpinCrpinR) at (9,12) {Lem.$\ref{lem:xinXcpinCrpinR}$};

    \node[draw] (lem_b_c_same_img_exist_a) at (-1,10) {Lem.$\ref{lem:b_c_same_img_exist_a}$};
    \node[draw] (kpcpxrp_po) at (9,10) {Lem.$\ref{kpcpxrp_po}$};
    
    \node[draw] (subgraph_counting_lem_decomp_w_u) at (0,8) {Lem.$\ref{subgraph_counting:lem:decomp_w_u}$};
    \node[draw] (lem_g_monic) at (6,8) {Lem.$\ref{lem:g_monic}$};

   \node[draw] (lem_h_hp_diff_g_gp_diff) at (2,6) {Lem.$\ref{lem:h_hp_diff_g_gp_diff}$};
  

    \node[draw] (subgraph_counting_lem_xlnlmxrnr) at (-2,4) {Lem.$\ref{subgraph_counting:lem:xlnlmxrnr}$};
    \node[draw] (subgraph_counting_lem_w_u_l_not_geq_r_not) at (10,4) {Lem.$\ref{subgraph_counting:lem:w_u_l_not_geq_r_not}$};

    \node[draw] (subgraph_counting_lem_w_g_geq_w_h_leq) at (0,2) {Lem.$\ref{subgraph_counting:lem:w_g_geq_w_h_leq}$};
    \node[draw,red] (antipattern_lem_xgm_xhmp_xl_xr_llsss) at (9,-1) {Lem.$\ref{antipattern:lem:xgm_xhmp_xl_xr_llsss}$};
    \node[draw,red] (antipattern_lem_xh_decomp) at (10,8) {Lem.$\ref{antipattern:lem:xh_decomp}$};

    \node[draw](subgraph_counting_thm_termination_grs) at (0,0) {Thm.$\ref{subgraph_counting:thm:termination_grs}$};

    \node[draw,red] (antipattern_lem_xgm_xhmp_xl_xr) at (9, -2) {Lem.$\ref{antipattern:lem:xgm_xhmp_xl_xr}$};

        \draw[->] (prop_pb_eq_po) to (lem_b_c_same_img_exist_a);
        \draw[->] (prop_graph_adhesive) to (kpcpxrp_po);
        \draw[->] (lem_xinXcpinCrpinR) to  (lem_g_monic);
        \draw[->] (lem_b_c_same_img_exist_a) -- (lem_g_monic);
        \draw[->] (prop_pb_eq_po) to (lem_g_monic);
        
        \draw[->] (kpcpck_pullback) to (lem_h_hp_diff_g_gp_diff);
        \draw[->] (kpcpxrp_po) to  (lem_h_hp_diff_g_gp_diff);
        \draw[->] (lem_g_monic) to  (lem_h_hp_diff_g_gp_diff);
        \draw[->] (lem_xinXcpinCrpinR) -- (lem_h_hp_diff_g_gp_diff);
        \draw[->] (lem_b_c_same_img_exist_a) to (lem_h_hp_diff_g_gp_diff);

        \draw[->] (prop_pb_eq_po) to (subgraph_counting_lem_decomp_w_u);

        \draw[->] (prop_pb_eq_po)to  (subgraph_counting_lem_w_u_l_not_geq_r_not);
        \draw[->] (kpcpck_pullback) to (subgraph_counting_lem_w_u_l_not_geq_r_not);
        \draw[->] (kpcpxrp_po) to (subgraph_counting_lem_w_u_l_not_geq_r_not);
        \draw[->] (lem_g_monic) to (subgraph_counting_lem_w_u_l_not_geq_r_not);
        \draw[->] (lem_h_hp_diff_g_gp_diff) to (subgraph_counting_lem_w_u_l_not_geq_r_not);

        \draw[->] (subgraph_counting_lem_decomp_w_u) to  (subgraph_counting_lem_w_g_geq_w_h_leq);
        \draw[->] (subgraph_counting_lem_decomp_w_u) to (subgraph_counting_lem_w_g_geq_w_h_leq);
        \draw[->] (subgraph_counting_lem_w_u_l_not_geq_r_not) to (subgraph_counting_lem_w_g_geq_w_h_leq);
        \draw[->] (subgraph_counting_lem_xlnlmxrnr) to (subgraph_counting_lem_w_g_geq_w_h_leq);

        \draw[->] (subgraph_counting_lem_w_g_geq_w_h_leq) to (subgraph_counting_thm_termination_grs);

        \draw[->] (subgraph_counting_lem_decomp_w_u) to (antipattern_lem_xgm_xhmp_xl_xr);
        \draw[->] (antipattern_lem_xgm_xhmp_xl_xr_llsss) to (antipattern_lem_xgm_xhmp_xl_xr);


        % \draw[->] (prop_pb_eq_po) to (antipattern_lem_xh_decomp);
        % \draw[->] (kpcpck_pullback) to (antipattern_lem_xh_decomp);
        % \draw[->] (kpcpxrp_po) to (antipattern_lem_xh_decomp);

        \draw[->] (prop_pb_eq_po) to (antipattern_lem_xh_decomp);
        \draw[->] (kpcpck_pullback) to (antipattern_lem_xh_decomp);
        \draw[->] (kpcpxrp_po) to (antipattern_lem_xh_decomp);

        \node[draw,red] (antipattern_lem_xglnotmlp_xhlnotmrp) at (3,-2) {Lem.$\ref{antipattern:lem:xglnotmlp_xhlnotmrp}$};
        % \draw[->] (subgraph_counting_lem_w_u_l_not_geq_r_not) to (antipattern_lem_xglnotmlp_xhlnotmrp);
        \draw[->] (prop_pb_eq_po) to (antipattern_lem_xglnotmlp_xhlnotmrp);
        \draw[->] (subgraph_counting_lem_decomp_w_u) to (antipattern_lem_xglnotmlp_xhlnotmrp);
        \draw[->] (antipattern_lem_xh_decomp) to (antipattern_lem_xglnotmlp_xhlnotmrp);
        \draw[->] (lem_g_monic) to (antipattern_lem_xglnotmlp_xhlnotmrp);
        % \draw[->] (lem_h_hp_diff_g_gp_diff) to (antipattern_lem_xglnotmlp_xhlnotmrp);
        % \draw[->] (prop_graph_adhesive) to (antipattern_lem_xglnotmlp_xhlnotmrp);

        \node[draw,red] (antipattern_lem_xglnotmlnotlp_xhlnotmrnotrp) at (9, 0) {Lem.$\ref{antipattern:lem:xglnotmlnotlp_xhlnotmrnotrp}$};
        \draw[->] (subgraph_counting_lem_w_u_l_not_geq_r_not) to (antipattern_lem_xglnotmlnotlp_xhlnotmrnotrp);
        \draw[->] (lem_g_monic) to (antipattern_lem_xglnotmlnotlp_xhlnotmrnotrp);
        \draw[->] (antipattern_lem_xh_decomp) to (antipattern_lem_xglnotmlnotlp_xhlnotmrnotrp);

        \draw[->] (lem_h_hp_diff_g_gp_diff) to (antipattern_lem_xglnotmlnotlp_xhlnotmrnotrp);

        \draw[->] (prop_graph_adhesive) to (antipattern_lem_xglnotmlnotlp_xhlnotmrnotrp);
        \draw[->] (prop_pb_eq_po) to (antipattern_lem_xglnotmlnotlp_xhlnotmrnotrp);
        
        \node[draw,red] (antipattern_lem_w_g_geq_w_h_leq) at (3,-3) {Lem.$\ref{antipattern:lem:w_g_geq_w_h_leq}$};
        \draw[->] (antipattern_lem_xglnotmlp_xhlnotmrp) to (antipattern_lem_w_g_geq_w_h_leq);
        \draw[->] (antipattern_lem_xglnotmlnotlp_xhlnotmrnotrp) to (antipattern_lem_w_g_geq_w_h_leq);
        \draw[->] (antipattern_lem_xgm_xhmp_xl_xr) to (antipattern_lem_w_g_geq_w_h_leq);

        \node[draw,red] (antipattern_thm_termination_grs) at (3,-4) {Thm.$\ref{antipattern:thm:termination_grs}$};
        \draw[->] (antipattern_lem_w_g_geq_w_h_leq) to  (antipattern_thm_termination_grs);
    \end{tikzpicture} 
}
    \caption{Dependency graph}
    \label{fig:dependency_graph}
\end{figure}


%%%%%%%%%%%%%%%%%%%%%%%%%%%%%%%%%%%%%%%%%%%%%%%%%%
%%%   Chapter : Non-Well-Founded Semirings   %%%%%
%%%%%%%%%%%%%%%%%%%%%%%%%%%%%%%%%%%%%%%%%%%%%%%%%%
\chapter{Termination of Graph Rewriting using Weighted Type Graphs over Non-Well-Founded Semirings} 
\label{chap:nwf}

We extend an existing method for termination of DPO graph rewriting systems, called type graph method~\footnote{The result has appeared in the workshop paper~\cite{qiu2025termination_nwf_v2_acceptedgcm}.}.

The weighted type graph method was first introduced by Hans, Konig and Bruggink~\cite{zantema2014termination} for cycle rewriting. Their subsequent work~\cite{bruggink2014termination} generalized it for DPO rewriting on edge-labeled directed multigraphs with injective rules and injective matches; later, it was extended to general DPO rewriting on edge-labeled multigraphs by Bruggink et al.~\cite{bruggink2015proving}; and it was adapted to broader categories and DPO variants by Endrullis and Overbeek~\cite{endrullis2024generalized_icgt}. 

The method assigns weights to morphisms targeting a weighted type graph over a well-founded semiring. The weight of a graph is defined
 as the sum of the weights of all morphisms from that graph to the type graph. Relative termination of $\mathcal{A}$ with respect to $\mathcal{B}$ is proven by ensuring that rewriting steps using rules in \( \mathcal{A} \) strictly decrease the weights of the host graphs, while rewriting steps using rules in \( \mathcal{B} \) do not increase them.
    
   Previous work~\cite{zantema2014termination,bruggink2014termination,bruggink2015proving} proposed three concrete semirings (on natural numbers): the natural tropical semiring $\mathfrak{T}$, the natural arctic semiring $\mathfrak{A}$, and the natural arithmetic semiring $\mathfrak{N}$.
    However, constructing weighted type graphs over these concrete semirings for DPO rewriting on edge-labeled directed multigraphs is difficult in general, because it requires quantifying over all edge-labeled directed multigraphs with weighted edges over $\mathbb{N}$.

    In response to this challenge, these works proposed to reduce the search space by fixing the number of nodes \( k\) and the maximum weight of edges of the weighted type graph, and then by constructing a weighted type graph as follows: 
    \begin{enumerate}
      \item Construct a graph with \( k \) nodes, with a directed and labeled edge per ordered pair of nodes and label from the finite set of edge labels $\Sigma$ of the rewriting system. 
      \item For each edge, decide if it exists in the weighted type graph;  if an edge exists, assign a weight (a natural number) to it.
      \item Check if the weighted type graph fulfils the required conditions.
    \end{enumerate}
    where $\Sigma$ is a finite set of edge labels that will be fixed throughout this chapter.
    Despite these constraints, searching for suitable weighted type graphs
    %  that witness termination of a DPO rewriting system on edge-labeled directed multigraphs
    remains challenging. 
    Let $n = k^2 \cdot | \Sigma |$:
    
    For a weighted type graph over the natural tropical semiring $\mathfrak{T}$ or the natural arctic semiring $\mathfrak{A}$, the problem amounts to
    checking the satisfiability of an existential Presburger arithmetic formula with $n$ binary variables and $n$ integer variables.
    While modern SMT solvers, such as Z3~\cite{de2008z3} (which incorporates the dedicated CutSat solver~\cite{z3ilp_cutsat}), can solve practical instances of this problem, its worst-case complexity remains exponential \( O(2^{2n}) \)~\cite{arithmetic2024z3}.

     For a weighted type graph over the natural arithmetic semiring $\mathfrak{N}$, the task amounts to
     checking the satisfiability of an existential Peano arithmetic formula with addition and multiplication, involving $n$ binary variables and $n$ integer variables. Though modern solvers like Z3 can tackle practical instances, it is a semi-decidable problem~\cite{matiyasevivc2003enumerable}.

    There are automated tools that implement the weighted type graph method for proving termination of DPO rewriting systems on edge-labeled directed multigraphs:  
    \texttt{TORPAcyc}~\cite{TORPAcyc} implements the weighted type graph method for cycle rewriting and uses the SMT solver \texttt{Yices}~\cite{yices} to solve constraint systems;
    \texttt{Grez}~\cite{grez} implements the weighted type graph method for DPO rewriting on edge-labeled directed multigraphs; \texttt{GraphTT-wtg}~\footnote{To the best of our knowledge, no implementation is publicly available.}~\cite{endrullis2024generalized_arxiv_v3} implements the weighted type graph method for DPO rewriting on many categories.
    \texttt{Grez} and \texttt{GraphTT-wtg}  
    use  the SMT solver \texttt{Z3} to solve constraint systems.

    In spite of the theoretical power of the weighted type graph method, its practical applicability is limited by the high computational complexity of searching for suitable weighted type graphs over well-founded semirings on 
    natural numbers. In fact, to reduce the search space, users of the above-mentioned tools are required
     to fix the size of the weighted type graph, as well as the maximum weight of edges in the weighted type graph. These are challenging tasks if not impossible to determine a priori even for experts in the field. 
    Furthermore, our experiments with the two publicly available ones \texttt{TORPAcyc} and \texttt{Grez} show that when the number of edge labels is $2$, they struggle to search for weighted type graphs with $4$ nodes and maximum edge weight of $2$, or to search for weighted type graphs with $3$ nodes when the maximum edge weight is larger than $3$~\footnote{Experiments were conducted on a laptop equipped with an i5-1038NG7 CPU, which features 4 cores, a base clock speed of 2 GHz, and a boost speed of 3.8 GHz.}. 
   Chapter~\ref{chap:nwf} addresses this limitation to make the weighted type graph method more accessible to non-expert users and to enhance its practical applicability.

To reduce the complexity of searching for weighted type graphs that witness termination of a DPO rewriting system on edge-labeled directed multigraphs, we propose extending the weighted type graph method to non-well-founded semirings and introduce three concrete semirings over the real numbers: the real tropical semiring $\mathfrak{T}'$, the real arctic semiring $\mathfrak{A}'$, and the real arithmetic semiring $\mathfrak{N}'$. The idea of using weights from non-well-founded domains to prove termination was used in the context of term rewriting by Lucas~\cite{lucas2006relative}.

Under the same constraints, searching for weighted type graphs over these semirings is computationally more tractable.
%  Specifically,
% for the real tropical semiring $\mathfrak{T}'$ or the real arctic semiring $\mathfrak{A}'$, the problem amounts to checking the satisfiability of an
% existential linear real arithmetic formula with $n$ binary variables and $n$ real variables for the real tropical semiring $\mathfrak{T}'$ or the real arctic semiring $\mathfrak{A}'$.
Z3 can solve practical instances of this problem more easily: Decidable for the real arithmetic semiring $\mathfrak{N}'$;
%  using CutSat solver, because there are only $n$ integer variables.
The worst-case complexity for solving this problem is $O(2^n)$ for the real tropical semiring $\mathfrak{T}'$ or the real arctic semiring $\mathfrak{A}'$;
% For a weighted type graph over the real arithmetic semiring $\mathfrak{N}'$, the task reduces to checking the satisfiability of
% an existential non-linear real arithmetic formula\textemdash a task that is decidable~\cite{collins1974quantifier,z3realarithmetic}.

The implementation of both our approach and the approach proposed by Endrullis et al.~\cite{endrullis2024generalized_icgt}, written in OCaml, will be presented in Chapter~\ref{chap:lyonparallel}.

We first introduce the type graph method~\cite{endrullis2024generalized_icgt} and then present our extension to non-well-founded semirings.

\section{Preliminaries}

\subsection{Well-Founded Semirings} 
\label{sec:well_founded_semiring}

\begin{definition}[Monoid]
    A \textbf{monoid} $(S, *, e)$ is a set $S$ equipped with a binary operation $* : S \times S \rightarrow S$ and having a spacial element $e$ such that:
    \begin{itemize}
        \item {$*$ is associative:} $ (x * y) * z = x * (y * z)$
        \item {$e$ is an identity element:} $  e * x = x = x * e$
    \end{itemize}
    A \textbf{commutative monoid} is a monoid where $*$ is commutative: $x * y = y * x$.
\end{definition}

\begin{definition}[Semiring]
    \label{def_semiring}
        A \textbf{semiring} is a 5-tuple $(S,\bigoplus,\bigotimes,0,1)$ satisfying the following conditions:
         \begin{itemize}
            \item {$(S,\bigoplus,0)$ is a commutative monoid} 
            \item {$(S,\bigotimes,1)$ is a monoid} 
            \item {$\bigotimes$ is distributive over $\bigoplus$:}
            \begin{itemize}
                \item $a \bigotimes ( b \bigoplus c) = (a \bigotimes b) \bigoplus (a \bigotimes c)$ 
                \item $(b \bigoplus c) \bigotimes a = (b \bigotimes a) \bigoplus (c \bigotimes a)$
            \end{itemize} 
            \item {0 is an annihilator for $\bigotimes$:} $x \bigotimes 0= 0 =  0 \bigotimes x $
        \end{itemize}
        A \textbf{commutative semiring} is a semiring where $\bigotimes$ is commutative.
\end{definition}

\begin{definition}[Well-founded Semiring]
    \label{def:well_founded_semiring}
    A \textbf{well-founded semiring} $(S, \oplus, \odot, 0, 1,\prec, \leq)$ consists of
    \begin{itemize}
        \item A semiring $(S, \oplus, \odot, 0, 1)$,
        \item non-empty orders $\prec, \leq \subseteq S \times S$,
    \end{itemize}
    such that the equivalent class $\prec / \leq$ is well-founded, $0 \neq 1$ and for all $x,y,z,w \in S$ we have
        \begin{align*}
            x \preceq x' \land y \preceq y' 
            &\Rightarrow
            x \oplus y \preceq x' \oplus y'
            &\tag{S1} \label{wfs:ax:s1} 
            \\   
            % x < y  
            % &\Rightarrow
            % x \oplus z \leq y \oplus z 
            % \tag{S1} \label{eq:ordered_semiring_plus_monotonic} 
            % \\ w
            x \prec x' \land y \prec y'  
            &\Rightarrow
            x \oplus y \prec x' \oplus y'
            &\tag{S2} \label{wfs:ax:s2} 
            \\
            x \preceq x' \land 1 \leq y
            &\Rightarrow 
            x \odot y \preceq x' \odot y \land y \odot x \preceq y \odot x' 
            &\tag{S3} \label{wfs:ax:s3} 
            \\
            x \prec x' \land 1 \leq y \neq 0 
            &\Rightarrow
            x \odot y \prec x' \odot y \land y \odot x \prec y \odot x'
            &\tag{S4} \label{wfs:ax:s4}
        \end{align*}
      The semiring is \textbf{strictly monotonic} if it additionally satisfies 
        \begin{flalign*}
                \hspace{5cm}x \prec x'  \land y \leq y'
                &\Rightarrow
                x \oplus y \prec x' \oplus y 
                &\tag{S8} \label{ax:s5} 
        \end{flalign*}
\end{definition}
% \begin{example} 
%     \label{example:real_semirings}
%     The natural tropical semiring: $\mathfrak{T} = (\mathbb{N} \cup \{+\infty\},\min,+,+\infty, 0, <, \leq)$
%     and the natural arctic semiring: $\mathfrak{A} = (\mathbb{N} \cup \{-\infty\},\max,+,-\infty, 0,<, \leq)$ are well-founded semirings.
%     The natural arithmetic semiring $\mathfrak{N} = (\mathbb{N},+,*,0,1,<, \leq)$ is a strictly monotonic.
% \end{example}
\begin{example}
    The natural tropical semiring $\mathfrak{T} = (\mathbb{N} \cup \{+\infty\},\min,+,+\infty, 0, <, \leq)$ has domain $\mathbb{N} \cup \{+\infty\}$, the binary function symbol $\oplus$ interpreted by $\min$ and the binary function symbol $\odot$ interpreted by $+$, the constant symbols $0_s$ and $1_s$ interpreted by $+\infty$ and $0_\mathbb{N}$, respectively, the binary relation symbols $\prec$ and $\leq$ interpreted by the canonical \enquote{strictly less than} and \enquote{less than or equal to} on $\mathbb{N} \cup \{+\infty\}$. It is a well-founded semiring. It is not strictly monotonic because $2 < 3$ but $2 \oplus 2 = \min(2,2) = 2 \not < 2 = \min(3,2) = 3 \oplus 2$.
\end{example}
    
\begin{example}
    The natural arctic semiring $\mathfrak{A} = (\mathbb{N} \cup \{-\infty\},\max,+,-\infty, 0,<,\leq)$ has domain $\mathbb{N} \cup \{-\infty\}$, the binary function symbol $\oplus$ interpreted by $\max$ and the binary function symbol $\odot$ interpreted by $+$, the constant symbols $0_s$ and $1_s$ interpreted by $-\infty$ and $0_\mathbb{N}$, respectively, the binary relation symbols $\prec$ and $\leq$ interpreted by the canonical \enquote{strictly less than} and \enquote{less than or equal to} on $\mathbb{N} \cup \{-\infty\}$. It is a well-founded semiring. It is not strictly monotonic because $2 < 3$ but $2 \oplus 3 = \max(2,3) = 3 \not < 3 = \max(3,3) = 3 \oplus 3$.
\end{example}

\begin{example}
    The natural arithmetic semiring $\mathfrak{N} = (\mathbb{N},+,*,0,1,<,\leq)$ has domain $\mathbb{N}$, the binary function symbol $\oplus$ interpreted by $+$ and the binary function symbol $\odot$ interpreted by $*$, the constant symbols $0_s$ and $1_s$ interpreted by $0_\mathbb{N}$ and $1_\mathbb{N}$, respectively, the binary relation symbols $\prec$ and $\leq$ interpreted by the canonical \enquote{strictly less than} and \enquote{less than or equal to} on $\mathbb{N}$. It is a strictly monotonic semiring.
\end{example}

\begin{notation} 
    \label{def:bigodot}
Let $(S, \oplus, \odot, 0_s, 1_s)$ be a semiring. We extend naturally the binary operations $\oplus$ and $\odot$ to finite sets $E \subseteq S$ by letting
    \begin{itemize}
        \item $\bigodot \emptyset \overset{\operatorname{def}}{=} 1_s$ and $\bigodot \left( E \cup \{x\} \right) \overset{\operatorname{def}}{=} \left( \bigodot E \right) \odot x$;
        \item $\bigoplus \emptyset \overset{\operatorname{def}}{=} 0_s$ and $\bigoplus \left( E \cup \{x\} \right) \overset{\operatorname{def}}{=} \left( \bigoplus E \right) \oplus x$.
    \end{itemize}
%  \begin{flalign*}
%     \bigodot \emptyset &\overset{\operatorname{def}}{=} 1_s
% \\
%     \bigodot \left( E \cup \{x\} \right) &\overset{\operatorname{def}}{=} \left( \bigodot E \right) \odot x
%     \\
%     \bigoplus \emptyset &\overset{\operatorname{def}}{=} 0_s
%     \\
%         \bigoplus \left( E \cup \{x\} \right) &\overset{\operatorname{def}}{=} \left( \bigoplus E \right) \oplus x
% \end{flalign*}
\end{notation}

% \begin{notation}
%     Let $S$ be a strongly monotonic measurable semiring. We write $0_s$, $1_s$ and $\mu_s$ to denote the additive and multiplicative neutral elements of $S$ and the homomorphism, respectively.
% \end{notation}
 
\subsection{Weighted Type Graph} 
\label{sec:weighted_type_graph}
Throughout this chapter, let \(\mathcal{C}\) be a fixed locally small category. A weighted type graph is an object \(T\) in \(\mathcal{C}\), called the type graph, equipped with a distinguished family of morphisms into \(T\), each endowed with a weight. 
A \enquote{type graph} is an object, following the terminology from the previous work by Bruggink~\cite{bruggink2014termination}.
 Given a weighted type graph and any object \(X\) in \(\mathcal{C}\), we interpret \(X\) as \(\operatorname{Hom}(X,T)\). The distinguished morphisms determine weight assignments on the elements of \(\operatorname{Hom}(X,T)\); the weight of \(X\) is then defined as the sum of the weights of the morphisms in \(\operatorname{Hom}(X,T)\). These assignments provide a basis for comparing objects of \(\mathcal{C}\) according to their weights.

\begin{definition}[\cite{endrullis2024generalized_arxiv_v2}]
    \label{wf:def:weighted_type_graph}
    A \textbf{weighted type graph} is a tuple~\(\mathcal{T} = (T, \mathbb{E}, \mathcal{S}, w)\) where
    \begin{itemize} 
        \item \(T\) is an object of $\mathcal{C}$, called \textbf{type graph},
        \item \(\mathbb{E}\) is a set of morphisms in $\mathcal{C}$ with codomain $T$, and the elements of \(\mathbb{E}\) are called \textbf{morphism-rulers},
        \item \(\mathcal{S}=(S, \oplus, \odot, 0, 1)\) is a commutative semiring,
        \item \(w : \mathbb{E} \to S \setminus \{0\}\) is a weight function such that for all $e \in \mathbb{E}, w(e) \geq 1$,
    \end{itemize}
    % The weighted type graph \(\mathcal{T}\) is \textbf{finitary} if for every \( (e :X \to T) \in \mathbb{E}\) and every object \(G\), \(\operatorname{Hom}(X, G)\) and \(\operatorname{Hom}(G, T)\) are finite sets.
    such that for every \( (e :X \to T) \in \mathbb{E}\) and every object \(G\), \(\operatorname{Hom}(X, G)\) and \(\operatorname{Hom}(G, T)\) are finite sets.
\end{definition}

\begin{example}
    \label{wf:example:weighted_type_graph}
     Consider the weighted type graph $\mathcal{T} = (T, \mathbb{E}, \mathcal{S}, w)$ where
    \begin{itemize}
        \item $T$ is the labeled graph shown below:
        % in Figure~\ref{fig:weighted_type_graph_instance_dfjlsjl},
        % \begin{figure}[H]
        %     \centering
        \begin{center}
            \begin{tikzpicture}
                    \graphbox{}{0mm}{0mm}{32mm}{28mm}{-10mm}{-14mm}{
                        \node[draw,circle] (1) at (0,0) {};
                        \node[draw,circle] (2) at (2,0) {};
                        \draw[->] (1) edge[loop above] node[midway, above] {$a$} (1) ;
                        \draw[->] (1) edge[loop below] node[midway, below] {$b$} (1) ;
                        \draw[->] (1) edge[bend left] node[midway, above] {$a$}  (2)  ;
                        \draw[->] (2) edge[bend left] node[midway, below] {$a$} (1)   ;
                    }
                \end{tikzpicture}
        \end{center}
            %     \caption{}
            %     \label{fig:weighted_type_graph_instance_dfjlsjl}
            % \end{figure}
        \item $\mathcal{S}$ is the natural arithmetic semiring $(\mathbb{N}, +_\mathbb{N}, *_\mathbb{N}, 0_\mathbb{N}, 1_\mathbb{N}, <_\mathbb{N},\leq_\mathbb{N})$,
        \item $\mathbb{E}$ consists of 4 morphism-rulers $e_{11a}, e_{12a}, e_{21a}, e_{11b}$ shown
            below:
        %  in Figure~\ref{fig:morphism_ruler__dfjsdlfjsadl}.
            % \begin{figure}[H]
            %    \centering
            \begin{center}
        \resizebox{0.69\textwidth}{!}{
            \begin{tikzpicture}
            \graphbox{\( K \)}{-50mm}{0mm}{40mm}{30mm}{2mm}{-6mm}{
                \coordinate (o) at (0mm,-10mm); 
                    \node[draw,circle] (l1) at ($(o)+(-10mm,0mm)$) {$1$};
                    \node[draw,circle] (l2) at ($(l1)+(2,0)$) {$2$};
                    \draw[->] (l1) -- (l2) node[midway,above] {$a$};
                } 
                \graphbox{$T$}{0mm}{0mm}{40mm}{30mm}{-10mm}{-15mm}{
                    \node[draw,circle] (1) at (0,0) {$1\ 2$};
                    \node[draw,circle] (2) at (2,0) {};
                    \draw[->] (1) edge[loop above] node[midway, above] {$a$} (1) ;
                    \draw[->] (1) edge[loop below] node[midway, below] {$b$} (1) ;
                    \draw[->] (1) edge[bend left] node[midway, above] {$a$}  (2)  ;
                    \draw[->] (2) edge[bend left] node[midway, below] {$a$} (1)   ;
                }
                \node () at (-5mm,-15mm) {$\overset{e_{11a}}{\to}$};
            \end{tikzpicture}
            } 

            \resizebox{0.69\textwidth}{!}{
            \begin{tikzpicture}
                \graphbox{\( K \)}{-50mm}{0mm}{40mm}{28mm}{2mm}{-6mm}{
                    \coordinate (o) at (0mm,-10mm); 
                    \node[draw,circle] (l1) at ($(o)+(-10mm,0mm)$) {$1$};
                    \node[draw,circle] (l2) at ($(l1)+(2,0)$) {$2$};
                    \draw[->] (l1) -- (l2) node[midway,above] {$b$};
                } 
                \graphbox{$T$}{0mm}{0mm}{40mm}{28mm}{-10mm}{-15mm}{
                    \node[draw,circle] (1) at (0,0) {$1\ 2$};
                    \node[draw,circle] (2) at (2,0) {};
                    \draw[->] (1) edge[loop above] node[midway, above] {$a$} (1) ;
                    \draw[->] (1) edge[loop below] node[midway, below] {$b$} (1) ;
                    \draw[->] (1) edge[bend left] node[midway, above] {$a$}  (2)  ;
                    \draw[->] (2) edge[bend left] node[midway, below] {$a$} (1)   ;
                }
                \node () at (-5mm,-15mm) {$\overset{e_{11b}}{\to}$};
            \end{tikzpicture}
            }

            \resizebox{0.69\textwidth}{!}{
                \begin{tikzpicture}
                    \graphbox{\( K \)}{-50mm}{0mm}{40mm}{28mm}{2mm}{-6mm}{
                        \coordinate (o) at (0mm,-10mm); 
                        \node[draw,circle] (l1) at ($(o)+(-10mm,0mm)$) {$1$};
                        \node[draw,circle] (l2) at ($(l1)+(2,0)$) {$2$};
                        \draw[->] (l1) -- (l2) node[midway,above] {$a$};
                    } 
                    \graphbox{$T$}{0mm}{0mm}{40mm}{26mm}{-10mm}{-15mm}{
                        \node[draw,circle] (1) at (0,0) {$1$};
                        \node[draw,circle] (2) at (2,0) {$2$};
                        \draw[->] (1) edge[loop above] node[midway, above] {$a$} (1) ;
                        \draw[->] (1) edge[loop below] node[midway, below] {$b$} (1) ;
                        \draw[->] (1) edge[bend left] node[midway, above] {$a$}  (2)  ;
                        \draw[->] (2) edge[bend left] node[midway, below] {$a$} (1)   ;
                    }
                    \node () at (-5mm,-15mm) {$\overset{e_{12a}}{\to}$};
                \end{tikzpicture}
            }

            \resizebox{0.69\textwidth}{!}{
            \begin{tikzpicture}
                \graphbox{\( K \)}{-50mm}{0mm}{40mm}{28mm}{2mm}{-6mm}{
                    \coordinate (o) at (0mm,-10mm); 
                    \node[draw,circle] (l1) at ($(o)+(-10mm,0mm)$) {$1$};
                    \node[draw,circle] (l2) at ($(l1)+(2,0)$) {$2$};
                    \draw[->] (l1) -- (l2) node[midway,above] {$a$};
                } 
                \graphbox{$T$}{0mm}{0mm}{40mm}{28mm}{-10mm}{-15mm}{
                    \node[draw,circle] (1) at (0,0) {2};
                    \node[draw,circle] (2) at (2,0) {1};
                    \draw[->] (1) edge[loop above] node[midway, above] {$a$} (1) ;
                    \draw[->] (1) edge[loop below] node[midway, below] {$b$} (1) ;
                    \draw[->] (1) edge[bend left] node[midway, above] {$a$}  (2)  ;
                    \draw[->] (2) edge[bend left] node[midway, below] {$a$} (1)   ;
                }
                \node () at (-5mm,-15mm) {$\overset{e_{21a}}{\to}$};
            \end{tikzpicture}
            }
        \end{center}
            %     \caption{}
            %     \label{fig:morphism_ruler__dfjsdlfjsadl}
            % \end{figure}
        \item $w(e) = 1_\mathbb{N}$ for all $e \in \mathbb{E}$.
         \end{itemize}
    %  with 
% $\mathbb{E}=\{e_{11a},e_{12a},e_{21a},e_{11b}\}$ as the set of morphism-rulers where 
%      $e_{uvl}$, for all edge from $u$ to $v$ labeled by $l$, has domain 
%      \tikz[baseline=-0.5ex]{
%         \node[draw,circle] (x) at (0,0) {};
%         \node[draw,circle] (y) at (1,0) {};
%         \draw[->] (x) -- (y) node[midway, above] {$l$};
%     } and image 
%     \begin{tikzpicture}[baseline=-2mm]
%         \node[draw, circle] (x) at (0,0) {$\mathrm{u}$};
%         \node[draw, circle] (y) at (1,0) {$\mathrm{v}$};
%         \draw[->]  (x) -- (y) node [midway,above] {$l$};
%     \end{tikzpicture} in the graph $T$,
    % and $w(e) = 1_\mathbb{N}$ for all $e \in \mathbb{E}$, is visualized in Figure~\ref{fig:weighted_type_graph_instance}. 
    To improve readability, this weighted type graph will visualized as shown below, where a label $l$ with superscript $n$ on an edge means that there is a morphism-ruler with weight $n$ from \raisebox{2pt}{
            \scalebox{0.7}{\tikz[baseline=-0.5ex]{
            \node [draw,circle] (x) at (0,0) {};
            \node[draw,circle] (y) at (1,0) {};
            \draw[->] (x)--(y) node[midway, above] {$l$};
        }}} to this edge in $\mathbb{E}$.
    % \begin{figure}[H]
    %     \centering
    \begin{center}
        \begin{tikzpicture}
            \graphbox{}{0mm}{0mm}{32mm}{28mm}{-10mm}{-14mm}{
                \node[draw,circle] (1) at (0,0) {1};
                \node[draw,circle] (2) at (2,0) {2};
                \draw[->] (1) edge[loop above] node[midway, above] {$a^{1}$} (1) ;
                \draw[->] (1) edge[loop below] node[midway, below] {$b^{1}$} (1) ;
                \draw[->] (1) edge[bend left] node[midway, above] {$a^{1}$}  (2)  ;
                \draw[->] (2) edge[bend left] node[midway, below] {$a^{1}$} (1)   ;
            }
        \end{tikzpicture}
    \end{center}
    %     \caption{}
    %     \label{fig:weighted_type_graph_instance}
    % \end{figure}
\end{example}

\subsection{Measuring objects by counting morphisms}
\label{sec:type_graph:wf:measuring_graphs}
We recall the definition of morphism weight and object weight relative to a type graph~\cite{endrullis2024generalized_icgt}.

\begin{notation}[\cite{endrullis2024generalized_arxiv_v2}]
Let \(A,B,C\) be objects. For morphism $\alpha\mathop{\colon} A \mathop{\to} B$, we introduce the notation 
          $$\set{ \alpha \mathop{\star} - \mathop{=} \gamma } \overset{\operatorname{def}}{=} \{ \beta \mathop{\in} \operatorname{Hom}(B, C) \mid \alpha \mathop{\star} \beta \mathop{=} \gamma \}.$$ 
For morphism $\beta\mathop{\colon} B \mathop{\to} C$, we introduce the notation
          $$\set{ - \mathop{\star} \beta \mathop{=} \gamma }  \overset{\operatorname{def}}{=} \{ \alpha \mathop{\in} \operatorname{Hom}(A, B) \mid \alpha \mathop{\star} \beta \mathop{=} \gamma \}.$$
For morphism set $E \mathop{\subseteq} \operatorname{Hom}(A,B)$ and morphism $\beta\mathop{\colon} B \mathop{\to} C$, we introduce the notation
          $$E \mathop{\star} \beta \overset{\operatorname{def}}{=} \set{ \alpha \mathop{\star} \beta \mid \alpha \mathop{\in} E }.$$
\end{notation}
Analogous to measuring a physical object, a morphism ruler measures a morphism as follows.
\begin{definition} 
    \label{def:measurement_of_a_morphism_relative_to_a_morphism_ruler}
    The \emph{measurement} of a morphism \( h:G \mathop{\to} T \) relative to a morphism-ruler \( e: X \mathop{\to} T \), denoted by $m_e(h)$, is defined as:
                \(
                m_e(h) 
                    \overset{\operatorname{def}}{=}
                \card{\{- \mathop{\star} h \mathop{=} e\}}.
                \)
\end{definition}


Combining the measurements of a morphism provided by all morphism-rulers with the help of a weight function $w$ gives the weight of the morphism (Definition~\ref{def:weight_of_a_morphism_relative_to_a_type_graph}).
\begin{definition} 
    \label{def:weight_of_a_morphism_relative_to_a_type_graph}
        Let $\mathcal{T}=(T,\mathbb{E},\mathcal{S},w)$ be a finitary weighted type graph.
         The \textbf{weight of a morphism $h: G \mathop{\rightarrow} T$ relative to a type graph $\mathcal{T}$} is defined as the semiring product of $w(e)^{m_e(h)}$ for all $e \mathop{\in} \mathbb{E}$:
        \[  w_{\mathcal{T}}(h) \overset{\operatorname{def}}{=} \underset{e \mathop{\in} \mathbb{E}}{\mathop{\bigodot}} 
                w(e)^{m_e(h)}.\]
\end{definition}
Finally, the weight of an object is defined as the semiring sum of the weights of all morphisms from the object to the underlying type graph $T$ of the weighted type graph $\mathcal{T}$.
\begin{definition}[Object weight]
    \label{def:weight_of_an_object_relative_to_a_type_graph}
       Let $\mathcal{T}=(T,\mathbb{E},\mathcal{S},w)$ be a finitary weighted type graph. The \textbf{weight of an object \( G \)} is defined as the semiring sum of $w_\mathcal{T}(h)$, for all \( h \mathop{\in} \operatorname{Hom}(G,T) \):
        \[ w_\mathcal{T}(G) \overset{\operatorname{def}}{=} \underset{h \mathop{\in} \operatorname{Hom}(G,T)}{\mathop{\bigoplus}}  w_\mathcal{T}(h). \]
\end{definition}

\begin{example}
    Consider the two morphisms shown below
    %  Figure~\ref{fig:example:two_weighted_type_graph_morphisms}
    % \begin{figure}[H]
    %     \centering
    \begin{center}
        \resizebox{0.7\textwidth}{!}{
        \begin{tikzpicture}
          \graphbox{\( L \)}{-50mm}{0mm}{40mm}{39mm}{2mm}{-6mm}{
            \coordinate (o) at (0mm,-10mm); 
            \node[draw,circle] (l1) at ($(o)+(-10mm,0mm)$) {1};
            \node[draw,circle] (l2) at ($(l1)+(2,0)$) {2};
            \node[draw,circle] (l3) at ($(l1)+(1,0)$) {3};
            \draw[] (l1) -- (l3) node[midway,above] {$a$};
            \draw[] (l3) -- (l2) node[midway,above] {$a$};
        } 
            \graphbox{$T$}{0mm}{0mm}{40mm}{39mm}{-10mm}{-17mm}{
                \node[draw,circle] (1) at (0,0) {$1\ 2$};
                \node[draw,circle] (2) at (2,0) {3};
                \draw[->] (1) edge[loop above] node[midway, above] {$a^{1}$} (1) ;
                \draw[->] (1) edge[loop below] node[midway, below] {$b^{1}$} (1) ;
                \draw[->] (1) edge[bend left] node[midway, above] {$a^{1}$}  (2)  ;
                \draw[->] (2) edge[bend left] node[midway, below] {$a^{1}$} (1)   ;
            }
            \node () at (-5mm,-15mm) {$\overset{h_{11}^1}{\to}$};
        \end{tikzpicture}
        }

        \resizebox{0.7\textwidth}{!}{
            \begin{tikzpicture}
              \graphbox{\(L\)}{-50mm}{0mm}{40mm}{39mm}{2mm}{-6mm}{
                \coordinate (o) at (0mm,-10mm); 
                \node[draw,circle] (l1) at ($(o)+(-10mm,0mm)$) {1};
                \node[draw,circle] (l2) at ($(l1)+(2,0)$) {2};
                \node[draw,circle] (l3) at ($(l1)+(1,0)$) {3};
                \draw[] (l1) -- (l3) node[midway,above] {$a$};
                \draw[] (l3) -- (l2) node[midway,above] {$a$};
            } 
                \graphbox{$T$}{0mm}{0mm}{40mm}{39mm}{-10mm}{-19mm}{
                    \node[draw,circle] (1) at (0,0) {$1\ 2\ 3$};
                    \node[draw,circle] (2) at (2,0) {};
                    \draw[->] (1) edge[loop above] node[midway, above] {$a^{1}$} (1) ;
                    \draw[->] (1) edge[loop below] node[midway, below] {$b^{1}$} (1) ;
                    \draw[->] (1) edge[bend left] node[midway, above] {$a^{1}$}  (2)  ;
                    \draw[->] (2) edge[bend left] node[midway, below] {$a^{1}$} (1)   ;(1)   ;
                }
                \node () at (-5mm,-15mm) {$\overset{h_{11}^2}{\to}$};
            \end{tikzpicture}
            }
    \end{center}
    %         \caption{}
    %         \label{fig:example:two_weighted_type_graph_morphisms}
    %   \end{figure}
       and the weighted type graph shown 
       below 
    %    in Figure~\ref{fig:weighted_type_graph_instance_sfsdsfs}.
    % \begin{figure}[H]
    %     \centering
    \begin{center}
        \begin{tikzpicture}
            \graphbox{}{0mm}{0mm}{32mm}{28mm}{-10mm}{-14mm}{
                \node[draw,circle] (1) at (0,0) {1};
                \node[draw,circle] (2) at (2,0) {2};
                \draw[->] (1) edge[loop above] node[midway, above] {$a^{1}$} (1) ;
                \draw[->] (1) edge[loop below] node[midway, below] {$b^{1}$} (1) ;
                \draw[->] (1) edge[bend left] node[midway, above] {$a^{1}$}  (2)  ;
                \draw[->] (2) edge[bend left] node[midway, below] {$a^{1}$} (1)   ;
            }
        \end{tikzpicture}
    %     \caption{}
    %     \label{fig:weighted_type_graph_instance_sfsdsfs}
    % \end{figure}
    \end{center}
    We have \begin{itemize}
        \item $ w_\mathcal{T}(h_{11}^1) \mathop{=} w(e_{13a})^{m_{e_{13a}}(h_{11}^1)} \mathop{\odot} w(e_{31a})^{m_{e_{31a}}(h_{11}^1)} =
     1^1 * 1^1 \mathop{=} 1$,
        \item $
        w_\mathcal{T}(h_{11}^2) 
        % \mathop{=} 1^1 * 1^1 
        \mathop{=} 1$.
    \end{itemize}
\end{example}   

  
\subsection{Estimating weights of pushout objects} 
\label{sec:type_graph:wf:weighing_pushout} 
% Let \( \mathcal{R} \mathop{=} \mathcal{A} \mathop{\cup} \mathcal{B} \) be a set of DPO rewriting rules and let $\mathcal{T}=(T,\mathbb{E},\mathcal{S},w)$ be a weighted type graph over a well-founded commutative semiring such that for every graph $G$ subject to rewriting, we have $\operatorname{Hom}(G,T)\mathop{\neq} \emptyset$. By~\autoref{rem:wf:weight_of_object_geq_1}, for all objects \( G \) subject to rewriting, \(w_\mathcal{T}(G) \mathop{\succeq}_S 1_\mathcal{S} \).
% The rewriting relation \( \mathop{\Rightarrow}_{\mathcal{A},\mathfrak{F}} \) is terminating relative to $\mathop{\Rightarrow}_{\mathcal{B},\mathfrak{F}}$ if (i) for all \(G \mathop{\Rightarrow}_{\mathcal{A},\mathfrak{F}} H\), \( w_\mathcal{T}(G) \mathop{\succ}_S w_\mathcal{T}(H)\), and (ii) for all \(G \mathop{\Rightarrow}_{\mathcal{B},\mathfrak{F}} H\), \( w_\mathcal{T}(G) \mathop{\succeq}_S w_\mathcal{T}(H) \). However, directly verifying all rewriting steps is infeasible due to the potentially infinite number of rewriting steps.

We first introduce the notions of \emph{morphism measurement excluding morphisms in a set}(Definition~\ref{def:weight_excluding_pre}) and \emph{weight of a morphism relative to a morphism-ruler set excluding some specific morphisms}(Definition~\ref{def:weight_excluding}),
to reduce notational overhead (following Endrullis and Overbeek~\cite{endrullis2024generalized_icgt}).
We then recall the notions of traceability of objects (Definition~\ref{def:traceability}) and relative monicity of morphisms (Definition~\ref{def:relative_monicity}),
which are used to define the notion of weighable pushout squares (Definition~\ref{def:weighable}). With the notion of a weighable pushout square, we recall a result from Endrullis and Overbeek~\cite{endrullis2024generalized_icgt} which shows that the weight of the pushout object of a weighable pushout square can be bounded by the weights of the other objects in the pushout square.

\todo{todo to do: pourquoi melange les lettres grecques et latines?}
For any set \( \Gamma \mathop{\subseteq} \operatorname{Hom}(A, G) \),
    we define $\Gamma'$ as the set consisting of all morphisms \( \iota : X \mathop{\to} G \) admitting morphisms \( \zeta \mathop{\colon} X \mathop{\to} A \) and \( \alpha \mathop{\in} \Gamma \) such that $\zeta \mathop{\star} \alpha \mathop{=} \iota$ holds, i.e., such that
     the diagram $\opn{XAG}$ shown below
    %  in Figure~\ref{fig:wf:measurement_of_a_morphism_relative_to_a_morphism_ruler_dfsdfsdsa} 
     is commutative.  Formally,  
    \(
    \Gamma' \overset{\operatorname{def}}{=} \left\{ \iota \mathop{\in} \operatorname{Hom}(X, G)~\middle|~\exists \alpha \mathop{\in} \Gamma,~\exists \zeta:X \mathop{\to} A,~\zeta \mathop{\star} \alpha \mathop{=} \iota \right\}. 
    \)

    \begin{center}
        % \resizebox{0.5\textwidth}{!}{
        \begin{tikzpicture}
            [scale=0.5]
            \node (X) at (0,0) {\(X \)};
            \node (A) at (4,2) {\( A \)};
            \node (G) at (8,0) {\( G \)}; 
            \draw[->] (X) -- (A) node[midway, above] {\( \zeta \)};
            \draw[ ->] (X) -- (G) node[midway, below] {\( \iota \)};
            \draw[->] (A) -- (G) node[midway, above] {\( \alpha \)};
            % \node at (2,0.5) {\( \mathop{=} \)};
            \node (T) at (4,-2) {\( T \)};
            \draw[ ->] (X) -- (T) node[midway, below ] {$e$};
            \draw[<-] (T) -- (G) node[midway, below] {$h$};
        \end{tikzpicture}
        % }
    \end{center} 
     

\begin{definition} 
    \label{def:weight_excluding_pre}
    Let \( \Gamma \mathop{\subseteq} \operatorname{Hom}(A, G) \).
    The \textbf{measurement of a morphism \( h:G \mathop{\to} T \) relative to a morphism-ruler \( e: X \mathop{\to} T \), excluding morphisms in \( \Gamma' \)}, denoted by $m_e(h-\Gamma)$, is defined as:
    \[
        m_e(h-\Gamma) \overset{\operatorname{def}}{=} 
            \card{\set{- \mathop{\star} h \mathop{=} e} \mathop{\setminus} \Gamma'}.
    \]
\end{definition}
\begin{definition}
    \label{def:weight_excluding}
    Let $\mathcal{T}=(T,\mathbb{E},\mathcal{S},w)$ be a weighted type graph. The \textbf{weight of a morphism $h: G \mathop{\to} T$ relative to a set $\mathbb{E}$ of morphism-rulers excluding morphisms in \( \Gamma' \)} is defined as the semiring product of $w(e)^{w_e(h-\Gamma)}$ for $e \mathop{\in} \mathbb{E}$:
    \[ 
        w_\mathcal{T}(h-\Gamma) \overset{\operatorname{def}}{=} \underset{e \mathop{\in} \mathbb{E}}{\mathop{\bigodot}} 
    w(e)^{m_e(h-\Gamma)}.
            \]
\end{definition}
% Consider the pushout square shown below.
% %  in Figure~\ref{fig:preliminaries:pushout_square_traceable_sdldfsfsfsdfkfjsladkj}. Informally, 
%  Any subgraph of the pushout object \(D\) is either the image of \(B\) or the image of \(C\), or is obtained by gluing together parts of the images of \(B\) and \(C\). This intuition is made precise in Definition~\ref{def:traceability}. \todo{todo to do: ????}
% \begin{figure}[H]
%     \centering
%      \resizebox{0.4\textwidth}{!}{
%     \begin{tikzpicture}
%       \node (A) at (0,0) {$A$};
%       \node (B) at (2,0) {$B$}; 
%       \node  (C) at (0,-2) {$C$}; 
%       \node  (D) at (2,-2) {$D$}; 
%       \node  (X) at (4,-2) {$X$};
%       \begin{scope}[nodes=rectangle]          
%       \draw [->] (A) to node [above,label,pos=0.5] {$\alpha$} (B);
%       \draw [->] (A) to node [left,label,pos=0.5] {$\beta$} (C);
%       \draw [->] (B) to node [right,label,pos=0.45] {$\beta'$} (D); 
%       \draw [->] (C) to node [below,label,pos=0.45] {$\alpha'$} (D);
%     %   \draw [->] (X) to node [below,label,pos=0.4] {$h$} (D);
%       \end{scope}
%       \node at ($(A)!.5!(D)$) {$\delta$};
%     \end{tikzpicture}
%     }
%     \caption{}
%     \label{fig:preliminaries:pushout_square_traceable_sdldfsfsfsdfkfjsladkj}
% \end{figure}

\begin{definition}[\cite{endrullis2024generalized_arxiv_v2}]
    \label{def:traceability}
Let $\Delta$ be a class of pushout squares. 
An object $X$ is said to be \textbf{traceable along $\Delta$} if for every diagram $\delta$ in $\Delta$, as shown below:
% in Figure~\ref{fig:preliminaries:pushout_square_traceable_sdlkfjsladkj}, and $h:X \mathop{\to} D$,      
% \begin{figure}[H]
%     \centering
\begin{center}
    %  \resizebox{0.5\textwidth}{!}{
    \begin{tikzpicture}
      \node (A) at (0,0) {$A$};
      \node (B) at (2,0) {$B$}; 
      \node  (C) at (0,-2) {$C$}; 
      \node  (D) at (2,-2) {$D$}; 
      \node  (X) at (4,-2) {$X$};
      \begin{scope}[nodes=rectangle]          
      \draw [->] (A) to node [above,label,pos=0.5] {$\alpha$} (B);
      \draw [->] (A) to node [left,label,pos=0.5] {$\beta$} (C);
      \draw [->] (B) to node [right,label,pos=0.45] {$\beta'$} (D); 
      \draw [->] (C) to node [below,label,pos=0.45] {$\alpha'$} (D);
      \draw [->] (X) to node [below,label,pos=0.4] {$h$} (D);
      \end{scope}
      \node at ($(A)!.5!(D)$) {$\delta$};
    \end{tikzpicture}
    % }
\end{center}
  one of the following conditions holds:
    \begin{enumerate}[label=(\alph*)]
        \item\label{traceable:a} there is a morphism $f : X \mathop{\to} B$ such that $h \mathop{=} f \mathop{\star} \beta'$, or
        \item\label{traceable:b} there is a morphism $g : X \mathop{\to} C$ such that $h \mathop{=} g \mathop{\star} \alpha'$.
    \end{enumerate}
    If additionally, 
    whenever \ref{traceable:a} and \ref {traceable:b} hold,
    \begin{enumerate}[label=(\alph*),resume]
        \item 
        there is a morphism $k : X \mathop{\to} A$ such that $h \mathop{=} k \mathop{\star} \alpha \mathop{\star} \beta' $,
    \end{enumerate}
    then we say that $X$ is \textbf{strongly traceable} along $\Delta$.
\end{definition}
Intuitively, in the category \textbf{Graph}, a graph $X$ is $\beta$-strongly traceable along a pushout square, if whenever $X$ occurs in $D$, it occurs either in $B$ or in $C$, and if it occurs in both, then it occurs in $A$ as well.

\begin{remark}[\cite{endrullis2024generalized_arxiv_v2}]
    \label{remark:traceability_graph}
    In \textbf{Graph}, the objects \tikz[baseline=-0.5ex]{
        \node[draw,circle] (x) at (0,0) {};
    } and
    %  \tikz[baseline=-0.5ex]{
    %     \node (x) at (0,0) {$\bullet$};
    %     \node (y) at (1,0) {$\bullet$};
    %     \draw[->] (x) -- (y) node[midway, above] {$x$};
    % } 
    \raisebox{2pt}{
            \scalebox{0.7}{\tikz[baseline=-0.5ex]{
            \node[draw,circle] (x) at (0,0) {};
            \node[draw,circle] (y) at (1,0) {};
            \draw[->] (x)--(y) node[midway, above] {$x$};
        }}}
    (for edge labels $x$) are the only (non-initial) objects that are (strongly) traceable along all pushout squares. 
    Other objects, such as loops~\tikz[baseline=-0.5ex]{
        \node[draw,circle] (x) at (0,0) {};
        \draw[->] (x) edge[loop right]  node[midway, right] {$x$}  (x)
    }, are strongly traceable if all morphisms in the square are monomorphisms.
\end{remark} 

\begin{example}
    Consider the following DPO diagram:
    %  shown in Figure~\ref{fig:nwf:two_pushout_squares}.
    \begin{center}
    % \begin{figure}[H]
    %     \centering 
      \resizebox{0.98\textwidth}{!}{
      \begin{tikzpicture}
          \graphbox{\( L \)}{0mm}{-3mm}{34mm}{12mm}{2mm}{2mm}{
              \coordinate (o) at (0mm,-8mm); 
              \node[draw,circle] (l1) at ($(o)+(-10mm,0mm)$) {1};
              \node[draw,circle] (l2) at ($(l1)+(2,0)$) {2};
              \node[draw,circle] (l3) at ($(l1)+(1,0)$) {3};
              \draw[] (l1) -- (l3) node[midway,above] {$a$};
              \draw[] (l3) -- (l2) node[midway,above] {$a$};
          } 
          \graphbox{\( K \)}{40mm}{-3mm}{34mm}{12mm}{2mm}{2mm}{
              \coordinate (o) at (0mm,-8mm); 
              \node[draw,circle] (l1) at ($(o)+(-10mm,0mm)$) {1};
              \node[draw,circle] (l2) at ($(l1)+(2,0)$) {2};
          }  
          \graphbox{\( R \)}{80mm}{-3mm}{45mm}{12mm}{2mm}{2mm}{
              \coordinate (o) at (-5mm,-8mm); 
              \node[draw,circle] (l1) at ($(o)+(-10mm,0mm)$) {1};
              \node[draw,circle] (l2) at ($(l1)+(3,0)$) {2};
              \node[draw,circle] (l3) at ($(l1)+(1,0)$) {4};
              \node[draw,circle] (l4) at ($(l1)+(2,0)$) {5};
              \draw[ ] (l1) -- (l3) node[midway,above] {$a$};
              \draw[ ] (l3) -- (l4) node[midway,above] {$a$};
              \draw[ ] (l4) -- (l2) node[midway,above] {$a$};
          }    
          \graphbox{\( G \)}{0mm}{-22mm}{34mm}{22mm}{2mm}{-3mm}{
              \coordinate (o) at (0mm,-3mm); 
              \node[draw,circle] (l1) at ($(o)+(-10mm,0mm)$) {1};
              \node[draw,circle] (l2) at ($(l1)+(2,0)$) {2};
              \node[draw,circle] (l3) at ($(l1)+(1,0)$) {3};
              \node[draw,circle] (l4) at ($(l2)+(0,-1)$) {6};
              \draw[] (l1) -- (l3) node[midway,above] {$a$};
              \draw[] (l3) -- (l2) node[midway,above] {$a$};
              \draw[ ] (l2) -- (l4) node[midway,right] {$a$};
              \node[draw,circle] (l6) at ($(l1)+(0,-1)$) {7};
              \draw[] (l1) -- (l6) node[midway,left] {$a$};
          }    
          \graphbox{\( C  \)}{40mm}{-22mm}{34mm}{22mm}{2mm}{-3mm}{
              \coordinate (o) at (0mm,-3mm); 
              \node[draw,circle] (l1) at ($(o)+(-10mm,0mm)$) {1};
              \node[draw,circle] (l2) at ($(l1)+(2,0)$) {2};
              \node[draw,circle] (l4) at ($(l2)+(0,-1)$) {6};
              \draw[ ] (l2) -- (l4) node[midway,right] {$a$};
              \node[ draw,circle] (l6) at ($(l1)+(0,-1)$) {7};
              \draw[ ] (l1) -- (l6) node[midway,left] {$a$};
          }    
          \graphbox{\( H \)}{80mm}{-22mm}{45mm}{22mm}{2mm}{-3mm}{
              \coordinate (o) at (-5mm,-3mm); 
              \node[draw,circle] (l1) at ($(o)+(-10mm,0mm)$) {1};
              \node[draw,circle] (l2) at ($(l1)+(3,0)$) {2};
              \node[draw,circle] (l3) at ($(l1)+(1,0)$) {4};
              \node[draw,circle] (l4) at ($(l1)+(2,0)$) {5};
              \node[ draw,circle] (l5) at ($(l2)+(0,-1)$) {6};
              \node[ draw,circle] (l6) at ($(l1)+(0,-1)$) {7};
              \draw[ ] (l1) -- (l6) node[midway,left] {$a$};
              \draw[] (l1) -- (l3) node[midway,above] {$a$};
              \draw[] (l3) -- (l4) node[midway,above] {$a$};
              \draw[ ] (l4) -- (l2) node[midway,above] {$a$};
              \draw[ ] (l2) -- (l5) node[midway,right] {$a$};
          }    
          \node () at (37mm,-8mm) {\( \leftarrowtail \)}; % K -> L
          \node () at (77mm,-8mm) {\( \rightarrowtail \)}; % K -> R
          \node () at (15mm,-18mm) {\( m\ \downarrowtail \)};
          \node () at (37mm,-33mm) {\( \leftarrowtail \)};
          \node () at (58mm,-18mm) {\( u\downarrowtail \)};
          \node () at (102mm,-18mm) {\( \downarrowtail \)};
          \node () at (77mm,-33mm) {\( \rightarrowtail \)}; % C -> H
      \end{tikzpicture}
      }
    \end{center}

   The graph 
   \tikz[baseline=-0.5ex]{
        \node[draw,circle] (x) at (0,0) {};
        \node[draw,circle] (y) at (1,0) {};
        \node[draw,circle] (z) at (2,0) {};
        \draw[->] (x) -- (y) node[midway, above] {$a$};
        \draw[<-] (z) -- (y) node[midway, above] {$a$};
    } is not traceable in both pushout squares. The graph \tikz[baseline=-0.5ex]{
        \node[draw,circle] (x) at (0,0) {};
        \node[draw,circle] (y) at (1,0) {};
        \draw[->] (x) -- (y) node[midway, above] {$a$};
    } is traceable (and strongly traceable) along both pushout squares.
\end{example}

% \begin{example} 
%     Consider the rewriting step shown in Figure~\ref{fig:nwf:rewriting_step_traceability}.
%     The graph \begin{tikzpicture}
%         \node[draw, circle] (x) at (0,0) {};
%         \node[draw, circle] (y) at (1,0) {};
%         \draw[->]  (x) -- (y) node [midway,above] {$c$};
%     \end{tikzpicture} 
%     is traceable along both pushout squares because for every morphism $h$ from this graph to $G$ (resp. to $H$), $h$ factors through $m$ or $l$ (resp. through $u$ or $r$).

%      It is $u$-strongly traceable along the left pushout square, because the only morphism $h:X \mathop{\to} G$ factorizable through $m$ and $l$ has as image the graph \begin{tikzpicture}
%         \node[draw, circle] (x) at (0,0) {z};
%         \node[draw, circle] (y) at (1,0) {y};
%         \draw[->]  (x) -- (y) node [midway,above] {$c$};
%      \end{tikzpicture} and there is a unique morphism $k:X \mathop{\to} A$ such that $h \mathop{=} k \mathop{\star} l \mathop{\star} m$.

%     It is not $u$-strongly traceable along the right pushout square
%     The graph \begin{tikzpicture}
%         \node[draw, circle] (x) at (0,0) {};
%         \node[draw, circle] (y) at (1,0) {};
%         \node[draw, circle] (z) at (2,0) {};
%         \draw[->]  (x) -- (y) node [midway,above] {$c$};
%         \draw[->]  (y) -- (z) node [midway,above] {$c$};
%     \end{tikzpicture} is not traceable along both pushout squares. \todo{todo: why}
% %     because    
% %     \begin{tikzpicture}
% %        \node[draw, circle] (x) at (0,0) {z};
% %        \node[draw, circle] (y) at (1,0) {x};
% %        \node[draw, circle] (z) at (2,0) {y};
% %        \draw[->]  (x) -- (y) node [midway,above] {$c$};
% %        \draw[->]  (y) -- (z) node [midway,above] {$c$};
% %    \end{tikzpicture} is in the pushout object $G$ but in neither $L$ nor $C$.
%    \begin{figure}[H]
%     \centering
%     \resizebox{0.7\textwidth}{!}{
%     \begin{tikzpicture}
%         \graphbox{$L$}{0mm}{0mm}{35mm}{35mm}{2mm}{-7mm}{
%         \node[draw,circle] (1) at (0,0) {x};
%         \node[draw,circle] (2) at (1,-2) {y};
%         \node[draw,circle] (3) at (-1,-2) {z};
%         \draw[->] (1) edge node[midway,right]  {$c$}  node[midway,left]{} (2);
%         \draw[->] (2) edge node[midway,above]  {$c$}  node[midway,below]{}  (3);
%         }
%         \graphbox{$K$}{45mm}{0mm}{35mm}{35mm}{2mm}{-7mm}{
%         % I
%             % \def\x{5};
%             % \def\y{0};
%             \node[draw,circle] (1) at (0,0) {x};
%             \node[draw,circle] (2) at (1,-2) {y};
%             \node[draw,circle] (3) at (-1,-2) {z};
%             \draw[->] (2) edge node[midway,above]  {$c$}  node[midway,below]{}  (3);
%         }
%         \node () at (40mm,-15mm)  {$\overset{l}{\leftarrowtail}$};

%         % \node () at (2.5,-2.5) {$PO$};
%         \graphbox{$R$}{90mm}{0mm}{35mm}{35mm}{2mm}{-7mm}{
%         % R
%             % \def\x{10};
%             % \def\y{-1.1};
%             \node[draw,circle] (1) at (0,-1) {x y z};
%             \draw[->] (1) edge[loop below]  node[midway, below] {$c$}  (1);
%         }
%         \node () at (85mm,-15mm)  {$\overset{r}{\rightarrow}$};

%         \graphbox{$G$}{0mm}{-45mm}{35mm}{35mm}{2mm}{-7mm}{
%         % \node () at (7.5,-2.5) {$PO$};
%         % G
%             % \def\x{0};
%             % \def\y{-5};
%             \node[draw,circle] (1) at (0,0) {x};
%             \node[draw,circle] (2) at (1,-2) {y};
%             \node[draw,circle] (3) at (-1,-2) {z};
%             \draw[->] (1) edge node[midway,right]  {$c$}  node[midway,left]{} (2);
%             \draw[->] (2) edge node[midway,above]  {$c$}  node[midway,below]{}  (3);
%             \draw[->] (3) edge node[midway,above]  {$c$}  node[midway,below]{}  (1);
%         }
%         % C
%         \graphbox{$C$}{45mm}{-45mm}{35mm}{35mm}{2mm}{-7mm}
%         {    
%             \node[draw,circle] (1) at (0,0) {x};
%             \node[draw,circle] (2) at (1,-2) {y};
%             \node[draw,circle] (3) at (-1,-2) {z};
%             % \draw[->] (1) edge node[midway,right]  {$c$}  node[midway,left]{} (2);
%             \draw[->] (3) edge node[midway,above]  {$c$}  node[midway,below]{} (1);
%             \draw[->] (2) edge node[midway,above]  {$c$}  node[midway,below]{} (3);
%             }
%         \node () at (40mm,-55mm)  {$\overset{l'}{\leftarrowtail}$};
%         % H
%         \graphbox{$H$}{90mm}{-45mm}{35mm}{35mm}{2mm}{-7mm}{
%         % R
%             % \def\x{10};
%             % \def\y{-1.1};
%             \node[draw,circle] (1) at (0,-1) {x y z};
%             \draw[->] (1) edge[loop below]  node[midway, below] {} node[midway, below] {$c$} (1);
%             \draw[->] (1) edge[loop above]  node[midway, below] {} node[midway, above] {$c$} (1);
%         }

%         \node () at (85mm,-55mm)  {$\overset{r'}{\rightarrow}$};
%         \node () at (17mm,-40mm) {$m\downarrowtail$};
%         \node () at (62mm,-40mm) {$u\downarrowtail$};
%         \node () at (107mm,-40mm) {$\downarrow m'$};
%     \end{tikzpicture}
%     }
%     \caption{}
%     \label{fig:nwf:rewriting_step_traceability}
%    \end{figure}
% \end{example} 
The following concepts are restricted versions of the concept of monicity.
\begin{definition}[\cite{endrullis2024generalized_arxiv_v2}]
    \label{def:relative_monicity}
    Let $A,B,C$ and $X$ be objects, $\Gamma$ a set of objects, $S$ a set of morphisms to $A$ and $u:C \mathop{\to} A$ a morphism. A morphism $f : A \mathop{\to} B$ is said to be
    \begin{itemize} 
        \item 
            \textbf{monic for $S$} 
            if $g \mathop{\star} f \mathop{=} h \mathop{\star} f$ implies $g \mathop{=} h$ for all $g, h \mathop{\in} S$;
        \item 
            \textbf{$X$-monic} if $f$ is monic for $\operatorname{Hom}(X, A)$;
        \item \textbf{$X$-monic outside of $u$}, if $f$ is monic for \( \operatorname{Hom}(X,A) \mathop{\setminus} \left ( \operatorname{Hom}(X,C) \mathop{\star} u \right ) \);
        \item  \textbf{$\Gamma$-monic} if $f$ is $X$-monic for every $X \mathop{\in} \Gamma$.
    \end{itemize}
\end{definition} 
\begin{example}[Edge-Monicity \cite{endrullis2024generalized_arxiv_v2}]
 
 Let \(f\mathop{\colon} G \mathop{\to} H\) be a morphism in \(\mathbf{Graph}\) that does not identify distinct edges, but may identify distinct nodes. Then \(f\) is not necessarily monic. Indeed, if \(u\mathop{\neq} v\) are nodes of \(G\) with \(f(u)=f(v)\), let \(K\) be the graph consisting of a single node and no edges, and let \(g,h\mathop{\colon} K \mathop{\to} G\) send that node to \(u\) and \(v\), respectively. Then \(g\mathop{\neq} h\) while \(f\circ g=f\circ h\), so \(f\) is not monic.

Nevertheless \(f\) is \(\Gamma\)-monic, where
\[
\Gamma \mathop{=} \left\{ \vcenter{\hbox{
\begin{tikzpicture}[baseline=-0.5ex]
\node[draw,circle] (A) at (0,0) {};
\node[draw,circle] (B) at (1,0) {};
\draw[->] (A) -- (B) node[midway, above] {$x$};
\end{tikzpicture}
}} \mathop{\mid} x \text{ is an edge label} \right\}.
\]
To see this, let \(X\in\Gamma\) and let \(g,h\mathop{\colon} X \mathop{\to} G\) satisfy \(g \mathop{\star} f \mathop{=} h \mathop{\star} f\). Let \(e\) denote the unique edge of \(X\). From \(f(g(e))\mathop{=}f(h(e))\) and the hypothesis that \(f\) is injective on edges we obtain \(g(e)\mathop{=}h(e)\). Since we are working with directed graphs, the images of the two nodes of \(X\) are determined by the image of \(e\), so \(g\) and \(h\) agree on both nodes. Hence \(g \mathop{=} h\), and \(f\) is \(\Gamma\)-monic.

\end{example}
\begin{definition}[\text{\cite[\textdef~4.9]{endrullis2024generalized_arxiv_v2}}]
    \label{def:weighable}
        Let  $\mathcal{T} \mathop{=} (T,\mathbb{E}, S, w)$ be a weighted type graph.
        Consider the pushout square $\delta$ shown below.
        \begin{center}
            \begin{tikzpicture}
            \node (A) at (0,0) {$A$};
            \node (B) at (2,0) {$B$}; 
            \node  (C) at (0,-2) {$C$}; 
            \node  (D) at (2,-2) {$D$}; 
            \begin{scope}[nodes=rectangle]          
            \draw [->] (A) to node [above,label,pos=0.5] {$\alpha$} (B);
            \draw [->] (A) to node [left,label,pos=0.5] {$\beta$} (C);
            \draw [->] (B) to node [right,label,pos=0.45] {$\beta'$} (D); 
            \draw [->] (C) to node [below,label,pos=0.45] {$\alpha'$} (D);
            \end{scope}
            \node at ($(A)!.5!(D)$) {$\delta$};
            \end{tikzpicture}
        \end{center}
        We say that $\delta$ is said to be
         \begin{enumerate}[label=(\alph*)]
        \item \textbf{weighable} with $\mathcal{T}$ if the following hold:
            \begin{enumerate}[label=(\roman*)]
                \item $dom(\mathbb{E})$ is strongly traceable along $\delta$,
                \item $\beta'$ is $dom(\mathbb{E})$-monic,
                \item $\alpha'$ is $dom(\mathbb{E})$-monic outside of $\beta$.
            \end{enumerate}
        \item \textbf{bounded-above} by $\mathcal{T}$ if $dom(\mathbb{E})$ is traceable along $\delta$.
    \end{enumerate}
   
\end{definition}


The following lemma shows that, under certain constraints, exact weights of pushout object can be precisely computed or upperbounded.
\begin{lemma}[\text{\cite[Lemma 4.13]{endrullis2024generalized_arxiv_v2}}]
    \label{lem_4d13}
    Let $\mathcal{T} \mathop{=} (T,\mathbb{E}, S, w)$ be a finitary weighted type graph. Consider the pushout square $\delta$ in Example~\ref{def:weighable}.
  We define
     \[k \mathop{=} \underset{t_A:A \mathop{\rightarrow} T}{\mathop{\bigoplus}}
            \left ( 
                \underset{\substack{t_C:C \mathop{\rightarrow} T\\
                                            t_A \mathop{=} \beta \mathop{\star} t_C }}{\mathop{\bigoplus}}
                        w_\mathcal{T}(t_C - \beta)     
                 \right ) 
            \mathop{\odot} 
                w_\mathcal{T}(\set{\alpha \mathop{\star} - \mathop{=} t_A})
    \]
    The following conditions hold:
    \begin{enumerate}[label=(\Alph*)]
        \item  $w_\mathcal{T}(D)=k$ if $\delta$ is weighable with $\mathcal{T}$.
        \item  $w_\mathcal{T}(D)\mathop{\preceq} k$ if $\delta$ is bounded above by $\mathcal{T}$  and \(w(e) \mathop{\succeq} 1_S\) for all $e \mathop{\in} \mathbb{E}$.
    \end{enumerate}
\end{lemma}

    Consider the following DPO diagram
    %  shown in Figure~\ref{fig:nwf:dpo_diagram_sdfslafgjsl} 
     that defines a rewriting step \( G \mathop{\Rightarrow}_{\rho,\mathfrak{F}} H \).  
    \begin{center}  
        \begin{tikzpicture}
            % [scale=1.5]
            % [node distance=11mm]
            \node (I) at (0,0) {$K$};
            \node (L) at (-2,0) {$L$};
            \node (R) at (2,0) {$R$};
            \node (G) at (-2,-2) {$G$};
            \node (C) at (0,-2) {$C$};
            \node (H) at (2,-2) {$H$};
            \draw [->] (I) to  node [midway,below] {$l$} (L);
            \draw [->] (I) to  node [midway,below] {$r$} (R);
            \draw [->] (L) to node [midway,right] {$m$} (G);
            \draw [->] (I) to node [midway,right] {$u$} (C);
            \draw [->] (R) to node [midway,left] {$m'$} (H);
            \draw [->] (C) to node [midway,above] {$l'$} (G);
            \draw [->] (C) to node [midway,above] {$r'$} (H);
            \node [at=($(I)!.5!(G)$)] {\normalfont PO};
            \node [at=($(I)!.5!(H)$)] {\normalfont PO};
        \end{tikzpicture}
    \end{center}
    If the left pushout square is weighable with $\mathcal{T}$ and the right pushout square is bounded above by $\mathcal{T}$, by Lemma~\ref{lem_4d13}, we obtain:

\begin{flalign*}
    w_\mathcal{T}(G) 
        & \mathop{=} \mathop{\bigoplus}_{t_K: K \mathop{\rightarrow} T} 
        \left ( \mathop{\bigoplus}_{\substack{t_C: C \mathop{\rightarrow} T\\ t_K \mathop{=} u \mathop{\star} t_C}}
          w_\mathcal{T}(t_C - u) \right ) 
          \mathop{\odot}
        \left (\mathop{\bigoplus}_{\substack{t_L: L \mathop{\rightarrow} T\\ t_K \mathop{=} l \mathop{\star} t_L}}
        w_\mathcal{T}(t_L) \right )
         \\
    w_\mathcal{T}(H) 
        &  \mathop{\preceq} \mathop{\bigoplus}_{t_K: K \mathop{\rightarrow} T} 
        \left ( \mathop{\bigoplus}_{\substack{t_C: C \mathop{\rightarrow} T\\ t_K \mathop{=} u \mathop{\star} t_C}}
         w_\mathcal{T}(t_C - u) \right ) 
         \mathop{\odot} 
         \left ( \mathop{\bigoplus}_{\substack{t_R: R \mathop{\rightarrow} T\\ t_K \mathop{=} r \mathop{\star} t_R}}
            w_\mathcal{T}(t_R) \right ) \\
\end{flalign*}
According to these results, in \textsection\ref{wf:sec:termination}, we establish a termination criterion by, for every $t_K: K \mathop{\rightarrow} T$, comparing
$\mathop{\bigoplus}_{\substack{t_L: L \mathop{\rightarrow} T\\ t_K \mathop{=} l \mathop{\star} t_L}}
        w_\mathcal{T}(t_L)$ and 
$\mathop{\bigoplus}_{\substack{t_R: R \mathop{\rightarrow} T\\ t_K \mathop{=} r \mathop{\star} t_R}} 
        w_\mathcal{T}(t_R)$.
Before doing so, we introduce the concept of context closure in the following section. 
   

It is crucial to ensure that the weight of any object subject to rewriting is not $0_S$, because \(0_S\) behaves unpredictably in strongly monotonic measurable semirings. For instance, in the natural and real tropical semirings \((\mathbb{N} \cup \{+\infty\}, \min, +, +\infty, 0_S, <, \leq)\), \(0_S\) is the greatest element \(+\infty\), while in the natural and real arctic semirings \((\mathbb{N} \cup \{-\infty\}, \max, +, -\infty, 0_S, <, \leq)\), \(0_S\) is the smallest element \(-\infty\).
This issue is addressed by the existence of a context closure. 

\subsection{Context Closure}
\label{sec:context_closure}
% \begin{remark} 
%   \label{remark:semiring_0_unpredictable}
%   The requirement \textquote{for all \(e \in \mathbb{E}, w(e) \neq 0_S\)} is necessary because \(0_S\) behaves unpredictably in strongly monotonic measurable semirings. For instance, in the natural and real tropical semirings \((\mathbb{N} \cup \{+\infty\}, \min, +)\), \(0_S\) is the greatest element \(+\infty\), while in the natural and real arctic semirings \((\mathbb{N} \cup \{-\infty\}, \max, +)\), \(0_S\) is the smallest element \(-\infty\).
% \end{remark} 

\begin{definition}[Context Closure~\cite{endrullis2024generalized_arxiv_v2}]
    \label{def:context_closure}  
    \ \newline 
\begin{minipage}{0.65\textwidth}
    Let $\mathcal{T}=(T,\mathbb{E},\mathcal{S},w)$ be a finitary weighted type graph, \(\rho = (L \overset{l}{\leftarrow} K \overset{r}{\rightarrow} R ) \) a DPO rewriting rule and $\mathfrak{F}$ a rewriting framework. 
    A \textbf{context closure} for $\rho$ and $\mathcal{T}$ in $\mathfrak{F}$ is a morphism $c:L \rightarrow T$ such that for every DPO diagram in $\mathfrak{F}(\rho)$ (shown on the right) 
    there exists $\alpha : G \rightarrow T$ such that $m \star \alpha = c$.
\end{minipage}
\begin{minipage}{0.35\textwidth}
    \begin{center}
        \begin{tikzpicture}[rotate=90]
          \node (I) {$K$}; 
          \node (L) [left of=I] {$L$};
          \node (R) [right of=I] {$R$};
          \node (G) [below of=L] {$G$};
          \node (C) [below of=I] {$C$};
          \node (H) [below of=R] {$H$};
          \node (T) [left=of $(L)!0.5!(G)$] {$T$};
          \draw [->] (L) to  node [label, above] {$c$}  (T);
          \draw [->] (G) to  node [label, below] {$\alpha$} (T);
          \draw [->] (I) to node [label, above] {$l$} (L);
          \draw [->] (I) to node [label,above] {$r$} (R);
          \draw [->] (L) to node [label, right] {$m$} (G);
          \draw [->] (I) to (C);
          \draw [->] (R) to (H);
          \draw [->] (C) to (G);
          \draw [->] (C) to (H);
        \end{tikzpicture}
      \end{center}
\end{minipage}
\end{definition}
\begin{example}
    \label{example:context_closure}
    Let the set of edge labels be $\Sigma = \{a,b\}$.
    Consider the DPO rule in \autoref{ex:grsaa}, the weighted type graph in \autoref{example:wft:weighted_type_graph}, 
   and the morphism $c$ illustrated below. Since for every match $m : L \to G$, there exists a morphism $\alpha : G \to T$ such that $m \star \alpha = c$, the morphism $c$ is a context closure for the DPO rule in the type graph.
  \begin{center}
    \begin{tikzpicture}
      \graphbox{\( L \)}{-50mm}{0mm}{40mm}{40mm}{2mm}{-8mm}{
        \coordinate (o) at (0mm,-10mm); 
        \node[draw,circle] (l1) at ($(o)+(-10mm,0mm)$) {1};
        \node[draw,circle] (l2) at ($(l1)+(2,0)$) {2};
        \node[draw,circle] (l3) at ($(l1) + (1,0)$) {3};
        \draw[] (l1) -- (l3) node[midway,above] {a};
        \draw[] (l3) -- (l2) node[midway,above] {a};
    } 
        \graphbox{$T$}{0mm}{0mm}{40mm}{40mm}{-10mm}{-17mm}{
            % \node[draw,circle] (1) at (0,0) {$1\ 2\ 3$};
            % \node[draw,circle] (2) at (2,0) {};
            \coordinate (o) at (2mm,-3mm); 
            \node[draw,circle] (1) at ($(o)+(0,0mm)$) {$1\ 2\ 3$};
            \node[draw,circle] (2) at ($(o)+(2,0)$) {};
            \draw[->] (1) edge[loop above] node[midway, above] {$a^{1}$} (1) ;
            \draw[->] (1) edge[loop below] node[midway, below] {$b^{1}$} (1) ;
            \draw[->] (1) edge[bend left] node[midway, above] {$a^{1}$}  (2)  ;
            \draw[->] (2) edge[bend left] node[midway, below] {$a^{1}$} (1)   ;
        }
        \node () at (-5mm,-15mm) {$\overset{c}{\to}$};
    \end{tikzpicture}
  \end{center}
\end{example} 

\subsection{Decreasing rules}
\label{sec:decreasing_rules}
The definition below adapts the concept of decreasing rules from~\cite{endrullis2024generalized_arxiv_v2}. The difference lies in the treatment of weight comparisons for uniformly decreasing and closure decreasing rules: instead of directly comparing weights, we evaluate the difference between weights relative to the homomorphism $\mu$ of the strongly monotonic measurable semirings.
%  This difference must exceed a fixed positive constant $\delta \in \mathbb{R}_{>0}$.

\begin{definition}[Decreasing Rules]
    \label{def:decreasing_rule}
    Let $\mathcal{T} = (T,\mathbb{E}, (S, \oplus, \odot, 0_S, 1_S, \prec, \leq),w)$ be a finitary weighted type graph, \(\mathfrak{F}\) a DPO rewriting framework, $\rho = (L \overset{l}{\leftarrow} K \overset{r}{\rightarrow} R)$ a DPO rewriting rule.

    \noindent
    The rule $\rho$ is \textbf{weakly decreasing} w.r.t. $\mathcal{T}$ in $\mathfrak{F}$ if 
            for every $t_K : K \to T$,
                $$ 
                  w_\mathcal{T}(\{l \star - = t_K\}) \geq w_\mathcal{T}(\{r\star - = t_K\})$$
           
    \noindent
    The rule $\rho$ is \textbf{uniformly decreasing} w.r.t. $\mathcal{T}$ in $\mathfrak{F}$ if the following hold:
        \begin{itemize}
            \item[]- there exists a context closure $c_\rho$ for $\rho$ and $\mathcal{T}$ in $\mathfrak{F}$, 
            \item[]- for every $t_K : K \to T$,
            \begin{itemize}
                \item[] $\bullet$ $\{l \star - = t_K\} = \emptyset = \{r \star - = t_K\}$, or
                \item[] $\bullet$ $w_\mathcal{T}(\{l \star - = t_K\}) >   w_\mathcal{T}(\{r \star - = t_K\}) $.
            \end{itemize}
        \end{itemize}  
         
    \noindent
   If $S$ is moreover strictly monotonic, we say the rule $\rho$ is
            \textbf{closure decreasing} w.r.t. $\mathcal{T}$ in $\mathfrak{F}$ if the following hold:
            \begin{itemize}
                \item[]- $\rho$ is weakly decreasing,
                \item[]- there exists a context closure $c_\rho$ for $\rho$ and $\mathcal{T}$ in $\mathfrak{F}$,
                \item[]- $w_\mathcal{T}(\{l \star - = t_K\})  >  w_\mathcal{T}(\{r \star - = t_K\})$ for $t_K = l \star c_\rho$.
            \end{itemize}
\end{definition}

\begin{example}
    \label{example:decreasing_rule}
    Consider the DPO rule in \autoref{ex:grsaa} and the weighted type graph in \autoref{example:wft:weighted_type_graph} over the natural arithmetic semiring $\mathfrak{N} = (\mathbb{N},+,*,0_\mathbb{N},1_\mathbb{N},<,\leq)$. There are $t_K^{11}, t_K^{12}, t_K^{21}, t_K^{22}:K \rightarrow T$ as depicted below:

    \begin{center}
        \resizebox{0.49\textwidth}{!}{
            \begin{tikzpicture}
            \graphbox{\( K \)}{-50mm}{0mm}{40mm}{30mm}{2mm}{-6mm}{
                \coordinate (o) at (0mm,-10mm); 
                \node[draw,circle] (l1) at ($(o)+(-10mm,0mm)$) {1};
                \node[draw,circle] (l2) at ($(l1)+(2,0)$) {2};
                % \node[draw,circle] (l3) at ($(l1) + (1,0)$) {3};
                % \draw[] (l1) -- (l3) node[midway,above] {a};
                % \draw[] (l3) -- (l2) node[midway,above] {a};
            } 
                \graphbox{$T$}{0mm}{0mm}{40mm}{30mm}{-10mm}{-15mm}{
                    \node[draw,circle] (1) at (0,0) {$1\ 2$};
                    \node[draw,circle] (2) at (2,0) {};
                    \draw[->] (1) edge[loop above] node[midway, above] {$a$} (1) ;
                    \draw[->] (1) edge[loop below] node[midway, below] {$b$} (1) ;
                    \draw[->] (1) edge[bend left] node[midway, above] {$a$}  (2)  ;
                    \draw[->] (2) edge[bend left] node[midway, below] {$a$} (1)   ;
                }
                \node () at (-5mm,-15mm) {$\overset{t_K^{11}}{\to}$};
            \end{tikzpicture}
            } 
            \resizebox{0.49\textwidth}{!}{
            \begin{tikzpicture}
                \graphbox{\( K \)}{-50mm}{0mm}{40mm}{30mm}{2mm}{-6mm}{
                \coordinate (o) at (0mm,-10mm); 
                \node[draw,circle] (l1) at ($(o)+(-10mm,0mm)$) {1};
                \node[draw,circle] (l2) at ($(l1)+(2,0)$) {2};
                % \node[draw,circle] (l3) at ($(l1) + (1,0)$) {3};
                % \draw[] (l1) -- (l3) node[midway,above] {a};
                % \draw[] (l3) -- (l2) node[midway,above] {a};
            } 
                \graphbox{$T$}{0mm}{0mm}{40mm}{30mm}{-10mm}{-15mm}{
                    \node[draw,circle] (1) at (0,0) {$1$};
                    \node[draw,circle] (2) at (2,0) {2};
                    \draw[->] (1) edge[loop above] node[midway, above] {$a$} (1) ;
                    \draw[->] (1) edge[loop below] node[midway, below] {$b$} (1) ;
                    \draw[->] (1) edge[bend left] node[midway, above] {$a$}  (2)  ;
                    \draw[->] (2) edge[bend left] node[midway, below] {$a$} (1)   ;
                }
                \node () at (-5mm,-15mm) {$\overset{t_K^{12}}{\to}$};
            \end{tikzpicture}
            }
            \hspace{5mm}
            
            \resizebox{0.49\textwidth}{!}{
            \begin{tikzpicture}
                \graphbox{\( K \)}{-50mm}{0mm}{40mm}{30mm}{2mm}{-6mm}{
                \coordinate (o) at (0mm,-10mm); 
                \node[draw,circle] (l1) at ($(o)+(-10mm,0mm)$) {1};
                \node[draw,circle] (l2) at ($(l1)+(2,0)$) {2};
                % \node[draw,circle] (l3) at ($(l1) + (1,0)$) {3};
                % \draw[] (l1) -- (l3) node[midway,above] {a};
                % \draw[] (l3) -- (l2) node[midway,above] {a};
            } 
                \graphbox{$T$}{0mm}{0mm}{40mm}{30mm}{-10mm}{-15mm}{
                    \node[draw,circle] (1) at (0,0) {2};
                    \node[draw,circle] (2) at (2,0) {1};
                    \draw[->] (1) edge[loop above] node[midway, above] {$a$} (1) ;
                    \draw[->] (1) edge[loop below] node[midway, below] {$b$} (1) ;
                    \draw[->] (1) edge[bend left] node[midway, above] {$a$}  (2)  ;
                    \draw[->] (2) edge[bend left] node[midway, below] {$a$} (1)   ;
                }
                \node () at (-5mm,-15mm) {$\overset{t_K^{21}}{\to}$};
            \end{tikzpicture}
            }
            \resizebox{0.49\textwidth}{!}{
            \begin{tikzpicture}
                \graphbox{\( K \)}{-50mm}{0mm}{40mm}{30mm}{2mm}{-6mm}{
                \coordinate (o) at (0mm,-10mm); 
                \node[draw,circle] (l1) at ($(o)+(-10mm,0mm)$) {1};
                \node[draw,circle] (l2) at ($(l1)+(2,0)$) {2};
                % \node[draw,circle] (l3) at ($(l1) + (1,0)$) {3};
                % \draw[] (l1) -- (l3) node[midway,above] {a};
                % \draw[] (l3) -- (l2) node[midway,above] {a};
            } 
                \graphbox{$T$}{0mm}{0mm}{40mm}{30mm}{-10mm}{-15mm}{
                    \node[draw,circle] (1) at (0,0) {};
                    \node[draw,circle] (2) at (2,0) {$1\ 2$};
                    \draw[->] (1) edge[loop above] node[midway, above] {$a$} (1) ;
                    \draw[->] (1) edge[loop below] node[midway, below] {$b$} (1) ;
                    \draw[->] (1) edge[bend left] node[midway, above] {$a$}  (2)  ;
                    \draw[->] (2) edge[bend left] node[midway, below] {$a$} (1)   ;
                }
                \node () at (-5mm,-15mm) {$\overset{t_K^{22}}{\to}$};
            \end{tikzpicture}
            }
      \end{center}
    The set $\{l \star - = t_K^{11}\}$ has two morphisms $h_{11}^1$ and $h_{11}^2$ as illustrated below:
    \begin{center}
        \resizebox{0.49\textwidth}{!}{
        \begin{tikzpicture}
          \graphbox{\( L \)}{-50mm}{0mm}{40mm}{40mm}{2mm}{-6mm}{
            \coordinate (o) at (0mm,-10mm); 
            \node[draw,circle] (l1) at ($(o)+(-10mm,0mm)$) {1};
            \node[draw,circle] (l2) at ($(l1)+(2,0)$) {2};
            \node[draw,circle] (l3) at ($(l1) + (1,0)$) {3};
            \draw[] (l1) -- (l3) node[midway,above] {a};
            \draw[] (l3) -- (l2) node[midway,above] {a};
        } 
            \graphbox{$T$}{0mm}{0mm}{40mm}{40mm}{-10mm}{-17mm}{
                \node[draw,circle] (1) at (0,0) {$1\ 2$};
                \node[draw,circle] (2) at (2,0) {3};
                \draw[->] (1) edge[loop above] node[midway, above] {$a^{1}$} (1) ;
                \draw[->] (1) edge[loop below] node[midway, below] {$b^{1}$} (1) ;
                \draw[->] (1) edge[bend left] node[midway, above] {$a^{1}$}  (2)  ;
                \draw[->] (2) edge[bend left] node[midway, below] {$a^{1}$} (1)   ;
            }
            \node () at (-5mm,-15mm) {$\overset{h_{11}^1}{\to}$};
        \end{tikzpicture}
        }
        \resizebox{0.49\textwidth}{!}{
            \begin{tikzpicture}
              \graphbox{\(L\)}{-50mm}{0mm}{40mm}{40mm}{2mm}{-10mm}{
                \coordinate (o) at (0mm,-10mm); 
                \node[draw,circle] (l1) at ($(o)+(-10mm,0mm)$) {1};
                \node[draw,circle] (l2) at ($(l1)+(2,0)$) {2};
                \node[draw,circle] (l3) at ($(l1) + (1,0)$) {3};
                \draw[] (l1) -- (l3) node[midway,above] {a};
                \draw[] (l3) -- (l2) node[midway,above] {a};
            } 
                \graphbox{$T$}{0mm}{0mm}{40mm}{40mm}{-10mm}{-20mm}{
                    \node[draw,circle] (1) at (0,0) {$1\ 2\ 3$};
                    \node[draw,circle] (2) at (2,0) {};
                    \draw[->] (1) edge[loop above] node[midway, above] {$a^{1}$} (1) ;
                    \draw[->] (1) edge[loop below] node[midway, below] {$b^{1}$} (1) ;
                    \draw[->] (1) edge[bend left] node[midway, above] {$a^{1}$}  (2)  ;
                    \draw[->] (2) edge[bend left] node[midway, below] {$a^{1}$} (1)   ;(1)   ;
                }
                \node () at (-5mm,-15mm) {$\overset{h_{11}^2}{\to}$};
            \end{tikzpicture}
            }
      \end{center}
    Therefore, we have \begin{flalign*}
        w_\mathcal{T}(\{l \star - = t_K^{11}\})
        =&w_\mathcal{T}(\{h_{11}^1, h_{11}^2\})\\
        \overset{\mathrm{def}}{=}&w_\mathcal{T}(h_{11}^1) + w_\mathcal{T}(h_{11}^2) \\
        =&(1^1 * 1^1) + (1^1 * 1^1)\\
        =&2
    \end{flalign*}
    The set $\{r \star - = t_K^{11}\}$ has one morphism $h_{11}^3$ as illustrated below:
    \begin{center}
        \resizebox{0.49\textwidth}{!}{
        \begin{tikzpicture}
          \graphbox{\( R \)}{-55mm}{0mm}{45mm}{44mm}{1mm}{-22mm}{
            \coordinate (o) at (-5mm,-3mm); 
            \node[draw,circle] (l1) at ($(o)+(-10mm,0mm)$) {1};
            \node[draw,circle] (l2) at ($(l1)+(3,0)$) {2};
            \node[draw,circle] (l3) at ($(l1) + (1,0)$) {4};
            \node[draw,circle] (l4) at ($(l1) + (2,0)$) {5};
            \draw[->] (l1) -- (l3) node[midway,above] {a};
            \draw[->] (l3) -- (l4) node[midway,above] {b};
            \draw[->] (l4) -- (l2) node[midway,above] {a};
        } 
            \graphbox{$T$}{0mm}{0mm}{40mm}{44mm}{-10mm}{-22mm}{
                \node[draw,circle] (1) at (0,0) {$1\ 2\ 4\ 5$};
                \node[draw,circle] (2) at (2,0) {};
                \draw[->] (1) edge[loop above] node[midway, above] {$a^{1}$} (1) ;
                \draw[->] (1) edge[loop below] node[midway, below] {$b^{1}$} (1) ;
                \draw[->] (1) edge[bend left] node[midway, above] {$a^{1}$}  (2)  ;
                \draw[->] (2) edge[bend left] node[midway, below] {$a^{1}$} (1)   ;
            }
            \node () at (-5mm,-19mm) {$\overset{h_{11}^3}{\to}$};
        \end{tikzpicture}
        }
      \end{center}
    Therefore, we have: $w_\mathcal{T}(\{r \star - = t_K^{11}\}) = w_\mathcal{T}(h_{11}^3) = 1^1 * 1^1 * 1 ^ 1 = 1$. Thus, we have $w_\mathcal{T}(\{l \star - = t_K^{11}\}) = 2 \geq 1 = w_\mathcal{T}(\{r \star - = t_K^{11}\})$.

    Similarly, we can check that $w_\mathcal{T}(\{l \star - = t_K^{12}\}) = 1 \geq 1 = w_\mathcal{T}(\{r \star - = t_K^{12}\})$,  $w_\mathcal{T}(\{l \star - = t_K^{21}\}) = 1 \geq 1 = w_\mathcal{T}(\{r \star - = t_K^{21}\})$, and $w_\mathcal{T}(\{l \star - = t_K^{22}\}) = 1 \geq 1 = w_\mathcal{T}(\{r \star - = t_K^{22}\})$. The rule is therefore weakly decreasing.

    There exists a context closure $c$ for the DPO rule in the weighted type graph, as shown in \autoref{example:context_closure}.
    Since we have additionally $t_K^{11} = l \star c$ and $w_\mathcal{T}(\{l \star - = t_K^{11}\}) = 2 > 1 = w_\mathcal{T}(\{r \star - = t_K^{11}\})$, the rule is closure decreasing since the semiring is strictly monotonic.
\end{example} 
\begin{theorem}[Termination of DPO rewriting system] 
    \label{thm:wf:termination_grs}
    Let $\mathcal{A}$ and $\mathcal{B}$ be sets of DPO rewriting rules, $\mathcal{T} = (T,\mathbb{E}, (S, \oplus, \odot, 0_S, 1_S, \prec, \leq), w)$ a finitary weighted type graph and $\mathfrak{F}$ a DPO rewriting framework such that

        \begin{itemize}
            \item \(\operatorname{left}(\Delta)\) is weighable with \(\mathcal{T}\),
            \item \(\operatorname{right}(\Delta)\) is bounded above by \(\mathcal{T}\). 
        \end{itemize}
    for every rule $\rho \in (\mathcal{A }\cup \mathcal{B })$ and every double pushout diagram  
        $\Delta \in \mathfrak{F}(\rho)$. If the following conditions hold:
    \begin{enumerate}
        \item either every $\rho \in \mathcal{A}$ is uniformly decreasing or every $\rho \in \mathcal{A}$ is closure decreasing, and
        \item every rule $\rho \in \mathcal{B}$ is weakly decreasing,
    \end{enumerate}
    then $\Rightarrow_{\mathcal{A},\mathfrak{F}}$ is \textbf{terminating} relative to $\Rightarrow_{\mathcal{B},\mathfrak{F}}$.
\end{theorem} 
\begin{example} 
    \label{example:wf:termination}
    Termination of the DPO rule in \autoref{ex:grsaa} can be established using \autoref{thm:termination_grs} together with the weighted type graph in \autoref{example:wft:weighted_type_graph} over the natural arithmetic semiring $\mathfrak{N} = (\mathbb{N},+,*,0_\mathbb{R},1_\mathbb{N},<,\leq)$. It is closure decreasing by \autoref{example:decreasing_rule}.
\end{example}

The following lemma guaranteeing that, under certain constraints: exact weights of host graphs can be computed, and upper bounds for result-graph weights can be derived. 
\begin{lemma}[\cite{endrullis2024generalized_arxiv_v2}]
    \label{lem_4d13}
\ \newline
\begin{minipage}{0.7\textwidth}
    Let $\mathcal{T} = (T,\mathbb{E}, (S, \oplus, \odot, 0_S, 1_S, \prec, \mu), w)$ be a finitary weighted type graph. Consider the pushout square $\delta$ illustrated on the right. We define
\end{minipage}
\begin{minipage}{0.3\textwidth}
    \begin{center}{\normalfont
        \begin{tikzpicture}[node distance=12mm]
            \node (A) {$A$};
            \node (B) [right of=A] {$B$};
            \node (C) [below of=A] {$C$};
            \node (D) [right of=C] {$D$};
            
            \draw [->] (A) to node [above, label] {$\alpha$} (B);
            \draw [->] (A) to node [left, label] {$\beta$} (C);
            \draw [->] (B) to node [right, label] {$\beta'$} (D);
            \draw [->] (C) to node [below, label] {$\alpha'$} (D);
            
            \node [at=($(A)!.5!(D)$)] {$\delta$};
        \end{tikzpicture}
    }\end{center}
\end{minipage}
     \[k = \underset{t_A:A \rightarrow T}{\bigoplus}
            \left ( 
                \underset{\substack{t_C:C \rightarrow T\\
                                            t_A = \beta \star t_C }}{\bigoplus}
                        w_\mathcal{T}(t_C - \beta)     
                 \right ) 
            \odot 
                w_\mathcal{T}(\set{\alpha \star - = t_A})
    \]
    The following conditions hold
    \begin{enumerate}[label=(\Alph*)]
        \item  $w_\mathcal{T}(D)=k$ if $\delta$ is weighable with $\mathcal{T}$.
        \item  $w_\mathcal{T}(D)\preceq k$ if $\delta$ is bounded above by $\mathcal{T}$  and \(w(e) \succeq 1_S\) for all $e \in \mathbb{E}$.
    \end{enumerate}
\end{lemma}
     
  
\section{Extension to Non-Well-Founded Semirings}


\subsection{Strongly monotonic measurable semiring}
\label{sec:strongly_monotonic_measurable_semiring}
A commutative semiring is an algebraic structure (see~\cite{bruggink2015proving}\cite{endrullis2024generalized_arxiv_v2}). In this work, we adapt the concept of an ordered semiring from~\cite{endrullis2024generalized_arxiv_v2}. 
The key difference is that the semiring is required to be equipped with a homomorphism to the extended real numbers instead of being well-founded.
Throughout the remainder of this article, $<$ and $\leq$ denote the canonical irreflexive and reflexive orders on the set of extended real numbers $\overline{\mathbb{R}} = \mathbb{R} \cup \{-\infty, +\infty\}$.
\begin{definition}[Strongly monotonic measurable semiring]
    \label{def:real_strongly_monotonic_semiring}
    A \textbf{strongly monotonic measurable semiring} $(S, \oplus, \odot, 0, 1, \prec, \mu)$ consists of
    \begin{itemize} 
        \item A commutative semiring $(S, \oplus, \odot, 0, 1)$,
        \item A non-empty irreflexive order $\prec$ on $S$,
        \item A homomorphism $\mu : (S, \prec) \to ( \overline{\mathbb{R}}, < )$,
    \end{itemize}
    such that $0 \neq 1$ and for all $x,y,z,w \in S$, for all $\delta \in \mathbb{R}_{\geq 0}$, we have
        \begin{align*}
            1 \preceq x \land 1 \preceq y 
            &\Rightarrow
            1 \preceq x \oplus y
            &\tag{S0} \label{ax:s0} 
            \\ 
            x \preceq x' \land y \preceq y' 
            &\Rightarrow
            x \oplus y \preceq x' \oplus y'
            &\tag{S1} \label{ax:s1} 
            \\   
            % x < y  
            % &\Rightarrow
            % x \oplus z \leq y \oplus z 
            % \tag{S1} \label{eq:ordered_semiring_plus_monotonic} 
            % \\ w
            x \prec x' \land y \prec y'  
            &\Rightarrow
            x \oplus y \prec x' \oplus y'
            &\tag{S2} \label{ax:s2} 
            \\
            \delta + \mu(x) < \mu(y) \land \delta + \mu(z) < \mu(w)
            &\Rightarrow
            \delta + \mu(x \oplus z) < \mu(y \oplus w)
            &\tag{S3} \label{ax:s2'}
            \\
            x \preceq x'
            &\Rightarrow 
            x \odot y \preceq x' \odot y 
            &\tag{S4} \label{ax:s3} 
            \\
            x \prec x' \land y \neq 0 
            &\Rightarrow
            x \odot y \prec x' \odot y
            &\tag{S5} \label{ax:s4}
            \\ 
            \delta + \mu(x) < \mu(y) \land 1 \preceq z \neq 0
            &\Rightarrow
            \delta + \mu(x \odot z) < \mu(y \odot z)
            &\tag{S6} \label{ax:s4'}
            \\
            \delta+ \mu(x) < \mu(x') \land y \neq 0
            &\Rightarrow
            \mu(x \odot y) < \mu(x' \odot y)
            &\tag{S7} \label{ax:s4''}
        %    \\
            % \\
            % 1 \leq z \neq 0 \land X < Y  
            % &\Rightarrow
            % \exists \mu(x * z) < \mu( y * z)
            % \tag{S101} \label{eq:strongly_ordered_measurable_semiring_lt_preserved_neq0_geq1}  
        %      \\     
        %     a + X < Y \land z \neq 0 
        %    &\Rightarrow
        %    \exists b> 0. b + \mu(x* z) < \mu(y * z) 
        %    \tag{S3} \label{eq:ordered_semiring_times_stable_under_mesure} 
        \end{align*}
        where $\preceq$ denotes the reflexive closure of $\prec$. The semiring is a \textbf{strictly monotonic measurable semiring} if it additionally satisfies 
    \begin{flalign*}
        \hspace{4.5cm} x \prec x' 
        &\Rightarrow
        x \oplus y \prec x' \oplus y 
        &\tag{S8} \label{ax:s5} 
        \\
        \delta + \mu(x) < \mu(x')
        &\Rightarrow
        \delta + \mu(x \oplus y) < \mu(x' \oplus y)
        &\tag{S9} \label{ax:s5'}
    \end{flalign*}
\end{definition} 
\begin{example} 
    % The real tropical semiring $\mathfrak{T}' = (\mathbb{R} \cup \{+\infty\}, \min,+,+\infty, 0_\mathbb{R},<,\operatorname{id}_{\mathbb{R} \cup \{+\infty\}})$ has domain $\mathbb{R} \cup \{+\infty\}$, the binary function symbol $\oplus$ interpreted by $\min$ and the binary function symbol $\odot$ interpreted by $+$, the constant symbols $0_s$ and $1_s$ interpreted by $+\infty$ and $0_\mathbb{R}$, respectively, the binary relation symbol $\prec$ interpreted by the canonical order $<$ on $\mathbb{R} \cup \{+\infty\}$, and the unary function symbol $\mu$ interpreted by the identity function on $\mathbb{R} \cup \{+\infty\}$. The real tropical semiring is a strongly monotonic measurable semiring. It is not strictly monotonic measurable because $2 < 3$ but $2 \oplus 2 = \min(2,2) = 2 \not < 2 = \min(3,2) = 3 \oplus 2$.

     The real tropical semiring $\mathfrak{T}' = (\mathbb{R} \cup \{+\infty\}, \min,+_\mathbb{R},+\infty, 0_\mathbb{R},<_\mathbb{R},\operatorname{id}_{\mathbb{R} \cup \{+\infty\}})$ is an instance of the strongly monotonic measurable semiring where
     \begin{flalign*}
         S & \mapsto \mathbb{R} \cup \{+\infty\}
         \\
         \oplus & \mapsto \min
         \\
         \odot & \mapsto +_\mathbb{R}
         \\
         0_s & \mapsto +\infty
         \\
         1_s & \mapsto 0_\mathbb{R}
         \\
         \prec & \mapsto <_\mathbb{R}
         \\
         \mu & \mapsto \operatorname{id}_{\mathbb{R} \cup \{+\infty\}}
     \end{flalign*}

    It is a strongly monotonic measurable semiring but not strictly monotonic measurable, because we have $2 <_\mathbb{R} 3$ but $2 \oplus 2 = \min(2,2) = 2 \not <_\mathbb{R} 2 = \min(3,2) = 3 \oplus 2$.
\end{example}
\begin{example}
    % The real arctic semiring $\mathfrak{A}' = (\mathbb{R} \cup \{-\infty\},\max,+,-\infty, 0_\mathbb{R},<,\operatorname{id}_{\mathbb{R} \cup \{-\infty\}})$ has domain $\mathbb{R} \cup \{-\infty\}$, the binary function symbol $\oplus$ interpreted by $\max$ and the binary function symbol $\odot$ interpreted by $+$, the constant symbols $0_s$ and $1_s$ interpreted by $-\infty$ and $0_\mathbb{R}$, respectively, the binary relation symbol $\prec$ interpreted by the canonical order $<$ on $\mathbb{R} \cup \{-\infty\}$, and the unary function symbol $\mu$ interpreted by the identity function on $\mathbb{R} \cup \{-\infty\}$. The real arctic semiring is a strongly monotonic measurable semiring. It is not strictly monotonic measurable because $2 < 3$ but $2 \oplus 3 = \max(2,3) = 3 \not < 3 = \max(3,3) = 3 \oplus 3$.
    The real arctic semiring $\mathfrak{A}' = (\mathbb{R} \cup \{-\infty\},\max,+_\mathbb{R},-\infty, 0_\mathbb{R},<_\mathbb{R},\operatorname{id}_{\mathbb{R} \cup \{-\infty\}})$ is an instance of the strongly monotonic measurable semiring where
    \begin{flalign*}
        S & \mapsto \mathbb{R} \cup \{-\infty\}
        \\
        \oplus & \mapsto \max
        \\
        \odot & \mapsto +_\mathbb{R}
        \\
        0_s & \mapsto -\infty
        \\
        1_s & \mapsto 0_\mathbb{R}
        \\
        \prec & \mapsto <_\mathbb{R}
        \\
        \mu & \mapsto \operatorname{id}_{\mathbb{R} \cup \{-\infty\}}
    \end{flalign*}
    It is a strongly monotonic measurable semiring but not strictly monotonic measurable, because $2 <_\mathbb{R} 3$ but $2 \oplus 3 = \max(2,3) = 3 \not <_\mathbb{R} 3 = \max(3,3) = 3 \oplus 3$.
%    The real arctic semiring: $\mathfrak{A}' = (\mathbb{R} \cup \{-\infty\},\max,+,-\infty, 0,<,\operatorname{id}_{\mathbb{R} \cup \{-\infty\}})$.
\end{example}
\begin{example}
    The real arithmetic semiring $\mathfrak{N}' = (\mathbb{R}^+,+_\mathbb{R},*_\mathbb{R},0_\mathbb{R},1_\mathbb{R},<,\operatorname{id}_{\mathbb{R}^+})$ is an instance of the strongly monotonic measurable semiring where
    \begin{flalign*}
        S & \mapsto \mathbb{R}^+
        \\
        \oplus & \mapsto +_\mathbb{R}
        \\
        \odot & \mapsto *_\mathbb{R}
        \\
        0_s & \mapsto 0_\mathbb{R}
        \\
        1_s & \mapsto 1_\mathbb{R}
        \\
        \prec & \mapsto <_\mathbb{R}
        \\
        \mu & \mapsto \operatorname{id}_{\mathbb{R}^+}
    \end{flalign*}
    It is a strictly monotonic measurable semiring. 
    % The real arithmetic semiring $\mathfrak{N}' = (\mathbb{R}^+,+,*,0_\mathbb{R},1_\mathbb{R},<,\operatorname{id}_{\mathbb{R}^+})$ has as domain the set $\mathbb{R}^+$ of positive real numbers, the binary function symbol $\oplus$ interpreted by $+$ and the binary function symbol $\odot$ interpreted by $*$, the constant symbols $0_s$ and $1_s$ interpreted by $0_\mathbb{R}$ and $1_\mathbb{R}$, respectively, the binary relation symbol $\prec$ interpreted by the canonical order $<$ on $\mathbb{R}^+$, and the unary function symbol $\mu$ interpreted by the identity function on $\mathbb{R}^+$. The real arithmetic semiring is a strictly monotonic measurable semiring. 
\end{example}
\begin{example} 
    \label{example:real_semirings}
    The natural tropical semiring $\mathfrak{T} = (\mathbb{N} \cup \{+\infty\},\min,+_{\mathbb{N}},+\infty, 0_\mathbb{N}, <_{\mathbb{N}} , \operatorname{id}_{\mathbb{N} \cup \{+\infty\}})$ is an instance of the strongly monotonic measurable semiring where
    \begin{flalign*}
        S & \mapsto \mathbb{N} \cup \{+\infty\}
        \\
        \oplus & \mapsto \min
        \\
        \odot & \mapsto +_\mathbb{N}
        \\
        0_s & \mapsto +\infty
        \\
        1_s & \mapsto 0_\mathbb{N}
        \\
        \prec & \mapsto <_\mathbb{N}
        \\
        \mu & \mapsto \operatorname{id}_{\mathbb{N} \cup \{+\infty\}}
    \end{flalign*}
    The natural arctic semiring $\mathfrak{A} = (\mathbb{N} \cup \{-\infty\},\max,+_{\mathbb{N}},-\infty, 0_\mathbb{N},<_{\mathbb{N}}, \operatorname{id}_{\mathbb{N} \cup \{-\infty\}})$ is an instance of the strongly monotonic measurable semiring where
    \begin{flalign*}
        S & \mapsto \mathbb{N} \cup \{-\infty\}
        \\
        \oplus & \mapsto \max
        \\
        \odot & \mapsto +_\mathbb{N}
        \\
        0_s & \mapsto -\infty
        \\
        1_s & \mapsto 0_\mathbb{N}
        \\
        \prec & \mapsto <_\mathbb{N}
        \\
        \mu & \mapsto \operatorname{id}_{\mathbb{N} \cup \{-\infty\}}
    \end{flalign*}  
    The natural arithmetic semiring $\mathfrak{N} = (\mathbb{N},+_\mathbb{N},*,0_\mathbb{N},1_\mathbb{N},<_\mathbb{N},\operatorname{id}_\mathbb{N})$ is an instance of the strictly monotonic measurable semiring where
    \begin{flalign*}
        S & \mapsto \mathbb{N}
        \\
        \oplus & \mapsto +_\mathbb{N}
        \\
        \odot & \mapsto *_\mathbb{N}
        \\
        0_s & \mapsto 0_\mathbb{N}
        \\
        1_s & \mapsto 1_\mathbb{N}
        \\
        \prec & \mapsto <_\mathbb{N}
        \\
        \mu & \mapsto \operatorname{id}_\mathbb{N}
    \end{flalign*}
\end{example}

% \begin{notation} 
%     \label{def:bigodot}
% Let $(S, \oplus, \odot, 0_s, 1_s)$ be a semiring. We extend naturally the binary operations $\oplus$ and $\odot$ to finite sets $E \subseteq S$ by letting
%     \begin{itemize}
%         \item $\bigodot \emptyset \overset{\operatorname{def}}{=} 1_s$ and $\bigodot \left( E \cup \{x\} \right) \overset{\operatorname{def}}{=} \left( \bigodot E \right) \odot x$;
%         \item $\bigoplus \emptyset \overset{\operatorname{def}}{=} 0_s$ and $\bigoplus \left( E \cup \{x\} \right) \overset{\operatorname{def}}{=} \left( \bigoplus E \right) \oplus x$.
%     \end{itemize}
% %  \begin{flalign*}
% %     \bigodot \emptyset &\overset{\operatorname{def}}{=} 1_s
% % \\
% %     \bigodot \left( E \cup \{x\} \right) &\overset{\operatorname{def}}{=} \left( \bigodot E \right) \odot x
% %     \\
% %     \bigoplus \emptyset &\overset{\operatorname{def}}{=} 0_s
% %     \\
% %         \bigoplus \left( E \cup \{x\} \right) &\overset{\operatorname{def}}{=} \left( \bigoplus E \right) \oplus x
% % \end{flalign*}
% \end{notation}

% \textcolor{blue}{\begin{remark}
%     \label{remark:diff_measurable_semiring}
% A strongly monotonic measurable semiring differs from a well-founded strongly monotonic semiring in~\cite{endrullis2024generalized} in four ways:
% \begin{enumerate}[label=(\arabic*),noitemsep]
%     \item Replacing well-foundedness with the homomorphism~$\mu$ and introducing Axioms \eqref{ax:s2'}, \eqref{ax:s4'}, and \eqref{ax:s4''}. This modification enables the use of non-well-founded semirings (e.g., as in~\autoref{example:real_semirings}).
%     % \item Removing the condition $1_S \preceq y$ from the original Axioms~(S3) and~(S4)\todo{This is not helpful}, resulting in our Axioms~(S4) and~(S5). This relaxation allows inclusion of elements smaller than $1_S$, as motivated in~\autoref{remark:greater_than_1}.
%     \item Adding Axiom~\eqref{ax:s0}. This technical adjustment ensures that if every $\mathcal{T}$-valued element of a type graph (see~\autoref{def:weighted_type_graph}) has a weight greater than $1_S$, then all objects subject to rewriting inherit this property.
%     \item Defining $\preceq$ as the reflexive closure of $\prec$ to simplify the theory. This is motivated by the fact that concrete semirings proposed in prior work and our paper satisfy this property.
% \end{enumerate} 
% \end{remark}}
\subsection{Measuring objects by counting morphisms}
\begin{example}
    \label{nwf:example:weighted_type_graph}
     In \textbf{Graph}, the weighted type graph $\mathcal{T} = (T, \mathbb{E}, \mathcal{S}, w)$ where $\mathcal{S}$ is the real arithmetic semiring $(\mathbb{R}^+, +, *, 0_{\mathbb{R}^+}, 1_{\mathbb{R}^+}, <, \operatorname{id}_{\mathbb{R}^+})$,
     $T$ as the graph illustrated below (without superscripts), $\mathbb{E}=\{e_{11a},e_{12a},e_{21a},e_{11b}\}$ as the set of morphism-rulers where 
     $e_{uvl}$ has domain 
     \tikz[baseline=-0.5ex]{
        \node (x) at (0,0) {$\bullet$};
        \node (y) at (1,0) {$\bullet$};
        \draw[->] (x) -- (y) node[midway, above] {$l$};
    } and image 
    \begin{tikzpicture}
        \node[draw, circle] (x) at (0,0) {$\mathrm{u}$};
        \node[draw, circle] (y) at (1,0) {$\mathrm{v}$};
        \draw[->]  (x) -- (y) node [midway,above] {$l$};
    \end{tikzpicture} in the graph $T$,
    and $w(e) = 1.0$ for all $e \in \mathbb{E}$, can be visualized as follows:
    \begin{center}
        \begin{tikzpicture}
            \graphbox{}{0mm}{0mm}{32mm}{28mm}{-10mm}{-14mm}{
                \node[draw,circle] (1) at (0,0) {1};
                \node[draw,circle] (2) at (2,0) {2};
                \draw[->] (1) edge[loop above] node[midway, above] {$a^{1.0}$} (1) ;
                \draw[->] (1) edge[loop below] node[midway, below] {$b^{1.0}$} (1) ;
                \draw[->] (1) edge[bend left] node[midway, above] {$a^{1.0}$}  (2)  ;
                \draw[->] (2) edge[bend left] node[midway, below] {$a^{1.0}$} (1)   ;
            }
        \end{tikzpicture}
    \end{center}
\end{example}

\begin{example}
    Consider the following morphisms and the weighted type graph $\mathcal{T}$ from \autoref{nwf:example:weighted_type_graph}:
    \begin{center}
        \resizebox{0.49\textwidth}{!}{
        \begin{tikzpicture}
          \graphbox{\( L \)}{-50mm}{0mm}{40mm}{39mm}{2mm}{-6mm}{
            \coordinate (o) at (0mm,-10mm); 
            \node[draw,circle] (l1) at ($(o)+(-10mm,0mm)$) {1};
            \node[draw,circle] (l2) at ($(l1)+(2,0)$) {2};
            \node[draw,circle] (l3) at ($(l1) + (1,0)$) {3};
            \draw[] (l1) -- (l3) node[midway,above] {a};
            \draw[] (l3) -- (l2) node[midway,above] {a};
        } 
            \graphbox{$T$}{0mm}{0mm}{40mm}{39mm}{-10mm}{-17mm}{
                \node[draw,circle] (1) at (0,0) {$1\ 2$};
                \node[draw,circle] (2) at (2,0) {3};
                \draw[->] (1) edge[loop above] node[midway, above] {$a^{1.0}$} (1) ;
                \draw[->] (1) edge[loop below] node[midway, below] {$b^{1.0}$} (1) ;
                \draw[->] (1) edge[bend left] node[midway, above] {$a^{1.0}$}  (2)  ;
                \draw[->] (2) edge[bend left] node[midway, below] {$a^{1.0}$} (1)   ;
            }
            \node () at (-5mm,-15mm) {$\overset{h_{11}^1}{\to}$};
        \end{tikzpicture}
        }
        \resizebox{0.49\textwidth}{!}{
            \begin{tikzpicture}
              \graphbox{\(L\)}{-50mm}{0mm}{40mm}{39mm}{2mm}{-6mm}{
                \coordinate (o) at (0mm,-10mm); 
                \node[draw,circle] (l1) at ($(o)+(-10mm,0mm)$) {1};
                \node[draw,circle] (l2) at ($(l1)+(2,0)$) {2};
                \node[draw,circle] (l3) at ($(l1) + (1,0)$) {3};
                \draw[] (l1) -- (l3) node[midway,above] {a};
                \draw[] (l3) -- (l2) node[midway,above] {a};
            } 
                \graphbox{$T$}{0mm}{0mm}{40mm}{39mm}{-10mm}{-19mm}{
                    \node[draw,circle] (1) at (0,0) {$1\ 2\ 3$};
                    \node[draw,circle] (2) at (2,0) {};
                    \draw[->] (1) edge[loop above] node[midway, above] {$a^{1.0}$} (1) ;
                    \draw[->] (1) edge[loop below] node[midway, below] {$b^{1.0}$} (1) ;
                    \draw[->] (1) edge[bend left] node[midway, above] {$a^{1.0}$}  (2)  ;
                    \draw[->] (2) edge[bend left] node[midway, below] {$a^{1.0}$} (1)   ;(1)   ;
                }
                \node () at (-5mm,-15mm) {$\overset{h_{11}^2}{\to}$};
            \end{tikzpicture}
            }
      \end{center}
    We have $ w_\mathcal{T}(h_{11}^1) = w(e_{13a})^{m_{e_{13a}}(h_{11}^1)} \odot w(e_{31a})^{m_{e_{31a}}(h_{11}^1)} =
     1.0^1 * 1.0^1 = 1.0$ and $
        w_\mathcal{T}(h_{11}^2) 
        % = 1.0^1 * 1.0^1 
        = 1.0$
\end{example}









\subsection{Estimating weights of pushout objects}
Let \( \mathcal{R} = \mathcal{A} \cup \mathcal{B} \) be a set of DPO rewriting rules and let \(\mathcal{T} = (T, \mathbb{E}, \mathcal{S}, w)\) be a weighted type graph over the strongly monotonic measurable semiring $\mathcal{S}$.

The following proposition ensures that every morphism to a type graph $(T,S,\mathbb{E}, w)$ has a weight different from $0_S$ and greater than or equal to $1_S$.
%  then the weight of a morphism to a type graph is greater than $1_S$ too. 
This will be used in \autoref{nwf:thm:termination_grs}.
\begin{proposition}  
    \label{prop_endrullis_2d7}
    Let $(S, \oplus, \odot, 0_S, 1_S, \prec, \mu)$ be a strongly monotonic measurable semiring. We have for all $x,y\in S$:
    \begin{align*}
        0_S \neq x \land 0_S \neq y 
        &\Rightarrow 0_S \neq x \odot y 
        \tag{S10} \label{eq:prop_neq0_mul_neq0}  
        \\
        1_S \preceq x \land 1_S \preceq y \land 0_S \neq x \land 0_S \neq y  
        &\Rightarrow
         1_S \preceq x \odot y 
         \tag{S11} \label{eq:prop_neg0_ge1_mul_ge1}  
         \\
         0_S \neq x \land 0_S \neq y   
         &\Rightarrow 0_S \neq x \oplus y
         \tag{S12} \label{eq:prop_neq0_plus_neq0}  
    \end{align*}
\end{proposition}

\begin{proof}
    \label{proof_prop_endrullis_2d7}
    Let $x,y \in S$ such that $x, y \neq 0_S$. By definition, $\prec$ is not empty, therefore there exist $a, b \in S$ such that $a \prec b$.
    \begin{itemize}
        \item $\prec$ is irreflexive, because $\mu: (S, \prec) \to (\overline{\mathbb{R}}, <_{\overline{\mathbb{R}}})$ is a homomorphism.
        \item \ref*{eq:prop_neq0_mul_neq0}:  
        Suppose $(x \odot y)=0_S$. 
        We have 
        \begin{flalign*}
             0_S &= a \odot 0_S & \text{$0_S$ is annihilator for $\odot$}\\
                 &= a \odot (x \odot y) &\text{by assumption on $x \odot y$}\\ 
                 &= a \odot x \odot y &\text{by associativity} \\
                 &\prec b \odot x \odot y &\text{by Axiom~\eqref{ax:s4} and $x,y\neq 0_S$}\\
                 &= b \odot (x \odot y)  &\text{by associativity}  \\
                 &= b \odot 0_S &\text{by assumption on $x \odot y$} \\
                 &= 0_S
        \end{flalign*}
         which contradicts the irreflexivity of $\prec$. 
        \item \ref*{eq:prop_neg0_ge1_mul_ge1}:
        Suppose
          $1_S \preceq x$ and $1_S \preceq y$. We have either $1_S = x$ or $1_S \prec x$. If $1_S = x$ then 
          \begin{flalign*}
            x \odot y &= 1_S \odot y & \\
                      &= y  & \\
                      & \succeq 1_S &\text{by assumption}
          \end{flalign*}
          If $1_S \prec x$ then $
        %   1_S \preceq y = 
          1_S \odot y \prec x \odot y$ by Axiom~\eqref{ax:s4} since $y \neq 0$ by assumption.
        \item \ref*{eq:prop_neq0_plus_neq0}:  
        By Axiom~\eqref{ax:s4}, we have $a \odot x \prec b \odot x$ and $a\odot y \prec b \odot y $. 
        \begin{flalign*}
            a \odot \left(x  \oplus y \right) &= \left( a \odot x \right)  \oplus \big(a \odot y \big)  & \text{by distributivity}\\
            & \prec \left(b \odot x \right)   \oplus \left( b \odot y \right)  & \text{by Axiom~\eqref{ax:s2}}\\
            & = b \odot \left(x  \oplus y \right) & \text{by distributivity}
        \end{flalign*} 
        if $x  \oplus y = 0_S$ then we have $0_S = a \odot 0_S = a \odot \left(x  \oplus y\right) \prec b \odot \left(x  \oplus y\right) = b \odot 0_S = 0_S$ which contradicts the irreflexivity of $\prec$. 
    \end{itemize}
\end{proof} 

% \begin{proposition}  
%     % Let $\mathcal{S} = (S, \oplus, \odot, 0_\mathcal{S}, 1_\mathcal{S}, \prec, \mu)$ be a strongly monotonic measurable semiring. We have 
%     For all $x,y$ in $\mathcal{S}$:
%     \begin{align*}
%         0_\mathcal{S} \neq x \land 0_\mathcal{S} \neq y 
%         &\Rightarrow 0_\mathcal{S} \neq x \odot y 
%         \\
%         1_\mathcal{S} \preceq x \land 1_\mathcal{S} \preceq y \land 0_\mathcal{S} \neq x \land 0_\mathcal{S} \neq y  
%         &\Rightarrow
%          1_\mathcal{S} \preceq x \odot y 
%          \\
%          0_\mathcal{S} \neq x \land 0_\mathcal{S} \neq y   
%          &\Rightarrow 0_\mathcal{S} \neq x \oplus y
%     \end{align*}
% \end{proposition}

    Suppose that for all $x \in \mathcal{S}$, if $ 1_\mathcal{S} \preceq x \neq 0_\mathcal{S}$ then $\mu(x) \geq_{\mathbb{R}^+} \mu(1_\mathcal{S})$ and $\mu(x) \in \mathbb{R}$. 
    By definition of a weighted type graph, $1_\mathcal{S} \leq w(e)\neq 0_\mathcal{S}$ for all $e\in\mathbb{E}$.
    Therefore, by the proposition above and the definition of the weight of a morphism (\autoref{def:weight_of_a_morphism_relative_to_a_type_graph}), for all morphisms \( h \) from a finite graph $G$ to $\mathcal{T}$, we have \( 0_\mathcal{S} \neq w_\mathcal{T}(h) \succeq 1_\mathcal{S} \). 
    By the above proposition, Axiom \eqref{ax:s0} and the definition of the weight of an object (\autoref{def:weight_of_an_object_relative_to_a_type_graph}), we conclude that for all objects \( G \) with $\operatorname{Hom}(G,T)\neq \emptyset$, \( 0_\mathcal{S} \neq w_\mathcal{T}(G) \succeq 1_\mathcal{S} \). By assumption, we have that for all objects \( G \) with $\operatorname{Hom}(G,T)\neq \emptyset$, \( \mu(w_\mathcal{T}(G)) \geq_{\mathbb{R}^+} \mu(1_\mathcal{S}) = 0 \) if $\mathcal{S}$ is the real tropical or arctic semiring, and \( \mu(w_\mathcal{T}(G)) \geq_{\mathbb{R}^+} 1 \) if $\mathcal{S}$ is the arithmetic semiring.

Suppose additionally that for every graph $G$ subject to rewriting, we have $\operatorname{Hom}(G,T)\neq \emptyset$. The rewriting relation \( \Rightarrow_{\mathcal{A},\mathfrak{F}} \) is terminating relative to $\Rightarrow_{\mathcal{B},\mathfrak{F}}$ if there exists a constant $\delta > 0$ such that: (i) for all \(G \Rightarrow_{\mathcal{A},\mathfrak{F}} H\), \( \mu(w_\mathcal{T}(G)) >_{\mathbb{R}^+} \mu(w_\mathcal{T}(H)) + \delta \), and (ii) for all \(G \Rightarrow_{\mathcal{B},\mathfrak{F}} H\), \( \mu(w_\mathcal{T}(G)) \geq \mu(w_\mathcal{T}(H)) \). However, directly verifying all rewriting steps is infeasible due to the potentially infinite number of rewriting steps.


\noindent
\begin{minipage}{0.7\textwidth}\setlength{\parindent}{1em}
    Consider a DPO diagram (illustrated on the right) that defines a rewriting step \( G \Rightarrow_{\rho,\mathfrak{F}} H \). 
    If the left pushout square is weighable with $\mathcal{T}$ and the right pushout square is bounded above by $\mathcal{T}$ (concepts defined in \autoref{def:weighable}), then, by \cite[Lemma 4.13]{endrullis2024generalized_arxiv_v2}, we obtain:
\end{minipage}%
\hfill
\begin{minipage}{0.29\textwidth}
    \hfill
    \resizebox{0.85\textwidth}{!}{
    \begin{tikzpicture}
        % [node distance=11mm]
        \node (I) at (0,0) {$K$};
        \node (L) at (-2,0) {$L$};
        \node (R) at (2,0) {$R$};
        \node (G) at (-2,-2) {$G$};
        \node (C) at (0,-2) {$C$};
        \node (H) at (2,-2) {$H$};
        \draw [->] (I) to  node [midway,below] {$l$} (L);
        \draw [->] (I) to  node [midway,below] {$r$} (R);
        \draw [->] (L) to node [midway,right] {$m$} (G);
        \draw [->] (I) to node [midway,right] {$u$} (C);
        \draw [->] (R) to node [midway,left] {$m'$} (H);
        \draw [->] (C) to node [midway,above] {$l'$} (G);
        \draw [->] (C) to node [midway,above] {$r'$} (H);
        \node [at=($(I)!.5!(G)$)] {\normalfont PO};
        \node [at=($(I)!.5!(H)$)] {\normalfont PO};
      \end{tikzpicture}
    }
\end{minipage}

% \begin{flalign*}
%     w_\mathcal{T}(G) 
%         & \overset{def}{=} \bigoplus_{t_G: G \rightarrow T} w_\mathcal{T}(t_G) 
%         \overset{def}{=} \bigoplus_{t_K: K \rightarrow T} \bigoplus_{\substack{t_G: G \rightarrow T\\ t_K = l \star m \star t_G}} w_\mathcal{T}(t_G) \\
%     w_\mathcal{T}(H) 
%         &\overset{def}{=} \bigoplus_{t_H: H \rightarrow T} w_\mathcal{T}(t_H) 
%          \overset{def}{=} \bigoplus_{t_K: K \rightarrow T} \bigoplus_{\substack{t_H: H \rightarrow T\\ t_K = r \star m' \star t_H}} w_\mathcal{T}(t_H)
% \end{flalign*}

% If the left pushout square is weighable with $\mathcal{T}$ and the right pushout square is bounded above by $\mathcal{T}$ (concepts defined in \autoref{def:weighable}), then, by \cite[Lemma 5.10]{endrullis2024generalized_lmcs}, we derive the following assertions:
% \begin{flalign*}
%     w_\mathcal{T}(G) 
%         & = \bigoplus_{t_K: K \rightarrow T} \bigoplus_{\substack{t_G: G \rightarrow T\\ t_K = l \star m \star t_G}} w_\mathcal{T}(m \star t_G) \odot (l' \star t_G - u) \\
%     w_\mathcal{T}(H) 
%         &  \preceq \bigoplus_{t_K: K \rightarrow T} \bigoplus_{\substack{t_H: H \rightarrow T\\ t_K = r \star m' \star t_H}} w_\mathcal{T}(m' \star t_H) \odot (r' \star t_H - u)
% \end{flalign*}
% By \cite[Theorem 5.12]{endrullis2024generalized_lmcs}, it follows:
% \begin{flalign*}
%     w_\mathcal{T}(G) 
%         & = \bigoplus_{t_K: K \rightarrow T} 
%         \bigoplus_{\substack{t_C: C \rightarrow T\\ t_K = u \star t_C}}
%         \bigoplus_{\substack{t_L: L \rightarrow T\\ t_K = l \star t_L}}
%             w_\mathcal{T}(t_L) \odot w_\mathcal{T}(t_C - u) \\
%     w_\mathcal{T}(H) 
%         &  \preceq \bigoplus_{t_K: K \rightarrow T} 
%         \bigoplus_{\substack{t_C: C \rightarrow T\\ t_K = u \star t_C}}
%         \bigoplus_{\substack{t_R: R \rightarrow T\\ t_K = r \star t_R}}
%             w_\mathcal{T}(t_R) \odot w_\mathcal{T}(t_C - u) \\
% \end{flalign*}
% By \cite[Theorem 5.13]{endrullis2024generalized_lmcs}, we have the following assertions: 
\begin{flalign*}
    w_\mathcal{T}(G) 
        & = \bigoplus_{t_K: K \rightarrow T} 
        \left ( \bigoplus_{\substack{t_C: C \rightarrow T\\ t_K = u \star t_C}}
          w_\mathcal{T}(t_C - u) \right ) 
          \odot
        \left (\bigoplus_{\substack{t_L: L \rightarrow T\\ t_K = l \star t_L}}
        w_\mathcal{T}(t_L) \right )
         \\
    w_\mathcal{T}(H) 
        &  \preceq \bigoplus_{t_K: K \rightarrow T} 
        \left ( \bigoplus_{\substack{t_C: C \rightarrow T\\ t_K = u \star t_C}}
         w_\mathcal{T}(t_C - u) \right ) 
         \odot 
         \left ( \bigoplus_{\substack{t_R: R \rightarrow T\\ t_K = r \star t_R}}
            w_\mathcal{T}(t_R) \right ) \\
\end{flalign*}
According to these results, a termination criterion will be established by comparing the weights 
$\bigoplus_{\substack{t_L: L \rightarrow T\\ t_K = l \star t_L}}
        w_\mathcal{T}(t_L)$ and 
$\bigoplus_{\substack{t_R: R \rightarrow T\\ t_K = r \star t_R}} 
        w_\mathcal{T}(t_R)$ for every $t_K: K \rightarrow T$.












        
\subsection{Decreasing rules}
\label{nwf:sec:decreasing_rules}
The definition below adapts the concept of decreasing rules by Endrullis and Overbeek~\cite{endrullis2024generalized_arxiv_v2}. The difference lies in the treatment of weight comparisons for uniformly decreasing and closure decreasing rules: instead of directly comparing weights, we evaluate the difference between weights relative to the homomorphism $\mu$ of the strongly monotonic measurable semirings.
%  This difference must exceed a fixed positive constant $\delta \in \mathbb{R}_{>0}$.

\begin{definition} 
    \label{nwf:def:decreasing_rule}
    Let $\mathcal{T} = (T,\mathbb{E}, (S, \oplus, \odot, 0_S, 1_S, \prec, \mu), w)$ be a finitary weighted type graph, \(\mathfrak{F}\) a DPO rewriting framework, $\rho = (L \overset{l}{\leftarrow} K \overset{r}{\rightarrow} R)$ a DPO rewriting rule, and $\delta \in \mathbb{R}_{>0}$. 

    \noindent
    The rule $\rho$ is \textbf{weakly decreasing} with respect to $\mathcal{T}$ in $\mathfrak{F}$ if 
            for every $t_K : K \to T$,
                $$ 
                  w_\mathcal{T}(\{l \star - = t_K\}) \succeq w_\mathcal{T}(\{r\star - = t_K\}).$$
           
    \noindent
    The rule $\rho$ is \textbf{$\delta$-uniformly decreasing} with respect to $\mathcal{T}$ in $\mathfrak{F}$ if the following hold:
        \begin{itemize}
            \item[]- there is a context closure $c_\rho$ for $\rho$ and $\mathcal{T}$ in $\mathfrak{F}$, 
            \item[]- for every $t_K : K \to T$,
            \begin{itemize}
                \item[] $\bullet$ $\{l \star - = t_K\} = \emptyset = \{r \star - = t_K\}$, or
                \item[] $\bullet$ $\mu(w_\mathcal{T}(\{l \star - = t_K\}))  >_{\mathbb{R}^+}   \mu(w_\mathcal{T}(\{r \star - = t_K\})) + \delta$.
            \end{itemize}
        \end{itemize}  
         
    \noindent
    The rule $\rho$ is
            \textbf{$\delta$-closure decreasing} with respect to $\mathcal{T}$ in $\mathfrak{F}$ if the following hold:
            \begin{itemize}
                \item[]- $S$ is strictly monotonic measurable semiring,
                \item[]- $\rho$ is weakly decreasing,
                \item[]- there is a context closure $c_\rho$ for $\rho$ and $\mathcal{T}$ in $\mathfrak{F}$,
                \item[]- $ \mu(w_\mathcal{T}(\{l \star - = t_K\}))  >_{\mathbb{R}^+}  \mu(w_\mathcal{T}(\{r \star - = t_K\}))  + \delta$ for $t_K = l \star c_\rho$.
            \end{itemize}
\end{definition}

\begin{example}
    \label{nwf:example:decreasing_rule}
    Consider the following DPO rule:
    % shown in Figure~\ref{fig:nwf:rewriting_step_tracsdfdsafdfsd}
     \begin{center}
      \resizebox{0.7\textwidth}{!}{
      \begin{tikzpicture}
          \graphbox{$L$}{0mm}{0mm}{34mm}{15mm}{2mm}{-5mm}{
              \coordinate (o) at (0mm,-3mm); 
              \node[draw,circle] (l1) at ($(o)+(-10mm,0mm)$) {1};
              \node[draw,circle] (l2) at ($(l1)+(2,0)$) {2};
              \node[draw,circle] (l3) at ($(l1) + (1,0)$) {3};
              \draw[->] (l1) -- (l3) node[midway,above] {$a$};
              \draw[->] (l3) -- (l2) node[midway,above] {$a$};
          }     
          \graphbox{$K$}{40mm}{0mm}{24mm}{15mm}{2mm}{-5mm}{
              \coordinate (o) at (5mm,-3mm); 
              \node[draw,circle] (l1) at ($(o)+(-10mm,0mm)$) {1};
              \node[draw,circle] (l2) at ($(l1)+(1,0)$) {2};
              % \node[draw,circle] (l3) at ($(l1) + (1,0)$) {$\ $};
              % \draw[->] (l1) -- (l3) node[midway,above] {$a$};
              % \draw[->] (l3) -- (l2) node[midway,above] {$a$};
          }    
          \graphbox{$R$}{70mm}{0mm}{45mm}{15mm}{2mm}{-5mm}{
              \coordinate (o) at (-5mm,-3mm); 
              \node[draw,circle] (l1) at ($(o)+(-10mm,0mm)$) {1};
              \node[draw,circle] (l2) at ($(l1)+(3,0)$) {2};
              \node[draw,circle] (l3) at ($(l1) + (1,0)$) {4};
              \node[draw,circle] (l4) at ($(l1) + (2,0)$) {5};
              \draw[->] (l1) -- (l3) node[midway,above] {$a$};
              \draw[->] (l3) -- (l4) node[midway,above] {$b$};
              \draw[->] (l4) -- (l2) node[midway,above] {$a$};
          }    
          \node () at (37mm,-8mm) {$\overset{l}{\leftarrowtail}$};
          \node () at (67mm,-8mm) {$\overset{r}{\rightarrowtail}$};
          % \draw[>->] (51mm,2mm) -- (52mm,3mm);
      \end{tikzpicture}
      }
  \end{center}
   and the following weighted type graph
    % in Example~\ref{nwf:example:weighted_type_graphd_grsaa_sdfsf} 
    over the real arithmetic semiring $\mathfrak{N}' = (\mathbb{R}^+,+,*,0_\mathbb{R},1_\mathbb{R},<,\operatorname{id}_{\mathbb{R}^+})$:
    \begin{center}
        \begin{tikzpicture}
            \graphbox{}{0mm}{0mm}{32mm}{28mm}{-10mm}{-14mm}{
                \node[draw,circle] (1) at (0,0) {};
                \node[draw,circle] (2) at (2,0) {};
                \draw[->] (1) edge[loop above] node[midway, above] {$a^{1.0}$} (1) ;
                \draw[->] (1) edge[loop below] node[midway, below] {$b^{1.0}$} (1) ;
                \draw[->] (1) edge[bend left] node[midway, above] {$a^{1.0}$}  (2)  ;
                \draw[->] (2) edge[bend left] node[midway, below] {$a^{1.0}$} (1)   ;
            }
        \end{tikzpicture}
    \end{center}
     There are 4 morphisms $t_K^{11}, t_K^{12}, t_K^{21}, t_K^{22}$ from $K$ to $T$:
    \begin{center} 
        \resizebox{0.6\textwidth}{!}{
            \begin{tikzpicture}
            \graphbox{\( K \)}{-50mm}{0mm}{40mm}{30mm}{2mm}{-6mm}{
                \coordinate (o) at (0mm,-10mm); 
                \node[draw,circle] (l1) at ($(o)+(-10mm,0mm)$) {1};
                \node[draw,circle] (l2) at ($(l1)+(2,0)$) {2};
                % \node[draw,circle] (l3) at ($(l1) + (1,0)$) {3};
                % \draw[] (l1) -- (l3) node[midway,above] {$a$};
                % \draw[] (l3) -- (l2) node[midway,above] {$a$};
            } 
                \graphbox{$T$}{0mm}{0mm}{40mm}{30mm}{-10mm}{-15mm}{
                    \node[draw,circle] (1) at (0,0) {$1\ 2$};
                    \node[draw,circle] (2) at (2,0) {};
                    \draw[->] (1) edge[loop above] node[midway, above] {$a$} (1) ;
                    \draw[->] (1) edge[loop below] node[midway, below] {$b$} (1) ;
                    \draw[->] (1) edge[bend left] node[midway, above] {$a$}  (2)  ;
                    \draw[->] (2) edge[bend left] node[midway, below] {$a$} (1)   ;
                }
                \node () at (-5mm,-15mm) {$\overset{t_K^{11}}{\to}$};
            \end{tikzpicture}
            } 

            \resizebox{0.6\textwidth}{!}{
            \begin{tikzpicture}
                \graphbox{\( K \)}{-50mm}{0mm}{40mm}{30mm}{2mm}{-6mm}{
                \coordinate (o) at (0mm,-10mm); 
                \node[draw,circle] (l1) at ($(o)+(-10mm,0mm)$) {1};
                \node[draw,circle] (l2) at ($(l1)+(2,0)$) {2};
                % \node[draw,circle] (l3) at ($(l1) + (1,0)$) {3};
                % \draw[] (l1) -- (l3) node[midway,above] {$a$};
                % \draw[] (l3) -- (l2) node[midway,above] {$a$};
            } 
                \graphbox{$T$}{0mm}{0mm}{40mm}{30mm}{-10mm}{-15mm}{
                    \node[draw,circle] (1) at (0,0) {$1$};
                    \node[draw,circle] (2) at (2,0) {2};
                    \draw[->] (1) edge[loop above] node[midway, above] {$a$} (1) ;
                    \draw[->] (1) edge[loop below] node[midway, below] {$b$} (1) ;
                    \draw[->] (1) edge[bend left] node[midway, above] {$a$}  (2)  ;
                    \draw[->] (2) edge[bend left] node[midway, below] {$a$} (1)   ;
                }
                \node () at (-5mm,-15mm) {$\overset{t_K^{12}}{\to}$};
            \end{tikzpicture}
            }
            
            \resizebox{0.6\textwidth}{!}{
            \begin{tikzpicture}
                \graphbox{\( K \)}{-50mm}{0mm}{40mm}{30mm}{2mm}{-6mm}{
                \coordinate (o) at (0mm,-10mm); 
                \node[draw,circle] (l1) at ($(o)+(-10mm,0mm)$) {1};
                \node[draw,circle] (l2) at ($(l1)+(2,0)$) {2};
                % \node[draw,circle] (l3) at ($(l1) + (1,0)$) {3};
                % \draw[] (l1) -- (l3) node[midway,above] {$a$};
                % \draw[] (l3) -- (l2) node[midway,above] {$a$};
            } 
                \graphbox{$T$}{0mm}{0mm}{40mm}{30mm}{-10mm}{-15mm}{
                    \node[draw,circle] (1) at (0,0) {2};
                    \node[draw,circle] (2) at (2,0) {1};
                    \draw[->] (1) edge[loop above] node[midway, above] {$a$} (1) ;
                    \draw[->] (1) edge[loop below] node[midway, below] {$b$} (1) ;
                    \draw[->] (1) edge[bend left] node[midway, above] {$a$}  (2)  ;
                    \draw[->] (2) edge[bend left] node[midway, below] {$a$} (1)   ;
                }
                \node () at (-5mm,-15mm) {$\overset{t_K^{21}}{\to}$};
            \end{tikzpicture}
            }

            \resizebox{0.6\textwidth}{!}{
            \begin{tikzpicture}
                \graphbox{\( K \)}{-50mm}{0mm}{40mm}{30mm}{2mm}{-6mm}{
                \coordinate (o) at (0mm,-10mm); 
                \node[draw,circle] (l1) at ($(o)+(-10mm,0mm)$) {1};
                \node[draw,circle] (l2) at ($(l1)+(2,0)$) {2};
                % \node[draw,circle] (l3) at ($(l1) + (1,0)$) {3};
                % \draw[] (l1) -- (l3) node[midway,above] {$a$};
                % \draw[] (l3) -- (l2) node[midway,above] {$a$};
            } 
                \graphbox{$T$}{0mm}{0mm}{40mm}{30mm}{-10mm}{-15mm}{
                    \node[draw,circle] (1) at (0,0) {};
                    \node[draw,circle] (2) at (2,0) {$1\ 2$};
                    \draw[->] (1) edge[loop above] node[midway, above] {$a$} (1) ;
                    \draw[->] (1) edge[loop below] node[midway, below] {$b$} (1) ;
                    \draw[->] (1) edge[bend left] node[midway, above] {$a$}  (2)  ;
                    \draw[->] (2) edge[bend left] node[midway, below] {$a$} (1)   ;
                }
                \node () at (-5mm,-15mm) {$\overset{t_K^{22}}{\to}$};
            \end{tikzpicture}
            }
            % \caption{}
            % \label{fig:nwf:example:decreasing_rule_morphisms_grsaa_sfsdjl}
      \end{center}
    The set $\{l \star - = t_K^{11}\}$ has two morphisms $h_{11}^1$ and $h_{11}^2$:
    \begin{center}
        \resizebox{0.6\textwidth}{!}{
        \begin{tikzpicture}
          \graphbox{\( L \)}{-50mm}{0mm}{40mm}{40mm}{2mm}{-6mm}{
            \coordinate (o) at (0mm,-10mm); 
            \node[draw,circle] (l1) at ($(o)+(-10mm,0mm)$) {1};
            \node[draw,circle] (l2) at ($(l1)+(2,0)$) {2};
            \node[draw,circle] (l3) at ($(l1) + (1,0)$) {3};
            \draw[] (l1) -- (l3) node[midway,above] {$a$};
            \draw[] (l3) -- (l2) node[midway,above] {$a$};
        } 
            \graphbox{$T$}{0mm}{0mm}{40mm}{40mm}{-10mm}{-17mm}{
                \node[draw,circle] (1) at (0,0) {$1\ 2$};
                \node[draw,circle] (2) at (2,0) {3};
                \draw[->] (1) edge[loop above] node[midway, above] {$a^{1.0}$} (1) ;
                \draw[->] (1) edge[loop below] node[midway, below] {$b^{1.0}$} (1) ;
                \draw[->] (1) edge[bend left] node[midway, above] {$a^{1.0}$}  (2)  ;
                \draw[->] (2) edge[bend left] node[midway, below] {$a^{1.0}$} (1)   ;
            }
            \node () at (-5mm,-15mm) {$\overset{h_{11}^1}{\to}$};
        \end{tikzpicture}
        }

        \resizebox{0.6\textwidth}{!}{
            \begin{tikzpicture}
              \graphbox{\(L\)}{-50mm}{0mm}{40mm}{40mm}{2mm}{-10mm}{
                \coordinate (o) at (0mm,-10mm); 
                \node[draw,circle] (l1) at ($(o)+(-10mm,0mm)$) {1};
                \node[draw,circle] (l2) at ($(l1)+(2,0)$) {2};
                \node[draw,circle] (l3) at ($(l1) + (1,0)$) {3};
                \draw[] (l1) -- (l3) node[midway,above] {$a$};
                \draw[] (l3) -- (l2) node[midway,above] {$a$};
            } 
                \graphbox{$T$}{0mm}{0mm}{40mm}{40mm}{-10mm}{-20mm}{
                    \node[draw,circle] (1) at (0,0) {$1\ 2\ 3$};
                    \node[draw,circle] (2) at (2,0) {};
                    \draw[->] (1) edge[loop above] node[midway, above] {$a^{1.0}$} (1) ;
                    \draw[->] (1) edge[loop below] node[midway, below] {$b^{1.0}$} (1) ;
                    \draw[->] (1) edge[bend left] node[midway, above] {$a^{1.0}$}  (2)  ;
                    \draw[->] (2) edge[bend left] node[midway, below] {$a^{1.0}$} (1)   ;(1)   ;
                }
                \node () at (-5mm,-15mm) {$\overset{h_{11}^2}{\to}$};
            \end{tikzpicture}
            }
        \end{center}
    Therefore, we have \begin{flalign*}
        w_\mathcal{T}(\{l \star - = t_K^{11}\})
        =&w_\mathcal{T}(\{h_{11}^1, h_{11}^2\})\\
        =&w_\mathcal{T}(h_{11}^1) + w_\mathcal{T}(h_{11}^2) \\
        =&(1.0^1 * 1.0^1) + (1.0^1 * 1.0^1)\\
        =&2.0.
    \end{flalign*}
    The set $\{r \star - = t_K^{11}\}$ has one morphism $h_{11}^3$ which is illustrated below:
    \begin{center}
        \resizebox{0.49\textwidth}{!}{
        \begin{tikzpicture}
          \graphbox{\( R \)}{-55mm}{0mm}{45mm}{44mm}{1mm}{-22mm}{
            \coordinate (o) at (-5mm,-3mm); 
            \node[draw,circle] (l1) at ($(o)+(-10mm,0mm)$) {1};
            \node[draw,circle] (l2) at ($(l1)+(3,0)$) {2};
            \node[draw,circle] (l3) at ($(l1) + (1,0)$) {4};
            \node[draw,circle] (l4) at ($(l1) + (2,0)$) {5};
            \draw[->] (l1) -- (l3) node[midway,above] {$a$};
            \draw[->] (l3) -- (l4) node[midway,above] {$b$};
            \draw[->] (l4) -- (l2) node[midway,above] {$a$};
        } 
            \graphbox{$T$}{0mm}{0mm}{40mm}{44mm}{-10mm}{-22mm}{
                \node[draw,circle] (1) at (0,0) {$1\ 2\ 4\ 5$};
                \node[draw,circle] (2) at (2,0) {};
                \draw[->] (1) edge[loop above] node[midway, above] {$a^{1.0}$} (1) ;
                \draw[->] (1) edge[loop below] node[midway, below] {$b^{1.0}$} (1) ;
                \draw[->] (1) edge[bend left] node[midway, above] {$a^{1.0}$}  (2)  ;
                \draw[->] (2) edge[bend left] node[midway, below] {$a^{1.0}$} (1)   ;
            }
            \node () at (-5mm,-19mm) {$\overset{h_{11}^3}{\to}$};
        \end{tikzpicture}
        }
      \end{center}
    Therefore, we have: $$w_\mathcal{T}(\{r \star - = t_K^{11}\}) = w_\mathcal{T}(h_{11}^3) = 1.0^1 * 1.0^1 * 1.0 ^ 1 = 1.0.$$
    The following inequality holds:
    $$w_\mathcal{T}(\{l \star - = t_K^{11}\}) = 2.0 \geq 1.0 = w_\mathcal{T}(\{r \star - = t_K^{11}\}).$$
    Similarly, we can check that 
        \begin{itemize}
            \item $w_\mathcal{T}(\{l \star - = t_K^{12}\}) = 1.0 \geq 1.0 = w_\mathcal{T}(\{r \star - = t_K^{12}\})$, and
            \item $w_\mathcal{T}(\{l \star - = t_K^{21}\}) = 1.0 \geq 1.0 = w_\mathcal{T}(\{r \star - = t_K^{21}\})$, and
            \item $w_\mathcal{T}(\{l \star - = t_K^{22}\}) = 1.0 \geq 1.0 = w_\mathcal{T}(\{r \star - = t_K^{22}\})$.
        \end{itemize}
    Therefore, the rule is therefore weakly decreasing.
    
    Moreover, there is a context closure $c$ for the DPO rule in the weighted type graph, as explained in Example~\ref{wf:example:context_closure}, which is illustrated below:
    \begin{center} 
    \resizebox{0.6\textwidth}{!}{
    \begin{tikzpicture}
      \graphbox{\( L \)}{-50mm}{0mm}{40mm}{40mm}{2mm}{-8mm}{
        \coordinate (o) at (0mm,-10mm); 
        \node[draw,circle] (l1) at ($(o)+(-10mm,0mm)$) {1};
        \node[draw,circle] (l2) at ($(l1)+(2,0)$) {2};
        \node[draw,circle] (l3) at ($(l1) + (1,0)$) {3};
        \draw[] (l1) -- (l3) node[midway,above] {$a$};
        \draw[] (l3) -- (l2) node[midway,above] {$a$};
    } 
        \graphbox{$T$}{0mm}{0mm}{40mm}{40mm}{-10mm}{-17mm}{
            % \node[draw,circle] (1) at (0,0) {$1\ 2\ 3$};
            % \node[draw,circle] (2) at (2,0) {};
            \coordinate (o) at (2mm,-3mm); 
            \node[draw,circle] (1) at ($(o)+(0,0mm)$) {$1\ 2\ 3$};
            \node[draw,circle] (2) at ($(o)+(2,0)$) {};
            \draw[->] (1) edge[loop above] node[midway, above] {$a^{1}$} (1) ;
            \draw[->] (1) edge[loop below] node[midway, below] {$b^{1}$} (1) ;
            \draw[->] (1) edge[bend left] node[midway, above] {$a^{1}$}  (2)  ;
            \draw[->] (2) edge[bend left] node[midway, below] {$a^{1}$} (1)   ;
        }
        \node () at (-5mm,-15mm) {$\overset{c}{\to}$};
    \end{tikzpicture}
    }
  \end{center}
    Additionally, we have 
    \begin{itemize}
        \item $t_K^{11} = l \star c$, and 
        \item $w_\mathcal{T}(\{l \star - = t_K^{11}\}) = 2.0 > 1.0 + \delta = w_\mathcal{T}(\{r \star - = t_K^{11}\}) + \delta$ for $0 < \delta < 1.0$.
    \end{itemize} 
    Thus, we conclude that the rule is $\delta$-closure decreasing with $0 < \delta < 1.0$ since the semiring is strictly monotonic measurable.
\end{example} 

The following lemma by Endrullis and Overbeek ~\cite{endrullis2024generalized_arxiv_v2} guarantees that, under certain constraints: exact weights of host graphs can be computed, and upper bounds for result-graph weights can be derived. 
\begin{lemma}
    \label{nwf:lem_4d13}
\ \newline
\begin{minipage}{0.7\textwidth}
    Let $\mathcal{T} = (T,\mathbb{E}, (S, \oplus, \odot, 0_S, 1_S, \prec, \mu), w)$ be a finitary weighted type graph. Consider the pushout square $\delta$ illustrated on the right. We define
\end{minipage}
\begin{minipage}{0.3\textwidth}
    \begin{center}{\normalfont
        \begin{tikzpicture}[node distance=12mm]
            \node (A) {$A$};
            \node (B) [right of=A] {$B$};
            \node (C) [below of=A] {$C$};
            \node (D) [right of=C] {$D$};
            
            \draw [->] (A) to node [above, label] {$\alpha$} (B);
            \draw [->] (A) to node [left, label] {$\beta$} (C);
            \draw [->] (B) to node [right, label] {$\beta'$} (D);
            \draw [->] (C) to node [below, label] {$\alpha'$} (D);
            
            \node [at=($(A)!.5!(D)$)] {$\delta$};
        \end{tikzpicture}
    }\end{center}
\end{minipage}
     \[k = \underset{t_A:A \rightarrow T}{\bigoplus}
            \left ( 
                \underset{\substack{t_C:C \rightarrow T\\
                                            t_A = \beta \star t_C }}{\bigoplus}
                        w_\mathcal{T}(t_C - \beta)     
                 \right ) 
            \odot 
                w_\mathcal{T}(\set{\alpha \star - = t_A})
    \]
    The following conditions hold
    \begin{enumerate}[label=(\Alph*)]
        \item  $w_\mathcal{T}(D)=k$ if $\delta$ is weighable with $\mathcal{T}$.
        \item  $w_\mathcal{T}(D)\preceq k$ if $\delta$ is bounded above by $\mathcal{T}$  and \(w(e) \succeq 1_S\) for all $e \in \mathbb{E}$.
    \end{enumerate}
\end{lemma}

Adapted from \cite[Theorem C.3]{endrullis2024generalized_arxiv_v2} by Endrullis and Overbeek, the lemma below relates decreasing rules to weight reduction under specific constraints. \todo{todo to do: ref 3.30 incorrect}
\begin{lemma}[Decreasing steps]
    \label{nwf:lem:decreasing_step}
    Let $\mathcal{T} = (T,\mathbb{E}, (S, \oplus, \odot, 0_S, 1_S, \prec, \mu), w)$ be a finitary weighted type graph, $\rho$ a rewriting rule and $\Delta \in \mathfrak{F}(\rho)$ the following DPO diagram
    \begin{center}
        \resizebox{0.3\textwidth}{!}{
            \begin{tikzpicture}[node distance=11mm]
            \node (I) {$K$};
            \node (L) [left of= I] {$L$};
            \node (R) [right of=I] {$R$}; 
            \node (G) [below of=L] {$G$};
            \node (C) [below of=I] {$C$};
            \node (H) [below of=R] {$H$};
            %   \node (T) [left=of $(L)!0.5!(G)$] {$T$};
            %   \draw [->] (L) to  node [label, above] {$c$}  (T);
            %   \draw [->] (G) to  node [label, below] {$\alpha$} (T);
            \draw [->] (I) to node [label, below] {$l$} (L);
            \draw [->] (I) to node [label, below] {$r$} (R);
            \draw [->] (L) to  (G);
            \draw [->] (I) to (C);
            \draw [->] (R) to (H);
            \draw [->] (C) to (G);
            \draw [->] (C) to (H);
            \end{tikzpicture}
        }
      \end{center}
    such that the following conditions hold:
 
   \begin{itemize}
       \item $\operatorname{left}(\Delta)$ is weighable with \(\mathcal{T}\), and
       \item $\operatorname{right}(\Delta)$ is bounded above by \(\mathcal{T}\), and
       \item $w(e) \succeq 1_S$ for all $e \in \mathbb{E}$.
   \end{itemize}

   \noindent
  We have:
   \begin{itemize}
       \item $\mu(w_\mathcal{T}(G)) \succeq \mu(w_\mathcal{T}(H))$ if $\rho$ is weakly decreasing,
       \item $\mu(w_\mathcal{T}(G)) >_{\mathbb{R}^+} \mu(w_\mathcal{T}(H)) + \delta$ if $\rho$ is $\delta$-uniformly or $\delta$-closure decreasing for some $\delta >_{\mathbb{R}^+} 0$ and $w(e) \succeq 1_S$ for all $e \in \mathbb{E}$.
   \end{itemize}
\end{lemma} 

\begin{proof}
    \label{nwf:proof:decreasing_step}
    \noindent For every \( t_K: K \rightarrow T \), we define
$
        S_{t_K} \overset{\operatorname{def}}{=}   
        \underset{\substack{t_C:C \rightarrow T \\
        t_K = h_{KC} \star t_C }}{\bigoplus} 
        w_\mathcal{T}(t_C - h_{KC})  
$.
    
    \noindent For all $t_K: K \to T$ and $X,Y \in S$, the following claims hold:
    \begin{enumerate}[label=(\alph*)] 
        \item \label{s_nz} $S_{t_K} \ne 0_S$ if there is $t_C$ with $ t_K = h_{KC} \star t_C$.  
        \begin{proof}
            By definition of weighted type graph, for all $e \in \mathbb{E}$, we have 
            \begin{flalign}
                w(e) \neq 0_S \label{eq_we_neq_0s1111}
            \end{flalign}
            For every $t_C:C \to T$, we have 
            \begin{flalign*}
                &w_\mathcal{T}(t_C - h_{KC}) \\
               =&\bigodot_{e\in \mathbb{E}} w_e(t_C - h_{KC}) 
                 \hspace{5.5cm}\text{by Definition~\ref{def:weight_excluding}}\\
               =&\bigodot_{e\in \mathbb{E}} 
                 \bigodot_{\substack{\alpha \in \{- * t_C = e\}\\
                    \alpha \notin \left\{ \iota \in \operatorname{Hom}(X, C)~\middle|~\exists \zeta:X \to K,~\zeta \star h_{KC} = \iota \right\}
                 }
                 } w(e)  \hspace{2cm} \text{by Definition~\ref{def:weight_excluding_pre}} 
                 \\
               \neq&0_S \hspace{1.3cm}\text{by Equation~\eqref{eq_we_neq_0s1111}, Proposition~\ref{prop_endrullis_2d7}~\eqref{eq:prop_neq0_mul_neq0} and Definition~\ref{def:bigodot}}  
            \end{flalign*}

            Therefore, $S_{t_K} \overset{\operatorname{def}}{=}   
            \underset{\substack{t_C:C \rightarrow T \\
            t_K = h_{KC} \star t_C }}{\bigoplus} 
            w_\mathcal{T}(t_C - h_{KC}) \neq 0_S$ if there is be a morphism such that $t_K = h_{KC} \star t_C$ by Proposition~\ref{prop_endrullis_2d7}~\eqref{eq:prop_neq0_plus_neq0}.
            % For every $t_C:C \to T$ such that $t_K = h_{KC} \star t_C$, we have $w_\mathcal{T}(t_C - h_{KC}) \neq 0_S$ by \eqref{eq:prop_neq0_mul_neq0}. 
        \end{proof}
        
        \item \label{s_ge1} $S_{t_K} \succeq 1_S$ if there is $t_C$ with $ t_K = h_{KC} \star t_C$ and $w_\mathcal{T}(e) \succeq 1_S$ for all $e \in \mathbb{E}$,
        \begin{proof}
            By the definition of weighted type graph, for all $e \in \mathbb{E}$, we have $w(e) \neq 0_S$.  
            By assumption, we have $w_\mathcal{T}(e) \succeq 1_S$ for all $e \in \mathbb{E}$. Thus, we have 
            \begin{flalign}
                1_S \preceq w(e) \neq 0_S \label{eq_we_neq_0s_geq1_0}
            \end{flalign}
            By Proposition~\ref{prop_endrullis_2d7}~\eqref{eq:prop_neg0_ge1_mul_ge1}, we have
            \begin{flalign}
                1_S \preceq w_\mathcal{T}(t_C - h_{KC}) \label{eq_we_neq_0s_geq1}
            \end{flalign}

            Therefore, $S_{t_K} \overset{\operatorname{def}}{=}   
            \underset{\substack{t_C:C \rightarrow T \\
            t_K = h_{KC} \star t_C }}{\bigoplus} 
            w_\mathcal{T}(t_C - h_{KC}) \succeq 1_S$ if there is be a morphism such that $t_K = h_{KC} \star t_C$ by Definition~\ref{def:real_strongly_monotonic_semiring}~\eqref{ax:s1}.
        \end{proof}
        
        % \item \label{claim:le} $Y \succeq X \implies  Y \odot S_{t_K} \succeq X \odot S_{t_K}$
        % \\ by Axiom \eqref{ax:s3}. \todo{to delete: inutile}
         
        \item \label{claim:st} if there is $t_C$ with $t_K = h_{KC} \star t_C$ then
        $$ \mu(Y) >_{\mathbb{R}^+} \mu(X) + \delta  \implies \mu(Y \odot S_{t_K}) >_{\mathbb{R}^+} \mu(X \odot S_{t_K})$$
                % $$\left (\exists \alpha \geq \delta.~\mu(Y) > \mu(X) + \alpha \right ) \implies (\mu(Y \odot S_{t_K}) > \mu(X \odot S_{t_K}))$$
        \begin{proof}
           Suppose that there is $t_C$ with $t_K = h_{KC} \star t_C$. We have $S_{t_K} \neq 0_S$ by Claim~\ref{s_nz}, and we conclude by Definition~\ref{def:real_strongly_monotonic_semiring}~\eqref{ax:s4''}.
        \end{proof}
    
        \item \label{claim:sh_{DT}elta} 
        if there is $t_C$ with $t_K = h_{KC} \star t_C$, and  $w_\mathcal{T}(e) \succeq 1_S$ for all $e \in \mathbb{E}$ then
        % $$\left (\exists \alpha \geq \delta.~\mu(Y) > \mu(X) +  \alpha \right ) \implies (\exists \beta \geq \delta. \mu(Y \odot S_{t_K}) > \mu(X \odot S_{t_K})  + \beta)$$
        $$\mu(Y) >_{\mathbb{R}^+} \mu(X) +  \delta \implies \mu(Y \odot S_{t_K}) >_{\mathbb{R}^+} \mu(X \odot S_{t_K})  + \delta $$
        \begin{proof}
            Suppose that there is $t_C$ with $t_K = h_{KC} \star t_C$. We have $1_S \preceq S_{t_K} \neq 0_S$ by Claim~\ref{s_nz} and Claim~\ref{s_ge1}, and we conclude by Definition~\ref{def:real_strongly_monotonic_semiring}~\eqref{ax:s4'}. 
        \end{proof}

        \item \label{claim:0} 
        if there is no $t_C$ with $t_K = h_{KC} \star t_C$ then  $S_{t_K} = 0_S$, thus
        $$Y \odot S_{t_K} = 0_S = X \odot S_{t_K} $$
    
        \item \label{claim:exist_st} 
        If there is a context closure $t_L$ for $\rho$ and $T$ in $\mathfrak{F}$ , then, let $t_K = l \star t_L$, we have
        $$ \mu(Y) >_{\mathbb{R}^+} \mu(X) + \delta \implies \mu(Y \odot S_{t_K}) >_{\mathbb{R}^+} \mu(X \odot S_{t_K})$$
        \begin{proof}
            
       By Definition~\ref{def:context_closure} of context closure, there is $t_G : G \rightarrow T$ such that 
        \begin{flalign*}
             t_L = h_{LG} \star t_G \tag{1} \label{eq_tl_hlg_tg}
        \end{flalign*}
      i.e. we have the following commutative diagram:
     
    \begin{center}
        \begin{tikzpicture}[node distance=11mm]
          \node (I) at (0,0) {$K$};
          \node (L) at (-2,0) {$L$};
          \node (R) at (2,0) {$R$};
          \node (G) at (-2,-1) {$G$};
          \node (C) at (0,-1) {$C$};
          \node (H) at (2,-1) {$H$};
          \node (T) at (-3.5,-0.5) {$T$};
          \draw [->] (L) to  node [label, above] {$t_L$}  (T);
          \draw [->] (G) to  node [label, below] {$t_G$} (T);
          \draw [->] (I) to node [label, above] {$l$} (L);
          \draw [->] (I) to node [label,above] {$r$} (R);
        %   \draw [->] (L) to node [label, right] {$m$} (G);
        \draw [->] (L) to node [label, right] {} (G);
          \draw [->] (I) to (C);
          \draw [->] (R) to (H);
          \draw [->] (C) to (G);
          \draw [->] (C) to (H);
        \end{tikzpicture}
      \end{center}
    
        Let $t_C \overset{\operatorname{def}}{=} h_{CG} \star t_G$ and $t_K \overset{\operatorname{def}}{=} l \star t_L$. We have\\
        \begin{flalign*}
              t_K  &=  l \star t_L &\text{by definition of $t_K$}
            \\ &=   l \star (h_{LG}  \star t_G) & \text{by Equation~\eqref{eq_tl_hlg_tg}\text{ above}}
            \\ &= (l \star h_{LG}) \star t_G &\text{by associativity }
            \\ &= (h_{KC} \star h_{CG}) \star t_G & \text{by commutativity of $\square KLGC$}
            \\ &= h_{KC} \star (h_{CG}  \star t_G) & \text{by associativity}
            \\ & = h_{KC} \star t_C &\text{by definition of $t_C$}
        \end{flalign*}
        and the claim follows from Claim~\ref{claim:st}, since $t_C$ is a morphism such that $t_K = h_{KC} \star t_C$.
    \end{proof}

        \item \label{claim:exist_sh_{DT}elta} 
        If there is a context closure $t_L$ for $\rho$ and $T$ in $\mathfrak{F}$, and $w_\mathcal{T}(e) \succeq 1_S$ for all $e \in \mathbb{E}$ then, let $t_K = l \star t_L$, we have 
            % $$\left (\exists \alpha \geq \delta. Y \succ X +  \alpha \right ) \implies (\exists \beta \geq \delta. Y \odot S_{t_K} \succ X \odot S_{t_K}  + \beta)$$
        $$\mu(Y) >_{\mathbb{R}^+} \mu(X) + \delta \implies \mu(Y \odot  S_{t_K}) \succ_{\mathbb{R}^+} \mu(X \odot  S_{t_K})  + \delta$$ 
        \begin{proof}
            The proof is analogous to the proof of Claim~\ref{claim:exist_st} but with Claim~\ref{claim:sh_{DT}elta} instead of Claim~\ref{claim:st} in the end.
        \end{proof} 
    \end{enumerate}
    
    \noindent For every \( t_K: K \rightarrow T \), let
    \begin{flalign*}
        \Lambda_{t_K} &\overset{\operatorname{def}}{=}  w_\mathcal{T}(\{l \star - = t_K\})
        \\
        \Omega_{t_K} &\overset{\operatorname{def}}{=}  w_\mathcal{T}(\{r \star - = t_K\})
    \end{flalign*}
  By Lemma~\ref{nwf:lem_4d13}, we have 
        \begin{flalign*} 
            w_\mathcal{T}(G) &=
                \underset{\substack{t_K: K \rightarrow T}}{\bigoplus}     \ \
            (S_{t_K} \odot \Lambda_{t_K}),
              \\
            w_\mathcal{T}(H) &\preceq
            \underset{\substack{t_K: K \rightarrow T}}{\bigoplus}     \ \
                (S_{t_K} \odot \Omega_{t_K}).
        \end{flalign*}

    \noindent We complete the proof by case analysis:
    \begin{enumerate}
        \item  If $\rho$ is weakly decreasing, then, by Definition~\ref{nwf:def:decreasing_rule} of weakly decreasing rule, we have $\Lambda_{t_K} \succeq \Omega_{t_K}$
        for every $t_K: K \rightarrow T$. 
        By Definition~\ref{def:real_strongly_monotonic_semiring}~\eqref{ax:s3}, for every  $ t_K : K \rightarrow T$, we have 
                \begin{flalign*} 
                    S_{t_K} \odot \Lambda_{t_K} \succeq S_{t_K} \odot \Omega_{t_K}. \tag{NE} \label{steps:weightC:ge} 
                \end{flalign*}
        Thus, we have $w_\mathcal{T}(G) \succeq w_\mathcal{T}(H)$, from Definition~\ref{def:real_strongly_monotonic_semiring}~\eqref{ax:s1}.
        % \item
        %     If $\rho$ is $\delta$-uniformly decreasing for some $\delta \in \mathbb{R}_{>0}$, then, by Definition~\ref{def:decreasing_rule} of $\delta$-uniformly decreasing rule,
        %     for all $t_K : K \to T$, we have 
        %                     \begin{itemize}                                
        %                         \item $\exists \alpha \geq \delta.~\mu(\Lambda_{t_K}) > \mu(\Omega_{t_K}) + \alpha$, or
        %                         \item $\{l \star - = t_K\} = \emptyset = \{r \star - = t_K\}$
        %                     \end{itemize}
        %     From \ref{claim:st} and \ref{claim:0}, for every \( t_K: K \rightarrow T \), we have
        %     \begin{enumerate}[label=(\roman*)]
        %         \item $S_{t_K} \odot \Lambda_{t_K} = 0_S =  S_{t_K} \odot \Omega_{t_K}$, or        
        %         \item  \label{it:strict} $ \mu(\Lambda_{t_K} \odot S_{t_K}) >  \mu(S_{t_K} \odot \Omega_{t_K})$
        %     \end{enumerate}
        %     To establish $ \mu(w_\mathcal{T}(G)) > \mu(w_\mathcal{T}(H))$, using \eqref{ax:s2}, 
        %     it suffices to show that we have case~\ref{it:strict} for some $t_K : K \to T$.
        %     This follows from \ref{claim:exist_st} since we have a context closure for $\rho$ and $\mathcal{T}$, by assumption.
        \item  
            Suppose that $\rho$ is $\delta$-uniformly decreasing for $\delta \in \mathbb{R}_{>0}$, and $w_\mathcal{T}(e) \succeq 1_S$ for all $e \in \mathbb{E}$. By Definition~\ref{nwf:def:decreasing_rule} of $\delta$-uniformly decreasing rule,
            for all $t_K : K \to T$, we have  
                            \begin{itemize}                                
                                \item $\mu(\Lambda_{t_K}) >_{\mathbb{R}^+} \mu(\Omega_{t_K}) + \delta$,
                                % $\exists \alpha \geq \delta  .~\mu(\Lambda_{t_K}) > \mu(\Omega_{t_K}) + \alpha$, 
                                 or
                                \item $\{l \star - = t_K\} = \emptyset = \{r \star - = t_K\}$
                            \end{itemize}
            From Claim~\ref{claim:sh_{DT}elta} and Claim~\ref{claim:0}, for every \( t_K: K \rightarrow T \), we obtain
            \begin{enumerate}[label=(\roman*)]
                \item $S_{t_K} \odot \Lambda_{t_K} = 0_S =  S_{t_K} \odot \Omega_{t_K}$, or
                \item  \label{it:strich_{DT}elta}  $\mu(\Lambda_{t_K} \odot S_{t_K}) >_{\mathbb{R}^+} \mu(\Omega_{t_K} \odot S_{t_K}) + \delta$
                % $\exists \beta \geq \delta  .\  \mu(\Lambda_{t_K} \odot S_{t_K}) > \mu(\Omega_{t_K} \odot S_{t_K}) + \beta$
            \end{enumerate}
            To establish $ \mu(w_\mathcal{T}(G)) >_{\mathbb{R}^+} \mu(w_\mathcal{T}(H)) + \delta$, using Definition~\ref{def:real_strongly_monotonic_semiring}~\eqref{ax:s2'}, 
            it suffices to show that we have Case~\ref{it:strich_{DT}elta} for some $t_K : K \to T$.
            This follows from Claim~\ref{claim:exist_sh_{DT}elta} since we have a context closure for $\rho$ and $\mathcal{T}$ by assumption.
            % \item
            % If $\rho$ is $\delta$-closure decreasing, 
            % then it is also weakly decreasing and we obtain \eqref{steps:weightC:ge} for every $t_K : K \to T$.
            % Since the semiring is strictly ordered museurable, it suffices to show that there exists some $t_K : K \to T$ such that
            % \begin{align}
            %    \mu( S_{t_K} \odot \Lambda_{t_K})  \succ \mu(S_{t_K} \odot \Omega_{t_K})
            %   \tag{$\star$} \label{steps:weightC:gt}
            % \end{align}
            % in order to conclude $\mu(w_\mathcal{T}(G)) > \mu(w_\mathcal{T}(H))$ by Equation~\eqref{ax:s5}.
            % By Definition~\ref{def:decreasing_rule}, there is a context closure $t_L$ for $\rho$ and $T$, and 
            % % $\exists \delta' \geq \delta  .\  \Lambda_{t_K} > \Omega_{t_K} + \delta'$
            % $\mu(\Lambda_{t_K}) > \mu(\Omega_{t_K}) + \delta$
            % for $t_K = l \star t_L$. Thus, we obtain  Equation~\eqref{steps:weightC:gt} by~\ref{claim:exist_st}.
        \item
            If $\rho$ is $\delta$-closure decreasing, and $w_\mathcal{T}(e) \succeq 1_S$ for all $e \in \mathbb{E}$ then it is also weakly decreasing and we obtain Inequality~\eqref{steps:weightC:ge} for every $t_K : K \to T$.
            Since the semiring is a strictly monotonic museurable semiring,  by \autoref{def:real_strongly_monotonic_semiring}~\eqref{ax:s2'} and Definition~\ref{def:real_strongly_monotonic_semiring}~\eqref{ax:s5'}, it suffices to show that there is some $t_K : K \to T$ such that 
            \begin{align}
                % \exists \alpha \geq \delta.~
                \mu(S_{t_K} \odot \Lambda_{t_K}) >_{\mathbb{R}^+} \mu(S_{t_K} \odot \Omega_{t_K}) + 
                \delta
                % \alpha
              \tag{$\star\star$}\label{steps:weightC:gh_{DT}elta}
            \end{align}
            in order to conclude $ \mu(w_\mathcal{T}(G)) >_{\mathbb{R}^+} \mu(w_\mathcal{T}(H)) + \delta$.
            There is a context closure $t_L$ for $\rho$ and $T$, and
            $\mu(\Lambda_{t_K}) >_{\mathbb{R}^+} \mu(\Omega_{t_K}) + \delta$
            % $\exists \delta' \geq \delta  .\  \Lambda_{t_K} > \Omega_{t_K} + \delta'$
            for $t_K = l \star t_L$. Thus, we obtain Inequality~\eqref{steps:weightC:gh_{DT}elta} by Claim~\ref{claim:exist_sh_{DT}elta}.
    \end{enumerate}
\end{proof}  




\subsection{Termination Criterion}
\label{nwf:sec:proving_termination}
Finally, we formulate a termination criterion leveraging type graphs weighted over non-well-founded semirings. 
\begin{theorem}[Termination of DPO rewriting system] 
    \label{nwf:thm:termination_grs}
    Let $\mathcal{A}$ and $\mathcal{B}$ be sets of DPO rewriting rules, $\mathcal{T} = (T,\mathbb{E}, (S, \oplus, \odot, 0_S, 1_S, \prec, \mu), w)$ a finitary weighted type graph and $\mathfrak{F}$ a DPO rewriting framework such that

     \begin{enumerate}[label=\roman*)]
        \item\label{thm1:hyp3} $w(e) \succeq 1_S$ for all $e \in \mathbb{E}$, and
        % \item\label{thm1:hyp4} $\{s \in S\mid 1_S \leq s \neq 0_S\} \subseteq \mathbb{R}_{>0}$ 
        % \item\label{thm1:hyp4} for all $x \in S$, if $ 1_S \preceq x \neq 0_S$ then $\mu(x) \geq \mu(1_S)$ and $\mu(x) \in \mathbb{R}$,
        \item\label{thm1:hyp4} for all $x \in S$, if $ 1_S \preceq x \neq 0_S$ then $\mu(x) \geq \mu(1_S)$ and $\mu(x) \in \mathbb{R}$, and
        \item for every rule $\rho \in (\mathcal{A }\cup \mathcal{B })$ and every double pushout diagram  
        $\Delta \in \mathfrak{F}(\rho)$ 
        \begin{itemize}
            \item \(\operatorname{left}(\Delta)\) is weighable with \(\mathcal{T}\),
            \item \(\operatorname{right}(\Delta)\) is bounded above by \(\mathcal{T}\). 
        \end{itemize}
    \end{enumerate}       

    \noindent If the following conditions hold:
    \begin{enumerate}
        \item there exists $\delta >_{\mathbb{R}^+} 0$ such that either every $\rho \in \mathcal{A}$ is $\delta$-uniformly decreasing, or every $\rho \in \mathcal{A}$ is $\delta$-closure decreasing, and
        \item every rule $\rho \in \mathcal{B}$ is weakly decreasing,
    \end{enumerate}
    then $\Rightarrow_{\mathcal{A},\mathfrak{F}}$ is \textbf{terminating} relative to $\Rightarrow_{\mathcal{B},\mathfrak{F}}$.
\end{theorem} 
% \begin{proof}
%     See the \hyperref[proof_termination_grs]{Appendix}.
% \end{proof}

\begin{proof} 
    \label{nwf:proof_termination_grs}
    From the definition of weighted type graph, we have 
    $$\text{for all}~e\in\mathbb{E}, w(e) \neq 0_S$$ 
    From Asumption~\ref{thm1:hyp3}, we have 
    $$\text{for all}~e\in\mathbb{E},1_S \preceq w(e)$$
    Therefore, we have 
    \begin{flalign}
        \text{for all}~e\in\mathbb{E},1_S \preceq w(e)\neq 0_S \label{thm_eq_we_neq0_geq1}
    \end{flalign} 
    Let $G$ be a graph admitting a match of a DPO rewriting rule. We have 
    \begin{flalign*}
        w_\mathcal{T}(G) &\overset{\operatorname{def}}{=} 
            \underset{h \in \operatorname{Hom}(G,T)}{\bigoplus}  w_\mathcal{T}(h) \\
        & \overset{\operatorname{def}}{=} 
        \underset{h \in \operatorname{Hom}(G,T)}{\bigoplus} 
            \left ( \underset{e \in \mathbb{E}}{\bigodot} 
            \left(  
                \underset{\alpha \in \{- \star h = e\}}{\bigodot}w(e) 
            \right)
            \right )\\
    \end{flalign*} 
    By Equation~\eqref{thm_eq_we_neq0_geq1},~\autoref{prop_endrullis_2d7},~\autoref{def:bigodot} and $1_S \neq 0_S$, for every $h \in \operatorname{Hom}(G,T)$, we have
    \begin{flalign}
        1_S \preceq 
        \underset{e \in \mathbb{E}}{\bigodot} 
                \left(  
                    \underset{\alpha \in \{- \star h = e\}}{\bigodot}w(e) 
                \right) 
        \neq 0_S
    \end{flalign}
    Since $G$ be a graph admitting a match of a DPO rewriting rule, by~\autoref{prop_endrullis_2d7} and~\autoref{def:real_strongly_monotonic_semiring}\eqref{ax:s0}, we have $$1_S \preceq w_\mathcal{T}(G) \neq 0_S$$
    % By \autoref{prop_endrullis_2d7} and \eqref{ax:s0}, we have 
    % $$\forall G\in\mathcal{C}_0, (\exists H\in \mathcal{C}_0. G \Rightarrow_\mathcal{R} H) \rightarrow (1_S \preceq w_\mathcal{T}(G) \neq 0_S)$$
    By Assumption~\ref{thm1:hyp4}, we have 
    % $$\forall G\in\mathcal{C}_0, (\exists H\in \mathcal{C}_0. G \Rightarrow_\mathcal{R} H) \rightarrow (\mu(w_\mathcal{T}(G)) \geq \mu(1_S))$$
      $$\mu(w_\mathcal{T}(G)) \geq_{\mathbb{R}^+} \mu(1_S)$$
 
    % Let $G \in \mathcal{C}_0$ be an object. By Assumption \ref{thm1:hyp4}, we have $\mu(w_\mathcal{T}(G)) \in \mathbb{R}$ and $\mu(w_\mathcal{T}(G)) \geq \mu(1_S)$.

    By~\autoref{nwf:lem:decreasing_step}, every rewriting step with rules in $\mathcal{A}$ strictly decreases the weight by at least $\delta$ and no rewriting step with rules in $\mathcal{B}$ increases the weight.
    Consequently, there is no infinite rewriting sequence with an infinite rewriting steps with rules in $\mathcal{A}$ from $G$.
\end{proof}
\begin{example}
    \label{nwf:example:termination}
    Termination of the DPO rule in \autoref{ex:grsaa} can be established using \autoref{nwf:thm:termination_grs} together with the weighted type graph in \autoref{nwf:example:weighted_type_graph} over the real arithmetic semiring $\mathfrak{N}' = (\mathbb{R}^+,+,*,0_\mathbb{R},1_\mathbb{R},<,\operatorname{id}_{\mathbb{R}^+})$. It is $\delta$-closure decreasing with $\delta = 0.5$ by \autoref{nwf:example:decreasing_rule}.
\end{example}

 

\subsection{Implementation}
\label{nwf:sec:type_graph:implementation}
% We have implemented both our approach and the one proposed by Endrullis and Overbeek~\cite{endrullis2024generalized_arxiv_v2} for edge-labeled DPO rewriting systems into a unified tool. These two approaches can be launched in parallel and collaborate to establish termination of a DPO rewriting system.
Chapter~\ref{chap:lyonparallel} provides a description of the implementation of our approach and the one proposed by Endrullis and Overbeek~\cite{endrullis2024generalized_arxiv_v2} in our tool, called LyonParallel.

\paragraph{Searching Strategy.}
Consider a DPO rewriting system $\mathcal{R}$. Given a fixed semiring $S$, a processor is spawned to iteratively search for a suitable weighted type graph over $S$.
The search starts with a weighted type graph containing 1 node, incrementally increasing the node count by 1 until a maximum of 4 nodes is reached.
For a fixed number of nodes,
 if $S$ is a semiring over the natural numbers, the maximum edge weight is initialized to 1 and incremented by 1 (up to a limit of 3) if no suitable weighted type graph is found;
if $S$ is a semiring over the real numbers, weights are constrained to be positive real numbers, and no upper bound is imposed.
Processors targeting different semirings can run in parallel to analyze the same DPO rewriting system.
When a processor discovers a weighted type graph that witnesses relative termination of a subset of rules $\mathcal{A}$ with respect to another subset of rules $\mathcal{B}$ such that $\mathcal{R} \mathop{=} \mathcal{A} \mathop{\cup} \mathcal{B}$, it broadcasts this result to all processors and waits for them to terminate. If $\mathcal{B} \mathop{=} \emptyset$, the system's termination is proven. Otherwise, we initiate a new search for the rules in $\mathcal{B}$ unless a timeout is reached.
 
For a fixed number of nodes $k$, we adopt the approach proposed in \cite{bruggink2015proving, bruggink2014termination,zantema2014termination} to reduce the search space. Specifically, we construct a weighted type graph with
% \( \mathcal{T}=(T, \mathbb{E}, S, w) \) where: (1) the type graph $T$ has 
$k$ nodes and no parallel edges of the same label.
%  and (2) the set $\mathbb{E}$ of morphism-rulers consists of identity morphisms from subgraphs of the type graph $T$ (each is a labeled edge with its incident nodes) to the type graph $T$.
The search proceeds as follows: 
\begin{enumerate}
    \item decide if $s \overset{l}{\to} t$ exists for every pair of nodes \( s, t\) and label \( l\);
    \item assign a weight to every existing edge;
    \item verify the existence of a constant $\delta >0$ and a partition of the rule set $\mathcal{R}$ into a non-empty subset $\mathcal{A}$ and a subset $\mathcal{B}$
 such that:
    \begin{itemize}
        \item either all rules in $\mathcal{A}$ are $\delta$-uniformly decreasing or all rules in $\mathcal{A}$ are $\delta$-closure decreasing, and
        \item all rules in $\mathcal{B}$ are weakly decreasing.
    \end{itemize}
\end{enumerate}

%  Check if the weighted type graph satisfy requirements
% (1) construct a graph with \( k \) nodes and no parallel edges of the same label, which has a directed edge from node $s$ to node $t$ labeled by $l$ for each ordered pair $(s,t)$ of nodes and label $l$ from the finite set of edge labels $\Sigma$ of the rewriting system; (2) decide whether each edge exists in the weighted type graph; (3) assign a weight (a natural number or a real number according to $S$) to every existing edge; (4) 

This procedure amounts to checking the satisfiability of an existential Presburger arithmetic, Peano arithmetic, linear real arithmetic, or non-linear real arithmetic formula depending on the semiring $S$ considered.

\paragraph{Z3 Modeling.}
The type graph $T$ is modeled in Z3 by defining 
a boolean variable $x_{u,v,l} \mathop{\in} \mathbb{B}$ for every directed labeled edge $u\overset{l}{\to} v$ in the type graph, where $u,v\in\{1,...,k\}$ are nodes and $l \mathop{\in} \Sigma$ is an edge label. 
The variable $x_{u,v,l}$ has the value \textit{true} if the directed edge $u\overset{l}{\to} v$ exists in the result type graph, \textit{false} otherwise. 

The weight function $w$ and the set $\mathbb{E}$ of morphism-rulers are modeled by defining 
a variable $y_{u,v,l}$ of Z3's \textit{Real} sort or \textit{Int} sort (the solver's default theory for real numbers and integers), depending on the semiring $\mathcal{S}$ considered,
 for every pair $u,v\in\{1,...,k\}$ of nodes and edge label $l \mathop{\in} \Sigma$. 
The variable $y_{u,v,l}$ represents the weight of the directed labeled edge $u\overset{l}{\to} v$ in the resulting weighted type graph, but only if edge $u\overset{l}{\to} v$ exists (i.e. $x_{u,v,l}$ has the value \textit{true}). 
They are constrained to be non-negative.

A variable $\delta \mathop{\in} \mathbb{R}_{>0}$ is defined to ensure that there is
a partition of the rule set $\mathcal{R}$ into a non-empty subset $\mathcal{A}$ and a subset $\mathcal{B}$ such that all rules in $\mathcal{A}$ are either $\delta$-uniformly decreasing or $\delta$-closure decreasing, and all rules in $\mathcal{B}$ are weakly decreasing.

For convenience, the following auxiliary variables are also defined:
\begin{itemize}
    \item a boolean variable $v_h \mathop{\in} \mathbb{B}$ for every morphism $h$ from $L$ or $R$ to $T$;
    \item a real-valued variable $v_{h'} \mathop{\in} \mathbb{R}_{\geq 0}$ for every morphism $h$ from $L$ or $R$ to $T$;
    \item a boolean variable $v_c \mathop{\in} \mathbb{B}$ for every morphism $c$ from $L$ to $T$.
\end{itemize}  
The value of $v_h$ has the value \textit{true} if the morphism $h$ exists in the result type graph (i.e. all edges in its image exist), and \textit{false} otherwise.
The variable $v_{h'}$ holds the weight of the morphism $h$, provided the morphism $h$ exists (i.e., $v_h$ has the value \textit{true}).
The variable $v_c$ has the value \textit{true} if the morphism $c$ exists in the result type graph and can serve as a context closure for the rule, and \textit{false} otherwise.

If there is an assignment of values to 
    \begin{itemize}
        \item $(x_{u,v,l})$, for all $u,v \mathop{\in} \{1,...,k\}$ and for all $l \mathop{\in} \Sigma$, and
        \item $(y_{u,v,l})$, for all $u,v \mathop{\in} \{1,...,k\}$ and for all $l \mathop{\in} \Sigma$, and
        \item $\delta$,
    \end{itemize} 
    such that all conditions of Theorem~\ref{nwf:thm:termination_grs} are satisfied, then a suitable weighted type graph that witnesses termination of the DPO rewriting system can be constructed, because
    \begin{itemize}
        \item the values of \( (x_{u,v,l})_{u,v \mathop{\in} \{1,...,k\}, l \mathop{\in} \Sigma} \) define the type graph $T$, and
        \item the values of \( (y_{u,v,l})_{u,v \mathop{\in} \{1,...,k\}, l \mathop{\in} \Sigma} \) define the weight function $w$ and the set $\mathbb{E}$ of morphism-rulers.
    \end{itemize}  

\subsection{Empirical Results}
\label{nwf:sec:type_graph:result}
We evaluated our OCaml implementation using 13 examples from~\cite{endrullis2024generalized_arxiv_v2,plump1995ontermination,plump2018modular,bruggink2015proving,bruggink2014termination}, excluding~\cite[Example 6.4]{endrullis2024generalized_arxiv_v2} which is a DPO rewriting system in an unsupported category, and~\cite[Example 6]{bruggink2014termination} because it duplicates~\cite[Example 4]{bruggink2015proving}. Our approach with semirings over the real numbers proves termination for 8 out of 13 systems tested. Table~\ref{tabular:benchmarks} summarizes runtime performance 
across benchmark examples under different semiring configurations~\footnote{Experiments were conducted on a laptop equipped with an i5-1038NG7 CPU, which features 4 cores, a base clock speed of 2 GHz, a boost speed of 3.8 GHz, and 16 GB RAM.}.


% try_type_graph 1 200 
\begin{table}[htb]   
    \renewcommand{\arraystretch}{1.2}
    \centering
\begin{tabular}{|c|c|c|c|c|c|c |}
    \hline
     &\;\;A\;\;&\;\;a\;\;&\;\;T\;\;&\;\;t\;\;&\;\; N\;\;&\;\;n\;\; \\
    \hline
   ~\cite[Example 6.2]{endrullis2024generalized_arxiv_v2} & & & & & 2.68 &1.15   \\
   ~\cite[Example 6.3]{endrullis2024generalized_arxiv_v2} & & & & & 2.74 &1.16   \\
    %~\cite[Example 6.4]{endrullis2024generalized_arxiv_v2}&  && & &  &   &  &  & our only Graph\\
   ~\cite[Example D.3]{endrullis2024generalized_arxiv_v2} &2.25 
    % (2.30+2.188+2.2637+2.2928+ 2.196) 
    & 1.18
    % (1.17+1.2549+1.165+ 1.1576\mathop{+}1.162)
    & & & 2.24& 1.18    \\
   ~\cite[Example 3.8]{plump1995ontermination}
 & 2.95& 1.90 & 2.94 &1.87  & 3.49  &1.87   \\
    %~\cite[Example 3]{plump2018modular}
    % & 3.45& 3.44 & 3.55 &2.38  & 3.96  &17.60 &   4.24 AT &  3.45 at& 4.07 t  \\
   ~\cite[Example 4]{plump2018modular} &4.26& 3.19&  4.24 & 3.13 &
    5.82
    % (5.77+5.83+5.775+5.84+5.818+5.89)
    & timeout  \\

   ~\cite[Example 5]{plump2018modular} & 5.54
    &5.55
    % (5.53 +5.660\mathop{+}5.5213+5.5369+5.484689\mathop{+}5.6404\mathop{+}5.4877\mathop{+}5.5026)
    & 5.53& 5.50& 9.11&  5.62  \\
   ~\cite[Example 6]{plump2018modular} &  &  & & &  &   \\
    % (* en haut j'ai enverse tropical et arctic *)
    %~\cite[Example 6]{plump2018modular} &  &6.93 & 6.38 & 6.30 & 7.87&  7.22  &  11.71 at & 7.71 AT& 13.60 ATNat \\
   ~\cite[Example 4]{bruggink2015proving} &
    2.44
    % (2.4415+ 2.4181+2.4466+2.4366+2.4624+2.4479)
    & 
    2.46
    % (2.4405+2.5180+2.4137+2.4539+2.4480+2.4750)
    & 
    2.47
    % (2.4660+2.4457+2.4559+2.4526+2.4588+2.4847+2.5731)/7
    &
    2.54
    % (2.4380+2.4557+2.5590+2.5631+2.7245+2.5041+2.5182)/7
    & 4.58 & 
    2.46
    % (2.3987+2.4521+2.4620+2.4558+2.5128+2.4507)/6
     \\
   ~\cite[Example 5]{bruggink2015proving} &  &  &&& 7.80& timeout  \\
   ~\cite[Example 6]{bruggink2015proving} &  &  &&& 9.75& timeout   \\
   ~\cite[Example 1]{bruggink2014termination} & 
     2.26
    %  (2.2887\mathop{+}2.2386 +2.2719\mathop{+}2.2735 +2.3025\mathop{+}2.1925)/6
     &1.18
    %  (1.1764+1.1837+1.1709+1.2222+1.1608+1.1645)/6
     & & &2.24
    %  (2.2365+2.2388+2.2124+2.2395+2.2550+2.2513)/6
     & 1.18  
    %  (1.2498+1.1710+1.1634+1.1655+1.1687+1.1733)/6
     \\
   ~\cite[Example 4]{bruggink2014termination} &  2.25
    % (2.2545+2.2512+2.2529+2.3498+2.2137+ 2.2257\mathop{+}2.2167\mathop{+}2.2413)/8
    & 1.22 & 2.24
    % (2.2981+2.2026\mathop{+}2.2229+2.2524+2.2421+2.2415)/6
    &1.18
    % (1.1261+1.1850+1.1848+1.1706+1.2052+1.1920)/6
    &2.25
    % (2.2347+2.3106+2.2130+2.2391+2.2466+2.2493)/6 
    & 1.19 

    \\
   ~\cite[Example 5]{bruggink2014termination} &  4.23 & 3.23  & 4.25 &3.28  & 5.82 & timeout \\
    %~\cite[Example 6]{bruggink2014termination} &  2.58& 1.63 & 2.61&1.65&2.67 & 1.50 \\
    %~\cite[Routing]{bruggink2014termination} &  & & & & &    
    % \\
    %~\autoref{nonwf} with $n=3$ & - & 1.90 &-& 1.95 &- & 2.00 & - & 2.47 a & 3.56 nT  \\
    \hline
    \end{tabular}
    % \begin{table}[htbp] 
    %     \renewcommand{\arraystretch}{1.2}
    %     \centering
    % \begin{tabular}{|c|c|c|c|c|c|c|c|c|c|}
    %     \hline
    %      &\;\;A\;\;&\;\;a\;\;&\;\;T\;\;&\;\;t\;\;&\;\; N\;\;&\;\;n\;\; &\;\;ATN\;\;& \;\;atn\;\;&\;\;ATNatn\;\;\\
    %     \hline
    %    ~\cite[Example 6.2]{endrullis2024generalized_arxiv_v2} & & & & & 2.80 &1.29 &   3.30 N&  2.38 n & 2.50 n \\
    %     %~\cite[Example 6.4]{endrullis2024generalized_arxiv_v2}&  && & &  &   &  &  & our only Graph\\
    %    ~\cite[Example D.3]{endrullis2024generalized_arxiv_v2} & 2.30 & 1.63 & & & 2.98& 1.33  &    2.83 N &  2.09 a& 2.19 a  \\
    %    ~\cite[Example 3.8]{plump1995ontermination}
    %  & 3.45& 3.44 & 3.55 &2.38  & 3.96  &17.60 &   4.24 AT &  3.45 at& 4.07 t  \\
    %    ~\cite[Example 3]{plump2018modular}
    %     & 3.45& 3.44 & 3.55 &2.38  & 3.96  &17.60 &   4.24 AT &  3.45 at& 4.07 t  \\
    %    ~\cite[Example 4]{plump2018modular} &  5.17 & 5.31  & 5.31 &4.18  & 6.95 & 21.68 & 7.10 ATN & 5.55 atn &7.15 atnA \\
    %    ~\cite[Example 5]{plump2018modular} & 6.27&6.93 & 6.38 & 6.30 & 7.87&  7.22  &  11.71 AT & 7.71 at& 13.60 atnAT \\
    %    ~\cite[Example 6]{plump2018modular} &  &  & & &  &  & & & \\
    %     % (* en haut j'ai enverse tropical et arctic *)
    %     %~\cite[Example 6]{plump2018modular} &  &6.93 & 6.38 & 6.30 & 7.87&  7.22  &  11.71 at & 7.71 AT& 13.60 ATNat \\
    %    ~\cite[Example 4]{bruggink2015proving} &2.76 & 2.91& 2.75 &2.79  & 5.05 & 2.69&  4.92 AT & 3.28 at &5.86 anT\\
    %    ~\cite[Example 5]{bruggink2015proving} &  &  &&& 10.17& timeout&  10.96 ATN& timeout& 21.07 atN  \\
    %    ~\cite[Example 6]{bruggink2015proving} &  &  &&& 21.50& timeout  &  25.39 ATN & timeout &78.28 atnAN\\
    %    ~\cite[Example 1]{bruggink2014termination} &  2.42&1.64& & &3.41& 1.33  &   3.08 A& 3.32 a &3.53 n \\
    %    ~\cite[Example 4]{bruggink2014termination} &  2.58& 1.63 & 2.61&1.65&2.67 & 1.50& 2.79 N& 1.66 a &2.57 t\\
    %    ~\cite[Example 5]{bruggink2014termination} &  5.17 & 5.31  & 5.31 &4.18  & 6.95 & 21.68 & 7.10 ATN & 5.55 atn &7.15 atnA \\
    %    ~\cite[Example 6]{bruggink2014termination} &  2.58& 1.63 & 2.61&1.65&2.67 & 1.50& 2.79 N& 1.66 a &2.57 t\\
    %    ~\cite[Routing]{bruggink2014termination} &  & & & & & & 11.38 ATN & 14.52 atn & 16.44 at  
    %     \\
    %     %~\autoref{nonwf} with $n=3$ & - & 1.90 &-& 1.95 &- & 2.00 & - & 2.47 a & 3.56 nT  \\
    %     \hline
    %     \end{tabular}%
    
    \medskip
    \caption{
    Runtime performance in seconds (200s timeout).
    Columns denote semirings:
    ``A'', ``T'', ``N'' denote the arctic, tropical, and arithmetic semirings over the (extended) natural numbers;
    ``a'', ``t'', ``n'' denote the arctic, tropical, and arithmetic semirings over the (extended) real numbers.
    % ``ATN", ``atn", ``ATNatn" denote combinations of different semirings. For these combined configurations, the specific semirings involved are also indicated. 
    Empty cells indicate termination not proven.
     }
    \label{tabular:benchmarks}
\end{table}

% \begin{table}[htbp]  
%     \renewcommand{\arraystretch}{1.2}
%     \centering
% \begin{tabular}{|c|c|c|c|c|c|c |}
%     \hline
%      &\;\;A\;\;&\;\;a\;\;&\;\;T\;\;&\;\;t\;\;&\;\; N\;\;&\;\;n\;\; \\
%     \hline
%    ~\cite[Example 6.2]{endrullis2024generalized_arxiv_v2} & & & & & 2.80 &1.29   \\
%    ~\cite[Example 6.3]{endrullis2024generalized_arxiv_v2} & & & & & 2.74 &1.16   \\
%     %~\cite[Example 6.4]{endrullis2024generalized_arxiv_v2}&  && & &  &   &  &  & our only Graph\\
%    ~\cite[Example D.3]{endrullis2024generalized_arxiv_v2} & 2.30 & 1.63 & & & 2.98& 1.33    \\
%    ~\cite[Example 3.8]{plump1995ontermination}
%  & 3.45& 3.44 & 3.55 &2.38  & 3.96  &17.60   \\
%     %~\cite[Example 3]{plump2018modular}
%     % & 3.45& 3.44 & 3.55 &2.38  & 3.96  &17.60 &   4.24 AT &  3.45 at& 4.07 t  \\
%    ~\cite[Example 4]{plump2018modular} &  5.17 & 5.31  & 5.31 &4.18  & 6.95 & 21.68   \\
%    ~\cite[Example 5]{plump2018modular} & 6.27&6.93 & 6.38 & 6.30 & 7.87&  7.22    \\
%    ~\cite[Example 6]{plump2018modular} &  &  & & &  &   \\
%     % (* en haut j'ai enverse tropical et arctic *)
%     %~\cite[Example 6]{plump2018modular} &  &6.93 & 6.38 & 6.30 & 7.87&  7.22  &  11.71 at & 7.71 AT& 13.60 ATNat \\
%    ~\cite[Example 4]{bruggink2015proving} &2.76 & 2.91& 2.75 &2.79  & 5.05 & 2.69 \\
%    ~\cite[Example 5]{bruggink2015proving} &  &  &&& 7.87& timeout  \\
%    ~\cite[Example 6]{bruggink2015proving} &  &  &&& 9.77& timeout   \\
%    ~\cite[Example 1]{bruggink2014termination} &  2.42&1.64& & &3.41& 1.33   \\
%    ~\cite[Example 4]{bruggink2014termination} &  2.58& 1.63 & 2.61&1.65&2.67 & 1.50 \\
%    ~\cite[Example 5]{bruggink2014termination} &  5.17 & 5.31  & 5.31 &4.18  & 6.95 & 21.68  \\
%    ~\cite[Example 6]{bruggink2014termination} &  2.58& 1.63 & 2.61&1.65&2.67 & 1.50 \\
%     %~\cite[Routing]{bruggink2014termination} &  & & & & &    
%     % \\
%     %~\autoref{nonwf} with $n=3$ & - & 1.90 &-& 1.95 &- & 2.00 & - & 2.47 a & 3.56 nT  \\
%     \hline
%     \end{tabular}
%     % \begin{table}[htbp] 
%     %     \renewcommand{\arraystretch}{1.2}
%     %     \centering
%     % \begin{tabular}{|c|c|c|c|c|c|c|c|c|c|}
%     %     \hline
%     %      &\;\;A\;\;&\;\;a\;\;&\;\;T\;\;&\;\;t\;\;&\;\; N\;\;&\;\;n\;\; &\;\;ATN\;\;& \;\;atn\;\;&\;\;ATNatn\;\;\\
%     %     \hline
%     %    ~\cite[Example 6.2]{endrullis2024generalized_arxiv_v2} & & & & & 2.80 &1.29 &   3.30 N&  2.38 n & 2.50 n \\
%     %     %~\cite[Example 6.4]{endrullis2024generalized_arxiv_v2}&  && & &  &   &  &  & our only Graph\\
%     %    ~\cite[Example D.3]{endrullis2024generalized_arxiv_v2} & 2.30 & 1.63 & & & 2.98& 1.33  &    2.83 N &  2.09 a& 2.19 a  \\
%     %    ~\cite[Example 3.8]{plump1995ontermination}
%     %  & 3.45& 3.44 & 3.55 &2.38  & 3.96  &17.60 &   4.24 AT &  3.45 at& 4.07 t  \\
%     %    ~\cite[Example 3]{plump2018modular}
%     %     & 3.45& 3.44 & 3.55 &2.38  & 3.96  &17.60 &   4.24 AT &  3.45 at& 4.07 t  \\
%     %    ~\cite[Example 4]{plump2018modular} &  5.17 & 5.31  & 5.31 &4.18  & 6.95 & 21.68 & 7.10 ATN & 5.55 atn &7.15 atnA \\
%     %    ~\cite[Example 5]{plump2018modular} & 6.27&6.93 & 6.38 & 6.30 & 7.87&  7.22  &  11.71 AT & 7.71 at& 13.60 atnAT \\
%     %    ~\cite[Example 6]{plump2018modular} &  &  & & &  &  & & & \\
%     %     % (* en haut j'ai enverse tropical et arctic *)
%     %     %~\cite[Example 6]{plump2018modular} &  &6.93 & 6.38 & 6.30 & 7.87&  7.22  &  11.71 at & 7.71 AT& 13.60 ATNat \\
%     %    ~\cite[Example 4]{bruggink2015proving} &2.76 & 2.91& 2.75 &2.79  & 5.05 & 2.69&  4.92 AT & 3.28 at &5.86 anT\\
%     %    ~\cite[Example 5]{bruggink2015proving} &  &  &&& 10.17& timeout&  10.96 ATN& timeout& 21.07 atN  \\
%     %    ~\cite[Example 6]{bruggink2015proving} &  &  &&& 21.50& timeout  &  25.39 ATN & timeout &78.28 atnAN\\
%     %    ~\cite[Example 1]{bruggink2014termination} &  2.42&1.64& & &3.41& 1.33  &   3.08 A& 3.32 a &3.53 n \\
%     %    ~\cite[Example 4]{bruggink2014termination} &  2.58& 1.63 & 2.61&1.65&2.67 & 1.50& 2.79 N& 1.66 a &2.57 t\\
%     %    ~\cite[Example 5]{bruggink2014termination} &  5.17 & 5.31  & 5.31 &4.18  & 6.95 & 21.68 & 7.10 ATN & 5.55 atn &7.15 atnA \\
%     %    ~\cite[Example 6]{bruggink2014termination} &  2.58& 1.63 & 2.61&1.65&2.67 & 1.50& 2.79 N& 1.66 a &2.57 t\\
%     %    ~\cite[Routing]{bruggink2014termination} &  & & & & & & 11.38 ATN & 14.52 atn & 16.44 at  
%     %     \\
%     %     %~\autoref{nonwf} with $n=3$ & - & 1.90 &-& 1.95 &- & 2.00 & - & 2.47 a & 3.56 nT  \\
%     %     \hline
%     %     \end{tabular}%
    
%     \medskip
%     \caption{
%     Runtime performance in seconds using the automated strategy (200s timeout).
%     Columns denote semirings:
%     ``A'', ``T'', ``N'' denote the arctic, tropical, and arithmetic semirings over (extended) natural numbers;
%     ``a'', ``t'', ``n'' denote the arctic, tropical, and arithmetic semirings over (extended) real numbers.
%     % ``ATN", ``atn", ``ATNatn" denote combinations of different semirings. For these combined configurations, the specific semirings involved are also indicated. 
%     Empty cells indicate termination not proven.
%      }
%     \label{tabular:benchmarks}
%     \end{table}
     

 

\subsection{Discussion}
\label{nwf:sec:type_graph:related_work}
\paragraph{Acceleration with the real tropical/arctic semiring.} 
Consider the examples in \autoref{tabular:benchmarks} whose termination is provable using the tropical or arctic semirings. For all such examples except \cite[Example 5]{plump2018modular} and \cite[Example 4]{bruggink2015proving}, the real-number semiring approach achieves better runtime performance than the natural-number semiring approach. The acceleration rate varies across examples due to Z3's internal heuristics. This result is notable because termination proofs for these systems only require weighted type graphs with small edge weights (e.g., values less than 2 for most examples and at most 3 for all), which theoretically favors the implementation of the natural-number semiring approach.

\paragraph{Timeout with the real arithmetic semiring.}
The systems whose termination is not proven by our tool are \cite[Example 6]{plump2018modular}, \cite[Examples 5 and 6]{bruggink2015proving}, \cite[Example 4]{plump2018modular}, and \cite[Example 5]{bruggink2014termination}.
The first case (\cite[Example 6]{plump2018modular}) is due to a limitation of the type graph method in general, as the system has a rule whose right-hand side graph can be embedded into its left-hand side graph, which is problematic for the type graph method as pointed out in \cite[Example D.4]{endrullis2024generalized_arxiv_v2}. For the remaining cases,
 the failure stems from the double-exponential time complexity~\cite{collins1974quantifier,z3realarithmetic} in the number of variables \( (x_{u,v,l})_{u,v \in \{1,...,k\}, l \in \Sigma} \) and \( (y_{u,v,l})_{u,v \in \{1,...,k\}, l \in \Sigma} \), which is also influenced by the length of the conditions of \autoref{thm:termination_grs} encoded in Z3.
\paragraph{Acceleration with the real arithmetic semiring.} For some examples, the approach with semirings on real numbers can achieve a better runtime performance than the approach with semirings on natural numbers. These are (1) examples which need weighted type graphs with 1 or 2 nodes, and the constraints can be expressed in Z3 with a short formula, or (2) examples which can first be simplified by eliminating some rules using weighted type graphs with 1 node, and then the remaining rules can be 
easily shown to terminate
as the first case.
 \paragraph{Interest of the approach with semiring on natural numbers.}
 In the cases where the termination cannot be proven with the semiring on real numbers,
 using the semiring on natural numbers is more efficient, as it allows one to restrict the search space to an extremely small size by bounding the maximum weight of edges, which is not possible for the semiring on real numbers. However, this advantage can be achieved only if the termination can be proven with a weighted type graph with small edge weights, as is the case for all examples from the previous work in the type graph method.
\paragraph{Limitations of our experiments.} 
While our search strategy is very close to the real-world scenario that a non-expert user would use to prove termination of DPO rewriting systems, 
it makes it difficult to estimate the runtime performance gain of our approach over the one proposed by Endrullis and Overbeek~\cite{endrullis2024generalized_arxiv_v2}.  
This is because the reported runtimes include (1) translation of constraints into Z3's input format, (2) the time spent in Z3 to solve the constraint system, and (3) the parallelism overhead. 
Another problem is that the examples in our experiments are limited to those whose termination can be proven with a weighted type graph whose edge weights are small (at most 3). These examples favor the approach with semirings on natural numbers, as the search space is extremely small when the maximum edge weight is bounded by 3.
\paragraph{Comparison with \texttt{GraphTT-wtg}.}
A comparison of our implementation of the approach with semirings on real numbers with the one proposed in \texttt{GraphTT-wtg} by Endrullis and Overbeek~\cite{endrullis2024generalized_arxiv_v3} is desired, but this tool is not publicly available and we do not think their article provides enough information of the implementation to allow us to do so.

\section{Conclusion}
\label{nwf:sec:type_graph:conclusion}
The type graph method is a technique for proving termination of DPO rewriting systems. 
To apply this method, one needs to construct a weighted type graph over a semiring that witnesses termination of the rewriting system, and three concrete well-founded semirings over the natural numbers have been proposed in prior work: the tropical semiring, the arctic semiring, and the arithmetic semiring.
However, constructing suitable weighted type graphs over these concrete semirings for DPO rewriting systems on edge-labeled directed multigraphs is difficult. 
To address this challenge, we investigate weighted type graphs over non-well-founded semirings, and propose three corresponding concrete semirings over the real numbers: the tropical semiring, the arctic semiring, and the arithmetic semiring. Construction of weighted type graphs over these concrete semirings is computationally easier.
    
We implemented both our approach and the one proposed by Endrullis et al.~\cite{endrullis2024generalized_arxiv_v2} for DPO rewriting systems on edge-labeled directed multigraphs into a unified tool. Experiments show that for the tropical and arctic semirings, the weighted type graphs over the real semiring can be constructed more easily, but for the arithmetic semiring, the weighted type graphs over the real semiring are more difficult to construct than those over the natural semiring.
% , because the complexity is extremely high and the search space cannot be restricted by bounding the maximum edge weight. 

A promising direction involves generalizing the type graph method to other algebraic graph rewriting formalisms, broadening its applicability in automated verification. 
  

%%%%%%%%%%%%%%%%%%%%%%%%%%%%%%%%%%%%%%%%%%%%%%%%%%%
%%%%   Chapter : morphism counting   %%%%%%%%%%%%%%
%%%%%%%%%%%%%%%%%%%%%%%%%%%%%%%%%%%%%%%%%%%%%%%%%%%
\chapter{Termination of Injective DPO Graph Rewriting
Systems using Morphism Counting}
\label{chap:subgraph_counting}

\section{Introduction}
\label{subgraph_counting:sec:intro}

In this chap, we present an new automated method for relative termination~\cite{geser1990relative} of double-pushout (DPO) graph rewriting systems with injective rules~\cite{corradini1997algebraic,habel2001double,konig2018atutorial}. The method, based on the interpretation method framework~\cite{nipkow1998term}, works as follows. Let \( \mathbb{X} \) denote a fixed set of graphs. We assign weights (natural numbers) to monomorphisms from graphs in \( \mathbb{X} \), and define the weight of a graph $G$ as the sum of the weights of all monomorphisms from graphs in \( \mathbb{X} \) to $G$. 

Let $\mathcal{A}$ and $\mathcal{B}$ be sets of DPO rewriting rules. Let $X \in \mathbb{X}$ and $G \Rightarrow H$ be a rewriting step defined by the following double-pushouts diagram.
\begin{center}
 \resizebox{0.4\textwidth}{!}{
     \begin{tikzpicture}
            % [node distance=15mm]
            \node (I) at (0,0) {$K$};
            \node (L)  at (-2,0) {$L$};
            \node (R)  at (2,0) {$R$};
            \node (G)  at (-2,-2) {$G$};
            \node (C)  at (0,-2) {$C$};
            \node (H)  at (2,-2) {$H$};
            \draw [>->] (I) to  node [midway,below] {$l$} (L);
            \draw [>->] (I) to  node [midway,below] {$r$} (R);
            \draw [>->] (L) to node [midway,right] {$m$} (G);
            \draw [>->] (I) to  node [midway,right] 
            % {$u$}
            {} (C);
            \draw [>->] (R) to  node [midway,right] 
            {$m'$}
            (H);
            \draw [>->] (C) to node [midway,above] {$l'$} (G);
            \draw [>->] (C) to node [midway,above] 
            {$r'$} 
            (H);
            % \node [at=($(I)!.5!(G)$)] {\normalfont PO};
            % \node [at=($(I)!.5!(H)$)] {\normalfont PO};
        \end{tikzpicture}
    }
\end{center}

For a monomorphism from $X$ to $G$, its image falls into exactly one of the following three mutually exclusive cases. First, the image may be fully included in $m(L)$. Second, the image may be fully included in $l'(C)$; Third, the image may be neither fully included in $m(L)$ nor fully included in $l'(C)$. Similary, for a monomorphism from $X$ to $H$, its image falls into exactly one of the following three mutually exclusive cases. First, the image may be fully included in $m'(R)$. Second, the image may be fully included in $r'(C)$; Third, the image may be neither fully included in $m'(R)$ nor fully included in $r'(C)$. The number of monomorphisms from $X$ to $G$ whose image is fully included in $m(L)$ can be precisely calculated, as can the number of monomorphisms from $X$ to $H$ whose image is fully included in $m'(R)$. Moreover, the number of monomorphisms from $X$ to $G$ equals to the that of monomorphisms to $H$. 

To estimate the change in weight when applying a rewriting rule, we focus on 
the monomorphisms in the third case. Specifically, our technique requires that, for every graph $X$ in $\mathbb{X}$ and every rewriting step $G \Rightarrow H$ using a rule from $\mathcal{A}$ or $\mathcal{B}$, there exists an injective mapping from the set of monomorphisms $x : X \to H$ (where $\opn{Im}(x)$ is neither fully included in the right-hand side graph nor fully included in the context) to the set of monomorphisms from $y : X \to G$ (where $\opn{Im}(y)$ is neither fully included in the left-hand side graph nor fully included in the context). 

Suppose, additionally, that the following conditions hold: (i) For every rule in \( \mathcal{A} \), the left-hand side graph's weight is strictly greater than the right-hand side graph's weight; (ii) For every rule in \( \mathcal{B} \), the left-hand side graph's weight is greater than or equal to the right-hand side graph's weight. 

Under these conditions, rewriting steps using rules in \( \mathcal{A} \) strictly decrease the weights of the host graphs, while rewriting steps using rules in \( \mathcal{B} \) do not increase them.
Consequently, the rewriting rules in \( \mathcal{A} \) can only be applied finitely many times, since the weight of a finite graph is a natural number and strictly decreases with each rewriting step using a rule in \( \mathcal{A} \), the process must terminate after finitely many iterations.  
 
Our method proves termination for systems such as the one in~\autoref{subgraph_counting:ex_contrib_variant}, which prior interpretation-based approaches~\cite{zantema2014termination,bruggink2014termination,bruggink2015proving,
endrullis2024generalized_arxiv_v2,
overbeek2024termination_lmcs} cannot handle. 
We also provide an implementation of our technique.  
   
The remainder of this chapter is organized as follows.
We first define the weight of a graph in~\textsection~\ref{subgraph_counting:sec:interpretation}. 
Using this definition, we outline the general approach and identify the key challenge in~\textsection~\ref{subgraph_counting:sec:general_idea} and introduce the concept of non-increasing rules in \textsection~\ref{subgraph_counting:sec:non-increasing}. 
A solution to the challenge is then proposed in~\textsection~\ref{subgraph_counting:sec:solution_to_the_key_challenge}, culminating in our termination criterion in~\textsection~\ref{subgraph_counting:sec:termination}.
\textsection~\ref{subgraph_counting:sec:examples} illustrates the method with some examples.
\textsection~\ref{subgraph_counting:sec:related_work} compares our approach with some existing methods.
Finally,~\textsection~\ref{subgraph_counting:sec:conclusion} concludes with remarks and future research directions. The missing proofs are available in 
\iflongversion
the appendix.
\else
the preprint version~\cite{qiu2025termination}.
\fi

\section{Preliminaries} 
\label{subgraph_counting:sec:pre} 
While Definition~\ref{def:rewriting_inclusion} of injective rewriting rules and rewriting steps may appear more restrictive, they are equivalent to those introduced in \textsection~\ref{sec:dpo}. This equivalence arises because algebraic graph rewriting via double pushouts is defined up to isomorphism. Our preference for these definitions stems from their practical advantage: they provide a clearer view of how the number of monomorphisms from specific graphs changes after a rewriting step. 
 
\begin{definition}
    \label{def:rewriting_inclusion}
    An \textbf{inclusion} is a function $f$ such that $\opn{dom}(f) \subseteq \opn{codom}(f)$ and $f(x) = x$ for every $x \in \opn{dom}(f)$.
    A rewriting rule $\varphi = (L \overset{l}{\leftarrowtail} K \overset{r}{\rightarrowtail} R)$ consists of two inclusions $l$ and $r$.  
    Two rules $\varphi' = (L' \overset{l'}{\leftarrowtail} K' \overset{r'}{\rightarrowtail} R')$ and $\varphi$ are \textbf{equivalent} if there exist isomorphisms $a,b,c$ such that 
    $l' \star a = b \star l$ and $r' \star c = b \star r$, i.e. the following diagram is commutative:    
    % the commutative diagram illustrated in Figure~\ref{fig:subgraph_counting:general_rewriting_step_sdfsl_sdfjlsdkjfl} can be constructed:
            \begin{center}
                    \resizebox{0.40\textwidth}{!}{
                        \begin{tikzpicture}
                                \node (I) at (0,0) {$K'$};
                                \node (L) at (-2,0) {$L'$};
                                \node (R) at (2,0) {$R'$};
                                \node (G) at (-2,-2) {$L$};
                                \node (C) at (0,-2) {$K$};
                                \node (H) at (2,-2) {$R$};
                                \draw [>->] (I) to  node [midway,below] {$l'$} (L);
                                \draw [>->] (I) to  node [midway,below] {$r'$} (R);
                                \draw [>->] (L) to node [midway,right] {\textbf{a}} (G);
                                \draw [>->] (I) to node [midway,right] {\textbf{b}} (C);
                                \draw [>->] (R) to node [midway,left] {\textbf{c}} (H);
                                \draw [>->] (C) to node [midway,above] {$l$} (G);
                                \draw [>->] (C) to node [midway,above] {$r$} (H);
                                \node [at=($(I)!.5!(G)$)] {$\mathop{=}$};
                                \node [at=($(I)!.5!(H)$)] {$\mathop{=}$};
                            \end{tikzpicture}
                    }
                    % \caption{}
                    % \label{fig:subgraph_counting:general_rewriting_step_sdfsl_sdfjlsdkjfl}
                \end{center}
    A rewriting step $G \Rightarrow_\varphi H$ is a double-pushout:
    %  in Figure~\ref{fig:subgraph_counting:gedfadsdsfdsfds}
%   with injective match $m :L' \rightarrowtail G$:
         \begin{center}
          \resizebox{0.40\textwidth}{!}{
              \begin{tikzpicture}
                    \node (I) at (0,0) {$K'$};
                    \node (L) at (-2,0) {$L'$};
                    \node (R) at (2,0) {$R'$};
                    \node (G) at (-2,-2) {$G$};
                    \node (C) at (0,-2) {$C$};
                    \node (H) at (2,-2) {$H$};
                    \draw [>->] (I) to  node [midway,below] {$l'$} (L);
                    \draw [>->] (I) to  node [midway,below] {$r'$} (R);
                    \draw [>->] (L) to node [midway,left] {} (G);
                    \draw [>->] (I) to node [midway,right] { } (C);
                    \draw [>->] (R) to node [midway,left] { } (H);
                    \draw [>->] (C) to node [midway,above] { } (G);
                    \draw [>->] (C) to node [midway,above] { } (H);
                    \node [at=($(I)!.5!(G)$)] {\normalfont PO};
                    \node [at=($(I)!.5!(H)$)] {\normalfont PO};
                  \end{tikzpicture}          }
      \end{center}
      where all arrows are inclusions and $L' \overset{l'}{\leftarrowtail} K' \overset{r'}{\rightarrowtail} R'$ is equivalent to $\varphi$.
\end{definition}

\subsection{pregraphs}
Graphs in double-pushouts that define rewriting steps in \textbf{Graph} can be decomposed into unions of pregraphs. pregraphs model finite directed edge-labeled multigraphs that may contain dangling edges.
\begin{definition}
    \label{def:pregraph}
    A pregraph $G$ consists of
    \begin{itemize}
        \item a finite set of nodes, denoted by $V(G)$,
        \item a subset $V_{\nexists}(G) \subseteq V(G)$ of missing nodes
        \item a finite set of edges, denoted by $E(G)$,
        \item a source function, denoted by \( \textbf{src}_G \), (resp. a target function, denoted by \( \textbf{dst}_G \)) that assigns to each edge an existing or missing node,
        \item an edge-labeling function, denoted by $l_G$, that assigns labels to edges.
    \end{itemize}
\end{definition} 
The set of existing nodes, $V(G) \setminus V_{\nexists}(G)$, is denoted by $V_\exists(G)$.

 \begin{example}
    \label{def:pregraph}
     A pregraph with node-set $\{1,2,3\}$, a missing node $2$, and 3 dangling edges is illustrated below:
     \begin{center}
        \resizebox{0.5\textwidth}{!}{
        \begin{tikzpicture}
            \graphbox{\(\)}{00mm}{-20mm}{45mm}{20mm}{2mm}{-5mm}{ 
                \coordinate (o) at (-5mm,-8mm); 
                \node[draw,circle] (l1) at ($(o)+(-10mm,0mm)$) {1};
                % \node[draw,circle] (l2) at ($(l1)+(3,0)$) {};
                \node[draw,circle,dashed,red] (l3) at ($(l1)+(1,0)$) {2};
                \node[draw,circle] (l4) at ($(l1)+(2,0)$) {3};
                \draw[->] (l1) edge[bend right]  node[midway,below] {$a$} (l3);
                \draw[->] (l1) edge[bend left] node[midway,above] {$a$}  (l3);
                \draw[->] (l3) -- (l4) node[midway,above] {$b$};
                % \draw[->] (l4) -- (l2) node[midway,above] {$a$};
            }  
        \end{tikzpicture} 
    }
    \end{center}
 \end{example} 
 
 A homomorphism of pregraphs is a mapping between two pregraphs that preserves structure and respects missing and existing nodes.
 \begin{definition} 
    \label{def:pregraph:homomorphism}
    Let \( G \) and \( H \) be pregraphs. A \textbf{homomorphism of pregraphs} $h: G \to H$ consists a pair of functions $h_V: V(G) \to V(H) $ and $h_E: E(G) \to E(H)$ such that
    \begin{itemize}
        \item for every edge \( a \colon s \to t \) in \( G \), the edge \( h_E(a) \) in $H$ has source \( h_V(s) \) and target \( h_V(t) \),
        \item for all $s \in V_{\nexists}(G)$, we have $h_V(s) \in V_{\nexists}(H)$,
        \item for all $t \in V_\exists(G)$, we have $h_V(t) \in V_\exists(H)$.
    \end{itemize} 
\end{definition}

\begin{definition} 
    \label{def:pregraph:composition}
    Let $f: (G,\lambda) \to (G',\lambda')$ and $g: (G',\lambda') \to (G'',\lambda'')$ be homomorphisms of pregraphs. Their composition, denoted by $g \circ f$, is the homomorphism defined by the component-wise composition of functions $g_V \circ f_V$ and $g_E \circ f_E$.
\end{definition}
\begin{definition} 
    \label{def:subpregraph} Let $G$ be a pregraph. A pregraph $H$ is a \textbf{subpregraph} of $G$, written $H \rightarrowtail G$, if all of the following conditions hold:
     \begin{itemize} 
            \item $V(H) \subseteq V(G)$, \item $V_{\nexists}(H) = V(H) \cap V_{\nexists}(G)$, 
            \item $E(H) \subseteq E(G)$, 
            \item the source and target functions of $H$ are the restrictions of those of $G$ to $E(H)$,
            \item the edge-labeling function of $H$ is the restriction $l_H = {l_G}_{|_{E(H)}}$. 
        \end{itemize} 
\end{definition}
We define two operations on pregraphs that will be used in the following discussion.
The \textbf{union} $C \cup R$ of two subpregraphs $C \subseteq G$ and $R \subseteq G$ is a pregraph with 
\begin{itemize}
    \item Node set: $ V(C) \cup V(R) $,
    \item Missing nodes:
        $(
            V_{\nexists}(C) \cup V_{\nexists}(R)) \setminus 
            (   V_\exists(C) \cup 
             V_\exists(R)  
                   )$
    \item Edge set: $ E(C) \cup E(R) $, 
    \item Source function: $ \opn{src}_C \cup \opn{src}_R $,  
    \item Target function: $ \opn{dst}_C \cup \opn{dst}_R $,  
    \item Edge-labeling function: $ l_C \cup l_R $.   
\end{itemize} 
The definition is illustrated below:
% in Figure~\ref{fig:subgraph_counting:union_of_pregraphs}.
% \begin{figure}[H]
%     \centering
\begin{center}
    \resizebox{\textwidth}{!}{
        \begin{tikzpicture}  
            \graphbox{\( R \)}{91mm}{-22mm}{45mm}{25mm}{2mm}{2mm}{
                \coordinate (o) at (-5mm,-8mm); 
                \node[draw,circle] (l1) at ($(o)+(-10mm,0mm)$) {1};
                \node[draw,circle] (l2) at ($(l1)+(3,0)$) {2};
                \node[draw,circle] (l3) at ($(l1) + (1,0)$) {4};
                \node[draw,circle] (l4) at ($(l1) + (2,0)$) {5};
                \draw[->] (l1) -- (l3) node[midway,above] {$a$};
                \draw[->] (l3) -- (l4) node[midway,above] {$b$};
                \draw[->] (l4) -- (l2) node[midway,above] {$a$};
            }    
        
            \graphbox{\( C  \)}{0mm}{-22mm}{34mm}{25mm}{2mm}{-3mm}{
                \coordinate (o) at (0mm,-6mm); 
                \node[draw,circle,dashed] (l1) at ($(o)+(-10mm,0mm)$) {1};
                \node[draw,circle,dashed] (l2) at ($(l1)+(2,0)$) {2};
                \node[draw,circle] (l4) at ($(l2) + (0,-1)$) {6};
                \draw[->] (l2) -- (l4) node[midway,right] {$a$};
                \draw[->] (l2) edge[out=-135,in=-45]node[midway,below] {$a$} (l1) ;
                % \node[ draw,circle] (l6) at ($(l1) + (0,-1)$) {7};
                % \draw[<-] (l1) -- (l6) node[midway,left] {$a$};
            }    

            \graphbox{\(C \cup R \)}{40mm}{-22mm}{45mm}{25mm}{2mm}{-3mm}{
                \coordinate (o) at (-5mm,-6mm); 
                \node[draw,circle] (l1) at ($(o)+(-10mm,0mm)$) {1};
                \node[draw,circle] (l2) at ($(l1)+(3,0)$) {2};
                \node[draw,circle] (l3) at ($(l1) + (1,0)$) {4};
                \node[draw,circle] (l4) at ($(l1) + (2,0)$) {5};
                \node[ draw,circle] (l5) at ($(l2) + (0,-1)$) {6};
                % \node[ draw,circle] (l6) at ($(l1) + (0,-1)$) {7};
                % \draw[<-] (l1) -- (l6) node[midway,left] {$a$};
                \draw[->] (l1) -- (l3) node[midway,above] {$a$};
                \draw[->] (l3) -- (l4) node[midway,above] {$b$};
                \draw[->] (l4) -- (l2) node[midway,above] {$a$};
                \draw[->] (l2) -- (l5) node[midway,right] {$a$};
                \draw[->] (l2) edge[out=-135,in=-45]node[midway,below] {$a$} (l1) ;
            }    
             
        \end{tikzpicture}
    }
    % \caption{Union of pregraphs}
    % \label{fig:subgraph_counting:union_of_pregraphs}
    % \end{figure}
    \end{center}

    % \textcolor{red}{Intersection} $G \cap H$: 
    % \vspace{1mm}

    % % \begin{center}
    %     \resizebox{0.95\textwidth}{!}{
    %     \begin{tikzpicture}
    %         \graphbox{\( G \)}{0mm}{-22mm}{34mm}{26mm}{2mm}{-3mm}{
    %             \coordinate (o) at (0mm,-6mm); 
    %             \node[draw,circle] (l1) at ($(o)+(-10mm,0mm)$) {1};
    %             \node[draw,circle] (l2) at ($(l1)+(2,0)$) {2};
    %             \node[draw,circle] (l3) at ($(l1) + (1,0)$) {3};
    %             \node[draw,circle] (l4) at ($(l2) + (0,-1)$) {6};
    %             \draw[->] (l1) -- (l3) node[midway,above] {$a$};
    %             \draw[->] (l3) -- (l2) node[midway,above] {$a$};
    %             \draw[->] (l2) -- (l4) node[midway,right] {$a$};
    %             \node[draw,circle] (l6) at ($(l1) + (0,-1)$) {7};
    %             \draw[<-] (l1) -- (l6) node[midway,left] {$a$};
    %             \draw[->] (l2) edge[out=-135,in=-45]node[midway,below] {$a$} (l1) ;
    %         }    
    
    %         \graphbox{ \textcolor{red}{\(G \cap H\)}  }{40mm}{-22mm}{34mm}{25mm}{2mm}{-3mm}{
    %             \coordinate (o) at (0mm,-6mm); 
    %             \node[draw,circle] (l1) at ($(o)+(-10mm,0mm)$) {1};
    %             \node[draw,circle] (l2) at ($(l1)+(2,0)$) {2};
    %             \node[draw,circle] (l4) at ($(l2) + (0,-1)$) {6};
    %             \draw[->] (l2) -- (l4) node[midway,right] {$a$};
    %             \draw[->] (l2) edge[out=-135,in=-45]node[midway,below] {$a$} (l1) ;
    %             \node[ draw,circle] (l6) at ($(l1) + (0,-1)$) {7};
    %             \draw[<-] (l1) -- (l6) node[midway,left] {$a$};
    %         }    
    
    %         \graphbox{\( H \)}{80mm}{-22mm}{45mm}{25mm}{2mm}{-6mm}{
    %             \coordinate (o) at (-5mm,-3mm); 
    %             \node[draw,circle] (l1) at ($(o)+(-10mm,0mm)$) {1};
    %             \node[draw,circle] (l2) at ($(l1)+(3,0)$) {2};
    %             \node[draw,circle] (l3) at ($(l1) + (1,0)$) {4};
    %             \node[draw,circle] (l4) at ($(l1) + (2,0)$) {5};
    %             \node[ draw,circle] (l5) at ($(l2) + (0,-1)$) {6};
    %             \node[ draw,circle] (l6) at ($(l1) + (0,-1)$) {7};
    %             \draw[<-] (l1) -- (l6) node[midway,left] {$a$};
    %             \draw[->] (l1) -- (l3) node[midway,above] {$a$};
    %             \draw[->] (l3) -- (l4) node[midway,above] {$b$};
    %             \draw[->] (l4) -- (l2) node[midway,above] {$a$};
    %             \draw[->] (l2) -- (l5) node[midway,right] {$a$};
    %             \draw[->] (l2) edge[out=-135,in=-45]node[midway,below] {$a$} (l1) ;
    %         }    
           
    %         % \node () at (37mm,-33mm) {\( \leftarrowtail \)};
    %         % \node () at (77mm,-33mm) {\( \rightarrowtail \)}; % C -> H
    %     \end{tikzpicture}
    %     }         
    % \end{center}
The \textbf{relative complement} of $R$ in $H$ where $R \subseteq H$, denoted by $H \setminus R$, is the pregraph with
\begin{itemize}
    \item Edge set: $X \isdef E(H) \setminus E(R)$,
    \item Node set: $Y \isdef \bigcup_{e \in X} \{\opn{src}(e), \opn{dst}(e)\}$,
    \item Set of missing nodes: $
   ( V_{\nexists}(H)  \cup
        (
                  V_\exists(H) 
                  \cap 
                  V_\exists(R)
                  )
   )
        \cap  
        Y
                  $,
    \item Source function: ${\opn{src}_H}$ restricted to $X$,
    \item Target function: ${\opn{dst}_H}$ restricted to $X$,
    \item Edge-labeling function : $l_H$ restricted to $X$.
\end{itemize}
An example is shown below where the graph $H$, $R$ and the relative complement $H \setminus R$ are illustrated:
% in Figure~\ref{fig:subgraph_counting:relative_complement_of_pregraphs}.
\begin{center}  
% \begin{figure}[H]
    %     \centering
  \resizebox{\textwidth}{!}{
        \begin{tikzpicture}  
            \phantom{
                \graphbox{\( G \)}{0mm}{-22mm}{34mm}{22mm}{2mm}{-3mm}{
                }  
            }
            \graphbox{\( R \)}{80mm}{-22mm}{45mm}{25mm}{2mm}{2mm}{
                \coordinate (o) at (-5mm,-8mm); 
                \node[draw,circle] (l1) at ($(o)+(-10mm,0mm)$) {1};
                \node[draw,circle] (l2) at ($(l1)+(3,0)$) {2};
                \node[draw,circle] (l3) at ($(l1) + (1,0)$) {4};
                \node[draw,circle] (l4) at ($(l1) + (2,0)$) {5};
                \draw[->] (l1) -- (l3) node[midway,above] {$a$};
                \draw[->] (l3) -- (l4) node[midway,above] {$b$};
                \draw[->] (l4) -- (l2) node[midway,above] {$a$};
            }    
        
            \graphbox{\( H \setminus R  \)}{40mm}{-22mm}{34mm}{25mm}{2mm}{-3mm}{
                \coordinate (o) at (0mm,-6mm); 
                \node[draw,dashed,circle] (l1) at ($(o)+(-10mm,0mm)$) {1};
                \node[draw,dashed,circle] (l2) at ($(l1)+(2,0)$) {2};
                \node[draw,circle] (l4) at ($(l2) + (0,-1)$) {6};
                \draw[->] (l2) -- (l4) node[midway,right] {$a$};
                \draw[->] (l2) edge[out=-135,in=-45]node[midway,below] {$a$} (l1) ;
                % \node[ draw,circle] (l6) at ($(l1) + (0,-1)$) {7};
                % \draw[<-] (l1) -- (l6) node[midway,left] {$a$};
            }    

            \graphbox{\( H \)}{-11mm}{-22mm}{45mm}{25mm}{2mm}{-3mm}{
                \coordinate (o) at (-5mm,-6mm); 
                \node[draw,circle] (l1) at ($(o)+(-10mm,0mm)$) {1};
                \node[draw,circle] (l2) at ($(l1)+(3,0)$) {2};
                \node[draw,circle] (l3) at ($(l1) + (1,0)$) {4};
                \node[draw,circle] (l4) at ($(l1) + (2,0)$) {5};
                \node[ draw,circle] (l5) at ($(l2) + (0,-1)$) {6};
                % \node[ draw,circle] (l6) at ($(l1) + (0,-1)$) {7};
                % \draw[<-] (l1) -- (l6) node[midway,left] {$a$};
                \draw[->] (l1) -- (l3) node[midway,above] {$a$};
                \draw[->] (l3) -- (l4) node[midway,above] {$b$};
                \draw[->] (l4) -- (l2) node[midway,above] {$a$};
                \draw[->] (l2) -- (l5) node[midway,right] {$a$};
                \draw[->] (l2) edge[out=-135,in=-45]node[midway,below] {$a$} (l1) ;
            }    
            % \node () at (102mm,-18mm) {\( \downarrowtail \)};
            % \node () at (77mm,-33mm) {\( \rightarrowtail \)}; % C -> H
        \end{tikzpicture}
    }
    % \caption{Relative complement of $R$ in $H$.}
    % \label{fig:subgraph_counting:relative_complement_of_pregraphs}
    % \end{figure}
\end{center}

 
\subsection{$X$-occurrences}

    Consider the following double pushout diagram in the category \textbf{Graph}, where all arrows are inclusions:
    % illustrated in Figure~\ref{fig:subgraph_counting:pushoutsquare} can be decomposed into the diagram illustrated in Figure~\ref{fig:subgraph_counting:decomposition_of_rewriting_step}.
    % \begin{figure}[H]
    %     \centering
    \begin{center}
       \resizebox{0.49\textwidth}{!}{
        \begin{tikzpicture}
            % [node distance=15mm]
            \node (I) at (0,0) {$K$};
            \node (L)  at (-2,0) {$L$};
            \node (R)  at (2,0) {$R$};
            \node (G)  at (-2,-2) {$G$};
            \node (C)  at (0,-2) {$C$};
            \node (H)  at (2,-2) {$H$};
            \draw [>->] (I) to  node [midway,below] { } (L);
            \draw [>->] (I) to  node [midway,below] { } (R);
            \draw [>->] (L) to node [midway,right] { } (G);
            \draw [>->] (I) to  node [midway,right] 
            % {$u$}
            {} (C);
            \draw [>->] (R) to  node [midway,right] 
            {}
            (H);
            \draw [>->] (C) to node [midway,above] {} (G);
            \draw [>->] (C) to node [midway,above] 
            {}
            (H);
            % \node [at=($(I)!.5!(G)$)] {\normalfont PO};
            % \node [at=($(I)!.5!(H)$)] {\normalfont PO};
        \end{tikzpicture}
    }
    % \caption{}
    % \label{fig:subgraph_counting:pushoutsquare}
    % \end{figure}
    \end{center}
    Since all arrows are inclusions, by Definition~\ref{def:rewriting_inclusion}, graphs can be decomposed as unions of pregraphs as illustrated below, 
     where $K, L', R',C'$ are pregraphs and $L = L' \cup K$, $R = R' \cup K$, $C = C' \cup K$, $G = L' \cup K \cup C'$, $H = R' \cup K \cup C'$.
    % \begin{figure}[H]
    %     \centering
    \begin{center}
    \resizebox{0.49\textwidth}{!}{
            \begin{tikzpicture} 
                \coordinate (k) at (0, 0);
                \draw[fill=white] ($(k)+(0,0)$) rectangle ($(k)+(0.5,0.5)$);
                \node () at ($(k)+(0.25,0.25)$) {\( \mathrm{K} \)};
            
                \coordinate (c) at (0, -2.2);
                \draw[fill=blue!20]
                ($(c)+(0,-0.5)$)
                -- ($(c)+(0,0.5)$) 
                -- ($(c)+(1,0.5)$) 
                arc[start angle=0, end angle=-90, radius=1]
                -- cycle;
                \node () at ($(c)+(0.75,0.25)$) {\( \mathrm{C'} \)};
                \draw[fill=white] ($(c)+(0,0)$) rectangle ($(c)+(0.5,0.5)$);
                \node () at ($(c)+(0.25,0.25)$) {\( \mathrm{K} \)};
            
                \coordinate (l) at (-3, 0);
                \draw[fill=orange!20] ($(l)+(-0.5,0)$) rectangle ($(l)+(0.5,1)$);
                \node () at ($(l)+(-0.23,0.25)$) {\( \mathrm{L'} \)};
                \draw[fill=white] ($(l)+(0,0)$) rectangle ($(l)+(0.5,0.5)$);
                \node () at ($(l)+(0.25,0.25)$) {\( \mathrm{K} \)};
            
                \coordinate (g) at (-3, -2.2);
                \draw[fill=blue!20]
                ($(g)+(0,-0.5)$)
                -- ($(g)+(0,0.5)$)
                -- ($(g)+(1,0.5)$) 
                arc[start angle=0, end angle=-90, radius=1]
                -- cycle;
                \draw[fill=orange!20] ($(g)+(-0.5,0)$) rectangle ($(g)+(0.5,1)$);
                \node () at ($(g)+(0.75,0.25)$) {\( \mathrm{C'} \)};
                \node () at ($(g)+(-0.23,0.25)$) {\( \mathrm{L'} \)};
                \draw[fill=white] ($(g)+(0,0)$) rectangle ($(g)+(0.5,0.5)$);
                \node () at ($(g)+(0.25,0.25)$) {\( \mathrm{K} \)};
            
                \coordinate (r) at (3,0);
                \draw[fill=red!20] ($(r)+(-0.5,0)$)
                -- ($(r)+(-0.5,0.5)$)
                -- ($(r)+(0,1)$)
                --  ($(r)+(0.5,1)$)
                -- ($(r)+(0.5,0)$)
                -- cycle;
                \node () at ($(r)+(-0.23,0.25)$) {\( \mathrm{R'} \)};
                \draw[fill=white] ($(r)+(0,0)$) rectangle ($(r)+(0.5,0.5)$);
                \node () at ($(r)+(0.25,0.25)$) {\( \mathrm{K} \)};
            
                \coordinate (h) at (3, -2.2);
                \draw[fill=blue!20]
                ($(h)+(0,-0.5)$)
                -- ($(h)+(0,0.5)$)
                -- ($(h)+(1,0.5)$) 
                arc[start angle=0, end angle=-90, radius=1]
                -- cycle;
                \draw[fill=red!20] ($(h)+(-0.5,0)$)
                -- ($(h)+(-0.5,0.5)$)
                -- ($(h)+(0,1)$)
                --  ($(h)+(0.5,1)$)
                -- ($(h)+(0.5,0)$)
                -- cycle;
            \node () at ($(h)+(0.75,0.25)$) {\( \mathrm{C'} \)};
            \draw[fill=white] ($(h)+(0,0)$) rectangle ($(h)+(0.5,0.5)$);
            \node () at ($(h)+(0.25,0.25)$) {\( \mathrm{K} \)};
            \node () at ($(h)+(-0.23,0.25)$) {\( \mathrm{R'} \)};
            
                \node[ font=\huge] (kl) at ($(k)!0.5!(l)+(0.25,0.25)$)
            %    {\( \overset{l}{\leftarrowtail} \)}
                {\( \leftarrowtail \)}
                ; 
                \node[ font=\huge] (kr) at ($(k)!0.5!(r)+(0.25,0.25)$)
                {\( \rightarrowtail \)}
                ;  
                \node[ font=\huge] (cg) at ($(c)!0.5!(g)+(0.25,0.25)$) 
                {\( \leftarrowtail \)}
            ;  
                \node[ font=\huge] (ch) at ($(c)!0.5!(h)+(0.25,0.25)$)
                {\( \rightarrowtail \)}
            ; 
                \node[ font=\huge] (kc) at ($(k)!0.5!(c)+(0.2,0.4)$) {\( \downarrowtail \)}; 
            %   \node[ font=\LARGE] () at ($(l)!0.5!(g)+(0.5,0.4)$) {$m$}; 
                \node[ font=\huge] (lg) at ($(l)!0.5!(g)+(0.1,0.4)$) {\( \downarrowtail \)}; 
                \node[ font=\huge] (rh) at ($(r)!0.5!(h)+(0.1,0.4)$) {\( \downarrowtail \)}; 
            %   \node[ font=\LARGE] () at ($(r)!0.5!(h)+(0.55,0.4)$) {$m'$}; 
            %   \node[ font=\LARGE] () at ($(k)!0.5!(c)+(0.5,0.4)$) {$u$}; 
            \end{tikzpicture}
          }
%         \caption{}
%         \label{fig:subgraph_counting:decomposition_of_rewriting_step}
%    \end{figure}
    \end{center}
\begin{definition}
    \label{def:x_occurrence}
     Let $X$ be a subgraph of $L$.  An \textbf{$X$-occurrence} in a graph $G$ is a \textbf{monomorphism} from $X$ to $G$. An $X$-occurrence in $G$ is 
    \begin{itemize}
        \item \textbf{shared} if its image is included in $C$,
        \item \textbf{explicit} if its image is  included in $L$,
        \item \textbf{implicit} if its image has elements in both $C'$ and $L'$.
    \end{itemize}  
    An $X$-occurrence in $H$ is 
    \begin{itemize}
        \item \textbf{shared} if its image is included in $C$,
        \item \textbf{explicit} if its image is included in $R$,
        \item \textbf{implicit} if its image has elements in both $C'$ and $R'$.
    \end{itemize} 

    Intuitively, the rewriting step implicitly creates the implicit $X$-occurrences in $H$ and implicitly destroys implicit $X$-occurrences in $G$. 
    
    Given an implicit $X$-occurrence $x$ in $H$, its image can be decomposed as the union $R_x \cup C_x$, where $R_x = \opn{Im}(x) \cap R$ (the part of the image overlapping the right-hand side R), and $C_x = \opn{Im}(x) \cap C$ (the part overlapping the context C).
    
    For every graph monomorphism $h : R_x \rightarrowtail L$ that preserves interface elements (i.e., \( h(k) = k \) for all \( k \in K \)), we define the \textbf{corresponding $X$-occurrence of $x$ relative to $h$} as the $X$-occurrence $y$ in $G$ given by:
        $$
        y(z) = 
        \begin{cases}
        z & \text{if } z \in x^{-1}(C_x), \\
        h(x(z)) & \text{if } z \in x^{-1}(R_x \setminus C_x).
        \end{cases}
        $$
   \end{definition} 



% \begin{definition}
%     \label{def:corresponding_occurrence}
%      An \textbf{$X$-occurrence} in $G$ is a subgraph of $G$ isomorphic to $X$. An \textbf{$X$-occurrence $x$ implicitly created by the rewriting step} is an \( X \)-occurrence in $H$ which is included neither in \( R \) nor in $C$. 
%       The occurrence $x$ can be decomposed as the union of graphs $x = R_x \cup C_x$ where $R_x = x \cap R$ and $C_x = x \cap C$. For every graph monomorphism $h : R_x \rightarrowtail L$ that preserves interface elements (i.e., \( h(k) = k \) for all \( k \in K \)), we define the \textbf{corresponding $X$-occurrence of $x$ relative to $h$} as the $X$-occurrence in $G$ given by $h(R_x) \cup C_x$.
%     An \textbf{$X$-occurrence implicitly destroyed by the rewriting step} is an \( X \)-occurrence in $G$ which is included neither in \( L \) nor in $C$.
% \end{definition}


\begin{example}
    Let $X$ be the graph
\tikz[baseline=-0.5ex]{ 
        \node[draw,circle] (x) at (0,0) {};  
        \node[draw,circle] (y) at (1,0) {};
        \node[draw,circle] (z) at (2,0) {};
        \draw[->] (x) -- node[midway,above] {$a$} (y) ;
        \draw[->] (y) -- node[midway,above] {$a$} (z) ;
}.
Consider the $X$-occurrences in the decomposition of a rewriting step using the rule in Example~\ref{ex:grsaa} and recalled below:
% in Figure~\ref{fig:subgraph_counting:occurrdlfdjslfkjsa}.
    % \begin{figure}[H]
    %   \centering
    \begin{center}
      \resizebox{0.7\textwidth}{!}{
      \begin{tikzpicture}
          \graphbox{$L$}{0mm}{0mm}{34mm}{15mm}{2mm}{-5mm}{
              \coordinate (o) at (0mm,-3mm); 
              \node[draw,circle] (l1) at ($(o)+(-10mm,0mm)$) {1};
              \node[draw,circle] (l2) at ($(l1)+(2,0)$) {2};
              \node[draw,circle] (l3) at ($(l1) + (1,0)$) {3};
              \draw[->] (l1) -- (l3) node[midway,above] {$a$};
              \draw[->] (l3) -- (l2) node[midway,above] {$a$};
          }     
          \graphbox{$K$}{40mm}{0mm}{24mm}{15mm}{2mm}{-5mm}{
              \coordinate (o) at (5mm,-3mm); 
              \node[draw,circle] (l1) at ($(o)+(-10mm,0mm)$) {1};
              \node[draw,circle] (l2) at ($(l1)+(1,0)$) {2};
              % \node[draw,circle] (l3) at ($(l1) + (1,0)$) {$\ $};
              % \draw[->] (l1) -- (l3) node[midway,above] {$a$};
              % \draw[->] (l3) -- (l2) node[midway,above] {$a$};
          }    
          \graphbox{$R$}{70mm}{0mm}{45mm}{15mm}{2mm}{-5mm}{
              \coordinate (o) at (-5mm,-3mm); 
              \node[draw,circle] (l1) at ($(o)+(-10mm,0mm)$) {1};
              \node[draw,circle] (l2) at ($(l1)+(3,0)$) {2};
              \node[draw,circle] (l3) at ($(l1) + (1,0)$) {4};
              \node[draw,circle] (l4) at ($(l1) + (2,0)$) {5};
              \draw[->] (l1) -- (l3) node[midway,above] {$a$};
              \draw[->] (l3) -- (l4) node[midway,above] {$b$};
              \draw[->] (l4) -- (l2) node[midway,above] {$a$};
          }    
          \node () at (37mm,-8mm) {$\overset{l}{\leftarrowtail}$};
          \node () at (67mm,-8mm) {$\overset{r}{\rightarrowtail}$};
          % \draw[>->] (51mm,2mm) -- (52mm,3mm);
      \end{tikzpicture}
      }
%       \caption{}
%       \label{fig:subgraph_counting:occurrdlfdjslfkjsa}
%   \end{figure}
    \end{center}
   We illustrate the aforementioned decomposition
    % in Figure~\ref{fig:subgraph_counting:occurrendfsdadsgces_of_x} 
    using the visual notation for graph homomorphisms introduced in Notation~\ref{notation:graph_homomorphism}, with distinct pregraphs highlighted in different colors.
% \begin{figure}[H]
%     \centering
    \begin{center}
    \resizebox{0.8\textwidth}{!}{
      \begin{tikzpicture}
              \graphbox{\( L = L' \cup K \)}{0mm}{5mm}{34mm}{20mm}{2mm}{-5mm}{
                  \coordinate (o) at (0mm,-8mm); 
                  \node[draw,circle] (l1) at ($(o)+(-10mm,0mm)$) {1};
                  \node[draw,circle] (l2) at ($(l1)+(2,0)$) {2};
                  \node[orange,draw,circle] (l3) at ($(l1) + (1,0)$) {3};
                  \draw[orange,->] (l1) -- (l3) node[midway,above] {$a$};
                  \draw[orange,->] (l3) -- (l2) node[midway,above] {$a$};
              } 

              \graphbox{\( K \)}{40mm}{5mm}{34mm}{20mm}{2mm}{-5mm}{
                  \coordinate (o) at (0mm,-8mm); 
                  \node[draw,circle] (l1) at ($(o)+(-10mm,0mm)$) {1};
                  \node[draw,circle] (l2) at ($(l1)+(2,0)$) {2};
              }  

              \graphbox{\( R = R' \cup K \)}{80mm}{5mm}{45mm}{20mm}{2mm}{-5mm}{
                  \coordinate (o) at (-5mm,-8mm); 
                  \node[draw,circle] (l1) at ($(o)+(-10mm,0mm)$) {1};
                  \node[draw,circle] (l2) at ($(l1)+(3,0)$) {2};
                  \node[red,draw,circle] (l3) at ($(l1) + (1,0)$) {4};
                  \node[red,draw,circle] (l4) at ($(l1) + (2,0)$) {5};
                  \draw[red,->] (l1) -- (l3) node[midway,above] {$a$};
                  \draw[red,->] (l3) -- (l4) node[midway,above] {$b$};
                  \draw[red,->] (l4) -- (l2) node[midway,above] {$a$};
              }    

              \graphbox{\( G = L' \cup K \cup C' \)}{0mm}{-22mm}{34mm}{30mm}{2mm}{-10mm}{
                  \coordinate (o) at (0mm,-3mm); 
                  \node[draw,circle] (l1) at ($(o)+(-10mm,0mm)$) {1};
                  \node[draw,circle] (l2) at ($(l1)+(2,0)$) {2};
                  \node[draw,circle,orange] (l3) at ($(l1) + (1,0)$) {3};
                  \node[blue, draw,circle] (l4) at ($(l2) + (0,-1)$) {6};
                  \draw[orange,->] (l1) -- (l3) node[midway,above] {$a$};
                  \draw[orange,->] (l3) -- (l2) node[midway,above] {$a$};
                  \draw[blue,->] (l2) -- (l4) node[midway,right] {$a$};
                %   \node[blue,draw,circle] (l6) at ($(l1) + (0,-1)$) {7};
                %   \draw[blue,<-] (l1) -- (l6) node[midway,left] {$a$};
                \node[blue,draw,circle] (l7) at ($(l4) + (-1,0)$) {7};
                  \draw[blue,->] (l4) -- (l7) node[midway,below] {$a$};
              }    

              \graphbox{\( C = K \cup C' \)}{40mm}{-22mm}{34mm}{30mm}{2mm}{-10mm}{
                  \coordinate (o) at (0mm,-3mm); 
                  \node[draw,circle] (l1) at ($(o)+(-10mm,0mm)$) {1};
                  \node[draw,circle] (l2) at ($(l1)+(2,0)$) {2};
                  \node[blue,draw,circle] (l4) at ($(l2) + (0,-1)$) {6};
                  \draw[blue,->] (l2) -- (l4) node[midway,right] {$a$};
                %   \node[blue,draw,circle] (l6) at ($(l1) + (0,-1)$) {7};
                %   \draw[blue,<-] (l1) -- (l6) node[midway,left] {$a$};
                    \node[blue,draw,circle] (l7) at ($(l4) + (-1,0)$) {7};
                  \draw[blue,->] (l4) -- (l7) node[midway,below] {$a$};
              }    

              \graphbox{\( H = R' \cup K \cup C' \)}{80mm}{-22mm}{45mm}{30mm}{2mm}{-10mm}{
                  \coordinate (o) at (-5mm,-3mm); 
                  \node[draw,circle] (l1) at ($(o)+(-10mm,0mm)$) {1};
                  \node[draw,circle] (l2) at ($(l1)+(3,0)$) {2};
                  \node[draw,circle,red] (l3) at ($(l1) + (1,0)$) {4};
                  \node[draw,circle,red] (l4) at ($(l1) + (2,0)$) {5};
                  \node[blue,draw,circle] (l5) at ($(l2) + (0,-1)$) {6};
                %   \node[blue,draw,circle] (l6) at ($(l1) + (0,-1)$) {7};
                %   \draw[blue,<-] (l1) -- (l6) node[midway,left] {$a$};
                  \draw[red,->] (l1) -- (l3) node[midway,above] {$a$};
                  \draw[red,->] (l3) -- (l4) node[midway,above] {$b$};
                  \draw[red,->] (l4) -- (l2) node[midway,above] {$a$};
                  \draw[blue,->] (l2) -- (l5) node[midway,right] {$a$};
                        \node[blue,draw,circle] (l7) at ($(l5) + (-1,0)$) {7};
                  \draw[blue,->] (l5) -- (l7) node[midway,below] {$a$};
              }    

              \node () at (37mm,-8mm) {\( \leftarrowtail \)}; % K -> L
              \node () at (77mm,-8mm) {\( \rightarrowtail \)}; % K -> R
              \node () at (17mm,-18mm) {\( m\ \downarrowtail \)};
              \node () at (37mm,-33mm) {\( \leftarrowtail \)};
              \node () at (52mm,-18mm) {\( \downarrowtail \)};
              \node () at (92mm,-18mm) {\( \downarrowtail \)};
              \node () at (77mm,-33mm) {\( \rightarrowtail \)}; % C -> H
      \end{tikzpicture}
      }
%       \caption{}
%       \label{fig:subgraph_counting:occurrendfsdadsgces_of_x}
%    \end{figure}
    \end{center}
    In this example, the graph $X$ has 4 distinct $X$-occurrences in $G$ and $H$.
   There are 
\begin{itemize}
    \item 1 occurrence shared by $G$ and $H$: 
     \tikz[baseline=-0.5ex]{ 
        \node[draw,circle] (x) at (0,0) {2};  
        \node[blue,draw,circle] (y) at (1,0) {6};
        \node[blue,draw,circle] (z) at (2,0) {7};
        \draw[blue,->] (x) -- node[midway,above] {$a$} (y) ;
        \draw[blue,->] (y) -- node[midway,above] {$a$} (z) ;
     }
    \item 1 explicit occurrence in $G$:
     \tikz[baseline=-0.5ex]{ 
        \node[draw,circle] (x) at (0,0) {1};  
        \node[orange,draw,circle] (y) at (1,0) {3};
        \node[draw,circle] (z) at (2,0) {2};
        \draw[orange,->] (x) -- node[midway,above] {$a$} (y) ;
        \draw[orange,->] (y) -- node[midway,above] {$a$} (z) ;
    }
 
    \item 1 implicit occurrences in $G$:
    \tikz[baseline=-0.5ex]{ 
        \node[orange,draw,circle] (x) at (0,0) {3};  
        \node[draw,circle] (y) at (1,0) {2};
        \node[blue,draw,circle] (z) at (2,0) {6};
        \draw[orange,->] (x) -- node[midway,above] {$a$} (y) ;
        \draw[blue,->] (y) -- node[midway,above] {$a$} (z) ;
    } 
 
    \item 0 explicit occurrence in $H$,
    \item 1 implicit occurrence in $H$: 
    \tikz[baseline=-0.5ex]{ 
        \node[red,draw,circle] (x) at (0,0) {5};  
        \node[draw,circle] (y) at (1,0) {2};
        \node[blue,draw,circle] (z) at (2,0) {6};
        \draw[red,->] (x) -- node[midway,above] {$a$} (y) ;
        \draw[blue,->] (y) -- node[midway,above] {$a$} (z) ;
}
 
\end{itemize}
    The corresponding X-occurrence of the implicit occurrences \tikz[baseline=-0.5ex]{ 
        \node[red,draw,circle] (x) at (0,0) {5};  
        \node[draw,circle] (y) at (1,0) {2};
        \node[blue,draw,circle] (z) at (2,0) {6};
        \draw[red,->] (x) -- node[midway,above] {$a$} (y) ;
        \draw[blue,->] (y) -- node[midway,above] {$a$} (z) ;
} in $H$ relative to the unique monomorphism from 
    \tikz[baseline=-0.5ex]{
      \node[red,draw,circle] (x) at (0,0) {5};
      \node[draw,circle] (y) at (1,0) {2};
      \draw[red,->] (x) -- (y) node[midway, above] {$a$};
  }
     to 
    \tikz[baseline=-0.5ex]{
      \node[orange,draw,circle] (x) at (0,0) {3};
      \node[draw,circle] (y) at (1,0) {2};
      \draw[orange,->] (x) -- (y) node[midway, above] {$a$};
  }
    is the implicit occurrences 
       \tikz[baseline=-0.5ex]{ 
        \node[orange,draw,circle] (x) at (0,0) {3};  
        \node[draw,circle] (y) at (1,0) {2};
        \node[blue,draw,circle] (z) at (2,0) {6};
        \draw[orange,->] (x) -- node[midway,above] {$a$} (y) ;
        \draw[blue,->] (y) -- node[midway,above] {$a$} (z) ;
    } 
    in $G$.
\end{example}

% In this section, we present a method for termination analysis of DPO graph rewriting systems with injective rules and monic matches. 
% Since it is an interpretation method \cite{nipkow1998term}, we first define the weight of a graph in \autoref{def:graph_weight}. 
% With this definition, we explain the general idea of the method and identify the key challenge in \autoref{sec:general_idea}. 
% In \autoref{sec:non-increasing}, we introduce the concept of non-increasing rules. 
% Using this concept, a solution to this challenge is then proposed in \autoref{sec:solution_to_the_key_challenge}. 
% Finally, in \autoref{sec:termination}, we state our termination criterion.

\section{Measuring graphs by counting subgraphs}
\label{sec:interpretation}
Let $\mathcal{R} = \mathcal{A} \cup \mathcal{B}$ be a set of injective rules.
Define $\mathrm{Sub}(\mathrm{left}(\mathcal{R}))$ as the set of all subgraphs of left-hand sides graphs of rules in $\mathcal{R}$.
We call elements of $\mathrm{Sub}(\mathrm{left}(\mathcal{R}))$ \textbf{ruler-graphs}. Analogous to measuring physical objects with rulers, we measure graphs by counting occurrences of these ruler-graphs.

\begin{definition}
    % Let \( X \) be a ruler-graph and \( G \) a graph. The \textbf{measurement} of \( G \) with respect to \( X \), denoted \( m_X(G) \), is 
    For a ruler-graph \( X \) and graph \( G \), the \textbf{measurement} \( m_X(G) \) is:
    \(
        m_X(G) = \card{\operatorname{Mono}(X, G)}
    \).
\end{definition} 
The occurrence count of ruler-graphs before and after a rewriting step provides a termination criterion: 
    If every rewriting step using a rule in $\mathcal{A}$ strictly decreases the number of $X$'s occurrences, and every rewriting step using a rule in $\mathcal{B}$ never increases it, then the rewriting relation induced by $\mathcal{A}$ is terminating relative to the rewriting relation induced by $\mathcal{B}$.
 
For some ruler-graphs, occurrence changes during a rewriting step using an injective rule are precisely computable by simply counting their occurrences in the left- and right-hand side graphs. Examples include: the ruler-graph \tikz[baseline=-0.5ex]{
        \node (x) at (0,0) {$\bullet$};
    } consisting of a single node; the ruler-graph \tikz[baseline=-0.5ex]{
        \node (x) at (0,0) {$\bullet$};
        \draw[->] (x) edge[loop right] node[midway, right] {$x$} (x) ;
    } consisting of a single node and a single loop labeled by $x\in \Sigma$; the ruler-graph \tikz[baseline=-0.5ex]{
        \node (x) at (0,0) {$\bullet$};
        \node (y) at (1,0) {$\bullet$ };
        \draw[->] (x) -- (y) node[midway, above] {$x$};
    } consisting of two nodes and a single labeled edge between them.  
    
However, systems like \autoref{ex:grsaa_rx} require more complex ruler-graphs such as \tikz[baseline=-0.5ex]{
        \node (x) at (0,0) {$\bullet$};
        \node (y) at (1,0) {$\bullet$ };
        \node (z) at (2,0) { $\bullet$};
        \draw[->] (x) -- (y) node[midway, above] {$a$};
        \draw[->] (y) -- (z) node[midway, above] {$a$};
    } as shown in \autoref{ex:termination:grsaa}.
    For such ruler-graph, occurrence changes during rewriting is unpredictable in general, as their occurrence creation and destruction depend on the context graph\textemdash a challenge we address in later sections.
    
To combine measurements from multiple ruler-graphs (analogous to combining measurements of physical objects from multiple rulers), we introduce:
\begin{definition}[Weight function]
    \label{def:weight_function}
    A \textbf{weight function} for a set of ruler-graphs \( \mathbb{X} \) is a map \( s_{\mathbb{X}} \colon \mathbb{X} \to \mathbb{N} \) assigning a weight in $\mathbb{N}$ to each ruler-graph.
\end{definition}
\autoref{ex:overbeek_5d6} demonstrates the usefulness of a weight function that assigns different weight to measurements from different ruler-graphs.
\begin{definition}[Graph weight]
    \label{def:graph_weight}  
    Let \( \mathbb{X} \) be a set of ruler-graphs, 
    \( s_{\mathbb{X}} \) a weight function, and \( G \) a graph. The \textbf{weight of graph} \( G \) relative to \( s_{\mathbb{X}} \), denoted \( w_{s_{\mathbb{X}}}(G) \), is defined as: 
    \[
        w_{s_{\mathbb{X}}}(G) = \sum_{X \in \mathbb{X}} s_{\mathbb{X}}(X) \cdot m_X(G).  
    \]  
\end{definition}


% \section{Measuring graphs with ruler-Graphs}
% \label{sec:interpretation}
% 
%     Let $\mathcal{R}$ be a set of injective DPO graph rewriting rules and $\mathrm{Sub}(\mathrm{left}(\mathcal{R}))$ the set of subgraphs of the left-hand-side graphs of the rules in $\mathcal{R}$. 
%     Analogous to measuring vectors using a basis within a vector space, 
%     or measuring physical objects using rulers, 
%     we measure graphs using ruler-graphs\textemdash special graphs in $\mathrm{Sub}(\mathrm{left}(\mathcal{R}))$ that act as reference structures. 
%     A ruler-graph measures a graph $G$, and provides a measurement of $G$ relative to that ruler-graph.  
%     Similarly, analogous to aggregating the coordinates of a vector to compute a weight, or combining measurements of physical objects provided by multiple rulers to obtain a weight, we summarize the measurements across multiple ruler-graphs as the weight of a graph.
% }
% \todo{A: all of this is a little bit over-formalized}
% \begin{definition}[Ruler-graph]
%     A \textbf{ruler-graph} is a graph in $\mathrm{Sub}(\mathrm{left}(\mathcal{R}))$.
% \end{definition} 
% \begin{definition}[Measurement]
%     Let \( X \) be a ruler-graph and \( G \) a graph. The \textbf{measurement} of \( G \) with respect to \( X \), denoted \( m_X(G) \), is 
%     \(
%         m_X(G) = \big|\!\operatorname{Mono}(X, G)\!\big| 
%     \).
% \end{definition}
% \begin{definition}[Weight function]
%     \label{def:weight_function}
%     A \textbf{weight function} for a set of ruler-graphs \( \mathbb{X} \) is a map \( s_{\mathbb{X}} \colon \mathbb{X} \to \mathbb{N} \) assigning a weight in $\mathbb{N}$ to each ruler-graph.
% \end{definition}
% \begin{definition}[Graph weight]
%     \label{def:graph_weight}  
%     % Let \( \mathcal{R} \) be a set of graph rewriting rules.
%     % Let $\mathbb{X}$ be a set of ruler-graphs,  
%     % \( s_\mathbb{X} \colon \mathbb{X} \rightarrow \mathbb{N} \) a weight function, and \( G \) a graph. 
%     % The \textbf{weight \( w_{s_\mathbb{X}}(G) \)} of \( G \) is the weighted sum of the measurements of \( G \) provided by rulers in \( \mathbb{X} \):
%     % \[
%     %     w_{s_\mathbb{X}}(G) = \sum_{X \in \mathbb{X}} s_\mathbb{X}(X) \cdot \big|\!\operatorname{Mono}(X, G)\!\big|.
%     % \]
%     Let \( \mathbb{X} \) be a set of ruler-graphs, 
%     \( s_{\mathbb{X}} \) a weight function, and \( G \) a graph. The \textbf{weight of graph} \( G \) relative to \( s_{\mathbb{X}} \), denoted \( w_{s_{\mathbb{X}}}(G) \) (or simply $w(G)$ when $s_{\mathbb{X}}$ is clear from context), is defined as: 
%     \[
%         w_{s_{\mathbb{X}}}(G) = \sum_{X \in \mathbb{X}} s_{\mathbb{X}}(X) \cdot m_X(G).  
%     \]  
% \end{definition}
% % \begin{remark}
% %     This work focuses on deriving sufficient conditions for termination, deferring the construction of ruler-graphs to future work.
% %     %  however, considering a set of subgraphs of the left-hand-side graphs of the rewriting rules is already fruitful, as demonstrated in \autoref{sec:related_work}.
% % \end{remark}
% % \begin{remark}
%     Trivial ruler-graphs include graphs with a single node and no edges, graphs with a single node and a single labeled loop, and graphs with two nodes and a single labeled edge between them. 
%     For any such trivial ruler-graph, the change in the number of its occurrences in a graph rewriting step using a rule in $\mathcal{R}$ can be precisely computed by simply counting its occurrences in the left- and right-hand-side graphs. }\todo{reviewer 1: The choice of the set of ruler graphs is rather opaque. Intuitively, it should be related to the left-hand sides of the rules in $A$ and $B$. A discussion on this would be valuable before introducing ruler graphs.}




\section{General idea and key challenge}
\label{sec:general_idea}
For the remainder of this section, we assume: a set \( \mathcal{R} = \mathcal{A} \cup \mathcal{B} \) of injective DPO rewriting rules, a set \( \mathbb{X} \) of ruler-graphs, and a weight function \( s_{\mathbb{X}} \).  Let \( \rho = (L \overset{l}{\leftarrowtail} K \overset{r}{\rightarrowtail} R) \) denote a rule in \( \mathcal{R} \) and \( X \in \mathbb{X} \).

The rewriting relation \( \Rightarrow_{\mathcal{A},\mathfrak{M}} \) is terminating relative to $\Rightarrow_{\mathcal{B},\mathfrak{M}}$ if the following hold: (i) for all \(G \Rightarrow_{\mathcal{A},\mathfrak{M}} H\), \( w_{s_\mathbb{X}}(G) > w_{s_\mathbb{X}}(H) \), and (ii) for all \(G \Rightarrow_{\mathcal{B},\mathfrak{M}} H\), \( w_{s_\mathbb{X}}(G) \geq w_{s_\mathbb{X}}(H) \). However, directly verifying all rewriting steps is infeasible due to their potential infiniteness. 
Thus, we aim to provide a rule-based sufficient condition for termination.
% \begin{remark}
%     To simplify our presentation, we treat the following three notions interchangeably when the monomorphism \( h: A \to B \) is understood: the domain graph \( A \), the monomorphism \( h: A \to B \), the image \( h(A) \subseteq B \) (an occurrence of \( A \) in \( B \)).
%     Consequently, in a rewriting step defined by the diagram in \autoref{def:rewriting_step}, we streamline descriptions involving subgraph images. Rather than stating that there is an occurrence of a ruler-graph \( X \) formed by the image of subgraphs \( R' \subseteq R \) and \( C' \subseteq C \) around the image of \( K' \subseteq K \), we simply say \( X \) is formed by \( R' \) and \( C' \) around \( K' \). Rather than stating that $b \notin h(A)$, we simply say that $b \notin A$.
% \end{remark}
\begin{notation}
    % The disjoint union of two sets \( S \) and \( S' \) will be denoted by \( S \uplus S' \). Let \( X, A, B, G \) be graphs, and let \( \alpha \colon A \to G \) and \( \beta \colon B \to G \) be morphisms. We define the following sets of monomorphisms of $X$ in $G$: 
\ \newline
\noindent
\begin{minipage}{0.69\textwidth}
\vspace{2mm}
\noindent  The disjoint union of two sets \( S \) and \( S' \) will be denoted by \( S \uplus S' \). Let \( X, A, B, G \) be graphs, and let \( \alpha \colon A \to G \) and \( \beta \colon B \to G \) be morphisms as illustrated on the right. We define the following sets of monomorphisms from $X$ in $G$ based on their interaction with $\alpha$ and $\beta$:    
\end{minipage}%
\begin{minipage}{0.30\textwidth}
    \hfill
        \resizebox{0.9\textwidth}{!}{
        % \begin{tikzpicture}
        %     \node (B) at (2,0) {$A$}; 
        %     \node (C) at (0,2) {$B$}; 
        %     \node (D) at (0,0) {$G$}; 
        %     \draw [>->] (B) to node [below,label,pos=0.45] {$\alpha$} (D); 
        %     \draw [>->] (C) to node [left,label,pos=0.45] {$\beta$} (D);
        % \end{tikzpicture}
            \begin{tikzpicture}
            \node (X) at (-1.5,-1) {$X$};
            \node (B) at (2,0) {$A$}; 
            \node (C) at (0,2) {$B$}; 
            \node (D) at (0,0) {$G$}; 
            \draw [>->] (B) to node [above,label,pos=0.45] {$\alpha$} (D); 
            \draw [>->] (C) to node [right,label,pos=0.45] {$\beta$} (D);
            \draw [>->,dashed] (X) to node [above,label,pos=0.45] {$\iota$} (D);
            \draw [>->,dashed] (X) to node [below,label,pos=0.45] {$\zeta$} (B);
            \draw [>->,dashed] (X) to node [above,label,pos=0.45] {$\eta$} (C);
        \end{tikzpicture}
        }
\end{minipage} 
    \begin{align*}
        \operatorname{Mono}(X,G,\alpha) &= \left\{ \iota \colon X \rightarrowtail G \;\middle|\; \exists \zeta \colon X \rightarrowtail A.\, \iota = \zeta \star \alpha \right\}
        \\
        \operatorname{Mono}(X,G,\lnot \alpha) &= \left\{ \iota \colon X \rightarrowtail G \;\middle|\; \nexists \zeta \colon X \rightarrowtail A.\, \iota = \zeta \star \alpha \right\}
        \\
        \operatorname{Mono}(X,G,\lnot \alpha, \beta) &= \left\{ 
            \iota \colon X \rightarrowtail G \;\middle|\; 
                \begin{aligned}  
                    &(\nexists \zeta \colon X \rightarrowtail A.\, \iota = \zeta \star \alpha) \\ 
                    &\land (\exists \eta \colon X \rightarrowtail B.\, \iota = \eta \star \beta)
                \end{aligned}
        \right\}
        \\
        \operatorname{Mono}(X,G,\lnot \alpha, \lnot \beta) &= \left\{ 
            \iota \colon X \rightarrowtail G \;\middle|\; 
                \begin{aligned}
                    &(\nexists \zeta \colon X \rightarrowtail A.\, \iota = \zeta \star \alpha) \\
                    &\land (\nexists \eta \colon X \rightarrowtail B.\, \iota = \eta \star \beta)
                \end{aligned}
        \right\}
    \end{align*}
\end{notation}
\noindent
\begin{minipage}{0.7\textwidth}\setlength{\parindent}{1em}
\noindent Consider a rewriting step \( G \Rightarrow_{\rho,\mathfrak{M}} H \) defined by the DPO diagram illustrated on the right. The sets of monomorphisms \( \operatorname{Mono}(X,G) \) and \( \operatorname{Mono}(X,H) \) can be decomposed into the following disjoint unions:
\end{minipage}%
\begin{minipage}{0.29\textwidth}
    \hfill
    % \begin{tikzpicture}
    %     % [node distance=15mm]
    %     \node (I) {$K$};
    %     \node (L) [left of=I] {$L$};
    %     \node (R) [right of=I] {$R$};
    %     \node (G) [below of=L] {$G$};
    %     \node (C) [below of=I] {$C$};
    %     \node (H) [below of=R] {$H$};
    %     \draw [>->] (I) to  node [midway,below] {$l$} (L);
    %     \draw [>->] (I) to  node [midway,below] {$r$} (R);
    %     \draw [>->] (L) to node [midway,right] {$m$} (G);
    %     \draw [>->] (I) to  node [midway,right] {} (C);
    %     \draw [>->] (R) to  node [midway,right] {$m'$} (H);
    %     \draw [>->] (C) to node [midway,above] {$l'$} (G);
    %     \draw [>->] (C) to node [midway,above] {$r'$} (H);
    %     % \node [at=($(I)!.5!(G)$)] {\normalfont PO};
    %     % \node [at=($(I)!.5!(H)$)] {\normalfont PO};
    %   \end{tikzpicture}
    \resizebox{0.9\textwidth}{!}{
     \begin{tikzpicture}
            % [node distance=15mm]
            \node (I) at (0,0) {$K$};
            \node (L)  at (-2,0) {$L$};
            \node (R)  at (2,0) {$R$};
            \node (G)  at (-2,-2) {$G$};
            \node (C)  at (0,-2) {$C$};
            \node (H)  at (2,-2) {$H$};
            \draw [>->] (I) to  node [midway,below] {$l$} (L);
            \draw [>->] (I) to  node [midway,below] {$r$} (R);
            \draw [>->] (L) to node [midway,right] {$m$} (G);
            \draw [>->] (I) to  node [midway,right] 
            % {$u$}
            {} (C);
            \draw [>->] (R) to  node [midway,right] 
            {$m'$}
            (H);
            \draw [>->] (C) to node [midway,above] {$l'$} (G);
            \draw [>->] (C) to node [midway,above] 
            {$r'$} 
            (H);
            % \node [at=($(I)!.5!(G)$)] {\normalfont PO};
            % \node [at=($(I)!.5!(H)$)] {\normalfont PO};
        \end{tikzpicture}
    }
\end{minipage}


\begin{flalign*}
    \operatorname{Mono}(X,G) &= 
    \operatorname{Mono}(X,G,l')
    \uplus
    \operatorname{Mono}(X,G,\lnot l',m) 
    \uplus
    \operatorname{Mono}(X,G,\lnot l',\lnot m)
    \\
    \operatorname{Mono}(X,H) &= 
    \operatorname{Mono}(X,H,r')
    \uplus
    \operatorname{Mono}(X,H,\lnot r',m') 
    \uplus
    \operatorname{Mono}(X,H,\lnot r',\lnot m')
\end{flalign*}
Thus, we have:
\begin{flalign*}
    \card{\operatorname{Mono}(X,G)} ={} &
    \card{\operatorname{Mono}(X,G,l')}
    +
    \card{\operatorname{Mono}(X,G,\lnot l',m)} \\
    &+
    \card{\operatorname{Mono}(X,G,\lnot l',\lnot m)}
    \\
    \card{\operatorname{Mono}(X,H)} ={} &
    \card{\operatorname{Mono}(X,H,r')} 
    +
    \card{\operatorname{Mono}(X,H,\lnot r',m')} \\
    &+\card{\operatorname{Mono}(X,H,\lnot r',\lnot m')}
\end{flalign*}
\noindent Using the following \autoref{lem:decomp_w_u} 
(see \iflongversion
Appendix \autoref{proof:dcomp_w_u} 
\else
\cite[Lemma 27]{qiu2025termination}
\fi for a proof), we derive:
\begin{flalign*}
&\card{\operatorname{Mono}(X,G)} =
\card{\operatorname{Mono}(X,C)}
+
\card{\operatorname{Mono}(X,L, \lnot l)}
+
\card{\operatorname{Mono}(X,G, \lnot l',\lnot m)} 
 \\
 &\card{\operatorname{Mono}(X,H)} =
 \card{\operatorname{Mono}(X,C)}
 +
 \card{ \operatorname{Mono}(X,R, \lnot r) }
 +
 \card{ \operatorname{Mono}(X,H, \lnot r',\lnot m')}  
\end{flalign*}
\begin{lemma}
    \label{lem:decomp_w_u}
    \ \newline
    \noindent
    \begin{minipage}{0.69\textwidth}
        Let $X$ be a ruler-graph. For a pushout square as shown on the right, we have 
        % $\card{\operatorname{Mono}(X, B)} = \card{\operatorname{Mono}(X, D, \beta')}$ and $\card{\operatorname{Mono}(X, C, \lnot \beta)} = \card{\operatorname{Mono}(X, D, \lnot \beta', \alpha')}$.
        \begin{flalign*}
            \card{\operatorname{Mono}(X, B)} &= \card{\operatorname{Mono}(X, D, \beta')}
            \\
            \card{\operatorname{Mono}(X, C, \lnot \beta)} &= \card{\operatorname{Mono}(X, D, \lnot \beta', \alpha')}
        \end{flalign*}
    \end{minipage}
    \hfill
    \begin{minipage}{0.29\textwidth}
        \hfill
        \begin{tikzpicture}
            \node (A) {$A$};
            \node [below of=A] (B) {$B$}; 
            \node [left of=A] (C) {$C$}; 
            \node [left of=B] (D) {$D$}; 
            \begin{scope}[nodes=rectangle]          
            \draw [>->] (A) to node [right,label,pos=0.5] {$\alpha$} (B);
            \draw [>->] (A) to node [above,label,pos=0.5] {$\beta$} (C);
            \draw [>->] (B) to node [below,label,pos=0.45] {$\beta'$} (D); 
            \draw [>->] (C) to node [left,label,pos=0.45] {$\alpha'$} (D);
            \end{scope}
        \end{tikzpicture}
    \end{minipage} 
\end{lemma}
Therefore, we obtain:
\begin{flalign*}
      &\card{\operatorname{Mono}(X,G)} -  \card{\operatorname{Mono}(X,H)} \\
     =& \card{\operatorname{Mono}(X,L, \lnot l)}- \card{\operatorname{Mono}(X,R, \lnot r)}  + \\ 
     &\card{\operatorname{Mono}(X,G, \lnot l',\lnot m)} - 
     \card{\operatorname{Mono}(X,H, \lnot r',\lnot m') } 
\end{flalign*}
By Lemma~\ref{lem:xlnlmxrnr} (see \iflongversion
    Appendix \autoref{proof:lem:xlnlmxrnr}
    \else
    \cite[Lemma 28]{qiu2025termination}
    \fi for a proof), it can be simplified to:
\begin{flalign*}
     &\card{\operatorname{Mono}(X,G)} -  \card{\operatorname{Mono}(X,H)} \\
    = & \card{\operatorname{Mono}(X,L)}- \card{\operatorname{Mono}(X,R)}  + \\ 
      & \card{\operatorname{Mono}(X,G, \lnot l',\lnot m)} - 
    \card{\operatorname{Mono}(X,H, \lnot r',\lnot m') } 
\end{flalign*}
\begin{lemma}
    \label{lem:xlnlmxrnr}
    Let $X$ be a graph. Let $L \overset{l}{\leftarrowtail} K \overset{r}{\rightarrowtail} R$ be an injective DPO graph rewriting rule. We have 
    \[
       \card{\operatorname{Mono}(X, L, \lnot l)}  - \card{\operatorname{Mono}(X, R, \lnot r)} 
       = 
       \card{\operatorname{Mono}(X, L)}  - \card{\operatorname{Mono}(X, R)} 
        \]
\end{lemma}
Consequently, the weight difference satisfies:
\begin{flalign*}
     &w_{s_\mathbb{X}}(G) - w_{s_\mathbb{X}}(H)\\
    \overset{\operatorname{def}}{=}&\sum_{X \in \mathbb{X}}^{}s_\mathbb{X}(X) * m_X(G) - \sum_{X \in \mathbb{X}}^{}s_\mathbb{X}(X) * m_X(H)\\
    \overset{\operatorname{def}}{=}&\sum_{X \in \mathbb{X}}^{}s_\mathbb{X}(X) * |\operatorname{Mono}(X,G)| - \sum_{X \in \mathbb{X}}^{}s_\mathbb{X}(X) * |\operatorname{Mono}(X,H)|\\
    =&\sum_{X \in \mathbb{X}}^{}s_\mathbb{X}(X) * \left( \card{\operatorname{Mono}(X,G)} -  \card{\operatorname{Mono}(X,H)} \right)\\
    =&\sum_{X \in \mathbb{X}}^{}s_\mathbb{X}(X) * \left(  \card{\operatorname{Mono}(X,L)}- \card{\operatorname{Mono}(X,R)} \right)+\\
       & \sum_{X \in \mathbb{X}}^{}s_\mathbb{X}(X) * \left( 
     \card{\operatorname{Mono}(X,G, \lnot l',\lnot m)} - 
    \card{\operatorname{Mono}(X,H, \lnot r',\lnot m') } \right)
\end{flalign*}
Since $\sum_{X \in \mathbb{X}}^{}s(X) * \left(  \card{\operatorname{Mono}(X,L)}- \card{\operatorname{Mono}(X,R)} \right)$ can be precisely computed, the key challenge is to establish a lower bound for
\begin{flalign*}
    % \label{ineq:key_challenge}
    \sum_{X \in \mathbb{X}}^{} s_\mathbb{X}(X) * 
         \left ( \card{\operatorname{Mono}(X,G, \lnot l',\lnot m)} - 
     \card{\operatorname{Mono}(X,H, \lnot r',\lnot m') } \right ) 
    % \label{ineq:ge0_sss}
\end{flalign*} 
\section{Non-increasing rules}
\label{sec:non-increasing}
% \noindent
% \begin{minipage}{0.7\textwidth}
%     % \setlength{\parindent}{1em}
%     Let \( \rho = (L \overset{l}{\leftarrowtail} K \overset{r}{\rightarrowtail} R) \) be an injective DPO rewriting rule. 
%     A DPO diagram as illustrated on the right defines a rewriting step using $\rho$ with injective match \( m \).
% \end{minipage}%
% \hfill 
% \begin{minipage}{0.3\textwidth}
%     \hfill
%     % \begin{figure}[htbp] 
%     %     \center
%         \begin{tikzpicture}
%             % [node distance=15mm]
%             \node (I) {$K$};
%             \node (L) [left of=I] {$L$};
%             \node (R) [right of=I] {$R$};
%             \node (G) [below of=L] {$G$};
%             \node (C) [below of=I] {$C$};
%             \node (H) [below of=R] {$H$};
%             \draw [>->] (I) to  node [midway,below] {$l$} (L);
%             \draw [>->] (I) to  node [midway,below] {$r$} (R);
%             \draw [>->] (L) to node [midway,right] {$m$} (G);
%             \draw [>->] (I) to  node [midway,right] 
%             % {$u$}
%             {} (C);
%             \draw [>->] (R) to  node [midway,right] 
%             % {$m'$}
%             {} (H);
%             \draw [>->] (C) to node [midway,above] {$l'$} (G);
%             \draw [>->] (C) to node [midway,above] 
%             % {$r'$} 
%             {}
%             (H);
%             % \node [at=($(I)!.5!(G)$)] {\normalfont PO};
%             % \node [at=($(I)!.5!(H)$)] {\normalfont PO};
%         \end{tikzpicture}
%             % \caption{}
%     %         \label{diag:set_dpo_intuition_1}
%     % \end{figure}
% \end{minipage}

% \noindent
% \begin{minipage}{0.5\textwidth}
%     % \setlength{\parindent}{1em}
%      It follows that the DPO diagram depicted on the right, where all morphisms are inclusion functions and $H'$ is a graph isomorphic to $H$, holds.
% \end{minipage}%
% \hfill 
% \begin{minipage}{0.5\textwidth}  
%     \hfill
%     \begin{tikzpicture}
%         [node distance=20mm]
%         \node (I) at (0,0) {$\operatorname{Im}(l \star m)$};
%         \node (L) at (-2,0) {$\operatorname{Im}(m)$};
%         \node (R) at (2,0) {$\operatorname{Im}(m')$};
%         \node (G) at (-2,-1) {$G$};
%         \node (C) at (0,-1) {$\operatorname{Im}(l')$};
%         \node (H) at (2,-1) {$H'$};
%         \draw [>->] (I) to  node [midway,above] {} (L);
%         \draw [>->] (I) to  node [midway,above] {} (R);
%         \draw [>->] (L) to node [midway,right] {} (G);
%         \draw [>->] (I) to  node [midway,right] {} (C);
%         \draw [>->] (R) to  node [midway,right] {} (H);
%         \draw [>->] (C) to node [midway,above] {} (G);
%         \draw [>->] (C) to node [midway,above] {} (H);
%         % \node [at=($(I)!.5!(G)$)] {\normalfont PO};
%         % \node [at=($(I)!.5!(H)$)] {\normalfont PO};
%       \end{tikzpicture}
% \end{minipage}


%  \noindent
%     \begin{minipage}{0.6\textwidth}
%         By viewing a graph as a set o f nodes and edges and graph monomorphisms as set injections, the above DPO diagram can be considered as a DPO diagram in the category of sets, illustrated on the right.
%     \end{minipage}
%     \hfill
%     \begin{minipage}{0.3\textwidth}
%           \hfill
%           \resizebox{\textwidth}{!}{
%             \begin{tikzpicture}
%               \coordinate (k) at (0, 0);
%               \draw[fill=white] ($(k)+(0,0)$) rectangle ($(k)+(0.5,0.5)$);
%               \node () at ($(k)+(0.25,0.25)$) {\( \mathrm{K} \)};
          
%               \coordinate (c) at (0, -2.2);
%               \draw[fill=blue!20]
%               ($(c)+(0,-0.5)$)
%               -- ($(c)+(0,0.5)$) 
%               -- ($(c)+(1,0.5)$) 
%               arc[start angle=0, end angle=-90, radius=1]
%               -- cycle;
%               \node () at ($(c)+(0.75,0.25)$) {\( \mathrm{C'} \)};
%               \draw[fill=white] ($(c)+(0,0)$) rectangle ($(c)+(0.5,0.5)$);
%               \node () at ($(c)+(0.25,0.25)$) {\( \mathrm{K} \)};
          
%               \coordinate (l) at (-3, 0);
%               \draw[fill=orange!20] ($(l)+(-0.5,0)$) rectangle ($(l)+(0.5,1)$);
%               \node () at ($(l)+(-0.23,0.25)$) {\( \mathrm{L'} \)};
%               \draw[fill=white] ($(l)+(0,0)$) rectangle ($(l)+(0.5,0.5)$);
%               \node () at ($(l)+(0.25,0.25)$) {\( \mathrm{K} \)};
          
%               \coordinate (g) at (-3, -2.2);
%               \draw[fill=blue!20]
%               ($(g)+(0,-0.5)$)
%               -- ($(g)+(0,0.5)$)
%               -- ($(g)+(1,0.5)$) 
%               arc[start angle=0, end angle=-90, radius=1]
%               -- cycle;
%               \draw[fill=orange!20] ($(g)+(-0.5,0)$) rectangle ($(g)+(0.5,1)$);
%               \node () at ($(g)+(0.75,0.25)$) {\( \mathrm{C'} \)};
%               \node () at ($(g)+(-0.23,0.25)$) {\( \mathrm{L'} \)};
%               \draw[fill=white] ($(g)+(0,0)$) rectangle ($(g)+(0.5,0.5)$);
%               \node () at ($(g)+(0.25,0.25)$) {\( \mathrm{K} \)};
          
%               \coordinate (r) at (3,0);
%               \draw[fill=red!20] ($(r)+(-0.5,0)$)
%                 -- ($(r)+(-0.5,0.5)$)
%                 -- ($(r)+(0,1)$)
%                 --  ($(r)+(0.5,1)$)
%                 -- ($(r)+(0.5,0)$)
%                 -- cycle;
%               \node () at ($(r)+(-0.23,0.25)$) {\( \mathrm{R'} \)};
%               \draw[fill=white] ($(r)+(0,0)$) rectangle ($(r)+(0.5,0.5)$);
%               \node () at ($(r)+(0.25,0.25)$) {\( \mathrm{K} \)};
          
%               \coordinate (h) at (3, -2.2);
%               \draw[fill=blue!20]
%               ($(h)+(0,-0.5)$)
%               -- ($(h)+(0,0.5)$)
%               -- ($(h)+(1,0.5)$) 
%               arc[start angle=0, end angle=-90, radius=1]
%               -- cycle;
%               \draw[fill=red!20] ($(h)+(-0.5,0)$)
%               -- ($(h)+(-0.5,0.5)$)
%               -- ($(h)+(0,1)$)
%               --  ($(h)+(0.5,1)$)
%               -- ($(h)+(0.5,0)$)
%               -- cycle;
%              \node () at ($(h)+(0.75,0.25)$) {\( \mathrm{C'} \)};
%              \draw[fill=white] ($(h)+(0,0)$) rectangle ($(h)+(0.5,0.5)$);
%              \node () at ($(h)+(0.25,0.25)$) {\( \mathrm{K} \)};
%              \node () at ($(h)+(-0.23,0.25)$) {\( \mathrm{R'} \)};
          
%               \node[ font=\huge] (kl) at ($(k)!0.5!(l)+(0.25,0.25)$)
%              %    {\( \overset{l}{\leftarrowtail} \)}
%                 {\( \leftarrowtail \)}
%                ; 
%               \node[ font=\huge] (kr) at ($(k)!0.5!(r)+(0.25,0.25)$)
%                {\( \rightarrowtail \)}
%                ;  
%               \node[ font=\huge] (cg) at ($(c)!0.5!(g)+(0.25,0.25)$) 
%               {\( \leftarrowtail \)}
%              ;  
%               \node[ font=\huge] (ch) at ($(c)!0.5!(h)+(0.25,0.25)$)
%                {\( \rightarrowtail \)}
%              ; 
%               \node[ font=\huge] (kc) at ($(k)!0.5!(c)+(0.2,0.4)$) {\( \downarrowtail \)}; 
%             %   \node[ font=\LARGE] () at ($(l)!0.5!(g)+(0.5,0.4)$) {$m$}; 
%               \node[ font=\huge] (lg) at ($(l)!0.5!(g)+(0.1,0.4)$) {\( \downarrowtail \)}; 
%               \node[ font=\huge] (rh) at ($(r)!0.5!(h)+(0.1,0.4)$) {\( \downarrowtail \)}; 
%             %   \node[ font=\LARGE] () at ($(r)!0.5!(h)+(0.55,0.4)$) {$m'$}; 
%             %   \node[ font=\LARGE] () at ($(k)!0.5!(c)+(0.5,0.4)$) {$u$}; 
%             \end{tikzpicture}
%         }
%     \end{minipage}
\begin{minipage}{0.59\textwidth}
    % \setlength{\parindent}{1em}
    Let \( \rho = (L \overset{l}{\leftarrowtail} K \overset{r}{\rightarrowtail} R) \) be an injective DPO rewriting rule. 
    A rewriting step using $\rho$ with injective match \( m \) is defined by a DPO diagram as shown on the upper right.
     % \setlength{\parindent}{1em}
     It follows that both squares in the diagram depicted on the middle right, where all morphisms are inclusion functions and $H'$ is a graph isomorphic to $H$, are pushouts.
     By viewing a graph as a set of nodes and edges and graph monomorphisms as set injections, the DPO diagram on the middle right can be considered as a DPO diagram in the category of sets, illustrated on the lower right. In
\end{minipage}%
\hfill 
\begin{minipage}{0.39\textwidth}
    % \hfill 
    %     \resizebox{0.7\textwidth}{!}{
    %         \begin{tikzpicture}
    %             % [node distance=15mm]
    %             \node (I) {$K$};
    %             \node (L) [left of=I] {$L$};
    %             \node (R) [right of=I] {$R$};
    %             \node (G) [below of=L] {$G$};
    %             \node (C) [below of=I] {$C$};
    %             \node (H) [below of=R] {$H$};
    %             \draw [>->] (I) to  node [midway,below] {$l$} (L);
    %             \draw [>->] (I) to  node [midway,below] {$r$} (R);
    %             \draw [>->] (L) to node [midway,right] {$m$} (G);
    %             \draw [>->] (I) to  node [midway,right] 
    %             % {$u$}
    %             {} (C);
    %             \draw [>->] (R) to  node [midway,right] 
    %             % {$m'$}
    %             {} (H);
    %             \draw [>->] (C) to node [midway,above] {$l'$} (G);
    %             \draw [>->] (C) to node [midway,above] 
    %             % {$r'$} 
    %             {}
    %             (H);
    %             % \node [at=($(I)!.5!(G)$)] {\normalfont PO};
    %             % \node [at=($(I)!.5!(H)$)] {\normalfont PO};
    %         \end{tikzpicture}
    %     }
    \centering
    \resizebox{0.7\textwidth}{!}{
        \begin{tikzpicture}
            % [node distance=15mm]
            \node (I) at (0,0) {$K$};
            \node (L)  at (-2,0) {$L$};
            \node (R)  at (2,0) {$R$};
            \node (G)  at (-2,-2) {$G$};
            \node (C)  at (0,-2) {$C$};
            \node (H)  at (2,-2) {$H$};
            \draw [>->] (I) to  node [midway,below] {$l$} (L);
            \draw [>->] (I) to  node [midway,below] {$r$} (R);
            \draw [>->] (L) to node [midway,right] {$m$} (G);
            \draw [>->] (I) to  node [midway,right] 
            % {$u$}
            {} (C);
            \draw [>->] (R) to  node [midway,right] 
            {$m'$}
            (H);
            \draw [>->] (C) to node [midway,above] {$l'$} (G);
            \draw [>->] (C) to node [midway,above] 
            % {$r'$} 
            {}
            (H);
            % \node [at=($(I)!.5!(G)$)] {\normalfont PO};
            % \node [at=($(I)!.5!(H)$)] {\normalfont PO};
        \end{tikzpicture}
    }

        \resizebox{0.7\textwidth}{!}{
            \begin{tikzpicture}
                [node distance=20mm]
                \node (I) at (0,0) {$\operatorname{Im}(l \star m)$};
                \node (L) at (-2,0) {$\operatorname{Im}(m)$};
                \node (R) at (2,0) {$\operatorname{Im}(m')$};
                \node (G) at (-2,-2) {$G$};
                \node (C) at (0,-2) {$\operatorname{Im}(l')$};
                \node (H) at (2,-2) {$H'$};
                \draw [>->] (I) to  node [midway,above] {} (L);
                \draw [>->] (I) to  node [midway,above] {} (R);
                \draw [>->] (L) to node [midway,right] {} (G);
                \draw [>->] (I) to  node [midway,right] {} (C);
                \draw [>->] (R) to  node [midway,right] {} (H);
                \draw [>->] (C) to node [midway,above] {} (G);
                \draw [>->] (C) to node [midway,above] {} (H);
                % \node [at=($(I)!.5!(G)$)] {\normalfont PO};
                % \node [at=($(I)!.5!(H)$)] {\normalfont PO};
            \end{tikzpicture}
        }

          \resizebox{0.7\textwidth}{!}{
            \begin{tikzpicture}
                \coordinate (k) at (0, 0);
                \draw[fill=white] ($(k)+(0,0)$) rectangle ($(k)+(0.5,0.5)$);
                \node () at ($(k)+(0.25,0.25)$) {\( \mathrm{K} \)};
            
                \coordinate (c) at (0, -2.2);
                \draw[fill=blue!20]
                ($(c)+(0,-0.5)$)
                -- ($(c)+(0,0.5)$) 
                -- ($(c)+(1,0.5)$) 
                arc[start angle=0, end angle=-90, radius=1]
                -- cycle;
                \node () at ($(c)+(0.75,0.25)$) {\( \mathrm{C'} \)};
                \draw[fill=white] ($(c)+(0,0)$) rectangle ($(c)+(0.5,0.5)$);
                \node () at ($(c)+(0.25,0.25)$) {\( \mathrm{K} \)};
            
                \coordinate (l) at (-3, 0);
                \draw[fill=orange!20] ($(l)+(-0.5,0)$) rectangle ($(l)+(0.5,1)$);
                \node () at ($(l)+(-0.23,0.25)$) {\( \mathrm{L'} \)};
                \draw[fill=white] ($(l)+(0,0)$) rectangle ($(l)+(0.5,0.5)$);
                \node () at ($(l)+(0.25,0.25)$) {\( \mathrm{K} \)};
            
                \coordinate (g) at (-3, -2.2);
                \draw[fill=blue!20]
                ($(g)+(0,-0.5)$)
                -- ($(g)+(0,0.5)$)
                -- ($(g)+(1,0.5)$) 
                arc[start angle=0, end angle=-90, radius=1]
                -- cycle;
                \draw[fill=orange!20] ($(g)+(-0.5,0)$) rectangle ($(g)+(0.5,1)$);
                \node () at ($(g)+(0.75,0.25)$) {\( \mathrm{C'} \)};
                \node () at ($(g)+(-0.23,0.25)$) {\( \mathrm{L'} \)};
                \draw[fill=white] ($(g)+(0,0)$) rectangle ($(g)+(0.5,0.5)$);
                \node () at ($(g)+(0.25,0.25)$) {\( \mathrm{K} \)};
            
                \coordinate (r) at (3,0);
                \draw[fill=red!20] ($(r)+(-0.5,0)$)
                -- ($(r)+(-0.5,0.5)$)
                -- ($(r)+(0,1)$)
                --  ($(r)+(0.5,1)$)
                -- ($(r)+(0.5,0)$)
                -- cycle;
                \node () at ($(r)+(-0.23,0.25)$) {\( \mathrm{R'} \)};
                \draw[fill=white] ($(r)+(0,0)$) rectangle ($(r)+(0.5,0.5)$);
                \node () at ($(r)+(0.25,0.25)$) {\( \mathrm{K} \)};
            
                \coordinate (h) at (3, -2.2);
                \draw[fill=blue!20]
                ($(h)+(0,-0.5)$)
                -- ($(h)+(0,0.5)$)
                -- ($(h)+(1,0.5)$) 
                arc[start angle=0, end angle=-90, radius=1]
                -- cycle;
                \draw[fill=red!20] ($(h)+(-0.5,0)$)
                -- ($(h)+(-0.5,0.5)$)
                -- ($(h)+(0,1)$)
                --  ($(h)+(0.5,1)$)
                -- ($(h)+(0.5,0)$)
                -- cycle;
            \node () at ($(h)+(0.75,0.25)$) {\( \mathrm{C'} \)};
            \draw[fill=white] ($(h)+(0,0)$) rectangle ($(h)+(0.5,0.5)$);
            \node () at ($(h)+(0.25,0.25)$) {\( \mathrm{K} \)};
            \node () at ($(h)+(-0.23,0.25)$) {\( \mathrm{R'} \)};
            
                \node[ font=\huge] (kl) at ($(k)!0.5!(l)+(0.25,0.25)$)
            %    {\( \overset{l}{\leftarrowtail} \)}
                {\( \leftarrowtail \)}
                ; 
                \node[ font=\huge] (kr) at ($(k)!0.5!(r)+(0.25,0.25)$)
                {\( \rightarrowtail \)}
                ;  
                \node[ font=\huge] (cg) at ($(c)!0.5!(g)+(0.25,0.25)$) 
                {\( \leftarrowtail \)}
            ;  
                \node[ font=\huge] (ch) at ($(c)!0.5!(h)+(0.25,0.25)$)
                {\( \rightarrowtail \)}
            ; 
                \node[ font=\huge] (kc) at ($(k)!0.5!(c)+(0.2,0.4)$) {\( \downarrowtail \)}; 
            %   \node[ font=\LARGE] () at ($(l)!0.5!(g)+(0.5,0.4)$) {$m$}; 
                \node[ font=\huge] (lg) at ($(l)!0.5!(g)+(0.1,0.4)$) {\( \downarrowtail \)}; 
                \node[ font=\huge] (rh) at ($(r)!0.5!(h)+(0.1,0.4)$) {\( \downarrowtail \)}; 
            %   \node[ font=\LARGE] () at ($(r)!0.5!(h)+(0.55,0.4)$) {$m'$}; 
            %   \node[ font=\LARGE] () at ($(k)!0.5!(c)+(0.5,0.4)$) {$u$}; 
            \end{tikzpicture}
          }
\end{minipage}
the lower right diagram, $L', C', R', K$ are mutually disjoint sets such that $K$, 
         $K \uplus L'$, 
         $K \uplus R'$, 
         $K \uplus C'$, 
         $L' \uplus K \uplus C'$ and 
         $R' \uplus K \uplus C'$ 
            are corresponding sets of 
         $\operatorname{Im}(l \star m)$, 
         $\operatorname{Im}(m)$, 
         $\operatorname{Im}(m')$,
         $\operatorname{Im}(l')$,
         $G$ and
         $H'$ respectively,
         and all morphisms as inclusion functions.
    % Up to isomorphism, we can further consider all morphisms as inclusion functions, illustrated on the right. In this visualization, $L',R',C',K$ are mutually disjoint sets with $L' \uplus K = L$, $R' \uplus K= R$, $C' \uplus K = C$, $G = L \cup C$ and $H = R \cup C$.
    Thus, any subgraph of $G$ (resp. $H$) can be considered as a union of subgraphs of $C$ and $L$ (resp. $R$), up to isomorphism.
    % \color{red}\begin{definition}
%         Consider a rewriting step defined by the DPO diagram above. Let $X$ be a ruler-graph and $G$ a graph. An \textbf{$X$-occurrence} in $G$ is a subgraph of $G$ isomorphic to $X$. An \textbf{$X$-occurrence $x$ implicitly created by the rewriting step} is an \( X \)-occurrence in $H$ which is included neither in \( R \) nor in $C$. 
%         Let $R' = x \cap R$, $C' = x \cap C$ and $h_{R'L}:R' \rightarrowtail L$ a monomorphism preserving interface elements. The \textbf{corresponding $X$-occurrence relative to $h_{R'L}$} of $x$ is the $X$-occurrence in $G$ defined by $h_{R'L}(R') \cup C'$.
%         An \textbf{$X$-occurrence $x$ implicitly destroyed by the rewriting step} is an \( X \)-occurrence in $G$ which is included neither in \( L \) nor in $C$. 
%     \end{definition}
% \color{black}

With the above DPO diagrams in mind, we define:

\begin{definition}
    \label{def:corresponding_occurrence}
     An \textbf{$X$-occurrence} in $G$ is a subgraph of $G$ isomorphic to $X$. An \textbf{$X$-occurrence $x$ implicitly created by the rewriting step} is an \( X \)-occurrence in $H$ which is included neither in \( R \) nor in $C$. 
      The occurrence $x$ can be decomposed as the union of graphs $x = R_x \cup C_x$ where $R_x = x \cap R$ and $C_x = x \cap C$. For every graph monomorphism $h : R_x \rightarrowtail L$ that preserves interface elements (i.e., \( h(k) = k \) for all \( k \in K \)), we define the \textbf{corresponding $X$-occurrence of $x$ relative to $h$} as the $X$-occurrence in $G$ given by $h(R_x) \cup C_x$.
    An \textbf{$X$-occurrence implicitly destroyed by the rewriting step} is an \( X \)-occurrence in $G$ which is included neither in \( L \) nor in $C$.
\end{definition}
We define the set \(D(R,X)\) of all subgraphs of \( R \) which, when glued along some common interface elements with a subgraph $C'$ of a context graph \( C \), can form an \( X \)-occurrence implicitly created by the rewriting step.

\begin{definition}
    \label{def:rx}
    % Let \(X\) be a graph and 
    % \(
    %     L \overset{l}{\leftarrowtail} K \overset{r}{\rightarrowtail} R
    % \) an injective DPO rewriting rule.
    % %  such that $l$ is edge-injective and $r$ is node-injective.
    % The set \(D(R,X)\) consists of all subgraphs \( R' \subseteq R \) such that there exists a diagram (shown on the right) where
    % \(h_{R'X} \colon R' \rightarrowtail X \) and \(h_{C'X} \colon C' \rightarrowtail X \) are not isomorphisms, and \(h_{K'K} \colon K' \rightarrowtail K \) and \(h_{R'R} \colon R' \rightarrowtail R \) are inclusion functions. The union of all subgraphs in \( D(R,X) \) is denoted \( R_X \).
    % Let \(X\) be a graph and 
    % \(
    %     L \overset{l}{\leftarrowtail} K \overset{r}{\rightarrowtail} R
    % \) an injective DPO rewriting rule.
    % %  such that $l$ is edge-injective and $r$ is node-injective.
    % The set \(D(R,X)\) consists of all subgraphs \( R' \subseteq R \) such that 
    % we can construct a diagram as shown on the right where \(h_{K'K} \colon K' \rightarrowtail K \) and \(h_{R'R} \colon R' \rightarrowtail R \) are inclusion functions, and \(h_{R'X} \colon R' \rightarrowtail X \) and \(h_{C'X} \colon C' \rightarrowtail X \) are not isomorphisms. The union of all subgraphs in \( D(R,X) \) is denoted \( R_X \).
    Let \(X\) be a graph and 
    \(
        L \overset{l}{\leftarrowtail} K \overset{r}{\rightarrowtail} R
    \) an injective DPO rewriting rule.
    %  such that $l$ is edge-injective and $r$ is node-injective.
    The set \(D(R,X)\) consists of all subgraphs \( R' \subseteq R \)
    satisfying the following conditions:
    \newline
    \noindent
    \begin{minipage}{0.69\textwidth}\setlength{\parindent}{1em}
    \begin{itemize}
        \item a diagram as shown on the right where \(h_{K'K} \colon K' \rightarrowtail K \) and \(h_{R'R} \colon R' \rightarrowtail R \) are inclusion functions can be constructed,
        \item $R'$ is not a subgraph of $r(K)$,
        \item $R'$ is not isomorphic to $X$.
    \end{itemize}
    \end{minipage}
    \begin{minipage}{0.29\textwidth}
        \hfill
        \resizebox{0.9\textwidth}{!}{
        \begin{tikzpicture}[node distance=11mm]
            \node (k) {K};
            \node (r) [above=of k] {R};
            \node (r') [right=of r] {$R'$};
            \node (k') [right=of k] {K'};
            \node (x)  [right=of r'] {X};
            \node (c') [below=of x] {C'};
            \node () [at=($(r)!0.5!(k')$)] {PB};
            \node () [at=($(x)!0.5!(k')$)] {PO};
            \draw[>->] (k) -- (r) node[pos= .5, left] {$r$}; 
            \draw[>->] (r') -- (r) node[pos= .5, above] {$h_{R'R}$};
            \draw[>->] (k') -- (k) node[pos= .5, below] {$h_{K'K}$};
            \draw[>->] (k') -- (r');
            \draw[>->] (k') -- (c');
            \draw[>->] (c') -- (x);
            \draw[>->] (r') -- (x);
        \end{tikzpicture}
        }
    \end{minipage}
    % The union of all subgraphs in \( D(R,X) \) is denoted \( R_X \).
\end{definition}  
The first condition of the definition ensures that any $R' \in D(R,X)$ can be glued along some common interface elements with a subgraph $C'$ of a context graph \( C \) to form an \( X \)-occurrence. The second and third conditions ensure that such $X$-occurrences are included neither in \( C \) nor in \( R \), thereby guaranteeing that they are implicitly created as defined in \autoref{def:corresponding_occurrence}.

% In another word, the set $D(R,X)$ contains all subgraphs $R''$ of $R$ which is not a subgraph of $K$ such that there exists a rewriting step defined by the diagram above such that $R'' \cup C''$ for a subgraph $C''$ of $C$ is an $X$-occurrence created by the rewriting step. 
\begin{notation}
    We use the notation from~\cite[Notation 1]{overbeek2023apbpotutorial} to visualize edge-labeled graph homomorphisms. Labeled graphs are enclosed in boxes with their names displayed in the top-left corner. Nodes and edges are assigned subsets of \(\mathbb{N}\) as identifiers, and these identifiers are chosen such that: (i) Each node or edge \( y \) in the codomain graph is assigned the union of the identifiers of all nodes or edges in the domain graph that are mapped to \( y \); (ii) The graph homomorphism is uniquely determined by this assignment.
    
    \noindent To further improve readability, we represent sets by listing their elements. Additionally, we omit identifiers when doing so does not cause confusion. This is illustrated in the following representation of a homomorphism \( h: G \to H \).
    
    \begin{center}
        \resizebox{0.45\textwidth}{!}{
        \begin{tikzpicture}
            \graphbox{\( G \)}{00mm}{-20mm}{45mm}{20mm}{2mm}{-5mm}{
                \coordinate (o) at (-5mm,-8mm); 
                \node[draw,circle] (l1) at ($(o)+(-10mm,0mm)$) {1};
                \node[draw,circle] (l2) at ($(l1)+(3,0)$) {2};
                \node[draw,circle] (l3) at ($(l1)+(1,0)$) {3};
                \node[draw,circle] (l4) at ($(l1)+(2,0)$) {4};
                \draw[->] (l1) -- (l3) node[midway,above] {a};
                \draw[->] (l3) -- (l4) node[midway,above] {b};
                \draw[->] (l4) -- (l2) node[midway,above] {a};
            }  
            \graphbox{\( H \)}{50mm}{-20mm}{34mm}{20mm}{2mm}{-5mm}{
                \coordinate (o) at (0mm,-8mm); 
                \node[draw,circle] (l1) at ($(o)+(-10mm,0mm)$) {1};
                \node[draw,circle] (l2) at ($(l1)+(2,0)$) {2};
                \node[draw,circle] (l3) at ($(l1)+(1,0)$) {3\ 4};
                \draw[->] (l1) -- (l3) node[midway,above] {a};
                \draw[->] (l3) edge[loop above] (l3) node[midway,above] {b};
                \draw[->] (l3) -- (l2) node[midway,above] {a};
            }      
            % \node () at (53mm,-30mm) {$\rightarrow$};
        \end{tikzpicture}
    }
    \end{center} 
    In this example, the sets \(\{1\}\), \(\{2\}\), \(\{3\}\), \(\{4\}\), and \(\{3,4\}\) are represented as \(1\), \(2\), \(3\), \(4\), and \(3\ 4\), respectively. Edge identifiers are omitted.
\end{notation} 

\begin{example}
    \label{ex:grsaa_rx}
    Consider the injective DPO rule from \cite[Example 6]{bruggink2014termination}:  
    \begin{center} 
        \resizebox{0.7\textwidth}{!}{
        \begin{tikzpicture}
            \graphbox{$L$}{0mm}{0mm}{34mm}{15mm}{2mm}{-5mm}{
                \coordinate (o) at (0mm,-3mm); 
                \node[draw,circle] (l1) at ($(o)+(-10mm,0mm)$) {1};
                \node[draw,circle] (l2) at ($(l1)+(2,0)$) {2};
                \node[draw,circle] (l3) at ($(l1) + (1,0)$) {3};
                \draw[->] (l1) -- (l3) node[midway,above] {a};
                \draw[->] (l3) -- (l2) node[midway,above] {a};
            }     
            \graphbox{$K$}{40mm}{0mm}{24mm}{15mm}{2mm}{-5mm}{
                \coordinate (o) at (5mm,-3mm); 
                \node[draw,circle] (l1) at ($(o)+(-10mm,0mm)$) {1};
                \node[draw,circle] (l2) at ($(l1)+(1,0)$) {2};
                % \node[draw,circle] (l3) at ($(l1) + (1,0)$) {$\ $};
                % \draw[->] (l1) -- (l3) node[midway,above] {a};
                % \draw[->] (l3) -- (l2) node[midway,above] {a};
            }    
            \graphbox{$R$}{70mm}{0mm}{45mm}{15mm}{2mm}{-5mm}{
                \coordinate (o) at (-5mm,-3mm); 
                \node[draw,circle] (l1) at ($(o)+(-10mm,0mm)$) {1};
                \node[draw,circle] (l2) at ($(l1)+(3,0)$) {2};
                \node[draw,circle] (l3) at ($(l1) + (1,0)$) {4};
                \node[draw,circle] (l4) at ($(l1) + (2,0)$) {5};
                \draw[->] (l1) -- (l3) node[midway,above] {a};
                \draw[->] (l3) -- (l4) node[midway,above] {b};
                \draw[->] (l4) -- (l2) node[midway,above] {a};
            }    

            \node () at (37mm,-8mm) {$\leftarrowtail$};
            \node () at (67mm,-8mm) {$\rightarrowtail$};

            % \draw[>->] (51mm,2mm) -- (52mm,3mm);
        \end{tikzpicture}
        }
    \end{center}
    Let $X$ be \tikz[baseline=-0.5ex]{
        \node (x) at (0,0) {$\bullet$};
        \node (y) at (1,0) {$\bullet$ };
        \node (z) at (2,0) { $\bullet$};
        \draw[->] (x) -- (y) node[midway, above] {$a$};
        \draw[->] (y) -- (z) node[midway, above] {$a$};
    } and $R' \in D(R,X)$. 
    $R'$ is neither the empty graph nor isomorphic to $X$, by Conditions 2 and 3 of \autoref{def:rx}. 
    Therefore, $R'$ must contain an interface node (either node 1 or node 2), because if it contains no interface node, Condition 1 of \autoref{def:rx} would be violated due to $X$'s connectivity.
    Additionally, $R'$ must contain node 4 or node 5 to avoid being a subgraph of $r(K)$ (Condition 2). 
    Therefore, $R'$ must contain either $1 \overset{a}{\to} 4$ or $5 \overset{a}{\to} 2$, because if node 4 or node 5 is isolated, Condition 1 would be violated due to $X$'s connectivity.
    Furthermore, $R'$ must not contain $4 \overset{b}{\to} 5$ by Condition 1.
    The following graph is not in $D(R,X)$ because it violates Condition 1.
    \begin{center}
        \resizebox{0.24\textwidth}{!}{
            \begin{tikzpicture}
                \graphbox{}{70mm}{0mm}{45mm}{15mm}{2mm}{-5mm}{
                    \coordinate (o) at (-5mm,-3mm); 
                    \node[draw,circle] (l1) at ($(o)+(-10mm,0mm)$) {1};
                    \node[draw,circle] (l3) at ($(l1) + (1,0)$) {4};
                    \node[draw,circle] (l2) at ($(l1)+(3,0)$) {2};
                    \node[draw,circle] (l4) at ($(l1) + (2,0)$) {5};
                    \draw[->] (l4) -- (l2) node[midway,above] {a};
                    \draw[->] (l1) -- (l3) node[midway,above] {a};
                }    
            \end{tikzpicture}
        } 
    \end{center}
    Thus, $D(R,X)$ consists of the four graphs $R'_1$, $R'_2$, $R'_3$ and $R'_4$ shown below.
    For each graph, the construction of the pullback square in Condition 1 is demonstrated in \autoref{example:grs_aa:has_more_left} and the construction of the pushout square in Condition 1 is straightforward. 
    \begin{center}
        \resizebox{0.24\textwidth}{!}{
            \begin{tikzpicture}
                \graphbox{$R'_1$}{70mm}{0mm}{45mm}{15mm}{2mm}{-5mm}{
                    \coordinate (o) at (-5mm,-3mm); 
                    \node[circle] (l1) at ($(o)+(-10mm,0mm)$) {};
                    \node[draw,circle] (l2) at ($(l1)+(3,0)$) {2};
                    \node[draw,circle] (l4) at ($(l1) + (2,0)$) {5};
                    \draw[->] (l4) -- (l2) node[midway,above] {a};
                }    
            \end{tikzpicture}
        }
        \resizebox{0.24\textwidth}{!}{
            \begin{tikzpicture}
                \graphbox{$R'_2$}{70mm}{0mm}{45mm}{15mm}{2mm}{-5mm}{
                    \coordinate (o) at (-5mm,-3mm); 
                    \node[draw,circle] (l1) at ($(o)+(-10mm,0mm)$) {1};
                    \node[draw,circle] (l2) at ($(l1)+(3,0)$) {2};
                    \node[draw,circle] (l4) at ($(l1) + (2,0)$) {5};
                    \draw[->] (l4) -- (l2) node[midway,above] {a};
                }    
            \end{tikzpicture}
        } 
        \resizebox{0.24\textwidth}{!}{
            \begin{tikzpicture}
                \graphbox{$R'_3$}{70mm}{0mm}{45mm}{15mm}{2mm}{-5mm}{
                    \coordinate (o) at (-5mm,-3mm); 
                    \node[draw,circle] (l1) at ($(o)+(-10mm,0mm)$) {1};
                    \node[draw,circle] (l3) at ($(l1) + (1,0)$) {4};
                    \draw[->] (l1) -- (l3) node[midway,above] {a};
                }   
        \end{tikzpicture}
        } 
        \resizebox{0.24\textwidth}{!}{
            \begin{tikzpicture}
                \graphbox{$R'_4$}{70mm}{0mm}{45mm}{15mm}{2mm}{-5mm}{
                    \coordinate (o) at (-5mm,-3mm); 
                    \node[draw,circle] (l1) at ($(o)+(-10mm,0mm)$) {1};
                    \node[draw,circle] (l2) at ($(l1)+(3,0)$) {2};
                    \node[draw,circle] (l3) at ($(l1) + (1,0)$) {4};
                    \draw[->] (l1) -- (l3) node[midway,above] {a};
                }    
            \end{tikzpicture}
        }
    \end{center}
\end{example} 
 
\begin{example}
    \label{ex:rx_counter_ex}
    Let $X$ be the graph 
    \tikz[baseline=-0.5ex]{ 
            \node (x) at (0,0) {$\bullet$}; 
            \node (y) at (1,0) {$\bullet$};
            \node (z) at (2,0) {$\bullet$};
            \draw[->] (x) -- (y)   {};
            \draw[->] (y) -- (z)   {};
    }. Consider the following rule.
    \begin{center}
        \resizebox{0.6\textwidth}{!}{
            \begin{tikzpicture}
                \graphbox{$L$}{0mm}{0mm}{34mm}{15mm}{2mm}{-5mm}{
                    \coordinate (o) at (0mm,-3mm); 
                    \node[draw,circle] (l1) at ($(o)+(-10mm,0mm)$) {1};
                    % \node[draw,circle] (l2) at ($(l1)+(2,0)$) {2};
                    \node[draw,circle] (l3) at ($(l1) + (1,0)$) {2};
                    \draw[->] (l3) -- (l1) node[midway,above] {};
                    % \draw[->] (l3) -- (l2) node[midway,above] {a};
                }     
                \graphbox{$K$}{40mm}{0mm}{24mm}{15mm}{2mm}{-5mm}{
                    \coordinate (o) at (5mm,-3mm); 
                    \node[draw,circle] (l1) at ($(o)+(-10mm,0mm)$) {1};
                    \node[draw,circle] (l2) at ($(l1)+(1,0)$) {2};
                    % \node[draw,circle] (l3) at ($(l1) + (1,0)$) {$\ $};
                    \draw[->] (l2) -- (l1) node[midway,above] {};
                    % \draw[->] (l3) -- (l2) node[midway,above] {a};
                }    
                \graphbox{$R$}{70mm}{0mm}{40mm}{15mm}{2mm}{-5mm}{
                    \coordinate (o) at (-5mm,-3mm); 
                    \node[draw,circle] (l1) at ($(o)+(-10mm,0mm)$) {1};
                    % \node[draw,circle] (l2) at ($(l1)+(3,0)$) {2};
                    \node[draw,circle] (l3) at ($(l1) + (1,0)$) {2};
                    \node[draw,circle] (l4) at ($(l1) + (2,0)$) {3};
                    \draw[->] (l3) -- (l1) node[midway,above] {};
                    \draw[->] (l3) -- (l4) node[midway,above] {};
                    % \draw[->] (l4) -- (l2) node[midway,above] {a};
                }    
                \node () at (37mm,-8mm) {$\leftarrowtail$};
                \node () at (67mm,-8mm) {$\rightarrowtail$};
            \end{tikzpicture}
        }
    \end{center}
    Let $R' \in D(R,X)$. $R'$ must contain node 3, otherwise it would be a subgraph of $r(K)$, violating Condition 2 of \autoref{def:rx}. Node 3 must be connected to an interface node, otherwise Condition 1 of \autoref{def:rx} would be violated. Therefore, $R'$ must include \raisebox{2pt}{\scalebox{0.6}{\tikz[baseline=-0.5ex]{
        \node [draw,circle] (x) at (0,0) {2};
        \node[draw,circle] (y) at (1,0) {3};
        \draw[->] (x) -- (y) {};
    }}}.
    The graph $R$ is not in $D(R,X)$ because it violates Condition 1. Thus, the set \( D(R,X) \) consists of $R'_1$:
    \raisebox{2pt}{\scalebox{0.6}{\tikz[baseline=-0.5ex]{
        \node [draw,circle] (x) at (0,0) {2};
        \node[draw,circle] (y) at (1,0) {3};
        \draw[->] (x) -- (y) {};
    }}} and $R'_2$:
    \raisebox{2pt}{\scalebox{0.6}{\tikz[baseline=-0.5ex]{
        \node [draw,circle] (node1) at (-1,0) {1};
        \node [draw,circle] (x) at (0,0) {2};
        \node[draw,circle] (y) at (1,0) {3};
        \draw[->] (x) -- (y) {};
    }}}.
\end{example}
We define the notion of \( X \)-non-increasing rules. Intuitively, if a rule $\rho$ is $X$-non-increasing, then for any rewriting step using $\rho$, there exists an injective mapping from the set of \( X \)-occurrences implicitly created by the step to the set of \( X \)-occurrences implicitly destroyed by the step. This will be proved in \autoref{lem:w_u_l_not_geq_r_not}.
\begin{definition}[Non-increasing rule]
    \label{def:creates_more_x_on_the_left}
    Let \(\rho = (L \overset{l}{\leftarrowtail} K \overset{r}{\rightarrowtail} R)\) be a rule and \(X\) a ruler-graph. 
    Rule \( \rho \) is \emph{\( X \)-non-increasing} if the following four conditions hold:
    \newline
    \vspace{1mm}
    \noindent
    \begin{minipage}{0.64\textwidth} 
    \begin{enumerate}
        \item For all $R' \in D(R,X)$, there exists a morphism \(h_{R'L}: R' \rightarrowtail L \) such that both squares in the diagram shown on the right are pullbacks, where all morphisms other than $h_{R'L}$, $l$ and $r$ are inclusion functions, 
        \item \label{def:non_increasing:non_clapse} For all $R' \in D(R,X)$, for all node or edge $r'$ in $R'$, if $ r' \notin \operatorname{Im}(r)$ then $h_{R'L}(r') \notin \operatorname{Im}(l)$,
        \item For all $R',R'' \in D(R,X)$, for all edges $x \in R'$ and $y \in R''$, if $x \neq y$, then $h_{R'L}(x) \neq h_{R''L}(y)$,
        \item If $X$ has isolated nodes, then for all $R',R'' \in D(R,X)$ and nodes $x \in R', y \in R''$, if $x \neq y$, then $h_{R'L}(x) \neq h_{R''L}(y)$.
    \end{enumerate} 
    \end{minipage}
    \hfill
    \begin{minipage}{0.4\textwidth}
        % \hfill
        \begin{center}
            % \resizebox{0.9\textwidth}{!}{ 
                % \begin{tikzpicture}[rotate=90]
                %     \node (k) {K}; 
                %     \node (k') [above=of k] {$K'$};
                %     \node (l) [left=of k] {$L$};
                %     \node (rb) [above=of l] {$R'$};
                %     \node (r) [right=of k] {$R$};
                %     \node (rb') [above=of r] {$R'$};
                %     \draw[<-<]  (l) -- (k) node [midway,below] {l};
                %     \draw[<-<]  (r) -- (k) node [midway,below] {r};
                %     \draw[>->]  (k') -- (rb');
                %     \draw[>->]  (rb') -- (r);
                %     \draw[>->]  (k') -- (rb);
                %     \draw[>->]  (k') -- (k);
                %     \draw[>->]  (rb) -- (l) node [midway,above,sloped] {$h_{R'L}$};
                %     \node () [at=($(k)!0.5!(rb)$)] {PB};
                %     \node () [at=($(k)!0.5!(rb')$)] {PB};
                % \end{tikzpicture}
                \begin{tikzpicture} 
                    \node (k) {K}; 
                    \node (k') [left=of k] {$K'$};
                    \node (l) [above=of k] {$L$};
                    \node (rb) [left=of l] {$R'$};
                    \node (r) [below=of k] {$R$};
                    \node (rb') [left=of r] {$R'$};
                    \draw[<-<]  (l) -- (k) node [midway,right] {l};
                    \draw[<-<]  (r) -- (k) node [midway,right] {r};
                    \draw[>->]  (k') -- (rb');
                    \draw[>->]  (rb') -- (r);
                    \draw[>->]  (k') -- (rb);
                    \draw[>->]  (k') -- (k);
                    \draw[>->]  (rb) -- (l) node [midway,above] {$h_{R'L}$};
                    \node () [at=($(k)!0.5!(rb)$)] {PB};
                    \node () [at=($(k)!0.5!(rb')$)] {PB};
                \end{tikzpicture}
            % }
        \end{center}
    \end{minipage}

    %  \begin{enumerate}
    %     % \item[(2)] \label{def:non_increasing:non_clapse} For all node or arrow $r'$ in $R'$, if $ r' \notin \operatorname{Im}(r)$ then $h_{R'L}(r') \notin \operatorname{Im}(l)$,
    %     % \item there exists a morphism $h_{R_XL}$ such that
    %     %     \( h_{R'R_X} \star h_{R_XL} = h_{R'L} \) for all $R' \in D(R,X)$;
    %     % \item[(3)] \label{def:non_increasing:edge_injective} For all $R',R'' \in D(R,X)$, for all edges $x \in R'$ and $y \in R''$, if $x \neq y$, then $h_{R'L}(x) \neq h_{R''L}(y)$,
    %     %  \item $h$  is edge-injective 
    %     % \item $\bigcup_{R' \in D(R,X)} h_{R'L}$ is an edge-injective morphism from $R_X$ to $L$;
    %     % \item[(4)] \label{def:non_increasing:isolated_nodes} If $X$ has isolated nodes, then for all $R',R'' \in D(R,X)$ and nodes $x \in R', y \in R''$, if $x \neq y$, then $h_{R'L}(x) \neq h_{R''L}(y)$.
    %     % \item $h$ is node-injective if $X$ has isolated nodes;
    %  \end{enumerate}
    %  \todo{est il util d'avoir "or arrow?" c'est une partie de edge injective, non ?}
\end{definition}
The first condition ensures that whenever $R' \in D(R,X)$ forms an $X$-occurrence in the result graph $H$, this occurrence has a corresponding $X$-occurrence relative to the morphism $h_{R'L}$ (as defined in \autoref{def:corresponding_occurrence}) in the host graph $G$.
The second condition guarantees that any $X$-occurrence implicitly created by the rewriting step has its corresponding $X$-occurrence not included in the context (Consequently, its corresponding $X$-occurrence is implicitly destroyed by the rewriting step). Counterexamples are provided in \autoref{ex:cond_2_necessaire} and \autoref{ex:cond_2_necessaire2}.
The third and fourth conditions ensure that distinct $X$-occurrences created implicitly have distinct corresponding $X$-occurrences.
Counterexamples for these conditions are given in \autoref{ex:cond3_necessaire} and \autoref{ex:cond4_necessaire}.

The following example shows an $X$-non-increasing rule. 
\begin{example} 
    \label{example:grs_aa:has_more_left}
    Let $X$ be \tikz[baseline=-0.5ex]{
        \node (x) at (0,0) {$\bullet$};
        \node (y) at (1,0) {$\bullet$};
        \node (z) at (2,0) {$\bullet$};
        \draw[->] (x) -- (y) node[midway, above] {$a$};
        \draw[<-] (z) -- (y) node[midway, above] {$a$};
    }. Consider the rule presented
    in \autoref{ex:grsaa_rx}. By \autoref{ex:grsaa_rx}, the set \( D(R,X) \) consists of exactly four graphs: $R'_1$, $R'_2$, $R'_3$ and $R'_4$.
    The rule is $X$-non-increasing because the following diagrams can be constructed, and $h_{R'_1L}, h_{R'_2L}, h_{R'_3L}$ and $h_{R'_4L}$ satisfy all conditions in \autoref{def:creates_more_x_on_the_left}. 
%1
\begin{center}
    \resizebox{0.7\textwidth}{!}{
        \begin{tikzpicture}
            \graphbox{$R'_1$}{0mm}{0mm}{34mm}{15mm}{2mm}{-5mm}{
                \coordinate (o) at (0mm,-3mm); 
                \node (l1) at ($(o)+(-10mm,0mm)$) {};
                \node[draw,circle] (l2) at ($(l1)+(2,0)$) {2};
                \node[draw,circle] (l3) at ($(l1) + (1,0)$) {5};
                % \draw[->] (l1) -- (l3) node[midway,above] {a};
                \draw[->] (l3) -- (l2) node[midway,above] {a};
            }     
            \graphbox{$K'$}{40mm}{0mm}{24mm}{15mm}{2mm}{-5mm}{
                \coordinate (o) at (5mm,-3mm); 
                \node (l1) at ($(o)+(-10mm,0mm)$) {};
                \node[draw,circle] (l2) at ($(l1)+(1,0)$) {2};
                % \node[draw,circle] (l3) at ($(l1) + (1,0)$) {$\ $};
                % \draw[->] (l1) -- (l3) node[midway,above] {a};
                % \draw[->] (l3) -- (l2) node[midway,above] {a};
            }    
            \graphbox{$R'_1$}{70mm}{0mm}{45mm}{15mm}{2mm}{-5mm}{
                \coordinate (o) at (-5mm,-3mm); 
                \node (l1) at ($(o)+(-10mm,0mm)$) {};
                \node[draw,circle] (l2) at ($(l1)+(3,0)$) {2};
                % \node[draw,circle] (l3) at ($(l1) + (1,0)$) {4};
                \node[draw,circle] (l4) at ($(l1) + (2,0)$) {5};
                % \draw[->] (l1) -- (l3) node[midway,above] {a};
                % \draw[->] (l3) -- (l4) node[midway,above] {b};
                \draw[->] (l4) -- (l2) node[midway,above] {a};
            }    
            \graphbox{$L$}{0mm}{-20mm}{34mm}{15mm}{2mm}{-5mm}{
                \coordinate (o) at (0mm,-3mm); 
                \node[draw,circle] (l1) at ($(o)+(-10mm,0mm)$) {1};
                \node[draw,circle] (l2) at ($(l1)+(2,0)$) {2};
                \node[draw,circle] (l3) at ($(l1) + (1,0)$) {5};
                \draw[->] (l1) -- (l3) node[midway,above] {a};
                \draw[->] (l3) -- (l2) node[midway,above] {a};
            }     
            \graphbox{$K$}{40mm}{-20mm}{24mm}{15mm}{2mm}{-5mm}{
                \coordinate (o) at (5mm,-3mm); 
                \node[draw,circle] (l1) at ($(o)+(-10mm,0mm)$) {1};
                \node[draw,circle] (l2) at ($(l1)+(1,0)$) {2};
                % \node[draw,circle] (l3) at ($(l1) + (1,0)$) {$\ $};
                % \draw[->] (l1) -- (l3) node[midway,above] {a};
                % \draw[->] (l3) -- (l2) node[midway,above] {a};
            }    
            \graphbox{$R$}{70mm}{-20mm}{45mm}{15mm}{2mm}{-5mm}{
                \coordinate (o) at (-5mm,-3mm); 
                \node[draw,circle] (l1) at ($(o)+(-10mm,0mm)$) {1};
                \node[draw,circle] (l2) at ($(l1)+(3,0)$) {2};
                \node[draw,circle] (l3) at ($(l1) + (1,0)$) {4};
                \node[draw,circle] (l4) at ($(l1) + (2,0)$) {5};
                \draw[->] (l1) -- (l3) node[midway,above] {a};
                \draw[->] (l3) -- (l4) node[midway,above] {b};
                \draw[->] (l4) -- (l2) node[midway,above] {a};
            }    
            \node () at (37mm,-8mm) {$\leftarrowtail$};
            \node () at (17mm,-17mm) {$\downarrowtail$};
            \node () at (11.5mm,-17.5mm) {$h_{R'_1L}$};
            \node () at (52mm,-17mm) {$\downarrowtail$};
            \node () at (92mm,-17mm) {$\downarrowtail$};
            \node () at (37mm,-18mm) {\text{PB}};
            \node () at (68mm,-18mm) {\text{PB}};
            \node () at (67mm,-8mm) {$\rightarrowtail$};
            \node () at (37mm,-28mm) {$\leftarrowtail$};
            \node () at (67mm,-28mm) {$\rightarrowtail$};
        \end{tikzpicture}
    } 
\end{center}   
% 2
\begin{center}
    \resizebox{0.7\textwidth}{!}{
        \begin{tikzpicture}
            \graphbox{$R'_2$}{0mm}{0mm}{34mm}{15mm}{2mm}{-5mm}{
                \coordinate (o) at (0mm,-3mm); 
                \node[draw,circle] (l1) at ($(o)+(-10mm,0mm)$) {1};
                \node[draw,circle] (l2) at ($(l1)+(2,0)$) {2};
                \node[draw,circle] (l3) at ($(l1) + (1,0)$) {5};
                % \draw[->] (l1) -- (l3) node[midway,above] {a};
                \draw[->] (l3) -- (l2) node[midway,above] {a};
            }     
            \graphbox{$K'$}{40mm}{0mm}{24mm}{15mm}{2mm}{-5mm}{
                \coordinate (o) at (5mm,-3mm); 
                \node[draw,circle] (l1) at ($(o)+(-10mm,0mm)$) {1};
                \node[draw,circle] (l2) at ($(l1)+(1,0)$) {2};
                % \node[draw,circle] (l3) at ($(l1) + (1,0)$) {$\ $};
                % \draw[->] (l1) -- (l3) node[midway,above] {a};
                % \draw[->] (l3) -- (l2) node[midway,above] {a};
            }    
            \graphbox{$R'_2$}{70mm}{0mm}{45mm}{15mm}{2mm}{-5mm}{
                \coordinate (o) at (-5mm,-3mm); 
                \node[draw,circle] (l1) at ($(o)+(-10mm,0mm)$) {1};
                \node[draw,circle] (l2) at ($(l1)+(3,0)$) {2};
                % \node[draw,circle] (l3) at ($(l1) + (1,0)$) {4};
                \node[draw,circle] (l4) at ($(l1) + (2,0)$) {5};
                % \draw[->] (l1) -- (l3) node[midway,above] {a};
                % \draw[->] (l3) -- (l4) node[midway,above] {b};
                \draw[->] (l4) -- (l2) node[midway,above] {a};
            }    
            \graphbox{$L$}{0mm}{-20mm}{34mm}{15mm}{2mm}{-5mm}{
                \coordinate (o) at (0mm,-3mm); 
                \node[draw,circle] (l1) at ($(o)+(-10mm,0mm)$) {1};
                \node[draw,circle] (l2) at ($(l1)+(2,0)$) {2};
                \node[draw,circle] (l3) at ($(l1) + (1,0)$) {5};
                \draw[->] (l1) -- (l3) node[midway,above] {a};
                \draw[->] (l3) -- (l2) node[midway,above] {a};
            }      
            \graphbox{$K$}{40mm}{-20mm}{24mm}{15mm}{2mm}{-5mm}{
                \coordinate (o) at (5mm,-3mm); 
                \node[draw,circle] (l1) at ($(o)+(-10mm,0mm)$) {1};
                \node[draw,circle] (l2) at ($(l1)+(1,0)$) {2};
                % \node[draw,circle] (l3) at ($(l1) + (1,0)$) {$\ $};
                % \draw[->] (l1) -- (l3) node[midway,above] {a};
                % \draw[->] (l3) -- (l2) node[midway,above] {a};
            }     
            \graphbox{$R$}{70mm}{-20mm}{45mm}{15mm}{2mm}{-5mm}{
                \coordinate (o) at (-5mm,-3mm); 
                \node[draw,circle] (l1) at ($(o)+(-10mm,0mm)$) {1};
                \node[draw,circle] (l2) at ($(l1)+(3,0)$) {2};
                \node[draw,circle] (l3) at ($(l1) + (1,0)$) {4};
                \node[draw,circle] (l4) at ($(l1) + (2,0)$) {5};
                \draw[->] (l1) -- (l3) node[midway,above] {a};
                \draw[->] (l3) -- (l4) node[midway,above] {b};
                \draw[->] (l4) -- (l2) node[midway,above] {a};
            }    
            \node () at (37mm,-8mm) {$\leftarrowtail$};
            \node () at (17mm,-17mm) {$\downarrowtail$};
            \node () at (11.5mm,-17.5mm) {$h_{R'_2L}$};
            \node () at (52mm,-17mm) {$\downarrowtail$};
            \node () at (92mm,-17mm) {$\downarrowtail$};
            \node () at (37mm,-18mm) {\text{PB}};
            \node () at (68mm,-18mm) {\text{PB}};
            \node () at (67mm,-8mm) {$\rightarrowtail$};
            \node () at (37mm,-28mm) {$\leftarrowtail$};
            \node () at (67mm,-28mm) {$\rightarrowtail$};
        \end{tikzpicture}
    } 
\end{center}
%3
\begin{center}
    \resizebox{0.7\textwidth}{!}{
        \begin{tikzpicture}
            \graphbox{$R'_3$}{0mm}{0mm}{34mm}{15mm}{2mm}{-5mm}{
                \coordinate (o) at (0mm,-3mm); 
                \node[draw,circle] (l1) at ($(o)+(-10mm,0mm)$) {1};
                % \node[draw,circle] (l2) at ($(l1)+(2,0)$) {2};
                \node[draw,circle] (l3) at ($(l1) + (1,0)$) {4};
                \draw[->] (l1) -- (l3) node[midway,above] {a};
                % \draw[->] (l3) -- (l2) node[midway,above] {a};
            }     
            \graphbox{$K'$}{40mm}{0mm}{24mm}{15mm}{2mm}{-5mm}{
                \coordinate (o) at (5mm,-3mm); 
                \node[draw,circle] (l1) at ($(o)+(-10mm,0mm)$) {1};
                % \node[draw,circle] (l2) at ($(l1)+(1,0)$) {2};
                % \node[draw,circle] (l3) at ($(l1) + (1,0)$) {$\ $};
                % \draw[->] (l1) -- (l3) node[midway,above] {a};
                % \draw[->] (l3) -- (l2) node[midway,above] {a};
            }     
            \graphbox{$R'_3$}{70mm}{0mm}{45mm}{15mm}{2mm}{-5mm}{
                \coordinate (o) at (-5mm,-3mm); 
                \node[draw,circle] (l1) at ($(o)+(-10mm,0mm)$) {1};
                % \node[draw,circle] (l2) at ($(l1)+(3,0)$) {4};
                \node[draw,circle] (l3) at ($(l1) + (1,0)$) {4};
                % \node[draw,circle] (l4) at ($(l1) + (2,0)$) {5};
                \draw[->] (l1) -- (l3) node[midway,above] {a};
                % \draw[->] (l3) -- (l4) node[midway,above] {b};
                % \draw[->] (l4) -- (l2) node[midway,above] {a};
            }    
            \graphbox{$L$}{0mm}{-20mm}{34mm}{15mm}{2mm}{-5mm}{
                \coordinate (o) at (0mm,-3mm); 
                \node[draw,circle] (l1) at ($(o)+(-10mm,0mm)$) {1};
                \node[draw,circle] (l2) at ($(l1)+(2,0)$) {2};
                \node[draw,circle] (l3) at ($(l1) + (1,0)$) {4};
                \draw[->] (l1) -- (l3) node[midway,above] {a};
                \draw[->] (l3) -- (l2) node[midway,above] {a};
            }     
            \graphbox{$K$}{40mm}{-20mm}{24mm}{15mm}{2mm}{-5mm}{
                \coordinate (o) at (5mm,-3mm); 
                \node[draw,circle] (l1) at ($(o)+(-10mm,0mm)$) {1};
                \node[draw,circle] (l2) at ($(l1)+(1,0)$) {2};
                % \node[draw,circle] (l3) at ($(l1) + (1,0)$) {$\ $};
                % \draw[->] (l1) -- (l3) node[midway,above] {a};
                % \draw[->] (l3) -- (l2) node[midway,above] {a};
            }    
            \graphbox{$R$}{70mm}{-20mm}{45mm}{15mm}{2mm}{-5mm}{
                \coordinate (o) at (-5mm,-3mm); 
                \node[draw,circle] (l1) at ($(o)+(-10mm,0mm)$) {1};
                \node[draw,circle] (l2) at ($(l1)+(3,0)$) {2};
                \node[draw,circle] (l3) at ($(l1) + (1,0)$) {4};
                \node[draw,circle] (l4) at ($(l1) + (2,0)$) {5};
                \draw[->] (l1) -- (l3) node[midway,above] {a};
                \draw[->] (l3) -- (l4) node[midway,above] {b};
                \draw[->] (l4) -- (l2) node[midway,above] {a};
            }    
            \node () at (37mm,-8mm) {$\leftarrowtail$};
            \node () at (17mm,-17mm) {$\downarrowtail$};
            \node () at (11.5mm,-17.5mm) {$h_{R'_3L}$};
            \node () at (52mm,-17mm) {$\downarrowtail$};
            \node () at (92mm,-17mm) {$\downarrowtail$};
            \node () at (37mm,-18mm) {\text{PB}};
            \node () at (68mm,-18mm) {\text{PB}};
            \node () at (67mm,-8mm) {$\rightarrowtail$};
            \node () at (37mm,-28mm) {$\leftarrowtail$};
            \node () at (67mm,-28mm) {$\rightarrowtail$};
        \end{tikzpicture}
    } 
\end{center}
%4
\begin{center}
    \resizebox{0.7\textwidth}{!}{
        \begin{tikzpicture}
            \graphbox{$R'_4$}{0mm}{0mm}{34mm}{15mm}{2mm}{-5mm}{
                \coordinate (o) at (0mm,-3mm); 
                \node[draw,circle] (l1) at ($(o)+(-10mm,0mm)$) {1};
                \node[draw,circle] (l2) at ($(l1)+(2,0)$) {2};
                \node[draw,circle] (l3) at ($(l1) + (1,0)$) {4};
                \draw[->] (l1) -- (l3) node[midway,above] {a};
                % \draw[->] (l3) -- (l2) node[midway,above] {a};
            }     
            \graphbox{$K'$}{40mm}{0mm}{24mm}{15mm}{2mm}{-5mm}{
                \coordinate (o) at (5mm,-3mm); 
                \node[draw,circle] (l1) at ($(o)+(-10mm,0mm)$) {1};
                \node[draw,circle] (l2) at ($(l1)+(1,0)$) {2};
                % \node[draw,circle] (l3) at ($(l1) + (1,0)$) {$\ $};
                % \draw[->] (l1) -- (l3) node[midway,above] {a};
                % \draw[->] (l3) -- (l2) node[midway,above] {a};
            }     
            \graphbox{$R'_4$}{70mm}{0mm}{45mm}{15mm}{2mm}{-5mm}{
                \coordinate (o) at (-5mm,-3mm); 
                \node[draw,circle] (l1) at ($(o)+(-10mm,0mm)$) {1};
                \node[draw,circle] (l2) at ($(l1)+(3,0)$) {2};
                \node[draw,circle] (l3) at ($(l1) + (1,0)$) {4};
                % \node[draw,circle] (l4) at ($(l1) + (2,0)$) {5};
                \draw[->] (l1) -- (l3) node[midway,above] {a};
                % \draw[->] (l3) -- (l4) node[midway,above] {b};
                % \draw[->] (l4) -- (l2) node[midway,above] {a};
            }    
            \graphbox{$L$}{0mm}{-20mm}{34mm}{15mm}{2mm}{-5mm}{
                \coordinate (o) at (0mm,-3mm); 
                \node[draw,circle] (l1) at ($(o)+(-10mm,0mm)$) {1};
                \node[draw,circle] (l2) at ($(l1)+(2,0)$) {2};
                \node[draw,circle] (l3) at ($(l1) + (1,0)$) {4};
                \draw[->] (l1) -- (l3) node[midway,above] {a};
                \draw[->] (l3) -- (l2) node[midway,above] {a};
            }     
            \graphbox{$K$}{40mm}{-20mm}{24mm}{15mm}{2mm}{-5mm}{
                \coordinate (o) at (5mm,-3mm); 
                \node[draw,circle] (l1) at ($(o)+(-10mm,0mm)$) {1};
                \node[draw,circle] (l2) at ($(l1)+(1,0)$) {2};
                % \node[draw,circle] (l3) at ($(l1) + (1,0)$) {$\ $};
                % \draw[->] (l1) -- (l3) node[midway,above] {a};
                % \draw[->] (l3) -- (l2) node[midway,above] {a};
            }    
            \graphbox{$R$}{70mm}{-20mm}{45mm}{15mm}{2mm}{-5mm}{
                \coordinate (o) at (-5mm,-3mm); 
                \node[draw,circle] (l1) at ($(o)+(-10mm,0mm)$) {1};
                \node[draw,circle] (l2) at ($(l1)+(3,0)$) {2};
                \node[draw,circle] (l3) at ($(l1) + (1,0)$) {4};
                \node[draw,circle] (l4) at ($(l1) + (2,0)$) {5};
                \draw[->] (l1) -- (l3) node[midway,above] {a};
                \draw[->] (l3) -- (l4) node[midway,above] {b};
                \draw[->] (l4) -- (l2) node[midway,above] {a};
            }    
            \node () at (37mm,-8mm) {$\leftarrowtail$};
            \node () at (17mm,-17mm) {$\downarrowtail$};
            \node () at (11.5mm,-17.5mm) {$h_{R'_4L}$};
            \node () at (52mm,-17mm) {$\downarrowtail$};
            \node () at (92mm,-17mm) {$\downarrowtail$}; 
            \node () at (37mm,-18mm) {\text{PB}};
            \node () at (68mm,-18mm) {\text{PB}};
            \node () at (67mm,-8mm) {$\rightarrowtail$};
            \node () at (37mm,-28mm) {$\leftarrowtail$};
            \node () at (67mm,-28mm) {$\rightarrowtail$};
        \end{tikzpicture}
    } 
\end{center}
To illustrate the $X$-non-increasing property, consider the rewriting step defined by the DPO diagram below. Whenever \( R'_i \) (for \( 1 \leq i \leq 4 \)) forms an \( X \)-occurrence $x$ that is included neither in \( R \) nor $C$ with a subgraph $C' \subseteq C \cap x \subseteq C$, the graph \( h_{R'_iL}(R'_i) \) forms a \( X \)-occurrence in $G$ with $C'$. Furthermore, distinct $X$-occurrences created in $H$ have distinct corresponding $X$-occurrences in $G$.
    \begin{center} 
        \resizebox{0.7\textwidth}{!}{
        \begin{tikzpicture}
            \graphbox{\( L \)}{0mm}{-3mm}{34mm}{12mm}{2mm}{2mm}{
                \coordinate (o) at (0mm,-8mm); 
                \node[draw,circle] (l1) at ($(o)+(-10mm,0mm)$) {1};
                \node[draw,circle] (l2) at ($(l1)+(2,0)$) {2};
                \node[draw,circle] (l3) at ($(l1) + (1,0)$) {3};
                \draw[->] (l1) -- (l3) node[midway,above] {a};
                \draw[->] (l3) -- (l2) node[midway,above] {a};
            } 

            \graphbox{\( K \)}{40mm}{-3mm}{34mm}{12mm}{2mm}{2mm}{
                \coordinate (o) at (0mm,-8mm); 
                \node[draw,circle] (l1) at ($(o)+(-10mm,0mm)$) {1};
                \node[draw,circle] (l2) at ($(l1)+(2,0)$) {2};
            }  

            \graphbox{\( R \)}{80mm}{-3mm}{45mm}{12mm}{2mm}{2mm}{
                \coordinate (o) at (-5mm,-8mm); 
                \node[draw,circle] (l1) at ($(o)+(-10mm,0mm)$) {1};
                \node[draw,circle] (l2) at ($(l1)+(3,0)$) {2};
                \node[draw,circle] (l3) at ($(l1) + (1,0)$) {4};
                \node[draw,circle] (l4) at ($(l1) + (2,0)$) {5};
                \draw[->] (l1) -- (l3) node[midway,above] {a};
                \draw[->] (l3) -- (l4) node[midway,above] {b};
                \draw[->] (l4) -- (l2) node[midway,above] {a};
            }    

            \graphbox{\( G \)}{0mm}{-22mm}{34mm}{22mm}{2mm}{-3mm}{
                \coordinate (o) at (0mm,-3mm); 
                \node[draw,circle] (l1) at ($(o)+(-10mm,0mm)$) {1};
                \node[draw,circle] (l2) at ($(l1)+(2,0)$) {2};
                \node[draw,circle] (l3) at ($(l1) + (1,0)$) {3};
                \node[draw,circle] (l4) at ($(l2) + (0,-1)$) {6};
                \draw[->] (l1) -- (l3) node[midway,above] {a};
                \draw[->] (l3) -- (l2) node[midway,above] {a};
                \draw[->] (l2) -- (l4) node[midway,right] {a};
                \node[draw,circle] (l6) at ($(l1) + (0,-1)$) {7};
                \draw[<-] (l1) -- (l6) node[midway,left] {a};
                \draw[->] (l2) edge[out=-135,in=-45]node[midway,below] {a} (l1) ;
            }    

            \graphbox{\( C  \)}{40mm}{-22mm}{34mm}{22mm}{2mm}{-3mm}{
                \coordinate (o) at (0mm,-3mm); 
                \node[draw,circle] (l1) at ($(o)+(-10mm,0mm)$) {1};
                \node[draw,circle] (l2) at ($(l1)+(2,0)$) {2};
                \node[draw,circle] (l4) at ($(l2) + (0,-1)$) {6};
                \draw[->] (l2) -- (l4) node[midway,right] {a};
                \draw[->] (l2) edge[out=-135,in=-45]node[midway,below] {a} (l1) ;
                \node[ draw,circle] (l6) at ($(l1) + (0,-1)$) {7};
                \draw[<-] (l1) -- (l6) node[midway,left] {a};
            }    

            \graphbox{\( H \)}{80mm}{-22mm}{45mm}{22mm}{2mm}{-3mm}{
                \coordinate (o) at (-5mm,-3mm); 
                \node[draw,circle] (l1) at ($(o)+(-10mm,0mm)$) {1};
                \node[draw,circle] (l2) at ($(l1)+(3,0)$) {2};
                \node[draw,circle] (l3) at ($(l1) + (1,0)$) {4};
                \node[draw,circle] (l4) at ($(l1) + (2,0)$) {5};
                \node[ draw,circle] (l5) at ($(l2) + (0,-1)$) {6};
                \node[ draw,circle] (l6) at ($(l1) + (0,-1)$) {7};
                \draw[<-] (l1) -- (l6) node[midway,left] {a};
                \draw[->] (l1) -- (l3) node[midway,above] {a};
                \draw[->] (l3) -- (l4) node[midway,above] {b};
                \draw[->] (l4) -- (l2) node[midway,above] {a};
                \draw[->] (l2) -- (l5) node[midway,right] {a};
                \draw[->] (l2) edge[out=-135,in=-45]node[midway,below] {a} (l1) ;
            }    

            \node () at (37mm,-8mm) {\( \leftarrowtail \)}; % K -> L
            \node () at (77mm,-8mm) {\( \rightarrowtail \)}; % K -> R
            \node () at (15mm,-18mm) {\( m\ \downarrowtail \)};
            \node () at (37mm,-33mm) {\( \leftarrowtail \)};
            \node () at (58mm,-18mm) {\( \downarrowtail \)};
            \node () at (102mm,-18mm) {\( \downarrowtail \)};
            \node () at (77mm,-33mm) {\( \rightarrowtail \)}; % C -> H
        \end{tikzpicture}
        }
    \end{center}
\end{example}

% cond 2 necessaire  1
\begin{example} 
    \label{ex:cond_2_necessaire}
    Let $X$ be the graph 
    \tikz[baseline=-0.5ex]{ 
            \node (x) at (0,0) {$\bullet$}; 
            \node (y) at (1,0) {$\bullet$};
            \node (z) at (2,0) {$\bullet$};
            \draw[->] (x) -- (y)   {};
            \draw[->] (y) -- (z)   {};
    }. The following rewriting rule is potentially $X$-increasing.
    \begin{center}
        \resizebox{0.7\textwidth}{!}{
            \begin{tikzpicture}
                \graphbox{$L$}{0mm}{0mm}{34mm}{15mm}{2mm}{-5mm}{
                    \coordinate (o) at (0mm,-3mm); 
                    \node[draw,circle] (l1) at ($(o)+(-10mm,0mm)$) {1};
                    % \node[draw,circle] (l2) at ($(l1)+(2,0)$) {2};
                    \node[draw,circle] (l3) at ($(l1) + (1,0)$) {2};
                    \draw[->] (l3) -- (l1) node[midway,above] {};
                    % \draw[->] (l3) -- (l2) node[midway,above] {a};
                }     
                \graphbox{$K$}{40mm}{0mm}{24mm}{15mm}{2mm}{-5mm}{
                    \coordinate (o) at (5mm,-3mm); 
                    \node[draw,circle] (l1) at ($(o)+(-10mm,0mm)$) {1};
                    \node[draw,circle] (l2) at ($(l1)+(1,0)$) {2};
                    % \node[draw,circle] (l3) at ($(l1) + (1,0)$) {$\ $};
                    \draw[->] (l2) -- (l1) node[midway,above] {};
                    % \draw[->] (l3) -- (l2) node[midway,above] {a};
                }    
                \graphbox{$R$}{70mm}{0mm}{45mm}{15mm}{2mm}{-5mm}{
                    \coordinate (o) at (-5mm,-3mm); 
                    \node[draw,circle] (l1) at ($(o)+(-10mm,0mm)$) {1};
                    % \node[draw,circle] (l2) at ($(l1)+(3,0)$) {2};
                    \node[draw,circle] (l3) at ($(l1) + (1,0)$) {2};
                    \node[draw,circle] (l4) at ($(l1) + (2,0)$) {3};
                    \draw[->] (l3) -- (l1) node[midway,above] {};
                    \draw[->] (l3) -- (l4) node[midway,above] {};
                    % \draw[->] (l4) -- (l2) node[midway,above] {a};
                }    
                \node () at (37mm,-8mm) {$\leftarrowtail$};
                \node () at (67mm,-8mm) {$\rightarrowtail$};
            \end{tikzpicture}
        }
    \end{center}
    As shown in \autoref{ex:rx_counter_ex}, \( D(R,X) \) contains the graph $R'$:
    \raisebox{2pt}{
        \scalebox{0.6}{\tikz[baseline=-0.5ex]{
        \node [draw,circle] (x) at (0,0) {2};
        \node[draw,circle] (y) at (1,0) {3};
        \draw[->] (x) -- (y) {};
    }}}. The unique monomorphism $h_{R'L}:R' \rightarrowtail L$ fails the second condition of \autoref{def:creates_more_x_on_the_left}. 
    %  as illustrated below 
    % \begin{center}
    %     \resizebox{0.6\textwidth}{!}{
    %     \begin{tikzpicture}
    %             \graphbox{$R'$}{40mm}{-20mm}{24mm}{15mm}{2mm}{-5mm}{
    %                 \coordinate (o) at (5mm,-3mm); 
    %                 \node[draw,circle] (l1) at ($(o)+(-10mm,0mm)$) {2};
    %                 \node[draw,circle] (l2) at ($(l1)+(1,0)$) {3};
    %                 % \node[draw,circle] (l3) at ($(l1) + (1,0)$) {$\ $};
    %                 \draw[->] (l1) -- (l2) node[midway,above] {};
    %                 % \draw[->] (l3) -- (l2) node[midway,above] {a};
    %             }    
    %             \graphbox{$L$}{70mm}{-20mm}{45mm}{15mm}{2mm}{-5mm}{
    %                 \coordinate (o) at (-5mm,-3mm); 
    %                 \node[draw,circle] (l1) at ($(o)+(-10mm,0mm)$) {3};
    %                 % \node[draw,circle] (l2) at ($(l1)+(3,0)$) {2};
    %                 \node[draw,circle] (l2) at ($(l1) + (1,0)$) {2};
    %                 \draw[->] (l2) -- (l1) node[midway,above] {};
    %                 % \draw[->] (l4) -- (l2) node[midway,above] {a};
    %             }    
    %             \node () at (67mm,-28mm) {$\rightarrowtail$};
    %         \end{tikzpicture}
    %     }
    % \end{center}
    For any rewriting step using this rule, the corresponding $X$-occurrence relative to $h_{R'L}$ of any implicitly created $X$-occurrence is included in the context.
\end{example}
% cond 2 necessaire  2
\begin{example}
    \label{ex:cond_2_necessaire2}
    Let $X$ be the graph 
    \tikz[baseline=-0.5ex]{ 
            \node (x) at (0,0) {$\bullet$}; 
            \node (y) at (1,0) {$\bullet$};
            \node (z) at (2,0) {$\bullet$};
            \draw[->] (x) -- (y)   {};
            % \draw[->] (y) -- (z)   {};
    }. The following rewriting rule is potentially $X$-increasing.
    \begin{center}
        \resizebox{0.7\textwidth}{!}{
            \begin{tikzpicture}
                \graphbox{$L$}{0mm}{0mm}{34mm}{15mm}{2mm}{-5mm}{
                    \coordinate (o) at (0mm,-3mm); 
                    \node[draw,circle] (l1) at ($(o)+(-10mm,0mm)$) {1};
                    \node[draw,circle] (l3) at ($(l1) + (1,0)$) {2};
                }     
                \graphbox{$K$}{40mm}{0mm}{24mm}{15mm}{2mm}{-5mm}{
                    \coordinate (o) at (5mm,-3mm); 
                    \node[draw,circle] (l1) at ($(o)+(-10mm,0mm)$) {1};
                    \node[draw,circle] (l2) at ($(l1)+(1,0)$) {2};
                }    
                \graphbox{$R$}{70mm}{0mm}{45mm}{15mm}{2mm}{-5mm}{
                    \coordinate (o) at (-5mm,-3mm); 
                    \node[draw,circle] (l1) at ($(o)+(-10mm,0mm)$) {1};
                    \node[draw,circle] (l3) at ($(l1) + (1,0)$) {2};
                    \node[draw,circle] (l4) at ($(l1) + (2,0)$) {3};
                }    
                \node () at (37mm,-8mm) {$\leftarrowtail$};
                \node () at (67mm,-8mm) {$\rightarrowtail$};
            \end{tikzpicture}
        }
    \end{center}
    The set \( D(R,X) \) contains $R'$:
    \raisebox{2pt}{ \scalebox{0.6}{\tikz[baseline=-0.5ex]{
        \node [draw,circle] (x) at (0,0) {1};
        \node[draw,circle] (y) at (1,0) {3};
    }}}.  
    Despite the existence of a unique monomorphism $h_{R'L}: R' \rightarrowtail L$ preserving interface elements, this rule fails the second condition of \autoref{def:creates_more_x_on_the_left}.
    For any rewriting step using this rule, the corresponding $X$-occurrence relative to $h_{R'L}$ of any implicitly created $X$-occurrence is included in the context.
\end{example}
\begin{example}
    \label{ex:cond3_necessaire}
    Let $X$ be the graph 
    \tikz[baseline=-0.5ex]{ 
            \node (x) at (0,0) {$\bullet$}; 
            \node (y) at (1,0) {$\bullet$};
            \node (z) at (2,0) {$\bullet$};
            \draw[->] (x) -- (y)   {};
            \draw[->] (y) -- (z) {};
    }. The following rewriting rule is potentially $X$-increasing.
    \begin{center}
        \resizebox{0.6\textwidth}{!}{
            \begin{tikzpicture}
                \graphbox{$L$}{0mm}{0mm}{30mm}{20mm}{0}{0}{
                    \node[draw,circle]  (x) at (-6mm,-12mm) {x};
                    \node[draw,circle] (y) at (6mm,-12mm) {};
                    % \node[draw,circle]  (z) at (6mm,0mm) {};
                    \draw[->]  (x) to (y);
                    % \draw[->] (y) to[bend right=20] (z);
                    % \draw[->]  (z) to[bend right=20] (y);
                }
                \graphbox{$K$}{40mm}{0mm}{30mm}{20mm}{0}{0}{
                    \node[draw,circle]  (x) at (-6mm,-12mm) {x};
                    % \node[draw,circle]  (y) at (6mm,-12mm) {y};
                }
                \graphbox{$R$}{80mm}{0mm}{30mm}{20mm}{0}{0}{
                    \node[draw,circle]  (x) at (-6mm,-12mm) {x};
                        \node[draw,circle]  (y) at (6mm,-8mm) {y};
                        \node[draw,circle]  (z) at (6mm,-16mm) {z};
                        \draw[->]  (x) to (y);
                        \draw[->]  (x) to (z);
                }
                \node () at (35mm,-12mm) {$\leftarrowtail$};
                \node () at (75mm,-12mm) {$\rightarrowtail$};
            \end{tikzpicture}
        }
    \end{center}
    The set \( D(R,X) \) contains exactly two elements $R'$:
    \raisebox{2pt}{ 
        \scalebox{0.6}{\tikz[baseline=-0.5ex]{
        \node [draw,circle] (x) at (0,0) {x};
        \node[draw,circle] (y) at (1,0) {y};
        \draw[->] (x) -- (y) {};
    }}} and $R''$:\raisebox{2pt}{ 
        \scalebox{0.6}{\tikz[baseline=-0.5ex]{
        \node [draw,circle] (x) at (0,0) {x};
        \node[draw,circle] (y) at (1,0) {z};
        \draw[->] (x) -- (y) {};
    }}}. 
    Despite unique monomorphisms \( h_{R'L}: R' \rightarrowtail L \) and \( h_{R''L}: R'' \rightarrowtail L \) preserving interface elements, they fail the third condition of \autoref{def:creates_more_x_on_the_left}. 
    For any rewriting step using this rule, implicitly created $X$-occurrences with the same subgraph of the context have the same corresponding $X$-occurrence. 
\end{example}
 
\begin{example}
    \label{ex:cond4_necessaire}
    Let $X$ be the graph 
    \tikz[baseline=-0.5ex]{ 
            \node (x) at (0,0) {$\bullet$}; 
            \node (y) at (1,0) {$\bullet$};
            \node (z) at (2,0) {$\bullet$};
            \draw[->] (x) -- (y)   {};
    }, which has an isolated node. The following rewriting rule is potentially $X$-increasing.
    \begin{center}
        \resizebox{0.6\textwidth}{!}{
            \begin{tikzpicture}
                \graphbox{$L$}{0mm}{0mm}{30mm}{20mm}{0}{0}{
                    \node[draw,circle]  (x) at (-6mm,-10mm) {x};
                    \node[draw,circle] (y) at (6mm,-10mm) {y};
                }
                \graphbox{$K$}{40mm}{0mm}{30mm}{20mm}{0}{0}{
                    \node[draw,circle]  (x) at (-6mm,-10mm) {x};
                }
                \graphbox{$R$}{80mm}{0mm}{30mm}{20mm}{0}{0}{
                    \node[draw,circle]  (x) at (-6mm,-10mm) {x};
                        \node[draw,circle]  (y) at (6mm,-6mm) {z};
                        \node[draw,circle]  (z) at (6mm,-14mm) {w};
                }
                \node () at (35mm,-10mm) {$\leftarrowtail$};
                \node () at (75mm,-10mm) {$\rightarrowtail$};
            \end{tikzpicture}
        }
    \end{center}
    The set \( D(R,X) \) contains exactly two elements $R'$:
    \raisebox{2pt}{\scalebox{0.6}{\tikz[baseline=-0.5ex]{
        \node [draw,circle] (x) at (0,0) {x};
        \node[draw,circle] (y) at (1,0) {z};
        % \draw[->] (x) -- (y) {};
    }}} and $R''$:\raisebox{2pt}{\scalebox{0.6}{\tikz[baseline=-0.5ex]{
        \node [draw,circle] (x) at (0,0) {x};
        \node[draw,circle] (y) at (1,0) {w};
        % \draw[->] (x) -- (y) {};
    }}}. 
    There are unique monomorphisms $h_{R'L}:R' \rightarrowtail L$ and $h_{R''L}:R'' \rightarrowtail L$ preserving interface elements, but they fail the fourth condition of \autoref{def:creates_more_x_on_the_left}.
    For any rewriting step using this rule, implicitly created $X$-occurrences with the same subgraph of the context have the same corresponding $X$-occurrence.
\end{example}


\section{Solution to the key challenge}
\label{sec:solution_to_the_key_challenge}
The following lemma (see \iflongversion
Appendix \hyperref[proof:lem:w_u_l_not_geq_r_not]{Section C}
\else
\cite[Lemma 40]{qiu2025termination}
\fi for a proof) ensures that if $\rho$ is $X$-non-increasing, then, for any rewriting step using $\rho$, 
more X-occurrences are implicitly destroyed than implicitly created by the rewriting step.
% the cardinality of the set of \( X \)-occurrences implicitly destroyed by the rewriting step is at least as large as the cardinality of the set of \( X \)-occurrences implicitly created by the rewriting step.
% There are more implicitly destroyed X-occurrences than implicitly created X-occurrences.

\begin{lemma} 
    \label{lem:w_u_l_not_geq_r_not}
    \ \newline 
    \noindent 
    \begin{minipage}{0.7\textwidth}
        Let $X$ be a ruler-graph and $\rho = (L \overset{l}{\leftarrowtail} K \overset{r}{\rightarrowtail} R)$ an injective DPO rewriting rule.
        Suppose that $\rho$ is $X$-non-increasing. For every rewriting step induced by the DPO diagram (shown on the right), we have 
        \[
            |\operatorname{Mono}(X, G, \lnot m, \lnot l')| \geq |\operatorname{Mono}(X, H, \lnot m', \lnot r')|
        \]
    \end{minipage}
    \hfill
    \begin{minipage}{0.29\textwidth}
        \hfill
        \resizebox{0.9\textwidth}{!}{
                \begin{tikzpicture}
            \node (k) at (0,1) {K};
            \node (l) at (-2,1) {L};
            \node (r) at (2,1) {R};
            \node (c) at (0,-1) {C};
            \node (g) at (-2,-1) {G};
            \node (h) at (2,-1) {H};
            \draw[<-<]  (l) -- (k) node [midway,below] {l};
            \draw[>->]  (k) -- (r) node [midway,below] {r};
            \draw[>->] (c) -- (g) node [midway, above] {l'};
            \draw[>->] (c) -- (h) node [midway,above] {r'};
            \draw[>->] (l) -- (g) node[midway, right] {m};
            \draw[>->] (r) -- (h) node[midway, left] {m'};
            \draw[>->] (k) -- (c) node[midway, left] {};
            \node () [at=($(l)!0.5!(c)$)] {PO};
            \node () [at=($(r)!0.5!(c)$)] {PO};
        \end{tikzpicture}
        }
    \end{minipage}
\end{lemma}
 
  
\section{Termination criterion}
\label{subgraph_counting:sec:termination} 
Finally, we present a lemma that allows us to overapproximate the change in weight of a graph when a rule is applied, followed by our termination criterion. 

\begin{lemma}[Decreasing step]
    \label{subgraph_counting:lem:w_g_geq_w_h_leq}
    Let $\rho = (L \overset{l}{\leftarrowtail} K \overset{r}{\rightarrowtail} R)$ be an injective DPO rewriting rule,
    \( \mathbb{X} \) a set of ruler-graphs,
    \( s_{\mathbb{X}} \colon \mathbb{X} \to \mathbb{N} \) a weight function,
    and \( G \Rightarrow_{\rho,\mathfrak{M}} H \) a rewriting step. 
    If $\rho$ is \( X \)-non-increasing for every ruler-graph \( X \in \mathbb{X} \), then $
        w_{s_\mathbb{X}}(G) - w_{s_\mathbb{X}}(H) 
        \geq 
        w_{s_\mathbb{X}}(L) - w_{s_\mathbb{X}}(R)
    $.
\end{lemma}
\begin{proof}
    \iflongversion
        See~\textsection~\ref{subgraph_counting:proof:lem:w_g_geq_w_h_leq}.
    \else
        See \cite[Lemma 41]{qiu2025termination}.
    \fi 
\end{proof} 
\begin{theorem}[Termination] 
    \label{subgraph_counting:thm:termination_grs}
    Let \(\mathcal{A}\) and \(\mathcal{B}\) be sets of injective DPO rewriting rules, $\mathbb{X}$ a set of ruler-graphs, and $s_\mathbb{X}$ a weight function. If the following hold:
    \begin{enumerate}
        \item  for every $\rho \in \mathcal{A} \cup \mathcal{B}$ and for every $X \in \mathbb{X}$, the rule $\rho$ is $X$-non-increasing,
        \item for every \(\rho \in \mathcal{A}\), we have \( w_{s_\mathbb{X}}(lhs(\rho)) > w_{s_\mathbb{X}}(rhs(\rho)) \),
        \item for every \(\rho \in \mathcal{B}\), we have \( w_{s_\mathbb{X}}(lhs(\rho)) \geq w_{s_\mathbb{X}}(rhs(\rho)) \).
    \end{enumerate}
    Then \(\Rightarrow_{\mathcal{A},\mathcal{M}}\) terminates relative to \(\Rightarrow_{\mathcal{B},\mathcal{M}}\).
\end{theorem}
\begin{proof}
    \iflongversion
        See~\textsection~\ref{subgraph_counting:proof:lem:w_g_geq_w_h_leq}.
    \else
        See \cite[Lemma 41]{qiu2025termination}.
    \fi 
\end{proof} 
\noindent Our technique is described for DPO rewriting with injective matches due to their expressiveness \cite{habel2001double}, but it also extends to non-injective matches.

\begin{definition}[Quotient Rule \cite{habel2001double}] 
    \ \newline 
    \noindent
    \begin{minipage}{0.7\textwidth}
        Given a rule $\rho = L \leftarrow K \rightarrow R$, a rule $\rho' = (L' \leftarrow K' \rightarrow R')$ is called a \textbf{quotient rule} of $\rho$ if there exists a DPO diagram (shown on the right)
    where all vertical graph morphisms are surjective. The set of quotient rules of $\rho$ is denoted by $Q(\rho)$, and 
    for a rule set $\mathcal{A}$, we write $Q(\mathcal{A}) = \bigcup_{\rho\in\mathcal{A}} Q(\rho)$.
    \end{minipage}
    \hfill
    \begin{minipage}{0.29\textwidth}
        \hfill
        \begin{tikzpicture}[scale=0.7]
            \node (k) at (0,0) {K};
            \node (l) at (-2,0) {L};
            \node (r) at (2,0) {R};
            \node (k') at (0,-2) {K'};
            \node (l') at (-2,-2) {L'};
            \node (r') at (2,-2)  {R'};
            \draw[->] (k) -> (l);
            \draw[->] (k) -> (r); 
            \draw[->] (k') -> (l'); 
            \draw[->] (k') -> (r'); 
            \draw[->] (k) -> (k'); 
            \draw[->] (l) -> (l'); 
            \draw[->] (r) -> (r'); 
            \node () [at=($(l)!0.5!(k')$)] {PO};
            \node () [at=($(r)!0.5!(k')$)] {PO};
        \end{tikzpicture}
    \end{minipage}
\end{definition}

\begin{lemma}[Quotient Lemma \cite{habel2001double}]
    Let $\rho$ be a DPO rewriting rule. For any graphs $G$ and $H$,
    $G \Rightarrow_{\rho,\mathfrak{F}} H$ if and only if $G \Rightarrow_{Q(\rho),\mathfrak{M}} H$.
\end{lemma}
\begin{corollary}
    \label{cor:termination}
    Let \(\mathcal{A}\) and \(\mathcal{B}\) be sets of injective DPO rewriting rules. 
    The rewriting relation $\Rightarrow_{\mathcal{A},\mathfrak{F}}$ is terminating relative to $\Rightarrow_{\mathcal{B},\mathfrak{F}}$ 
    if 
    $\Rightarrow_{Q(\mathcal{A}),\mathfrak{M}}$ is terminating relative to $\Rightarrow_{Q(\mathcal{B}),\mathfrak{M}}$.
\end{corollary} 

\section{Examples} 
\label{subgraph_counting:sec:examples}
While the type graph method~\cite{bruggink2014termination, bruggink2015proving,endrullis2024generalized_arxiv_v2} can prove termination for this example, it requires constructing an appropriate weighted type graph, a task that is difficult in general~\cite[\textsection 6]{bruggink2015proving}. 
In contrast, our criterion requires checking simpler conditions.
\begin{example}
    \label{subgraph_counting:ex:termination:grsaa}
    Consider the rewriting rule $\rho$ in Example~\ref{subgraph_counting:ex:grsaa_rx}, let $X$ be the graph \tikz[baseline=-0.5ex]{
        \node[draw,circle] (x) at (0,0) { };
        \node[draw,circle] (y) at (1,0) { }; 
        \node[draw,circle] (z) at (2,0) { };
        \draw[->] (x) -- (y) node[midway, above] {$a$};
        \draw[<-] (z) -- (y) node[midway, above] {$a$};
    }, $\mathbb{X} = \{X\}$ and $s_\mathbb{X}(X) = 1$. The rule is $X$-non-increasing as shown in Example~\ref{example:grs_aa:has_more_left}. 
    Since \(w_{s_\mathbb{X}}(\operatorname{lhs}(\rho)) = 1 > 0 = w_{s_\mathbb{X}}(\operatorname{rhs}(\rho)) \),
    it terminates by Theorem~\ref{subgraph_counting:thm:termination_grs}.
\end{example}
We consider an example for which the techniques from \cite{bruggink2014termination,bruggink2015proving,endrullis2024generalized_arxiv_v2,plump2018modular,overbeek2024termination_lmcs} fail.
\begin{example} 
    \label{subgraph_counting:ex_contrib_variant}
    The rewriting rule shown below 
    % in Figure~\ref{fig:subgraph_counting:ex_confdkjfakljl} 
    is a variant of a rule presented in \cite[Example 6]{plump2018modular}, obtained by removing all edges in the interface:
     
    % \begin{figure}[H]
    %     \centering
    \begin{center}
        \resizebox{0.6\textwidth}{!}{
            \begin{tikzpicture}
                \graphbox{$L$}{0mm}{0mm}{35mm}{35mm}{2mm}{-5mm}{
                    \coordinate (delta) at (0,-18mm);
                    \node[draw,circle] (l1) at ($(delta) + (-1,1.5)$) {$1$};
                    \node[draw,circle] (l2) at ($(delta) + (1,1.5)$) {$2$};
                    \node[draw,circle] (l3) at ($(delta) + (0,0)$) {3};
                    \draw[->] (l1) -- (l3) node[midway,left] {$s$};
                    \draw[->] (l2) -- (l3) node[midway,right] {$s$};
                    \draw[->] (l3) edge [loop below] node {$0$} (l3);
                }
                \graphbox{$K$}{40mm}{0mm}{35mm}{35mm}{2mm}{-5mm}{
                    \coordinate (delta) at (0,-18mm);
                    \coordinate (interfaceorigin) at ($(delta) +(5,0)$);
                    \node[draw,circle] (r1) at ($(delta) +(-1,1.5)$) {$1$};
                    \node[draw,circle] (r2) at ($(delta) +(0.5,1.5)$) {$2$};
                    \node[draw,circle] (r3) at ($(delta) + (0,0)$) {3};
                    % \draw[->] (r1) -- (r3) node[midway,left] {$s$};
                    % \draw[->] (r3) edge [loop below] node {$0$} (r3);
                }
                \graphbox{$R$}{80mm}{0mm}{50mm}{35mm}{2mm}{-5mm}{
                    \coordinate (delta) at (-10mm,-18mm);
                    \node[draw,circle] (r1) at ($(delta) + (-1,1.5)$) {$1$};
                    \node[draw,circle] (r2) at ($(delta) + (0.5,1.5)$) {$2$};
                    \node[draw,circle] (r3) at ($(delta) + (0,0)$) {3};
                    \node[draw,circle] (r4) at ($(delta) + (1,0)$) {4};
                    \draw[->] (r1) -- (r3) node[midway,left] {$s$};
                    \draw[->] (r2) -- (r4) node[midway,right] {$s$};
                    \draw[->] (r4) edge [loop below] node {$0$} (r4);
                    \draw[->] (r3) edge [loop below] node {$0$} (r3);
                    \node[draw,circle] (r5) at ($(r2) + (1.5,0)$) {};
                    \draw[->] (r5) edge [loop below] node {$0$} (r5);
                    \draw[->] (r5) edge [loop right] node {$0$} (r5);
                    \draw[->] (r5) edge [loop left] node {$0$} (r5);
                }
                % \graphbox{$R_x$}{40mm}{40mm}{35mm}{35mm}{2mm}{-5mm}{
                %     \coordinate (delta) at (0,-18mm);
                %     \coordinate (rxorigin) at ($(interfaceorigin)+(0,6)$);
                %     \node[draw,circle] (r1) at ($(delta) + (-1,1.5)$) {$1$};
                %     \node[draw,circle] (r2) at ($(delta) +  (0.5,1.5)$) {$2$};
                %     \node[draw,circle] (r3) at ($(delta) +  (0,0)$) {3};
                %     \draw[->] (r1) -- (r3) node[midway,left] {$s$};
                %     % \draw[->] (r3) edge [loop below] node {$0$} (r3);
                % }
                \node () at (38mm,-18mm) {$\leftarrowtail$};
                \node () at (77mm,-18mm) {$\rightarrowtail$};
                % \node () at (57mm,2mm) {$\uparrowtail$};
                % \node () at (38mm,2mm) {$\swarrowtail$};
                % \node () at (79mm,2mm) {$\searrowtail$};
            \end{tikzpicture}
            }
          \end{center}
    %         \caption{}
    %         \label{fig:subgraph_counting:ex_confdkjfakljl}
    % \end{figure}
        %     \caption{Diagram of Example~\ref{ex_contrib_variant}}
        %     \label{fig:contrib_variant}
        % \end{figure}
        Let $X$ be the ruler-graph 
        \tikz[baseline=-0.5ex]{ 
                \node[draw,circle] (x) at (0,0) { }; 
                \node[draw,circle] (y) at (1,0) { };
                \node[draw,circle] (z) at (2,0) { };
                \draw[->] (x) -- (y) node[midway, above] {$s$};
                \draw[->] (z) -- (y) node[midway, above] {$s$};
        }, $\mathbb{X} = \{X\}$ and $s_\mathbb{X}(X)=1$.
        The set \( D(R,X) \) consists of two elements $R'_1$:
        \raisebox{2pt}{
            \scalebox{0.7}{\tikz[baseline=-0.5ex]{
            \node [draw,circle] (x) at (0,0) {$1$};
            \node[draw,circle] (y) at (1,0) {3};
            \draw[->] (x) -- (y) node[midway, above] {$s$};
        }}} and $R'_2$:
        \raisebox{2pt}{
            \scalebox{0.7}{\tikz[baseline=-0.5ex]{
            \node [draw,circle] (z) at (-1,0) {$2$};
            \node [draw,circle] (x) at (0,0) {$1$};
            \node[draw,circle] (y) at (1,0) {3};
            \draw[->] (x) -- (y) node[midway, above] {$s$};
        }}}. 
        Each \( R'_i \) admits a unique monomorphism \( h_{R'_i L} \colon R'_i \rightarrowtail L \) preserving interface elements. 
        Therefore, the rule is $X$-non-increasing under $h_{R'_1L}$ and $h_{R'_2L}$. This is because all conditions in Definition~\ref{subgraph_counting:def:creates_more_x_on_the_left} are satisfied.
        Since \(w_{s_\mathbb{X}}(L) = 1 > 0 = w_{s_\mathbb{X}}(R)\).
        It follows that the system terminates by Theorem~\ref{subgraph_counting:thm:termination_grs}.
\end{example}

The following example demonstrates the usefulness of a weight function that assigns distinct weight to measurements from different ruler-graphs.
\begin{example} 
    \label{ex:overbeek_5d6}
    Consider the rewriting rules presented in \cite[Example 5.6]{overbeek2024termination_lmcs} illustrated 
    below.
    % in Figure~\ref{fig:subgraph_counting:ex_overbdfjkahsdfkjfhew}.
    % \begin{figure}[H]
    %       \centering
    \begin{center}
        \begin{center} 
          $\rho = $\scalebox{0.7}{{
\begin{tikzpicture}[baseline=-10mm]
    \graphbox{$L$}{0mm}{0mm}{31mm}{20mm}{2mm}{-13mm}{
      \node [draw, circle] (x) at (-7mm,0mm) {1};
      \node [draw, circle] (y) at (0mm,0mm) {2};
      \draw[->] (x) edge[loop above] node  {$a$} (x);
      \draw[->] (y) edge [loop above] node {$a$} (y);
    }
    \graphbox{$K$}{32mm}{-0mm}{18mm}{20mm}{0mm}{-10mm}{
    }
    \begin{scope}[opacity=1]        
    \graphbox{$R$}{51mm}{-0mm}{38mm}{20mm}{2mm}{-13mm}{
      \node [draw, circle] (x) at (-7mm,0mm) {3};
      \node [draw, circle] (y) at (0mm,0mm) {4};
      \node [draw, circle] (z) at (7mm,0mm) {5};
      \draw[->] (x) edge[loop above] node  {$b$} (x);
      \draw[->] (y) edge[loop above] node  {$b$} (y);
      \draw[->] (z) edge[loop above] node  {$b$} (z);
    }
    \end{scope}
\end{tikzpicture}
}}
        \end{center}
        \begin{center}
        $\tau = $\scalebox{0.7}{{\begin{tikzpicture}[baseline=-10mm]
    \graphbox{L}{0mm}{0}{31mm}{20mm}{2mm}{-13mm}{
        \node[draw,circle] (x) at (-7mm,0mm) {1};
        \node[draw,circle] (y) at (0mm,0mm) {2};
        \draw[->] (x) edge[loop above] node {$b$} (x);
        \draw[->] (y) edge [loop above] node {$b$} (y);
      }
      \graphbox{K}{32mm}{0}{18mm}{20mm}{2mm}{-13mm}{
      }
      \begin{scope}[opacity=1]        
      \graphbox{R}{51mm}{0}{25mm}{20mm}{1mm}{-13mm}{
        \node[draw,circle] (x) at (-3.5mm,0mm) {3};
        \draw[->] (x) edge[loop above] node {$a$} (x);
      }
      \end{scope}
  \end{tikzpicture}}
        }
        \end{center}
  %       \caption{}
  %       \label{fig:subgraph_counting:ex_overbdfjkahsdfkjfhew}
  % \end{figure}
      \end{center}
     Let $X$ be the ruler-graph
    \tikz[baseline=-0.5ex]{
        \node[draw,circle] (x) at (0,0) { };
        \draw[->] (x) edge [loop right] node {$a$} (x);
    }, $Y$ the ruler-graph 
    \tikz[baseline=-0.5ex]{  
        \node[draw,circle] (x) at (0,0) { };
        \draw[->] (x) edge [loop right] node {$b$} (x);
    } and $\mathbb{X} = \{X, Y\}$.
    The sets $D(R_\rho, X)$, $D(R_\rho, Y)$, $D(R_\tau, X)$ and $D(R_\tau, Y)$ are all empty. Therefore, both rules $\rho$ and $\tau$ are $X$- and $Y$-non-increasing as all conditions in Definition~\ref{subgraph_counting:def:creates_more_x_on_the_left} are met.
    
    For $s_\mathbb{X}(X) = 5$ and $s_\mathbb{X}(Y) = 3$, we have $
    w_{s_\mathbb{X}}(L_\rho) = 10 > 9 = w_{s_\mathbb{X}}(R_\rho)
    $ and $
    w_{s_\mathbb{X}}(L_\tau) = 6 > 5 = w_{s_\mathbb{X}}(R_\tau)$.
    Therefore, the system terminates by Theorem~\ref{subgraph_counting:thm:termination_grs}.

  However, for $s_\mathbb{X}'(X) = 1$ and $s_\mathbb{X}'(Y) = 1$, we have
    \(
        w_{s_\mathbb{X}'}(L_\rho) = 2 \not > 3 = w_{s_\mathbb{X}'}(R_\rho)
    \) which prevents us from applying Theorem~\ref{subgraph_counting:thm:termination_grs} to prove termination of the rewriting system.
\end{example}

% \begin{example}[Limitation] 
%     \label{ex:plump95_4d1}
%    Consider a rule presented in \cite[Example 4.1]{Plump1995}:
 
%     \begin{center}
%         \resizebox{0.7\textwidth}{!}{
%       $\rho = $  \begin{tikzpicture}[baseline=-20mm]
%         \graphbox{$L$}{0mm}{0}{32mm}{40mm}{0}{0}{
%             \node[draw,circle]  (n1) at (0,-6mm) {$1$};
%             \node[draw,circle]   (n2) at (0mm,-26mm) {$2$};
%             \node[draw,circle]   (n3) at (-10mm,-26mm) {3};
%             \node[draw,circle]   (n4) at (10mm,-26mm) {4};
%             \draw[->]  (n2) edge [loop below] node  {$a$} (n2);
%             \draw[->]  (n3) edge [loop below] node  {$b$} (n3);
%             \draw[->]  (n1) to node [right] {$f$} (n2) ;
%             \draw[->]  (n1) to node [right] {$f$}  (n3);
%             \draw[->]  (n1) to node [right] {$f$}  (n4);
%           }
%           \graphbox{$K$}{33mm}{0}{32mm}{40mm}{0}{-0}{
%             \node[draw,circle]  (n1) at (0,-6mm) {$1$};
%             \node[draw,circle]   (n2) at (0mm,-26mm) {$2$};
%             \node[draw,circle]   (n3) at (-10mm,-26mm) {3};
%             \node[draw,circle]   (n4) at (10mm,-26mm) {4};
%             \draw[->]  (n2) edge [loop below] node  {$a$} (n2);
%             \draw[->]  (n3) edge [loop below] node  {$b$} (n3);
%           }
%           \graphbox{$R$}{66mm}{0}{32mm}{40mm}{0}{-0}{
%             \node[draw,circle]  (n1) at (0,-6mm) {$1$};
%             \node[draw,circle]   (n2) at (0mm,-26mm) {$2$};
%             \node[draw,circle]   (n3) at (-10mm,-26mm) {3};
%             \node[draw,circle]   (n4) at (10mm,-26mm) {4};
%             \draw[->]  (n2) edge [loop below] node  {$a$} (n2);
%             \draw[->]  (n3) edge [loop below] node  {$b$} (n3);
%             \draw[->]  (n1) edge [bend left] node [right] {$f$} (n4) ;
%             \draw[->]  (n1) edge [bend right] node [left] {$f$}  (n4);
%             \draw[->]  (n1) to node [right] {$f$}  (n4);
%           }   
%     \end{tikzpicture}
%         }
%   \end{center}
%   \noindent
%   \begin{minipage}{0.7\textwidth}
%     To prove its termination with our method, a ruler-graph $X$ containing an edge labeled by $f$ must be used, because the number of occurrences of every graph containing only edges labeled by $a$ and $b$ does not change by any rewriting step using $\rho$. But in this case, condition \eqref{def:non_increasing:edge_injective} of Definition~\ref{def:creates_more_x_on_the_left} cannot be satisfied.
%   \end{minipage}
%   \begin{minipage}{0.3\textwidth}
%     \hfill
%   \begin{center}
%     \begin{tikzpicture} 
%         \graphbox{}{0}{0}{30mm}{20mm}{-10mm}{-10mm}{
%             \node[draw,circle]  (n1) at (0,0mm) {$1$};
%             % \node[draw,circle]   (n2) at (0mm,-26mm) {$2$};
%             % \node[draw,circle]   (n3) at (-10mm,-26mm) {3};
%             \node[draw,circle]   (n4) at (20mm,0mm) {4};
%             % \draw[->]  (n2) edge [loop below] node  {$a$} (n2);
%             % \draw[->]  (n3) edge [loop below] node  {$b$} (n3);
%             \draw[->]  (n1) edge [bend left] node [above] {$f$} (n4) ;
%             \draw[->]  (n1) edge [bend right] node [below] {$f$}  (n4);
%             \draw[->]  (n1) to node {$f$}  (n4);
%           }  
%     \end{tikzpicture}
%   \end{center}
% \end{minipage}
  
%   \end{example}
% Finally, we present a limitation of our approach.
 

\begin{example}
  \label{ex_endrullis_6d3_endrullis_5d8}
  Consider the following rule from~\cite[Example 6.3]{endrullis2024generalized_arxiv_v2}: 
  %  in Figure~\ref{fig:subgraph_counting:ex_endrullis_6d3_endrullis_5d8}.
  % \begin{figure}[H]
  %     \centering
  \begin{center}
      \resizebox{0.6\textwidth}{!}{
          \begin{tikzpicture}
              \graphbox{$L$}{0mm}{0mm}{35mm}{25mm}{2mm}{-8mm}{
                  \coordinate (delta) at (0,-18mm);
                  \node[draw,circle] (x) at (-6mm,-10mm) {x};
                  \node[draw,circle] (y) at (6mm,-10mm) {y};
                  \node[draw,circle]  (z) at (6mm,0mm) {z};
                  \draw[->]  (x) to (y);
                  \draw[->] (y) to[bend right=20] (z);
                  \draw[->]  (z) to[bend right=20] (y);
              }
              \graphbox{$K$}{40mm}{0mm}{35mm}{25mm}{2mm}{-8mm}{
                  \node[draw,circle]  (x) at (-6mm,-10mm) {x};
                  \node[draw,circle]  (y) at (6mm,-10mm) {y};
              }
              \graphbox{$R$}{80mm}{0mm}{35mm}{25mm}{2mm}{-8mm}{
                  \node[draw,circle]  (x) at (-6mm,-10mm) {x};
                      \node[draw,circle]  (y) at (6mm,-10mm) {y};
                      \node[draw,circle]  (z) at (0mm,0mm) {w};
                      \draw[->]  (x) to (y);
                      \draw[->]  (y) to [bend left=0] (z);
                      \draw[->]  (z) to [bend left=0] (x);
              }
              % \graphbox{$R_x$}{40mm}{20mm}{35mm}{15mm}{2mm}{2mm}{
              %     \node[draw,circle]  (x) at (-6mm,-10mm) {x};
              %     \node[draw,circle]  (y) at (6mm,-10mm) {y};
              %     \draw[->]  (x) to (y);
              % }
              \node () at (37mm,-12mm) {$\leftarrowtail$};
              \node () at (78mm,-12mm) {$\rightarrowtail$};
              % \node () at (57mm,2mm) {$\uparrowtail$};
              % \node () at (38mm,2mm) {$\swarrowtail$};
              % \node () at (79mm,2mm) {$\searrowtail$};
          \end{tikzpicture}
      }
  %     \caption{}
  %     \label{fig:subgraph_counting:ex_endrullis_6d3_endrullis_5d8}
  % \end{figure}
  \end{center}

  Let $X$ be the graph 
  \tikz[baseline=-0.5ex]{
      \node[draw,circle] (x) at (0,0) { };
      \node[draw,circle] (y) at (1,0) { };
      \draw[->] (x) to[bend right=20] (y);
      \draw[->] (y) to[bend right=20] (x); 
  } and $\mathbb{X} = \{X\}$. Let $s_\mathbb{X}$ be the weight function associating the weight of $1$ to $X$. There is only one element 
  \scalebox{0.7}{\tikz[baseline=-0.5ex]{
          \node [draw,circle] (x) at (0,0) {x};
          \node[draw,circle] (y) at (1,0) {y};
          \draw[->] (x) -- (y) node[midway, above] {};
      }}
  in $D(X, R)$. 
   Since $w_{s_\mathbb{X}}(L) = 1 > 0 = w_{s_\mathbb{X}}(R)
      $, the rule terminates by Theorem~\ref{subgraph_counting:thm:termination_grs}.
\end{example}


\begin{example} 
  \label{ex:overbeek_5d5}
  Consider the following rewriting rul from~\cite[Example 5.5]{overbeek2024termination_lmcs}:
  \begin{center}
  % \begin{figure}[H]
  %   \centering
      \resizebox{0.6\textwidth}{!}{
          \begin{tikzpicture}
              \graphbox{$L$}{0mm}{0mm}{35mm}{25mm}{2mm}{-8mm}{
                  \coordinate (delta) at (0,-18mm);
                  \node[draw,circle] (x) at (-6mm,-10mm) {x};
                  \node[draw,circle] (y) at (6mm,-10mm) {y};
                  \node[draw,circle]  (z) at (6mm,0mm) {z};
                  \draw[->]  (x) to (y);
                  \draw[->] (y) to[bend right=20] (z);
                  \draw[->]  (z) to[bend right=20] (y);
              }
              \graphbox{$K$}{40mm}{0mm}{35mm}{25mm}{2mm}{-8mm}{
                  \node[draw,circle]  (x) at (-6mm,-10mm) {x};
                  \node[draw,circle]  (y) at (6mm,-10mm) {y};
                  \draw[->] (x) to (y);
              }
              \graphbox{$R$}{80mm}{0mm}{35mm}{25mm}{2mm}{-8mm}{
                  \node[draw,circle]  (x) at (-6mm,-10mm) {x};
                      \node[draw,circle]  (y) at (6mm,-10mm) {y};
                      \node[draw,circle]  (z) at (0mm,0mm) {w};
                      \draw[->]  (x) to (y);
                      \draw[->]  (y) to [bend left=0] (z);
                      \draw[->]  (z) to [bend left=0] (x);
              }
              % \graphbox{$R_x$}{40mm}{20mm}{35mm}{15mm}{2mm}{2mm}{
              %     \node[draw,circle]  (x) at (-6mm,-10mm) {x};
              %     \node[draw,circle]  (y) at (6mm,-10mm) {y};
              %     \draw[->]  (x) to (y);
              % }
              \node () at (37mm,-12mm) {$\leftarrowtail$};
              \node () at (78mm,-12mm) {$\rightarrowtail$};
              % \node () at (57mm,2mm) {$\uparrowtail$};
              % \node () at (38mm,2mm) {$\swarrowtail$};
              % \node () at (79mm,2mm) {$\searrowtail$};
          \end{tikzpicture}
      }
  %     \caption{}
  %     \label{fig:subgraph_counting:overbeek_5d5}
  % \end{figure}
  \end{center}
    The termination of the system in the DPO rewriting framework with monic matches can be established using our approach; the argument is identical to that of Example~\ref{ex_endrullis_6d3_endrullis_5d8}.
\end{example}

\begin{example}
  \label{ex:plump_ex4}
  Our method cannot prove the termination of the DPO rewriting system with monic matches from~\cite[Example 4]{plump2018modular}:
  %  in Figure~\ref{fig:subgraph_counting:ex_plump_ex4}.
  % \begin{figure}[H]
\begin{center}
  $r_1 = ${ 
   \begin{tikzpicture}[baseline=-20mm]
         \graphbox{$L$}{0mm}{-11mm}{35mm}{15mm}{0.5mm}{-9mm}{
           \node [draw, circle] (x) at (-10mm,0mm) {$1$};
           \node [draw, circle] (y) at (0mm,0mm) {$2$};
           \node [draw, circle] (z) at (10mm,0mm) {3};
           \draw[->] (x) to node [above] {$0$} (y);
           \draw[->] (y) to node [above] {$L$} (z);
         }
         \graphbox{$K$}{36mm}{-11mm}{35mm}{15mm}{0.5mm}{-9mm}{
           \node [draw, circle] (x) at (-10mm,0mm) {$1$};
           %\node [draw, circle] (y) at (0mm,0mm) {$2$};
           \node [draw, circle] (z) at (10mm,0mm) {3};
         }
         \begin{scope}[opacity=1]        
         \graphbox{$R$}{72mm}{-11mm}{40mm}{15mm}{-4mm}{-9mm}{
           \node [draw, circle] (x) at (-10mm,0mm) {$1$};
           \node [draw, circle] (y) at (0mm,0mm) {$2$};
           \node [draw, circle] (z) at (10mm,0mm) {3};
           \node [draw, circle] (w) at (20mm,0mm) {3};
           \draw[->] (x) to node [above] {$L$} (y);
           \draw[->] (y) to node [above] {$1$} (z);
           \draw[->] (z) to node [above] {$1$} (w);
         }
         \end{scope}
       \end{tikzpicture}
  }
 \end{center}
 
 \begin{center}
   $r_2 = ${ 
    \begin{tikzpicture}[baseline=-20mm]
       \graphbox{$L$}{0mm}{-11mm}{35mm}{15mm}{0.5mm}{-9mm}{
           \node [draw, circle] (x) at (-10mm,0mm) {$1$};
           \node [draw, circle] (y) at (0mm,0mm) {$2$};
           \node [draw, circle] (z) at (10mm,0mm) {3};
           \draw[->] (x) to node [above] {$R$} (y);
           \draw[->] (y) to node [above] {$1$} (z);
       }
       \graphbox{$K$}{36mm}{-11mm}{35mm}{15mm}{0.5mm}{-9mm}{
           \node [draw, circle] (x) at (-10mm,0mm) {$1$};
           %\node [draw, circle] (y) at (0mm,0mm) {$2$};
           \node [draw, circle] (z) at (10mm,0mm) {3};
       }
       \begin{scope}[opacity=1]        
       \graphbox{$R$}{72mm}{-11mm}{35mm}{15mm}{0.5mm}{-9mm}{
           \node [draw, circle] (x) at (-10mm,0mm) {$1$};
           \node [draw, circle] (y) at (0mm,0mm) {$2$};
           \node [draw, circle] (z) at (10mm,0mm) {3};
           \draw[->] (x) to node [above] {$0$} (y);
           \draw[->] (y) to node [above] {$R$} (z);
       }
       \end{scope}
       \end{tikzpicture}
   }
   \end{center}

\begin{center}
  $r_3 = ${ 
   \begin{tikzpicture}[baseline=-20mm]
         \graphbox{$L$}{0mm}{-11mm}{35mm}{15mm}{0.5mm}{-9mm}{
           \node [draw, circle] (x) at (-10mm,0mm) {$1$};
           \node [draw, circle] (y) at (0mm,0mm) {$2$};
           \node [draw, circle] (z) at (10mm,0mm) {3};
           \draw[->] (x) to node [above] {$B$} (y);
           \draw[->] (y) to node [above] {$L$} (z);
         }
         \graphbox{$K$}{36mm}{-11mm}{35mm}{15mm}{0.5mm}{-9mm}{
           \node [draw, circle] (x) at (-10mm,0mm) {$1$};
           %\node [draw, circle] (y) at (0mm,0mm) {$2$};
           \node [draw, circle] (z) at (10mm,0mm) {3};
         }
         \begin{scope}[opacity=1]        
         \graphbox{$R$}{72mm}{-11mm}{35mm}{15mm}{0.5mm}{-9mm}{
           \node [draw, circle] (x) at (-10mm,0mm) {$1$};
           \node [draw, circle] (y) at (0mm,0mm) {$2$};
           \draw[->] (x) to node [above] {$R$} (y);
         }
         \end{scope}
       \end{tikzpicture}
  }
 \end{center}
 
 \begin{center}
   $r_4 = ${ 
    \begin{tikzpicture}[baseline=-20mm]
       \graphbox{$L$}{0mm}{-11mm}{35mm}{15mm}{0.5mm}{-9mm}{
           \node [draw, circle] (x) at (-10mm,0mm) {$1$};
           \node [draw, circle] (y) at (0mm,0mm) {$2$};
           \node [draw, circle] (z) at (10mm,0mm) {3};
           \draw[->] (x) to node [above] {$R$} (y);
           \draw[->] (y) to node [above] {$B$} (z);
       }
       \graphbox{$K$}{36mm}{-11mm}{35mm}{15mm}{0.5mm}{-9mm}{
           \node [draw, circle] (x) at (-10mm,0mm) {$1$};
           %\node [draw, circle] (y) at (0mm,0mm) {$2$};
           \node [draw, circle] (z) at (10mm,0mm) {3};
       }
       \begin{scope}[opacity=1]        
       \graphbox{$R$}{72mm}{-11mm}{35mm}{15mm}{0.5mm}{-9mm}{
           \node [draw, circle] (x) at (-10mm,0mm) {$1$};
           \node [draw, circle] (y) at (0mm,0mm) {$2$};
           \node [draw, circle] (z) at (10mm,0mm) {3};
           \draw[->] (x) to node [above] {$L$} (y);
           \draw[->] (y) to node [above] {$B$} (z);
       }
       \end{scope}
       \end{tikzpicture}
   }
   \end{center}
  %  \caption{}
  %  \label{fig:subgraph_counting:ex_plump_ex4}
  % \end{figure}
  The rule $r_3$ can be eliminated by considering occurrences of \tikz[baseline=-0.5ex]{
    \node[draw,circle] (x) at (0,0) { };
    \node[draw,circle] (y) at (1,0) { };
    \draw[->] (x) -- (y) node[midway, above] {$B$};
  }, and the rule $r_4$ can be eliminated by considering occurrences of \tikz[baseline=-0.5ex]{
  \node[draw,circle] (x) at (0,0) { };
  \node[draw,circle] (y) at (1,0) { };
  \draw[->] (x) -- (y) node[midway, above] {$R$};
}.

However, our method cannot eliminate $r_1$ and $r_2$. Yet, combined with the technique proposed by Plump~\cite{plump2018modular}, we can prove the termination of this rewriting system: $\{r_1\}$ and $\{r_2\}$ are both terminating, and their union is terminating~\cite{plump2018modular}.
\end{example}


\begin{example}
  \label{ex:overbeek_5d8_plump1995_3d8_plump2018_3_overbeek_5d8}
  Consider the rewriting rules from~\cite[Example~3.8]{plump1995ontermination}:
  % in Figure~\ref{fig:subgraph_counting:overbeek_5d8_plump1995_3d8_plump2018_3_overbeek_5d8}.

  % \begin{figure}[H]
  %   \centering
\begin{center}
 $\rho = ${ 
  \begin{tikzpicture}[baseline=-20mm]
        \graphbox{$L$}{0mm}{-11mm}{35mm}{15mm}{0.5mm}{-9mm}{
          \node [draw, circle] (x) at (-10mm,0mm) {$1$};
          \node [draw, circle] (y) at (0mm,0mm) {$2$};
          \node [draw, circle] (z) at (10mm,0mm) {3};
          \draw[->] (x) to node [above] {$a$} (y);
          \draw[->] (y) to node [above] {$b$} (z);
        }
        \graphbox{$K$}{36mm}{-11mm}{35mm}{15mm}{0.5mm}{-9mm}{
          \node [draw, circle] (x) at (-10mm,0mm) {$1$};
          %\node [draw, circle] (y) at (0mm,0mm) {$2$};
          \node [draw, circle] (z) at (10mm,0mm) {3};
        }
        \begin{scope}[opacity=1]        
        \graphbox{$R$}{72mm}{-11mm}{35mm}{15mm}{0.5mm}{-9mm}{
          \node [draw, circle] (x) at (-10mm,0mm) {$1$};
          \node [draw, circle] (y) at (0mm,0mm) {$2$};
          \node [draw, circle] (z) at (10mm,0mm) {3};
          \draw[->] (x) to node [above] {$a$} (y);
          \draw[->] (y) to node [above] {$c$} (z);
        }
        \end{scope}
      \end{tikzpicture}
 }
\end{center}

\begin{center}
  $\tau = ${ 
   \begin{tikzpicture}[baseline=-20mm]
      \graphbox{$L$}{0mm}{-11mm}{35mm}{15mm}{0.5mm}{-9mm}{
          \node [draw, circle] (x) at (-10mm,0mm) {$1$};
          \node [draw, circle] (y) at (0mm,0mm) {$2$};
          \node [draw, circle] (z) at (10mm,0mm) {3};
          \draw[->] (x) to node [above] {$c$} (y);
          \draw[->] (y) to node [above] {$d$} (z);
      }
      \graphbox{$K$}{36mm}{-11mm}{35mm}{15mm}{0.5mm}{-9mm}{
          \node [draw, circle] (x) at (-10mm,0mm) {$1$};
          %\node [draw, circle] (y) at (0mm,0mm) {$2$};
          \node [draw, circle] (z) at (10mm,0mm) {3};
      }
      \begin{scope}[opacity=1]        
      \graphbox{$R$}{72mm}{-11mm}{35mm}{15mm}{0.5mm}{-9mm}{
          \node [draw, circle] (x) at (-10mm,0mm) {$1$};
          \node [draw, circle] (y) at (0mm,0mm) {$2$};
          \node [draw, circle] (z) at (10mm,0mm) {3};
          \draw[->] (x) to node [above] {$d$} (y);
          \draw[->] (y) to node [above] {$b$} (z);
      }
      \end{scope}
      \end{tikzpicture}
  }
  \end{center}
%   \caption{}
%   \label{fig:subgraph_counting:overbeek_5d8_plump1995_3d8_plump2018_3_overbeek_5d8}
% \end{figure}
  Let $X$ be  \tikz[baseline=-0.5ex]{
      \node[draw,circle] (x) at (0,0) { };
      \node[draw,circle] (y) at (1,0) { };
      \node[draw,circle] (z) at (2,0) { };
      \draw[->] (x) -- (y) node[midway, above] {$a$};
      \draw[<-] (z) -- (y) node[midway, above] {$b$};
  } 
  We have $D(R_\rho,X) = \emptyset$ and $D(R_\tau,X) = \emptyset$. Therefore, both rules are $X$-non-increasing. 
  We have 
  $|\operatorname{Mono}(X,L)| = 1 > 0 = |\operatorname{Mono}(X,R)|$ for $\rho$, and 
  $|\operatorname{Mono}(X,L)| = 0 = 0 = |\operatorname{Mono}(X,R)|$ for $\tau$, therefore $\rho$ can be eliminated.
  
  Let $Y$ be  \tikz[baseline=-0.5ex]{
      \node[draw,circle] (x) at (0,0) {};
      \node[draw,circle] (y) at (1,0) {};
      \draw[->] (x) -- (y) node[midway, above] {$c$};
  }. We have $D(R_\tau,Y) = \emptyset$. Therefore, $\tau$ is $Y$-non-increasing. Moreover, we have $|\operatorname{Mono}(Y,L)| = 1 > 0 = |\operatorname{Mono}(Y,R)|$.
  Thus, this rewriting system terminates by Theorem~\ref{subgraph_counting:thm:termination_grs}.
\end{example} 


\begin{example}
  \label{ex:overbeek_5d3}
  Consider the following rewriting rules from~\cite[Example 5.3]{overbeek2024termination_lmcs}:
  % in Figure~\ref{fig:subgraph_counting:ex_overbeek_5d3}.
  % \begin{figure}[H]
  %   \centering
\begin{center}
  $\rho = ${ 
  \begin{tikzpicture} 
      \graphbox{$L$}{0mm}{0mm}{23mm}{12mm}{2mm}{-5mm}{
          \node[draw,circle] (x) at (0mm,0mm) {$1$};
          \draw[->] (x) edge [loop right] node {} (x);
      }
      \graphbox{$K$}{24mm}{0mm}{23mm}{12mm}{2mm}{-5mm}{
        \node[draw,circle] (x) at (0mm,0mm) {$1$};
      }
      \begin{scope}[opacity=1]        
      \graphbox{$R$}{48mm}{0mm}{30mm}{12mm}{-1mm}{-5mm}{
        \node[draw,circle] (x) at (0mm,0mm) {$1$};
        \node[draw,circle] (y) at (10mm,0mm) {$2$};
      }
      \end{scope}
    \end{tikzpicture}
  }
  \end{center}

  \begin{center}
  $\tau = $
  \begin{tikzpicture}
      \graphbox{$L$}{0mm}{-11mm}{32mm}{12mm}{2mm}{-8mm}{
      \node[draw,circle]  (x) at (-6mm,0mm) {$1$};
      \node[draw,circle] (y) at (6mm,0mm) {$2$};
      \draw[->]  (x) to (y);
      }
      \graphbox{$K$}{33mm}{-11mm}{32mm}{12mm}{2.5mm}{-8mm}{
          \node[draw,circle]  (x) at (-6mm,0mm) {$1$};
          \node[draw,circle]  (y) at (6mm,0mm) {$2$};
      }
      \begin{scope}[opacity=1]        
      \graphbox{$R$}{66mm}{-11mm}{40mm}{12mm}{-2mm}{-8mm}{
          \node[draw,circle]  (x) at (-6mm,0mm) {$1$};
          \node[draw,circle]  (y) at (6mm,0mm) {$2$};
          \node[draw,circle]  (z) at (17mm,0mm) {3};
      }
      \end{scope}
    \end{tikzpicture}
  \end{center}
%   \caption{}
%   \label{fig:subgraph_counting:ex_overbeek_5d3}
% \end{figure}
  Let $X$ be 
  \tikz[baseline=-0.5ex]{
      \node[draw,circle] (x) at (0,0) { };
      \node[draw,circle] (y) at (1,0) { };
      \draw[->] (x) to (y);
  }. Both rules are $X$-non-increasing, and $|\operatorname{Mono}(X,L)| = 1 > 0 = |\operatorname{Mono}(X,R)|$ holds for both rules. Thus, this rewriting system terminates by Theorem~\ref{subgraph_counting:thm:termination_grs}.
\end{example}

\begin{example} 
  \label{ex:endrullis2024_6d2}  
  Consider the rewriting rule from~\cite[Example 6.2]{endrullis2024generalized_arxiv_v2}:
  %  in Figure~\ref{fig:subgraph_counting:ex_endrullis2024_6d2}.

  % \begin{figure}[H]
  %   \centering
  \begin{center}
  \begin{tikzpicture}  
     \graphbox{$L$}{0mm}{-11mm}{32mm}{15mm}{2mm}{-8mm}{  
         \node[draw,circle]  (x) at (-6mm,0mm) {$1$};  
        \draw[->] (x) edge[loop above] node  {} (x);

         \node[draw,circle]  (y) at (6mm,0mm) {$2$};  
       }  
       \graphbox{$K$}{33mm}{-11mm}{32mm}{15mm}{2mm}{-8mm}{  
        \node[draw,circle]  (x) at (-6mm,0mm) {$1$};  
        \node[draw,circle]  (y) at (6mm,0mm) {$2$};  
       }  
       \graphbox{$R$}{66mm}{-11mm}{32mm}{15mm}{1mm}{-8mm}{  
        \node[draw,circle]  (x) at (-6mm,0mm) {$1$};  
        \node[draw,circle]  (y) at (6mm,0mm) {$2$};  
        \draw[->]  (x) to (y);  
       }    
 \end{tikzpicture}  
\end{center}
%  \caption{}
%  \label{fig:subgraph_counting:ex_endrullis2024_6d2}
% \end{figure}

 Our method can prove its termination with the ruler-graph \tikz[baseline=-0.5ex]{
  \node[draw,circle] (x) at (0,0) {};
  \draw[->] (x) edge[loop above] node  {} (x);
} of weight $1$.
\end{example} 

\begin{example}[Limitation]
  \label{ex:plump95_4d1}
 Consider the rewriting rule from~\cite[Example 4.1]{plump1995ontermination}:
%  in Figure~\ref{fig:subgraph_counting:ex_plump95_4d1}.
  % \begin{figure}[H]
  %   \centering
\begin{center}
      \resizebox{0.7\textwidth}{!}{
    $\rho = $  \begin{tikzpicture}[baseline=-20mm]
      \graphbox{$L$}{0mm}{0}{32mm}{40mm}{0}{0}{
          \node[draw,circle]  (n1) at (0,-6mm) {$1$};
          \node[draw,circle]   (n2) at (0mm,-26mm) {$2$};
          \node[draw,circle]   (n3) at (-10mm,-26mm) {3};
          \node[draw,circle]   (n4) at (10mm,-26mm) {4};
          \draw[->]  (n2) edge [loop below] node  {$a$} (n2);
          \draw[->]  (n3) edge [loop below] node  {$b$} (n3);
          \draw[->]  (n1) to node [right] {$f$} (n2) ;
          \draw[->]  (n1) to node [right] {$f$}  (n3);
          \draw[->]  (n1) to node [right] {$f$}  (n4);
        }
        \graphbox{$K$}{33mm}{0}{32mm}{40mm}{0}{-0}{
          \node[draw,circle]  (n1) at (0,-6mm) {$1$};
          \node[draw,circle]   (n2) at (0mm,-26mm) {$2$};
          \node[draw,circle]   (n3) at (-10mm,-26mm) {3};
          \node[draw,circle]   (n4) at (10mm,-26mm) {4};
          \draw[->]  (n2) edge [loop below] node  {$a$} (n2);
          \draw[->]  (n3) edge [loop below] node  {$b$} (n3);
        }
        \graphbox{$R$}{66mm}{0}{32mm}{40mm}{0}{-0}{
          \node[draw,circle]  (n1) at (0,-6mm) {$1$};
          \node[draw,circle]   (n2) at (0mm,-26mm) {$2$};
          \node[draw,circle]   (n3) at (-10mm,-26mm) {3};
          \node[draw,circle]   (n4) at (10mm,-26mm) {4};
          \draw[->]  (n2) edge [loop below] node  {$a$} (n2);
          \draw[->]  (n3) edge [loop below] node  {$b$} (n3);
          \draw[->]  (n1) edge [bend left] node [right] {$f$} (n4) ;
          \draw[->]  (n1) edge [bend right] node [left] {$f$}  (n4);
          \draw[->]  (n1) to node [right] {$f$}  (n4);
        }   
  \end{tikzpicture}
      }
  %   \caption{}
  %   \label{fig:subgraph_counting:ex_plump95_4d1}
  % \end{figure}
\end{center}
  To prove its termination with our method, a ruler-graph $X$ containing an edge labeled by $f$ must be used, because the number of occurrences of every graph containing only edges labeled by $a$ and $b$ does not change by any rewriting step using $\rho$. But in this case, condition~\eqref{def:non_increasing_rule_img_edges_distinct} of Definition~\ref{subgraph_counting:def:creates_more_x_on_the_left} cannot be satisfied.
\todo{to do todo: more details}
%   \begin{center}
%   \begin{tikzpicture} 
%       \graphbox{}{0}{0}{30mm}{20mm}{-10mm}{-10mm}{
%           \node[draw,circle]  (n1) at (0,0mm) {$1$};
%           % \node[draw,circle]   (n2) at (0mm,-26mm) {$2$};
%           % \node[draw,circle]   (n3) at (-10mm,-26mm) {3};
%           \node[draw,circle]   (n4) at (20mm,0mm) {4};
%           % \draw[->]  (n2) edge [loop below] node  {$a$} (n2);
%           % \draw[->]  (n3) edge [loop below] node  {$b$} (n3);
%           \draw[->]  (n1) edge [bend left] node [above] {$f$} (n4) ;
%           \draw[->]  (n1) edge [bend right] node [below] {$f$}  (n4);
%           \draw[->]  (n1) to node {$f$}  (n4);
%         }  
%   \end{tikzpicture}
% \end{center}
\end{example}

\begin{example}[Limitation]
    \label{ex:endrullis:d3:limitation}  
    Our method fails to prove termination of the following rewriting rule from \cite[Example D.3]{endrullis2024generalized_arxiv_v2}:
    % \begin{figure}[H]
    %   \centering
    \begin{center}
        \resizebox{0.7\textwidth}{!}{
    \begin{tikzpicture}  
       \graphbox{$L$}{0mm}{-11mm}{32mm}{15mm}{2mm}{-8mm}{  
           \node[draw,circle]  (x) at (-6mm,0mm) {$1$};  
           \node[draw,circle]  (y) at (6mm,0mm) {};  
         }  
         \graphbox{$K$}{33mm}{-11mm}{32mm}{15mm}{2mm}{-8mm}{  
           \node[draw,circle]  (x) at (-6mm,0mm) {$1$};  
         }  
         \graphbox{$R$}{66mm}{-11mm}{32mm}{15mm}{1mm}{-8mm}{  
          \node[draw,circle]  (x) at (-6mm,0mm) {$1$};  
          \node[draw,circle]  (y) at (6mm,0mm) {};  
          \draw[->]  (x) to (y);  
         }    
   \end{tikzpicture}
        }  
    %     \caption{}
    %     \label{fig:subgraph_counting:ex_endrullis_d3_limitation}
    % \end{figure}
  \end{center}
    %  in Figure~\ref{fig:subgraph_counting:ex_endrullis_d3_limitation}.
    % because our approach cannot take into account antipatterns(see~\cite[Remark 6.2]{overbeek2024termination_lmcs}).
    This happens because, for any system $\mathcal{R}$, rewriting step $G \Rightarrow_{\mathcal{R},\mathfrak{M}} H$, set of ruler-graphs $\mathbb{X}$, and weight function $s_\mathbb{X}$, we have $w_{s_\mathbb{X}}(G) \geq w_{s_\mathbb{X}}(H)$ because there is an injective graph homomorphism from $G$ to $H$.
  \end{example} 

\begin{example}[Limitation]
    \label{ex:overbeek:5d2:limitation}
    Consider a rewriting rule given by Overbeek and Endrullis in~\cite[Example 5.2]{overbeek2024termination_lmcs} shown below.
    % in Figure~\ref{fig:subgraph_counting:ex_overbeek_5d2}.
  % \begin{figure}[H]
  %   \centering
  \begin{center} 
    \begin{tikzpicture}
        \graphbox{$L$}{0mm}{-11mm}{32mm}{15mm}{2mm}{-8mm}{
            \node[draw,circle]  (x) at (-6mm,0mm) {$1$};
            \node[draw,circle]   (y) at (6mm,0mm) {$2$};
            \draw[->]  (x) to (y);
          }
          \graphbox{$K$}{33mm}{-11mm}{32mm}{15mm}{2mm}{-8mm}{
            \node[draw,circle]  (x) at (-6mm,0mm) {$1$};
            \node[draw,circle]  (y) at (6mm,0mm) {$2$};
            \draw[->]  (x) to (y);
          }
          \graphbox{$R$}{66mm}{-11mm}{32mm}{15mm}{1mm}{-8mm}{
            \node[draw,circle]  (x) at (0mm,0mm) {1 2};
            \draw[->]  (x) edge [loop right] (x);
          }  
    \end{tikzpicture} 
  %   \caption{}
  %   \label{fig:subgraph_counting:ex_overbeek_5d2}
  % \end{figure}
  \end{center}
    The termination of this rule is not in the scope of our method due to its non-node-injective right-hand side morphism. It can be proved using the method proposed by Overbeek et al.~\cite{overbeek2024termination_lmcs}.
\end{example}
 
\begin{example}[Limitation]
  \label{ex:overbeek:5d6_bis:limitation} 
  Consider the rewriting rule with monic matches from~\cite[Example 5.6]{overbeek2024termination_lmcs}:
  \begin{center}
    $\rho' = $\scalebox{0.9} { {
\begin{tikzpicture}[baseline=-10mm]
    \graphbox{$L$}{0mm}{0mm}{38mm}{20mm}{2mm}{-13mm}{
      \node [draw, circle] (x) at (-7mm,0mm) {1};
      \node [draw, circle] (y) at (3mm,0mm) {2 3};
      \draw[->] (x) edge[loop above] node  {$a$} (x);
      \draw[->] (y) edge [loop above] node {$a$} (y);
    }
    \graphbox{$K$}{42mm}{-0mm}{38mm}{20mm}{0mm}{-10mm}{
        \node [draw, circle] (x) at (-7mm,0mm) {1};
        \node [draw, circle] (y) at (0mm,0mm) {2};
        \node [draw, circle] (z) at (7mm,0mm) {3};    
    }
    \begin{scope}[opacity=1]        
    \graphbox{$R$}{85mm}{-0mm}{38mm}{20mm}{2mm}{-13mm}{
      \node [draw, circle] (x) at (-7mm,0mm) {1};
      \node [draw, circle] (y) at (0mm,0mm) {2};
      \node [draw, circle] (z) at (7mm,0mm) {3};
      \draw[->] (x) edge[loop above] node  {$b$} (x);
      \draw[->] (y) edge[loop above] node  {$b$} (y);
      \draw[->] (z) edge[loop above] node  {$b$} (z);
    }
    \end{scope}
    \node () at (40mm,-10mm) {$\leftarrow$};
    \node () at (83mm,-10mm) {$\rightarrow$};
\end{tikzpicture}
}

}
    % \caption{}
    % \label{fig:subgraph_counting:ex_overbeek_5d6_bis}
  \end{center} 
  Its termination is not in the scope of our method due to its non-node-injective left-hand side morphism. 
\end{example}

  
\section{Related work}
\label{subgraph_counting:sec:related_work} 
A comparative analysis of termination techniques for DPO graph rewriting systems, drawn from prior work~\cite{plump1995ontermination,plump2018modular,bruggink2014termination,bruggink2015proving,endrullis2024generalized_arxiv_v2,overbeek2024termination_lmcs}, is summarized in~\autoref{tab:comparison:subgraph_counting}. Our approach successfully proves termination for 14 of these systems. In the remainder of this section, we compare our method with some existing methods in more detail.


{% local group
  \setlength{\tabcolsep}{3pt}
  \renewcommand{\arraystretch}{1}
\begin{table}[H]
   \centering
  \small % Reduce font size
  \caption{Applicability of termination techniques to DPO rewriting examples.
  The symbol \ding{51} indicates termination can be proved by the technique,
  \ding{55} indicates it cannot be proved, and 
  $-$ denotes irrelevance or out-of-scope cases.
        }
 \label{tab:comparison:subgraph_counting}
% \setlength{\tabcolsep}{4pt} % Reduce horizontal padding
   \begin{NiceTabular}{ccccccccc}[vlines] % <-- 9 columns now (was 7)
    \Hline
                % \diagbox{\enskip \textbf{Examples}}{\textbf{Techniques}} 
    \Block{1-2}{\diagbox{\enskip \textbf{Examples}}{\textbf{Techniques}}} & 
    &
    \RowStyle{\rotate}
     \makecell{Forward closure \cite{plump1995ontermination}} % NEW column #1
    & \RowStyle{\rotate}
     \makecell{Modular criterion \cite{plump2018modular}} % NEW column #2
    & \RowStyle{\rotate}
     \makecell{Type graph \cite{bruggink2014termination}}  
    & \RowStyle{\rotate}
     \makecell{Type graph \cite{bruggink2015proving}} 
    & \RowStyle{\rotate}
     \makecell{Type graph \cite{endrullis2024generalized_arxiv_v2}} 
    & \RowStyle{\rotate}
     \makecell{Subgraph 
            counting \cite{overbeek2024termination_lmcs}} 
    & \RowStyle{\rotate}
     \makecell{Morphism
                Counting 
                \\(Chapter~\ref{chap:subgraph_counting})} \\
    \Hline
    \Hline
% ok examples
Chapter~\ref{chap:subgraph_counting} 
&~\autoref{subgraph_counting:ex_contrib_variant}
  & \ding{51} & \ding{55} & \ding{55} & \ding{55} & \ding{55} & \ding{55} & \ding{51} \\ \Hline
  
\cite{plump1995ontermination} &
\hyperref[ex:overbeek_5d8_plump1995_3d8_plump2018_3_overbeek_5d8]{Example 3.8}
             & \ding{51} & -- & -- & -- & -- &
          --
             & \ding{51}\\ 
\hline

\Block{2-1}{\cite{plump2018modular}} & \hyperref[ex:overbeek_5d8_plump1995_3d8_plump2018_3_overbeek_5d8]{Example 3} 
          & -- & \ding{51} &  -- & -- & -- & 
          --
          & \ding{51}\\ 

\Hline

%~\cite{plump2018modular} 
& Example 5 &  -- &  \ding{51} &   -- & -- & -- &  
            --
          & \ding{51}\\  
\Hline

\cite{bruggink2014termination} & \hyperref[ex:termination:grsaa]{Example 4 and 6}  
  & -- & -- & \ding{51} & -- & -- & 
            --
          & \ding{51} \\ \Hline

\Block{2-1}{\cite{bruggink2015proving}} & \hyperref[ex:termination:grsaa]{Example 2}  
  & -- & -- & -- & \ding{51} & -- & 
  -- &  \ding{51}\\ \Hline
  
%~\cite{bruggink2015proving}
 & Example 4 
  & -- & -- & -- & \ding{51} & -- & 
  --& \ding{51} \\ \Hline


% ----- from endrullis2024generalized_arxiv_v2 -----
\Block{3-1}{\cite{endrullis2024generalized_arxiv_v2}} & \hyperref[ex:endrullis2024_6d2]{Example 6.2}  
  & -- & -- & -- & -- & \ding{51} & -- & \ding{51}\\ \Hline

%~\cite{endrullis2024generalized_arxiv_v2}
 & \hyperref[ex_endrullis_6d3_endrullis_5d8]{Example 6.3}
  & -- & -- & -- & -- & \ding{51} &% \ding{55} 
  \ding{55} & \ding{51}\\ \Hline

%~\cite{endrullis2024generalized_arxiv_v2} 
& \hyperref[ex:overbeek_5d8_plump1995_3d8_plump2018_3_overbeek_5d8]{Example D.1}
  & -- & -- & -- & -- & \ding{51} & -- & \ding{51}\\ \Hline

  \Block{4-1}{\cite{overbeek2024termination_lmcs}} & \hyperref[ex:overbeek_5d3]{Example 5.3}
  & -- & -- & -- & -- & -- & \ding{51} & \ding{51}\\ \Hline

%~\cite{overbeek2024termination_lmcs} \hyperref[ex:overbeek_5d3]{Example 5.3 monic matches}
%     & -- & -- & -- & -- & -- & \ding{51} & \ding{51}\\ \Hline
%~\cite{overbeek2024termination_lmcs} 
& \hyperref[ex:overbeek_5d5]{Example 5.5} 
  & -- & -- & -- & -- & -- & \ding{51} & \ding{51}\\ \Hline

%~\cite{overbeek2024termination_lmcs} 
& \hyperref[ex:overbeek_5d6]{Example 5.6}
  & -- & -- & -- & -- & -- & \ding{51} & \ding{51} \\ \Hline

%~\cite{overbeek2024termination_lmcs} \hyperref[ex:overbeek_5d6]{Example 5.6 bis}
%     & -- & -- & -- & -- & -- & \ding{51} & -- \\ \Hline


%~\cite{overbeek2024termination_lmcs}  
& \hyperref[ex:overbeek_5d8_plump1995_3d8_plump2018_3_overbeek_5d8]{Example 5.8}
  & -- & -- & -- & -- & -- & \ding{51} & \ding{51}\\ \Hline

      % not supported examples  
     ~\cite{plump2018modular} &  \hyperref[ex:plump2018_ex6_endrullis_d4]{Example 6} &  -- & \ding{51} & -- & -- & -- & 
      --
          & -- \\
      \Hline

\Block{3-1}{\cite{endrullis2024generalized_arxiv_v2}}
 & Example 6.4  
      & -- & -- & -- & -- & \ding{51} & -- & -- \\ \Hline

%  ~\cite{endrullis2024generalized_arxiv_v2}
  &  Example 6.5  
      & -- & -- & -- & -- &  \ding{51} & -- & -- \\ \Hline

    %  ~\cite{endrullis2024generalized_arxiv_v2}
       &\hyperref[ex:plump2018_ex6_endrullis_d4]{Example D.4} 
      & -- & -- & -- & -- & \ding{51} & -- & --\\ \Hline

   % ----- from overbeek2024termination_lmcs -----
   \Block{3-1}{\cite{overbeek2024termination_lmcs}} & Example 5.2
      & -- & -- & -- & -- & -- & \ding{51} & -- \\ \Hline

    %  ~\cite{overbeek2024termination_lmcs} 
      & Example 5.7 
      & -- & -- & -- & -- & -- & \ding{51} & -- \\ \Hline
      
%  ~\cite{overbeek2024termination_lmcs} 
  & Example 5.9 
      & -- & -- & -- & -- & -- & \ding{51} & --\\ \Hline
 

    % not ok examples
   ~\cite{plump1995ontermination} & \hyperref[ex:plump95_4d1]{Example 4.1} & \ding{51} & -- & -- & -- & -- & 
              \ding{51}
              
              & \ding{55}\\ 
   \Hline
  ~\cite{plump2018modular} & \hyperref[ex:plump_ex4]{Example 4} &  -- &  \ding{51} &  -- & -- & -- & 
               --
               & \ding{55}\\ 
   \Hline

   \Block{3-1}{\cite{bruggink2014termination}} & Example 1 
   & -- & -- & \ding{51} & -- & -- & 
                 --
               &  \ding{55}\\ 
   \Hline

%   ~\cite{bruggink2014termination} 
   & Routing Protocol
       & -- & -- & \ding{51} & -- & -- & 
           --
           &  \ding{55}\\  
           \Hline
% \Block{2-1}{\cite{bruggink2014termination}}
 & \hyperref[ex:plump_ex4]{Example 5}
   & -- & -- & \ding{51} & -- & -- & -- &  \ding{55}\\ 
\Hline

\Block{2-1}{\cite{bruggink2015proving}} & \hyperref[ex:bruggink2015_ex5]{Example 5}
   & -- & -- & -- & \ding{51} & -- &  
   -- &  \ding{55}\\
   \Hline

%   ~\cite{bruggink2015proving} 
   & \hyperref[ex:bruggink2015_ex6_endrullis2024_d2]{Example 6} 
   & -- & -- & -- & \ding{51} & -- &  
   --&  \ding{55}\\  
   \Hline

   \Block{2-1}{\cite{endrullis2024generalized_arxiv_v2}} & \hyperref[ex:bruggink2015_ex6_endrullis2024_d2]{Example D.2} 
   & -- & -- & -- & -- & \ding{51} & -- & \ding{55}\\ 
   \Hline

%   ~\cite{endrullis2024generalized_arxiv_v2}
   & \hyperref[rem:d3_limitation]{Example D.3}
   & -- & -- & -- & -- & \ding{51} & \ding{55} & \ding{55}\\ \Hline

  \end{NiceTabular}
  % \vspace{2mm}
  % \caption{Applicability of termination techniques to DPO rewriting examples.
  %  The symbol \ding{51} indicates termination can be proved by the technique,
  %  \ding{55} indicates it cannot be proved, and 
  %  $-$ denotes irrelevance or out-of-scope cases.
  %        }
  % \label{tab:comparison}
  \end{table}
} 


The subgraph-counting method in~\cite{overbeek2024termination_lmcs} by Overbeek and Endrullis is designed for the PBPO+~\cite{overbeek2023graph, overbeek2023apbpotutorial}—a rewriting formalism capable of simulating left-injective DPO rewriting. It can prove termination for systems like~\cite[Examples 5.2, 5.7, 5.9]{overbeek2024termination_lmcs} and~\cite[Example 6]{plump2018modular}, which lie beyond the scope of our technique. Additionally, this method can be applied for many categories. 

It is very closely related to our work in the setting where they are both defined. Both methods weigh objects by summing weighted morphisms targeting them, and, for both methods, the key challenge lies in estimating weights for morphisms whose images partially overlap with the rewriting context, because morphisms fully embedded in the context are shared between host and resulting graphs, while those entirely within left- or right-hand-side graphs are trivial to quantify.  

The two methods employ distinct strategies to overcome this challenge. The subgraph-counting method relies on type-morphisms (see~\cite[page 9, remark 4.11, Lemma 4.23]{overbeek2024termination_lmcs}), whereas our approach requires injective mappings from (i) subgraph occurrences partially overlapping the rewriting context in the result graph to (ii) those in the host graph.

Both approaches have limitations and strengths. Their method suffers from discrete interface, as mentioned in~\cite[Example 5.5]{overbeek2024termination_lmcs}. For instance, it proves termination for~\cite[Example 5.5]{overbeek2024termination_lmcs} but fails for~\cite[Example 6.3]{endrullis2024generalized_arxiv_v2}, where the rules differ only by a discrete interface. 
Similarly, it fails to prove termination for~\autoref{subgraph_counting:ex_contrib_variant}, but if an edge labeled by ``s'' from node $1$ to node $3$ is added, then it succeeds. Our approach, however, can handle all the aforementioned cases. Nevertheless, their method can address~\cite[Example 4.1]{plump1995ontermination}, which remains beyond the scope of our technique. Finally, both approaches cannot handle~\cite[Example D.3]{endrullis2024generalized_arxiv_v2} (see~\cite[Remark 6.2]{overbeek2024termination_lmcs} and~\autoref{ex:endrullis:d3:limitation}).

The type graph method, which weighs an object by summing the weights of morphisms from the object to a type graph, was initially introduced by Zantema, K{\"o}nig and Bruggink~\cite{zantema2014termination} for cycle-rewriting systems. 
This method has since been generalized to edge-labeled multigraphs by Bruggink et al.~\cite{bruggink2014termination} for DPO rewriting with monic matches and injective rules, later extended to DPO rewriting in general by Bruggink et al.~\cite{bruggink2015proving}, and further adapted to more categories and different DPO variants by Endrullis et al.~\cite{endrullis2024generalized_arxiv_v2}. 
These methods are not directly comparable with our technique in general.

On the one hand, the termination of~\autoref{subgraph_counting:ex_contrib_variant} can be proved by our method but not by the type graph methods due to the exsitence of a surjection from the output graph to the input graph as explained in~\cite[Example D.4]{endrullis2024generalized_arxiv_v2}. On the other hand, our method cannot prove the termination of~\cite[Example 1, 5 and Ad-hoc Routing Protocol]{bruggink2014termination}, nor~\cite[Example 5, 6]{bruggink2015proving}, nor~\cite[Examples D2 and D3]{endrullis2024generalized_arxiv_v2}.

Plump~\cite{plump1995ontermination} introduced a necessary and sufficient termination condition for left-injective DPO hypergraph rewriting via forward closure, though verifying this condition is undecidable. 
While this method proves termination of~\autoref{subgraph_counting:ex_contrib_variant}, our approach succeeds on~\cite[Example 3.8]{plump1995ontermination} but fails in proving termination of~\cite[Example 4.1]{plump1995ontermination}. 

Plump~\cite{plump2018modular} later proposed a modular critical pair-based strategy for left-injective DPO hypergraph rewriting with monic matches. 
Our method complements this: while modularity reduces global complexity, each subsystem requires individual termination proofs. For example, the measure based on the indegree proposed in~\cite{plump2018modular} cannot prove the termination of~\autoref{subgraph_counting:ex_contrib_variant} due to the additional loops. Specifically, if the context graph is identical to the interface graph, then we have $3^k < 2^k+2^k+3^k$ for all $k \in \mathbb{N}$. In contrast, our method succeeds.
Additionally, ur method can proe termination for~\cite[Examples 1 and 5]{plump2018modular} but not~\cite[Example 4]{plump2018modular}, and non-injective rules (e.g.,~\cite[Example 6]{plump2018modular}) fall outside our scope.

In~\cite{levendovszky2007termination}, a terminaiton criterion for DPO rewriting with monic matches, injective rules and negative application conditions on finite typed attributed graphs is proposed. It is based on the fact that if the application of every infinite sequence of rules requires the initial graph to be infinite, then the system terminates. This technique is theoretically very interesting, but it is hard to automatically check the termination condition.

~\cite{bottoni2005termination} presents a termination criterion for DPO and SPO rewriting on high-level replacement units. The high-level replacement units that they consider are rewriting systems with very restrictive external control mechanisms. The method relies on a mesearing function satisfying a very strong constrainte, and the instance of such a mesearing function proposed is node and edge counting, which is subsumed by our method in the setting of DPO rewriting with injective rules.

~\cite{bottoni2010atermination} presents a criterion for termination of DPO rewriting with monic matches, injective rules and negative application conditions, based on the construction of a labeled transition system.
whose states represent overlaps between the negative application condition and the right hand side that can give rise to cycles.

Levendovszky et al.~\cite{levendovszky2007termination} propose a termination criterion for DPO rewriting (monic matches, injective rules, negative application condition), though automated verification is hard as explained in~\cite[\textsection 6]{levendovszky2007termination}. 

Bottoni et al.~\cite{bottoni2005termination} propose a termination criterion for DPO/SPO rewriting on high-level replacement units. Their method imposes a strongly constrained measuring function and the only concretes measuring function proposed are node-counting and edge-counting.

Bottoni et al.~\cite{bottoni2010atermination} presents a criterion for termination of DPO rewriting with monic matches, injective rules and negative application conditions, based on the construction of a labeled transition system. 



\section{Conclusion and future work}
\label{subgraph_counting:sec:conclusion} 
We have presented a machine-checkable sufficient condition for relative termination of DPO systems with injective rules on edge-labeled multigraphs.  
This method resolves cases where prior interpretation-based techniques \cite{zantema2014termination,bruggink2014termination,bruggink2015proving,endrullis2024generalized_arxiv_v2,
overbeek2024termination_lmcs} fail (e.g.,Example~\ref{subgraph_counting:ex_contrib_variant}) and, unlike the subgraph-counting approach in \cite{overbeek2024termination_lmcs}, avoids limitations arising from interfaces with no edges.
The method has been implemented in an OCaml tool. 

An interesting direction for future work is to extend our method to DPO systems with negative application conditions and to address the limitation highlighted in Example~\ref{ex:endrullis:d3:limitation}.

\section{Appendix}
\subsection*{Adhesive Category and VK-square}
\begin{proposition}[\text{\cite[Lemma 13]{lack2004adhesive}}]
    \label{prop:pb_eq_po}
    In the category \textbf{Graph}, pushouts along monomorphisms are also pullbacks. 
\end{proposition}

\begin{definition}[Van Kampen Square \text{\cite[Definition 2]{lack2004adhesive}}]
    \label{def:vk_square}
    A \textit{van Kampen (VK) square} (i) is a pushout square which satisfies the following condition: given a commutative cube (ii) of which (i) forms the bottom face and the back faces are pullbacks, the front faces are pullbacks if and only if the top face is a pushout.

    
    \begin{enumerate}
        \item[(i)]
        \begin{tikzcd}
            & C \arrow[ld, "m"'] \arrow[rd, "f"]  &\\
            A \arrow[rd,"g"'] & & B \arrow[ld,"n"]\\
            &D &  
        \end{tikzcd}
    
        \item[(ii)]
        \begin{tikzcd}
            A' \arrow[dd, "a" description, near start] \arrow[dr, "g'" description] & & C' \arrow[ll, "m'" description, near start]  \arrow[dd, "c" description, near start] \arrow[dr, "f'" description] & \\
            & D'  \arrow[dd, "d" description, near start] & & B' \arrow[ll, "n'" description, near start] \arrow[dd, "b" description, near start] \\
            A  \arrow[dr, "g" description] & & C \arrow[dr, "f" description] \arrow[ll, "m" description, near start] & \\
            & D  & & B \arrow[ll, "n" description]
        \end{tikzcd}
    \end{enumerate}
\end{definition} 

\begin{definition}[Adhesive category\text{\cite[Definition 5]{lack2004adhesive}}] 
    \label{def:adhesive_cat}
    A category $\mathcal{C}$ is said to be \textbf{adhesive} if 
    \begin{itemize}
        \item $\mathcal{C}$ has pushouts along  monomorphisms, and
        \item $\mathcal{C}$ has pullbacks, and
        \item pushouts along monomorphisms are VK squares.
    \end{itemize}
\end{definition}

\begin{proposition}[\text{\cite[Proposition 8]{lack2004adhesive}}]
    \label{prop:graph_adhesive}
    The category \textbf{Graph} of directed labeled multigraphs is adhesive.
\end{proposition}
\subsection*{Auxiliary lemmas}
\begin{lemma}
    \label{lem:b_c_same_img_exist_a}
    Consider the pushout square in $\mathbf{Graph}$ of form
    \begin{center} 
        \begin{tikzpicture}[node distance=12mm]
          \node (A) at (0,0) {$A$};
          \node (B) at (2,0) {$B$}; 
          \node (C) at (0,-2) {$C$}; 
          \node (D) at (2,-2) {$D$}; 
          \node () [at=($(A)!0.5!(D)$)] {\normalfont PO};
          \draw [>->] (A) to (B);
          \draw [>->] (A) to (C);
          \draw [>->] (B) to (D); 
          \draw [>->] (C) to (D);
        \end{tikzpicture}
    \end{center} 
    The following statements hold:
    \begin{enumerate}
        \item  If $ b\in B $ and $ c \mathop{\in} C $ such that $ h_{BD}(b) \mathop{=} h_{CD}(c) $ then there is $ a \mathop{\in} A $ such that $ h_{AB}(a) \mathop{=} b $ and $ h_{AC}(a)= c $;
        \item  For all \( b \mathop{\in} B \mathop{\setminus} h_{AB}(A) \) and \( c \mathop{\in} C \), we have \( \beta'(b) \mathop{\neq} \alpha'(c) \).
    \end{enumerate}
\end{lemma}
\begin{proof}
     Let \( b \mathop{\in} B \) and \( c \mathop{\in} C \) such that \( h_{BD}(b) \mathop{=} h_{CD}(c) \). 
     
     There are two cases: if $b$ and $c$ are two nodes then let \( S \) be the singleton graph with a unique node \( u \), otherwise\textemdash$b$ and $c$ must be two edges with the same label $l$\textemdash let \( S \) be the graph with two distinct nodes and a unique edge \( u \) labeled by $l$.
     
     Define morphisms \( h_{SB}: S \mathop{\to} B \) and \( h_{SC}: S \mathop{\to} C \) such that \( h_{SB}(u) \mathop{=} b \) and \( h_{SC}(u) \mathop{=} c \). We have \( h_{SB} \mathop{\star} h_{BD} \mathop{=} h_{SC} \mathop{\star} h_{CD} \). The pushout square \( ACDB \) along monomorphisms is also a pullback square by Proposition~\ref{prop:pb_eq_po}, and the universal property of the pullback guarantees the existence of a unique morphism \( h_{SA}: S \mathop{\to} A \) such that \( h_{SA} \mathop{\star} h_{AB} \mathop{=} h_{SB} \) and \( h_{SA} \mathop{\star} h_{AC} \mathop{=} h_{SC} \). Let \( a \mathop{=} h_{SA}(u) \). We have \( a \mathop{\in} A \), and by composition we obtain: 
     \begin{itemize}
        \item \(h_{AB}(a) \mathop{=} (h_{SA} \mathop{\star} h_{AB})(u) \mathop{=} h_{SB}(u) \mathop{=} b \), and
        \item \(h_{AC}(a) \mathop{=} (h_{SA} \mathop{\star} h_{AC})(u) \mathop{=} h_{SC}(u) \mathop{=} c\).
     \end{itemize} 

    The second assertion is a direct consequentce of the first. Suppose, for contradiction, there are \( b \mathop{\in} B \mathop{\setminus} h_{AB}(A) \) and \( c \mathop{\in} C \) such that \( h_{BD}(b) \mathop{=} h_{CD}(c) \). Since the first assertion holds, there must exist \( a \mathop{\in} A \) with \( h_{AB}(a) \mathop{=} b \) and \( h_{AC}(a) \mathop{=} c \). However, this contradicts the assumption that \( b \notin h_{AB}(A) \).
    \qed
\end{proof}

\begin{lemma}
    \label{kpcpck_pullback}
    Consider the following commutative diagram in \textbf{Graph}.
    \begin{center}
        % \resizebox{6cm}{!}{ 
            \begin{tikzpicture}[scale=1.2]
                \node (k) at (0,0) {K};
                \node (r) at (4,0) {R};
                \node (c) at (0,-3) {C};
                \node (h) at (4,-3) {H};
                \draw[<-<]  (r) -- (k) node [midway,above] {};
                \draw[>->] (c) -- (h) node [midway, below] {};
                \draw[>->] (r) -- (h) node[midway, left] {};
                \draw[>->] (k) -- (c) node[midway, left] {};
                \node (k') at (1,-1) {K'};
                \node (r') at (2,-1) {R'};
                \node (c') at (1,-2) {C'};
                \node () at (1.5,-1.5) {$PB$};
                \node () at (3,-1.5) {$PB$};
                \node () at (1.5,-2.5) {$PB$};
                \node (x) at (2,-2) {X};
                \draw [>->] (c') -- (x);
                \draw [>->] (r') -- (x);
                \draw [>->] (k') -- (r');
                \draw [>->] (k') -- (c');
                \draw [>->] (c') -- (c);
                \draw[>->] (r') -- (r);
                \draw[>->] (x) -- (h);
                \draw[->] (k') -- (k);
            \end{tikzpicture}
        % }
        \end{center}  
        The squares $K'KCC'$ and $K'KRR'$ are pulllback square.
\end{lemma}
\begin{proof}
    Let $K \mathop{\leftarrow} A \mathop{\rightarrow} C'$ be a span such that 
    \begin{flalign}
        h_{AK} \mathop{\star} h_{KC} \mathop{=} h_{AC'} \mathop{\star} h_{C'C}. \label{hakhkchacp}
    \end{flalign}
    We have the following commutative diagram 

    \begin{center}
        % \resizebox{0.3\textwidth}{!}{
            \begin{tikzpicture}[scale=1.2]
                \node (k) at (0,0) {K};
                \node (r) at (4,0) {R};
                \node (c) at (0,-3) {C};
                \node (h) at (4,-3) {H};
                \node (h) at (4,-3) {H};
                \node (a) at (-1,0) {A};
                \draw[<-<]  (r) -- (k) node [midway,above] {};
                \draw[>->] (c) -- (h) node [midway, below] {};
                \draw[>->] (r) -- (h) node[midway, left] {};
                \draw[>->] (k) -- (c) node[midway, left] {};
                \node (k') at (1,-1) {K'};
                \node (r') at (2,-1) {R'};
                \node (c') at (1,-2) {C'};
                \node () at (1.5,-1.5) {$PB$};
                \node () at (3,-1.5) {$PB$};
                \node () at (1.5,-2.5) {$PB$};
                \node (x) at (2,-2) {X};
                \draw [>->] (c') -- (x);
                \draw [>->] (r') -- (x);
                \draw [>->] (k') -- (r');
                \draw [>->] (k') -- (c');
                \draw [>->] (c') -- (c);
                \draw[>->] (r') -- (r);
                \draw[>->] (x) -- (h);
                \draw[->] (k') -- (k);
                \draw[->,red] (a) -- (c');
                \draw[->,red] (a) -- (k);
            \end{tikzpicture}
        % }
        \end{center} 
    We show first the existence of a morphism $h_{AK'} \mathop{\colon} A \mathop{\to} K'$ such that the following statements hold:
    \begin{flalign*}
        h_{AK'} \mathop{\star} h_{K'C'} \mathop{=} h_{AC'},  
        \\
        h_{AK'} \mathop{\star} h_{K'K} \mathop{=} h_{AK}. 
    \end{flalign*}
    From Equation~\eqref{hakhkchacp}, we obtain 
    \begin{flalign}
        h_{AK} \mathop{\star} h_{KC} \mathop{\star} h_{CH}= h_{AC'} \mathop{\star} h_{C'C} \mathop{\star} h_{CH}. \label{hakhkchchhacp}
    \end{flalign}
    From Equation~\eqref{hakhkchchhacp} and commutativity of $C'CHX$ and $KCHR$, we have $h_{KC} \mathop{\star} h_{CH} \mathop{=} h_{KR} \mathop{\star} h_{RH}$ and $h_{C'C} \mathop{\star} h_{CH} \mathop{=} h_{C'X} \mathop{\star} h_{XH}$ and deduce  
     \begin{flalign}
        h_{AK} \mathop{\star} h_{KR} \mathop{\star} h_{RH}= h_{AC'} \mathop{\star} h_{C'X} \mathop{\star} h_{XH}. \label{hakhkrhrhhacp}
     \end{flalign}
    Therefore, the following commutative diagram holds 

    \begin{center}
        % \resizebox{0.3\textwidth}{!}{
            \begin{tikzpicture}[scale=1.2]
                \node (k) at (0,0) {K};
                \node (r) at (4,0) {R};
                \node (h) at (4,-3) {H};
                \node (h) at (4,-3) {H};
                \node (a) at (-1,0) {A};
                \draw[<-<]  (r) -- (k) node [midway,above] {};
                \draw[>->] (r) -- (h) node[midway, left] {};
                \node (r') at (2,-1) {R'};
                \node (c') at (1,-2) {C'};
                \node () at (3,-1.5) {$PB$};
                \node (x) at (2,-2) {X};
                \draw [>->] (c') -- (x);
                \draw [>->] (r') -- (x);
                \draw[>->] (r') -- (r);
                \draw[>->] (x) -- (h);
                \draw[->] (a) -- (c');
                \draw[->] (a) -- (k);
            \end{tikzpicture}
        % }
        \end{center} 
    Since $R'XHR$ is a pullback square, by the universal property Definition~\ref{def:cat:pb}, there is a unique morphism $h_{AR'} \mathop{\colon} A \mathop{\to} R'$ such that 
    \begin{flalign}
        h_{AR'} \mathop{\star} h_{R'R} \mathop{=} h_{AK} \mathop{\star} h_{KR},  \label{harprprakkr}
        \\
        h_{AR'} \mathop{\star} h_{R'X} \mathop{=} h_{AC'} \mathop{\star} h_{C'X}. \label{harprpxacpcpx}
    \end{flalign}
    We have thus the following commutative diagram 

        \begin{center}
            % \resizebox{0.2\textwidth}{!}{
                \begin{tikzpicture}[scale=1.2]
                    \node (a) at (-1,0) {A};
                    \node (k') at (1,-1) {K'};
                    \node (r') at (2,-1) {R'};
                    \node (c') at (1,-2) {C'};
                    \node () at (1.5,-1.5) {$PB$};
                    \node (x) at (2,-2) {X};
                    \draw [>->] (c') -- (x);
                    \draw [>->] (r') -- (x);
                    \draw [>->] (k') -- (r');
                    \draw [>->] (k') -- (c');
                    \draw[->] (a) -- (c');
                    \draw[->] (a) -- (r');
                \end{tikzpicture}
            % }
        \end{center} 
    Since $K'C'XR'$ is a pullback square, there is a unique morphism $h_{AK'} \mathop{\colon} A \mathop{\to} K'$ such that 
    \begin{flalign}
        h_{AK'} \mathop{\star} h_{K'C'} \mathop{=} h_{AC'},  \label{hakphkpcphacp}
        \\
        h_{AK'} \mathop{\star} h_{K'R'} \mathop{=} h_{AR'}. \label{hakpkprparp}
    \end{flalign}
    We have $h_{AK'} \mathop{\star} h_{K'C'} \mathop{=} h_{AC'}$ by Equation~\eqref{hakphkpcphacp}, and $h_{AK'} \mathop{\star} h_{K'K} \mathop{=} h_{AK}$ because $h_{KR}$ is injective and the following holds:
    \begin{flalign*}
        (h_{AK'} \mathop{\star} h_{K'K}) \mathop{\star} h_{KR} &= h_{AK} \mathop{\star} h_{K'R'} \mathop{\star} h_{R'R} \hspace{0.5cm} &\text{by commutativity of $K'KRR'$}\\
        &= h_{AR'} \mathop{\star} h_{R'R} &\text{by Equation~\eqref{hakpkprparp}} \\
        &= h_{AK} \mathop{\star} h_{KR}. &\text{by Equation~\eqref{harprprakkr}}
    \end{flalign*}
    Now, we show the unicity of such a morphism from $A$ to $K'$. 
    Let $h_{AK'}' \mathop{\colon} A \mathop{\to} K'$ be a morphism  such that 
    \begin{flalign*}
        h_{AK'}' \mathop{\star} h_{K'C'} \mathop{=} h_{AC'},  
        \\
        h_{AK'}' \mathop{\star} h_{K'K} \mathop{=} h_{AK}.
    \end{flalign*}
    From $h_{AK'} \mathop{\star} h_{K'C'} \mathop{=} h_{AC'}$, $h_{AK'}' \mathop{\star} h_{K'C'} \mathop{=} h_{AC'}$ and injectivity of $h_{K'C'}$, we deduce $h_{AK'} \mathop{=} h_{AK'}'$.

    \qed
\end{proof}


\begin{lemma}
    \label{kpcpxrp_po}
    Consider the following commutative diagram in \textbf{Graph} where $KCHR$ is a pushout.
    \begin{center}
        % \resizebox{6cm}{!}{
            \begin{tikzpicture}
                \node (k) at (0,0) {K};
                \node (r) at (4,0) {R};
                \node (c) at (0,-3) {C};
                \node (h) at (4,-3) {H};
                % \node (rb) at ($\scl*(1.5,-0.5)$) {$R_X$};
                % \node (h') at ($\scl*(1.5,-1.5)$) {$H'$};
                % \draw[>->]  (rb) -- (h') node [midway,above] {};
                % \draw[>->]  (c) -- (h') node [midway,above] {};
                \draw[<-<]  (r) -- (k) node [midway,above] {};
                \draw[>->] (c) -- (h) node [midway, below] {};
                \draw[>->] (r) -- (h) node[midway, left] {};
                \draw[>->] (k) -- (c) node[midway, left] {};
                % \draw[->] (rb) to (l);
                % \draw[<-<] (rb) to (k);
                \node (k') at (1,-1) {K'};
                \node (r') at (2,-1) {R'};
                \node (c') at (1,-2) {C'};
                % \node () at (1.5,-1.5) {$PB$};
                \node () at (3,-1.5) {$PB$};
                \node () at (1.5,-2.5) {$PB$};
                \node () at (1.5,-0.5) {$PB$};
                \node () at (0.5,-1.5) {$PB$};
                \node (x) at (2,-2) {X};
                % \draw [->] (x) -- (h') node[midway] {!};
                \draw [>->] (c') -- (x);
                \draw [>->] (r') -- (x);
                \draw [>->] (k') -- (r');
                \draw [>->] (k') -- (c');
                \draw [>->] (c') -- (c);
                % \draw [->] (r') -- (rb);
                % \draw [->] (r') -- (rb);
                % \draw [->] (h') -- (h);
                \draw[>->] (r') -- (r);
                \draw[>->] (x) -- (h) node[midway,right] {$h$};
                % \node (rb) at ($\scl*(\sclx*1.5,-0.5)$) {$R_X$};
                % \node (h') at ($\scl*(\sclx*1.5,-1.2)$) {$H'$};
                % \draw[>->]  (rb) -- (h') node [midway,above] {};
                % \draw[>->]  (c) -- (h') node [midway,above] {};
                % \draw[>->]  (rb) -- (h') node [midway,above] {};
                % \draw[>->] (rb) to (r);
                % \draw[<-<] (rb) to (k);
                \draw[>->] (k') -- (k) ;
            \end{tikzpicture}
        % }
        \end{center} 
        The square $K'C'XR'$ is a pushout.
\end{lemma}
\begin{proof}
    \textbf{Graph} is adhesive by Proposition~\ref{prop:graph_adhesive}. By Definition~\ref{def:adhesive_cat}, pushouts along monomorphisms are VK squares in \textbf{Graph}. The square $KCHR$ is a pushout along monomorphisms in \textbf{Graph}. Hence $KCHR$ is VK square. By Definition~\ref{def:vk_square}, the square $K'C'XR'$ is pushout if the squares $C'CHX$ and $R'XHR$ are pullbacks. By assumption, the squares $C'CHX$ and $R'XHR$ are pullbacks. Therefore, $K'C'XR'$ is pushout.
\end{proof} 

\begin{lemma}
    \label{lem:xinXcpinCrpinR}
        Consider the following pushout square in \textbf{Graph}.
    \begin{center}
        \begin{tikzpicture}
            \node (k') at (0,0) {K'};
            \node (r') at (-2,0) {R'};
            \node (c') at (0,-2) {C'};
            \node (x) at (-2,-2) {X};
            \node () at (-1,-1) {$PO$};
            \draw [>->] (c') -- (x);
            \draw [>->] (r') -- (x);
            \draw [>->] (k') -- (r');
            \draw [>->] (k') -- (c');
        \end{tikzpicture}
    \end{center} 
     Let $x$ be a node (resp. an edge) in $X$. There is either a node (resp. an edge) $r'$ in $R'$ such that $h_{R'X}(r') \mathop{=} x$ or a node (resp. an edge) $c'$ in $C'$ such that $h_{C'X}(c') \mathop{=} x$.
\end{lemma}
\begin{proof}
    Suppose, by contradiction, that there is neither a node (resp. an edge) $r'$ in $R'$ such that $h_{R'X}(r') \mathop{=} x$ nor a node (resp. an edge) $c'$ in $C'$ such that $h_{C'X}(c') \mathop{=} x$.

    Let $G$ be a graph with two distinct nodes (resp. edges) $y$ and $y'$.
    By the universal property of pushout square Definition~\ref{def:cat:po}, there is a unique morphism $h_{XG}:X \rightarrowtail G$ such that $h_{C'G} \mathop{=} h_{C'X} \mathop{\star} h_{XG}$ and  $h_{R'G} \mathop{=} h_{R'X} \mathop{\star} h_{XG}$.
    \begin{center}
        \begin{tikzpicture}
            \node (g) at (-4,-4) {$G$};
            \node (k') at (-1,-1) {$K'$};
            \node (r') at (-2,-1) {$R'$};
            \node (c') at (-1,-2) {$C'$};
            \node (x) at (-2,-2) {$X$};
            \node () at (-1.5,-1.5) {$PO$};
            \draw [>->] (c') -- (x);
            \draw [>->] (r') -- (x);
            \draw [>->] (k') -- (r');
            \draw [>->] (k') -- (c');
            \draw [>->] (c') -- (g);
            \draw [>->] (r') -- (g);
            \draw [->] (x) -- (g) node[midway] {!};
        \end{tikzpicture}
    \end{center}
    However, if we define two morphisms $f$ and $g$ by 
    % \begin{flalign*}
    %     f(z) &= h_{XG}(z) \text{if $z \mathop{\neq} x$}\\
    %     f(x) &= y \\
    %     g(z) &= h_{XG}(z) \text{if $z \mathop{\neq} x$}\\
    %     g(x) &= y' \\
    % \end{flalign*}
    \[
    f(z)=
    \begin{cases}
    h_{XG}(z), & \text{if $z \mathop{\neq} x$},\\
    y,         & \text{if $z \mathop{=} x$},
    \end{cases}
    \]

    \[
    g(z)=
    \begin{cases}
    h_{XG}(z), & \text{if $z \mathop{\neq} x$},\\
    y',        & \text{if $z \mathop{=} x$}.
    \end{cases}
    \]
    we have $f \mathop{\neq} g$ and 
    \begin{flalign*}
        h_{C'G} \mathop{=} h_{C'X} \mathop{\star} f,\\
        h_{R'G} \mathop{=} h_{R'X} \mathop{\star} f,\\
        h_{C'G} \mathop{=} h_{C'X} \mathop{\star} g,\\
        h_{R'G} \mathop{=} h_{R'X} \mathop{\star} g.
    \end{flalign*}
    which is a contradiction.
\end{proof}

\begin{lemma}
    \label{lem:g_monic}
    Consider the following commutative diagram in \textbf{Graph}.
    \begin{center}
        \begin{tikzpicture}
            \node (k) at (0,0) {K};
            \node (c) at (0,-4) {C};
            \node (l) at (-4,0) {$L$};
            \node (g) at (-4,-4) {$G$};
            \draw[>->] (l) -- (g) node [midway,above] {};
            \draw[>->] (c) -- (g) node [midway,above] {};
            \draw[>->] (k) -- (c) node[midway, left] {};
            \draw[<-<] (l) to (k);
            \node (k') at (-1,-1) {$K'$};
            \node (r') at (-2,-1) {$R'$};
            \node (c') at (-1,-2) {$C'$};
            \node (x) at (-2,-2) {$X$};
            \node () at (-1.5,-1.5) {$PO$};
            \node () at (-0.5,-1.5) {$PB$};
            \draw [>->] (c') -- (x);
            \draw [>->] (r') -- (x);
            \draw [>->] (k') -- (r');
            \draw [>->] (k') -- (c');
            \draw [>->] (c') -- (c);
            \draw [>->] (r') -- (l);
            \draw [>->] (k') -- (k);
        \end{tikzpicture}
    \end{center}
    there is a unique monomorphism $h_{XG}:X \rightarrowtail G$ such that
    $C'XGC$ and $R'XGL$ are commutative squares.
    %  $h_{C'C} \mathop{\star} h_{CG} \mathop{=} h_{C'X} \mathop{\star} h_{XG}$ and  $h_{R'L} \mathop{\star} h_{LG} \mathop{=} h_{R'X} \mathop{\star} h_{XG}$.
\end{lemma}
\begin{proof}
    there is a unique morphism $h_{XG}:X \mathop{\to} G$, by Definition~\ref{def:cat:po}, because $K'C'XR'$ is a pushout\todo{$K'C'XR'$ PO} and 
    $h_{K'C'} \mathop{\star} h_{C'C} \mathop{\star} h_{CG}
    =h_{K'R'} \mathop{\star} h_{R'L} \mathop{\star} h_{LG}$.
    % the following hold
    % \begin{flalign*}
    %     &h_{K'C'} \mathop{\star} h_{C'C} \mathop{\star} h_{CG} &\\ 
    %    =&h_{K'K} \mathop{\star} h_{KC} \mathop{\star} h_{CG} &\hspace{2cm}\text{by commutativity of $K'C'CK$}\\
    %    =&h_{K'K} \mathop{\star} h_{KL} \mathop{\star} h_{LG} &\text{by the commutativity of $KCGL$} \\
    %    =&h_{K'R'} \mathop{\star} h_{R'L} \mathop{\star} h_{LG} &\text{by the commutativity of $K'KLR'$}
    % \end{flalign*}
    Therefore, the following commutative diagram holds
    \begin{center}
        \begin{tikzpicture}
            \node (k) at (0,0) {$K$};
            \node (c) at (0,-4) {$C$};
            \node (l) at (-4,0) {$L$};
            \node (g) at (-4,-4) {$G$};
            \draw[>->] (l) -- (g) node [midway,above] {};
            \draw[>->] (c) -- (g) node [midway,above] {};
            \draw[>->] (k) -- (c) node[midway, left] {};
            \draw[<-<] (l) to (k);
            \node (k') at (-1,-1) {$K'$};
            \node (r') at (-2,-1) {$R'$};
            \node (c') at (-1,-2) {$C'$};
            \node (x) at (-2,-2) {$X$};
            \node () at (-1.5,-1.5) {$PO$};
            \node () at (-0.5,-1.5) {$PB$};
            \draw [>->] (c') -- (x);
            \draw [>->] (r') -- (x);
            \draw [>->] (k') -- (r');
            \draw [>->] (k') -- (c');
            \draw [>->] (c') -- (c);
            \draw [>->] (r') -- (l);
            \draw [>->] (k') -- (k);
            \draw [->] (x) -- (g) node[midway,above] {!};
        \end{tikzpicture}
    \end{center}
    It remains to show that $h_{XG}$ is injective.
    
    Let $x_1,x_2$ be two elements in the graph $X$ such that $x_1 \mathop{\neq} x_2$. Since $K'R'XC'$ is a pushout\todo{$K'R'XC'$ P.O.} square, by Lemma~\ref{lem:xinXcpinCrpinR}, there are 4 possible cases:
    \begin{enumerate}
        \item there are $c_1, c_2 \mathop{\in} C'$ such that $h_{C'X}(c_1) \mathop{=} x_1$ and $h_{C'X}(c_2) \mathop{=} x_2$;
        \item there are $r_1, r_2 \mathop{\in} R'$ such that $h_{R'X}(r_1) \mathop{=} x_1$ and $h_{R'X}(r_2) \mathop{=} x_2$;
        \item there are $r' \mathop{\in} R'$ and $c' \mathop{\in} C'$ such that $h_{C'X}(c') \mathop{=} x_1$ and $h_{R'X}(r') \mathop{=} x_2$;
        \item there are $r' \mathop{\in} R'$ and $c' \mathop{\in} C'$ such that $h_{C'X}(c') \mathop{=} x_2$ and $h_{R'X}(r') \mathop{=} x_1$.
    \end{enumerate}
    Since the second and fourth cases are symmetric to the first and third case, respectively, it suffices to show that in the first and third cases, the inequality $h_{XG}(x_1) \mathop{\neq} h_{XG}(x_2)$ holds. 
    \begin{itemize}
        \item[Case (1)] Suppose that there are $c_1, c_2 \mathop{\in} C'$ such that $h_{C'X}(c_1) \mathop{=} x_1$ and $h_{C'X}(c_2) \mathop{=} x_2$. 
        
        We have $c_1 \mathop{\neq} c_2$ since $x_1 \mathop{\neq} x_2$ and $h_{C'X}$ is injective.
        
        Therefore we have 
        \begin{flalign*}
            h_{XG}(x_1)&= h_{XG}(h_{C'X}(c_1)) & \text{by definition of $c_1$}\\
                        &= (h_{C'X} \mathop{\star} h_{XG})(c_1)  \\
                        &= (h_{C'C} \mathop{\star} h_{CG})(c_1) & \text{by commutativity of $C'CGX$} \\
                        &\mathop{\neq} (h_{C'C} \mathop{\star} h_{CG})(c_2) & \text{by injectivity of $h_{C'C} \mathop{\star} h_{CG}$ and $c_1 \mathop{\neq} c_2$} \\
                        &= (h_{C'X} \mathop{\star} h_{XG})(c_2) & \text{by commutativity of $C'CGX$} \\
                        &= h_{XG}(h_{C'X}(c_2)) \\
                        &= h_{XG}(x_2). & \text{by definition of $c_2$}
        \end{flalign*}
        
        % \begin{center}
        %     \begin{tikzpicture}
        %         \node (k) at (0,0) {K};
        %         \node (c) at (0,-4) {C};
        %         \node (l) at (-4,0) {$L$};
        %         \node (g) at (-4,-4) {$G$};
        %         \draw[>->] (l) -- (g) node [midway,above] {};
        %         \draw[>->] (c) -- (g) node [midway,above] {};
        %         \draw[>->] (k) -- (c) node[midway, left] {};
        %         \draw[<-<] (l) to (k);
        %         \node (k') at (-1,-1) {K'};
        %         \node (r') at (-2,-1) {R'};
        %         \node (c') at (-1,-2) {C'};
        %         \node (x) at (-2,-2) {X};
        %         \node () at (-1.5,-1.5) {$PO$};
        %         \node () at (-0.5,-1.5) {$PB$};
        %         \draw [>->] (c') -- (x);
        %         \draw [>->] (r') -- (x);
        %         \draw [>->] (k') -- (r');
        %         \draw [>->] (k') -- (c');
        %         \draw [>->] (c') -- (c);
        %         \draw [>->] (r') -- (l);
        %         \draw [>->] (k') -- (k);
        %         \draw [->] (x) -- (g) node[midway,above] {!};
        %     \end{tikzpicture}
        % \end{center}
        % \item[(2)] Suppose that there are $r_1, r_2 \mathop{\in} R'$ such that $h_{R'X}(r_1) \mathop{=} x_1$ and $h_{R'X}(r_2) \mathop{=} x_2$. The proof of this case is analogous to the previous case.

        % We have $r_1 \mathop{\neq} r_2$ because $x_1 \mathop{\neq} x_2$ and $h_{R'X}$ is injective, and
        % \begin{flalign*}
        %     h_{XG}(x_1) &= h_{XG}(h_{R'X}(r_1)) & \text{by definition of $r_1$}\\
        %                  &= (h_{R'X} \mathop{\star} h_{XG})(r_1)  \\
        %                  &= (h_{R'L} \mathop{\star} h_{LG})(r_1) & \text{by commutativity of $R'XGL$} \\
        %                  &\mathop{\neq} (h_{R'L} \mathop{\star} h_{LG})(r_2) & \text{by injectivity of $h_{R'L} \mathop{\star} h_{LG}$ and $r_1 \mathop{\neq} r_2$} \\
        %                  &= (h_{R'X} \mathop{\star} h_{XG})(r_2) & \text{by commutativity of $R'XGL$} \\
        %                  &= h_{XG}(h_{R'X}(r_2)) \\
        %                  &= h_{XG}(x_2) & \text{by definition of $r_2$}
        % \end{flalign*}
        
        % \begin{center}
        %     \begin{tikzpicture}
        %         \node (k) at (0,0) {K};
        %         \node (c) at (0,-4) {C};
        %         \node (l) at (-4,0) {$L$};
        %         \node (g) at (-4,-4) {$G$};
        %         \draw[>->] (l) -- (g) node [midway,above] {};
        %         \draw[>->] (c) -- (g) node [midway,above] {};
        %         \draw[>->] (k) -- (c) node[midway, left] {};
        %         \draw[<-<] (l) to (k);
        %         \node (k') at (-1,-1) {K'};
        %         \node (r') at (-2,-1) {R'};
        %         \node (c') at (-1,-2) {C'};
        %         \node (x) at (-2,-2) {X};
        %         \node () at (-1.5,-1.5) {$PO$};
        %         \node () at (-0.5,-1.5) {$PB$};
        %         \draw [>->] (c') -- (x);
        %         \draw [>->] (r') -- (x);
        %         \draw [>->] (k') -- (r');
        %         \draw [>->] (k') -- (c');
        %         \draw [>->] (c') -- (c);
        %         \draw [>->] (r') -- (l);
        %         \draw [>->] (k') -- (k);
        %         \draw [->] (x) -- (g) node[midway,above] {!};
        %     \end{tikzpicture}
        % \end{center}
        \item[Case (3)] Suppose that there are $r' \mathop{\in} R'$ and $c' \mathop{\in} C'$ such that $h_{C'X}(c') \mathop{=} x_1$ and $h_{R'X}(r') \mathop{=} x_2$.
        % , or $h_{C'X}(c') \mathop{=} x_2$ and $h_{R'X}(r') \mathop{=} x_1$. By symmetry, we suppose that the first case holds. Let $c \mathop{=} h_{C'C}(c')$ and $l \mathop{=} h_{R'L}(r')$
        
        We have three cases:
            \begin{itemize}
                \item[(3.1)] Suppose that $h_{C'C}(c') \notin \operatorname{Im}(h_{KC})$. We have
                    \begin{flalign*}
                        h_{XG}(x_1) &= h_{XG}(h_{C'X}(c')) & \text{by definition of $c'$}\\
                                     &= (h_{C'X} \mathop{\star} h_{XG})(c') \\
                                     &= (h_{C'C} \mathop{\star} h_{CG})(c') & \text{by commutativity of $C'CGX$} \\
                                     &= h_{CG}(h_{C'C}(c')) & \\
                                    %  &= h_{CG}(c) & \\
                                     &\mathop{\neq} 
                                    %  h_{LG}(l) & \text{by Lemma~\ref{lem:b_c_same_img_exist_a} and $c\notin \operatorname{Im}(h_{KC})$}  
                                    h_{LG}(h_{R'L}(r')) & \text{by Lemma~\ref{lem:b_c_same_img_exist_a} and $h_{C'C}(c') \notin \operatorname{Im}(h_{KC})$}   
                                    \\
                                     &= (h_{R'L} \mathop{\star} h_{LG})(r') 
                                     \\
                                     &= (h_{R'X} \mathop{\star} h_{XG})(r') & \text{by commutativity of $R'XGL$} 
                                     \\
                                     &= h_{XG}(h_{R'X}(r'))
                                     \\
                                     &= h_{XG}(x_2). & \text{by definition of $r'$}
                    \end{flalign*}
                    
                    % \begin{center}
                    %     \begin{tikzpicture}
                    %         \node (k) at (0,0) {K};
                    %         \node (c) at (0,-4) {C};
                    %         \node (l) at (-4,0) {$L$};
                    %         \node (g) at (-4,-4) {$G$};
                    %         \draw[>->] (l) -- (g) node [midway,above] {};
                    %         \draw[>->] (c) -- (g) node [midway,above] {};
                    %         \draw[>->] (k) -- (c) node[midway, left] {};
                    %         \draw[<-<] (l) to (k);
                    %         \node (k') at (-1,-1) {K'};
                    %         \node (r') at (-2,-1) {R'};
                    %         \node (c') at (-1,-2) {C'};
                    %         \node (x) at (-2,-2) {X};
                    %         \node () at (-1.5,-1.5) {$PO$};
                    %         \node () at (-0.5,-1.5) {$PB$};
                    %         \draw [>->] (c') -- (x);
                    %         \draw [>->] (r') -- (x);
                    %         \draw [>->] (k') -- (r');
                    %         \draw [>->] (k') -- (c');
                    %         \draw [>->] (c') -- (c);
                    %         \draw [>->] (r') -- (l);
                    %         \draw [>->] (k') -- (k);
                    %         \draw [->] (x) -- (g) node[midway,above] {!};
                    %     \end{tikzpicture}
                    % \end{center}
                \item[(3.2)] Suppose $h_{R'L}(r') \notin \operatorname{Im}(h_{KL})$. This case is symmetric to the previous case.
                % where $c \notin \operatorname{Im}(h_{KC})$.
                \item[(3.3)] Suppose that there are $k_1, k_2 \mathop{\in} K$ such that $h_{KC}(k_1) \mathop{=} h_{C'C}(c')$ and $h_{KL}(k_2) \mathop{=} h_{R'L}(r')$. 
                
                We cannot have $k_1=k_2$. 
                
                Suppose, by contradiction, that we have 
                \begin{flalign}
                    k_1=k_2. \label{k1eqk2}
                \end{flalign} 
                
                % \begin{center}
                %     \begin{tikzpicture}
                %         \node (k) at (0,0) {K};
                %         \node (c) at (0,-4) {C};
                %         \node (l) at (-4,0) {$L$};
                %         \node (g) at (-4,-4) {$G$};
                %         \draw[>->] (l) -- (g) node [midway,above] {};
                %         \draw[>->] (c) -- (g) node [midway,above] {};
                %         \draw[>->] (k) -- (c) node[midway, left] {};
                %         \draw[<-<] (l) to (k);
                %         \node (k') at (-1,-1) {K'};
                %         \node (r') at (-2,-1) {R'};
                %         \node (c') at (-1,-2) {C'};
                %         \node (x) at (-2,-2) {X};
                %         \node () at (-1.5,-1.5) {$PO$};
                %         \node () at (-0.5,-1.5) {$PB$};
                %         \draw [>->] (c') -- (x);
                %         \draw [>->] (r') -- (x);
                %         \draw [>->] (k') -- (r');
                %         \draw [>->] (k') -- (c');
                %         \draw [>->] (c') -- (c);
                %         \draw [>->] (r') -- (l);
                %         \draw [>->] (k') -- (k);
                %         \draw [->] (x) -- (g) node[midway,above] {!};
                %     \end{tikzpicture}
                % \end{center}
                Since $K'C'CK$ is also a pullback square by Proposition~\ref{prop:pb_eq_po} and $h_{KC}(k_1) \mathop{=} h_{C'C}(c')$ and $h_{C'C}(c') \mathop{=}  h_{C'C}(c')$, there is $k' \mathop{\in} K'$ such that $h_{K'K}(k') \mathop{=} k_1$ and $h_{K'C'}(k') \mathop{=} c'$.
 
                We have $h_{K'R'}(k') \mathop{=} r'$ because $h_{R'L}$ is injective, $h_{R'L}(r') \mathop{=} h_{R'L}(r')$ and 
                \begin{flalign*}
                   &h_{R'L}(h_{K'R'}(k')) \hspace{2cm}&\\
                  =&(h_{K'R'} \mathop{\star} h_{R'L})(k') &\\
                  =&(h_{K'K} \mathop{\star} h_{KL})(k') & \text{by commutativity of $K'R'LK$}\\
                  =&(h_{KL})(h_{K'K}(k'))\\
                  =&(h_{KL})(k_1)\\
                  =&(h_{KL})(k_2). & \text{by assumption~\eqref{k1eqk2}}
                %  \\
                %   =& l
                \end{flalign*}
 
             Therefore, we have the contradiction:
                     \begin{flalign*}
                         x_1 &= h_{C'X}(c')  & \text{by definition of $c'$}\\
                             &= h_{C'X}(h_{K'C'}(k')) & \text{by definition of $k'$}\\
                             &= (h_{K'C'} \mathop{\star} h_{C'X})(k') &  \\
                             &= (h_{K'R'} \mathop{\star} h_{R'X})(k') & \text{by the commutativity of $K'C'XR'$} \\
                             &= h_{R'X}(h_{K'R'}(k')) &  \\
                             &= h_{R'X}(r') &  \text{by definition of $r'$}\\
                             &= x_2.& \text{by definition of $x_2$}
                     \end{flalign*}

                Thus, the following inequality holds 
                \begin{flalign}
                    k_1 \mathop{\neq} k_2. \label{k1neqk2}
                \end{flalign}

                \begin{flalign*}
                    h_{XG}(x_1) &= h_{XG}(h_{C'X}(c')) & \text{by definition of $c'$} \\
                            &= (h_{C'X} \mathop{\star} h_{XG})(c') &  \\
                            &= (h_{C'C} \mathop{\star} h_{CG})(c') & \text{by commutativity of $C'CGX$} \\
                            % &= (h_{CG})(c) & \text{by definition of $c$} \\
                            &= (h_{CG})(h_{KC}(k_1)) & \text{by definition of $k_1$} 
                            \\
                            &= (h_{KC} \mathop{\star} h_{CG})(k_1) & 
                            \\
                            &\mathop{\neq} (h_{KC} \mathop{\star} h_{CG})(k_2)  & \text{by injectivity of $h_{KC}\star h_{CG}$ and~\eqref{k1neqk2}} 
                            \\
                            &= (h_{KL} \mathop{\star} h_{LG})(k_2) & \text{by commutativity of $KLGC$}
                            \\
                            &= h_{LG}(h_{KL}(k_2)) & 
                            \\
                            % &= h_{LG}(l) & \text{by the definition of $k_2$}
                            % \\
                            &= h_{LG}(h_{R'L}(r')) & \text{by the definition of $k_2$}
                            \\
                            &= (h_{R'L} \mathop{\star} h_{LG})(r') & 
                            \\
                            &= (h_{R'X} \mathop{\star} h_{XG})(r') & \text{by commutativity of $R'XLG$} \\
                            &= h_{XG}(h_{R'X}(r'))\\
                                &= h_{XG}(x_2). & \text{by definition of $r'$}
                \end{flalign*}
                
                % \begin{center}
                %     \begin{tikzpicture}
                %         \node (k) at (0,0) {K};
                %         \node (c) at (0,-4) {C};
                %         \node (l) at (-4,0) {$L$};
                %         \node (g) at (-4,-4) {$G$};
                %         \draw[>->] (l) -- (g) node [midway,above] {};
                %         \draw[>->] (c) -- (g) node [midway,above] {};
                %         \draw[>->] (k) -- (c) node[midway, left] {};
                %         \draw[<-<] (l) to (k);
                %         \node (k') at (-1,-1) {K'};
                %         \node (r') at (-2,-1) {R'};
                %         \node (c') at (-1,-2) {C'};
                %         \node (x) at (-2,-2) {X};
                %         \node () at (-1.5,-1.5) {$PO$};
                %         \node () at (-0.5,-1.5) {$PB$};
                %         \draw [>->] (c') -- (x);
                %         \draw [>->] (r') -- (x);
                %         \draw [>->] (k') -- (r');
                %         \draw [>->] (k') -- (c');
                %         \draw [>->] (c') -- (c);
                %         \draw [>->] (r') -- (l);
                %         \draw [>->] (k') -- (k);
                %         \draw [->] (x) -- (g) node[midway,above] {!};
                %     \end{tikzpicture}
                % \end{center}
            \end{itemize}            
    \end{itemize}
\end{proof}



\begin{lemma}
    \label{lem:h_hp_diff_g_gp_diff}
    
    Let $X$ be a graph. Let $\rho=(L \overset{l}{\leftarrowtail} K \overset{r}{\rightarrowtail} R)$ be an injective DPO rewriting rule which is \(X\)-non-increasing under $\Psi$. Consider the following DPO diagram:

    \begin{center}
        \begin{tikzpicture}
            \node (k) at (0,1) {$K$};
            \node (l) at (-2,1) {$L$};
            \node (r) at (2,1) {$R$};
            \node (c) at (0,-1) {$C$};
            \node (g) at (-2,-1) {$G$};
            \node (h) at (2,-1) {$H$};
            % \node (rb) at (1.5,-0.5) {$R_X$};
            % \node (h') at (1.5,-1.5) {$H'$};
            % \draw[>->]  (rb) -- (h') node [midway,above] {};
            % \draw[>->]  (c) -- (h') node [midway,above] {};
            \draw[<-<]  (l) -- (k) node [midway,below] {$l$};
            \draw[>->]  (k) -- (r) node [midway,below] {$r$};
            \draw[>->] (c) -- (g) node [midway, above] {$l'$};
            \draw[>->] (c) -- (h) node [midway,above] {$r'$};
            \draw[>->] (l) -- (g) node[midway, right] {$m$};
            \draw[>->] (r) -- (h) node[midway, left] {$m'$};
            \draw[>->] (k) -- (c) node[midway, left] {};
            \node () [at=($(l)!0.5!(c)$)] {$PO$};
            \node () [at=($(r)!0.5!(c)$)] {$PO$};
            % \draw[>->] (rb) to (r);
            % \draw[->] (rb) to (l);
            % \draw[<-<] (rb) to (k);
        \end{tikzpicture}
    \end{center}
   
    Let $h_{XH}',h_{XH}'':X \rightarrowtail H$ be monomorphisms.
    In the category \textbf{Graph}, we can construct the following commutative diagram
    \begin{center}
        \begin{tikzpicture} 
            \node (k) at (0,0) {K};
            \node (c) at (0,-4) {C};
            \node (l) at (-4,0) {$R$};
            \node (g) at (-4,-4) {$H$};
            \node (r) at (4,0) {$L$};
            \node (h) at (4,-4) {$G$};
            \draw [>->] (k) -- (r); 
            \draw [>->] (c) -- (h);
            \draw [>->] (r) -- (h);
            \draw[>->] (l) -- (g) node [midway,above] {};
            \draw[>->] (c) -- (g) node [midway,above] {};
            \draw[>->] (k) -- (c) node[midway, left] {};
            \draw[<-<] (l) to (k);
            \node (k') at (-2,-1) {K'};
            \node (r') at (-3.5,-1.5) {R'};
            \node (c') at (-1.5,-2.2) {C'};
            \node (x) at (-3,-3) {X};
            \draw [>->] (c') -- (x);
            \draw [>->] (r') -- (x); 
            \draw [>->] (k') -- (r');
            \draw [>->] (k') -- (c');
            \draw [>->] (c') -- (c);
            \draw [>->] (r') -- (l);
            \draw [>->] (k') -- (k);

            \node (k'') at (-0.5,-1.5) {K''};
            \node (r'') at (-2.5,-2) {R''};
            \node (c'') at (-0.5,-2.5) {C''};
            \draw [>->] (c'') -- (x);
            \draw [>->] (r'') -- (x);
            \draw [>->] (k'') -- (r'');
            \draw [>->] (k'') -- (c'');
            \draw [>->] (k'') -- (k);
            \draw [>->] (c'') -- (c);
            \draw [>->] (r'') -- (l);

            \draw [>->] (x) edge[bend left] node[pos=0.1,below]{$h_{XH}'$} (g);
            \draw [>->] (x) edge[bend right] node[midway,above]{$h_{XH}''$} (g);

        \end{tikzpicture}
    \end{center}
    where $K'R'XC'$, $K''R''XC''$, $K'KCC'$, $K''KCC''$, $R'RHX$ and $R''RHX$ are pullbacks.

    By Lemma~\ref{kpcpck_pullback}, $K'KRR'$ and $K''KCC''$ are pullbacks. 

    By Lemma~\ref{kpcpxrp_po}, $K'R'XC'$ and $K''R''XC''$ are pushouts.

    Since $\rho$ is $X$-non-increasing under $\Psi$, there are $\Psi(R'): R' \rightarrowtail L$ and $\Psi(R''):R'' \rightarrowtail L$ such that $K'R'LK$ and $K''R''LK$ are pullbacks.

    By Lemma~\ref{lem:g_monic}, we have the following commutative diagram
    \begin{center}
        \begin{tikzpicture} 
            \node (k) at (0,0) {K};
            \node (c) at (0,-4) {C};
            \node (l) at (-4,0) {L};
            \node (g) at (-4,-4) {G};
            \node (r) at (4,0) {L};
            \node (h) at (4,-4) {G};
            \draw [>->] (k) -- (r); 
            \draw [>->] (c) -- (h);
            \draw [>->] (r) -- (h);
            \draw[>->] (l) -- (g) node [midway,above] {};
            \draw[>->] (c) -- (g) node [midway,above] {};
            \draw[>->] (k) -- (c) node[midway, left] {};
            \draw[<-<] (l) to (k);
            \node (k') at (-2,-1) {K'};
            \node (r') at (-3.5,-1.5) {R'};
            \node (c') at (-1.5,-2.2) {C'};
            \node (x) at (-3,-3) {X};
            \draw [>->] (c') -- (x);
            \draw [>->] (r') -- (x); 
            \draw [>->] (k') -- (r');
            \draw [>->] (k') -- (c');
            \draw [>->] (c') -- (c);
            \draw [>->] (r') -- (l);
            \draw [>->] (k') -- (k);

            \node (k'') at (-0.5,-1.5) {K''};
            \node (r'') at (-2.5,-2) {R''};
            \node (c'') at (-0.5,-2.5) {C''};
            \draw [>->] (c'') -- (x);
            \draw [>->] (r'') -- (x);
            \draw [>->] (k'') -- (r'');
            \draw [>->] (k'') -- (c'');
            \draw [>->] (k'') -- (k);
            \draw [>->] (c'') -- (c);
            \draw [>->] (r'') -- (l);

            \draw [>->] (x) edge[bend left] node[pos=0.1,below]{$h_{XG}'$} (g);
            \draw [>->] (x) edge[bend right] node[midway,above]{$h_{XG}''$} (g);

        \end{tikzpicture}
    \end{center}
    where  

    \begin{itemize}
        \item $K'R'XC'$ is pushout and pullback;
        \item $K''R''XC''$ is pushout and pullback;
        \item $K'KCC'$ and $K''KCC''$ are pullback;
        \item $h_{XG}':X \mathop{\to} G$ is the unique morphism such that $h_{C'C} \mathop{\star} h_{CG} \mathop{=} h_{C'X} \mathop{\star} h_{XG}'$ and  $h_{R'L} \mathop{\star} h_{LG} \mathop{=} h_{R'X} \mathop{\star} h_{XG}'$ where $h_{R'L} \mathop{=} \Psi(R')$;
        \item $h_{XG}'':X \mathop{\to} G$ is the unique morphism such that $h_{C''C} \mathop{\star} h_{CG} \mathop{=} h_{C''X} \mathop{\star} h_{XG}''$ and  $h_{R''L} \mathop{\star} h_{LG} \mathop{=} h_{R''X} \mathop{\star} h_{XG}''$ where $h_{R''L} \mathop{=} \Psi(R'')$.
    \end{itemize}
    If $h_{XH}' \mathop{\neq} h_{XH}''$ then $h_{XG}' \mathop{\neq} h_{XG}''$.
\end{lemma}
\begin{proof}
    Suppose that $h_{XH}' \mathop{\neq} h_{XH}''$. We are going to show $h_{XG}' \mathop{\neq} h_{XG}''$.

    From $h_{XH}' \mathop{\neq} h_{XH}''$, we deduce there is an element $x \mathop{\in} X$ such that the inequality $h_{XH}(x) \mathop{\neq} h_{XH}''(x)$ holds. We can suppose that 
        \begin{flalign}
            \text{$x$ is either an isolated nodes or an edge,} \label{x_isolated_or_edge}
        \end{flalign} 
         \todo{$x$ is either an isolated node or an edge}
    because otherwise $x$ would be a non-isolated node and we can take an incident edge of $x$.

    \noindent
    Since $K'R'XC'$ and $K''R''XC''$ are pushouts, by Lemma~\ref{lem:xinXcpinCrpinR}, there are 4 possible cases:
    \begin{enumerate}
        \item there are $c' \mathop{\in} C', c'' \mathop{\in} C''$ such that $h_{C'X}(c') \mathop{=} h_{C'X}(c'') \mathop{=} x$;
        \item  there are $r' \mathop{\in} R', r'' \mathop{\in} R''$ such that $h_{R'X}(r') \mathop{=} h_{R'X}(r'') \mathop{=} x$;
        \item there are arrows $r' \mathop{\in} R'$ and $c'' \mathop{\in} C''$ such that 
        \begin{flalign*}
            h_{C''X}(c'') \mathop{=} x
            \\
            h_{R'X}(r') \mathop{=} x
            \\
            \nexists c' \mathop{\in} C'. h_{C'X}(c') \mathop{=} x 
            \\
            \nexists r'' \mathop{\in} R''. h_{R''X}(r'') \mathop{=} x;
        \end{flalign*}
        \item there are arrows $r'' \mathop{\in} R''$ and $c' \mathop{\in} C'$ such that 
        \begin{flalign*}
            h_{C'X}(c') \mathop{=} x
            \\
            h_{R''X}(r'') \mathop{=} x
            \\
            \nexists c'' \mathop{\in} C''. h_{C''X}(c'') \mathop{=} x 
            \\
            \nexists r' \mathop{\in} R'. h_{R'X}(r') \mathop{=} x.  
        \end{flalign*}
    \end{enumerate}

    We are going to show that in each case we have $h_{XG}'(x) \mathop{\neq} h_{XG}''(x)$. Since the fourth case is symmetric to the third case, it suffices to analyse the first three cases.

    \begin{itemize}
        \item[(1)] Suppose that there are $c' \mathop{\in} C', c'' \mathop{\in} C''$ such that $h_{C'X}(c') \mathop{=} h_{C'X}(c'') \mathop{=} x$. 
        Suppose, by contradiction, that the following equality holds
        \begin{flalign}
            h_{C'C}(c') \mathop{=} h_{C''C}(c''). \label{cpccp_eq_cppccpp}
        \end{flalign}
         We have the following contradiction:
        \begin{flalign*}
            &h_{XH}'(x) \\
           =&h_{XH}'(h_{C'X}(c')) \\
           =&(h_{C'X} \mathop{\star} h_{XH}')(c') \\
           =&(h_{C'C} \mathop{\star} h_{CH})(c') &\text{By commutativity of $C'XHC$}\\
           =&h_{CH}(h_{C'C}(c')) \\
           =&h_{CH}(h_{C''C}(c'')) & \text{ by assumption~\eqref{cpccp_eq_cppccpp}}\\
           =&(h_{C''C} \mathop{\star} h_{CH})(c'') \\
           =&(h_{C''X} \mathop{\star} h_{XH}'')(c'') &\text{By commutativity of $C''CHX$}\\
           =&h_{XH}''(h_{C''X}(c'')) \\
           \mathop{=} &h_{XH}''(x). \\
        \end{flalign*}
        
        % \begin{center}
        %     \begin{tikzpicture} 
        %         \node (k) at (0,0) {K};
        %         \node (c) at (0,-4) {C};
        %         \node (l) at (-4,0) {$L$};
        %         \node (g) at (-4,-4) {$G$};
        %         \draw[>->] (l) -- (g) node [midway,above] {};
        %         \draw[>->] (c) -- (g) node [midway,above] {};
        %         \draw[>->] (k) -- (c) node[midway, left] {};
        %         \draw[<-<] (l) to (k);
        %         \node (k') at (-2,-1) {K'};
        %         \node (r') at (-3.5,-1.5) {R'};
        %         \node (c') at (-1.5,-2.2) {C'};
        %         \node (x) at (-3,-3) {X};
        %         \draw [>->] (c') -- (x);
        %         \draw [>->] (r') -- (x); 
        %         \draw [>->] (k') -- (r');
        %         \draw [>->] (k') -- (c');
        %         \draw [>->] (c') -- (c);
        %         \draw [>->] (r') -- (l);
        %         \draw [>->] (k') -- (k);
    
        %         \node (k'') at (-0.5,-1.5) {K''};
        %         \node (r'') at (-2.5,-2) {R''};
        %         \node (c'') at (-0.5,-2.5) {C''};
        %         \draw [>->] (c'') -- (x);
        %         \draw [>->] (r'') -- (x);
        %         \draw [>->] (k'') -- (r'');
        %         \draw [>->] (k'') -- (c'');
        %         \draw [>->] (k'') -- (k);
        %         \draw [>->] (c'') -- (c);
        %         \draw [>->] (r'') -- (l);
    
        %         \draw [>->] (x) edge[bend left] node[pos=0.3,below]{!g''} (g);
        %         \draw [>->] (x) edge[bend right] node[midway,above]{!g'} (g);
    
        %     \end{tikzpicture}
        % \end{center}
        Therefore, we have 
        \begin{flalign}
            h_{C'C}(c') \mathop{\neq} h_{C''C}(c''). \label{cpccp_neq_cppccpp}
        \end{flalign} 
        \begin{flalign*}
            h_{XG}'(x)&= h_{XG}'(h_{C'X}(c')) & \text{by definition of $c'$}\\
                        &= (h_{C'X} \mathop{\star} h_{XG}')(c')  \\
                        &= (h_{C'C} \mathop{\star} h_{CG})(c') & \text{by commutativity of $C'CGX$} \\
                        &= h_{CG}(h_{C'C}(c')) &  \\
                        &\mathop{\neq} h_{CG}(h_{C''C}(c'')) & \text{by~\eqref{cpccp_neq_cppccpp} and the injectivity of $h_{CG}$} \\
                        &= (h_{C''C} \mathop{\star} h_{CG})(c'') &  \\
                        &= (h_{C''X} \mathop{\star} h_{XG}'')(c'') & \text{by commutativity of $C''CGX$} \\
                        &= h_{XG}''(h_{C''X}(c'')) \\
                        &= h_{XG}''(x). & \text{by definition of $c''$}
        \end{flalign*}
  
        \item[(2)] Suppose that there are $r' \mathop{\in} R', r'' \mathop{\in} R''$ such that $h_{R'X}(r') \mathop{=} h_{R'X}(r'') \mathop{=} x$. We have 
        \begin{flalign}
            h_{R'R}(r') \mathop{\neq} h_{R''R}(r'') \label{rprrppr}
        \end{flalign}  
        by an argument analogous to the argument for $h_{C'C}(c') \mathop{\neq} h_{C''C}(c'')$ in the previous case. Since $\rho$ is $X$-non-increasing, by~\eqref{rprrppr} and~\eqref{x_isolated_or_edge}, we have \todo{use : x is either an isolated node or an edge}       
        \begin{flalign}
            h_{R'L}(r') \mathop{\neq} h_{R''L}(r''). \label{rplrprpplrpp}
        \end{flalign}  

        % Therefore, we have  
        % \begin{flalign*}
        %     h_{XG}'(x)&= h_{XG}'(h_{R'X}(r')) & \text{by definition of $r'$}\\
        %         &= (h_{R'X} \mathop{\star} h_{XG}')(r')  \\
        %         &= (h_{R'L} \mathop{\star} h_{LG})(r') & \text{by commutativity of $R'LGX$} \\
        %         &= h_{LG}(h_{R'L}(r')) &  \\
        %         &\mathop{\neq} h_{LG}(h_{R''L}(r'')) & \text{by $h_{R'R}(r') \mathop{\neq} h_{R''R}(r'')$ and the injectivity of $h_{RG}$} \\
        %         &= (h_{R''L} \mathop{\star} h_{LG})(r'') &  \\
        %         &= (h_{R''X} \mathop{\star} h_{XG}'')(r'') & \text{by commutativity of $R''LGX$} \\
        %         &= h_{XG}''(h_{R''X}(r'')) \\
        %         &= h_{XG}''(x) & \text{by definition of $r''$}
        % \end{flalign*}
         
        % \todo{$h_{R_XL}$ edge-injective} Since $h_{R_XL}$ is edge-injective by assumption, because it is $X$-non-increasing by assumption. Therefore, we have $h_{R'L}(r') \mathop{\neq} h_{R''L}(r'')$ because \( h_{R'R}(r') \mathop{\neq} h_{R''R}(r'') \). Hence, from~\ref{rprneqrppr}, we have 
        % \begin{flalign}
        %     h_{R'L}(r') \mathop{\neq} h_{R''L}(r'') \label{rplrprpplrpp}
        % \end{flalign}
        Hence, the following inequality holds
        
        % \begin{center}
        %     \begin{tikzpicture} 
        %         \node (k) at (0,0) {K};
        %         \node (c) at (0,-4) {C};
        %         \node (l) at (-4,0) {$L$};
        %         \node (g) at (-4,-4) {$G$};
        %         \draw[>->] (l) -- (g) node [midway,above] {};
        %         \draw[>->] (c) -- (g) node [midway,above] {};
        %         \draw[>->] (k) -- (c) node[midway, left] {};
        %         \draw[<-<] (l) to (k);
        %         \node (k') at (-2,-1) {K'};
        %         \node (r') at (-3.5,-1.5) {R'};
        %         \node (c') at (-1.5,-2.2) {C'};
        %         \node (x) at (-3,-3) {X};
        %         \draw [>->] (c') -- (x);
        %         \draw [>->] (r') -- (x); 
        %         \draw [>->] (k') -- (r');
        %         \draw [>->] (k') -- (c');
        %         \draw [>->] (c') -- (c);
        %         \draw [>->] (r') -- (l);
        %         \draw [>->] (k') -- (k);
    
        %         \node (k'') at (-0.5,-1.5) {K''};
        %         \node (r'') at (-2.5,-2) {R''};
        %         \node (c'') at (-0.5,-2.5) {C''};
        %         \draw [>->] (c'') -- (x);
        %         \draw [>->] (r'') -- (x);
        %         \draw [>->] (k'') -- (r'');
        %         \draw [>->] (k'') -- (c'');
        %         \draw [>->] (k'') -- (k);
        %         \draw [>->] (c'') -- (c);
        %         \draw [>->] (r'') -- (l);
    
        %         \draw [>->] (x) edge[bend left] node[pos=0.3,below]{!g''} (g);
        %         \draw [>->] (x) edge[bend right] node[midway,above]{!g'} (g);
    
        %     \end{tikzpicture}
        % \end{center}
        \begin{flalign*}
            h_{XG}'(x) &= h_{XG}'(h_{R'X}(r')) & \text{by definition of $r'$}\\
                         &= (h_{R'X} \mathop{\star} h_{XG}')(r')  \\
                         &= (h_{R'L} \mathop{\star} h_{LG})(r') & \text{by commutativity of $R'XGL$} \\
                         &= h_{LG}(h_{R'L}(r')) &  \\
                         &\mathop{\neq} h_{LG}(h_{R''L}(r'')) &  \text{By injectivity of $h_{LG}$ and~\eqref{rplrprpplrpp}}\\
                         &= (h_{R''L} \mathop{\star} h_{LG})(r'') &  \\
                         &= (h_{R''X} \mathop{\star} h_{XG}'')(r'') & \text{by commutativity of $R''XGL$} \\
                         &= h_{XG}''(h_{R''X}(r'')) \\
                         &= h_{XG}''(x). & \text{by definition of $r''$}
        \end{flalign*}  
        
        % \begin{center}
        %     \begin{tikzpicture} 
        %         \node (k) at (0,0) {K};
        %         \node (c) at (0,-4) {C};
        %         \node (l) at (-4,0) {$L$};
        %         \node (g) at (-4,-4) {$G$};
        %         \draw[>->] (l) -- (g) node [midway,above] {};
        %         \draw[>->] (c) -- (g) node [midway,above] {};
        %         \draw[>->] (k) -- (c) node[midway, left] {};
        %         \draw[<-<] (l) to (k);
        %         \node (k') at (-2,-1) {K'};
        %         \node (r') at (-3.5,-1.5) {R'};
        %         \node (c') at (-1.5,-2.2) {C'};
        %         \node (x) at (-3,-3) {X};
        %         \draw [>->] (c') -- (x);
        %         \draw [>->] (r') -- (x); 
        %         \draw [>->] (k') -- (r');
        %         \draw [>->] (k') -- (c');
        %         \draw [>->] (c') -- (c);
        %         \draw [>->] (r') -- (l);
        %         \draw [>->] (k') -- (k);
    
        %         \node (k'') at (-0.5,-1.5) {K''};
        %         \node (r'') at (-2.5,-2) {R''};
        %         \node (c'') at (-0.5,-2.5) {C''};
        %         \draw [>->] (c'') -- (x);
        %         \draw [>->] (r'') -- (x);
        %         \draw [>->] (k'') -- (r'');
        %         \draw [>->] (k'') -- (c'');
        %         \draw [>->] (k'') -- (k);
        %         \draw [>->] (c'') -- (c);
        %         \draw [>->] (r'') -- (l);
    
        %         \draw [>->] (x) edge[bend left] node[pos=0.3,below]{!g''} (g);
        %         \draw [>->] (x) edge[bend right] node[midway,above]{!g'} (g);
    
        %     \end{tikzpicture}
        % \end{center}
        \item[(3)] Suppose, that there are arrows $r' \mathop{\in} R'$ and $c'' \mathop{\in} C''$ such that the following statements hold:
        \begin{flalign}
            h_{C''X}(c'') \mathop{=} x, \nonumber
            \\
            h_{R'X}(r') \mathop{=} x, \nonumber
            \\
            \nexists c' \mathop{\in} C'. h_{C'X}(c') \mathop{=} x, \nonumber
            \\
            \nexists r'' \mathop{\in} R''. h_{R''X}(r'') \mathop{=} x.  \label{assumption_norpp}
        \end{flalign}

        Suppose, by contradiction, that we have \begin{flalign}
            h_{C''C}(c'') \mathop{\in} \operatorname{Im}(h_{KC}).  \label{assump_c_in_imhkc}
        \end{flalign} 
        Then,
        there are $k_2 \mathop{\in} K$ such that $h_{KC}(k_2) \mathop{=} h_{C''C}(c'')$.

        There is $k'' \mathop{\in} K''$ such that $h_{K''K}(k'') \mathop{=} k_2$ and $h_{K''C''}(k'') \mathop{=} c''$ becase $K''C''CK$ is a pullback square and $h_{KC}(k_2) \mathop{=} h_{C''C}(c'')$.

        We have $h_{K''R''}(k'') \mathop{\in} R''$ and 
        \begin{flalign*}
          &h_{R''X}(h_{K''R''}(k''))\\
         =&(h_{K''R''} \mathop{\star} h_{R''X})(k'') \\
         =&(h_{K''C''} \mathop{\star} h_{C''X})(k'') & \text{by commutativity of $K''R''XC''$}\\
         =&h_{C''X}(h_{K''C''}(k''))\\
         =&h_{C''X}(c'')\\
         =&x.
        \end{flalign*}
        which contradicts the assumption~\eqref{assumption_norpp}.
        
        Thus, the following holds
        \begin{flalign}
            h_{C''C}(c'') \notin \operatorname{Im}(h_{KC}).  \label{assump_c_notin_imhkc}
        \end{flalign} 

        We deduce
        \begin{flalign*}
            h_{XG}'(x) &= h_{XG}'(h_{C''X}(c'')) & \text{by definition of $c''$}\\
                         &= (h_{C''X} \mathop{\star} h_{XG}')(c'') \\
                         &= (h_{C''C} \mathop{\star} h_{CG})(c'') & \text{by commutativity of $C''CGX$} \\
                         &= h_{CG}(h_{C''C}(c'')) &  \\
                         &\mathop{\neq}  h_{LG}(h_{R'L}(r')) & \text{by Lemma~\ref{lem:b_c_same_img_exist_a} and the assumption~\eqref{assump_c_notin_imhkc}} \\
                         &= (h_{R'L} \mathop{\star} h_{LG})(r') \\
                         &= (h_{R'X} \mathop{\star} h_{XG}'')(r') & \text{by commutativity of $R'XGL$} \\
                         &= h_{XG}''(h_{R'X}(r'))\\
                         &= h_{XG}''(x). & \text{by definition of $r'$}
        \end{flalign*} 

    %     We analyse the two cases where $h_{C''C}(c'') \notin \operatorname{Im}(h_{KC}) $ and 
    %         \begin{itemize}
    %             \item[(3.1)] Suppose, that the following holds
    %                 \begin{flalign}
    %                     h_{C''C}(c'') \notin \operatorname{Im}(h_{KC})  \label{assump_c_notin_imhkc}
    %                 \end{flalign} then, we have 
    %                 \begin{flalign*}
    %                     h_{XG}'(x) &= h_{XG}'(h_{C''X}(c'')) & \text{by definition of $c''$}\\
    %                                  &= (h_{C''X} \mathop{\star} h_{XG}')(c'') \\
    %                                  &= (h_{C''C} \mathop{\star} h_{CG})(c'') & \text{by commutativity of $C''CGX$} \\
    %                                  &= h_{CG}(h_{C''C}(c'')) &  \\
    %                                  &\mathop{\neq}  h_{LG}(h_{R'L}(r')) & \text{by Lemma~\ref{lem:b_c_same_img_exist_a} and the assumption Equation~\eqref{assump_c_notin_imhkc}} \\
    %                                  &= (h_{R'L} \mathop{\star} h_{LG})(r') \\
    %                                  &= (h_{R'X} \mathop{\star} h_{XG}'')(r') & \text{by commutativity of $R'XGL$} \\
    %                                  &= h_{XG}''(h_{R'X}(r'))\\
    %                                  &= h_{XG}''(x) & \text{by definition of $r'$}
    %                 \end{flalign*} 
    %             \item[(3.2)] Suppose, that the following holds
    %             \begin{flalign}
    %                 h_{R'R}(r') \notin \operatorname{Im}(h_{KR}) \label{rprrp_im_kr}
    %             \end{flalign} 
    %             Since $\rho$ is $X$-non-increasing, by the assumption Equation~\eqref{rprrp_im_kr}, (Definition~\ref{def:creates_more_x_on_the_left})$h_{R'L}$ does not identify non-interface elements with interface elements. Therefore, the following holds 
    %             \begin{flalign}
    %                 h_{R'L}(r') \notin \operatorname{Im}(h_{KL}) \label{lnotinhkl}
    %             \end{flalign} 
    %             \todo{assumption: does not identify non-interface elements with interface elements} 
    %             Therefore, 
    %     
    %     % \begin{center}
    %     %     \begin{tikzpicture} 
    %     %         \node (k) at (0,0) {K};
    %     %         \node (c) at (0,-4) {C};
    %     %         \node (l) at (-4,0) {$L$};
    %     %         \node (g) at (-4,-4) {$G$};
    %     %         \draw[>->] (l) -- (g) node [midway,above] {};
    %     %         \draw[>->] (c) -- (g) node [midway,above] {};
    %     %         \draw[>->] (k) -- (c) node[midway, left] {};
    %     %         \draw[<-<] (l) to (k);
    %     %         \node (k') at (-2,-1) {K'};
    %     %         \node (r') at (-3.5,-1.5) {R'};
    %     %         \node (c') at (-1.5,-2.2) {C'};
    %     %         \node (x) at (-3,-3) {X};
    %     %         \draw [>->] (c') -- (x);
    %     %         \draw [>->] (r') -- (x); 
    %     %         \draw [>->] (k') -- (r');
    %     %         \draw [>->] (k') -- (c');
    %     %         \draw [>->] (c') -- (c);
    %     %         \draw [>->] (r') -- (l);
    %     %         \draw [>->] (k') -- (k);
    
    %     %         \node (k'') at (-0.5,-1.5) {K''};
    %     %         \node (r'') at (-2.5,-2) {R''};
    %     %         \node (c'') at (-0.5,-2.5) {C''};
    %     %         \draw [>->] (c'') -- (x);
    %     %         \draw [>->] (r'') -- (x);
    %     %         \draw [>->] (k'') -- (r'');
    %     %         \draw [>->] (k'') -- (c'');
    %     %         \draw [>->] (k'') -- (k);
    %     %         \draw [>->] (c'') -- (c);
    %     %         \draw [>->] (r'') -- (l);
    
    %     %         \draw [>->] (x) edge[bend left] node[pos=0.3,below]{!g''} (g);
    %     %         \draw [>->] (x) edge[bend right] node[midway,above]{!g'} (g);
    
    %     %     \end{tikzpicture}
    %     % \end{center}
    %                 \begin{flalign*}
    %                     h_{XG}'(x) &= h_{XG}'(h_{R'X}(r')) & \text{by definition of $r'$}\\
    %                                  &= (h_{R'X} \mathop{\star} h_{XG}')(r') \\
    %                                  &= (h_{R'L} \mathop{\star} h_{LG})(r') & \text{by commutativity of $R'LGX$} \\
    %                                  &= h_{LG}(h_{R'L}(r')) &  \\
    %                                  &= h_{CG}(h_{C''C}(c'')) & \text{by Lemma~\ref{lem:b_c_same_img_exist_a} and the assumption Equation~\eqref{lnotinhkl}}  \\
    %                                  &= (h_{C''C} \mathop{\star} h_{CG})(c'') \\
    %                                  &= (h_{C''X} \mathop{\star} h_{XG}'')(c'') & \text{by commutativity of $C''XGC$} \\
    %                                  &= h_{XG}''(h_{C''X}(c''))\\
    %                                  &= h_{XG}''(x) & \text{by definition of $c''$}
    %                 \end{flalign*}
    %                  \newpage  
    %          
    % \begin{center}
    %         \begin{tikzpicture} 
    %             \node (k) at (0,0) {K};
    %             \node (c) at (0,-4) {C};
    %             \node (l) at (-4,0) {$L$};
    %             \node (g) at (-4,-4) {$G$};
    %             \draw[>->] (l) -- (g) node [midway,above] {};
    %             \draw[>->] (c) -- (g) node [midway,above] {};
    %             \draw[>->] (k) -- (c) node[midway, left] {};
    %             \draw[<-<] (l) to (k);
    %             \node (k') at (-2,-1) {K'};
    %             \node (r') at (-3.5,-1.5) {R'};
    %             \node (c') at (-1.5,-2.2) {C'};
    %             \node (x) at (-3,-3) {X};
    %             \draw [>->] (c') -- (x);
    %             \draw [>->] (r') -- (x); 
    %             \draw [>->] (k') -- (r');
    %             \draw [>->] (k') -- (c');
    %             \draw [>->] (c') -- (c);
    %             \draw [>->] (r') -- (l);
    %             \draw [>->] (k') -- (k);
    
    %             \node (k'') at (-0.5,-1.5) {K''};
    %             \node (r'') at (-2.5,-2) {R''};
    %             \node (c'') at (-0.5,-2.5) {C''};
    %             \draw [>->] (c'') -- (x);
    %             \draw [>->] (r'') -- (x);
    %             \draw [>->] (k'') -- (r'');
    %             \draw [>->] (k'') -- (c'');
    %             \draw [>->] (k'') -- (k);
    %             \draw [>->] (c'') -- (c);
    %             \draw [>->] (r'') -- (l);
     
    %             \draw [>->] (x) edge[bend left] node[pos=0.3,below]{!g''} (g);
    %             \draw [>->] (x) edge[bend right] node[midway,above]{!g'} (g);
    
    %         \end{tikzpicture}
    %     \end{center}

    %             \item[(3.3)]  
    %             Suppose, 
    %             \color{red}
    %             that \begin{flalign}
    %                 h_{R'R}(r') \mathop{\in} h_{KR}(K) 
    %                 \\
    %                 h_{C''C}(c'') \mathop{\in} h_{KC}(K)
    %             \end{flalign}
    %             Thus,
    %             there are $k_1, k_2 \mathop{\in} K$ such that $h_{KR}(k_1) \mathop{=} h_{R'R}(r')$ and $h_{KC}(k_2) \mathop{=} h_{C''C}(c'')$.
    %             \color{black}

    %             There is $k'' \mathop{\in} K''$ such that $h_{K''K}(k'') \mathop{=} k_2$ and $h_{K''C''}(k'') \mathop{=} c''$ becase $K''C''CK$ is a pullback square and $h_{KC}(k_2) \mathop{=} c \mathop{=} h_{C''C}(c'')$.

    %             We cannot have $k_1=k_2$. Suppose, by contradiction, that we have $k_1=k_2$. 
    %             \color{red} 
    %             Let $r'' \mathop{=} h_{K''R''}(k'')$. We have 
    %             \begin{flalign}
    %                 &h_{R''R}(r'')\\
    %                =&h_{R''R}(h_{K''R''}(k''))\\
    %                =&(h_{K''R''} \mathop{\star} h_{R''R})(k'') & \text{by definition of $r''$}\\
    %                =&(h_{K''K} \mathop{\star} h_{KR})(k'') &\text{by commutativity of $K''R''RK$}\\
    %                =&h_{KR}(h_{K''K}(k''))  \\
    %                =&h_{KR}(k_1) & \text{by definition of $k_1$} \\
    %                =&r
    %             \end{flalign}
    %             \color{black}
    %             % We have therefore $h_{K''C''}(k'') \mathop{=} c''$ because $h_{C''C}$ is injective, $h_{C''C}(c'') \mathop{=} c$ and 
    %             % \begin{flalign*}
    %             %    &h_{C''C}(h_{K''C''}(k'')) \\
    %             %   =&(h_{K''C''} \mathop{\star} h_{C''C})(k'')\\
    %             %   =&(h_{K''K} \mathop{\star} h_{KC})(k'')\\
    %             %   =&h_{KC}(h_{K''K}(k''))\\
    %             %   =&h_{KC}(k_2)\\
    %             %   =&c
    %             % \end{flalign*}
    %             and 
    %             \begin{flalign*}
    %               &h_{R''X}(r'')\\
    %              =&h_{R''X}(h_{K''R''}(k'')) \\
    %              =&(h_{K''R''} \mathop{\star} h_{R''X})(k'') \\
    %              =&(h_{K''C''} \mathop{\star} h_{C''X})(k'') & \text{by commutativity of $K''R''XC''$}\\
    %              =&h_{C''X}(h_{K''C''}(k''))\\
    %              =&h_{C''X}(c'')\\
    %              =&x
    %             \end{flalign*}


    %             \begin{flalign*}
    %                =&h_{R''X}(h_{K''R''}(k'')) \\
    %                =&(h_{K''R''} \mathop{\star} h_{R''X})(k'') \\
    %                =&(h_{K''C''} \mathop{\star} h_{C''X})(k'') & \text{by commutativity of $K''R''XC''$}\\
    %                =&h_{C''X}(h_{K''C''}(k''))\\
    %                =&h_{C''X}(c'')\\
    %                =&x
    %               \end{flalign*}

    %                     \begin{flalign*}
    %                         x &= h_{C''X}(c'')  & \text{by definition of $r''$}\\
    %                             &= h_{C''X}(h_{K'C''}(k')) & \text{by definition of $k'$}\\
    %                             &= (h_{K'C''} \mathop{\star} h_{C''X})(k') &  \\
    %                             &= (h_{K'R'} \mathop{\star} h_{R'X})(k') & \text{by the commutativity of $K'C'XR'$} \\
    %                             &= h_{R'X}(h_{K'R'}(k')) &  \\
    %                             &= h_{R'X}(r') &  \text{by definition of $k'$}\\
    %                             &= x& \text{by definition of $r'$}
    %                     \end{flalign*}
    %             which contradic the assumption Equation~\eqref{assumption_norpp}.

    %             Hence \begin{flalign}
    %                 k_1 \mathop{\neq} k_2 \label{k1_neq_k2}
    %             \end{flalign}

    %             From Equation~\eqref{k1_neq_k2}, we have 
    %             \begin{flalign}
    %                 h_{KR}(k_1) \mathop{\neq} h_{KR}(k_2) \label{hkrk1_neq_hkrk2}
    %             \end{flalign} 

    %             Since Equation~\eqref{hkrk1_neq_hkrk2} holds and $\rho$ is $X$-non-increasing, we have 
    %             \begin{flalign}
    %                 h_{KL}(k_1) \mathop{\neq} h_{KR}(k_2) \label{hkrk1_neq_hkrk2}
    %             \end{flalign} 

    %             By assumption on $\rho$ \todo{assumption: non-increasing}, we have $h_{KL}(k_2) \mathop{=} h_{R'L}(r') \mathop{=} l$ for some $k_2 \mathop{\in} K$.
                
    %             \begin{flalign*}
    %                 h_{XG}(x) &= h_{XG}(h_{C''X}(c'')) & \text{by definition of $c''$} \\
    %                         &= (h_{C''X} \mathop{\star} h_{XG})(c') &  \\
    %                         &= (h_{C''C} \mathop{\star} h_{CG})(c'') & \text{by commutativity of $C''CGX$} \\
    %                         &= h_{CG}(h_{C''C}(c'')) &  \\
    %                         &= h_{CG}(c) & \text{by definition of $c$} \\
    %                         &= h_{CG}(h_{KC}(k_1)) & \text{by definition of $k_1$} 
    %                         \\
    %                         &= (h_{KC} \mathop{\star} h_{CG})(k_1) & 
    %                         \\
    %                         &\mathop{\neq} (h_{KC} \mathop{\star} h_{CG})(k_2) & \text{by Equation~\eqref{k1_neq_k2} and injectivity of $h_{KC} \mathop{\star} h_{CG}$} 
    %                         \\
    %                         &= (h_{KL} \mathop{\star} h_{LG})(k_2) & 
    %                         \\
    %                         &= h_{LG}(h_{KL}(k_2)) & 
    %                         \\
    %                         &= h_{LG}(h_{R'L}(r')) & 
    %                         \\
    %                         &= (h_{R'L} \mathop{\star} h_{LG})(r') & 
    %                         \\
    %                         &= (h_{R'X} \mathop{\star} h_{XG})(r') & \text{by commutativity of $R'LGX$} \\
    %                         &= h_{XG}(h_{R'X}(r'))\\
    %                             &= h_{XG}(x) & \text{by definition of $r'$}
    %             \end{flalign*}
              
    %         \end{itemize}            
    \end{itemize}
\end{proof}
\subsection*{Proofs of~\autoref{subgraph_counting:lem:decomp_w_u},~\autoref{subgraph_counting:lem:xlnlmxrnr},~\autoref{subgraph_counting:lem:w_u_l_not_geq_r_not} and~\autoref{subgraph_counting:lem:w_g_geq_w_h_leq}}
\label{subgraph_counting:sec:appendix:a}
\begin*{\textbf{\autoref{subgraph_counting:lem:decomp_w_u}}}
\ \newline
\noindent
\begin{minipage}{0.7\textwidth}
    Let $X$ be a graph. For a pushout square as shown on the right, we have 
    % $\card{\operatorname{Mono}(X, B)} = \card{\operatorname{Mono}(X, D, \beta')}$ and $\card{\operatorname{Mono}(X, C, \lnot \beta)} = \card{\operatorname{Mono}(X, D, \lnot \beta', \alpha')}$.
    \begin{flalign*}
        \card{\operatorname{Mono}(X, B)} &= \card{\operatorname{Mono}(X, D, \beta')}
        \\
        \card{\operatorname{Mono}(X, C, \lnot \beta)} &= \card{\operatorname{Mono}(X, D, \lnot \beta', \alpha')}
    \end{flalign*}
    
\end{minipage}
\hfill
\begin{minipage}{0.3\textwidth}
    \hfill
    \begin{tikzpicture}
        \node (A) {$A$};
        \node [below of=A] (B) {$B$}; 
        \node [left of=A] (C) {$C$}; 
        \node [left of=B] (D) {$D$}; 
        \begin{scope}[nodes=rectangle]          
        \draw [>->] (A) to node [right,label,pos=0.5] {$\alpha$} (B);
        \draw [>->] (A) to node [above,label,pos=0.5] {$\beta$} (C);
        \draw [>->] (B) to node [below,label,pos=0.45] {$\beta'$} (D); 
        \draw [>->] (C) to node [right,label,pos=0.45] {$\alpha'$} (D);
        \end{scope}
    \end{tikzpicture}
\end{minipage}
\end*{}
\begin{proof}
    \label{subgraph_counting:proof:dcomp_w_u}
%    Since $\operatorname{Mono}(X, B) = \operatorname{Mono}(X, D, \beta')$ by definition of $\operatorname{Mono}(X, B)$ and $\operatorname{Mono}(X, D, \beta')$, the equality $\card{\operatorname{Mono}(X, B)} = \card{\operatorname{Mono}(X, D, \beta')}$ holds.
\begin{claim}
    $\card{\operatorname{Mono}(X, B)} = \card{\operatorname{Mono}(X, D, \beta')}$
\end{claim}
\begin{itemize} 
    \item Inclusion \(\operatorname{Mono}(X, B) \star \beta' \subseteq \operatorname{Mono}(X, D, \beta')\) holds by the definition of $\operatorname{Mono}(X, B)$ and $\operatorname{Mono}(X, D, \beta')$. 
    In fact, if \(\iota  \in \operatorname{Mono}(X, B) \star \beta' \), then there is a monomorphism \(\zeta : X \rightarrowtail B\) satisfying \(\iota = \zeta \star \beta'\). Since $\zeta \star \beta'$ is also a monomorphism\todo{monicity of $\beta'$}, \(\iota\) is also an element of \(\operatorname{Mono}(X, D, \beta')\).
    \item Inclusion \(\operatorname{Mono}(X, B) \star \beta' \supseteq \operatorname{Mono}(X, D, \beta')\) holds by the definition of $\operatorname{Mono}(X, B)$ and $\operatorname{Mono}(X, D, \beta')$.
    Suppose \(\iota \in \operatorname{Mono}(X, D, \beta')\). There is a monomorphism \(\zeta : X \rightarrowtail B\) satisfying \(\iota = \zeta \star \beta'\). It follows that \(\iota = \zeta \star \beta' \in \operatorname{Mono}(X, B) \star \beta'\).
    \item Hence, we have \(\operatorname{Mono}(X, B) \star \beta' = \operatorname{Mono}(X, D, \beta')\).
    \item $\card{\operatorname{Mono}(X, B)} \leq \card{\operatorname{Mono}(X, D, \beta')}$ follows by the injectivity of $\beta'$\todo{monicity of $\beta'$}.
\end{itemize}
\begin{claim}
     $\card{\operatorname{Mono}(X, C, \lnot \beta)} = \card{\operatorname{Mono}(X, D, \lnot \beta', \alpha')}$
   \end{claim}
    \begin{itemize}
        % \item The inclusion \(\operatorname{Mono}(X, C, \lnot \beta) \star \alpha'  \subseteq \operatorname{Mono}(X, D, \lnot \beta', \alpha')\) can be justified as follows. Suppose that \(\iota : X \to D\) is an element of \(\operatorname{Mono}(X, C, \lnot \beta) \star \alpha'\). According to the definition of \(\operatorname{Mono}(X, C, \lnot \beta) \star \alpha'\), there exists \(\eta : X \to C\) in \(\operatorname{Mono}(X, C, \lnot \beta)\) such that \(\iota = \eta \star \alpha'\). We need to show that no \(\zeta : X \to B\) exists such that \(\iota = \zeta \star \beta'\). Assuming the contrary, that such a \(\zeta\) exists. Under this assumption, the following commutative diagram holds:
        \item The inclusion \(\operatorname{Mono}(X, C, \lnot \beta) \star \alpha'  \subseteq \operatorname{Mono}(X, D, \lnot \beta', \alpha')\) can be justified as follows: Suppose \(
            \eta \star \alpha' \in \operatorname{Mono}(X, C, \lnot \beta) \star \alpha'\).
        %  By the definition of \(\operatorname{Mono}(X, C, \lnot \beta) \star \alpha'\), there exists a morphism \(\eta : X \to C\) in \(\operatorname{Mono}(X, C, \lnot \beta)\) such that \(\iota = \). 
         Since $\eta \star \alpha'$ is a monomorphism\todo{monicity of $\alpha'$}, it suffices to show that there is no \(\zeta : X \rightarrowtail B\) such that \(\eta \star \alpha' = \zeta \star \beta'\). Suppose, by contradiction, that such a \(\zeta\) exists, then the following commutative diagram holds:
        \begin{center}
            \begin{tikzpicture}[node distance=10mm]
                \node (A) {A}; 
                \node (B) [above right=of A] {B};
                \node (X) [left = of A] {X};
                \node (D) [below right=of B] {D};
                \node (C) [below right=of A] {C};
                \draw[>->] (A) to node[pos=.7, below] {$\alpha$} (B) ;
                \draw[>->] (C) to  node[pos=0.7, below] {$\alpha'$} (D);
                \draw[>->] (A) -- (C) node[pos=.4, right] {$\beta$};
                \draw[>->] (B) -- (D) node[pos=.4, right] {$\beta'$};
                \draw[>->] (X) -- node[above] {$\zeta$} (B);
                \draw[>->] (X) -- node[below] {$\eta$} (C);
            \end{tikzpicture}
        \end{center} 
        The pushout square \(\square ABDC\) is also a pullback square, by~\autoref{prop:pb_eq_po}. The universal property of the pullback provides a morphism \(\gamma : X \rightarrow A\) such that \(\eta = \gamma \star \beta\). \(\gamma\) is a monomorphism, because \(\eta = \gamma \star \beta\) and $\eta$ is a monomorphism. Therefore, the existence of $\gamma$ contradicts the assumption that \(\eta \in \operatorname{Mono}(X, C, \lnot \beta)\). Thus, \(\iota\) is also an element of \(\operatorname{Mono}(X, D, \lnot \beta', \alpha')\). 
        \item The inclusion \(\operatorname{Mono}(X, C, \lnot \beta) \star \alpha'  \supseteq \operatorname{Mono}(X, D, \lnot \beta', \alpha')\) can be justified as follows. Suppose that \(\iota : X \rightarrowtail D\) is an element of \(\operatorname{Mono}(X, D, \lnot \beta', \alpha')\). According to the definition of \(\operatorname{Mono}(X, D, \lnot \beta', \alpha')\), there exists \(\eta : X \rightarrowtail C\) making 
            \begin{flalign}
                \iota = \eta \star \alpha' \label{eq:etastaralphap}
            \end{flalign}
        We show that \(\eta\) is an element of 
        \(\operatorname{Mono}(X, C, \lnot \beta)\) by contradiction.
        
        Suppose that there is a monomorphism \(\zeta : X \rightarrowtail A\) such that 
        \begin{flalign}
            \eta = \zeta \star \beta \label{eq:zetastarbeta}
        \end{flalign}
        and observe:
        \begin{flalign*}
           \hspace{2cm} \iota &\overset{\operatorname{def}}{=} \eta \star \alpha' &\text{by \eqref{eq:etastaralphap}} 
           \\
                  &\overset{\operatorname{def}}{=} (\zeta \star \beta) \star \alpha' & \text{by \eqref{eq:zetastarbeta}}
                  \\
                  &= \zeta \star (\beta \star \alpha') & \text{by associativity} 
                  \\
                  &= \zeta \star (\alpha \star \beta') &\text{by commutativity of $ABDC$}\\
                  &= (\zeta \star \alpha) \star \beta'
        \end{flalign*}
          There is a contradiction because no such factorization of \(\iota\) should exist as \(\iota \in \operatorname{Mono}(X, D, \lnot \beta', \alpha')\).
        %    Therefore, there is no \(\zeta : X \rightarrow A\) such that \(\eta = \zeta \star \beta\), confirming \(\eta\) as an element of \(\operatorname{Mono}(X, C, \lnot \beta)\). Consequently, \(\iota = \eta \star \alpha'\) is an element of \(\operatorname{Mono}(X, C, \lnot \beta) \star \alpha'\).
        \item We conclude by the monicity of \(\alpha'\). \todo{monicity of $\alpha'$}
        % From \(\operatorname{Mono}(X, C, \lnot \beta) \star \alpha'  \subseteq \operatorname{Mono}(X, D, \lnot \beta', \alpha')\) and \(\operatorname{Mono}(X, C, \lnot \beta) \star \alpha'  \supseteq \operatorname{Mono}(X, D, \lnot \beta', \alpha')\), we conclude \(\operatorname{Mono}(X, C, \lnot \beta) \star \alpha' = \operatorname{Mono}(X, D, \lnot \beta', \alpha')\) and \(\card{\operatorname{Mono}(X, C, \lnot \beta) \star \alpha'} = \card{\operatorname{Mono}(X, D, \lnot \beta', \alpha')}\).
        % Since \(\alpha'\) is monic, \(\card{\operatorname{Mono}(X, C, \lnot \beta)} = \card{\operatorname{Mono}(X, C, \lnot \beta) \star \alpha'}\). Therefore, \(\card{\operatorname{Mono}(X, C, \lnot \beta)} = \card{\operatorname{Mono}(X, D, \lnot \beta', \alpha')}\).
    \end{itemize}
\end{proof} 



\noindent
\begin*{\textbf{\autoref{subgraph_counting:lem:xlnlmxrnr}}}
Let $X$ be a graph. Let $L \overset{l}{\leftarrow} K \overset{r}{\rightarrow} R$ be an injective DPO graph rewriting rule. We have 
\[
   \card{\operatorname{Mono}(X, L, \lnot l)}  - \card{\operatorname{Mono}(X, R, \lnot r)} 
   = 
   \card{\operatorname{Mono}(X, L)}  - \card{\operatorname{Mono}(X, R)} 
    \]
\end*{}
\begin{proof}
    \label{subgraph_counting:proof:lem:xlnlmxrnr}
    \begin{claim}
       $\card{\operatorname{Mono}(X, L, l)} = \card{\operatorname{Mono}(X, R, r)}$
    \end{claim}
    \begin{itemize}
        \item We prove $\card{\operatorname{Mono}(X, L, l)} \leq \card{\operatorname{Mono}(X, R, r)}$ by constructing an injection from $\operatorname{Mono}(X, L, l)$ to $\operatorname{Mono}(X, R, r)$. Let $h \in \operatorname{Mono}(X, L, l)$. By definition, we have $h = g \star l$ for some $g: X \rightarrowtail K$. Thus, $g \star r \in \operatorname{Mono}(X, L, r)$. Let $h' \in \operatorname{Mono}(X, L, l)$ such that $h \neq h'$. We have $h' = g' \star l$ for some $g':X \rightarrowtail K$. We have $g \neq g'$ because otherwise we would have $h = g \star l = g' \star l = h'$. From the monicity of $r$, we conclude $g \star r \neq g' \star r$.
        \item Analoguously, we can prove $\card{\operatorname{Mono}(X, L, l)} \geq \card{\operatorname{Mono}(X, R, r)}$ by constructing an injection from $\operatorname{Mono}(X, R, r)$ to $\operatorname{Mono}(X, L, l)$.
        %  Let $h \in \operatorname{Mono}(X, R, r)$. By definition, we have $h = g \star r$ for some $g: X \rightarrowtail K$. Thus, $g \star l \in \operatorname{Mono}(X, L, l)$.
        % Let $h' \in \operatorname{Mono}(X, R, r)$ such that $h \neq h'$. We have $h' = g' \star l$ for some $g':X \rightarrowtail K$. Since $r$ is monic, we deduce $g \neq g'$. From the monicity of $l$, we conclude $g \star r \neq g' \star r$.
    \end{itemize}


    % we have \[\card{\operatorname{Mono}(X, L, l)} - \card{\operatorname{Mono}(X, R, r)} = 0 \] becase 
    
    Hence, the following equality holds
    \begin{flalign*}
         &\card{\operatorname{Mono}(X, L, \lnot l)}  - \card{\operatorname{Mono}(X, R, \lnot r)}\\
        =&\card{\operatorname{Mono}(X, L) \setminus \operatorname{Mono}(X, L, l)}  - 
            \card{\operatorname{Mono}(X, R) \setminus \operatorname{Mono}(X, R, r)}\\
        =&(\card{\operatorname{Mono}(X, L)} - \card{\operatorname{Mono}(X, L, l)}) - (\card{\operatorname{Mono}(X, R)} - \card{\operatorname{Mono}(X, R, r)}) \\
        =&(\card{\operatorname{Mono}(X, L)} - \card{\operatorname{Mono}(X, R)})
         - (\card{\operatorname{Mono}(X, L, l)} - \card{\operatorname{Mono}(X, R, r)})\\
        = &\card{\operatorname{Mono}(X, L)} - \card{\operatorname{Mono}(X, R)}
    \end{flalign*}
\end{proof}

\noindent
\begin*{\textbf{\autoref{subgraph_counting:lem:w_u_l_not_geq_r_not}}}
\ \newline 
    \noindent 
    \begin{minipage}{0.7\textwidth}
        Let $X$ be a ruler-graph and $\rho = (L \overset{l}{\leftarrowtail} K \overset{r}{\rightarrowtail} R)$ an injective DPO rewriting rule.
        Suppose that $\rho$ is $X$-non-increasing. For every rewriting step induced by the DPO diagram (shown on the right), we have 
        \[
            |\operatorname{Mono}(X, G, \lnot m, \lnot l')| \geq |\operatorname{Mono}(X, H, \lnot m', \lnot r')|
        \]
    \end{minipage}
    \hfill
    \begin{minipage}{0.29\textwidth}
        \hfill
        \resizebox{0.9\textwidth}{!}{
                \begin{tikzpicture}
            \node (k) at (0,1) {K};
            \node (l) at (-2,1) {L};
            \node (r) at (2,1) {R};
            \node (c) at (0,-1) {C};
            \node (g) at (-2,-1) {G};
            \node (h) at (2,-1) {H};
            % \node (rb) at (1.5,-0.5) {$R_X$};
            % \node (h') at (1.5,-1.5) {$H'$};
            % \draw[>->]  (rb) -- (h') node [midway,above] {};
            % \draw[>->]  (c) -- (h') node [midway,above] {};
            \draw[<-<]  (l) -- (k) node [midway,below] {$l$};
            \draw[>->]  (k) -- (r) node [midway,below] {$r$};
            \draw[>->] (c) -- (g) node [midway, above] {$l'$};
            \draw[>->] (c) -- (h) node [midway,above] {$r'$};
            \draw[>->] (l) -- (g) node[midway, right] {$m$};
            \draw[>->] (r) -- (h) node[midway, left] {$m'$};
            \draw[>->] (k) -- (c) node[midway, left] {};
            \node () [at=($(l)!0.5!(c)$)] {$PO$};
            \node () [at=($(r)!0.5!(c)$)] {$PO$};
            % \draw[>->] (rb) to (r);
            % \draw[->] (rb) to (l);
            % \draw[<-<] (rb) to (k);
        \end{tikzpicture}
        }
    \end{minipage}
\end*{}
\begin{proof}
    \label{proof:lem:w_u_l_not_geq_r_not}
    We are going to show that for two arbitrary distinct monomorphisms $h_{XH}', h_{XH}'' \in \operatorname{Mono}(X, H, \lnot m', \lnot r')$, there are two distinct monomorphisms $h_{XG}', h_{XG}'' \in \operatorname{Mono}(X, G, \lnot m, \lnot l')$.

    Let $h_{XH}'\in \operatorname{Mono}(X, H, \lnot m', \lnot r')$ be a monomorphism. In the category \textbf{Graph}, we can construct the following commutative diagram
    \begin{center}
        \resizebox{6cm}{!}{
            \begin{tikzpicture}
                \node (k) at (0,0) {K};
                \node (r) at (4,0) {R};
                \node (c) at (0,-3) {C};
                \node (h) at (4,-3) {H};
                % \node (rb) at ($\scl*(1.5,-0.5)$) {$R_X$};
                
                % \node (h') at ($\scl*(1.5,-1.5)$) {$H'$};
                % \draw[>->]  (rb) -- (h') node [midway,above] {};
                % \draw[>->]  (c) -- (h') node [midway,above] {};
    
                \draw[<-<]  (r) -- (k) node [midway,above] {};
                \draw[>->] (c) -- (h) node [midway, below] {};
                \draw[>->] (r) -- (h) node[midway, left] {};
                \draw[>->] (k) -- (c) node[midway, left] {};
    
                % \draw[->] (rb) to (l);
                % \draw[<-<] (rb) to (k);
            
                \node (k') at (1,-1) {K'};
                \node (r') at (2,-1) {R'};
                \node (c') at (1,-2) {C'};
                \node () at (1.5,-1.5) {$PB$};
                \node () at (3,-1.5) {$PB$};
                \node () at (1.5,-2.5) {$PB$};
                \node (x) at (2,-2) {X};
                % \draw [->] (x) -- (h') node[midway] {!};
                \draw [>->] (c') -- (x);
                \draw [>->] (r') -- (x);
                \draw [>->] (k') -- (r');
                \draw [>->] (k') -- (c');
                \draw [>->] (c') -- (c);
                % \draw [->] (r') -- (rb);
                % \draw [->] (r') -- (rb);
                % \draw [->] (h') -- (h);
    
                \draw[>->] (r') -- (r);
                \draw[>->] (x) -- (h) node[midway,right] {$h$};
                % \node (rb) at ($\scl*(\sclx*1.5,-0.5)$) {$R_X$};
                % \node (h') at ($\scl*(\sclx*1.5,-1.2)$) {$H'$};
                % \draw[>->]  (rb) -- (h') node [midway,above] {};
                % \draw[>->]  (c) -- (h') node [midway,above] {};
                % \draw[>->]  (rb) -- (h') node [midway,above] {};
                % \draw[>->] (rb) to (r);
                % \draw[<-<] (rb) to (k);
            \end{tikzpicture}
        }
        \end{center} 
    
    $KCHR$ is also pullback, by~\autoref{prop:pb_eq_po}, because it is pushout. 
    Therefore, since $KCHR$ is also pullback and \(
        h_{K'C'} \star h_{C'C} \star h_{CH} =h_{K'R'} \star h_{R'R} \star h_{R'R}
    \) holds, there exists $h_{K'K}:K' \rightarrowtail K$ such that 
    \begin{flalign}
        h_{K'C'} \star h_{C'C} &= h_{K'K} \star h_{KC} \label{kpcpcpckpkkc}
        \\
        h_{K'R'} \star h_{R'R} &= h_{K'K} \star h_{KR} \label{kprprprkpk}
    \end{flalign} 
    
    The square $K'KCC'$ and $K'R'RK$ are pullbacks, by~\autoref{kpcpck_pullback}. Hence, the following commutative diagram holds
    \begin{center}
        \resizebox{6cm}{!}{
            \begin{tikzpicture}
                \node (k) at (0,0) {K};
                \node (r) at (4,0) {R};
                \node (c) at (0,-3) {C};
                \node (h) at (4,-3) {H};
                % \node (rb) at ($\scl*(1.5,-0.5)$) {$R_X$};
                % \node (h') at ($\scl*(1.5,-1.5)$) {$H'$};
                % \draw[>->]  (rb) -- (h') node [midway,above] {};
                % \draw[>->]  (c) -- (h') node [midway,above] {};
                \draw[<-<]  (r) -- (k) node [midway,above] {};
                \draw[>->] (c) -- (h) node [midway, below] {};
                \draw[>->] (r) -- (h) node[midway, left] {};
                \draw[>->] (k) -- (c) node[midway, left] {};
                % \draw[->] (rb) to (l);
                % \draw[<-<] (rb) to (k);
                \node (k') at (1,-1) {K'};
                \node (r') at (2,-1) {R'};
                \node (c') at (1,-2) {C'};
                \node () at (1.5,-1.5) {$PB$};
                \node () at (3,-1.5) {$PB$};
                \node () at (1.5,-2.5) {$PB$};
                \node () at (1.5,-0.5) {$PB$};
                \node () at (0.5,-1.5) {$PB$};
                \node (x) at (2,-2) {X};
                % \draw [->] (x) -- (h') node[midway] {!};
                \draw [>->] (c') -- (x);
                \draw [>->] (r') -- (x);
                \draw [>->] (k') -- (r');
                \draw [>->] (k') -- (c');
                \draw [>->] (c') -- (c);
                % \draw [->] (r') -- (rb);
                % \draw [->] (r') -- (rb);
                % \draw [->] (h') -- (h);
                \draw[>->] (r') -- (r);
                \draw[>->] (x) -- (h) node[midway,right] {$h$};
                % \node (rb) at ($\scl*(\sclx*1.5,-0.5)$) {$R_X$};
                % \node (h') at ($\scl*(\sclx*1.5,-1.2)$) {$H'$};
                % \draw[>->]  (rb) -- (h') node [midway,above] {};
                % \draw[>->]  (c) -- (h') node [midway,above] {};
                % \draw[>->]  (rb) -- (h') node [midway,above] {};
                % \draw[>->] (rb) to (r);
                % \draw[<-<] (rb) to (k);
                \draw[>->] (k') -- (k) ;
            \end{tikzpicture}
        }
        \end{center} 

    \todo{use: assumption: non increasing}
    Since the rule is by assumption $X$-non-increasing, there is a morphism $h_{R'L}$ such that $K'KLR'$ is commutative.
    
    $K'C'XR'$ is a pushout by~\autoref{kpcpxrp_po}. 
    
    Thus, the following commutative diagram holds
    \begin{center}
        \resizebox{6cm}{!}{
            \begin{tikzpicture}
                \node (k) at (0,0) {K};
                \node (r) at (4,0) {L};
                \node (c) at (0,-3) {C};
                \node (h) at (4,-3) {G};
                % \node (rb) at ($\scl*(1.5,-0.5)$) {$R_X$};
                
                % \node (h') at ($\scl*(1.5,-1.5)$) {$H'$};
                % \draw[>->]  (rb) -- (h') node [midway,above] {};
                % \draw[>->]  (c) -- (h') node [midway,above] {};
    
                \draw[<-<]  (r) -- (k) node [midway,above] {};
                \draw[>->] (c) -- (h) node [midway, below] {};
                \draw[>->] (r) -- (h) node[midway, left] {};
                \draw[>->] (k) -- (c) node[midway, left] {};
    
                % \draw[->] (rb) to (l);
                % \draw[<-<] (rb) to (k);
            
                \node (k') at (1,-1) {K'};
                \node (r') at (2,-1) {R'};
                \node (c') at (1,-2) {C'};
                \node () at (1.5,-1.5) {$PO$};
                % \node () at (1.5,-0.5) {$PB$};
                \node () at (0.5,-1.5) {$PB$};
                % \node () at (3,-1.5) {$PB$};
                % \node () at (1.5,-2.5) {$PB$};
                \node (x) at (2,-2) {X};
                % \draw [->] (x) -- (h') node[midway] {!};
                \draw [>->] (c') -- (x);
                \draw [>->] (r') -- (x);
                \draw [>->] (k') -- (r');
                \draw [>->] (k') -- (c');
                \draw [>->] (c') -- (c);
                \draw [->] (k') -- (k);
                % \draw [->] (r') -- (rb);
                % \draw [->] (r') -- (rb);
                % \draw [->] (h') -- (h);
    
                \draw[>->] (r') -- (r);
                % \draw[>->] (x) -- (h) node[midway,right] {$h$};
                % \node (rb) at ($\scl*(\sclx*1.5,-0.5)$) {$R_X$};
                % \node (h') at ($\scl*(\sclx*1.5,-1.2)$) {$H'$};
                % \draw[>->]  (rb) -- (h') node [midway,above] {};
                % \draw[>->]  (c) -- (h') node [midway,above] {};
                % \draw[>->]  (rb) -- (h') node [midway,above] {};
                % \draw[>->] (rb) to (r);
                % \draw[<-<] (rb) to (k);
            \end{tikzpicture}
        }
        \end{center} 

    By~\autoref{lem:g_monic}, there is a unique monomorphism $h_{XG}':X \rightarrowtail G$ such that $h_{C'X} \star h_{XG}' = h_{C'C} \star h_{CG}$ and $h_{R'X} \star h_{XG}' = h_{R'L} \star h_{LG}$.

    The monomorphism $h_{XG}'$ is in the set $\operatorname{Mono}(X, G, \lnot m, \lnot l')$, by $h_{XH}'\in \operatorname{Mono}(X, H, \lnot m', \lnot r')$ and~\autoref{subgraph_counting:def:creates_more_x_on_the_left}. 
    Intuitively, $h_{XH}'\in \operatorname{Mono}(X, H, \lnot m', \lnot r')$ implies that some elements in $X$ are mapped onto $h_{RH}(R) \setminus (r \star h_{RH})(K)$. By the first assumption of~\autoref{subgraph_counting:def:creates_more_x_on_the_left}, these elements are mapped onto $h_{LG}(L) \setminus (l \star h_{LG})(K)$.

    
    The monomorphism \( h_{XG}' \) belongs to the set \( \operatorname{Mono}(X, G, \lnot m, \lnot l') \), due to \( h_{XH}' \in \operatorname{Mono}(X, H, \lnot m', \lnot r') \) and~\autoref{subgraph_counting:def:creates_more_x_on_the_left}. Intuitively, \( h_{XH}' \in \operatorname{Mono}(X, H, \lnot m', \lnot r') \) implies that certain elements of \( X \) are mapped into \( h_{RH}(R) \setminus (r \star h_{RH})(K) \). By the second assumption of~\autoref{subgraph_counting:def:creates_more_x_on_the_left}, these elements are also mapped into \( h_{LG}(L) \setminus (l \star h_{LG})(K) \).

     \todo{use; def non-increase; condition 1}

    Let $h_{XH}'':X \rightarrowtail H$ be a monomorphism such that 
    \begin{flalign}
        h_{XH}' \neq h_{XH}'' \label{hpneqh}
    \end{flalign}
    
    In the category \textbf{Graph}, we can construct the following commutative diagram:
    \begin{center}
        \resizebox{6cm}{!}{
            \begin{tikzpicture}
                \node (k) at (0,0) {K};
                \node (r) at (4,0) {L};
                \node (c) at (0,-3) {C};
                \node (h) at (4,-3) {G};
                % \node (rb) at ($\scl*(1.5,-0.5)$) {$R_X$};
                
                % \node (h') at ($\scl*(1.5,-1.5)$) {$H'$};
                % \draw[>->]  (rb) -- (h') node [midway,above] {};
                % \draw[>->]  (c) -- (h') node [midway,above] {};
    
                \draw[<-<]  (r) -- (k) node [midway,above] {};
                \draw[>->] (c) -- (h) node [midway, below] {};
                \draw[>->] (r) -- (h) node[midway, left] {};
                \draw[>->] (k) -- (c) node[midway, left] {};
    
                % \draw[->] (rb) to (l);
                % \draw[<-<] (rb) to (k);
            
                \node (k') at (1,-1) {K'};
                \node (r') at (2,-1) {R'};
                \node (c') at (1,-2) {C'};
                \node () at (1.5,-1.5) {$PO$};
                % \node () at (1.5,-0.5) {$PB$};
                \node () at (0.5,-1.5) {$PB$};
                % \node () at (3,-1.5) {$PB$};
                % \node () at (1.5,-2.5) {$PB$};
                \node (x) at (2,-2) {X};
                % \draw [->] (x) -- (h') node[midway] {!};
                \draw [>->] (c') -- (x);
                \draw [>->] (r') -- (x);
                \draw [>->] (k') -- (r');
                \draw [>->] (k') -- (c');
                \draw [>->] (c') -- (c);
                \draw [->] (k') -- (k);
                % \draw [->] (r') -- (rb);
                % \draw [->] (r') -- (rb);
                % \draw [->] (h') -- (h);
    
                \draw[>->] (r') -- (r);
                % \draw[>->] (x) -- (h) node[midway,right] {$h$};
                % \node (rb) at ($\scl*(\sclx*1.5,-0.5)$) {$R_X$};
                % \node (h') at ($\scl*(\sclx*1.5,-1.2)$) {$H'$};
                % \draw[>->]  (rb) -- (h') node [midway,above] {};
                % \draw[>->]  (c) -- (h') node [midway,above] {};
                % \draw[>->]  (rb) -- (h') node [midway,above] {};
                % \draw[>->] (rb) to (r);
                % \draw[<-<] (rb) to (k);
            \end{tikzpicture}
        }
        \end{center} 

    Analoguously, there exists a unique monomorphism $h_{XG}'':X \rightarrowtail G$ such that $h_{C''X} \star h_{XG} ''= h_{C''C} \star h_{CG}$ and $h_{R''X} \star h_{XG}'' = h_{R''L} \star h_{LG}$.
    
    By \eqref{hpneqh} and~\autoref{lem:h_hp_diff_g_gp_diff}, we have $h_{XG} \neq h_{XG}'$.
\end{proof}



\noindent
\begin*{\textbf{\autoref{subgraph_counting:lem:w_g_geq_w_h_leq}}}[Decreasing step]
Let $\rho = (L \overset{l}{\leftarrowtail} K \overset{r}{\rightarrowtail} R)$ be an injective DPO rewriting rule,
\( \mathbb{X} \) a set of ruler-graphs,
\( s_{\mathbb{X}} \colon \mathbb{X} \to \mathbb{N} \) a weight function,
and \( G \Rightarrow_{\rho,\mathfrak{M}} H \) a rewriting step. 
If $\rho$ is \( X \)-non-increasing for every ruler-graph \( X \in \mathbb{X} \), then:\[
    w_{s_\mathbb{X}}(G) - w_{s_\mathbb{X}}(H) 
    \geq 
    w_{s_\mathbb{X}}(L) - w_{s_\mathbb{X}}(R)
\]
\end*{}
\begin{proof}
    \label{subgraph_counting:proof:lem:w_g_geq_w_h_leq}
    Let the following diagram be the witness diagram of the rewriting step
    \begin{center}
        \begin{tikzpicture}
            \node (k) at (0,1) {K};
            \node (l) at (-2,1) {L};
            \node (r) at (2,1) {R};
            \node (c) at (0,-1) {C};
            \node (g) at (-2,-1) {G};
            \node (h) at (2,-1) {H};
            \draw[<-<]  (l) -- (k) node [midway,below] {$l$};
            \draw[>->]  (k) -- (r) node [midway,below] {$r$};
            \draw[>->] (c) -- (g) node [midway, above] {$l'$};
            \draw[>->] (c) -- (h) node [midway,above] {$r'$};
            \draw[>->] (l) -- (g) node[midway, right] {$m$};
            \draw[>->] (r) -- (h) node[midway, left] {$m'$};
            \draw[>->] (k) -- (c) node[midway, left] {};
            \node () [at=($(l)!0.5!(c)$)] {$PO$};
            \node () [at=($(r)!0.5!(c)$)] {$PO$};
        \end{tikzpicture}
    \end{center}
    Let $X \in \mathbb{X}$.

    We have 
    \begin{flalign*}
        & \card{\operatorname{Mono}(X,G)}
            % \overset{\operatorname{def}}{=} 
            % w_{T_\Sigma}(t_G) 
        \\
        % =& \Psi_{G}\\
        =& 
        \card{\operatorname{Mono}(X,G,l')
        \uplus
         \operatorname{Mono}(X,G,\lnot l',m)
        \uplus
         \operatorname{Mono}(X,G,\lnot l',\lnot m)}\\
        =& 
        \card{\operatorname{Mono}(X,G,l')}
        +
         \card{\operatorname{Mono}(X,G,\lnot l',m)}
        +
        \card{\operatorname{Mono}(X,G,\lnot l',\lnot m)}
        \\
    =& \card{\operatorname{Mono}(X,C)}
            +
            \card{\operatorname{Mono}(X, L, \lnot l)} 
            +
            \card{\operatorname{Mono}(X, G,\lnot l', \lnot m)} 
        &  \text{by~\autoref{subgraph_counting:lem:decomp_w_u}}
    \end{flalign*}
    and
    \begin{flalign*}
        % ligne 1
        &\card{\operatorname{Mono}(X,H)}
        % \overset{\operatorname{def}}{=}  
        %     w_{T_\Sigma}(t_H )
        \\
        =& 
        \card{\operatorname{Mono}(X,H,r')
        \uplus
            \operatorname{Mono}(X,H,\lnot r',m')
        \uplus
            \operatorname{Mono}(X,H,\lnot r',\lnot m')}
        \\
        =& 
        \card{\operatorname{Mono}(X,H,r')}
        +
         \card{\operatorname{Mono}(X,H,\lnot r',m')}
        +
        \card{\operatorname{Mono}(X,H,\lnot r',\lnot m')}
        \\
        % = & \Psi_{H} \\
        % ligne 2
        = &
            \card{\operatorname{Mono}(X,C)}
            +
            \card{\operatorname{Mono}(X, R, \lnot r)} 
            +
            \card{\operatorname{Mono}(X, H, \lnot r', \lnot m')}
        % ( 
        %             w_{T_\Sigma} (h_{RH} \star t_H ) + w_{T_\Sigma} (h_{CH} \star t_H  -h_{KC}) 
        % )
        %         +  w_{T_\Sigma}(t_H  - \{h_{RH},h_{CH} \})  
        &  \text{by~\autoref{subgraph_counting:lem:decomp_w_u}}
    \end{flalign*}
    Therefore, the following inequality holds
     \begin{flalign*}
           &\card{\operatorname{Mono}(X,G)} - \card{\operatorname{Mono}(X,H)} &\\
           =&(\card{\operatorname{Mono}(X, L, \lnot l)} - \card{\operatorname{Mono}(X, R, \lnot r)}) + 
            (\card{\operatorname{Mono}(X, G, \lnot l', \lnot m)} - \card{\operatorname{Mono}(X, H, \lnot r', \lnot m')}) \\
       \geq& \card{\operatorname{Mono}(X, L, \lnot l)} - \card{\operatorname{Mono}(X, R, \lnot r)} \hspace{3cm} 
       \text{by~\autoref{subgraph_counting:lem:w_u_l_not_geq_r_not}} \\
          =& \card{\operatorname{Mono}(X, L)} - \card{\operatorname{Mono}(X, R)} \hspace{3cm} 
          \text{by~\autoref{subgraph_counting:lem:xlnlmxrnr}}
     \end{flalign*}
     Thus, the following inequality holds
     \begin{flalign*}
          &w_{s_\mathbb{X}}(G) - w_{s_\mathbb{X}}(H)
          \\
         =&\sum_{X \in \mathbb{X}}^{}w(X) * |\operatorname{Mono}(X,G)| - \sum_{X \in \mathbb{X}}^{}w(X) * |\operatorname{Mono}(X,H)|
         \\
         =&\sum_{X \in \mathbb{X}}^{}w(X) * \left( \card{\operatorname{Mono}(X,G)} -  \card{\operatorname{Mono}(X,H)} \right)
         \\
         \geq&\sum_{X \in \mathbb{X}}^{}w(X) * \left(  \card{\operatorname{Mono}(X,L)}- \card{\operatorname{Mono}(X,R)} \right)
         \\
         = &\sum_{X \in \mathbb{X}}^{}w(X) * \card{\operatorname{Mono}(X,L)} -  \sum_{X \in \mathbb{X}}^{}w(X) * \card{\operatorname{Mono}(X,R)} 
         \\
         = & w_{s_\mathbb{X}}(L) - w_{s_\mathbb{X}}(R)
     \end{flalign*}
    \qed
\end{proof}

\noindent
\begin*{\textbf{\autoref{subgraph_counting:thm:termination_grs}}}[Termination]
    Let \(\mathcal{A}\) and \(\mathcal{B}\) be sets of injective DPO rewriting rules, $\mathbb{X}$ a set of ruler-graphs and $s_\mathbb{X}$ a weight function. If the following hold
    \begin{enumerate}
        \item\label{subgraph_counting:thm:termination_grs:assump:1} $\rho$ is $X$-non-increasing for every rule $\rho \in \mathcal{A} \cup \mathcal{B}$ and for every ruler-graph $X \in \mathbb{X}$,
        \item\label{thm:termination_grs:assump:2} \( w_{s_\mathbb{X}}(lhs(\rho)) > w_{s_\mathbb{X}}(rhs(\rho)) \) for every rule \(\rho \in \mathcal{A}\),
        \item\label{thm:termination_grs:assump:3} \( w_{s_\mathbb{X}}(lhs(\rho)) \geq w_{s_\mathbb{X}}(rhs(\rho)) \) for every rule \(\rho \in \mathcal{B}\).
    \end{enumerate}
    Then \(\Rightarrow_{\mathcal{A},\mathcal{M}}\) terminates relative to \(\Rightarrow_{\mathcal{B},\mathcal{M}}\).
\end*{}
\begin{proof}
    \label{subgraph_counting:proof:thm:termination_grs}
    
    By assumption \eqref{subgraph_counting:thm:termination_grs:assump:1} and ~\autoref{subgraph_counting:lem:w_g_geq_w_h_leq}, the inequality \(
        w_{s_\mathbb{X}}(G) - w_{s_\mathbb{X}}(H) 
        \geq 
        w_{s_\mathbb{X}}(\operatorname{lhs}(\rho)) - w_{s_\mathbb{X}}(\operatorname{rhs}(\rho))
    \) holds for all rules $\rho \in \mathcal{A} \cup \mathcal{B}$ and for all rewriting steps $G \Rightarrow_{\rho, \mathcal{M}} H$.
    
    \noindent From assumptions \eqref{thm:termination_grs:assump:2} and \eqref{thm:termination_grs:assump:3}, we deduce 
    \begin{itemize}
        \item \( w_{s_\mathbb{X}}(G) > w_{s_\mathbb{X}}(H) \) for every rule \(\rho \in \mathcal{A}\);
        \item  \( w_{s_\mathbb{X}}(G) \geq w_{s_\mathbb{X}}(H) \) for every rule \(\rho \in \mathcal{B}\).
    \end{itemize}
    Since $w_{s_\mathbb{X}}(G) \in \mathbb{N}$ for all graph $G$. Every rewriting chain can only have a finite number of rewriting steps using rules in $\mathcal{A}$.
\end{proof} 


%%%%%%%%%%%%%%%%%%%%%%%%%%%%%%%%%%%%%%%%%%%%%%%%%%%%%%%%%%%
%%%   Chapter : morphism counting Antipattern  %%%%%%%%%%%%%%
%%%%%%%%%%%%%%%%%%%%%%%%%%%%%%%%%%%%%%%%%%%%%%%%%%%%%%%%%%%

\chapter{Termination of Injective DPO Graph Rewriting
Systems using Morphism Counting with antipatterns}
\label{chap:antipattern}
\section{Introduction}
\label{antipattern:sec:intro}
Chapter~\ref{chap:antipattern} extends the subgraph counting method presented in Chapter~\ref{chap:subgraph_counting}. 
Consider the DPO rewriting system on edge-labeld directed multi-graphs shown
below.
%  \begin{figure}[H]
%     \centering
\begin{center}
    \begin{tikzpicture}
      \graphbox{$L$}{0mm}{0mm}{34mm}{20mm}{2mm}{-5mm}{
          \coordinate (o) at (0mm,-3mm); 
          \node[draw,circle] (l1) at ($(o)+(-10mm,0mm)$) {1};
          \node[draw,circle] (l2) at ($(l1)+(2,0)$) {2};
          \node[draw,circle] (l3) at ($(l1) + (1,0)$) {3};
          \draw[->] (l1) -- (l3) node[midway,above] {$a$};
          \draw[->] (l3) -- (l2) node[midway,above] {$a$};
      }     
      \graphbox{$K$}{40mm}{0mm}{24mm}{20mm}{2mm}{-5mm}{
          \coordinate (o) at (5mm,-3mm); 
          \node[draw,circle] (l1) at ($(o)+(-10mm,0mm)$) {1};
          \node[draw,circle] (l2) at ($(l1)+(1,0)$) {2};
      }    
      \graphbox{$R$}{70mm}{0mm}{45mm}{20mm}{2mm}{-5mm}{
        \coordinate (o) at (0mm,-3mm); 
        \node[draw,circle] (l1) at ($(o)+(-10mm,0mm)$) {1};
        \node[draw,circle] (l2) at ($(l1)+(2,0)$) {2};
        \node[draw,circle] (l3) at ($(l1) + (1,0)$) {3};
        \draw[->] (l1) -- (l3) node[midway,above] {$a$};
        \draw[->] (l3) -- (l2) node[midway,above] {$a$};
        \draw[->] (l3) edge [loop below] node {$c$} (l3);
      }    

      \node () at (37mm,-10mm) {$\leftarrowtail$};
      \node () at (67mm,-10mm) {$\rightarrowtail$};
  \end{tikzpicture}
%   \caption{}
%   \label{fig:intro:graph_transformation_rule_anti_pattern_}
%  \end{figure} 
\end{center}
%  in Figure~\ref{fig:intro:graph_transformation_rule_anti_pattern_}. 
In this system, the number of injective graph homomorphisms from every pattern graph increases with each application of the rewriting rule. Therefore, the method introduced in Chapter~\ref{chap:subgraph_counting} cannot be applied. However, the number of injective graph homomorphisms from the graph
\raisebox{2pt}{
            \scalebox{0.7}{\tikz[baseline=-0.5ex]{
            \node [draw,circle] (z) at (-1,0) {};
            \node [draw,circle] (x) at (0,0) {};
            \node[draw,circle] (y) at (1,0) {};
            \draw[->] (z)--(x) node[midway, above] {$a$};
            \draw[->] (x)--(y) node[midway, above] {$a$};
        }}} 
    whose images are not included in occurrences of the graph
        \raisebox{2pt}{
            \scalebox{0.7}{\tikz[baseline=-0.5ex]{
            \node [draw,circle] (z) at (-1,0) {};
            \node [draw,circle] (x) at (0,0) {};
            \node[draw,circle] (y) at (1,0) {};
            \draw[->] (z)--(x) node[midway, above] {$a$};
            \draw[->] (x)--(y) node[midway, above] {$a$};
            \draw[->] (x) edge [loop below] node {$c$} (l3);
        }}}
decreases with each application of the rewriting rule.

This extension can handle systems (e.g. Example~\ref{antipattern:ex:grs_aca} and Example~\ref{antipattern:ex:endrullis:d3:termination}) which cannot be handled by the
prior interpretation-based approaches~\cite{zantema2014termination,bruggink2014termination,bruggink2015proving,endrullis2024generalized_arxiv_v2,overbeek2024termination_lmcs} and the morphism counting method presented in Chapter~\ref{chap:subgraph_counting}.
An implementation of this extension is also provided in our tool~\textbf{LyonParallel}.
  
We are going to proceed as follows:~\textsection~\ref{antipattern:sec:termination_criterion} presents the extension.~\textsection~\ref{antipattern:sec:examples} provides some examples,~\textsection~\ref{antipattern:sec:conclusion} concludes. Missing proofs and definitions are in the appendix of this chapter.
  
\section{A Termination Criterion} 
\label{antipattern:sec:termination_criterion}
% We define graph weight in Section\ref{sec:subgraph_counting_antipattern:interpretation} and introduce non-increasing rules in Section~\ref{sec:subgraph_counting_antipattern:non-increasing}. Using these concepts, Section~\ref{sec:subgraph_counting_antipattern:termination} outlines the methodology and termination criterion.

% \subsection{Generalized ruler-Graphs}
% \label{sec:subgraph_counting_antipattern:interpretation}
We generalize the concept of the ruler-graph presented in Chapter~\ref{chap:subgraph_counting}.

\begin{definition}
    \label{antipattern:def:ruler_graph}
    A \textbf{ruler-graph}\index{Ruler-graph} is an order pair \( (X, \mathcal{C}) \) where $X$ is the \textbf{underlying graph} and $\mathcal{C}$ is a set of monomorphisms from $X$ with $\card{\mathcal{C}}\mathop{\leq} 1$. Each $f \in \mathcal{C}$ is called a \textbf{forbidden context}\index{Forbidden context}.
\end{definition}
In this chapter we restrict attention to the case $\card{\mathcal{C}}\mathop{\leq} 1$, and defer an analysis for general finite $\mathcal{C}$ to future work.

For $f:X \rightarrowtail F \in \mathcal{C}$, we also call $F$ a forbidden context for simplicity.
The definition of measurement needs some modifications to take the forbidden context into account. When $\mathcal{C} \mathop{=} \emptyset$, the definition of measurement is the same as that in Definition~\ref{subgraph_counting:def:measurement} in Chapter~\ref{chap:subgraph_counting}. When $\mathcal{C} = \set{f}$ with $f:X \rightarrowtail F$, we exclude those monomorphisms whose images overlap with the forbidden context. These excluded monomorphisms are those that can be extended to a morphism $g: F \rightarrowtail G$. This leads us to the following definition.
\begin{definition} 
    \label{antipattern:def:measurement}
    \index{Measurement}
    Let \( \mathcal{X} = (X,\mathcal{C}) \) be a ruler-graph and \( G \) a graph. The \textbf{measurement} of \( G \) with respect to \( \mathcal{X}\), denoted \( m_\mathcal{X}(G) \), is defined as 
    \[
        m_\mathcal{X}(G) \isdef 
        \card{\{h \mathop{\in} \operatorname{Mono}(X,G) \mathop{\mid} \forall f \in \mathcal{C}. \nexists g\in \operatorname{Mono}(F,G).\,f \mathop{\star} g \mathop{=} h \}}
        % \begin{cases}
        %     \card{\{h \mathop{\in} \operatorname{Mono}(X,G) \mathop{\mid} \nexists g.\,f \mathop{\star} g \mathop{=} h \}}, & \text{if } \mathcal{C} \mathop{=}  \set{f}, \\
        %     \card{\operatorname{Mono}(X,G)}, & \text{if } \mathcal{C} \mathop{=} \emptyset.
        % \end{cases}
    \]
\end{definition}
The definitions of weight function and graph weights are identical to those in Definition~\ref{subgraph_counting:def:weight_function} and Definition~\ref{subgraph_counting:def:graph_weight} in Chapter~\ref{chap:subgraph_counting}, respectively, except that we use the new definition of measurement.
% \subsection{Non-increasing rules}
% \label{sec:subgraph_counting_antipattern:non-increasing}
% % We use \( S \uplus S' \) to denote the disjoint union of sets \( S \) and \( S' \). Let \( \rho \mathop{=} (L \overset{l}{\leftarrowtail} K \overset{r}{\rightarrowtail} R) \) be a rule.
% \newline
% \noindent
%     \begin{minipage}{0.65\textwidth}
%       The DPO diagram for \( G \mathop{\Rightarrow}_\rho H \) with injective match \( m \) can be seen, up to graph renaming, as a diagram in the category of sets where all morphisms are inclusions. The diagram is shown on the right.
%     \end{minipage}
%     \hfill
%     \begin{minipage}{0.34\textwidth}
%           \hfill
%           \resizebox{0.9\textwidth}{!}{
%             \begin{tikzpicture}
%               \coordinate (k) at (0, 0);
%               \draw[fill=white] ($(k)+(0,0)$) rectangle ($(k)+(0.5,0.5)$);
%               \node () at ($(k)+(0.25,0.25)$) {\( K \)};
          
%               \coordinate (c) at (0, -2);
%               \draw[fill=blue!20]
%               ($(c)+(0,-0.5)$)
%               -- ($(c)+(0,0.5)$) 
%               -- ($(c)+(1,0.5)$) 
%               arc[start angle=0, end angle=-90, radius=1]
%               -- cycle;
%               \node () at ($(c)+(0.7,0.2)$) {\( C' \)};
%               \draw[fill=white] ($(c)+(0,0)$) rectangle ($(c)+(0.5,0.5)$);
%               \node () at ($(c)+(0.25,0.25)$) {\( K \)};
          
%               \coordinate (l) at (-3, 0);
%               \draw[fill=orange!20] ($(l)+(-0.5,0)$) rectangle ($(l)+(0.5,1)$);
%               \node () at ($(l)+(-0.2,0.2)$) {\( L' \)};
%               \draw[fill=white] ($(l)+(0,0)$) rectangle ($(l)+(0.5,0.5)$);
%               \node () at ($(l)+(0.25,0.25)$) {\( K \)};
          
%               \coordinate (g) at (-3, -2);
%               \draw[fill=blue!20]
%               ($(g)+(0,-0.5)$)
%               -- ($(g)+(0,0.5)$)
%               -- ($(g)+(1,0.5)$) 
%               arc[start angle=0, end angle=-90, radius=1]
%               -- cycle;
%               \draw[fill=orange!20] ($(g)+(-0.5,0)$) rectangle ($(g)+(0.5,1)$);
%               \node () at ($(g)+(0.7,0.2)$) {\( C' \)};
%               \node () at ($(g)+(-0.2,0.2)$) {\( L' \)};
%               \draw[fill=white] ($(g)+(0,0)$) rectangle ($(g)+(0.5,0.5)$);
%               \node () at ($(g)+(0.25,0.25)$) {\( K \)};
          
%               \coordinate (r) at (3,0);
%               \draw[fill=red!20] ($(r)+(-0.5,0)$)
%                 -- ($(r)+(-0.5,0.5)$)
%                 -- ($(r)+(0,1)$)
%                 --  ($(r)+(0.5,1)$)
%                 -- ($(r)+(0.5,0)$)
%                 -- cycle;
%               \node () at ($(r)+(-0.2,0.2)$) {\( R' \)};
%               \draw[fill=white] ($(r)+(0,0)$) rectangle ($(r)+(0.5,0.5)$);
%               \node () at ($(r)+(0.25,0.25)$) {\( K \)};
          
%               \coordinate (h) at (3, -2);
%               \draw[fill=blue!20]
%               ($(h)+(0,-0.5)$)
%               -- ($(h)+(0,0.5)$)
%               -- ($(h)+(1,0.5)$) 
%               arc[start angle=0, end angle=-90, radius=1]
%               -- cycle;
%               \draw[fill=red!20] ($(h)+(-0.5,0)$)
%               -- ($(h)+(-0.5,0.5)$)
%               -- ($(h)+(0,1)$)
%               --  ($(h)+(0.5,1)$)
%               -- ($(h)+(0.5,0)$)
%               -- cycle;
%              \node () at ($(h)+(0.7,0.2)$) {\( C' \)};
%              \draw[fill=white] ($(h)+(0,0)$) rectangle ($(h)+(0.5,0.5)$);
%              \node () at ($(h)+(0.25,0.25)$) {\( K \)};
%              \node () at ($(h)+(-0.2,0.2)$) {\( R' \)};
          
%               \node (kl) at ($(k)!0.5!(l)+(0.1,0.2)$) {\( \overset{l}{\leftarrowtail} \)}; 
%               \node (kr) at ($(k)!0.5!(r)+(0.1,0.2)$) {\( \overset{r}{\rightarrowtail} \)};  
%               \node (cg) at ($(c)!0.5!(g)+(0.1,0.2)$) {\( \leftarrowtail \)}; 
%               \node (ch) at ($(c)!0.5!(h)+(0.1,0.2)$) {\( \rightarrowtail \)}; 
%               \node (kc) at ($(k)!0.5!(c)+(0.2,0.4)$) {\( \downarrowtail \)}; 
%               \node () at ($(l)!0.5!(g)+(-0.2,0.4)$) {\( m \)}; 
%               \node (lg) at ($(l)!0.5!(g)+(0.1,0.4)$) {\( \downarrowtail \)}; 
%               \node (rh) at ($(r)!0.5!(h)+(0.1,0.4)$) {\( \downarrowtail \)}; 
%             \end{tikzpicture}
%         % \end{center}
%         }
%     \end{minipage}
% \( L', R', C', K \) are mutually disjoint sets such that \( L' \uplus K \mathop{=} L \), \( R' \uplus K \mathop{=} R \), and \( C' \uplus K \mathop{=} C \). 
% Additionally, \( G \mathop{=} L \mathop{\cup} C \) and \( H \mathop{=} R \mathop{\cup} C' \). Thus, any subgraph of \( H \) (resp. \( G \)) can be seen as a union of subgraphs of \( C \) and \( R \) (resp. \( L \)).
% We provide the following definition to improve readability.
% % \begin{definition}
% %         \label{def:corresponding_occurrence}
% %     Let $X$ be a graph and \( G \) a graph. An \textbf{\( X \)-occurrence} in \( G \) is a subgraph of $G$ isomorphic to \( X \). 
% %     Since there is a bijection between \( X \)-occurrences in \( G \) and monomorphisms \( X \rightarrowtail G \), we use the terms \enquote{\( X \)-occurrence} and \enquote{monomorphism from \( X \) to \( G \)} interchangeably.  
% %     An \textbf{implicitly created \( X \)-occurrence} is an \( X \)-occurrence in \( H \) that is included neither in \( R \) nor in \( C \). 
% %     For such an occurrence \( x \), let \( R' \mathop{=} x \mathop{\cap} R \), \( C' \mathop{=} x \mathop{\cap} C \), and let \( h_{R'L}: R' \rightarrowtail L \) be a monomorphism preserving interface elements. The \textbf{corresponding \( X \)-occurrence} relative to \( h_{R'L} \) is \( h_{R'L}(R') \mathop{\cup} C' \).  
% %     An \textbf{implicitly destroyed \( X \)-occurrence} is an \( X \)-occurrence in \( G \) that is not in \( L \) or \( C \).
% % \end{definition}
% \begin{definition}
%         \label{def:subgraph_counting_antipattern:corresponding_occurrence}
%     Let $X$ be a graph and \( G \) a graph. An \textbf{\( X \)-occurrence} in \( G \) is an isomorphism from \( X \) to \( G \).
%     An \textbf{implicitly created \( X \)-occurrence} is an \( X \)-occurrence whose image in \( H \) is included neither in \( R \) nor in \( C \). 
%     For such an occurrence \( x \), let \( R' \mathop{=} x \mathop{\cap} R \), \( C' \mathop{=} x \mathop{\cap} C \), and let \( h_{R'L}: R' \rightarrowtail L \) be a monomorphism preserving interface elements. The \textbf{corresponding \( X \)-occurrence} relative to \( h_{R'L} \) is \( h_{R'L}(R') \mathop{\cup} C' \).  
%     An \textbf{implicitly destroyed \( X \)-occurrence} is an \( X \)-occurrence whose image in \( G \) is neither in \( L \) nor in \( C \).
% \end{definition} 
% Let $X$ be a graph. We define the set \(D(R,X)\) of all subgraphs of \( R \) which can form an implicitly created \( X \)-occurrence with a subgraph of the context graph.
% \begin{definition}[Distinguished subgraphs of the right-hand side graph]
%     \label{def:subgraph_counting_antipattern:rx}
%     % Let \(X\) be a graph and 
%     % \(
%     %     L \overset{l}{\leftarrowtail} K \overset{r}{\rightarrowtail} R
%     % \) an injective DPO rewriting rule.
%     % %  such that $l$ is edge-injective and $r$ is node-injective.
%     % The set \(D(R,X)\) consists of all subgraphs \( R' \mathop{\subseteq} R \) such that there exists a diagram (shown on the right) where
%     % \(h_{R'X} \mathop{\colon} R' \rightarrowtail X \) and \(h_{C'X} \mathop{\colon} C' \rightarrowtail X \) are not isomorphisms, and \(h_{K'K} \mathop{\colon} K' \rightarrowtail K \) and \(h_{R'R} \mathop{\colon} R' \rightarrowtail R \) are inclusion functions. The union of all subgraphs in \( D(R,X) \) is denoted \( R_X \).
%     % Let \(X\) be a graph and 
%     % \(
%     %     L \overset{l}{\leftarrowtail} K \overset{r}{\rightarrowtail} R
%     % \) an injective DPO rewriting rule.
%     % %  such that $l$ is edge-injective and $r$ is node-injective.
%     % The set \(D(R,X)\) consists of all subgraphs \( R' \mathop{\subseteq} R \) such that 
%     % we can construct a diagram as shown on the right where \(h_{K'K} \mathop{\colon} K' \rightarrowtail K \) and \(h_{R'R} \mathop{\colon} R' \rightarrowtail R \) are inclusion functions, and \(h_{R'X} \mathop{\colon} R' \rightarrowtail X \) and \(h_{C'X} \mathop{\colon} C' \rightarrowtail X \) are not isomorphisms. The union of all subgraphs in \( D(R,X) \) is denoted \( R_X \).
%     Let \(\mathcal{X}\) be a ruler-graph with underlying graph \(X\) and
%     \(
%         L \overset{l}{\leftarrowtail} K \overset{r}{\rightarrowtail} R
%     \) an injective DPO rewriting rule.
%     %  such that $l$ is edge-injective and $r$ is node-injective.
%     The set \textbf{\(D(R,X)\)} consists of all subgraphs \( R' \mathop{\subseteq} R \)
%     satisfying the following conditions:
%     \newline
%     \noindent
%     \begin{minipage}{0.7\textwidth}\setlength{\parindent}{1em}
%     \begin{itemize}
%         \item \label{def:drx:rp_not_sub} $R'$ is not a subgraph of $r(K)$,
%         \item \label{def:drx:rp_not_iso} $R'$ is not isomorphic to $X$,
%         \item  \label{def:drx:rp_diagram} a diagram as shown on the right where \(h_{K'K} \mathop{\colon} K' \rightarrowtail K \) and \(h_{R'R} \mathop{\colon} R' \rightarrowtail R \) are inclusion functions can be constructed.
%     \end{itemize}
%     \end{minipage}
%     \todo{condition~\ref{def:drx:rp_diagram} est trop restrictive que la terminaison de [ad hoc protocol, bruggink 2014]}
%     \begin{minipage}{0.299\textwidth}
%         \hfill
%         \resizebox{0.9\textwidth}{!}{
%       \begin{tikzpicture}[node distance=11mm]
%             \node (k) {K};
%             \node (r) [above=of k] {R};
%             \node (r') [right=of r] {$R'$};
%             \node (k') [right=of k] {K'};
%             \node (x)  [right=of r'] {X};
%             \node (c') [below=of x] {C'};
%             \node () [at=($(r)!0.5!(k')$)] {$PB$};
%             \node () [at=($(x)!0.5!(k')$)] {$PO$};
%             \draw[>->] (k) -- (r) node[pos= .5, left] {$r$}; 
%             \draw[>->] (r') -- (r) node[pos= .5, above] {$h_{R'R}$};
%             \draw[>->] (k') -- (k) node[pos= .5, below] {$h_{K'K}$};
%             \draw[>->] (k') -- (r');
%             \draw[>->] (k') -- (c');
%             \draw[>->] (c') -- (x);
%             \draw[>->] (r') -- (x);
%         \end{tikzpicture}
%         }
%     \end{minipage}
%     % The union of all subgraphs in \( D(R,X) \) is denoted \( R_X \).
% \end{definition}  

% The first condition of the definition ensures that any $R' \mathop{\in} D(R,X)$ can be glued along some common interface elements with a subgraph $C'$ of a context graph \( C \) to form an \( X \)-occurrence. The second and third conditions ensure that such $X$-occurrences are included neither in \( C \) nor in \( R \), thereby guaranteeing that they are implicitly created as defined in Definition~\ref{def:corresponding_occurrence}.

% \begin{notation}
% We adopt the visualization conventions from~\cite[Example 9]{overbeek2023apbpotutorial}: labeled graphs are boxed, with their names in the top-left corner; nodes and edges are assigned subsets of \(\mathbb{N}\) as identifiers; identifiers are chosen in such a way that: a node or an edge \( y \) in the codomain graph inherits the union of identifiers from all preimages in the domain, and this assignment uniquely determines the homomorphism. To further improve readability, sets are represented as lists (e.g., \(\{3,4\}\) becomes \(3\ 4\)), and identifiers are omitted when unambiguous. An example is shown below. 
% \begin{center}
%     \resizebox{0.45\textwidth}{!}{
%     \begin{tikzpicture}
%         \graphbox{\( G \)}{00mm}{-20mm}{45mm}{20mm}{2mm}{-5mm}{
%             \coordinate (o) at (-5mm,-8mm); 
%             \node[draw,circle] (l1) at ($(o)+(-10mm,0mm)$) {1};
%             \node[draw,circle] (l2) at ($(l1)+(3,0)$) {2};
%             \node[draw,circle] (l3) at ($(l1)+(1,0)$) {3};
%             \node[draw,circle] (l4) at ($(l1)+(2,0)$) {4};
%             \draw[->] (l1) -- (l3) node[midway,above] {$a$};
%             \draw[->] (l3) -- (l4) node[midway,above] {$b$};
%             \draw[->] (l4) -- (l2) node[midway,above] {$a$};
%         }  
%         \graphbox{\( H \)}{50mm}{-20mm}{34mm}{20mm}{2mm}{-5mm}{
%             \coordinate (o) at (0mm,-8mm); 
%             \node[draw,circle] (l1) at ($(o)+(-10mm,0mm)$) {1};
%             \node[draw,circle] (l2) at ($(l1)+(2,0)$) {2};
%             \node[draw,circle] (l3) at ($(l1)+(1,0)$) {3\ 4};
%             \draw[->] (l1) -- (l3) node[midway,above] {$a$};
%             \draw[->] (l3) edge[loop above] (l3) node[midway,above] {$b$};
%             \draw[->] (l3) -- (l2) node[midway,above] {$a$};
%         }      
%         % \node () at (53mm,-30mm) {$\rightarrow$};
%     \end{tikzpicture}
%     }
% \end{center}  
% \end{notation}


% \begin{example}
%     \label{ex:grsaa_rx}
%     Consider the injective DPO rule from \cite[Example 6]{bruggink2014termination}:  
%     \begin{center} 
%         \resizebox{0.7\textwidth}{!}{
%         \begin{tikzpicture}
%             \graphbox{$L$}{0mm}{0mm}{34mm}{15mm}{2mm}{-5mm}{
%                 \coordinate (o) at (0mm,-3mm); 
%                 \node[draw,circle] (l1) at ($(o)+(-10mm,0mm)$) {1};
%                 \node[draw,circle] (l2) at ($(l1)+(2,0)$) {2};
%                 \node[draw,circle] (l3) at ($(l1)+(1,0)$) {3};
%                 \draw[->] (l1) -- (l3) node[midway,above] {$a$};
%                 \draw[->] (l3) -- (l2) node[midway,above] {$a$};
%             }     
%             \graphbox{$K$}{40mm}{0mm}{24mm}{15mm}{2mm}{-5mm}{
%                 \coordinate (o) at (5mm,-3mm); 
%                 \node[draw,circle] (l1) at ($(o)+(-10mm,0mm)$) {1};
%                 \node[draw,circle] (l2) at ($(l1)+(1,0)$) {2};
%                 % \node[draw,circle] (l3) at ($(l1)+(1,0)$) {$\ $};
%                 % \draw[->] (l1) -- (l3) node[midway,above] {$a$};
%                 % \draw[->] (l3) -- (l2) node[midway,above] {$a$};
%             }    
%             \graphbox{$R$}{70mm}{0mm}{45mm}{15mm}{2mm}{-5mm}{
%                 \coordinate (o) at (-5mm,-3mm); 
%                 \node[draw,circle] (l1) at ($(o)+(-10mm,0mm)$) {1};
%                 \node[draw,circle] (l2) at ($(l1)+(3,0)$) {2};
%                 \node[draw,circle] (l3) at ($(l1)+(1,0)$) {4};
%                 \node[draw,circle] (l4) at ($(l1)+(2,0)$) {5};
%                 \draw[->] (l1) -- (l3) node[midway,above] {$a$};
%                 \draw[->] (l3) -- (l4) node[midway,above] {$b$};
%                 \draw[->] (l4) -- (l2) node[midway,above] {$a$};
%             }    

%             \node () at (37mm,-8mm) {$\leftarrowtail$};
%             \node () at (67mm,-8mm) {$\rightarrowtail$};

%             % \draw[>->] (51mm,2mm) -- (52mm,3mm);
%         \end{tikzpicture}
%         }
%     \end{center}
%     Let $X$ be \tikz[baseline=-0.5ex]{
%         \node (x) at (0,0) {$\bullet$};
%         \node (y) at (1,0) {$\bullet$ };
%         \node (z) at (2,0) { $\bullet$};
%         \draw[->] (x) -- (y) node[midway, above] {$a$};
%         \draw[->] (y) -- (z) node[midway, above] {$a$};
%     } and $R' \mathop{\in} D(R,X)$. 
%     $R'$ is neither the empty graph nor isomorphic to $X$, by Conditions 2 and 3 of Definition~\ref{def:rx}. 
%     Therefore, $R'$ must contain an interface node (either node 1 or node 2), because if it contains no interface node, Condition 1 of Definition~\ref{def:rx} would be violated due to $X$'s connectivity.
%     Additionally, $R'$ must contain node 4 or node 5 to avoid being a subgraph of $r(K)$ (Condition 2). 
%     Therefore, $R'$ must contain either $1 \overset{a}{\to} 4$ or $5 \overset{a}{\to} 2$, because if node 4 or node 5 is isolated, Condition 1 would be violated due to $X$'s connectivity.
%     Furthermore, $R'$ must not contain $4 \overset{b}{\to} 5$ by Condition 1.
%     The following graph is not in $D(R,X)$ because it violates Condition 1.
%    \begin{center}
%         \resizebox{0.24\textwidth}{!}{
%             \begin{tikzpicture}
%                 \graphbox{}{70mm}{0mm}{45mm}{15mm}{2mm}{-5mm}{
%                     \coordinate (o) at (-5mm,-3mm); 
%                     \node[draw,circle] (l1) at ($(o)+(-10mm,0mm)$) {1};
%                     \node[draw,circle] (l3) at ($(l1)+(1,0)$) {4};
%                     \node[draw,circle] (l2) at ($(l1)+(3,0)$) {2};
%                     \node[draw,circle] (l4) at ($(l1)+(2,0)$) {5};
%                     \draw[->] (l4) -- (l2) node[midway,above] {$a$};
%                     \draw[->] (l1) -- (l3) node[midway,above] {$a$};
%                 }    
%             \end{tikzpicture}
%         } 
%     \end{center}
%     Thus, $D(R,X)$ consists of the four graphs $R'_1$, $R'_2$, $R'_3$ and $R'_4$ shown below.
%     For each graph, the construction of the pullback square in Condition 1 is demonstrated in~\autoref{example:grs_aa:has_more_left} and the construction of the pushout square in Condition 1 is straightforward. 
%     \begin{center}
%         \resizebox{0.24\textwidth}{!}{
%             \begin{tikzpicture}
%                 \graphbox{$R'_1$}{70mm}{0mm}{45mm}{15mm}{2mm}{-5mm}{
%                     \coordinate (o) at (-5mm,-3mm); 
%                     \node[circle] (l1) at ($(o)+(-10mm,0mm)$) {};
%                     \node[draw,circle] (l2) at ($(l1)+(3,0)$) {2};
%                     \node[draw,circle] (l4) at ($(l1)+(2,0)$) {5};
%                     \draw[->] (l4) -- (l2) node[midway,above] {$a$};
%                 }    
%             \end{tikzpicture}
%         }
%         \resizebox{0.24\textwidth}{!}{
%             \begin{tikzpicture}
%                 \graphbox{$R'_2$}{70mm}{0mm}{45mm}{15mm}{2mm}{-5mm}{
%                     \coordinate (o) at (-5mm,-3mm); 
%                     \node[draw,circle] (l1) at ($(o)+(-10mm,0mm)$) {1};
%                     \node[draw,circle] (l2) at ($(l1)+(3,0)$) {2};
%                     \node[draw,circle] (l4) at ($(l1)+(2,0)$) {5};
%                     \draw[->] (l4) -- (l2) node[midway,above] {$a$};
%                 }    
%             \end{tikzpicture}
%         } 
%         \resizebox{0.24\textwidth}{!}{
%             \begin{tikzpicture}
%                 \graphbox{$R'_3$}{70mm}{0mm}{45mm}{15mm}{2mm}{-5mm}{
%                     \coordinate (o) at (-5mm,-3mm); 
%                     \node[draw,circle] (l1) at ($(o)+(-10mm,0mm)$) {1};
%                     \node[draw,circle] (l3) at ($(l1)+(1,0)$) {4};
%                     \draw[->] (l1) -- (l3) node[midway,above] {$a$};
%                 }   
%         \end{tikzpicture}
%         } 
%         \resizebox{0.24\textwidth}{!}{
%             \begin{tikzpicture}
%                 \graphbox{$R'_4$}{70mm}{0mm}{45mm}{15mm}{2mm}{-5mm}{
%                     \coordinate (o) at (-5mm,-3mm); 
%                     \node[draw,circle] (l1) at ($(o)+(-10mm,0mm)$) {1};
%                     \node[draw,circle] (l2) at ($(l1)+(3,0)$) {2};
%                     \node[draw,circle] (l3) at ($(l1)+(1,0)$) {4};
%                     \draw[->] (l1) -- (l3) node[midway,above] {$a$};
%                 }    
%             \end{tikzpicture}
%         }
%     \end{center}
% \end{example} 
 
% \begin{example}
%     \label{ex:rx_counter_ex}
%     Let $X$ be the graph 
%     \tikz[baseline=-0.5ex]{ 
%             \node (x) at (0,0) {$\bullet$}; 
%             \node (y) at (1,0) {$\bullet$};
%             \node (z) at (2,0) {$\bullet$};
%             \draw[->] (x) -- (y)   {};
%             \draw[->] (y) -- (z)   {};
%     }. Consider the following rule.
%     \begin{center}
%         \resizebox{0.6\textwidth}{!}{
%             \begin{tikzpicture}
%                 \graphbox{$L$}{0mm}{0mm}{34mm}{15mm}{2mm}{-5mm}{
%                     \coordinate (o) at (0mm,-3mm); 
%                     \node[draw,circle] (l1) at ($(o)+(-10mm,0mm)$) {1};
%                     % \node[draw,circle] (l2) at ($(l1)+(2,0)$) {2};
%                     \node[draw,circle] (l3) at ($(l1)+(1,0)$) {2};
%                     \draw[->] (l3) -- (l1) node[midway,above] {};
%                     % \draw[->] (l3) -- (l2) node[midway,above] {$a$};
%                 }     
%                 \graphbox{$K$}{40mm}{0mm}{24mm}{15mm}{2mm}{-5mm}{
%                     \coordinate (o) at (5mm,-3mm); 
%                     \node[draw,circle] (l1) at ($(o)+(-10mm,0mm)$) {1};
%                     \node[draw,circle] (l2) at ($(l1)+(1,0)$) {2};
%                     % \node[draw,circle] (l3) at ($(l1)+(1,0)$) {$\ $};
%                     \draw[->] (l2) -- (l1) node[midway,above] {};
%                     % \draw[->] (l3) -- (l2) node[midway,above] {$a$};
%                 }    
%                 \graphbox{$R$}{70mm}{0mm}{40mm}{15mm}{2mm}{-5mm}{
%                     \coordinate (o) at (-5mm,-3mm); 
%                     \node[draw,circle] (l1) at ($(o)+(-10mm,0mm)$) {1};
%                     % \node[draw,circle] (l2) at ($(l1)+(3,0)$) {2};
%                     \node[draw,circle] (l3) at ($(l1)+(1,0)$) {2};
%                     \node[draw,circle] (l4) at ($(l1)+(2,0)$) {3};
%                     \draw[->] (l3) -- (l1) node[midway,above] {};
%                     \draw[->] (l3) -- (l4) node[midway,above] {};
%                     % \draw[->] (l4) -- (l2) node[midway,above] {$a$};
%                 }    
%                 \node () at (37mm,-8mm) {$\leftarrowtail$};
%                 \node () at (67mm,-8mm) {$\rightarrowtail$};
%             \end{tikzpicture}
%         }
%     \end{center}
%     Let $R' \mathop{\in} D(R,X)$. $R'$ must contain node 3, otherwise it would be a subgraph of $r(K)$, violating Condition 2 of Definition~\ref{def:rx}. Node 3 must be connected to an interface node, otherwise Condition 1 of Definition~\ref{def:rx} would be violated. Therefore, $R'$ must include \raisebox{2pt}{\scalebox{0.6}{\tikz[baseline=-0.5ex]{
%         \node [draw,circle] (x) at (0,0) {2};
%         \node[draw,circle] (y) at (1,0) {3};
%         \draw[->] (x) -- (y) {};
%     }}}.
%     The graph $R$ is not in $D(R,X)$ because it violates Condition 1. Thus, the set \( D(R,X) \) consists of $R'_1$:
%     \raisebox{2pt}{\scalebox{0.6}{\tikz[baseline=-0.5ex]{
%         \node [draw,circle] (x) at (0,0) {2};
%         \node[draw,circle] (y) at (1,0) {3};
%         \draw[->] (x) -- (y) {};
%     }}} and $R'_2$:
%     \raisebox{2pt}{\scalebox{0.6}{\tikz[baseline=-0.5ex]{
%         \node [draw,circle] (node1) at (-1,0) {1};
%         \node [draw,circle] (x) at (0,0) {2};
%         \node[draw,circle] (y) at (1,0) {3};
%         \draw[->] (x) -- (y) {};
%     }}}.
% \end{example}
% % \begin{remark}
% %     Consider a rewriting step defined by the diagram in Definition~\ref{def:rewriting_step}. Let $X$ be a ruler-graph. The set $\operatorname{Mono}(X,H,\lnot m', \lnot r')$ (resp. $\operatorname{Mono}(X,H,\lnot m', \lnot r')$) is the set of all monomorphisms whose image is an occurrence of $X$ in $H$ (resp. $G$) formed by the image of a graph $R' \mathop{\in} D(R,X)$ and the image of a graph $C' \mathop{\subseteq} C$ around the image of a graph $K' \mathop{\subseteq} K$.
% % \end{remark}
% % \begin{definition}
% %    Let $X$ be a ruler-graph. Consider a rewriting step defined by the diagram in Definition~\ref{def:rewriting_step}. An occurrence $h_{XH}:X \mathop{\rightarrow} H$ of $X$ in $H$ is a \textbf{created occurrence} of $X$ in $H$ if the image of $h_{XH}$ is neither fully include in $m'(R)$ nor fully include in $r'(C)$. For a created occurrence $h_{XH}$, let $R' \mathop{\subseteq} R$ such that $m'(R') \mathop{=} h_{XH}(X) \mathop{\cap} m'(R)$, and let $h_{R'L}: R' \rightarrowtail L$ be a monomorphism, we say that an $h_{XG}$ 
% % \end{definition}

% % \newpage
% % \begin{definition}[Non-increasing rule]
% %     \label{def:creates_more_x_on_the_left}
% %     \newline
% %     \noindent
% %     \begin{minipage}{0.7\textwidth}\setlength{\parindent}{1em}
% %     Let \(\rho \mathop{=} (L \overset{l}{\leftarrowtail} K \overset{r}{\rightarrowtail} R)\) be a rule, and let $X$ be a graph.
% %     Let \( \Psi \) be a function associating $R' \mathop{\in} D(R,X)$ to a homomorphism in $\operatorname{Mono}(R',L)$.
% %     We say that \textbf{the number of occurrences of $X$ is non-increasing from left to right in any $\rho$-rewriting step under the mapping \(\Psi\)} (or simply $\rho$ is $X$-non-increasing) if the following four conditions hold:
% %     \end{minipage}
% %     \begin{minipage}{0.3\textwidth}
% %         \hfill
% %         \begin{center}
% %             \resizebox{0.9\textwidth}{!}{
% %                 \begin{tikzpicture}
% %                     \node (k) {K}; 
% %                     \node (k') [above=of k] {$K'$};
% %                     \node (l) [left=of k] {$L$};
% %                     \node (rb) [above=of l] {$R'$};
% %                     \node (r) [right=of k] {$R$};
% %                     \node (rb') [above=of r] {$R'$};
% %                     \draw[<-<]  (l) -- (k) node [midway,above] {$l$};
% %                     \draw[<-<]  (r) -- (k) node [midway,above] {$r$};
% %                     \draw[>->]  (k') -- (rb') node [midway,above] {};
% %                     \draw[>->]  (rb') -- (r) node [midway,above] {};
% %                     \draw[>->]  (k') -- (rb) node [midway, right] {};
% %                     \draw[>->]  (k') -- (k);
% %                     \draw[>->]  (rb) -- (l) node [midway,left] {$\Psi(R')$};
% %                     \node () [at=($(k)!0.5!(rb)$)] {$PB$};
% %                     \node () [at=($(k)!0.5!(rb')$)] {$PB$};
% %                 \end{tikzpicture}
% %             }
% %         \end{center}
% %     \end{minipage}
% %      \begin{enumerate}
% %         \item[(1)] \label{def:non_increasing:cond_1} For all $R' \mathop{\in} D(R,X)$ the diagram displayed above on the right holds, where $h_{R'R}$ is the inclusion function;
% %         \item[(2)] \label{def:non_increasing:non_clapse} For all $R' \mathop{\in} D(R,X)$ and for all nodes and arrows $x$ in $R'$, if $ x \notin \operatorname{Im}(r)$ then $\Psi(R')(x) \notin \operatorname{Im}(l)$;
% %         % \item there exists a morphism $h_{R_XL}$ such that
% %         %     \( h_{R'R_X} \mathop{\star} h_{R_XL} \mathop{=} h_{R'L} \) for all $R' \mathop{\in} D(R,X)$;
% %         \item[(3)] \label{def:non_increasing:edge_injective} For all $R',R'' \mathop{\in} D(R,X)$, for all edges $x \mathop{\in} R'$ and $y \mathop{\in} R''$, if $x \mathop{\neq} y$, then $\Psi(R')(x) \mathop{\neq} \Psi(R'')(y)$;
% %         %  \item $h$  is edge-injective 
% %         % \item $\bigcup_{R' \mathop{\in} D(R,X)} h_{R'L}$ is an edge-injective morphism from $R_X$ to $L$;
% %         \item[(4)] \label{def:non_increasing:isolated_nodes} If $X$ has isolated nodes, then for all $R',R'' \mathop{\in} D(R,X)$ and nodes $x \mathop{\in} R', y \mathop{\in} R''$, if $x \mathop{\neq} y$, then $\Psi(R')(x) \mathop{\neq} \Psi(R'')(y)$.
% %         % \item $h$ is node-injective if $X$ has isolated nodes;
% %      \end{enumerate}
% %      \todo{condition~\ref{def:non_increasing:cond_1} est trop restrictive que la terminaison de [ad hoc protocol, bruggink2014termination]}
% % \end{definition}
% We define the notion of \( X \)-non-increasing rules. Intuitively, if a rule $\rho$ is $X$-non-increasing, then for any rewriting step using $\rho$, there exists an injective mapping from the set of \( X \)-occurrences implicitly created by the step to the set of \( X \)-occurrences implicitly destroyed by the step. This is proved in \cite{qiu2025termination_icgt}.
% \begin{definition}[Non-increasing rule]
%     \label{def:creates_more_x_on_the_left}
%     Let \(\rho \mathop{=} (L \overset{l}{\leftarrowtail} K \overset{r}{\rightarrowtail} R)\) be a rule and \(X\) a graph. 
%     Let \( \Psi \) be a function associating $R' \mathop{\in} D(R,X)$ to a homomorphism in $\operatorname{Mono}(R',L)$.
%     We say that \textbf{the number of occurrences of $X$ is non-increasing from left to right in any $\rho$-rewriting step under the mapping \(\Psi\)}
%     (or that $\rho$ is $X$-\textbf{non-increasing} under $\Psi$. When $\Psi$ is clear from context or irrelevant to the discussion, we may simply say that $\rho$ is $X$-\textbf{non-increasing}.)
%     % (or simply $\rho$ is $X$-non-increasing under \(\Psi\), or $\rho$ is $X$-non-increasing when \(\Psi\) is clear from the context or irrelevant to the discussion)
%     if the following four conditions hold:
%     \newline
%     \vspace{1mm}
%     \noindent
%     \begin{minipage}{0.64\textwidth} 
%     \begin{enumerate}
%         \item 
%         For all $R' \mathop{\in} D(R,X)$, in the diagram shown on the right, where $h_{R'R}$ is the inclusion function, both squares are pullbacks,
%         % For all $R' \mathop{\in} D(R,X)$, there exists a morphism \(h_{R'L}: R' \rightarrowtail L \) such that \trackedtext{both squares in the diagram shown on the right are pullbacks}, where all morphisms other than $h_{R'L}$, $l$ and $r$ are inclusion functions,
%         \item \label{def:non_increasing:non_clapse} For all $R' \mathop{\in} D(R,X)$, for all nodes and edges $x$ in $R'$, if $ x \notin \operatorname{Im}(r)$ then 
%         $\Psi(R')(x) \notin \operatorname{Im}(l)$,
%         % $h_{R'L}(r') \notin \operatorname{Im}(l)$,
%         \item \label{def:non_increasing_rule_img_edges_distinct} For all $R',R'' \mathop{\in} D(R,X)$, for all edges $x$ in $R'$ and $y$ in $R''$, if $x \mathop{\neq} y$, then $\Psi(R')(x) \mathop{\neq} \Psi(R'')(y)$
%         \item If $X$ has isolated nodes, then for all $R',R'' \mathop{\in} D(R,X)$ and nodes $x$ and $y$ in $R'$ and $R''$ respectively, if $x \mathop{\neq} y$, then $\Psi(R')(x) \mathop{\neq} \Psi(R'')(y)$.
%     \end{enumerate} 
%     \end{minipage}
%     \hfill
%     \begin{minipage}{0.4\textwidth}
%         % \hfill
%         \begin{center}
%             % \resizebox{0.9\textwidth}{!}{ 
%                 % \begin{tikzpicture}[rotate=90]
%                 %     \node (k) {K}; 
%                 %     \node (k') [above=of k] {$K'$};
%                 %     \node (l) [left=of k] {$L$};
%                 %     \node (rb) [above=of l] {$R'$};
%                 %     \node (r) [right=of k] {$R$};
%                 %     \node (rb') [above=of r] {$R'$};
%                 %     \draw[<-<]  (l) -- (k) node [midway,below] {$l$};
%                 %     \draw[<-<]  (r) -- (k) node [midway,below] {$r$};
%                 %     \draw[>->]  (k') -- (rb');
%                 %     \draw[>->]  (rb') -- (r);
%                 %     \draw[>->]  (k') -- (rb);
%                 %     \draw[>->]  (k') -- (k);
%                 %     \draw[>->]  (rb) -- (l) node [midway,above,sloped] {$h_{R'L}$};
%                 %     \node () [at=($(k)!0.5!(rb)$)] {$PB$};
%                 %     \node () [at=($(k)!0.5!(rb')$)] {$PB$};
%                 % \end{tikzpicture}
%                 \begin{tikzpicture} 
%                     \node (k) {K}; 
%                     \node (k') [left=of k] {$K'$};
%                     \node (l) [above=of k] {$L$};
%                     \node (rb) [left=of l] {$R'$};
%                     \node (r) [below=of k] {$R$};
%                     \node (rb') [left=of r] {$R'$};
%                     \draw[<-<]  (l) -- (k) node [midway,right] {$l$};
%                     \draw[<-<]  (r) -- (k) node [midway,right] {$r$};
%                     \draw[>->]  (k') -- (rb');
%                     \draw[>->]  (rb') -- (r);
%                     \draw[>->]  (k') -- (rb);
%                     \draw[>->]  (k') -- (k);
%                     \draw[>->]  (rb) -- (l) node [midway,above] {$\Psi(R')$};
%                     \node () [at=($(k)!0.5!(rb)$)] {$PB$};
%                     \node () [at=($(k)!0.5!(rb')$)] {$PB$}; 
%                 \end{tikzpicture}
%             % }
%         \end{center}
%     \end{minipage}
% \end{definition}
\begin{example}
    \label{example:grs_aa:has_more_left}
    Consider the rule presented
    in Example~\ref{ex:grsaa_rx}. Let $X$ be \tikz[baseline=-0.5ex]{
        \node (x) at (0,0) {$\bullet$};
        \node (y) at (1,0) {$\bullet$};
        \node (z) at (2,0) {$\bullet$};
        \draw[->] (x) -- (y) node[midway, above] {$a$};
        \draw[<-] (z) -- (y) node[midway, above] {$a$};
    }.
    By Example~\ref{ex:grsaa_rx}, the set \( D(R,X) \) consists of exactly four graphs: $R'_1$, $R'_2$, $R'_3$ and $R'_4$.
    The rule is $X$-non-increasing because the following diagrams can be constructed, and $h_{R'_1L}, h_{R'_2L}, h_{R'_3L}$ and $h_{R'_4L}$ satisfy all conditions in Definition~\ref{def:creates_more_x_on_the_left}. 
    %1
    \begin{center}
        \resizebox{0.7\textwidth}{!}{
            \begin{tikzpicture}
                \graphbox{$R'_1$}{0mm}{0mm}{34mm}{15mm}{2mm}{-5mm}{
                    \coordinate (o) at (0mm,-3mm); 
                    \node (l1) at ($(o)+(-10mm,0mm)$) {};
                    \node[draw,circle] (l2) at ($(l1)+(2,0)$) {2};
                    \node[draw,circle] (l3) at ($(l1)+(1,0)$) {5};
                    % \draw[->] (l1) -- (l3) node[midway,above] {$a$};
                    \draw[->] (l3) -- (l2) node[midway,above] {$a$};
                }     
                \graphbox{$K'$}{40mm}{0mm}{24mm}{15mm}{2mm}{-5mm}{
                    \coordinate (o) at (5mm,-3mm); 
                    \node (l1) at ($(o)+(-10mm,0mm)$) {};
                    \node[draw,circle] (l2) at ($(l1)+(1,0)$) {2};
                    % \node[draw,circle] (l3) at ($(l1)+(1,0)$) {$\ $};
                    % \draw[->] (l1) -- (l3) node[midway,above] {$a$};
                    % \draw[->] (l3) -- (l2) node[midway,above] {$a$};
                }    
                \graphbox{$R'_1$}{70mm}{0mm}{45mm}{15mm}{2mm}{-5mm}{
                    \coordinate (o) at (-5mm,-3mm); 
                    \node (l1) at ($(o)+(-10mm,0mm)$) {};
                    \node[draw,circle] (l2) at ($(l1)+(3,0)$) {2};
                    % \node[draw,circle] (l3) at ($(l1)+(1,0)$) {4};
                    \node[draw,circle] (l4) at ($(l1)+(2,0)$) {5};
                    % \draw[->] (l1) -- (l3) node[midway,above] {$a$};
                    % \draw[->] (l3) -- (l4) node[midway,above] {$b$};
                    \draw[->] (l4) -- (l2) node[midway,above] {$a$};
                }    
                \graphbox{$L$}{0mm}{-20mm}{34mm}{15mm}{2mm}{-5mm}{
                    \coordinate (o) at (0mm,-3mm); 
                    \node[draw,circle] (l1) at ($(o)+(-10mm,0mm)$) {1};
                    \node[draw,circle] (l2) at ($(l1)+(2,0)$) {2};
                    \node[draw,circle] (l3) at ($(l1)+(1,0)$) {5};
                    \draw[->] (l1) -- (l3) node[midway,above] {$a$};
                    \draw[->] (l3) -- (l2) node[midway,above] {$a$};
                }     
                \graphbox{$K$}{40mm}{-20mm}{24mm}{15mm}{2mm}{-5mm}{
                    \coordinate (o) at (5mm,-3mm); 
                    \node[draw,circle] (l1) at ($(o)+(-10mm,0mm)$) {1};
                    \node[draw,circle] (l2) at ($(l1)+(1,0)$) {2};
                    % \node[draw,circle] (l3) at ($(l1)+(1,0)$) {$\ $};
                    % \draw[->] (l1) -- (l3) node[midway,above] {$a$};
                    % \draw[->] (l3) -- (l2) node[midway,above] {$a$};
                }    
                \graphbox{$R$}{70mm}{-20mm}{45mm}{15mm}{2mm}{-5mm}{
                    \coordinate (o) at (-5mm,-3mm); 
                    \node[draw,circle] (l1) at ($(o)+(-10mm,0mm)$) {1};
                    \node[draw,circle] (l2) at ($(l1)+(3,0)$) {2};
                    \node[draw,circle] (l3) at ($(l1)+(1,0)$) {4};
                    \node[draw,circle] (l4) at ($(l1)+(2,0)$) {5};
                    \draw[->] (l1) -- (l3) node[midway,above] {$a$};
                    \draw[->] (l3) -- (l4) node[midway,above] {$b$};
                    \draw[->] (l4) -- (l2) node[midway,above] {$a$};
                }    
                \node () at (37mm,-8mm) {$\leftarrowtail$};
                \node () at (17mm,-17mm) {$\downarrowtail$};
                \node () at (52mm,-17mm) {$\downarrowtail$};
                \node () at (92mm,-17mm) {$\downarrowtail$};
                \node () at (37mm,-18mm) {\text{PB}};
                \node () at (68mm,-18mm) {\text{PB}};
                \node () at (67mm,-8mm) {$\rightarrowtail$};
                \node () at (37mm,-28mm) {$\leftarrowtail$};
                \node () at (67mm,-28mm) {$\rightarrowtail$};
            \end{tikzpicture}
        } 
    \end{center}   
    % 2
    \begin{center}
        \resizebox{0.7\textwidth}{!}{
            \begin{tikzpicture}
                \graphbox{$R'_2$}{0mm}{0mm}{34mm}{15mm}{2mm}{-5mm}{
                    \coordinate (o) at (0mm,-3mm); 
                    \node[draw,circle] (l1) at ($(o)+(-10mm,0mm)$) {1};
                    \node[draw,circle] (l2) at ($(l1)+(2,0)$) {2};
                    \node[draw,circle] (l3) at ($(l1)+(1,0)$) {5};
                    % \draw[->] (l1) -- (l3) node[midway,above] {$a$};
                    \draw[->] (l3) -- (l2) node[midway,above] {$a$};
                }     
                \graphbox{$K'$}{40mm}{0mm}{24mm}{15mm}{2mm}{-5mm}{
                    \coordinate (o) at (5mm,-3mm); 
                    \node[draw,circle] (l1) at ($(o)+(-10mm,0mm)$) {1};
                    \node[draw,circle] (l2) at ($(l1)+(1,0)$) {2};
                    % \node[draw,circle] (l3) at ($(l1)+(1,0)$) {$\ $};
                    % \draw[->] (l1) -- (l3) node[midway,above] {$a$};
                    % \draw[->] (l3) -- (l2) node[midway,above] {$a$};
                }    
                \graphbox{$R'_2$}{70mm}{0mm}{45mm}{15mm}{2mm}{-5mm}{
                    \coordinate (o) at (-5mm,-3mm); 
                    \node[draw,circle] (l1) at ($(o)+(-10mm,0mm)$) {1};
                    \node[draw,circle] (l2) at ($(l1)+(3,0)$) {2};
                    % \node[draw,circle] (l3) at ($(l1)+(1,0)$) {4};
                    \node[draw,circle] (l4) at ($(l1)+(2,0)$) {5};
                    % \draw[->] (l1) -- (l3) node[midway,above] {$a$};
                    % \draw[->] (l3) -- (l4) node[midway,above] {$b$};
                    \draw[->] (l4) -- (l2) node[midway,above] {$a$};
                }    
                \graphbox{$L$}{0mm}{-20mm}{34mm}{15mm}{2mm}{-5mm}{
                    \coordinate (o) at (0mm,-3mm); 
                    \node[draw,circle] (l1) at ($(o)+(-10mm,0mm)$) {1};
                    \node[draw,circle] (l2) at ($(l1)+(2,0)$) {2};
                    \node[draw,circle] (l3) at ($(l1)+(1,0)$) {5};
                    \draw[->] (l1) -- (l3) node[midway,above] {$a$};
                    \draw[->] (l3) -- (l2) node[midway,above] {$a$};
                }     
                \graphbox{$K$}{40mm}{-20mm}{24mm}{15mm}{2mm}{-5mm}{
                    \coordinate (o) at (5mm,-3mm); 
                    \node[draw,circle] (l1) at ($(o)+(-10mm,0mm)$) {1};
                    \node[draw,circle] (l2) at ($(l1)+(1,0)$) {2};
                    % \node[draw,circle] (l3) at ($(l1)+(1,0)$) {$\ $};
                    % \draw[->] (l1) -- (l3) node[midway,above] {$a$};
                    % \draw[->] (l3) -- (l2) node[midway,above] {$a$};
                }    
                \graphbox{$R$}{70mm}{-20mm}{45mm}{15mm}{2mm}{-5mm}{
                    \coordinate (o) at (-5mm,-3mm); 
                    \node[draw,circle] (l1) at ($(o)+(-10mm,0mm)$) {1};
                    \node[draw,circle] (l2) at ($(l1)+(3,0)$) {2};
                    \node[draw,circle] (l3) at ($(l1)+(1,0)$) {4};
                    \node[draw,circle] (l4) at ($(l1)+(2,0)$) {5};
                    \draw[->] (l1) -- (l3) node[midway,above] {$a$};
                    \draw[->] (l3) -- (l4) node[midway,above] {$b$};
                    \draw[->] (l4) -- (l2) node[midway,above] {$a$};
                }    
                \node () at (37mm,-8mm) {$\leftarrowtail$};
                \node () at (17mm,-17mm) {$\downarrowtail$};
                \node () at (52mm,-17mm) {$\downarrowtail$};
                \node () at (92mm,-17mm) {$\downarrowtail$};
                \node () at (37mm,-18mm) {\text{PB}};
                \node () at (68mm,-18mm) {\text{PB}};
                \node () at (67mm,-8mm) {$\rightarrowtail$};
                \node () at (37mm,-28mm) {$\leftarrowtail$};
                \node () at (67mm,-28mm) {$\rightarrowtail$};
            \end{tikzpicture}
        } 
    \end{center}
    %3
    \begin{center}
        \resizebox{0.7\textwidth}{!}{
            \begin{tikzpicture}
                \graphbox{$R'_3$}{0mm}{0mm}{34mm}{15mm}{2mm}{-5mm}{
                    \coordinate (o) at (0mm,-3mm); 
                    \node[draw,circle] (l1) at ($(o)+(-10mm,0mm)$) {1};
                    % \node[draw,circle] (l2) at ($(l1)+(2,0)$) {2};
                    \node[draw,circle] (l3) at ($(l1)+(1,0)$) {4};
                    \draw[->] (l1) -- (l3) node[midway,above] {$a$};
                    % \draw[->] (l3) -- (l2) node[midway,above] {$a$};
                }     
                \graphbox{$K'$}{40mm}{0mm}{24mm}{15mm}{2mm}{-5mm}{
                    \coordinate (o) at (5mm,-3mm); 
                    \node[draw,circle] (l1) at ($(o)+(-10mm,0mm)$) {1};
                    % \node[draw,circle] (l2) at ($(l1)+(1,0)$) {2};
                    % \node[draw,circle] (l3) at ($(l1)+(1,0)$) {$\ $};
                    % \draw[->] (l1) -- (l3) node[midway,above] {$a$};
                    % \draw[->] (l3) -- (l2) node[midway,above] {$a$};
                }     
                \graphbox{$R'_3$}{70mm}{0mm}{45mm}{15mm}{2mm}{-5mm}{
                    \coordinate (o) at (-5mm,-3mm); 
                    \node[draw,circle] (l1) at ($(o)+(-10mm,0mm)$) {1};
                    % \node[draw,circle] (l2) at ($(l1)+(3,0)$) {4};
                    \node[draw,circle] (l3) at ($(l1)+(1,0)$) {4};
                    % \node[draw,circle] (l4) at ($(l1)+(2,0)$) {5};
                    \draw[->] (l1) -- (l3) node[midway,above] {$a$};
                    % \draw[->] (l3) -- (l4) node[midway,above] {$b$};
                    % \draw[->] (l4) -- (l2) node[midway,above] {$a$};
                }    
                \graphbox{$L$}{0mm}{-20mm}{34mm}{15mm}{2mm}{-5mm}{
                    \coordinate (o) at (0mm,-3mm); 
                    \node[draw,circle] (l1) at ($(o)+(-10mm,0mm)$) {1};
                    \node[draw,circle] (l2) at ($(l1)+(2,0)$) {2};
                    \node[draw,circle] (l3) at ($(l1)+(1,0)$) {4};
                    \draw[->] (l1) -- (l3) node[midway,above] {$a$};
                    \draw[->] (l3) -- (l2) node[midway,above] {$a$};
                }     
                \graphbox{$K$}{40mm}{-20mm}{24mm}{15mm}{2mm}{-5mm}{
                    \coordinate (o) at (5mm,-3mm); 
                    \node[draw,circle] (l1) at ($(o)+(-10mm,0mm)$) {1};
                    \node[draw,circle] (l2) at ($(l1)+(1,0)$) {2};
                    % \node[draw,circle] (l3) at ($(l1)+(1,0)$) {$\ $};
                    % \draw[->] (l1) -- (l3) node[midway,above] {$a$};
                    % \draw[->] (l3) -- (l2) node[midway,above] {$a$};
                }    
                \graphbox{$R$}{70mm}{-20mm}{45mm}{15mm}{2mm}{-5mm}{
                    \coordinate (o) at (-5mm,-3mm); 
                    \node[draw,circle] (l1) at ($(o)+(-10mm,0mm)$) {1};
                    \node[draw,circle] (l2) at ($(l1)+(3,0)$) {2};
                    \node[draw,circle] (l3) at ($(l1)+(1,0)$) {4};
                    \node[draw,circle] (l4) at ($(l1)+(2,0)$) {5};
                    \draw[->] (l1) -- (l3) node[midway,above] {$a$};
                    \draw[->] (l3) -- (l4) node[midway,above] {$b$};
                    \draw[->] (l4) -- (l2) node[midway,above] {$a$};
                }    
                \node () at (37mm,-8mm) {$\leftarrowtail$};
                \node () at (17mm,-17mm) {$\downarrowtail$};
                \node () at (52mm,-17mm) {$\downarrowtail$};
                \node () at (92mm,-17mm) {$\downarrowtail$};
                \node () at (37mm,-18mm) {\text{PB}};
                \node () at (68mm,-18mm) {\text{PB}};
                \node () at (67mm,-8mm) {$\rightarrowtail$};
                \node () at (37mm,-28mm) {$\leftarrowtail$};
                \node () at (67mm,-28mm) {$\rightarrowtail$};
            \end{tikzpicture}
        } 
    \end{center}
    %4
    \begin{center}
        \resizebox{0.7\textwidth}{!}{
            \begin{tikzpicture}
                \graphbox{$R'_4$}{0mm}{0mm}{34mm}{15mm}{2mm}{-5mm}{
                    \coordinate (o) at (0mm,-3mm); 
                    \node[draw,circle] (l1) at ($(o)+(-10mm,0mm)$) {1};
                    \node[draw,circle] (l2) at ($(l1)+(2,0)$) {2};
                    \node[draw,circle] (l3) at ($(l1)+(1,0)$) {4};
                    \draw[->] (l1) -- (l3) node[midway,above] {$a$};
                    % \draw[->] (l3) -- (l2) node[midway,above] {$a$};
                }     
                \graphbox{$K'$}{40mm}{0mm}{24mm}{15mm}{2mm}{-5mm}{
                    \coordinate (o) at (5mm,-3mm); 
                    \node[draw,circle] (l1) at ($(o)+(-10mm,0mm)$) {1};
                    \node[draw,circle] (l2) at ($(l1)+(1,0)$) {2};
                    % \node[draw,circle] (l3) at ($(l1)+(1,0)$) {$\ $};
                    % \draw[->] (l1) -- (l3) node[midway,above] {$a$};
                    % \draw[->] (l3) -- (l2) node[midway,above] {$a$};
                }     
                \graphbox{$R'_4$}{70mm}{0mm}{45mm}{15mm}{2mm}{-5mm}{
                    \coordinate (o) at (-5mm,-3mm); 
                    \node[draw,circle] (l1) at ($(o)+(-10mm,0mm)$) {1};
                    \node[draw,circle] (l2) at ($(l1)+(3,0)$) {2};
                    \node[draw,circle] (l3) at ($(l1)+(1,0)$) {4};
                    % \node[draw,circle] (l4) at ($(l1)+(2,0)$) {5};
                    \draw[->] (l1) -- (l3) node[midway,above] {$a$};
                    % \draw[->] (l3) -- (l4) node[midway,above] {$b$};
                    % \draw[->] (l4) -- (l2) node[midway,above] {$a$};
                }    
                \graphbox{$L$}{0mm}{-20mm}{34mm}{15mm}{2mm}{-5mm}{
                    \coordinate (o) at (0mm,-3mm); 
                    \node[draw,circle] (l1) at ($(o)+(-10mm,0mm)$) {1};
                    \node[draw,circle] (l2) at ($(l1)+(2,0)$) {2};
                    \node[draw,circle] (l3) at ($(l1)+(1,0)$) {4};
                    \draw[->] (l1) -- (l3) node[midway,above] {$a$};
                    \draw[->] (l3) -- (l2) node[midway,above] {$a$};
                }     
                \graphbox{$K$}{40mm}{-20mm}{24mm}{15mm}{2mm}{-5mm}{
                    \coordinate (o) at (5mm,-3mm); 
                    \node[draw,circle] (l1) at ($(o)+(-10mm,0mm)$) {1};
                    \node[draw,circle] (l2) at ($(l1)+(1,0)$) {2};
                    % \node[draw,circle] (l3) at ($(l1)+(1,0)$) {$\ $};
                    % \draw[->] (l1) -- (l3) node[midway,above] {$a$};
                    % \draw[->] (l3) -- (l2) node[midway,above] {$a$};
                }    
                \graphbox{$R$}{70mm}{-20mm}{45mm}{15mm}{2mm}{-5mm}{
                    \coordinate (o) at (-5mm,-3mm); 
                    \node[draw,circle] (l1) at ($(o)+(-10mm,0mm)$) {1};
                    \node[draw,circle] (l2) at ($(l1)+(3,0)$) {2};
                    \node[draw,circle] (l3) at ($(l1)+(1,0)$) {4};
                    \node[draw,circle] (l4) at ($(l1)+(2,0)$) {5};
                    \draw[->] (l1) -- (l3) node[midway,above] {$a$};
                    \draw[->] (l3) -- (l4) node[midway,above] {$b$};
                    \draw[->] (l4) -- (l2) node[midway,above] {$a$};
                }    
                \node () at (37mm,-8mm) {$\leftarrowtail$};
                \node () at (17mm,-17mm) {$\downarrowtail$};
                \node () at (52mm,-17mm) {$\downarrowtail$};
                \node () at (92mm,-17mm) {$\downarrowtail$};
                \node () at (37mm,-18mm) {\text{PB}};
                \node () at (68mm,-18mm) {\text{PB}};
                \node () at (67mm,-8mm) {$\rightarrowtail$};
                \node () at (37mm,-28mm) {$\leftarrowtail$};
                \node () at (67mm,-28mm) {$\rightarrowtail$};
            \end{tikzpicture}
        } 
    \end{center}
    % with the following morphisms $h_{R'L}$ and $h_{R''L}$:
    % \begin{center}
    %     \resizebox{0.6\textwidth}{!}{
    %         \begin{tikzpicture}
    %             \graphbox{$R'$}{0}{0}{44mm}{15mm}{2mm}{-5mm}{
    %                 \coordinate (o) at (-5mm,-3mm); 
    %                 \node[draw,circle] (l1) at ($(o)+(-10mm,0mm)$) {1};
    %                 \node[draw,circle] (l3) at ($(l1)+(1,0)$) {4};
    %                 \draw[->] (l1) -- (l3) node[midway,above] {$a$};
    %             } 
    %             \graphbox{$L$}{60mm}{0mm}{34mm}{15mm}{2mm}{-5mm}{
    %                 \coordinate (o) at (0mm,-3mm); 
    %                 \node[draw,circle] (l1) at ($(o)+(-10mm,0mm)$) {1};
    %                 \node[draw,circle] (l2) at ($(l1)+(2,0)$) {};
    %                 \node[draw,circle] (l3) at ($(l1)+(1,0)$) {4};
    %                 \draw[->] (l1) -- (l3) node[midway,above] {$a$};
    %                 \draw[->] (l3) -- (l2) node[midway,above] {$a$};
    %             }    
    %         \node () at (52mm,-8mm) {$\rightarrowtail$};
    %         \end{tikzpicture}
    %      }
    %     \end{center}
    %     \begin{center}
    %         \resizebox{0.6\textwidth}{!}{
    %             \begin{tikzpicture}
    %                 \graphbox{$R''$}{0}{0}{44mm}{15mm}{2mm}{-5mm}{
    %                     \coordinate (o) at (-5mm,-3mm); 
    %                     \node[] (l1) at ($(o)+(-10mm,0mm)$) {};
    %                     \node[draw,circle] (l2) at ($(l1)+(3,0)$) {2};
    %                     \node[draw,circle] (l4) at ($(l1)+(2,0)$) {5};
    %                     \draw[->] (l4) -- (l2) node[midway,above] {$a$};
    %                 } 
    %                 \graphbox{$L$}{60mm}{0mm}{34mm}{15mm}{2mm}{-5mm}{
    %                     \coordinate (o) at (0mm,-3mm); 
    %                     \node[draw,circle] (l1) at ($(o)+(-10mm,0mm)$) {};
    %                     \node[draw,circle] (l2) at ($(l1)+(2,0)$) {2};
    %                     \node[draw,circle] (l3) at ($(l1)+(1,0)$) {5};
    %                     \draw[->] (l1) -- (l3) node[midway,above] {$a$};
    %                     \draw[->] (l3) -- (l2) node[midway,above] {$a$};
    %                 }    
    %             \node () at (52mm,-8mm) {$\rightarrowtail$};
    %             \end{tikzpicture}
    %         }
    %     \end{center}
    % To illustrate the property captured by Definition~\ref{def:creates_more_x_on_the_left},
    % consider the rewriting step defined by the DPO diagram below. Whenever an element $R' \mathop{\in} D(R,X)$ forms an \( X \)-occurrence in $H$ that is included neither in \( R \) nor $C$ with $C' \mathop{\subseteq} C$, the graph
    % \( \operatorname{\Psi(R')} \mathop{\subseteq} L\) forms an \( X \)-occurrence in $G$ with $C'$ that is included neither in \( L \) nor $C$.
    % Furthermore, distinct $X$-occurrences in $H$ that are included neither in \( R \) nor $C$ have distinct corresponding $X$-occurrences in $G$ that are included neither in \( L \) nor $C$.
    To illustrate the $X$-non-increasing property captured by Definition~\ref{def:creates_more_x_on_the_left},
    consider the rewriting step defined by the DPO diagram below. Whenever \( R'_i \) (for \( 1 \leq i \leq 4 \)) forms an \( X \)-occurrence $x$ in $H$ that is included neither in \( R \) nor $C$ with a subgraph $C' \mathop{\subseteq} C \mathop{\cap} x \mathop{\subseteq} C$, the graph
    \( \operatorname{\Psi(R')} \mathop{\subseteq} L\) forms an \( X \)-occurrence in $G$ with $C'$ that is included neither in \( L \) nor $C$.
    Furthermore, distinct $X$-occurrences in $H$ that are included neither in \( R \) nor $C$ have distinct corresponding $X$-occurrences in $G$ that are included neither in \( L \) nor $C$.
    \begin{center} 
        \resizebox{0.7\textwidth}{!}{
        \begin{tikzpicture}
            \graphbox{\( L \)}{0mm}{5mm}{34mm}{20mm}{2mm}{-5mm}{
                \coordinate (o) at (0mm,-8mm); 
                \node[draw,circle] (l1) at ($(o)+(-10mm,0mm)$) {1};
                \node[draw,circle] (l2) at ($(l1)+(2,0)$) {2};
                \node[draw,circle] (l3) at ($(l1)+(1,0)$) {3};
                \draw[] (l1) -- (l3) node[midway,above] {$a$};
                \draw[] (l3) -- (l2) node[midway,above] {$a$};
            } 

            \graphbox{\( K \)}{40mm}{5mm}{34mm}{20mm}{2mm}{-5mm}{
                \coordinate (o) at (0mm,-8mm); 
                \node[draw,circle] (l1) at ($(o)+(-10mm,0mm)$) {1};
                \node[draw,circle] (l2) at ($(l1)+(2,0)$) {2};
            }  

            \graphbox{\( R \)}{80mm}{5mm}{45mm}{20mm}{2mm}{-5mm}{
                \coordinate (o) at (-5mm,-8mm); 
                \node[draw,circle] (l1) at ($(o)+(-10mm,0mm)$) {1};
                \node[draw,circle] (l2) at ($(l1)+(3,0)$) {2};
                \node[draw,circle] (l3) at ($(l1)+(1,0)$) {4};
                \node[draw,circle] (l4) at ($(l1)+(2,0)$) {5};
                \draw[ ] (l1) -- (l3) node[midway,above] {$a$};
                \draw[ ] (l3) -- (l4) node[midway,above] {$b$};
                \draw[ ] (l4) -- (l2) node[midway,above] {$a$};
            }    

            \graphbox{\( G \)}{0mm}{-22mm}{34mm}{30mm}{2mm}{-10mm}{
                \coordinate (o) at (0mm,-3mm); 
                \node[draw,circle] (l1) at ($(o)+(-10mm,0mm)$) {1};
                \node[draw,circle] (l2) at ($(l1)+(2,0)$) {2};
                \node[draw,circle] (l3) at ($(l1)+(1,0)$) {3};
                \node[draw,circle] (l4) at ($(l2)+(0,-1)$) {6};
                \draw[] (l1) -- (l3) node[midway,above] {$a$};
                \draw[] (l3) -- (l2) node[midway,above] {$a$};
                \draw[ ] (l2) -- (l4) node[midway,right] {$a$};
                \node[draw,circle] (l6) at ($(l1)+(0,-1)$) {7};
                \draw[] (l1) -- (l6) node[midway,left] {$a$};
            }    

            \graphbox{\( C  \)}{40mm}{-22mm}{34mm}{30mm}{2mm}{-10mm}{
                \coordinate (o) at (0mm,-3mm); 
                \node[draw,circle] (l1) at ($(o)+(-10mm,0mm)$) {1};
                \node[draw,circle] (l2) at ($(l1)+(2,0)$) {2};
                \node[draw,circle] (l4) at ($(l2)+(0,-1)$) {6};
                \draw[ ] (l2) -- (l4) node[midway,right] {$a$};
                \node[ draw,circle] (l6) at ($(l1)+(0,-1)$) {7};
                \draw[ ] (l1) -- (l6) node[midway,left] {$a$};
            }    

            \graphbox{\( H  \)}{80mm}{-22mm}{45mm}{30mm}{2mm}{-10mm}{
                \coordinate (o) at (-5mm,-3mm); 
                \node[draw,circle] (l1) at ($(o)+(-10mm,0mm)$) {1};
                \node[draw,circle] (l2) at ($(l1)+(3,0)$) {2};
                \node[draw,circle] (l3) at ($(l1)+(1,0)$) {4};
                \node[draw,circle] (l4) at ($(l1)+(2,0)$) {5};
                \node[ draw,circle] (l5) at ($(l2)+(0,-1)$) {6};
                \node[ draw,circle] (l6) at ($(l1)+(0,-1)$) {7};
                \draw[ ] (l1) -- (l6) node[midway,left] {$a$};
                \draw[] (l1) -- (l3) node[midway,above] {$a$};
                \draw[] (l3) -- (l4) node[midway,above] {$b$};
                \draw[ ] (l4) -- (l2) node[midway,above] {$a$};
                \draw[ ] (l2) -- (l5) node[midway,right] {$a$};
            }    
            \node () at (37mm,-18mm) {\text{PO}};
            \node () at (78mm,-18mm) {\text{PO}};
            \node () at (37mm,-8mm) {\( \leftarrowtail \)}; % K -> L
            \node () at (77mm,-8mm) {\( \rightarrowtail \)}; % K -> R
            \node () at (17mm,-18mm) {\( m\ \downarrowtail \)};
            \node () at (37mm,-33mm) {\( \leftarrowtail \)};
            \node () at (52mm,-18mm) {\( \downarrowtail \)};
            \node () at (92mm,-18mm) {\( \downarrowtail \)};
            \node () at (77mm,-33mm) {\( \rightarrowtail \)}; % C -> H
        \end{tikzpicture}
        } 
    \end{center}
\end{example}
% The first condition ensures that 
% if \( R' \) (resp. \( R'' \)) forms an occurrence of \( X \) with a subgraph $C'$ of the context in the result graph (denote this occurrence \( x \)) around some interface elements, then \( h_{R'L}(R') \) (resp. \( h_{R'L}(R'') \)) forms a corresponding occurrence of \( X \) with $C'$ around the same interface elements in the host graph (denote this occurrence \( x' \)).  

The second condition guarantees that any $X$-occurrence implicitly created by the rewriting step has its corresponding $X$-occurrence not included in the context (Consequently, its corresponding $X$-occurrence is implicitly destroyed by the rewriting step). Counterexamples are provided in Example~\ref{ex:cond_2_necessaire} and Example~\ref{ex:cond_2_necessaire2}.  

% cond 2 necessaire  1
\begin{example} 
    \label{ex:cond_2_necessaire}
    Let $X$ be the graph 
    \tikz[baseline=-0.5ex]{ 
            \node (x) at (0,0) {$\bullet$}; 
            \node (y) at (1,0) {$\bullet$};
            \node (z) at (2,0) {$\bullet$};
            \draw[->] (x) -- (y)   {};
            \draw[->] (y) -- (z)   {};
    }. The following rewriting rule is potentially $X$-increasing.
    \begin{center}
        \resizebox{0.7\textwidth}{!}{
            \begin{tikzpicture}
                \graphbox{$L$}{0mm}{0mm}{34mm}{15mm}{2mm}{-5mm}{
                    \coordinate (o) at (0mm,-3mm); 
                    \node[draw,circle] (l1) at ($(o)+(-10mm,0mm)$) {1};
                    % \node[draw,circle] (l2) at ($(l1)+(2,0)$) {2};
                    \node[draw,circle] (l3) at ($(l1)+(1,0)$) {2};
                    \draw[->] (l3) -- (l1) node[midway,above] {};
                    % \draw[->] (l3) -- (l2) node[midway,above] {$a$};
                }     
                \graphbox{$K$}{40mm}{0mm}{24mm}{15mm}{2mm}{-5mm}{
                    \coordinate (o) at (5mm,-3mm); 
                    \node[draw,circle] (l1) at ($(o)+(-10mm,0mm)$) {1};
                    \node[draw,circle] (l2) at ($(l1)+(1,0)$) {2};
                    % \node[draw,circle] (l3) at ($(l1)+(1,0)$) {$\ $};
                    \draw[->] (l2) -- (l1) node[midway,above] {};
                    % \draw[->] (l3) -- (l2) node[midway,above] {$a$};
                }    
                \graphbox{$R$}{70mm}{0mm}{45mm}{15mm}{2mm}{-5mm}{
                    \coordinate (o) at (-5mm,-3mm); 
                    \node[draw,circle] (l1) at ($(o)+(-10mm,0mm)$) {1};
                    % \node[draw,circle] (l2) at ($(l1)+(3,0)$) {2};
                    \node[draw,circle] (l3) at ($(l1)+(1,0)$) {2};
                    \node[draw,circle] (l4) at ($(l1)+(2,0)$) {3};
                    \draw[->] (l3) -- (l1) node[midway,above] {};
                    \draw[->] (l3) -- (l4) node[midway,above] {};
                    % \draw[->] (l4) -- (l2) node[midway,above] {$a$};
                }    
                \node () at (37mm,-8mm) {$\leftarrowtail$};
                \node () at (67mm,-8mm) {$\rightarrowtail$};
            \end{tikzpicture}
        }
    \end{center}
    As shown in Example~\ref{ex:rx_counter_ex}, \( D(R,X) \) contains the graph $R'$:
    \raisebox{2pt}{
        \scalebox{0.6}{\tikz[baseline=-0.5ex]{
        \node [draw,circle] (x) at (0,0) {2};
        \node[draw,circle] (y) at (1,0) {3};
        \draw[->] (x) -- (y) {};
    }}}. The unique monomorphism $h_{R'L}:R' \rightarrowtail L$ fails the second condition of Definition~\ref{def:creates_more_x_on_the_left}. 
    %  as illustrated below 
    % \begin{center}
    %     \resizebox{0.6\textwidth}{!}{
    %     \begin{tikzpicture}
    %             \graphbox{$R'$}{40mm}{-20mm}{24mm}{15mm}{2mm}{-5mm}{
    %                 \coordinate (o) at (5mm,-3mm); 
    %                 \node[draw,circle] (l1) at ($(o)+(-10mm,0mm)$) {2};
    %                 \node[draw,circle] (l2) at ($(l1)+(1,0)$) {3};
    %                 % \node[draw,circle] (l3) at ($(l1)+(1,0)$) {$\ $};
    %                 \draw[->] (l1) -- (l2) node[midway,above] {};
    %                 % \draw[->] (l3) -- (l2) node[midway,above] {$a$};
    %             }    
    %             \graphbox{$L$}{70mm}{-20mm}{45mm}{15mm}{2mm}{-5mm}{
    %                 \coordinate (o) at (-5mm,-3mm); 
    %                 \node[draw,circle] (l1) at ($(o)+(-10mm,0mm)$) {3};
    %                 % \node[draw,circle] (l2) at ($(l1)+(3,0)$) {2};
    %                 \node[draw,circle] (l2) at ($(l1)+(1,0)$) {2};
    %                 \draw[->] (l2) -- (l1) node[midway,above] {};
    %                 % \draw[->] (l4) -- (l2) node[midway,above] {$a$};
    %             }    
    %             \node () at (67mm,-28mm) {$\rightarrowtail$};
    %         \end{tikzpicture}
    %     }
    % \end{center}
    For any rewriting step using this rule, the corresponding $X$-occurrence relative to $h_{R'L}$ of any implicitly created $X$-occurrence is included in the context.
\end{example}
% cond 2 necessaire  2
\begin{example}
    \label{ex:cond_2_necessaire2}
    Let $X$ be the graph 
    \tikz[baseline=-0.5ex]{ 
            \node (x) at (0,0) {$\bullet$}; 
            \node (y) at (1,0) {$\bullet$};
            \node (z) at (2,0) {$\bullet$};
            \draw[->] (x) -- (y)   {};
            % \draw[->] (y) -- (z)   {};
    }. The following rewriting rule is potentially $X$-increasing.
    \begin{center}
        \resizebox{0.7\textwidth}{!}{
            \begin{tikzpicture}
                \graphbox{$L$}{0mm}{0mm}{34mm}{15mm}{2mm}{-5mm}{
                    \coordinate (o) at (0mm,-3mm); 
                    \node[draw,circle] (l1) at ($(o)+(-10mm,0mm)$) {1};
                    \node[draw,circle] (l3) at ($(l1)+(1,0)$) {2};
                }     
                \graphbox{$K$}{40mm}{0mm}{24mm}{15mm}{2mm}{-5mm}{
                    \coordinate (o) at (5mm,-3mm); 
                    \node[draw,circle] (l1) at ($(o)+(-10mm,0mm)$) {1};
                    \node[draw,circle] (l2) at ($(l1)+(1,0)$) {2};
                }    
                \graphbox{$R$}{70mm}{0mm}{45mm}{15mm}{2mm}{-5mm}{
                    \coordinate (o) at (-5mm,-3mm); 
                    \node[draw,circle] (l1) at ($(o)+(-10mm,0mm)$) {1};
                    \node[draw,circle] (l3) at ($(l1)+(1,0)$) {2};
                    \node[draw,circle] (l4) at ($(l1)+(2,0)$) {3};
                }    
                \node () at (37mm,-8mm) {$\leftarrowtail$};
                \node () at (67mm,-8mm) {$\rightarrowtail$};
            \end{tikzpicture}
        }
    \end{center}
    The set \( D(R,X) \) contains $R'$:
    \raisebox{2pt}{ \scalebox{0.6}{\tikz[baseline=-0.5ex]{
        \node [draw,circle] (x) at (0,0) {1};
        \node[draw,circle] (y) at (1,0) {3};
    }}}. \todo{D(R,X) !!!!!}
    Despite the existence of a unique monomorphism $h_{R'L}: R' \rightarrowtail L$ preserving interface elements, this rule fails the second condition of Definition~\ref{def:creates_more_x_on_the_left}.
    For any rewriting step using this rule, the corresponding $X$-occurrence relative to $h_{R'L}$ of any implicitly created $X$-occurrence is included in the context.
\end{example}
\begin{example}
    \label{ex:cond3_necessaire}
    Let $X$ be the graph 
    \tikz[baseline=-0.5ex]{ 
            \node (x) at (0,0) {$\bullet$}; 
            \node (y) at (1,0) {$\bullet$};
            \node (z) at (2,0) {$\bullet$};
            \draw[->] (x) -- (y)   {};
            \draw[->] (y) -- (z) {};
    }. The following rewriting rule is potentially $X$-increasing.
    \begin{center}
        \resizebox{0.6\textwidth}{!}{
            \begin{tikzpicture}
                \graphbox{$L$}{0mm}{0mm}{30mm}{20mm}{0}{0}{
                    \node[draw,circle]  (x) at (-6mm,-12mm) {x};
                    \node[draw,circle] (y) at (6mm,-12mm) {};
                    % \node[draw,circle]  (z) at (6mm,0mm) {};
                    \draw[->]  (x) to (y);
                    % \draw[->] (y) to[bend right=20] (z);
                    % \draw[->]  (z) to[bend right=20] (y);
                }
                \graphbox{$K$}{40mm}{0mm}{30mm}{20mm}{0}{0}{
                    \node[draw,circle]  (x) at (-6mm,-12mm) {x};
                    % \node[draw,circle]  (y) at (6mm,-12mm) {y};
                }
                \graphbox{$R$}{80mm}{0mm}{30mm}{20mm}{0}{0}{
                    \node[draw,circle]  (x) at (-6mm,-12mm) {x};
                        \node[draw,circle]  (y) at (6mm,-8mm) {y};
                        \node[draw,circle]  (z) at (6mm,-16mm) {z};
                        \draw[->]  (x) to (y);
                        \draw[->]  (x) to (z);
                }
                \node () at (35mm,-12mm) {$\leftarrowtail$};
                \node () at (75mm,-12mm) {$\rightarrowtail$};
            \end{tikzpicture}
        }
    \end{center}
    The set \( D(R,X) \) contains exactly two elements $R'$:
    \raisebox{2pt}{ 
        \scalebox{0.6}{\tikz[baseline=-0.5ex]{
        \node [draw,circle] (x) at (0,0) {x};
        \node[draw,circle] (y) at (1,0) {y};
        \draw[->] (x) -- (y) {};
    }}} and $R''$:\raisebox{2pt}{ 
        \scalebox{0.6}{\tikz[baseline=-0.5ex]{
        \node [draw,circle] (x) at (0,0) {x};
        \node[draw,circle] (y) at (1,0) {z};
        \draw[->] (x) -- (y) {};
    }}}. 
    Despite unique monomorphisms \( h_{R'L}: R' \rightarrowtail L \) and \( h_{R''L}: R'' \rightarrowtail L \) preserving interface elements, they fail the third condition of Definition~\ref{def:creates_more_x_on_the_left}. 
    For any rewriting step using this rule, implicitly created $X$-occurrences with the same subgraph of the context have the same corresponding $X$-occurrence. 
\end{example}
 
\begin{example}
    \label{ex:cond4_necessaire}
    Let $X$ be the graph 
    \tikz[baseline=-0.5ex]{ 
            \node (x) at (0,0) {$\bullet$}; 
            \node (y) at (1,0) {$\bullet$};
            \node (z) at (2,0) {$\bullet$};
            \draw[->] (x) -- (y)   {};
    }, which has an isolated node. The following rewriting rule is potentially $X$-increasing.
    \begin{center}
        \resizebox{0.6\textwidth}{!}{
            \begin{tikzpicture}
                \graphbox{$L$}{0mm}{0mm}{30mm}{20mm}{0}{0}{
                    \node[draw,circle]  (x) at (-6mm,-10mm) {x};
                    \node[draw,circle] (y) at (6mm,-10mm) {y};
                }
                \graphbox{$K$}{40mm}{0mm}{30mm}{20mm}{0}{0}{
                    \node[draw,circle]  (x) at (-6mm,-10mm) {x};
                }
                \graphbox{$R$}{80mm}{0mm}{30mm}{20mm}{0}{0}{
                    \node[draw,circle]  (x) at (-6mm,-10mm) {x};
                        \node[draw,circle]  (y) at (6mm,-6mm) {z};
                        \node[draw,circle]  (z) at (6mm,-14mm) {w};
                }
                \node () at (35mm,-10mm) {$\leftarrowtail$};
                \node () at (75mm,-10mm) {$\rightarrowtail$};
            \end{tikzpicture}
        }
    \end{center}
    The set \( D(R,X) \) contains exactly two elements $R'$:
    \raisebox{2pt}{\scalebox{0.6}{\tikz[baseline=-0.5ex]{
        \node [draw,circle] (x) at (0,0) {x};
        \node[draw,circle] (y) at (1,0) {z};
        % \draw[->] (x) -- (y) {};
    }}} and $R''$:\raisebox{2pt}{\scalebox{0.6}{\tikz[baseline=-0.5ex]{
        \node [draw,circle] (x) at (0,0) {x};
        \node[draw,circle] (y) at (1,0) {w};
        % \draw[->] (x) -- (y) {};
    }}}. 
    There are unique monomorphisms $h_{R'L}:R' \rightarrowtail L$ and $h_{R''L}:R'' \rightarrowtail L$ preserving interface elements, but they fail the fourth condition of Definition~\ref{def:creates_more_x_on_the_left}.
    For any rewriting step using this rule, implicitly created $X$-occurrences with the same subgraph of the context have the same corresponding $X$-occurrence.
\end{example}
% \subsection{Solution to the key challenge}
% \label{sec:solution_to_the_key_challenge}
% The following lemma ensures that if $\rho$ is $X$-non-increasing, then, for any rewriting step using $\rho$, the cardinality of the set of \( X \)-occurrences implicitly destroyed by the rewriting step is at least as large as the cardinality of the set of \( X \)-occurrences implicitly created by the rewriting step. 
% \subsection{Termination}
% \label{sec:subgraph_counting_antipattern:termination}
We define the following sets of monomorphisms from $X$ to $G$ to facilitate the discussion of the termination criterion. They are almost identical to the ones defined in~\autoref{subgraph_counting:notation:mono_sets}. The only difference is that the first parameter is a ruler-graph $\mathcal{X}=(X,f)$ defined in~\autoref{antipattern:def:ruler_graph} instead of the underlying graph $X$.
\begin{notation}
    \label{antipattern:notation:mono_sets}
     \ \newline
\noindent
\begin{minipage}{0.69\textwidth}
\vspace{2mm}
\noindent 
  Let \( \mathcal{X}\) be a ruler-graph with underlying graph $X$. The disjoint union of two sets \( S \) and \( S' \) will be denoted by \( S \uplus S' \). Let \( X, A, B, G \) be graphs, and let \( \alpha \colon A \to G \) and \( \beta \colon B \to G \) be morphisms as illustrated on the right. We define the following sets of monomorphisms from $X$ in $G$ based on their interaction with $\alpha$ and $\beta$:    
\end{minipage}% 
\begin{minipage}{0.30\textwidth}
    \hfill
        \resizebox{0.9\textwidth}{!}{
        % \begin{tikzpicture}
        %     \node (B) at (2,0) {$A$}; 
        %     \node (C) at (0,2) {$B$}; 
        %     \node (D) at (0,0) {$G$}; 
        %     \draw [>->] (B) to node [below,label,pos=0.45] {$\alpha$} (D); 
        %     \draw [>->] (C) to node [left,label,pos=0.45] {$\beta$} (D);
        % \end{tikzpicture}
            \begin{tikzpicture}
            \node (X) at (-1.5,-1) {$X$};
            \node (B) at (2,0) {$A$}; 
            \node (C) at (0,2) {$B$}; 
            \node (D) at (0,0) {$G$}; 
            \draw [>->] (B) to node [above,label,pos=0.45] {$\alpha$} (D); 
            \draw [>->] (C) to node [right,label,pos=0.45] {$\beta$} (D);
            \draw [>->,dashed] (X) to node [above,label,pos=0.45] {$\iota$} (D);
            \draw [>->,dashed] (X) to node [below,label,pos=0.45] {$\zeta$} (B);
            \draw [>->,dashed] (X) to node [above,label,pos=0.45] {$\eta$} (C);
        \end{tikzpicture}
        }
\end{minipage} 
    \begin{align*}
        \operatorname{Mono}(\mathcal{X},G) &= \operatorname{Mono}(X,G), 
        \\
        \operatorname{Mono}(\mathcal{X},G,\alpha) &= \left\{ \iota \colon X \rightarrowtail G 
        \;\middle|\; 
        \exists \zeta \colon X \rightarrowtail A.\, \iota = \zeta \star \alpha \right\}, 
        \\
        \operatorname{Mono}(\mathcal{X},G,\lnot \alpha) &= \left\{ \iota \colon X \rightarrowtail G 
        \;\middle|\; 
        \nexists \zeta \colon X \rightarrowtail A.\, \iota = \zeta \star \alpha \right\}, 
        \\
        \operatorname{Mono}(\mathcal{X},G,\lnot \alpha, \beta) &= \left\{ 
            \iota \colon X \rightarrowtail G \;\middle|\; 
                \begin{aligned}  
                    &(\nexists \zeta \colon X \rightarrowtail A.\, \iota = \zeta \star \alpha) \\ 
                    &\land (\exists \eta \colon X \rightarrowtail B.\, \iota = \eta \star \beta)
                \end{aligned}
        \right\},
        \\
        \operatorname{Mono}(\mathcal{X},G,\lnot \alpha, \lnot \beta) &= \left\{ 
            \iota \colon X \rightarrowtail G \;\middle|\; 
                \begin{aligned}
                    &(\nexists \zeta \colon X \rightarrowtail A.\, \iota = \zeta \star \alpha) \\
                    &\land (\nexists \eta \colon X \rightarrowtail B.\, \iota = \eta \star \beta)
                \end{aligned}
        \right\}.
    \end{align*}
    For a set $\operatorname{Mono}(\mathcal{X},\dots)$, we define $\operatorname{Mono}(\mathcal{X},\dots)_{\operatorname{NF}}$ as the subset of \( X \)-occurrences whose images are not included in any occurrence of its forbidden context if exists, and  $\operatorname{Mono}(\mathcal{X},\dots)_{\operatorname{F}}$ as the subset of \( X \)-occurrences whose images are included in some occurrences of its forbidden context if exists. 
\end{notation}
\noindent
\begin{minipage}{0.7\textwidth}\setlength{\parindent}{1em}
    Let $\mathcal{X}$ be a ruler-graph with underlying graph $X$ and \( \rho = (L \overset{l}{\leftarrowtail} K \overset{r}{\rightarrowtail} R) \) be a rule. 
    Consider the DPO diagram illustrated on the right. 
\end{minipage}%
\begin{minipage}{0.29\textwidth}
    \hfill
    \begin{tikzpicture}
        % [node distance=15mm]
        \node (I) {$K$};
        \node (L) [left of=I] {$L$};
        \node (R) [right of=I] {$R$};
        \node (G) [below of=L] {$G$};
        \node (C) [below of=I] {$C$};
        \node (H) [below of=R] {$H$}; 
        \draw [>->] (I) to  node [midway,below] {$l$} (L);
        \draw [>->] (I) to  node [midway,below] {$r$} (R);
        \draw [>->] (L) to node [midway,right] {$m$} (G);
        \draw [>->] (I) to  node [midway,right] {$u$} (C); 
        \draw [>->] (R) to  node [midway,right] {$m'$} (H);
        \draw [>->] (C) to node [midway,above] {$l'$} (G);
        \draw [>->] (C) to node [midway,above] {$r'$} (H);
        % \node [at=($(I)!.5!(G)$)] {\normalfont PO};
        % \node [at=($(I)!.5!(H)$)] {\normalfont PO};
      \end{tikzpicture}
\end{minipage} 

\vspace{1mm}
% The following decompositions hold:
% \begin{flalign*}
%     \operatorname{Mono}(\mathcal{X},G)_{\operatorname{NF}} &= 
%     \operatorname{Mono}(\mathcal{X},G,m)_{\operatorname{NF}}
%     \uplus
%     \operatorname{Mono}(\mathcal{X},G,\lnot m, l')_{\operatorname{NF}} 
%     \uplus
%     \operatorname{Mono}(\mathcal{X},G,\lnot m, \lnot l')_{\operatorname{NF}}
%     \\
%     \operatorname{Mono}(\mathcal{X},H)_{\operatorname{NF}} &= 
%     \operatorname{Mono}(\mathcal{X},H,m')_{\operatorname{NF}}
%     \uplus
%     \operatorname{Mono}(\mathcal{X},H,\lnot m', r')_{\operatorname{NF}}
%     \uplus
%     \operatorname{Mono}(\mathcal{X},H,\lnot m', \lnot r')_{\operatorname{NF}}
% \end{flalign*}

The set $\operatorname{Mono}(\mathcal{X},G)_{\operatorname{NF}}$ can be partitioned into three disjoint sets:
$$
    \operatorname{Mono}(\mathcal{X},G,m)_{\operatorname{NF}}
    \uplus
    \operatorname{Mono}(\mathcal{X},G,\lnot m, l')_{\operatorname{NF}} 
    \uplus
    \operatorname{Mono}(\mathcal{X},G,\lnot m, \lnot l')_{\operatorname{NF}}
$$
and $\operatorname{Mono}(\mathcal{X},H)_{\operatorname{NF}}$ can be partitioned into three disjoint sets:
$$
\operatorname{Mono}(\mathcal{X},H,m')_{\operatorname{NF}}
    \uplus
    \operatorname{Mono}(\mathcal{X},H,\lnot m', r')_{\operatorname{NF}}
    \uplus
    \operatorname{Mono}(\mathcal{X},H,\lnot m', \lnot r')_{\operatorname{NF}}
$$
Thus, we have 
\begin{flalign}
    &\card{\operatorname{Mono}(\mathcal{X},G)_{\operatorname{NF}}} - 
    \card{\operatorname{Mono}(\mathcal{X},H)_{\operatorname{NF}}} \nonumber
    \\=
    &(\card{\operatorname{Mono}(\mathcal{X},G,m)_{\operatorname{NF}}} 
        -  
    \card{\operatorname{Mono}(\mathcal{X},H,m')_{\operatorname{NF}}}) 
    +  \nonumber
    \\
    &(
        \card{\operatorname{Mono}(\mathcal{X},G,\lnot m, l')_{\operatorname{NF}}}
             - 
        \card{\operatorname{Mono}(\mathcal{X},H,\lnot m', r')_{\operatorname{NF}}}) +  \nonumber \\ 
    &(
        \card{\operatorname{Mono}(\mathcal{X},G,\lnot m, \lnot l')_{\operatorname{NF}}} 
            - 
        \card{\operatorname{Mono}(\mathcal{X},H,\lnot m', \lnot r')_{\operatorname{NF}}} \nonumber
    )
\end{flalign}
In the remainder of this section, we suppose that $\rho^{-1}$ is $F$-non-increasing if $\mathcal{X}= (X,f:X \rightarrowtail F)$ and, $\rho$ and $\rho^{-1}$ are $X$-non-increasing. To estimate $\card{\operatorname{Mono}(\mathcal{X},G)_{\operatorname{NF}}} - 
    \card{\operatorname{Mono}(\mathcal{X},H)_{\operatorname{NF}}}$, we establish the following lemmas.
% The following lemmas can be used to approximate $\card{\operatorname{Mono}(\mathcal{X},G)_{\operatorname{NF}}} - 
%     \card{\operatorname{Mono}(\mathcal{X},H)_{\operatorname{NF}}}$
\begin{lemma}
    \label{lem:xglnotmlp_xhlnotmrp}
        Let $\mathcal{X}$ be a ruler-graph with underlying graph $X$ and \( \rho = (L \overset{l}{\leftarrowtail} K \overset{r}{\rightarrowtail} R) \) be a rule. 
        Suppose that $\rho^{-1}$ is $F$-non-increasing if $\mathcal{X} = (X,f:X \rightarrowtail F)$.
        We have 
    $$\card{\operatorname{Mono}(\mathcal{X},G,\lnot m, l')_{\operatorname{NF}}} \geq
        \card{\operatorname{Mono}(\mathcal{X},H,\lnot m', r')_{\operatorname{NF}}}$$
\end{lemma} 
See~\textsection~\ref{antipattern:proof:lem:xglnotmlp_xhlnotmrp} for the proof of this lemma.
\begin{lemma}
    \label{lem:xglnotmlnotlp_xhlnotmrnotrp}
        Let $\mathcal{X}$ be a ruler-graph with underlying graph $X$ and let \( \rho = (L \overset{l}{\leftarrowtail} K \overset{r}{\rightarrowtail} R) \) be a rule. Suppose that $\rho$ and $\rho^{-1}$ are $X$-non-increasing. We have
    $$ 
        \card{\operatorname{Mono}(\mathcal{X},G,\lnot m, \lnot l')_{\operatorname{NF}}} \geq
        \card{\operatorname{Mono}(\mathcal{X},H,\lnot m', \lnot r')_{\operatorname{NF}}}
    $$
\end{lemma}
See~\textsection~\ref{antipattern:proof:lem:xglnotmlnotlp_xhlnotmrnotrp} for the proof of this lemma.
Under the assumptions of the above two lemmas, we can estimate the following inequality.
 \begin{flalign*}
     \card{\operatorname{Mono}(\mathcal{X},G)_{\operatorname{NF}}} - 
     \card{\operatorname{Mono}(\mathcal{X},H)_{\operatorname{NF}}} 
     \geq
     \card{\operatorname{Mono}(\mathcal{X},G,m)_{\operatorname{NF}}} - \card{\operatorname{Mono}(\mathcal{X},H,m')_{\operatorname{NF}}}
 \end{flalign*}

\begin{definition}
    \label{antipattern:def:gamma_l_rho_x}
    Let $\mathcal{X}$ be a ruler-graph with underlying graph $X$. 
    Let \( \rho = (L \overset{l}{\leftarrowtail} K \overset{r}{\rightarrowtail} R) \) be a rule.
    Define $\Gamma(\operatorname{Mono}(\mathcal{X},L)_{\operatorname{NF}})$ as the subset of $\operatorname{Mono}(\mathcal{X},L)_{\operatorname{NF}}$ consisting of $X$-occurrences whose images are not included in any occurrence of the forbidden context in $L$,
    but, in some rewriting step $G \Rightarrow_\rho^m H$, their images may possibly be included in some occurrence of the forbidden context in the host graph $G$.
    Formally,  $\Gamma(\operatorname{Mono}(\mathcal{X},L)_{\operatorname{NF}}) = \emptyset$ if $\mathcal{X} = (X)$, and if $\mathcal{X} = (X, f:X \rightarrowtail F)$, then a monomorphism $h_{XL}:X \rightarrowtail L$ is in $\Gamma(\operatorname{Mono}(\mathcal{X},L)_{\operatorname{NF}})$ iff the following hold:
    \begin{itemize}
        \item there does not exist $h_{FL}:F \rightarrowtail L$ such that $f \star h_{FL} = h_{XL}$,
         \item there exists a graph $C$ and a monomorphism $h_{KC}:K \overset{u}{\rightarrowtail} C$ such that, let $L \overset{m}{\rightarrowtail} G \overset{l'}{\leftarrowtail} C$ be the pushout of the span $L \overset{l}{\leftarrowtail} K \overset{u}{\rightarrowtail} C$, as illustrated below,
         there exist a monomorphism $h_{FG} : F \rightarrowtail G$ such that 
        %  $h_{XL} \star h_{LG} = h_{XF} \star h_{FG}$
         $h_{XL} \star m = h_{XF} \star h_{FG}$. 

         \resizebox{0.3\textwidth}{!}{
            \begin{tikzpicture}
                    \node (k) at (0,1) {K};
                    \node (l) at (-2,1) {L};
                    \node (c) at (0,-1) {C};
                    \node (g) at (-2,-1) {G};
                    \node (x) at (-4,-2) {X};
                    \node (f) at (-2,-2) {F};
                    \draw[<-<]  (l) -- (k) node [midway,below] {l};
                    \draw[>->] (c) -- (g) node [midway, above] {l'};
                    \draw[>->] (l) -- (g) node[midway, right] {m};
                    \draw[>->] (k) -- (c) node[midway, left] {u};
                    \draw[>->] (x) -- (l) node[midway, left] {$h_{XL}$};
                    \draw[>->] (x) -- (f) node[midway, above] {$f$};
                    \draw[>->] (f) -- (g) node[midway, right] {$h_{FG}$};
                    \node () [at=($(l)!0.5!(c)$)] {PO};
                \end{tikzpicture}
         }
    \end{itemize}
    Furthermore, we define $\Lambda(\mathcal{X},\rho) \overset{\operatorname{def}}{=} (\card{\operatorname{Mono}(\mathcal{X},L)_{\operatorname{NF}}} - 
    \card{\Gamma(\operatorname{Mono}(\mathcal{X},L)_{\operatorname{NF}})}) -
   \card{\operatorname{Mono}(\mathcal{X},R)_{\operatorname{NF}}}$ to improve the readability of the following discussion.
%    \trackedtext{Intuitively, $\Lambda(\mathcal{X},\rho)$ is the number of $X$-occurrences in $L$ ...}
\end{definition}
Note that in category \textbf{Graph}, the pushout of two arrows always exists~\cite[p.188]{corradini1997algebraic}. Which justifies the existence of the pushout square in the second condition of the above definition.
Furthermore, since $X$ and the forbidden context (if exists) are finite graphs, $\Lambda(\mathcal{X},\rho)$ can be precisely computed. Consequently, it provides a lower bound for the change in the number of $X$-occurrences in a rewriting step using $\rho$, as shown in the following lemma whose proof is given in~\textsection~\ref{antipattern:proof:lem:xgm_xhmp_xl_xr}.

% \begin{lemma}
%     \label{lem:decomp_w_u}
%     \ \newline
%     \noindent
%     \begin{minipage}{0.7\textwidth}
%         Let $X$ be a ruler-graph. For a pushout square as shown on the right, we have 
%         % $\card{\operatorname{Mono}(X, B)} = \card{\operatorname{Mono}(X, D, \beta')}$ and $\card{\operatorname{Mono}(X, C, \lnot \beta)} = \card{\operatorname{Mono}(X, D, \lnot \beta', \alpha')}$.
%         \begin{flalign*}
%             \card{\operatorname{Mono}(X, B)} &= \card{\operatorname{Mono}(X, D, \beta')}
%             \\
%             \card{\operatorname{Mono}(X, C, \lnot \beta)} &= \card{\operatorname{Mono}(X, D, \lnot \beta', \alpha')}
%         \end{flalign*}
%     \end{minipage}
%     \hfill
%     \begin{minipage}{0.3\textwidth}
%         \hfill
%         \begin{tikzpicture}
%             \node (A) {$A$};
%             \node [below of=A] (B) {$B$}; 
%             \node [left of=A] (C) {$C$}; 
%             \node [left of=B] (D) {$D$}; 
%             \begin{scope}[nodes=rectangle]          
%             \draw [>->] (A) to node [right,label,pos=0.5] {$\alpha$} (B);
%             \draw [>->] (A) to node [above,label,pos=0.5] {$\beta$} (C);
%             \draw [>->] (B) to node [below,label,pos=0.45] {$\beta'$} (D); 
%             \draw [>->] (C) to node [left,label,pos=0.45] {$\alpha'$} (D);
%             \end{scope}
%         \end{tikzpicture}
%     \end{minipage} 
% \end{lemma}
% \begin{proof}
%     See the Appendix,~\autoref{proof:dcomp_w_u}.
%  \end{proof}
% \begin{lemma}
%     \label{lem:xgm_xhmp_xl_xr}
%     % If the number of $X$-occurrences that are not  included in any occurrences of forbidden context $F \in F_x$ in a $\rho$-rewriting step is predictable, then
%     $
%         \card{\operatorname{Mono}(\mathcal{X},G,m)_{\operatorname{NF}}} - 
%         \card{\operatorname{Mono}(\mathcal{X},H,m')_{\operatorname{NF}}} 
%         \geq 
%         \card{\operatorname{Mono}(X,L)} -
%         \card{\operatorname{Mono}(X,R)}
%     $
% \end{lemma}
% \begin{proof}
%     See the Appendix, \textsection~\ref{proof:lem:xgm_xhmp_xl_xr}.
%  \end{proof}

\begin{lemma}
    \label{antipattern:lem:xgm_xhmp_xl_xr}
     Let $\mathcal{X} = (X, f:X \rightarrowtail F)$ be a ruler-graph and \( \rho = (L \overset{l}{\leftarrowtail} K \overset{r}{\rightarrowtail} R) \) be a rule. 
   We have
    \begin{flalign*}
        &\card{\operatorname{Mono}(\mathcal{X},G,m)_{\operatorname{NF}}} - 
        \card{\operatorname{Mono}(\mathcal{X},H,m')_{\operatorname{NF}}} 
        \geq
        \Lambda(\mathcal{X},\rho)
    \end{flalign*}
\end{lemma}
 Thus, by the unnumbered equation preceding~\autoref{antipattern:def:gamma_l_rho_x},  we can conclude that
 \begin{flalign}
         \card{\operatorname{Mono}(\mathcal{X},G)_{\operatorname{NF}}} - 
     \card{\operatorname{Mono}(\mathcal{X},H)_{\operatorname{NF}}} 
     \geq 
    \Lambda(\mathcal{X},\rho)
     \label{eq:mono_x_g_nf_mono_x_h_nf_geq}
 \end{flalign}
Consequently, we have  
\begin{flalign*}
    &w_{s_\mathbb{X}}(G) - w_{s_\mathbb{X}}(H)
    \\
   \overset{\operatorname{def}}{=}&\sum_{\mathcal{X} \in \mathbb{X}}^{}s_\mathbb{X}(X) * m_X(G) - \sum_{\mathcal{X} \in \mathbb{X}}^{}s_\mathbb{X}(X) * m_X(H)
   \\
   \overset{\operatorname{def}}{=}&\sum_{\mathcal{X} \in \mathbb{X}}^{}s_\mathbb{X}(X) * |\operatorname{Mono}(\mathcal{X},G)_{\operatorname{NF}}| - \sum_{\mathcal{X} \in \mathbb{X}}^{}s_\mathbb{X}(X) * |\operatorname{Mono}(\mathcal{X},H)_{\operatorname{NF}}|
   \\
   =&\sum_{\mathcal{X} \in \mathbb{X}}^{}s_\mathbb{X}(X) * \left( \card{\operatorname{Mono}(\mathcal{X},G)_{\operatorname{NF}}} - 
   \card{\operatorname{Mono}(\mathcal{X},H)_{\operatorname{NF}}} \right)
   \\
   \geq&\sum_{\mathcal{X} \in \mathbb{X}}^{}s_\mathbb{X}(X) * \Lambda(\mathcal{X},\rho)
   & \text{by~\eqref{eq:mono_x_g_nf_mono_x_h_nf_geq}}
%    \\
%    =&\sum_{\mathcal{X} \in \mathbb{X}}^{}s_\mathbb{X}(X) * \card{\operatorname{Mono}(X,G)} -  
%    \sum_{\mathcal{X} \in \mathbb{X}}^{}s_\mathbb{X}(X) *  \card{\operatorname{Mono}(X,H)}  
%     \\
%     \overset{\operatorname{def}}{=}&\sum_{\mathcal{X} \in \mathbb{X}}^{}s_\mathbb{X}(X) * m_X(L) - \sum_{\mathcal{X} \in \mathbb{X}}^{}s_\mathbb{X}(X) * m_X(R)
%     \\
%     \overset{\operatorname{def}}{=}& w_{s_\mathbb{X}}(L) - w_{s_\mathbb{X}}(R)
\end{flalign*} 
Which proves the following key lemma of this section.
\begin{lemma}[Decreasing step]
    \label{antipattern:lem:w_g_geq_w_h_leq}
    Let $\rho = (L \overset{l}{\leftarrowtail} K \overset{r}{\rightarrowtail} R)$ be an injective DPO rewriting rule,
    \( \mathbb{X} \) a set of ruler-graphs,
    \( s_{\mathbb{X}} \colon \mathbb{X} \to \mathbb{N} \) a weight function,
    and \( G \Rightarrow_{\rho,\mathfrak{M}} H \) a rewriting step. 
    Suppose that for every \( \mathcal{X} \in \mathbb{X} \) with underlying graph $X$, 
    % the number of $X$-occurrences that are not included in any occurrences of forbidden context $F \in F_x$ in a $\rho$-rewriting step is predictable,
    $\rho$ is $X$-non-increasing, and if $\mathcal{X}= (X,f:X \rightarrowtail F)$ then $\rho^{-1}$ is $X$-non-increasing and $F$-non-increasing. We have
     $$
        w_{s_\mathbb{X}}(G) - w_{s_\mathbb{X}}(H) 
        \geq 
        \sum_{\mathcal{X} \in \mathbb{X}}^{}s_\mathbb{X}(X) * \Lambda(\mathcal{X},\rho)
    $$
\end{lemma}
Finally, our main result follows.
\begin{theorem}[Termination] 
    \label{antipattern:thm:termination_grs} 
    Let \(\mathcal{A}\) and \(\mathcal{B}\) be sets of injective DPO rewriting rules, $\mathbb{X}$ a set of ruler-graphs and $s_\mathbb{X}$ a weight function. If the following hold:
    \begin{enumerate}
        \item  for every $\rho \in \mathcal{A} \cup \mathcal{B}$ and for every \( \mathcal{X} \in \mathbb{X} \) with underlying graph $X$, 
        % the number of $X$-occurrences that are not included in any occurrences of the forbidden context $F \in F_x$ in a $\rho$-rewriting step is predictable,
        $\rho$ is $X$-non-increasing, and if $\mathcal{X}= (X,f:X \rightarrowtail F)$ then $\rho^{-1}$ is $X$-non-increasing and $F$-non-increasing,
        % $\rho^{-1}$ is $F$-non-increasing if $\mathcal{X}= (X,f:X \rightarrowtail F)$ and, $\rho$ and $\rho^{-1}$ are $X$-non-increasing
        \item for every \(\rho \in \mathcal{A}\), we have
        % \( w_{s_\mathbb{X}}(lhs(\rho)) > w_{s_\mathbb{X}}(rhs(\rho)) \),
        $ \sum_{\mathcal{X} \in \mathbb{X}}^{}s_\mathbb{X}(X) * 
            \Lambda(\mathcal{X},\rho) > 0 $
        \item for every \(\rho \in \mathcal{B}\), we have   
        % \( w_{s_\mathbb{X}}(lhs(\rho)) \geq w_{s_\mathbb{X}}(rhs(\rho)) \).
        $ 
            \sum_{\mathcal{X} \in \mathbb{X}}^{}s_\mathbb{X}(X) * \Lambda(\mathcal{X},\rho) \geq 0 
        $
    \end{enumerate}
    Then \(\Rightarrow_{\mathcal{A},\mathcal{M}}\) is terminating relative to \(\Rightarrow_{\mathcal{B},\mathcal{M}}\).
\end{theorem}
\begin{proof}
    See~\textsection~\ref{antipattern:proof:thm:termination_grs}.
\end{proof}
\begin{remark}
    This work focuses on deriving sufficient conditions for termination, deferring the construction of ruler-graphs to future work.
\end{remark}  

\section{Examples}
\label{antipattern:sec:examples}
% \begin{example}
%   \label{antipattern:example:contrib_variant}
%   The following rewriting rule is obtained from a rewriting rule in~\cite[Example 6]{plump2018modular} by removing all edges of the interface graph and adding a node with the identifier $5$ and three self-loops labeled $0$ to the right-hand-side graph.
%   \begin{center}
%       \resizebox{0.8\textwidth}{!}{
%           \begin{tikzpicture}
%               \graphbox{$L$}{0mm}{0mm}{35mm}{35mm}{2mm}{-5mm}{
%                   \coordinate (delta) at (0,-18mm);
%                   \node[draw,circle] (l1) at ($(delta)+(-1,1.5)$) {1};
%                   \node[draw,circle] (l2) at ($(delta)+(1,1.5)$) {2};
%                   \node[draw,circle] (l3) at ($(delta)+(0,0)$) {3};
%                   \draw[->] (l1) -- (l3) node[midway,left] {$s$};
%                   \draw[->] (l2) -- (l3) node[midway,right] {$s$};
%                   \draw[->] (l3) edge [loop below] node {0} (l3);
%               }
%               \graphbox{$K$}{40mm}{0mm}{35mm}{35mm}{2mm}{-5mm}{
%                   \coordinate (delta) at (0,-18mm);
%                   \coordinate (interfaceorigin) at ($(delta) +(5,0)$);
%                   \node[draw,circle] (r1) at ($(delta) +(-1,1.5)$) {1};
%                   \node[draw,circle] (r2) at ($(delta) +(0.5,1.5)$) {2};
%                   \node[draw,circle] (r3) at ($(delta)+(0,0)$) {3};
%               }
%               \graphbox{$R$}{80mm}{0mm}{50mm}{35mm}{2mm}{-5mm}{
%                   \coordinate (delta) at (-10mm,-18mm);
%                   \node[draw,circle] (r1) at ($(delta)+(-1,1.5)$) {1};
%                   \node[draw,circle] (r2) at ($(delta)+(0.5,1.5)$) {2};
%                   \node[draw,circle] (r3) at ($(delta)+(0,0)$) {3};
%                   \node[draw,circle] (r4) at ($(delta)+(1,0)$) {4};
%                   \draw[->] (r1) -- (r3) node[midway,left] {$s$};
%                   \draw[->] (r2) -- (r4) node[midway,right] {$s$};
%                   \draw[->] (r4) edge [loop below] node {0} (r4);
%                   \draw[->] (r3) edge [loop below] node {0} (r3);
%                   \node[draw,circle] (r5) at ($(r2)+(1.5,0)$) {5};
%                   \draw[->] (r5) edge [loop below] node {0} (r5);
%                   \draw[->] (r5) edge [loop right] node {0} (r5);
%                   \draw[->] (r5) edge [loop left] node {0} (r5);
%               }
%               \node () at (38mm,-18mm) {$\leftarrowtail$};
%               \node () at (77mm,-18mm) {$\rightarrowtail$};
%           \end{tikzpicture}
%           }
%         \end{center}
  
%       Consider the ruler-graph $\mathcal{X} \mathop{=} (X)$, where $X$ is the graph
%       \tikz[baseline=-0.5ex]{ 
%               \node[draw,circle] (x) at (0,0) { }; 
%               \node[draw,circle] (y) at (1,0) { };
%               \node[draw,circle] (z) at (2,0) { };
%               \draw[->] (x) -- (y) node[midway, above] {$s$};
%               \draw[->] (z) -- (y) node[midway, above] {$s$};
%       } and there is no forbidden context. $\mathcal{X}$ has weight $1$ and $\mathbb{X} \mathop{=} \{X\}$.
%       The set \( D(R,X) \) contains a unique element $R'$:
%       \raisebox{2pt}{
%           \scalebox{0.7}{\tikz[baseline=-0.5ex]{
%           \node [draw,circle] (x) at (0,0) {1};
%           \node[draw,circle] (y) at (1,0) {3};
%           \draw[->] (x) -- (y) node[midway, above] {$s$};
%       }}}.
%       \textcolor{red}{The rule is $X$-non-increasing under the function that maps the subgraph $R'$ of $R$ in $R$ to the same subgraph in $L$. }\todo{to do todo:n'importe quoi}
%     %   Since $F_X \mathop{=} \emptyset$, the number of $X$-occurrences not included in any occurrences of forbidden context in any rewriting step is predictable.
%       By Theorem~\ref{antipattern:thm:termination_grs}, the rule terminates because \(|Mono(X,L)| - |Mono(X,R)| \mathop{=} 1 - 0 \mathop{=} 1 \mathop{>} 0 \). 
% \end{example}
\begin{example}
  \label{antipattern:ex:grs_aca}
  Consider the rule $\rho$ shown below.
  \begin{center}
  \resizebox{0.8\textwidth}{!}{ 
  \begin{tikzpicture}   
      \graphbox{$L$}{0mm}{0mm}{34mm}{20mm}{2mm}{-5mm}{
          \coordinate (o) at (0mm,-3mm); 
          \node[draw,circle] (l1) at ($(o)+(-10mm,0mm)$) {1};
          \node[draw,circle] (l2) at ($(l1)+(2,0)$) {2};
          \node[draw,circle] (l3) at ($(l1)+(1,0)$) {3};
          \draw[->] (l1) -- (l3) node[midway,above] {$a$};
          \draw[->] (l3) -- (l2) node[midway,above] {$a$};
      }     
      \graphbox{$K$}{40mm}{0mm}{24mm}{20mm}{2mm}{-5mm}{
          \coordinate (o) at (5mm,-3mm); 
          \node[draw,circle] (l1) at ($(o)+(-10mm,0mm)$) {1};
          \node[draw,circle] (l2) at ($(l1)+(1,0)$) {2};
          % \node[draw,circle] (l3) at ($(l1)+(1,0)$) {$\ $};
          % \draw[->] (l1) -- (l3) node[midway,above] {$a$};
          % \draw[->] (l3) -- (l2) node[midway,above] {$a$};
      }    
      \graphbox{$R$}{70mm}{0mm}{45mm}{20mm}{2mm}{-5mm}{
        \coordinate (o) at (0mm,-3mm); 
        \node[draw,circle] (l1) at ($(o)+(-10mm,0mm)$) {1};
        \node[draw,circle] (l2) at ($(l1)+(2,0)$) {2};
        \node[draw,circle] (l3) at ($(l1)+(1,0)$) {3};
        \draw[->] (l1) -- (l3) node[midway,above] {$a$};
        \draw[->] (l3) -- (l2) node[midway,above] {$a$};
        \draw[->] (l3) edge [loop below] node {$c$} (l3);
      }    

      \node () at (37mm,-10mm) {$\leftarrowtail$};
      \node () at (67mm,-10mm) {$\rightarrowtail$};

      % \draw[>->] (51mm,2mm) -- (52mm,3mm);
  \end{tikzpicture}
  }
\end{center}
Consider the ruler-graph $\mathcal{X} \mathop{=} (L, \set{f})$ where the underlying graph is $L$ and $f:L \rightarrowtail R$ is the inclusion function. Let $s_\mathbb{X}$ be the weight function with $s_\mathbb{X}(\mathcal{X})=1$.
% The number of $L$-occurrences not included in any occurrences of the forbidden context $R$ in any rewriting step using this rule is predictable:
% the condition~\ref{hyp:inj_mono_x_l_to_r} is satisfied because $Mono(L,L)_F \mathop{=} \emptyset$; the condition~\ref{hyp:inj_mono_f_l_to_r} is satisfied because $Mono(R,L) \mathop{=} \emptyset$; the conditions~\ref{hyp:x_non_increasing} and~\ref{hyp:f_non_increasing} are satisfied under identity function; the condition~\ref{hyp:f_non_increasing} is satisfied because for the condition~\ref{hyp:f_non_increasing} we employ identity function. 

Since $D(R,X)$ is the empty set, all conditions in Definition~\ref{subgraph_counting:def:creates_more_x_on_the_left} are straightforward to verify, and rule $\rho$ is $X$-non-increasing. 
Similarly, rule $\rho^{-1}$ is $X$-non-increasing and $F$-non-increasing, because $D(L,X) \mathop{=} \emptyset$ and $D(L,R) \mathop{=} \emptyset$. 

We have 
\begin{flalign*}
\Lambda(\mathcal{X},\rho) \overset{\operatorname{def}}{\mathop{=}}&\card{\operatorname{Mono}(X,L)_{\operatorname{NF}}} - 
\card{
    \Gamma(\operatorname{Mono}(\mathcal{X},L)_{\operatorname{NF}})
    }   -
\card{\operatorname{Mono}(X,R)_{\operatorname{NF}}} \\
\mathop{=}&1 - 0 - 0 \\
\mathop{=}&1 
\end{flalign*}
Therefore, $s_\mathbb{X}(\mathcal{X}) * \Lambda(\mathcal{X},\rho) \mathop{=} 1 \mathop{>} 0$, and the rule terminates by Theorem~\ref{antipattern:thm:termination_grs}.
\end{example} 
% Finally, we present a limitation of our approach.
\begin{example}
    \label{antipattern:ex:endrullis:d3:termination}  
    Our method can prove termination of the rewriting rule $\rho$ shown below.
    %  in Figure~\ref{fig:antipattern:ex_endrullis_d3} from~\cite[Example D.3]{endrullis2024generalized_arxiv_v2}.
    \begin{center}
        \resizebox{0.8\textwidth}{!}{
    \begin{tikzpicture}  
       \graphbox{$L$}{0mm}{-11mm}{32mm}{15mm}{2mm}{-8mm}{  
           \node[draw,circle]  (x) at (-6mm,0mm) {1};  
           \node[draw,circle]  (y) at (6mm,0mm) {};  
         }  
         \graphbox{$K$}{33mm}{-11mm}{32mm}{15mm}{2mm}{-8mm}{  
           \node[draw,circle]  (x) at (-6mm,0mm) {1};  
         }  
         \graphbox{$R$}{66mm}{-11mm}{32mm}{15mm}{1mm}{-8mm}{  
          \node[draw,circle]  (x) at (-6mm,0mm) {1};  
          \node[draw,circle]  (y) at (6mm,0mm) {};  
          \draw[->]  (x) to (y);  
         }    
   \end{tikzpicture}
    }
    \end{center}
    Let $\mathcal{X} \mathop{=} (L,\set{f})$ be the ruler-graph where $f:L \rightarrowtail R$ is the unique interface-preserving monomorphism from $L$ to $R$. Let $s_\mathbb{X}$ be the weight function with $s_\mathbb{X}(\mathcal{X})=1$.
    % The number of $X$-occurrences not included in any occurrences of the forbidden context $R$ in any $\rho$-rewriting step is predictable. 
    Since $D(R,X)$ is the empty set, rule $\rho$ is $X$-non-increasing.
    Rule $\rho^{-1}$ is $X$-non-increasing and $R$-non-increasing, because $D(L,X)$ and $D(L,R)$ are empty sets.
    Since we have 
    \begin{flalign*}
      \Lambda(\mathcal{X},\rho) \overset{\operatorname{def}}{\mathop{=}}&  
      \card{\operatorname{Mono}(X,L)_{\operatorname{NF}}} - 
      \card{ \Gamma(\operatorname{Mono}(\mathcal{X},L)_{\operatorname{NF}})} -
      \card{\operatorname{Mono}(X,R)_{\operatorname{NF}}}
      \\
      \mathop{=}&1 - 0 - 0 
      \\
      \mathop{=}&1
    \end{flalign*}
    Therefore, $s_\mathbb{X}(\mathcal{X}) * \Lambda(\mathcal{X},\rho) \mathop{=} 1 \mathop{>} 0$, and the rule terminates by Theorem~\ref{antipattern:thm:termination_grs}.
  \end{example} 

% \begin{example}[Limitation]
%   The termination of the rewriting rule presented in~\cite[Ad-hoc routing protocol]{bruggink2014termination} cannot be proven by our method, because the property defined in Definition~\ref{def:predictable} is too restrictive.
% \end{example} 

\section{Empirical Results}
\label{antipattern:sec:related_work} 
% \newpage 
% ici Table~\ref{tab:antipattern:comparison:old}
% \begin{table}[H]
%    \centering
% \small % Reduce font size
% % \setlength{\tabcolsep}{4pt} % Reduce horizontal padding
%    \begin{NiceTabular}{ccccccccc}[vlines] %
%     \Hline
%     \diagbox{\enskip \textbf{Examples}}{\textbf{Techniques}} &
%     \RowStyle{\rotate}
%      \makecell{Forward closure~\cite{plump1995ontermination}} % NEW column #1
%     & \RowStyle{\rotate}
%      \makecell{Modular criterion~\cite{plump2018modular}} % NEW column #2
%     & \RowStyle{\rotate}
%      \makecell{Type graph~\cite{bruggink2014termination}}  
%     & \RowStyle{\rotate}
%      \makecell{Type graph~\cite{bruggink2015proving}} 
%     & \RowStyle{\rotate}
%      \makecell{Type graph~\cite{endrullis2024generalized_arxiv_v2}} 
%     & \RowStyle{\rotate}
%      \makecell{Subgraph counting~\cite{overbeek2024termination_lmcs}} 
%     & \RowStyle{\rotate}
%      \makecell{Subgraph counting(Chapter~\ref{chap:subgraph_counting})}
%      & \RowStyle{\rotate}
%      \makecell{Subgraph counting(Chapter~\ref{chap:antipattern})} \\
%     \Hline
%     \Hline
%     % not ok examples
%    ~\cite{plump1995ontermination}~\hyperref[ex:plump95_4d1]{Example 4.1} & \ding{51} & -- & -- & -- & -- & 
%               \ding{51}
              
%               & \ding{55}& \ding{55}\\  
%    \Hline
%   ~\cite{plump2018modular}~\hyperref[ex:plump_ex4]{Example 4} &  -- &  \ding{51} &  -- & -- & -- & 
%                --
%                & \ding{55}& \ding{55}\\ 
%    \Hline

%   ~\cite{bruggink2014termination} Routing Protocol
%        & -- & -- & \ding{51} & -- & -- & 
%            --
%           & \ding{55} &  \ding{55}\\ \Hline

%   ~\cite{bruggink2014termination}\hyperref[ex:plump_ex4]{Example 5}
%    & -- & -- & \ding{51} & -- & -- & -- 
%    & \ding{55} &  \ding{55}\\ 
% \Hline

%   ~\cite{bruggink2015proving}~\hyperref[ex:bruggink2015_ex5]{Example 5}
%    & -- & -- & -- & \ding{51} & -- &  
%    --
%    & \ding{55} &  \ding{55}\\
%    \Hline

%   ~\cite{bruggink2015proving}~\hyperref[ex:bruggink2015_ex6_endrullis2024_d2]{Example 6} 
%    & -- & -- & -- & \ding{51} & -- &  
%    --
%    & \ding{55} &  \ding{55}\\ 
%    \Hline

%   ~\cite{endrullis2024generalized_arxiv_v2}~\hyperref[ex:bruggink2015_ex6_endrullis2024_d2]{Example D.2} 
%    & -- & -- & -- & -- & \ding{51} & -- 
%    & \ding{55} & \ding{55}\\ 
%    \Hline

%       % not supported examples  
%      ~\cite{plump2018modular}~\hyperref[ex:plump2018_ex6_endrullis_d4]{Example 6} &  -- & \ding{51} & -- & -- & -- & 
%       --
%           & -- & -- \\
%       \Hline

%      ~\cite{endrullis2024generalized_arxiv_v2} Example 6.4  
%       & -- & -- & -- & -- & \ding{51} & -- 
%        & -- & -- \\ \Hline

%  ~\cite{endrullis2024generalized_arxiv_v2} Example 6.5  
%       & -- & -- & -- & -- &  \ding{51} & --
%        & -- & -- \\ \Hline

%      ~\cite{endrullis2024generalized_arxiv_v2}~\hyperref[ex:plump2018_ex6_endrullis_d4]{Example D.4} 
%       & -- & -- & -- & -- & \ding{51} & -- 
%        & -- & --\\ \Hline

%    % ----- from overbeek2024termination_lmcs -----
%  ~\cite{overbeek2024termination_lmcs} Example 5.2
%       & -- & -- & -- & -- & -- & \ding{51} 
%        & -- & -- \\ \Hline

%      ~\cite{overbeek2024termination_lmcs} Example 5.7 
%       & -- & -- & -- & -- & -- & \ding{51} 
%        & -- & -- \\ \Hline
      
%  ~\cite{overbeek2024termination_lmcs} Example 5.9 
%       & -- & -- & -- & -- & -- & \ding{51} 
%        & -- & --\\ \Hline
 
%       % ok examples

%    ~\cite{plump1995ontermination}~\hyperref[ex:overbeek_5d8_plump1995_3d8_plump2018_3_overbeek_5d8]{Example 3.8}
%                   & \ding{51} & -- & -- & -- & -- &
%                --
%                  & \ding{51}  & \ding{51}\\ 
%      \hline
     
%     ~\cite{plump2018modular}~\hyperref[ex:overbeek_5d8_plump1995_3d8_plump2018_3_overbeek_5d8]{Example 3} 
%                & -- & \ding{51} &  -- & -- & -- & 
%                --
%                 & \ding{51} & \ding{51}\\ 

%     \Hline
%    ~\cite[Example 5]{plump2018modular}  &  -- &  \ding{51} &   -- & -- & -- &  
%                  --
%                & \ding{51} & \ding{51}\\ 
%     \Hline

%   ~\cite[Example 1]{bruggink2014termination} 
%        & -- & -- & \ding{51} & -- & -- & 
%                      \ding{55}
%                   & \ding{55}  &  \ding{51}\\ 
%    \Hline

%   ~\cite{bruggink2014termination}~\hyperref[ex:termination:grsaa]{Example 4 and 6}  
%        & -- & -- & \ding{51} & -- & -- & 
%                  --
%                & \ding{51} & \ding{51} \\ \Hline

%   ~\cite{bruggink2015proving}~\hyperref[ex:termination:grsaa]{Example 2}  
%        & -- & -- & -- & \ding{51} & -- & 
%        -- 
%         & \ding{51} &  \ding{51}\\ \Hline
       
%   ~\cite{bruggink2015proving} Example 4 
%        & -- & -- & -- & \ding{51} & -- & 
%        --
%         & \ding{51} & \ding{51} \\ \Hline


%     % ----- from endrullis2024generalized_arxiv_v2 -----
%    ~\cite{endrullis2024generalized_arxiv_v2}~\hyperref[ex:endrullis2024_6d2]{Example 6.2}  
%        & -- & -- & -- & -- & \ding{51} & -- & \ding{51}& \ding{51}\\ \Hline

%   ~\cite{endrullis2024generalized_arxiv_v2}~\hyperref[ex_endrullis_6d3_endrullis_5d8]{Example 6.3}
%        & -- & -- & -- & -- & \ding{51} &% \ding{55} 
%        \ding{55} 
%         & \ding{51} & \ding{51}\\ \Hline

%   ~\cite{endrullis2024generalized_arxiv_v2}\hyperref[ex:overbeek_5d8_plump1995_3d8_plump2018_3_overbeek_5d8]{Example D.1}
%        & -- & -- & -- & -- & \ding{51} & --
%         & \ding{51} & \ding{51}\\ \Hline

%   ~\cite{endrullis2024generalized_arxiv_v2}\hyperref[rem:d3_limitation]{Example D.3}
%        & -- & -- & -- & -- & \ding{51} & \ding{55} 
%         & \ding{55} & \ding{51}\\ \Hline

%   ~\cite{overbeek2024termination_lmcs}~\hyperref[ex:overbeek_5d3]{Example 5.3}
%        & -- & -- & -- & -- & -- & \ding{51} 
%         & \ding{51} & \ding{51}\\ \Hline

%    %~\cite{overbeek2024termination_lmcs}~\hyperref[ex:overbeek_5d3]{Example 5.3 monic matches}
%    %     & -- & -- & -- & -- & -- & \ding{51} & \ding{51}\\ \Hline
%  ~\cite{overbeek2024termination_lmcs}~\hyperref[ex:overbeek_5d5]{Example 5.5} 
%        & -- & -- & -- & -- & -- & \ding{51} 
%         & \ding{51} & \ding{51}\\ \Hline

%   ~\cite{overbeek2024termination_lmcs}~\hyperref[ex:overbeek_5d6]{Example 5.6}
%        & -- & -- & -- & -- & -- & \ding{51} 
%         & \ding{51} & \ding{51} \\ \Hline

%    %~\cite{overbeek2024termination_lmcs}~\hyperref[ex:overbeek_5d6]{Example 5.6 bis}
%    %     & -- & -- & -- & -- & -- & \ding{51} & -- \\ \Hline


%   ~\cite{overbeek2024termination_lmcs}~\hyperref[ex:overbeek_5d8_plump1995_3d8_plump2018_3_overbeek_5d8]{Example 5.8}
%        & -- & -- & -- & -- & -- & \ding{51} 
%         & \ding{51} & \ding{51}\\ \Hline

% Example~\ref{ex:overbeek_5d6}
%        & \ding{51} & \ding{55} & \ding{55} & \ding{55} & \ding{55} & \ding{55} 
%         & \ding{51} & \ding{51} \\ \Hline
%   \end{NiceTabular}
%   \caption{Applicability of termination techniques to DPO rewriting examples.
%    The symbol \ding{51} indicates termination can be proved by the technique,
%    \ding{55} indicates it cannot be proved, and 
%    $-$ denotes irrelevance or out-of-scope cases.
%          }
%   \label{tab:antipattern:comparison:old}
%   \end{table} 
% \newpage
A comparative analysis of termination techniques for DPO graph rewriting systems, drawn from prior work~\cite{plump1995ontermination,plump2018modular,bruggink2014termination,bruggink2015proving,endrullis2024generalized_arxiv_v2,
overbeek2024termination_lmcs} and Chapter~\ref{chap:subgraph_counting}, is summarized in Table~\ref{tab:comparison:antipattern}. 
Our tool successfully proves termination for 16 systems, including~\cite[Example D.3]{endrullis2024generalized_arxiv_v2} and~\cite[Example 1]{bruggink2014termination}. These two cases could not be handled by the subgraph counting methods introduced in~\cite{overbeek2024termination_lmcs} and Chapter~\ref{chap:subgraph_counting}, but are now provable using our extension.
{% local group
  \setlength{\tabcolsep}{3pt}
  \renewcommand{\arraystretch}{1}
\begin{table}[H]
   \centering
  \small % Reduce font size
  \caption{Applicability of termination techniques to DPO rewriting examples.
  The symbol \ding{51} indicates termination can be proved by the technique,
  \ding{55} indicates it cannot be proved, and 
  $-$ denotes irrelevance or out-of-scope cases.
        }
 \label{tab:comparison:antipattern}
% \setlength{\tabcolsep}{4pt} % Reduce horizontal padding
   \begin{NiceTabular}{cccccccccc}[vlines] % <-- 9 columns now (was 7)
    \Hline
                % \diagbox{\enskip \textbf{Examples}}{\textbf{Techniques}} 
    \Block{1-2}{\diagbox{\enskip \textbf{Examples}}{\textbf{Techniques}}} & 
    &
    \RowStyle{\rotate}
     \makecell{Forward closure~\cite{plump1995ontermination}} % NEW column #1
    & \RowStyle{\rotate}
     \makecell{Modular criterion~\cite{plump2018modular}} % NEW column #2
    & \RowStyle{\rotate}
     \makecell{Type graph~\cite{bruggink2014termination}}  
    & \RowStyle{\rotate}
     \makecell{Type graph~\cite{bruggink2015proving}} 
    & \RowStyle{\rotate}
     \makecell{Type graph~\cite{endrullis2024generalized_arxiv_v2}} 
    & \RowStyle{\rotate}
     \makecell{Subgraph counting~\cite{overbeek2024termination_lmcs}} 
    & \RowStyle{\rotate}
     \makecell{Morphism Counting \\
       (Chapter~\ref{chap:subgraph_counting})} 
    & \RowStyle{\rotate}
     \makecell{Morphism Counting \\
     (Chapter~\ref{chap:antipattern})} 
    \\
    \Hline
    \Hline
% ok examples
Chapter~\ref{chap:subgraph_counting} &Example~\ref{subgraph_counting:ex_contrib_variant}
  & \ding{51} & \ding{55} & \ding{55} & \ding{55} & \ding{55} & \ding{55} & \ding{51} & \ding{51}\\ \Hline
  
\cite{plump1995ontermination} &
\hyperref[ex:overbeek_5d8_plump1995_3d8_plump2018_3_overbeek_5d8]{Example 3.8}
             & \ding{51} & -- & -- & -- & -- & --  & \ding{51} & \ding{51}\\ 
\hline

\Block{2-1}{\cite{plump2018modular}} & \hyperref[ex:overbeek_5d8_plump1995_3d8_plump2018_3_overbeek_5d8]{Example 3} 
          & -- & \ding{51} &  -- & -- & -- &   -- & \ding{51} & \ding{51} \\ 

\Hline

%~\cite{plump2018modular} 
& Example 5 &  -- &  \ding{51} &   -- & -- & -- &  
            --
          & \ding{51} & \ding{51} \\  
\Hline
    \Block{2-1}{\cite{bruggink2014termination}}& Example 1 
   & -- & -- & \ding{51} & -- & -- & 
                 --
               &  \ding{55} & \ding{51}\\ 
   \Hline
  & \hyperref[ex:termination:grsaa]{Example 4 and 6}  
  & -- & -- & \ding{51} & -- & -- &  --  & \ding{51}& \ding{51}  \\ \Hline

\Block{2-1}{\cite{bruggink2015proving}} & \hyperref[ex:termination:grsaa]{Example 2}  
  & -- & -- & -- & \ding{51} & -- &  -- &  \ding{51} & \ding{51} \\ \Hline
  
%~\cite{bruggink2015proving}
 & Example 4 
  & -- & -- & -- & \ding{51} & -- & 
  --& \ding{51}& \ding{51}  \\ \Hline


% ----- from endrullis2024generalized_arxiv_v2 -----
\Block{4-1}{\cite{endrullis2024generalized_arxiv_v2}} & \hyperref[ex:endrullis2024_6d2]{Example 6.2}  
  & -- & -- & -- & -- & \ding{51} & -- & \ding{51} & \ding{51} \\ \Hline

%~\cite{endrullis2024generalized_arxiv_v2}
 & \hyperref[ex_endrullis_6d3_endrullis_5d8]{Example 6.3}
  & -- & -- & -- & -- & \ding{51} &% \ding{55} 
  \ding{55} & \ding{51} & \ding{51} \\ \Hline

%~\cite{endrullis2024generalized_arxiv_v2} 
& \hyperref[ex:overbeek_5d8_plump1995_3d8_plump2018_3_overbeek_5d8]{Example D.1}
  & -- & -- & -- & -- & \ding{51} & -- & \ding{51} & \ding{51} 
  \\ \Hline
     & \hyperref[rem:d3_limitation]{Example D.3}
   & -- & -- & -- & -- & \ding{51} & \ding{55} & \ding{55} & \ding{51}
   \\ \Hline
  \Block{4-1}{\cite{overbeek2024termination_lmcs}} & \hyperref[ex:overbeek_5d3]{Example 5.3}
  & -- & -- & -- & -- & -- & \ding{51} & \ding{51} & \ding{51} \\ \Hline

%~\cite{overbeek2024termination_lmcs} \hyperref[ex:overbeek_5d3]{Example 5.3 monic matches}
%     & -- & -- & -- & -- & -- & \ding{51} & \ding{51}\\ \Hline
%~\cite{overbeek2024termination_lmcs} 
& \hyperref[ex:overbeek_5d5]{Example 5.5} 
  & -- & -- & -- & -- & -- & \ding{51} & \ding{51} & \ding{51} \\ \Hline

%~\cite{overbeek2024termination_lmcs} 
& \hyperref[ex:overbeek_5d6]{Example 5.6}
  & -- & -- & -- & -- & -- & \ding{51} & \ding{51} & \ding{51} \\ \Hline

%~\cite{overbeek2024termination_lmcs} \hyperref[ex:overbeek_5d6]{Example 5.6 bis}
%     & -- & -- & -- & -- & -- & \ding{51} & -- \\ \Hline


%~\cite{overbeek2024termination_lmcs}  
& \hyperref[ex:overbeek_5d8_plump1995_3d8_plump2018_3_overbeek_5d8]{Example 5.8}
  & -- & -- & -- & -- & -- & \ding{51} & \ding{51}& \ding{51} \\ \Hline

      % not supported examples  
     ~\cite{plump2018modular} &  \hyperref[ex:plump2018_ex6_endrullis_d4]{Example 6} &  -- & \ding{51} & -- & -- & -- & -- & -- \\
      \Hline

\Block{3-1}{\cite{endrullis2024generalized_arxiv_v2}}
 & Example 6.4  
      & -- & -- & -- & -- & \ding{51} & -- & -- \\ \Hline

%  ~\cite{endrullis2024generalized_arxiv_v2}
  &  Example 6.5  
      & -- & -- & -- & -- &  \ding{51} & -- & -- \\ \Hline

    %  ~\cite{endrullis2024generalized_arxiv_v2}
       &\hyperref[ex:plump2018_ex6_endrullis_d4]{Example D.4} 
      & -- & -- & -- & -- & \ding{51} & -- & --\\ \Hline

   % ----- from overbeek2024termination_lmcs -----
   \Block{3-1}{\cite{overbeek2024termination_lmcs}} & Example 5.2
      & -- & -- & -- & -- & -- & \ding{51} & -- \\ \Hline

    %  ~\cite{overbeek2024termination_lmcs} 
      & Example 5.7 
      & -- & -- & -- & -- & -- & \ding{51} & -- \\ \Hline
      
%  ~\cite{overbeek2024termination_lmcs} 
  & Example 5.9 
      & -- & -- & -- & -- & -- & \ding{51} & --\\ \Hline
 

    % not ok examples
   ~\cite{plump1995ontermination} & \hyperref[ex:plump95_4d1]{Example 4.1} & \ding{51} & -- & -- & -- & -- & 
              \ding{51}
              
              & \ding{55}& \ding{55}\\ 
   \Hline
  ~\cite{plump2018modular} & \hyperref[ex:plump_ex4]{Example 4} &  -- &  \ding{51} &  -- & -- & -- & 
               --
               & \ding{55}& \ding{55}\\ 
   \Hline

%   ~\cite{bruggink2014termination} 
  \Block{2-1}{\cite{bruggink2014termination}}  & Routing Protocol
       & -- & -- & \ding{51} & -- & -- & 
           --
           &  \ding{55}& \ding{55}\\  
           \Hline
% \Block{2-1}{\cite{bruggink2014termination}}
 & \hyperref[ex:plump_ex4]{Example 5}
   & -- & -- & \ding{51} & -- & -- & -- &  \ding{55}& \ding{55}\\ 
\Hline

\Block{2-1}{\cite{bruggink2015proving}} & \hyperref[ex:bruggink2015_ex5]{Example 5}
   & -- & -- & -- & \ding{51} & -- &  
   -- &  \ding{55}& \ding{55}\\
   \Hline

%   ~\cite{bruggink2015proving} 
   & \hyperref[ex:bruggink2015_ex6_endrullis2024_d2]{Example 6} 
   & -- & -- & -- & \ding{51} & -- &  
   --&  \ding{55}& \ding{55}\\  
   \Hline

   \Block{1-1}{\cite{endrullis2024generalized_arxiv_v2}} & \hyperref[ex:bruggink2015_ex6_endrullis2024_d2]{Example D.2} 
   & -- & -- & -- & -- & \ding{51} & -- & \ding{55}& \ding{55}\\ 
   \Hline
  \end{NiceTabular}
  % \vspace{2mm}
  % \caption{Applicability of termination techniques to DPO rewriting examples.
  %  The symbol \ding{51} indicates termination can be proved by the technique,
  %  \ding{55} indicates it cannot be proved, and 
  %  $-$ denotes irrelevance or out-of-scope cases.
  %        }
  % \label{tab:comparison}
  \end{table}
}

\section{Conclusion and future work} 
\label{antipattern:sec:conclusion}
We have extended our subgraph-counting method to handle cases such as 
Examples~\ref{antipattern:ex:grs_aca} and~\ref{antipattern:ex:endrullis:d3:termination} where both the method of Overbeek and Endrullis~\cite{overbeek2024termination_lmcs} and the morphism counting method of Chapter~\ref{chap:subgraph_counting} fail. These failures arise when there exists an injection from a rule's left-hand side graph into its right-hand side graph. The extended method has been implemented in the tool described in Chapter~\ref{chap:lyonparallel}.
% Interesting directions for future work are to generalize our subgraph counting method to count subgraphs with multiple forbidden contexts, to understand the difference between the subgraph counting method presented in this chapter and that in~\cite{overbeek2024termination_lmcs} to provide a more general termination criterion that leverages the strengths of both methods for DPO rewriting systems, and to extend it to DPO systems with negative application conditions. 

\section{Appendix: Proofs}
% \subsection{}
% \label{sec:new_contrib_proof_antipattern}
% \begin{lemma}
%     \label{antipattern:lem:leq_gamma_mono_x_l}
%     \ \newline
%     \noindent
%     \begin{minipage}{0.7\textwidth}\setlength{\parindent}{1em}
%     Let $\mathcal{X} \mathop{=} (X, f:X \rightarrowtail F)$ be a ruler-graph and \( \rho \mathop{=} (L \overset{l}{\leftarrowtail} K \overset{r}{\rightarrowtail} R) \) be a rule. Consider the DPO diagram illustrated on the right. 
% \end{minipage}%
% \begin{minipage}{0.29\textwidth}
%     \hfill
%     \begin{tikzpicture}
%         % [node distance=15mm]
%         \node (I) {$K$};
%         \node (L) [left of=I] {$L$};
%         \node (R) [right of=I] {$R$};
%         \node (G) [below of=L] {$G$};
%         \node (C) [below of=I] {$C$};
%         \node (H) [below of=R] {$H$}; 
%         \draw [>->] (I) to  node [midway,below] {$l$} (L);
%         \draw [>->] (I) to  node [midway,below] {$r$} (R);
%         \draw [>->] (L) to node [midway,right] {$m$} (G);
%         \draw [>->] (I) to  node [midway,right] {$u$} (C); 
%         \draw [>->] (R) to  node [midway,right] {$m'$} (H);
%         \draw [>->] (C) to node [midway,above] {$l'$} (G);
%         \draw [>->] (C) to node [midway,above] {$r'$} (H);
%         % \node [at=($(I)!.5!(G)$)] {\normalfont PO};
%         % \node [at=($(I)!.5!(H)$)] {\normalfont PO};
%       \end{tikzpicture}
% \end{minipage} 

% We have
%         $  \{
%              h_{XL} \mathop{\in} \operatorname{Mono}(\mathcal{X},L)_{\operatorname{NF}} \mathop{\mid} 
%              \exists h_{FG}. h_{XL} \mathop{\star} m \mathop{=} f \mathop{\star} h_{FG}
%          \}
%          \mathop{\subseteq}
%             \Gamma(\operatorname{Mono}(\mathcal{X},L)_{\operatorname{NF}})
%             $ 
% \end{lemma}
% \begin{proof}


% \todo{useful?}

%     By Definition~\ref{antipattern:def:gamma_l_rho_x} of $\Gamma$.
% \end{proof}

\begin{lemma}
    \label{antipattern:lem:xgm_xhmp_xl_xr_llsss}
    Let $\mathcal{X} \mathop{=} (X, \set{f:X \rightarrowtail F})$ be a ruler-graph and \( \rho \mathop{=} (L \overset{l}{\leftarrowtail} K \overset{r}{\rightarrowtail} R) \) be a rule. Consider the DPO diagram shown below.
        \begin{center}
            \begin{tikzpicture}
                % [node distance=15mm]
                \node (I) {$K$};
                \node (L) [left of=I] {$L$};
                \node (R) [right of=I] {$R$};
                \node (G) [below of=L] {$G$};
                \node (C) [below of=I] {$C$};
                \node (H) [below of=R] {$H$}; 
                \draw [>->] (I) to  node [midway,below] {$l$} (L);
                \draw [>->] (I) to  node [midway,below] {$r$} (R);
                \draw [>->] (L) to node [midway,right] {$m$} (G);
                \draw [>->] (I) to  node [midway,right] {$u$} (C); 
                \draw [>->] (R) to  node [midway,right] {$m'$} (H);
                \draw [>->] (C) to node [midway,above] {$l'$} (G);
                \draw [>->] (C) to node [midway,above] {$r'$} (H);
                % \node [at=($(I)!.5!(G)$)] {\normalfont PO};
                % \node [at=($(I)!.5!(H)$)] {\normalfont PO};
            \end{tikzpicture}
        \end{center}

We hav the following equalities:
\begin{flalign*}
    \card{\operatorname{Mono}(\mathcal{X},G,m)_{\operatorname{NF}}} \mathop{=} 
    &\card{\operatorname{Mono}(\mathcal{X},L)_{\operatorname{NF}}} - 
    \\
    &\card{\{
                h_{XL} \mathop{\in} \operatorname{Mono}(\mathcal{X},L)_{\operatorname{NF}} \mathop{\mid} 
                \exists h_{FG}. h_{XL} \mathop{\star} m \mathop{=} f \mathop{\star} h_{FG}
            \}},
    \\
    \card{\operatorname{Mono}(\mathcal{X},H,m')_{\operatorname{NF}}} \mathop{=} 
    &\card{\operatorname{Mono}(\mathcal{X},R)_{\operatorname{NF}}} - 
    \\
    &\card{\{
                h_{XR} \mathop{\in} \operatorname{Mono}(\mathcal{X},R)_{\operatorname{NF}} \mathop{\mid} 
                \exists h_{FH}. h_{XR} \mathop{\star} m' \mathop{=} f \mathop{\star} h_{FH} 
                \}}.
\end{flalign*}
\end{lemma}
\begin{proof}
    By symmetry, it suffices to prove the first equality and the second equality follows. To prove the first equality, we show that there is a bijection between $\operatorname{Mono}(\mathcal{X},L)_{\operatorname{NF}} \mathop{\setminus} \{
                h_{XL} \mathop{\in} \operatorname{Mono}(\mathcal{X},L)_{\operatorname{NF}} \mathop{\mid} 
                \exists h_{FG}. h_{XL} \mathop{\star} m \mathop{=} f \mathop{\star} h_{FG}
            \}$ and $\operatorname{Mono}(\mathcal{X},G,m)_{\operatorname{NF}}$.

    Let $h_{XG} \mathop{\in} \operatorname{Mono}(\mathcal{X},G,m)_{\operatorname{NF}}$. By definition, there is a monomorphism $h_{XL}$ from $X$ to $L$ such that $h_{XG} \mathop{=} h_{XL} \mathop{\star} m$ and there is no monomorphism $h_{FG}$ from $F$ to $G$ such that $f \mathop{\star} h_{FG} \mathop{=} h_{XG}$. We must show $h_{XL} \mathop{\in} \operatorname{Mono}(\mathcal{X},L)_{\operatorname{NF}}$ and there is no $h_{FG}:f \rightarrowtail G$ such that $h_{XL} \mathop{\star} m \mathop{=} f \mathop{\star} h_{FG}$. These can be visualized by the following commutative diagram where edges that cannot exist are in red and dashed:

    \begin{center}
        \begin{tikzpicture}
            \node (k) at (0,1) {K};
            \node (l) at (-2,1) {L};
            \node (c) at (0,-1) {C};
            \node (g) at (-2,-1) {G};
            \node (x) at (-4,-3) {X};
            \node (f) at (-2,-3) {F};
            \draw[<-<]  (l) -- (k) node [midway,below] {$l$};
            \draw[>->] (c) -- (g) node [midway, above] {$l'$};
            \draw[>->] (l) -- (g) node[midway, right] {$m$};
            \draw[>->] (k) -- (c) node[midway, left] {};
            \draw[>->] (x) -- (g) node[midway,below] { };
            \draw[red,dashed,>->] (f) -- (g) node[midway] {$\times$};
            \draw[>->] (x) -- (l) node[midway,left] { };
            \draw[>->] (x) -- (f) node[midway,below] {$f$};
            \node () [at=($(l)!0.5!(c)$)] {$PO$};
        \end{tikzpicture}
    \end{center}
    
    We show $h_{XL} \mathop{\in} \operatorname{Mono}(\mathcal{X},L)_{\operatorname{NF}}$ by contradiction. Suppose that the contrary holds. There is a morphism $h_{FL}: F \rightarrowtail L$ such that $f \mathop{\star} h_{FL} \mathop{=} h_{XL}$.
    The image of $h_{XG}$ is, therefore, included in the image of $h_{FL} \mathop{\star} m$, as shown in the commutative diagram below, which contradicts the assumption that $h_{XG}$ is in $\operatorname{Mono}(\mathcal{X},G,m)_{\operatorname{NF}}$. 

    \begin{center}
        \begin{tikzpicture}
            \node (k) at (0,1) {K};
            \node (l) at (-2,1) {L};
            \node (c) at (0,-1) {C};
            \node (g) at (-2,-1) {G};
            \node (x) at (-4,-3) {X};
            \node (f) at (-4,-1) {F};
            \draw[<-<]  (l) -- (k) node [midway,below] {$l$};
            \draw[>->] (c) -- (g) node [midway, above] {$l'$};
            \draw[>->] (l) -- (g) node[midway, right] {$m$};
            \draw[>->] (k) -- (c) node[midway, left] {};
            \draw[>->] (x) -- (g) node[midway,below] { };
            \draw[>->] (f) -- (l) node[midway] {};
            \draw[>->] (x) -- (l) node[midway,left] { };
            \draw[>->] (x) -- (f) node[midway,left] {$f$};
            \node () [at=($(l)!0.5!(c)$)] {$PO$};
        \end{tikzpicture}
    \end{center}

    
    We show that there is no $h_{FG}:f \rightarrowtail G$ such that $h_{XL} \mathop{\star} m \mathop{=} f \mathop{\star} h_{FG}$ by contradiction. Suppose that the contrary holds. Under this assumption, we have the following commutative diagram. The image of $h_{XG}$ is, thus, included in the image of $f \mathop{\star} h_{FG}$ for some morphism $f$ which contradicts the assumption that $h_{XG}$ is in $\operatorname{Mono}(\mathcal{X},G,m)_{\operatorname{NF}}$.

    \begin{center}
        \begin{tikzpicture}
            \node (k) at (0,1) {K};
            \node (l) at (-2,1) {L};
            \node (c) at (0,-1) {C};
            \node (g) at (-2,-1) {G};
            \node (x) at (-4,-3) {X};
            \node (f) at (-2,-3) {F};
            \draw[<-<]  (l) -- (k) node [midway,below] {$l$};
            \draw[>->] (c) -- (g) node [midway, above] {$l'$};
            \draw[>->] (l) -- (g) node[midway, right] {$m$};
            \draw[>->] (k) -- (c) node[midway, left] {};
            \draw[>->] (x) -- (g) node[midway,below] { };
            \draw[>->] (f) -- (g) node[midway,right] {};
            \draw[>->] (x) -- (l) node[midway,left] { };
            \draw[>->] (x) -- (f) node[midway,below] {$f$};
            \node () [at=($(l)!0.5!(c)$)] {$PO$};
        \end{tikzpicture}
    \end{center}

    Let $h_{XG}, h_{XG}' \mathop{\in} \operatorname{Mono}(\mathcal{X},G,m)_{\operatorname{NF}}$ such that $h_{XG} \mathop{\neq} h_{XG}'$. By definition, there are morphisms $h_{XL}, h_{XL}': X \rightarrowtail L$ such that $h_{XG} \mathop{=} h_{XL} \mathop{\star} m$, $h_{XG}' \mathop{=} h_{XL}' \mathop{\star} m$, and no morphism $h_{FG}, h_{FG}': F \rightarrowtail G$ such that $f \mathop{\star} h_{FG} \mathop{=} h_{XG}$ and $f \mathop{\star} h_{FG}' \mathop{=} h_{XG}'$. Since $m$ is injection, we have $h_{XL} \mathop{\neq} h_{XL}'$.

    There is, therefore, an injection from the set $\operatorname{Mono}(\mathcal{X},G,m)_{\operatorname{NF}}$ to the set $\operatorname{Mono}(\mathcal{X},L)_{\operatorname{NF}} \mathop{\setminus} \{
                h_{XL} \mathop{\in} \operatorname{Mono}(\mathcal{X},L)_{\operatorname{NF}} \mathop{\mid} 
                \exists h_{FG}. h_{XL} \mathop{\star} m \mathop{=} f \mathop{\star} h_{FG}
            \}$.

    Conversely, let $h_{XL}$ be a morphism of the following set $$\operatorname{Mono}(\mathcal{X},L)_{\operatorname{NF}}\mathop{\setminus} \{
                h_{XL} \mathop{\in} \operatorname{Mono}(\mathcal{X},L)_{\operatorname{NF}} \mathop{\mid} 
                \exists h_{FG}. h_{XL} \mathop{\star} m \mathop{=} f \mathop{\star} h_{FG}
            \}.$$ 
    We have the following commutative diagram:
    
    \begin{center}
        \begin{tikzpicture}
            \node (k) at (0,1) {K};
            \node (l) at (-2,1) {L};
            \node (c) at (0,-1) {C};
            \node (g) at (-2,-1) {G};
            \node (x) at (-4,-3) {X};
            \node (f) at (-4,-1) {F};
            \draw[<-<]  (l) -- (k) node [midway,below] {$l$};
            \draw[>->] (c) -- (g) node [midway, above] {$l'$};
            \draw[>->] (l) -- (g) node[midway, right] {$m$};
            \draw[>->] (k) -- (c) node[midway, left] {};
            \draw[red,dashed,>->] (f) -- (l) node[midway] {$\times$};
            \draw[>->] (x) -- (l) node[midway,left] { };
            \draw[>->] (x) -- (f) node[midway,left] {$f$};
            \draw[red,dashed,>->] (f) -- (g) node[pos=0.75] {$\times$};
            \node () [at=($(l)!0.5!(c)$)] {$PO$};
        \end{tikzpicture}
    \end{center}
    To prove $h_{XL} \mathop{\star} m \mathop{\in} \operatorname{Mono}(\mathcal{X},G,m)_{\operatorname{NF}}$, we show that $h_{XL} \mathop{\star} m$ is an element of $\operatorname{Mono}(\mathcal{X},G,m)$ and that there is no morphism $h_{FG}: F \rightarrowtail G$ such that $f \mathop{\star} h_{FG} \mathop{=} h_{XL} \mathop{\star} m$. We have $h_{XL} \mathop{\star} m \mathop{\in} \operatorname{Mono}(\mathcal{X},G,m)$, because the equality $h_{XL} \mathop{\star} m \mathop{=} f \mathop{\star} h_{FG}$ holds by definition. There is no morphism $h_{FG}: F \rightarrowtail G$ such that $f \mathop{\star} h_{FG} \mathop{=} h_{XL} \mathop{\star} m$, by definition of $h_{XL}$.

    Let $h_{XL}, h_{XL}'$ be morphisms in the following set
     $$\operatorname{Mono}(\mathcal{X},L)_{\operatorname{NF}}\mathop{\setminus} \{
                h_{XL} \mathop{\in} \operatorname{Mono}(\mathcal{X},L)_{\operatorname{NF}} \mathop{\mid} 
                \exists h_{FG}. h_{XL} \mathop{\star} m \mathop{=} f \mathop{\star} h_{FG}
    \}$$
    such that $h_{XL} \mathop{\neq} h_{XL}'$.
    Since $m$ is injective, we have $h_{XL} \mathop{\star} m \mathop{\neq} h_{XL}' \mathop{\star} m$. 
 
    There is, thus, an injection from the following set
 $$  (\operatorname{Mono}(\mathcal{X},L)_{\operatorname{NF}}\mathop{\setminus} \{
                h_{XL} \mathop{\in} \operatorname{Mono}(\mathcal{X},L)_{\operatorname{NF}} \mathop{\mid} 
                \exists h_{FG}. h_{XL} \mathop{\star} m \mathop{=} f \mathop{\star} h_{FG}
            \}) \mathop{\star} m$$
    to $\operatorname{Mono}(\mathcal{X},G,m)_{\operatorname{NF}}$.
    Since the following two sets are in bijection
    \begin{itemize}
        \item $(\operatorname{Mono}(\mathcal{X},L)_{\operatorname{NF}}\mathop{\setminus} \{
                h_{XL} \mathop{\in} \operatorname{Mono}(\mathcal{X},L)_{\operatorname{NF}} \mathop{\mid} 
                \exists h_{FG}. h_{XL} \mathop{\star} m \mathop{=} f \mathop{\star} h_{FG}
            \}) \mathop{\star} m$, 
        \item $\operatorname{Mono}(\mathcal{X},L)_{\operatorname{NF}}\mathop{\setminus} \{
                h_{XL} \mathop{\in} \operatorname{Mono}(\mathcal{X},L)_{\operatorname{NF}} \mathop{\mid} 
                \exists h_{FG}. h_{XL} \mathop{\star} m \mathop{=} f \mathop{\star} h_{FG}
            \}$,
    \end{itemize} 
    we conclude that there is an injection from the set $A$ to $\operatorname{Mono}(\mathcal{X},G,m)_{\operatorname{NF}}$, where $A \isdef \operatorname{Mono}(\mathcal{X},L)_{\operatorname{NF}}\mathop{\setminus} \{
                h_{XL} \mathop{\in} \operatorname{Mono}(\mathcal{X},L)_{\operatorname{NF}} \mathop{\mid} 
                \exists h_{FG}. h_{XL} \mathop{\star} m \mathop{=} f \mathop{\star} h_{FG}
            \}$.

    Finally, we conclude.
\end{proof}



\noindent
\begin*{\textbf{Lemma~\ref{antipattern:lem:xgm_xhmp_xl_xr}}}
%     \newline
%     \noindent
% \begin{minipage}{0.7\textwidth}\setlength{\parindent}{1em}
%     Let $\mathcal{X}=(X,\mathcal{C})$
%     % $\mathcal{X} \mathop{=} (X, f:X \rightarrowtail F)$
%      be a ruler-graph and \( \rho \mathop{=} (L \overset{l}{\leftarrowtail} K \overset{r}{\rightarrowtail} R) \) be a rule. Consider the DPO diagram illustrated on the right. 
% \end{minipage}%
% \begin{minipage}{0.29\textwidth}
%     \hfill
%     \begin{tikzpicture}
%         % [node distance=15mm]
%         \node (I) {$K$};
%         \node (L) [left of=I] {$L$};
%         \node (R) [right of=I] {$R$};
%         \node (G) [below of=L] {$G$};
%         \node (C) [below of=I] {$C$};
%         \node (H) [below of=R] {$H$}; 
%         \draw [>->] (I) to  node [midway,below] {$l$} (L);
%         \draw [>->] (I) to  node [midway,below] {$r$} (R);
%         \draw [>->] (L) to node [midway,right] {$m$} (G);
%         \draw [>->] (I) to  node [midway,right] {$u$} (C); 
%         \draw [>->] (R) to  node [midway,right] {$m'$} (H);
%         \draw [>->] (C) to node [midway,above] {$l'$} (G);
%         \draw [>->] (C) to node [midway,above] {$r'$} (H);
%         % \node [at=($(I)!.5!(G)$)] {\normalfont PO};
%         % \node [at=($(I)!.5!(H)$)] {\normalfont PO};
%       \end{tikzpicture}
% \end{minipage} 

% We have 
%     $$
%         \card{\operatorname{Mono}(\mathcal{X},G,m)_{\operatorname{NF}}} - 
%         \card{\operatorname{Mono}(\mathcal{X},H,m')_{\operatorname{NF}}} \geq
%         \Lambda(\mathcal{X}, \rho)
%     $$
%     where $\Lambda(\mathcal{X},\rho) \overset{\operatorname{def}}{=} \card{\operatorname{Mono}(\mathcal{X},L)_{\operatorname{NF}}} - 
%     \card{\Gamma(\operatorname{Mono}(\mathcal{X},L)_{\operatorname{NF}})} -
%    \card{\operatorname{Mono}(\mathcal{X},R)_{\operatorname{NF}}}$.

  Let $\mathcal{X} \mathop{=} (X, \set{f:X \rightarrowtail F})$ be a ruler-graph, and let \( \rho \mathop{=} (L \overset{l}{\leftarrowtail} K \overset{r}{\rightarrowtail} R) \) be a rule. 
  Consider the DPO diagram shown below.
    \begin{tikzpicture}
        % [node distance=15mm]
        \node (I) {$K$};
        \node (L) [left of=I] {$L$};
        \node (R) [right of=I] {$R$};
        \node (G) [below of=L] {$G$};
        \node (C) [below of=I] {$C$};
        \node (H) [below of=R] {$H$}; 
        \draw [>->] (I) to  node [midway,below] {$l$} (L);
        \draw [>->] (I) to  node [midway,below] {$r$} (R);
        \draw [>->] (L) to node [midway,right] {$m$} (G);
        \draw [>->] (I) to  node [midway,right] {$u$} (C); 
        \draw [>->] (R) to  node [midway,right] {$m'$} (H);
        \draw [>->] (C) to node [midway,above] {$l'$} (G);
        \draw [>->] (C) to node [midway,above] {$r'$} (H);
        % \node [at=($(I)!.5!(G)$)] {\normalfont PO};
        % \node [at=($(I)!.5!(H)$)] {\normalfont PO};
      \end{tikzpicture}
We have 
    $$
        \card{\operatorname{Mono}(\mathcal{X},G,m)_{\operatorname{NF}}} - 
        \card{\operatorname{Mono}(\mathcal{X},H,m')_{\operatorname{NF}}} \geq
        \Lambda(\mathcal{X}, \rho)
    $$
    where $\Lambda(\mathcal{X},\rho) \overset{\operatorname{def}}{=} \card{\operatorname{Mono}(\mathcal{X},L)_{\operatorname{NF}}} - 
    \card{\Gamma(\operatorname{Mono}(\mathcal{X},L)_{\operatorname{NF}})} -
   \card{\operatorname{Mono}(\mathcal{X},R)_{\operatorname{NF}}}$.
\end*{}

\begin{proof}
    \label{antipattern:proof:lem:xgm_xhmp_xl_xr}
    There are two cases to consider: either $\mathcal{C}=\emptyset$ or $\mathcal{C}=\set{f:X \rightarrowtail F}$. For the first case, since $\Gamma(\operatorname{Mono}(\mathcal{X},L)_{\operatorname{NF}}) \mathop{=} \emptyset$ by definition, the inequality is equivalent to the following:
    $$
        \card{\operatorname{Mono}(\mathcal{X},G,m)_{\operatorname{NF}}} - 
        \card{\operatorname{Mono}(\mathcal{X},H,m')_{\operatorname{NF}}} \geq
        \card{\operatorname{Mono}(\mathcal{X},L)_{\operatorname{NF}}} - 
        \card{\operatorname{Mono}(\mathcal{X},R)_{\operatorname{NF}}}.
    $$
    By definition, it is equivalent to the following inequality:
    \begin{flalign*}
        \card{\operatorname{Mono}(X,G,m)} - 
        \card{\operatorname{Mono}(X,H,m')} \mathop{\geq}  
        \card{\operatorname{Mono}(X,L)} - 
   \card{\operatorname{Mono}(X,R)}
    \end{flalign*}
    because the following equalities hold by Lemma~\ref{subgraph_counting:lem:decomp_w_u}:
    \begin{itemize}
        \item $\card{\operatorname{Mono}(X,G,m)} \mathop{=} \card{\operatorname{Mono}(X,L)}$,
        \item $\card{\operatorname{Mono}(X,H,m')} =
   \card{\operatorname{Mono}(X,R)}$.
    \end{itemize} 

   For the second case, by Lemma~\ref{antipattern:lem:xgm_xhmp_xl_xr_llsss}, we have the following:
    \begin{flalign*}
        &\card{\operatorname{Mono}(\mathcal{X},G,m)_{\operatorname{NF}}} - 
        \card{\operatorname{Mono}(\mathcal{X},H,m')_{\operatorname{NF}}} 
        \\
        \mathop{=} &(\card{\operatorname{Mono}(\mathcal{X},L)_{\operatorname{NF}}} - 
            \card{\{
                h_{XL} \mathop{\in} \operatorname{Mono}(\mathcal{X},L)_{\operatorname{NF}} \mathop{\mid} 
                \exists h_{FG}. h_{XL} \mathop{\star} m \mathop{=} f \mathop{\star} h_{FG}
            \}}
        ) - 
        \\
            &
           (
            \card{\operatorname{Mono}(\mathcal{X},R)_{\operatorname{NF}}}
            -  
            \card{\{
                h_{XR} \mathop{\in} \operatorname{Mono}(\mathcal{X},R)_{\operatorname{NF}} \mathop{\mid} 
                \exists h_{FH}. h_{XR} \mathop{\star} m' \mathop{=} f \mathop{\star} h_{FH}
            \}}
           )
    \end{flalign*}
    where 
    $\{
                h_{XL} \mathop{\in} \operatorname{Mono}(\mathcal{X},L)_{\operatorname{NF}} \mathop{\mid} 
                \exists h_{FG}. h_{XL} \mathop{\star} m \mathop{=} h_{XF} \mathop{\star} h_{FG}
            \}$ is the set of occurrences of $X$ in $L$ that are not included in any occurrences of its forbidden context in $L$ but included in an occurrence of its forbidden context in $G$, and, analogously, $\{
                h_{XR} \mathop{\in} \operatorname{Mono}(\mathcal{X},R)_{\operatorname{NF}} \mathop{\mid} 
                \exists h_{FH}. h_{XR} \mathop{\star} m' \mathop{=} h_{XF} \mathop{\star} h_{FH}
            \}$ is the set of $X$-occurrences in $R$ that are not included in any occurrences of its forbidden context in $R$ but included in an occurrence of its forbidden context in $H$.

    The following equalities and inequalities hold:
    \begin{flalign*}
        &\card{\operatorname{Mono}(\mathcal{X},G,m)_{\operatorname{NF}}} \mathop{-} 
        \card{\operatorname{Mono}(\mathcal{X},H,m')_{\operatorname{NF}}} 
        \\
           \mathop{=} &(\card{\operatorname{Mono}(\mathcal{X},L)_{\operatorname{NF}}} \mathop{-} 
           \card{\operatorname{Mono}(\mathcal{X},R)_{\operatorname{NF}}}
       )\mathop{+}
       \\
           &
          (
            \card{\{
               h_{XR} \mathop{\in} \operatorname{Mono}(\mathcal{X},R)_{\operatorname{NF}} \mathop{\mid} 
               \exists h_{FH}. h_{XR} \mathop{\star} m' \mathop{=} f \mathop{\star} h_{FH}
           \}}
           \mathop{-} \\&
           \card{\{
                h_{XL} \mathop{\in} \operatorname{Mono}(\mathcal{X},L)_{\operatorname{NF}} \mathop{\mid} 
                \exists h_{FG}. h_{XL} \mathop{\star} m \mathop{=} f \mathop{\star} h_{FG}
            \}}         
          ) 
        \\ 
        \mathop{\geq}& (\card{\operatorname{Mono}(\mathcal{X},L)_{\operatorname{NF}}} \mathop{-} 
        \card{\operatorname{Mono}(\mathcal{X},R)_{\operatorname{NF}}}
        ) \mathop{-}
         \\
         & \card{\{
             h_{XL} \mathop{\in} \operatorname{Mono}(\mathcal{X},L)_{\operatorname{NF}} \mathop{\mid} 
             \exists h_{FG}. h_{XL} \mathop{\star} m \mathop{=} h_{XF} \mathop{\star} h_{FG}
         \}}        
       \\ 
       \mathop{=}& (\card{\operatorname{Mono}(\mathcal{X},L)_{\operatorname{NF}}} \mathop{-} 
       \card{\operatorname{Mono}(\mathcal{X},R)_{\operatorname{NF}}}
       ) \mathop{-} 
       \\
       &
    %   (
    %     0 \mathop{-}   
       \card{
        \Gamma(\operatorname{Mono}(\mathcal{X},L)_{\operatorname{NF}})
        }      
    %   )  
        \text{\hspace{1cm}by Definition~\ref{antipattern:def:gamma_l_rho_x}}
      \\
      \mathop{=}& 
      \card{\operatorname{Mono}(\mathcal{X},L)_{\operatorname{NF}}} \mathop{-} 
      \card{
        \Gamma(\operatorname{Mono}(\mathcal{X},L)_{\operatorname{NF}})
        } \mathop{-}
      \card{\operatorname{Mono}(\mathcal{X},R)_{\operatorname{NF}}}
      \\
      \mathop{=}& \Lambda(\mathcal{X}, \rho).
    \end{flalign*}
\end{proof}

\begin{lemma}
    \label{antipattern:lem:xh_decomp}
     Consider the following pushout diagram in the category \textbf{Graph}:
     \begin{center}
        \begin{tikzpicture}
            \node (k) at (0,1) {K};
            \node (r) at (2,1) {R};
            \node (c) at (0,-1) {C};
            \node (h) at (2,-1) {H};
            \draw[>->]  (k) -- (r) node [midway,below] {$r$};
            \draw[>->] (c) -- (h) node [midway,above] {$r'$};
            \draw[>->] (r) -- (h) node[midway, left] {$m'$};
            \draw[>->] (k) -- (c) node[midway, left] {};
            \node () [at=($(r)!0.5!(c)$)] {$PO$};
        \end{tikzpicture}
    \end{center}
    Let $X$ be a graph, and $h_{XH}:X \rightarrowtail H$ a monomorphism. Then, we can construct the following commutative diagram in the category \textbf{Graph}: 

    \begin{center}
        \resizebox{0.3\textwidth}{!}{
            \begin{tikzpicture}
                \node (k) at (0,0) {K};
                \node (r) at (4,0) {R};
                \node (c) at (0,-3) {C};
                \node (h) at (4,-3) {H};
                \draw[<-<]  (r) -- (k) node [midway,above] {};
                \draw[>->] (c) -- (h) node [midway, below] {};
                \draw[>->] (r) -- (h) node[midway, left] {};
                \draw[>->] (k) -- (c) node[midway, left] {};
                \node (k') at (1,-1) {K'};
                \node (r') at (2,-1) {R'};
                \node (c') at (1,-2) {C'};
                \node () at (1.5,-1.5) {$PO$};
                \node () at (3,-1.5) {$PB$};
                \node () at (1.5,-2.5) {$PB$};
                \node () at (1.5,-0.5) {$PB$};
                \node () at (0.5,-1.5) {$PB$};
                \node (x) at (2,-2) {X};
                \draw [>->] (c') -- (x);
                \draw [>->] (r') -- (x);
                \draw [>->] (k') -- (r');
                \draw [>->] (k') -- (c');
                \draw [>->] (c') -- (c);
                \draw[>->] (r') -- (r);
                \draw[>->] (x) -- (h) node[midway,right] {};
                \draw[>->] (k') -- (k) ;
            \end{tikzpicture}
        }
        \end{center} 
\end{lemma}
\begin{proof}    
         In the category \textbf{Graph}, we can construct the following commutative diagram \todo{should I give more details?}:
        \begin{center}
            \resizebox{0.3\textwidth}{!}{
                \begin{tikzpicture}
                    \node (k) at (0,0) {K};
                    \node (r) at (4,0) {R};
                    \node (c) at (0,-3) {C};
                    \node (h) at (4,-3) {H};
                    % \node (rb) at ($\scl*(1.5,-0.5)$) {$R_X$};
                    
                    % \node (h') at ($\scl*(1.5,-1.5)$) {$H'$};
                    % \draw[>->]  (rb) -- (h') node [midway,above] {};
                    % \draw[>->]  (c) -- (h') node [midway,above] {};
        
                    \draw[<-<]  (r) -- (k) node [midway,above] {};
                    \draw[>->] (c) -- (h) node [midway, below] {};
                    \draw[>->] (r) -- (h) node[midway, left] {};
                    \draw[>->] (k) -- (c) node[midway, left] {};
        
                    % \draw[->] (rb) to (l);
                    % \draw[<-<] (rb) to (k);
                
                    \node (k') at (1,-1) {K'};
                    \node (r') at (2,-1) {R'};
                    \node (c') at (1,-2) {C'};
                    \node () at (1.5,-1.5) {$PB$};
                    \node () at (3,-1.5) {$PB$};
                    \node () at (1.5,-2.5) {$PB$};
                    \node (x) at (2,-2) {X};
                    % \draw [->] (x) -- (h') node[midway] {!};
                    \draw [>->] (c') -- (x);
                    \draw [>->] (r') -- (x);
                    \draw [>->] (k') -- (r');
                    \draw [>->] (k') -- (c');
                    \draw [>->] (c') -- (c);
                    % \draw [->] (r') -- (rb);
                    % \draw [->] (r') -- (rb);
                    % \draw [->] (h') -- (h);
        
                    \draw[>->] (r') -- (r);
                    \draw[>->] (x) -- (h) node[midway,right] {$h$};
                    % \node (rb) at ($\scl*(\sclx*1.5,-0.5)$) {$R_X$};
                    % \node (h') at ($\scl*(\sclx*1.5,-1.2)$) {$H'$};
                    % \draw[>->]  (rb) -- (h') node [midway,above] {};
                    % \draw[>->]  (c) -- (h') node [midway,above] {};
                    % \draw[>->]  (rb) -- (h') node [midway,above] {};
                    % \draw[>->] (rb) to (r);
                    % \draw[<-<] (rb) to (k);
                \end{tikzpicture}
            }
            \end{center} 
        
        Since $KCHR$ is a pushout, it is also a pullback, by Proposition~\ref{prop:pb_eq_po}. 
        Therefore, from \(
            h_{K'C'} \mathop{\star} h_{C'C} \mathop{\star} h_{CH} =h_{K'R'} \mathop{\star} h_{R'R} \mathop{\star} h_{R'R}
        \) and the universal property of pullback, there is a unique morphism $h_{K'K}:K' \rightarrowtail K$ such that 
        \begin{flalign*}
            h_{K'C'} \mathop{\star} h_{C'C} &= h_{K'K} \mathop{\star} h_{KC}, 
            % \label{antipattern:kpcpcpckpkkc}
            \\
            h_{K'R'} \mathop{\star} h_{R'R} &= h_{K'K} \mathop{\star} h_{KR}.
            % \label{antipattern:kprprprkpk}
         \end{flalign*}
        The square $K'KCC'$ and $K'R'RK$ are pullbacks, by Lemma~\ref{kpcpck_pullback}. Furthermore, $K'C'XR'$ is a pushout by Lemma~\ref{kpcpxrp_po}. Hence, the following commutative diagram holds:
        \begin{center}
            \resizebox{0.3\textwidth}{!}{
                \begin{tikzpicture}
                    \node (k) at (0,0) {K};
                    \node (r) at (4,0) {R};
                    \node (c) at (0,-3) {C};
                    \node (h) at (4,-3) {H};
                    % \node (rb) at ($\scl*(1.5,-0.5)$) {$R_X$};
                    % \node (h') at ($\scl*(1.5,-1.5)$) {$H'$};
                    % \draw[>->]  (rb) -- (h') node [midway,above] {};
                    % \draw[>->]  (c) -- (h') node [midway,above] {};
                    \draw[<-<]  (r) -- (k) node [midway,above] {};
                    \draw[>->] (c) -- (h) node [midway, below] {};
                    \draw[>->] (r) -- (h) node[midway, left] {};
                    \draw[>->] (k) -- (c) node[midway, left] {};
                    % \draw[->] (rb) to (l);
                    % \draw[<-<] (rb) to (k);
                    \node (k') at (1,-1) {K'};
                    \node (r') at (2,-1) {R'};
                    \node (c') at (1,-2) {C'};
                    \node () at (1.5,-1.5) {$PB$};
                    \node () at (3,-1.5) {$PB$};
                    \node () at (1.5,-2.5) {$PB$};
                    \node () at (1.5,-0.5) {$PB$};
                    \node () at (0.5,-1.5) {$PB$};
                    \node (x) at (2,-2) {X};
                    % \draw [->] (x) -- (h') node[midway] {!};
                    \draw [>->] (c') -- (x);
                    \draw [>->] (r') -- (x);
                    \draw [>->] (k') -- (r');
                    \draw [>->] (k') -- (c');
                    \draw [>->] (c') -- (c);
                    % \draw [->] (r') -- (rb);
                    % \draw [->] (r') -- (rb);
                    % \draw [->] (h') -- (h);
                    \draw[>->] (r') -- (r);
                    \draw[>->] (x) -- (h) node[midway,right] {$h$};
                    % \node (rb) at ($\scl*(\sclx*1.5,-0.5)$) {$R_X$};
                    % \node (h') at ($\scl*(\sclx*1.5,-1.2)$) {$H'$};
                    % \draw[>->]  (rb) -- (h') node [midway,above] {};
                    % \draw[>->]  (c) -- (h') node [midway,above] {};
                    % \draw[>->]  (rb) -- (h') node [midway,above] {};
                    % \draw[>->] (rb) to (r);
                    % \draw[<-<] (rb) to (k);
                    \draw[>->] (k') -- (k) ;
                \end{tikzpicture}
            }
            \end{center} 
        \todo{use: assumption: non increasing}
        % Since the rule is by assumption $X$-non-increasing, there is a morphism $h_{R'L}$ such that $K'KLR'$ is commutative.
\end{proof}

\noindent
\begin*{\textbf{Lemma~\ref{antipattern:lem:xglnotmlp_xhlnotmrp}}}
        Let \( \rho \mathop{=} (L \overset{l}{\leftarrowtail} K \overset{r}{\rightarrowtail} R) \) be a rule and $\mathcal{X}=(X,\mathcal{C})$ be a ruler-graph with underlying graph $X$.
        Suppose that $\rho^{-1}$ is $F$-non-increasing if $\mathcal{C} \mathop{=} \set{f:X \rightarrowtail F}$. For the DPO shown below,
        \begin{center}
        \begin{tikzpicture}
            % [node distance=15mm]
            \node (I) {$K$};
            \node (L) [left of=I] {$L$};
            \node (R) [right of=I] {$R$};
            \node (G) [below of=L] {$G$}; 
            \node (C) [below of=I] {$C$};
            \node (H) [below of=R] {$H$}; 
            \draw [>->] (I) to  node [midway,below] {$l$} (L);
            \draw [>->] (I) to  node [midway,below] {$r$} (R);
            \draw [>->] (L) to node [midway,right] {$m$} (G);
            \draw [>->] (I) to  node [midway,right] {$u$} (C); 
            \draw [>->] (R) to  node [midway,right] {$m'$} (H);
            \draw [>->] (C) to node [midway,above] {$l'$} (G);
            \draw [>->] (C) to node [midway,above] {$r'$} (H);
            % \node [at=($(I)!.5!(G)$)] {\normalfont PO};
            % \node [at=($(I)!.5!(H)$)] {\normalfont PO};
        \end{tikzpicture}
        \end{center}
        the following inequality holds:
    $$\card{\operatorname{Mono}(\mathcal{X},G,\mathop{\lnot} m, l')_{\operatorname{NF}}} \geq
        \card{\operatorname{Mono}(\mathcal{X},H,\mathop{\lnot} m', r')_{\operatorname{NF}}}.$$
\end*{}
\begin{proof}
    \label{antipattern:proof:lem:xglnotmlp_xhlnotmrp}
    There are two cases to consider: either $\mathcal{C} \mathop{=}\emptyset$ or $\mathcal{C} \mathop{=} \set{f:X \rightarrowtail F}$. 
     In the first case, since there is no forbidden context, the inequality is equivalent to $
        \card{\operatorname{Mono}(\mathcal{X},G,\lnot m, l')} \geq
        \card{\operatorname{Mono}(\mathcal{X},H,\lnot m', r')}$.
    To show this inequality, we are going to establish an injection from $\operatorname{Mono}(\mathcal{X},H,\lnot m', r')$ to $\operatorname{Mono}(\mathcal{X},G,\lnot m, l')$. Let $h_{XH} \mathop{\in} \operatorname{Mono}(\mathcal{X},H,\lnot m', r')$. By definition, there is a monomorphism $h_{XC}:X \rightarrowtail C$ such that $h_{XH} \mathop{=} h_{XC} \mathop{\star} r'$. Thus we have the following commutative diagram:

     \begin{center}
    \begin{tikzpicture}
        \node (k) at (0,1) {K};
        \node (l) at (-2,1) {L};
        \node (r) at (2,1) {R};
        \node (c) at (0,-1) {C};
        \node (g) at (-2,-1) {G};
        \node (h) at (2,-1) {H};
        \node (x) at (4,-2) {X};
        \draw[>->] (x) -- (h) node [midway,above] {};
        \draw[>->] (x) edge (c) node [midway,above] {};
        \draw[>->]  (k) -- (r) node [midway,below] {$r$};
        \draw[>->]  (k) -- (l) node [midway,below] {$l$};
        \draw[>->] (c) -- (g) node [midway, above] {$l'$};
        \draw[>->] (c) -- (h) node [midway,above] {$r'$};
        \draw[>->] (l) -- (g) node[midway, right] {$m$};
        \draw[>->] (r) -- (h) node[midway, left] {$m'$};
        \draw[>->] (k) -- (c) node[midway, left] {$u$};
        \node () [at=($(l)!0.5!(c)$)] {$PO$};
        \node () [at=($(r)!0.5!(c)$)] {$PO$};
    \end{tikzpicture}
    \end{center}
    We must show $h_{XC} \mathop{\star} l' \mathop{\in} \operatorname{Mono}(\mathcal{X},G,\lnot m, l')$ which is equivalent to showing that there is no $h_{XL}$ such that 
    $h_{XL} \mathop{\star} m \mathop{=} h_{XC} \mathop{\star} l'$.
    Suppose that the contrary holds.
     Since $KLGC$ is a pullback by Proposition~\ref{prop:pb_eq_po}, there is a morphism $h_{XK}:X \rightarrowtail K$ such that 
            \begin{flalign}
         h_{XK} \mathop{\star} u \mathop{=} h_{XC}. \label{alabeldkfsjlk}
        \end{flalign} 
    Therefore, we have 
    \begin{flalign*}
        h_{XH} 
        &= h_{XC} \mathop{\star} r' \\
        &= (h_{XK} \mathop{\star} u) \mathop{\star} r' \hspace{1cm} by Equation~\eqref{alabeldkfsjlk}\\
        &= h_{XK} \mathop{\star} (r \mathop{\star} m') \\ 
        &= (h_{XK} \mathop{\star} r) \mathop{\star} m'.
    \end{flalign*}
    This contradicts the fact that $h_{XH} \mathop{\in} \operatorname{Mono}(\mathcal{X},H,\lnot m', r')$. Therefore, we have $h_{XC} \mathop{\star} l' \mathop{\in} \operatorname{Mono}(\mathcal{X},G,\lnot m, l')$.

    Furthermore, if $h_{XH}, h_{XH}' \mathop{\in} \operatorname{Mono}(\mathcal{X},H,\lnot m', r')$ with $h_{XH} \mathop{\neq} h_{XH}'$, then there are $h_{XC}, h_{XC}'$ such that $h_{XH} \mathop{=} h_{XC} \mathop{\star} r'$ and $h_{XH}' \mathop{=} h_{XC}' \mathop{\star} r'$. Since $r'$ is injective, we deduce that $h_{XC} \mathop{\neq} h_{XC}'$. Therefore, from the injectivity of $l'$, we have $h_{XC} \mathop{\star} l' \mathop{\neq} h_{XC}' \mathop{\star} l'$. Thus, there is an injection from $\operatorname{Mono}(\mathcal{X},H,\lnot m', r')$ to $\operatorname{Mono}(\mathcal{X},G,\lnot m, l')$, which proves the inequality.
        
        Now, let us focus on the second case, where $\mathcal{X} \mathop{=} (X, f:X \rightarrowtail F)$. We have the following equalities:
        \begin{flalign*}
            &\card{\operatorname{Mono}(\mathcal{X},G,\lnot m, l')_{\operatorname{NF}}} - 
            \card{\operatorname{Mono}(,H,\lnot m', r')_{\operatorname{NF}}} 
           \\
            \mathop{=} &(\card{\operatorname{Mono}(X,G,\lnot m, l')} - \card{\operatorname{Mono}(\mathcal{X},G,\lnot m, l')_{\operatorname{F}}}) - \\
              &(\card{\operatorname{Mono}(X,H,\lnot m', r')} - \card{\operatorname{Mono}(\mathcal{X},H,\lnot m', r')_{\operatorname{F}}})
            \\
            \mathop{=} &(\card{\operatorname{Mono}(X,G,\lnot m, l')} - \card{\operatorname{Mono}(X,H,\lnot m', r')}) +\\ 
              &(\card{\operatorname{Mono}(\mathcal{X},H,\lnot m', r')_{\operatorname{F}}} - \card{\operatorname{Mono}(\mathcal{X},G,\lnot m, l')_{\operatorname{F}}})
               \\
            \mathop{=} & (\card{\operatorname{Mono}(X,C,\lnot u)} - \card{\operatorname{Mono}(X,C,\lnot u)})\mathop{+}\\
              & (\card{\operatorname{Mono}(\mathcal{X},H,\lnot m', r')_{\operatorname{F}}} - \card{\operatorname{Mono}(\mathcal{X},G,\lnot m, l')_{\operatorname{F}}})
            \\
            &\hspace{5cm}\text{by Lemma~\ref{subgraph_counting:lem:decomp_w_u}}
            \\
            \mathop{=} & \card{\operatorname{Mono}(\mathcal{X},H,\lnot m', r')_{\operatorname{F}}} - \card{\operatorname{Mono}(\mathcal{X},G,\lnot m, l')_{\operatorname{F}}}
        \end{flalign*}
    We prove this lemma by show that the function $\phi$ from $\operatorname{Mono}(\mathcal{X},G,\lnot m, l')_{\operatorname{F}}$ to $\operatorname{Mono}(\mathcal{X},H,\lnot m', r')_{\operatorname{NF}}$
    defined by $\phi(h_{XG}) \mathop{=} h_{XC} \mathop{\star} r'$ for all element $h_{XG}$ in $\operatorname{Mono}(\mathcal{X},G,\lnot m, l')_{\operatorname{F}}$, where $h_{XC}:X \rightarrowtail C$ is the uniquemonomorphism such that $h_{XG} \mathop{=} h_{XC} \mathop{\star} l'$, is injective.

    We first show that $\phi$ is well-defined.

    Let $h_{XG} \mathop{\in} \operatorname{Mono}(\mathcal{X},G,\lnot m, l')_{\operatorname{F}}$. By definition, there are $h_{FG}:F \rightarrowtail G$ and $h_{XC}:X \rightarrowtail C$ such that
        \begin{flalign}
            h_{XG} &= h_{XC} \mathop{\star} l', \label{antipattern:lem:hxghxclp} \\
            h_{XF} \mathop{\star} h_{FG} &= H_{XG}. \label{antipattern:lem:hxfhfghxg}
        \end{flalign}
    Therefore, the following commutative diagram holds:

    \begin{center}
    \begin{tikzpicture}
        \node (k) at (0,1) {K};
        \node (l) at (-2,1) {L};
        \node (r) at (2,1) {R};
        \node (c) at (0,-1) {C};
        \node (g) at (-2,-1) {G};
        \node (h) at (2,-1) {H};
        \node (x) at (-4,-1) {X};
        \node (f) at  (-4,-2) {F};
        \draw[>->] (x) -- (g) node [midway,above] {};
        \draw[>->] (x) -- (f) node [midway,above] {};
        \draw[>->]  (f) -- (g) node [midway,above] {};
        % \draw[>->] (x) -- (l) node [midway,above] {};
        \draw[>->] (x) edge [bend right](c) node [midway,above] {};
        \draw[>->]  (k) -- (r) node [midway,below] {$r$};
        \draw[>->]  (k) -- (l) node [midway,below] {$l$};
        \draw[>->] (c) -- (g) node [midway, above] {$l'$};
        \draw[>->] (c) -- (h) node [midway,above] {$r'$};
        \draw[>->] (l) -- (g) node[midway, right] {$m$};
        \draw[>->] (r) -- (h) node[midway, left] {$m'$};
        \draw[>->] (k) -- (c) node[midway, left] {$u$};
        \node () [at=($(l)!0.5!(c)$)] {$PO$};
        \node () [at=($(r)!0.5!(c)$)] {$PO$};
    \end{tikzpicture}
    \end{center}


   We prove that $\phi(h_{XG})$ cannot be factorized by $m'$, by contradiction. Suppose that the contrary holds, i.e. there is $h_{XR}:X \mathop{\rightarrow} R$ such that $\phi(h_{XG}) \overset{def}{=} h_{XC} \mathop{\star} r' \mathop{=} h_{XR} \mathop{\star} m'$. The morphism $h_{XR}$ is injective, because $\phi(h_{XG})$ is injective. Therefore, under this assumption, we have the following commutative diagram:
    \begin{center}
    \begin{tikzpicture}
       \node (k) at (0,1) {K};
       \node (l) at (-2,1) {L};
       \node (r) at (2,1) {R};
       \node (c) at (0,-1) {C};
       \node (g) at (-2,-1) {G};
       \node (h) at (2,-1) {H};
       \node (x) at (-4,-1) {X};
       \node (f) at  (-4,-2) {F};
       \draw[>->] (x) -- (g) node [midway,above] {};
       \draw[>->] (x) -- (f) node [midway,above] {};
       \draw[>->]  (f) -- (g) node [midway,above] {};
       % \draw[>->] (x) -- (l) node [midway,above] {};
       \draw[>->] (x) edge [bend right](c) node [midway,above] {};
       \draw[>->] (x) edge [out=90,in=135](r) node [midway,above] {};
       \draw[>->] (f) edge [bend right](c) node [midway,above] {};
       \draw[>->]  (k) -- (r) node [midway,below] {$r$};
       \draw[>->]  (k) -- (l) node [midway,below] {$l$};
       \draw[>->] (c) -- (g) node [midway, above] {$l'$};
       \draw[>->] (c) -- (h) node [midway,above] {$r'$};
       \draw[>->] (l) -- (g) node[midway, right] {$m$};
       \draw[>->] (r) -- (h) node[midway, left] {$m'$};
       \draw[>->] (k) -- (c) node[midway, left] {$u$};
       \node () [at=($(l)!0.5!(c)$)] {$PO$};
       \node () [at=($(r)!0.5!(c)$)] {$PO$};
   \end{tikzpicture}
    \end{center}
   
   \noindent
    The square $KCHR$ is also a pullback by Proposition~\ref{prop:pb_eq_po}. 
    By the universal property of pullback, there is a unique morphism $h_{XK} \mathop{\colon} X \mathop{\rightarrow} K$, which is injective because $h_{XR}$ is injective, such that $h_{XK} \mathop{\star} r \mathop{=} h_{XR}$ and 
    \begin{flalign}
        h_{XK} \mathop{\star} u \mathop{=} h_{XC}. \label{antipattern:eq:hxk_u}
    \end{flalign} 
    Therefore, the following commutative diagram holds:
    \begin{center}
       \begin{tikzpicture}
            \node (k) at (0,1) {K};
            \node (l) at (-2,1) {L};
            \node (r) at (2,1) {R};
            \node (c) at (0,-1) {C};
            \node (g) at (-2,-1) {G};
            \node (h) at (2,-1) {H};
            \node (x) at (-4,-1) {X};
            \node (f) at  (-4,-2) {F};
            \draw[>->] (x) -- (g) node [midway,above] {};
            \draw[>->] (x) -- (f) node [midway,above] {};
            \draw[>->]  (f) -- (g) node [midway,above] {};
            % \draw[>->] (x) -- (l) node [midway,above] {};
            \draw[>->] (x) edge [bend right](c) node [midway,above] {};
            \draw[>->] (x) edge [out=90,in=135](r) node [midway,above] {};
            \draw[>->] (x) edge [out=90,in=135](k) node [midway,above] {};
            \draw[>->] (f) edge [bend right](c) node [midway,above] {};
            \draw[>->]  (k) -- (r) node [midway,below] {$r$};
            \draw[>->]  (k) -- (l) node [midway,below] {$l$};
            \draw[>->] (c) -- (g) node [midway, above] {$l'$};
            \draw[>->] (c) -- (h) node [midway,above] {$r'$};
            \draw[>->] (l) -- (g) node[midway, right] {$m$};
            \draw[>->] (r) -- (h) node[midway, left] {$m'$};
            \draw[>->] (k) -- (c) node[midway, left] {$u$};
            \node () [at=($(l)!0.5!(c)$)] {$PO$};
            \node () [at=($(r)!0.5!(c)$)] {$PB$};
        \end{tikzpicture}
    \end{center}
   
   \noindent 
   Thus, we have \begin{flalign*}
        h_{XG} &= h_{XC} \mathop{\star} l' \\
        &= (h_{XK} \mathop{\star} u) \mathop{\star} l' &\text{by Equation~\eqref{antipattern:eq:hxk_u}}\\
        &= h_{XK} \mathop{\star} (l \mathop{\star} m).  &\text{by commutativity of $KLGC$}
    \end{flalign*}
    This is a contradiction because $h_{XG}$ is an element of $\operatorname{Mono}(\mathcal{X},G,\lnot m, l')_{\operatorname{F}}$ and, by definition, cannot be factorized by $m$. Thus, the assumption is false, and $\phi(h_{XG})$ cannot be factorized by $m'$.

    % In order to show that $\phi(h_{XG})$ is an element in $\operatorname{Mono}(\mathcal{X},H,\lnot m', r')_{\operatorname{F}}$, since $\phi(h_{XG})$ can be factorized by $r'$ by definition, 
    Consider the commutative diagram shown below.
    It remains to show that $\phi(h_{XG})$ can be decomposed as $x \mathop{\star} y$ with $x \mathop{\colon} X \rightarrowtail F$ and $y \mathop{\colon} F \rightarrowtail H$. There are three mutually exclusive cases according to whether or not the morphism $h_{FG}$ can be factorized by $m$ or $l'$.
        \begin{center}
    \begin{tikzpicture}
        \node (k) at (0,1) {K};
        \node (l) at (-2,1) {L};
        \node (r) at (2,1) {R};
        \node (c) at (0,-1) {C};
        \node (g) at (-2,-1) {G};
        \node (h) at (2,-1) {H};
        \node (x) at (-4,-1) {X};
        \node (f) at  (-4,-2) {F};
        \draw[>->] (x) -- (g) node [midway,above] {};
        \draw[>->] (x) -- (f) node [midway,above] {};
        \draw[>->]  (f) -- (g) node [midway,above] {};
        % \draw[>->] (x) -- (l) node [midway,above] {};
        \draw[>->] (x) edge [bend right](c) node [midway,above] {};
        \draw[>->]  (k) -- (r) node [midway,below] {$r$};
        \draw[>->]  (k) -- (l) node [midway,below] {$l$};
        \draw[>->] (c) -- (g) node [midway, above] {$l'$};
        \draw[>->] (c) -- (h) node [midway,above] {$r'$};
        \draw[>->] (l) -- (g) node[midway, right] {$m$};
        \draw[>->] (r) -- (h) node[midway, left] {$m'$};
        \draw[>->] (k) -- (c) node[midway, left] {$u$};
        \node () [at=($(l)!0.5!(c)$)] {$PO$};
        \node () [at=($(r)!0.5!(c)$)] {$PO$};
    \end{tikzpicture}
    \end{center}
    %  $h_{FG}$ can be factorized by $m$ or $l'$:
    \begin{itemize}
        \item Case 1 : there is $h_{FL} \mathop{\colon} F \mathop{\rightarrow} L$ such that $h_{FG} \mathop{=} h_{FL} \mathop{\star} m$. 
        This case is impossible, because if it were true, then we would have the following commutative diagram:

        \begin{tikzpicture}
            \node (k) at (0,1) {K};
            \node (l) at (-2,1) {L};
            \node (r) at (2,1) {R};
            \node (c) at (0,-1) {C};
            \node (g) at (-2,-1) {G};
            \node (h) at (2,-1) {H};
            \node (x) at (-4,-1) {X};
            \node (f) at  (-4,-2) {F};
            \draw[>->] (x) -- (g) node [midway,above] {};
            \draw[>->] (x) -- (f) node [midway,above] {};
            \draw[>->]  (f) -- (g) node [midway,above] {};
            \draw[>->] (f) -- (l) node [midway,above] {};
            \draw[>->] (x) edge [bend right](c) node [midway,above] {};
            \draw[>->]  (k) -- (r) node [midway,below] {$r$};
            \draw[>->]  (k) -- (l) node [midway,below] {$l$};
            \draw[>->] (c) -- (g) node [midway, above] {$l'$};
            \draw[>->] (c) -- (h) node [midway,above] {$r'$};
            \draw[>->] (l) -- (g) node[midway, right] {$m$};
            \draw[>->] (r) -- (h) node[midway, left] {$m'$};
            \draw[>->] (k) -- (c) node[midway, left] {$u$};
            \node () [at=($(l)!0.5!(c)$)] {$PO$};
            \node () [at=($(r)!0.5!(c)$)] {$PO$};
        \end{tikzpicture}
        
        Therefore, we would have 
        \begin{flalign*}
             h_{XC} \mathop{\star} l' \mathop{=} &h_{XG}  & \text{by Equation~\eqref{antipattern:lem:hxghxclp}}\\
             \mathop{=} &f \mathop{\star} h_{FG} 
             &\text{by Equation~\eqref{antipattern:lem:hxfhfghxg}} \\
             \mathop{=} &f \mathop{\star} h_{FL} \mathop{\star} m.
        \end{flalign*} 
        This would imply that $h_{XC}$ is not in $\operatorname{Mono}(X,G,\lnot m, l')$.
        \item Case 2 : there does not exist $h_{FL} \mathop{\colon} F \mathop{\rightarrow} L$ such that $h_{FG} \mathop{=} h_{FL} \mathop{\star} m$, but there is $h_{FC} \mathop{\colon} F \mathop{\rightarrow} C$ such that $h_{FG} \mathop{=} h_{FC} \mathop{\star} l'$. 
        We have the following commutative diagram:

        \begin{tikzpicture}
            \node (k) at (0,1) {K};
            \node (l) at (-2,1) {L};
            \node (r) at (2,1) {R};
            \node (c) at (0,-1) {C};
            \node (g) at (-2,-1) {G};
            \node (h) at (2,-1) {H};
            \node (x) at (-4,-1) {X};
            \node (f) at  (-4,-2) {F};
            \draw[>->] (x) -- (g) node [midway,above] {};
            \draw[>->] (x) -- (f) node [midway,above] {};
            \draw[>->]  (f) -- (g) node [midway,above] {};
            % \draw[>->] (x) -- (l) node [midway,above] {};
            \draw[>->] (x) edge [bend right](c) node [midway,above] {};
            \draw[>->] (f) edge [bend right](c) node [midway,above] {};
            \draw[>->]  (k) -- (r) node [midway,below] {$r$};
            \draw[>->]  (k) -- (l) node [midway,below] {$l$};
            \draw[>->] (c) -- (g) node [midway, above] {$l'$};
            \draw[>->] (c) -- (h) node [midway,above] {$r'$};
            \draw[>->] (l) -- (g) node[midway, right] {$m$};
            \draw[>->] (r) -- (h) node[midway, left] {$m'$};
            \draw[>->] (k) -- (c) node[midway, left] {};
            \node () [at=($(l)!0.5!(c)$)] {$PO$};
            \node () [at=($(r)!0.5!(c)$)] {$PO$};
        \end{tikzpicture}

        We have therefore $\phi(h_{XG}) \mathop{=} h_{XC} \mathop{\star} r' \mathop{=} h_{XF} \mathop{\star} (h_{FC} \mathop{\star} r')$.

    \item Case 3 : there does not exist $h_{FL} \mathop{\colon}  F \mathop{\rightarrow} L$ such that $h_{FG} \mathop{=} h_{FL} \mathop{\star} m$,
    and there does not exist $h_{FC} \mathop{\colon} F \mathop{\rightarrow} C$ such that $h_{FG} \mathop{=} h_{FC} \mathop{\star} l'$.
     By Lemma~\ref{antipattern:lem:xh_decomp}, we can construct the following commutative diagram:

    \begin{center}
        \resizebox{0.3\textwidth}{!}{
            \begin{tikzpicture}
                \node (k) at (0,0) {K};
                \node (r) at (4,0) {L};
                \node (c) at (0,-3) {C};
                \node (h) at (4,-3) {G};
                \draw[<-<]  (r) -- (k) node [midway,above] {};
                \draw[>->] (c) -- (h) node [midway, below] {};
                \draw[>->] (r) -- (h) node[midway, left] {};
                \draw[>->] (k) -- (c) node[midway, left] {};
                \node (k') at (1,-1) {K'};
                \node (r') at (2,-1) {L'};
                \node (c') at (1,-2) {C'};
                \node () at (1.5,-1.5) {$PO$};
                \node () at (3,-1.5) {$PB$};
                \node () at (1.5,-2.5) {$PB$};
                \node () at (1.5,-0.5) {$PB$};
                \node () at (0.5,-1.5) {$PB$};
                \node (x) at (2,-2) {F};
                \draw [>->] (c') -- (x);
                \draw [>->] (r') -- (x);
                \draw [>->] (k') -- (r');
                \draw [>->] (k') -- (c');
                \draw [>->] (c') -- (c);
                \draw[>->] (r') -- (r);
                \draw[>->] (x) -- (h) node[midway,right] {};
                \draw[>->] (k') -- (k) ;
            \end{tikzpicture}
        }
        \end{center} 

    There is a unique monomorphism $h_{XC'} \mathop{\colon} X \rightarrowtail C'$ 
    such that $h_{XC'} \mathop{\star} h_{C'C} \mathop{=} h_{XC}$ and $h_{XC'} \mathop{\star} h_{C'F} \mathop{=} h_{XF}$,
    because we have $h_{XF} \mathop{\star} h_{FG} \mathop{=} h_{XG}$, $h_{XC} \mathop{\star} h_{CG} \mathop{=} h_{XG}$ and $C'CGF$ is a pullback.


    Thus, we have the following commutative diagram:
    \begin{center}
        \resizebox{0.5\textwidth}{!}{
            \begin{tikzpicture}
                \node (k) at (0,0) {K};
                \node (r) at (4,0) {L};
                \node (c) at (0,-3) {C};
                \node (h) at (4,-3) {G};
                % \node (rb) at ($\scl*(1.5,-0.5)$) {$R_X$};
                % \node (h') at ($\scl*(1.5,-1.5)$) {$H'$};
                % \draw[>->]  (rb) -- (h') node [midway,above] {};
                % \draw[>->]  (c) -- (h') node [midway,above] {};
                \draw[<-<]  (r) -- (k) node [midway,above] {};
                \draw[>->] (c) -- (h) node [midway, below] {$$};
                \draw[>->] (r) -- (h) node[midway, left] {$$};
                \draw[>->] (k) -- (c) node[midway, left] {};
                % \draw[->] (rb) to (l);
                % \draw[<-<] (rb) to (k);
                \node (k') at (1,-1) {K'};
                \node (r') at (2,-1) {L'};
                \node (x') at (-2,-2) {X};
                \node (c') at (1,-2) {C'};
                \node () at (1.5,-1.5) {$PO$};
                \node () at (3,-1.5) {$PB$};
                \node () at (1.5,-2.5) {$PB$};
                \node () at (1.5,-0.5) {$PB$};
                \node () at (0.5,-1.5) {$PB$};
                \node (x) at (2,-2) {F};
                % \draw [->] (x) -- (h') node[midway] {!};
                \draw [>->] (c') -- (x);
                \draw [>->] (x') -- (c');
                \draw [>->] (x') -- (c);
                \draw [>->] (x') edge[out=-15,in=210] (x);
                \draw [>->] (r') -- (x);
                \draw [>->] (k') -- (r');
                \draw [>->] (k') -- (c');
                \draw [>->] (c') -- (c);
                % \draw [->] (r') -- (rb);
                % \draw [->] (r') -- (rb);
                % \draw [->] (h') -- (h);
                \draw[>->] (r') -- (r) node[midway,below] {};
                \draw[>->] (x) -- (h) node[midway,right] {};
                % \node (rb) at ($\scl*(\sclx*1.5,-0.5)$) {$R_X$};
                % \node (h') at ($\scl*(\sclx*1.5,-1.2)$) {$H'$};
                % \draw[>->]  (rb) -- (h') node [midway,above] {};
                % \draw[>->]  (c) -- (h') node [midway,above] {};
                % \draw[>->]  (rb) -- (h') node [midway,above] {};
                % \draw[>->] (rb) to (r);
                % \draw[<-<] (rb) to (k);
                \draw[>->] (k') -- (k) ;
            \end{tikzpicture}
        }
        \end{center} 
    
     By the $F$-non-increasing property of $\rho^{-1}$,  
     \todo{assumption: the $F$-non-increasing property of $\rho^{-1}$} 
     there is a morphism $h_{L'R} \mathop{\colon} L' \rightarrowtail R$ such that $K'L'RK$ is a pullback.
        By Lemma~\ref{lem:g_monic}, there is a unique monomorphism $h_{FH} \mathop{\colon} F \rightarrowtail H$ such that $L'FHR$ and $C'FHC$ are commutative squares.
    %      $h_{L'R} \mathop{\star} h_{RH} \mathop{=} h_{L'F} \mathop{\star} h_{FH}$ and $h_{C'F} \mathop{\star} h_{FH} \mathop{=} h_{C'C} \mathop{\star} h_{CH}$, because $K'C'FL'$ is a pushout, and 
    % \begin{flalign*}
    %      &h_{K'L'} \mathop{\star} h_{L'R} \mathop{\star} h_{RH} \\
    %     =&h_{K'K} \mathop{\star} h_{KR} \mathop{\star} h_{RH} \\
    %     =&h_{K'K} \mathop{\star} h_{KC} \mathop{\star} h_{CH} \\
    %     =&h_{K'C'} \mathop{\star} h_{C'C} \mathop{\star} h_{CH}
    % \end{flalign*} 

    Thus, we have the following commutative diagram:
    \begin{center}
        \resizebox{0.5\textwidth}{!}{
            \begin{tikzpicture}
                \node (k) at (0,0) {K};
                \node (r) at (4,0) {R};
                \node (c) at (0,-3) {C};
                \node (h) at (4,-3) {H};
                % \node (rb) at ($\scl*(1.5,-0.5)$) {$R_X$};
                % \node (h') at ($\scl*(1.5,-1.5)$) {$H'$};
                % \draw[>->]  (rb) -- (h') node [midway,above] {};
                % \draw[>->]  (c) -- (h') node [midway,above] {};
                \draw[<-<]  (r) -- (k) node [midway,above] {};
                \draw[>->] (c) -- (h) node [midway, below] {$r'$};
                \draw[>->] (r) -- (h) node[midway, left] {$m'$};
                \draw[>->] (k) -- (c) node[midway, left] {};
                % \draw[->] (rb) to (l);
                % \draw[<-<] (rb) to (k);
                \node (k') at (1,-1) {K'};
                \node (r') at (2,-1) {L'};
                \node (x') at (-2,-2) {X};
                \node (c') at (1,-2) {C'};
                \node () at (1.5,-1.5) {$PO$};
                % \node () at (3,-1.5) {$PB$};
                % \node () at (1.5,-2.5) {$PB$};
                \node () at (1.5,-0.5) {$PB$};
                \node () at (0.5,-1.5) {$PB$};
                \node (x) at (2,-2) {F};
                % \draw [->] (x) -- (h') node[midway] {!};
                \draw [>->] (c') -- (x);
                \draw [>->] (x') -- (c');
                \draw [>->] (x') -- (c);
                \draw [>->] (x') edge[out=-15,in=210] (x);
                \draw [>->] (r') -- (x);
                \draw [>->] (k') -- (r');
                \draw [>->] (k') -- (c');
                \draw [>->] (c') -- (c);
                % \draw [->] (r') -- (rb);
                % \draw [->] (r') -- (rb);
                % \draw [->] (h') -- (h);
                \draw[>->] (r') -- (r) node[midway,below] {};
                \draw[>->] (x) -- (h) node[midway,right] {};
                % \node (rb) at ($\scl*(\sclx*1.5,-0.5)$) {$R_X$};
                % \node (h') at ($\scl*(\sclx*1.5,-1.2)$) {$H'$};
                % \draw[>->]  (rb) -- (h') node [midway,above] {};
                % \draw[>->]  (c) -- (h') node [midway,above] {};
                % \draw[>->]  (rb) -- (h') node [midway,above] {};
                % \draw[>->] (rb) to (r);
                % \draw[<-<] (rb) to (k);
                \draw[>->] (k') -- (k) ;
            \end{tikzpicture}
        }
        \end{center} 

        Finally, we have:
        \begin{flalign*}
            \phi(h_{XG}) &= h_{XC} \mathop{\star} r' \\
            &= h_{XC'} \mathop{\star} h_{C'C} \mathop{\star} r' \\
            &= (h_{XC'} \mathop{\star} h_{C'F}) \mathop{\star} h_{FH}.
        \end{flalign*} 
    \end{itemize}
    We have thus shown that $\phi$ is well defined.
    
    Proving its injectivity is straightforward. 
    Let $h_{XG}, h_{XG}' \mathop{\in} \operatorname{Mono}(\mathcal{X},G,\lnot m, l')_{\operatorname{F}}$ such that $h_{XG} \mathop{\neq} h_{XG}'$.
    There are morphisms $h_{XC}$ and $h_{XC}'$ such that $h_{XG} \mathop{=} h_{XC} \mathop{\star} l'$ and $h_{XG}' \mathop{=} h_{XC}' \mathop{\star} l'$.
    By injectivity of $l'$, we have $h_{XC} \mathop{\neq} h_{XC}'$. Therefore, since $r'$ is injective, we have $\phi(h_{XG}) \mathop{=} h_{XC} \mathop{\star} r' \mathop{\neq} h_{XC}' \mathop{\star} r' \mathop{=} \phi(h_{XG}')$.
\end{proof} 

\noindent
\begin*{\textbf{Lemma~\ref{antipattern:lem:xglnotmlnotlp_xhlnotmrnotrp}}}
        let \( \rho \mathop{=} (L \overset{l}{\leftarrowtail} K \overset{r}{\rightarrowtail} R) \) be a rule, and let $\mathcal{X}=(X,\mathcal{C})$ be a ruler-graph. Suppose that $\rho$ and $\rho^{-1}$ are $X$-non-increasing. For the DPO diagram shown below,
        \begin{center}
        \begin{tikzpicture}
            % [node distance=15mm]
            \node (I) {$K$};
            \node (L) [left of=I] {$L$};
            \node (R) [right of=I] {$R$};
            \node (G) [below of=L] {$G$};
            \node (C) [below of=I] {$C$};
            \node (H) [below of=R] {$H$}; 
            \draw [>->] (I) to  node [midway,below] {$l$} (L);
            \draw [>->] (I) to  node [midway,below] {$r$} (R);
            \draw [>->] (L) to node [midway,right] {$m$} (G);
            \draw [>->] (I) to  node [midway,right] {$u$} (C); 
            \draw [>->] (R) to  node [midway,right] {$m'$} (H);
            \draw [>->] (C) to node [midway,above] {$l'$} (G);
            \draw [>->] (C) to node [midway,above] {$r'$} (H);
            % \node [at=($(I)!.5!(G)$)] {\normalfont PO};
            % \node [at=($(I)!.5!(H)$)] {\normalfont PO};
        \end{tikzpicture}
    \end{center}
    the following inequality holds:
    $$ 
        \card{\operatorname{Mono}(\mathcal{X},G,\mathop{\lnot} m, \mathop{\lnot} l')_{\operatorname{NF}}} \geq
        \card{\operatorname{Mono}(\mathcal{X},H,\mathop{\lnot} m', \mathop{\lnot} r')_{\operatorname{NF}}}.
    $$ 
\end*{}
\begin{proof} 
      \label{antipattern:proof:lem:xglnotmlnotlp_xhlnotmrnotrp}
    There are two cases to consider: $\mathcal{C} \mathop{=} \emptyset$ or $\mathcal{C} \mathop{=} \set{f:X \rightarrowtail F}$. For the first case, the inequality is equivalent to the following inequality:
 $$
        \card{\operatorname{Mono}(\mathcal{X},G,\lnot m, \lnot l')} \geq
        \card{\operatorname{Mono}(\mathcal{X},H,\lnot m', \lnot r')},
    $$
     and the result follows from Lemma~\ref{subgraph_counting:lem:w_u_l_not_geq_r_not}. 
    Let us now consider the second case.
     \todo{assumption for Lemma~\ref{subgraph_counting:lem:w_u_l_not_geq_r_not}: $\rho$ is $X$-non-increasing} The following equalities and inequalities hold:
    \begin{flalign*}
        & \card{\operatorname{Mono}(\mathcal{X},G,\lnot m, \lnot l')_{\operatorname{NF}}} - 
        \card{\operatorname{Mono}(\mathcal{X},H,\lnot m', \lnot r')_{\operatorname{NF}}} 
        \\
        \mathop{=} &(\card{\operatorname{Mono}(X,G,\lnot m, \lnot l')} - \card{\operatorname{Mono}(\mathcal{X},G,\lnot m, \lnot l')_{\operatorname{F}}}) -
            \\ 
           &(\card{\operatorname{Mono}(X,H,\lnot m', \lnot r')} - \card{\operatorname{Mono}(\mathcal{X},H,\lnot m', \lnot r')_{\operatorname{F}}})
        \\
        \mathop{=} &(\card{\operatorname{Mono}(X,G,\lnot m, \lnot l')} - \card{\operatorname{Mono}(X,H,\lnot m', \lnot r')})\mathop{+}
        \\ 
        &
           (\card{\operatorname{Mono}(\mathcal{X},H,\lnot m', \lnot r')_{\operatorname{F}}} - 
           \card{\operatorname{Mono}(\mathcal{X},G,\lnot m, \lnot l')_{\operatorname{F}}})
           \\
        \mathop{\geq} & 
            \card{\operatorname{Mono}(\mathcal{X},H,\lnot m', \lnot r')_{\operatorname{F}}} 
            - 
            \card{\operatorname{Mono}(\mathcal{X},G,\lnot m, \lnot l')_{\operatorname{F}}}.
        &\text{by Lemma~\ref{subgraph_counting:lem:w_u_l_not_geq_r_not}}
    \end{flalign*}    

    % We prove this lemma by show that the function $\phi: \operatorname{Mono}(\mathcal{X},G,\lnot m, \lnot l')_{\operatorname{F}} \rightarrowtail \operatorname{Mono}(X,H,\lnot r', \lnot r')_{\operatorname{NF}}$
    % defined by $\phi(h_{XG}) \mathop{=} h_{XC} \mathop{\star} \lnot r'$ is injective where $h_{XC}:X \rightarrowtail C$ is a monomorphism such that $h_{XG} \mathop{=} h_{XC} \mathop{\star} \lnot l'$.

    We prove $\card{\operatorname{Mono}(\mathcal{X},H,\lnot m', \lnot r')_{\operatorname{F}}} \mathop{\geq} \card{\operatorname{Mono}(\mathcal{X},G,\lnot m, \lnot l')_{\operatorname{F}}}$ by showing
    that we can construct an injection $\phi \mathop{\colon}  \operatorname{Mono}(\mathcal{X},G,\lnot m, \lnot l')_{\operatorname{F}} \mathop{\to} \operatorname{Mono}(\mathcal{X},H,\lnot m', \lnot r')_{\operatorname{F}}$.

    Let $h_{XG} \mathop{\in} \operatorname{Mono}(\mathcal{X},G,\lnot m, \lnot l')_{\operatorname{F}}$. 
    By Lemma~\ref{antipattern:lem:xh_decomp}, we can construct the following commutative diagram:

    \begin{center}
        \resizebox{0.3\textwidth}{!}{
            \begin{tikzpicture}
                \node (k) at (0,0) {K};
                \node (r) at (4,0) {L};
                \node (c) at (0,-3) {C};
                \node (h) at (4,-3) {G};
                % \node (rb) at ($\scl*(1.5,-0.5)$) {$R_X$};
                % \node (h') at ($\scl*(1.5,-1.5)$) {$H'$};
                % \draw[>->]  (rb) -- (h') node [midway,above] {};
                % \draw[>->]  (c) -- (h') node [midway,above] {};
                \draw[<-<]  (r) -- (k) node [midway,above] {};
                \draw[>->] (c) -- (h) node [midway, below] {};
                \draw[>->] (r) -- (h) node[midway, left] {};
                \draw[>->] (k) -- (c) node[midway, left] {};
                % \draw[->] (rb) to (l);
                % \draw[<-<] (rb) to (k);
                \node (k') at (1,-1) {K'};
                \node (r') at (2,-1) {L'};
                \node (c') at (1,-2) {C'};
                \node () at (1.5,-1.5) {$PO$};
                \node () at (3,-1.5) {$PB$};
                \node () at (1.5,-2.5) {$PB$};
                \node () at (1.5,-0.5) {$PB$};
                \node () at (0.5,-1.5) {$PB$};
                \node (x) at (2,-2) {X};
                % \draw [->] (x) -- (h') node[midway] {!};
                \draw [>->] (c') -- (x);
                \draw [>->] (r') -- (x);
                \draw [>->] (k') -- (r');
                \draw [>->] (k') -- (c');
                \draw [>->] (c') -- (c);
                % \draw [->] (r') -- (rb);
                % \draw [->] (r') -- (rb);
                % \draw [->] (h') -- (h);
                \draw[>->] (r') -- (r);
                \draw[>->] (x) -- (h) node[midway,right] {};
                % \node (rb) at ($\scl*(\sclx*1.5,-0.5)$) {$R_X$};
                % \node (h') at ($\scl*(\sclx*1.5,-1.2)$) {$H'$};
                % \draw[>->]  (rb) -- (h') node [midway,above] {};
                % \draw[>->]  (c) -- (h') node [midway,above] {};
                % \draw[>->]  (rb) -- (h') node [midway,above] {};
                % \draw[>->] (rb) to (r);
                % \draw[<-<] (rb) to (k);
                \draw[>->] (k') -- (k) ;
            \end{tikzpicture}
        }
        \end{center} 

         \todo{assmption: $\rho^{-1}$ is $X$-non-increasing} There is a monomorphism $h_{L'R} \mathop{\colon} L' \rightarrowtail R$ because $\rho^{-1}$ is $X$-non-increasing by assumption. By Lemma~\ref{lem:g_monic}, we can construct the following commutative diagram:

        \begin{center}
            \resizebox{0.3\textwidth}{!}{
                \begin{tikzpicture}
                    \node (k) at (0,0) {K};
                    \node (r) at (4,0) {R};
                    \node (c) at (0,-3) {C};
                    \node (h) at (4,-3) {H};
                    % \node (rb) at ($\scl*(1.5,-0.5)$) {$R_X$};
                    % \node (h') at ($\scl*(1.5,-1.5)$) {$H'$};
                    % \draw[>->]  (rb) -- (h') node [midway,above] {};
                    % \draw[>->]  (c) -- (h') node [midway,above] {};
                    \draw[<-<]  (r) -- (k) node [midway,above] {};
                    \draw[>->] (c) -- (h) node [midway, below] {$r'$};
                    \draw[>->] (r) -- (h) node[midway, left] {$m'$};
                    \draw[>->] (k) -- (c) node[midway, left] {};
                    % \draw[->] (rb) to (l);
                    % \draw[<-<] (rb) to (k);
                    \node (k') at (1,-1) {K'};
                    \node (r') at (2,-1) {L'};
                    \node (c') at (1,-2) {C'};
                    \node () at (1.5,-1.5) {$PO$};
                    % \node () at (3,-1.5) {$PB$};
                    % \node () at (1.5,-2.5) {$PB$};
                    \node () at (1.5,-0.5) {$PB$};
                    \node () at (0.5,-1.5) {$PB$};
                    \node (x) at (2,-2) {X};
                    % \draw [->] (x) -- (h') node[midway] {!};
                    \draw [>->] (c') -- (x);
                    \draw [>->] (r') -- (x);
                    \draw [>->] (k') -- (r');
                    \draw [>->] (k') -- (c');
                    \draw [>->] (c') -- (c);
                    % \draw [->] (r') -- (rb);
                    % \draw [->] (r') -- (rb);
                    % \draw [->] (h') -- (h);
                    \draw[>->] (r') -- (r) node[midway,below] {};
                    \draw[>->] (x) -- (h) node[midway,right] {};
                    % \node (rb) at ($\scl*(\sclx*1.5,-0.5)$) {$R_X$};
                    % \node (h') at ($\scl*(\sclx*1.5,-1.2)$) {$H'$};
                    % \draw[>->]  (rb) -- (h') node [midway,above] {};
                    % \draw[>->]  (c) -- (h') node [midway,above] {};
                    % \draw[>->]  (rb) -- (h') node [midway,above] {};
                    % \draw[>->] (rb) to (r);
                    % \draw[<-<] (rb) to (k);
                    \draw[>->] (k') -- (k) ;
                \end{tikzpicture}
            }
            \end{center} 

       
        We define $\phi(h_{XG}) \overset{def}{=} h_{XH}$. The function $\phi$ is injective, by Lemma~\ref{lem:h_hp_diff_g_gp_diff} \todo{ Lemma~\ref{lem:h_hp_diff_g_gp_diff} uses the assumption: $\rho$ is $X$-non-increasing}. It remains to show that $\phi$ is well defined which requires to show that $\phi(h_{XG})$ is an element in $\operatorname{Mono}(X,H,\lnot m', \lnot r')_{\operatorname{F}}$, which is equivalent to show the following three conditions:
        \begin{enumerate}
            \item $h_{XH}$ cannot be factorized by $r'$,
            \item $h_{XH}$ cannot be factorized by $m'$,
            \item there is a monomorphism $h_{FH} \mathop{\colon} F \rightarrowtail H$ such that $h_{XH} \mathop{=} f \mathop{\star} h_{FH}$.
            % $h_{XH}$ can be decomposed as $x \mathop{\star} y$ where $x$ is a monomorphism from $X$ to $F$ and $y$ is a monomorphism from $F$ to $H$.
        \end{enumerate}

        By Proposition~\ref{prop:graph_adhesive}, $KCHR$ is a VK-square. Therefore, we have the following commutative diagram:

        \begin{center}
            \resizebox{0.3\textwidth}{!}{
                \begin{tikzpicture}
                    \node (k) at (0,0) {K};
                    \node (r) at (4,0) {R};
                    \node (c) at (0,-3) {C};
                    \node (h) at (4,-3) {H};
                    % \node (rb) at ($\scl*(1.5,-0.5)$) {$R_X$};
                    % \node (h') at ($\scl*(1.5,-1.5)$) {$H'$};
                    % \draw[>->]  (rb) -- (h') node [midway,above] {};
                    % \draw[>->]  (c) -- (h') node [midway,above] {};
                    \draw[<-<]  (r) -- (k) node [midway,above] {};
                    \draw[>->] (c) -- (h) node [midway, below] {$r'$};
                    \draw[>->] (r) -- (h) node[midway, left] {$m'$};
                    \draw[>->] (k) -- (c) node[midway, left] {};
                    % \draw[->] (rb) to (l);
                    % \draw[<-<] (rb) to (k);
                    \node (k') at (1,-1) {K'};
                    \node (r') at (2,-1) {L'};
                    \node (c') at (1,-2) {C'};
                    \node () at (1.5,-1.5) {$PO$};
                    \node () at (3,-1.5) {$PB$};
                    \node () at (1.5,-2.5) {$PB$};
                    \node () at (1.5,-0.5) {$PB$};
                    \node () at (0.5,-1.5) {$PB$};
                    \node (x) at (2,-2) {X};
                    % \draw [->] (x) -- (h') node[midway] {!};
                    \draw [>->] (c') -- (x);
                    \draw [>->] (r') -- (x);
                    \draw [>->] (k') -- (r');
                    \draw [>->] (k') -- (c');
                    \draw [>->] (c') -- (c);
                    % \draw [->] (r') -- (rb);
                    % \draw [->] (r') -- (rb);
                    % \draw [->] (h') -- (h);
                    \draw[>->] (r') -- (r) node[midway,below] {};
                    \draw[>->] (x) -- (h) node[midway,right] {};
                    % \node (rb) at ($\scl*(\sclx*1.5,-0.5)$) {$R_X$};
                    % \node (h') at ($\scl*(\sclx*1.5,-1.2)$) {$H'$};
                    % \draw[>->]  (rb) -- (h') node [midway,above] {};
                    % \draw[>->]  (c) -- (h') node [midway,above] {};
                    % \draw[>->]  (rb) -- (h') node [midway,above] {};
                    % \draw[>->] (rb) to (r);
                    % \draw[<-<] (rb) to (k);
                    \draw[>->] (k') -- (k) ;
                \end{tikzpicture}
            }
            \end{center} 

        We show that $h_{XH}$ cannot be factorized by $m'$ by contradiction.
        Assume that the contrary, i.e. there is a morphism $h_{XR} \mathop{\colon} X \rightarrowtail R$ such that $h_{XH} \mathop{=} h_{XR} \mathop{\star} m'$.

        There is a unique monomorphism $h_{XL'} \mathop{\colon} X \rightarrowtail L'$ such that $Id_X \mathop{=} h_{XL'} \mathop{\star} h_{L'X}$ and $h_{XR} \mathop{=} h_{XL'} \mathop{\star} h_{L'R}$ because $L'XHR$ is a pullback and $Id_X \mathop{\star} h_{XH} \mathop{=} h_{XR} \mathop{\star} h_{RH}$.
        Thus, we have the following equalities
         \begin{flalign*}
            h_{XG} \mathop{=} &Id_X  \mathop{\star} h_{XG} 
            \\
            \mathop{=} &h_{XL'} \mathop{\star} (h_{L'X} \mathop{\star} h_{XG})
            \\
            \mathop{=} &(h_{XL'} \mathop{\star} h_{L'L}) \mathop{\star} m
        \end{flalign*}
        which contradicts the fact that $h_{XG}$ cannot be factorized by $m$ for it is an element of $\operatorname{Mono}(\mathcal{X},G,\lnot m, \lnot l')_{\operatorname{F}}$.
        
        Symmetrically, $\phi(h_{XG})$ cannot be factorized by $r'$.

        It remains to show that there is a monomorphism $h_{FH} \mathop{\colon} F \rightarrowtail H$ such that $h_{XH} \mathop{=} f \mathop{\star} h_{FH}$. 
        
        Since $h_{XG} \mathop{\in} \operatorname{Mono}(\mathcal{X},G,\lnot m, \lnot l')_{\operatorname{F}}$, there is a morphism $h_{FG} \mathop{\colon} F \rightarrowtail G$ such that 
            the following equality holds:
        \begin{flalign}
            h_{XG} \mathop{=} f \mathop{\star} h_{FG}. \label{antipattern:eq:xxxhxg}
        \end{flalign}
        Thus, we have the following commutative diagram:

        \begin{center}
            \resizebox{0.3\textwidth}{!}{
                \begin{tikzpicture}
                    \node (k) at (0,0) {K};
                    \node (r) at (4,0) {L};
                    \node (c) at (0,-3) {C};
                    \node (x) at (2,-2) {X};
                    \node (h) at (4,-3) {G};
                    \node (f) at (3,-2.5) {F};
                    % \node (rb) at ($\scl*(1.5,-0.5)$) {$R_X$};
                    % \node (h') at ($\scl*(1.5,-1.5)$) {$H'$};
                    % \draw[>->]  (rb) -- (h') node [midway,above] {};
                    % \draw[>->]  (c) -- (h') node [midway,above] {};
                    \draw[<-<]  (r) -- (k) node [midway,above] {};
                    \draw[>->] (c) -- (h) node [midway, below] {$l'$};
                    \draw[>->] (r) -- (h) node[midway, right] {$m$};
                    \draw[>->] (k) -- (c) node[midway, left] {};
                    % \draw[->] (rb) to (l);
                    % \draw[<-<] (rb) to (k);
                    \node (k') at (1,-1) {K'};
                    \node (r') at (2,-1) {L'};
                    \node (c') at (1,-2) {C'};
                    \node () at (1.5,-1.5) {$PO$};
                    \node () at (3,-1.5) {$PB$};
                    \node () at (1.5,-2.5) {$PB$};
                    \node () at (1.5,-0.5) {$PB$};
                    \node () at (0.5,-1.5) {$PB$};
                    % \draw [->] (x) -- (h') node[midway] {!};
                    \draw [>->] (c') -- (x);
                    \draw [>->] (r') -- (x);
                    \draw [>->] (k') -- (r');
                    \draw [>->] (k') -- (c');
                    \draw [>->] (c') -- (c);
                    % \draw [->] (h') -- (h);
                    \draw[>->] (r') -- (r) node[midway,below] {};
                    \draw[>->] (x) -- (f) node[midway,right] {};
                    \draw[>->] (f) -- (h) node[midway,right] {};
                    % \node (rb) at ($\scl*(\sclx*1.5,-0.5)$) {$R_X$};
                    % \node (h') at ($\scl*(\sclx*1.5,-1.2)$) {$H'$};
                    % \draw[>->]  (rb) -- (h') node [midway,above] {};
                    % \draw[>->]  (c) -- (h') node [midway,above] {};
                    % \draw[>->]  (rb) -- (h') node [midway,above] {};
                    % \draw[>->] (rb) to (r);
                    % \draw[<-<] (rb) to (k);
                    \draw[>->] (k') -- (k) ;
                \end{tikzpicture}
            }
            \end{center} 


    Since we work in \textbf{Graph}, we can construct the following commutative diagram where $L''LGF$ and $C''CGF$ are pullbacks: \todo{do we need more details?}
     \begin{center}
            \resizebox{0.3\textwidth}{!}{
                \begin{tikzpicture}
                    \node (k) at (0,0) {K};
                    \node (r) at (4,0) {L};
                    \node (c) at (0,-3) {C};
                    \node (x) at (2,-2) {X};
                    \node (h) at (4,-3) {G};
                    \node (f) at (3,-2.5) {F};
                    % \node (rb) at ($\scl*(1.5,-0.5)$) {$R_X$};
                    % \node (h') at ($\scl*(1.5,-1.5)$) {$H'$};
                    % \draw[>->]  (rb) -- (h') node [midway,above] {};
                    % \draw[>->]  (c) -- (h') node [midway,above] {};
                    \draw[<-<]  (r) -- (k) node [midway,above] {};
                    \draw[>->] (c) -- (h) node [midway, below] {$l'$};
                    \draw[>->] (r) -- (h) node[midway, right] {$m$};
                    \draw[>->] (k) -- (c) node[midway, left] {};
                    % \draw[->] (rb) to (l);
                    % \draw[<-<] (rb) to (k);
                    \node (k') at (1,-1) {K'};
                    \node (r') at (2,-1) {L'};
                    \node (c') at (1,-2) {C'};
                    \node () at (1.5,-1.5) {$PO$};
                    \node (l'') at (3,-1.5) {$L''$};
                    \node (c'') at (1.5,-2.5) {$C''$};
                    \draw[>->] (c'') -- (c);
                    \draw[>->] (c'') --(f);
                    \draw[>->] (l'') --(f);
                    \draw[>->] (l'') --(r);
                    \node () at (1.5,-0.5) {$PB$};
                    \node () at (0.5,-1.5) {$PB$};
                    % \draw [->] (x) -- (h') node[midway] {!};
                    \draw [>->] (c') -- (x);
                    \draw [>->] (r') -- (x);
                    \draw [>->] (k') -- (r');
                    \draw [>->] (k') -- (c');
                    \draw [>->] (c') -- (c);
                    % \draw [->] (h') -- (h);
                    \draw[>->] (r') -- (r) node[midway,below] {};
                    \draw[>->] (x) -- (f) node[midway,right] {};
                    \draw[>->] (f) -- (h) node[midway,right] {};
                    % \node (rb) at ($\scl*(\sclx*1.5,-0.5)$) {$R_X$};
                    % \node (h') at ($\scl*(\sclx*1.5,-1.2)$) {$H'$};
                    % \draw[>->]  (rb) -- (h') node [midway,above] {};
                    % \draw[>->]  (c) -- (h') node [midway,above] {};
                    % \draw[>->]  (rb) -- (h') node [midway,above] {};
                    % \draw[>->] (rb) to (r);
                    % \draw[<-<] (rb) to (k);
                    \draw[>->] (k') -- (k) ;
                \end{tikzpicture}
            }
            \end{center} 

   Since $L''LGF$ is pullback and $h_{L'X} \mathop{\star} h_{XF} \mathop{\star} h_{FG} \mathop{=} h_{L'L''} \mathop{\star} h_{L''L} \mathop{\star} h_{LG}$, by the universal property of pullback, there is a unique morphism $h_{L'L''} \mathop{\colon} L' \rightarrowtail L''$ such that 
   \begin{flalign}
         h_{L'L''} \mathop{\star} h_{L''L} &= h_{L'L} \label{antipattern:lplpplppllpl}, \\
         h_{L'X} \mathop{\star} h_{XF} & \mathop{=} h_{L'L''} \mathop{\star} h_{L''F}.
         \nonumber
        %  \label{antipattern:hlpxhxfhlplpplppf}
   \end{flalign}
   $L'L''L$ and $L'L''FX$ are commutative diagrams. Analogously, since $C''CGF$ is pullback and $C'CHFX$ is a commutative diagram, there is a unique morphism $h_{C'C''} \mathop{\colon} C' \rightarrowtail C''$ such that 
   \begin{flalign}
            h_{C'C''} \mathop{\star} h_{C''C} &= h_{C'C}, \label{antipattern:cpcppc}\\
            h_{C'C''} \mathop{\star} h_{C''F} &= h_{C'X} \mathop{\star} h_{XF}.
            \nonumber 
            %  \label{antipattern:cpxf}
   \end{flalign}
   Thus, the following diagram holds:  
    \begin{center}
            \resizebox{0.5\textwidth}{!}{
                \begin{tikzpicture}
                    \node (k) at (-2,0) {K};
                    \node (r) at (4,0) {L};
                    \node (x) at (2,-2) {X};
                    \node (h) at (4,-3) {G};
                    \node (f) at (3,-2.5) {F};
                    \node (k') at (1,-1) {K'};
                    \node (r') at (2,-1) {L'};
                    \node (c) at (-2,-3) {C};
                    \node (c') at (1,-2) {C'};
                    \node (c'') at (-0.5,-2.5) {$C''$};
                    % \node () at (1.5,-1.5) {$PO$};
                    % \node () at (2.5,-1.5) {$PB$};
                    % \node () at (3.5,-1.5) {$PB$};
                    % \node () at (1.5,-2.3) {$PB$};
                    % \node () at (1.5,-2.8) {$PB$};
                    \node (l'') at (3,-0.5) {$L''$};
                    % \node (rb) at ($\scl*(1.5,-0.5)$) {$R_X$};
                    % \node (h') at ($\scl*(1.5,-1.5)$) {$H'$};
                    % \draw[>->]  (rb) -- (h') node [midway,above] {};
                    % \draw[>->]  (c) -- (h') node [midway,above] {};
                    \draw[<-<]  (r) -- (k) node [midway,above] {};
                    \draw[>->] (c) -- (h) node [midway, below] {$l'$};
                    \draw[>->] (r) -- (h) node[midway, right] {$m$};
                    \draw[>->] (k) -- (c) node[midway, left] {};
                    % \draw[->] (rb) to (l);
                    % \draw[<-<] (rb) to (k);
                    \draw[>->] (c'') -- (c);
                    \draw[>->] (c'') --(f);
                    \draw[>->] (l'') --(f);
                    \draw[>->] (l'') --(r);
                    % \node () at (1.5,-0.5) {$PB$};
                    % \node () at (-0.5,-1.5) {$PB$};
                    % \draw [->] (x) -- (h') node[midway] {!};
                    \draw [>->] (c') -- (x);
                    \draw [>->] (r') -- (x);
                    \draw [>->] (k') -- (r');
                    \draw [>->] (k') -- (c');
                    \draw [>->] (c') -- (c'');
                    % \draw [->] (h') -- (h);
                    \draw[>->] (r') -- (l'') node[midway,below] {};
                    \draw[>->] (x) -- (f) node[midway,right] {};
                    \draw[>->] (f) -- (h) node[midway,right] {};
                    % \node (rb) at ($\scl*(\sclx*1.5,-0.5)$) {$R_X$};
                    % \node (h') at ($\scl*(\sclx*1.5,-1.2)$) {$H'$};
                    % \draw[>->]  (rb) -- (h') node [midway,above] {};
                    % \draw[>->]  (c) -- (h') node [midway,above] {};
                    % \draw[>->]  (rb) -- (h') node [midway,above] {};
                    % \draw[>->] (rb) to (r);
                    % \draw[<-<] (rb) to (k);
                    \draw[>->] (k') -- (k) ;
                \end{tikzpicture}
            } 
            \end{center} 

   Since we work in \textbf{Graph}, we can construct $C'' \leftarrowtail K'' \rightarrowtail L''$ the pullback of $C'' \rightarrowtail F \leftarrowtail L''$, and the following commutative diagram holds:
    \begin{center}
            \resizebox{0.5\textwidth}{!}{
                \begin{tikzpicture}
                    \node (k) at (-2,0) {K};
                    \node (r) at (4,0) {L};
                    \node (x) at (2,-2) {X};
                    \node (h) at (4,-3) {G};
                    \node (f) at (3,-2.5) {F};
                    \node (k') at (1,-1) {K'};
                    \node (r') at (2,-1) {L'};
                    \node (c) at (-2,-3) {C};
                    \node (c') at (1,-2) {C'};
                    \node (c'') at (-0.5,-2.5) {$C''$};
                    \node (k'') at (-0.7,-0.8) {$K''$};
                    % \node () at (1.5,-1.5) {$PO$};
                    % \node () at (2.5,-1.5) {$PB$};
                    % \node () at (3.5,-1.5) {$PB$};
                    % \node () at (1.5,-2.3) {$PB$};
                    % \node () at (1.5,-2.8) {$PB$};
                    % \node () at (1.5,-0.5) {$PB$};
                    % \node () at (-0.5,-1.5) {$PB$};
                    \node (l'') at (3,-0.5) {$L''$};
                    % \node (rb) at ($\scl*(1.5,-0.5)$) {$R_X$};
                    % \node (h') at ($\scl*(1.5,-1.5)$) {$H'$};
                    % \draw[>->]  (rb) -- (h') node [midway,above] {};
                    % \draw[>->]  (c) -- (h') node [midway,above] {};
                    % \draw[>->] (k'') -- (k) node[midway,above] {};
                    % \draw[>->] (k') -- (k'') node[midway,above] {};
                    \draw[>->] (k'') -- (c'') node[midway,above] {};
                    \draw[>->] (k'') -- (l'') node[midway,above] {};
                    \draw[<-<]  (r) -- (k) node [midway,above] {};
                    \draw[>->] (c) -- (h) node [midway, below] {$l'$};
                    \draw[>->] (r) -- (h) node[midway, right] {$m$};
                    \draw[>->] (k) -- (c) node[midway, left] {};
                    % \draw[->] (rb) to (l);
                    % \draw[<-<] (rb) to (k);
                    \draw[>->] (c'') -- (c);
                    \draw[>->] (c'') --(f);
                    \draw[>->] (l'') --(f);
                    \draw[>->] (l'') --(r);

                    % \draw [->] (x) -- (h') node[midway] {!};
                    \draw [>->] (c') -- (x);
                    \draw [>->] (r') -- (x);
                    \draw [>->] (k') -- (r');
                    \draw [>->] (k') -- (c');
                    \draw [>->] (c') -- (c'');
                    % \draw [->] (h') -- (h);
                    \draw[>->] (r') -- (l'') node[midway,below] {};
                    \draw[>->] (x) -- (f) node[midway,right] {};
                    \draw[>->] (f) -- (h) node[midway,right] {};
                    % \node (rb) at ($\scl*(\sclx*1.5,-0.5)$) {$R_X$};
                    % \node (h') at ($\scl*(\sclx*1.5,-1.2)$) {$H'$};
                    % \draw[>->]  (rb) -- (h') node [midway,above] {};
                    % \draw[>->]  (c) -- (h') node [midway,above] {};
                    % \draw[>->]  (rb) -- (h') node [midway,above] {};
                    % \draw[>->] (rb) to (r);
                    % \draw[<-<] (rb) to (k);
                    \draw[>->] (k') -- (k) ;
                \end{tikzpicture}
            }
            \end{center} 

            Since $K''C''FL''$ is pullback and $K'C'C''FL''L'$ is a commutative diagram, by the universal property of pullback, there is a unique morphism $h_{K'K''} \mathop{\colon} K' \rightarrowtail K''$ such that $K'C'C''K''$, $K'L'L''K''$ in the following diagram are commutative:
            \begin{center}
            \resizebox{0.5\textwidth}{!}{
                \begin{tikzpicture}
                    \node (k) at (-2,0) {K};
                    \node (r) at (4,0) {L};
                    \node (x) at (2,-2) {X};
                    \node (h) at (4,-3) {G};
                    \node (f) at (3,-2.5) {F};
                    \node (k') at (1,-1) {K'};
                    \node (r') at (2,-1) {L'};
                    \node (c) at (-2,-3) {C};
                    \node (c') at (1,-2) {C'};
                    \node (c'') at (-0.5,-2.5) {$C''$};
                    \node (k'') at (-0.7,-0.8) {$K''$};
                    % \node () at (1.5,-1.5) {$PO$};
                    % \node () at (2.5,-1.5) {$PB$};
                    % \node () at (3.5,-1.5) {$PB$};
                    % \node () at (1.5,-2.3) {$PB$};
                    % \node () at (1.5,-2.8) {$PB$};
                    % \node () at (1.5,-0.5) {$PB$};
                    % \node () at (-0.5,-1.5) {$PB$};
                    \node (l'') at (3,-0.5) {$L''$};
                    % \node (rb) at ($\scl*(1.5,-0.5)$) {$R_X$};
                    % \node (h') at ($\scl*(1.5,-1.5)$) {$H'$};
                    % \draw[>->]  (rb) -- (h') node [midway,above] {};
                    % \draw[>->]  (c) -- (h') node [midway,above] {};
                    % \draw[>->] (k'') -- (k) node[midway,above] {};
                    % \draw[>->] (k') -- (k'') node[midway,above] {};
                    \draw[>->] (k'') -- (c'') node[midway,above] {};
                    \draw[>->] (k'') -- (l'') node[midway,above] {};
                    \draw[<-<]  (r) -- (k) node [midway,above] {};
                    \draw[>->] (c) -- (h) node [midway, below] {$l'$};
                    \draw[>->] (r) -- (h) node[midway, right] {$m$};
                    \draw[>->] (k) -- (c) node[midway, left] {};
                    \draw[>->] (k') -- (k'') node[midway, left] {};
                    % \draw[>->] (k'') -- (k) node[midway, left] {};
                    % \draw[->] (rb) to (l);
                    % \draw[<-<] (rb) to (k);
                    \draw[>->] (c'') -- (c);
                    \draw[>->] (c'') --(f);
                    \draw[>->] (l'') --(f);
                    \draw[>->] (l'') --(r);

                    % \draw [->] (x) -- (h') node[midway] {!};
                    \draw [>->] (c') -- (x);
                    \draw [>->] (r') -- (x);
                    \draw [>->] (k') -- (r');
                    \draw [>->] (k') -- (c');
                    \draw [>->] (c') -- (c'');
                    % \draw [->] (h') -- (h);
                    \draw[>->] (r') -- (l'') node[midway,below] {};
                    \draw[>->] (x) -- (f) node[midway,right] {};
                    \draw[>->] (f) -- (h) node[midway,right] {};
                    % \node (rb) at ($\scl*(\sclx*1.5,-0.5)$) {$R_X$};
                    % \node (h') at ($\scl*(\sclx*1.5,-1.2)$) {$H'$};
                    % \draw[>->]  (rb) -- (h') node [midway,above] {};
                    % \draw[>->]  (c) -- (h') node [midway,above] {};
                    % \draw[>->]  (rb) -- (h') node [midway,above] {};
                    % \draw[>->] (rb) to (r);
                    % \draw[<-<] (rb) to (k);
                    \draw[>->] (k') -- (k) ;
                \end{tikzpicture}
            }
            \end{center} 

          Similarly, since $KCGL$ is pullback and $K''C''CGLL''$ is a commutative diagram, by the universal property of pullback, there is a unique morphisms $h_{K''K} \mathop{\colon} K'' \rightarrowtail K$ such that $K''C''CK$, $K''L''LK$ in the following diagram are commutative:
 \begin{center}
            \resizebox{0.5\textwidth}{!}{
                \begin{tikzpicture}
                    \node (k) at (-2,0) {K};
                    \node (r) at (4,0) {L};
                    \node (x) at (2,-2) {X};
                    \node (h) at (4,-3) {G};
                    \node (f) at (3,-2.5) {F};
                    \node (k') at (1,-1) {K'};
                    \node (r') at (2,-1) {L'};
                    \node (c) at (-2,-3) {C};
                    \node (c') at (1,-2) {C'};
                    \node (c'') at (-0.5,-2.5) {$C''$};
                    \node (k'') at (-0.7,-0.8) {$K''$};
                    % \node () at (1.5,-1.5) {$PO$};
                    % \node () at (2.5,-1.5) {$PB$};
                    % \node () at (3.5,-1.5) {$PB$};
                    % \node () at (1.5,-2.3) {$PB$};
                    % \node () at (1.5,-2.8) {$PB$};
                    % \node () at (1.5,-0.5) {$PB$};
                    % \node () at (-0.5,-1.5) {$PB$};
                    \node (l'') at (3,-0.5) {$L''$};
                    % \node (rb) at ($\scl*(1.5,-0.5)$) {$R_X$};
                    % \node (h') at ($\scl*(1.5,-1.5)$) {$H'$};
                    % \draw[>->]  (rb) -- (h') node [midway,above] {};
                    % \draw[>->]  (c) -- (h') node [midway,above] {};
                    \draw[>->] (k'') -- (k) node[midway,above] {};
                    % \draw[>->] (k') -- (k'') node[midway,above] {};
                    \draw[>->] (k'') -- (c'') node[midway,above] {};
                    \draw[>->] (k'') -- (l'') node[midway,above] {};
                    \draw[<-<]  (r) -- (k) node [midway,above] {};
                    \draw[>->] (c) -- (h) node [midway, below] {$l'$};
                    \draw[>->] (r) -- (h) node[midway, right] {$m$};
                    \draw[>->] (k) -- (c) node[midway, left] {};
                    \draw[>->] (k') -- (k'') node[midway, left] {};
                    % \draw[>->] (k'') -- (k) node[midway, left] {};
                    % \draw[->] (rb) to (l);
                    % \draw[<-<] (rb) to (k);
                    \draw[>->] (c'') -- (c);
                    \draw[>->] (c'') --(f);
                    \draw[>->] (l'') --(f);
                    \draw[>->] (l'') --(r);

                    % \draw [->] (x) -- (h') node[midway] {!};
                    \draw [>->] (c') -- (x);
                    \draw [>->] (r') -- (x);
                    \draw [>->] (k') -- (r');
                    \draw [>->] (k') -- (c');
                    \draw [>->] (c') -- (c'');
                    % \draw [->] (h') -- (h);
                    \draw[>->] (r') -- (l'') node[midway,below] {};
                    \draw[>->] (x) -- (f) node[midway,right] {};
                    \draw[>->] (f) -- (h) node[midway,right] {};
                    % \node (rb) at ($\scl*(\sclx*1.5,-0.5)$) {$R_X$};
                    % \node (h') at ($\scl*(\sclx*1.5,-1.2)$) {$H'$};
                    % \draw[>->]  (rb) -- (h') node [midway,above] {};
                    % \draw[>->]  (c) -- (h') node [midway,above] {};
                    % \draw[>->]  (rb) -- (h') node [midway,above] {};
                    % \draw[>->] (rb) to (r);
                    % \draw[<-<] (rb) to (k);
                    \draw[>->] (k') -- (k) ;
                \end{tikzpicture}
            }
            \end{center} 

        There is a monomorphism $h_{L''R} \mathop{\colon} L'' \rightarrowtail R$ such that $K''L''RK$ is pullback, because $\rho^{-1}$ is $X$-non-increasing by assump.\todo{$\rho^{-1}$ is $X$-non-increasing} We have, thus, the following commutative diagram:

\begin{center}
            \resizebox{0.5\textwidth}{!}{
                \begin{tikzpicture}
                    \node (k) at (-2,0) {K};
                    \node (r) at (4,0) {R};
                    \node (x) at (2,-2) {X};
                    \node (h) at (4,-3) {H};
                    \node (f) at (3,-2.5) {F};
                    \node (k') at (1,-1) {K'};
                    \node (r') at (2,-1) {L'};
                    \node (c) at (-2,-3) {C};
                    \node (c') at (1,-2) {C'};
                    \node (c'') at (-0.5,-2.5) {$C''$};
                    \node (k'') at (-0.5,-0.5) {$K''$};
                    % \node () at (1.5,-1.5) {$PO$};
                    % \node () at (2.5,-1.5) {$PB$};
                    % \node () at (3.5,-1.5) {$PB$};
                    % \node () at (1.5,-2.3) {$PB$};
                    % \node () at (1.5,-2.8) {$PB$};
                    % \node () at (1.5,-0.5) {$PB$};
                    % \node () at (-0.5,-1.5) {$PB$};
                    \node (l'') at (3,-0.5) {$L''$};
                    % \node (rb) at ($\scl*(1.5,-0.5)$) {$R_X$};
                    % \node (h') at ($\scl*(1.5,-1.5)$) {$H'$};
                    % \draw[>->]  (rb) -- (h') node [midway,above] {};
                    % \draw[>->]  (c) -- (h') node [midway,above] {};
                    % \draw[>->] (k'') -- (k) node[midway,above] {};
                    % \draw[>->] (k') -- (k'') node[midway,above] {};
                    \draw[>->] (k'') -- (c'') node[midway,above] {};
                    \draw[>->] (k'') -- (l'') node[midway,above] {};
                    \draw[<-<]  (r) -- (k) node [midway,above] {};
                    \draw[>->] (c) -- (h) node [midway, below] {$r'$};
                    \draw[>->] (r) -- (h) node[midway, right] {$m'$};
                    \draw[>->] (k) -- (c) node[midway, left] {};
                    \draw[>->] (k') -- (k'') node[midway, left] {};
                    \draw[>->] (k'') -- (k) node[midway, left] {};
                    % \draw[->] (rb) to (l);
                    % \draw[<-<] (rb) to (k);
                    \draw[>->] (c'') -- (c);
                    \draw[>->] (c'') --(f);
                    \draw[>->] (l'') --(f);
                    \draw[>->] (l'') --(r);

                    % \draw [->] (x) -- (h') node[midway] {!};
                    \draw [>->] (c') -- (x);
                    \draw [>->] (r') -- (x);
                    \draw [>->] (k') -- (r');
                    \draw [>->] (k') -- (c');
                    \draw [>->] (c') -- (c'');
                    % \draw [->] (h') -- (h);
                    \draw[>->] (r') -- (l'') node[midway,below] {};
                    \draw[>->] (x) -- (f) node[midway,right] {};
                    % \draw[>->] (f) -- (h) node[midway,right] {};
                    % \node (rb) at ($\scl*(\sclx*1.5,-0.5)$) {$R_X$};
                    % \node (h') at ($\scl*(\sclx*1.5,-1.2)$) {$H'$};
                    % \draw[>->]  (rb) -- (h') node [midway,above] {};
                    % \draw[>->]  (c) -- (h') node [midway,above] {};
                    % \draw[>->]  (rb) -- (h') node [midway,above] {};
                    % \draw[>->] (rb) to (r);
                    % \draw[<-<] (rb) to (k);
                \end{tikzpicture}
            }
            \end{center} 
            
    The square $K''C''FL''$ is pushout because it is pullback, by Proposition~\ref{prop:pb_eq_po}. Furthermore, we have 
    \begin{flalign*}
         &(h_{K''C''} \mathop{\star} h_{C''C}) \mathop{\star} h_{CH} \\
        =&(h_{K''K} \mathop{\star} h_{KC}) \mathop{\star} h_{CH} \\
        \mathop{=} &h_{K''K} \mathop{\star} (h_{KR} \mathop{\star} h_{RH}) \\
        \mathop{=} &h_{K''L''} \mathop{\star} h_{L''R} \mathop{\star} h_{RH}.
    \end{flalign*}
    Therefore, by the universal property of pushout, 
     there is a unique morphism $h_{FH}$ such that $L''FHR$ and $C''FHC$ are commutative diagrams. Thus, we have the following commutative diagram:


\begin{center}
            \resizebox{0.5\textwidth}{!}{
                \begin{tikzpicture}
                    \node (k) at (-2,0) {K};
                    \node (r) at (4,0) {R};
                    \node (x) at (2,-2) {X};
                    \node (h) at (4,-3) {H};
                    \node (f) at (3,-2.5) {F};
                    \node (k') at (1,-1) {K'};
                    \node (r') at (2,-1) {L'};
                    \node (c) at (-2,-3) {C};
                    \node (c') at (1,-2) {C'};
                    \node (c'') at (-0.5,-2.5) {$C''$};
                    \node (k'') at (-0.5,-0.5) {$K''$};
                    % \node () at (1.5,-1.5) {$PO$};
                    % \node () at (2.5,-1.5) {$PB$};
                    % \node () at (3.5,-1.5) {$PB$};
                    % \node () at (1.5,-2.3) {$PB$};
                    % \node () at (1.5,-2.8) {$PB$};
                    % \node () at (1.5,-0.5) {$PB$};
                    % \node () at (-0.5,-1.5) {$PB$};
                    \node (l'') at (3,-0.5) {$L''$};
                    % \node (rb) at ($\scl*(1.5,-0.5)$) {$R_X$};
                    % \node (h') at ($\scl*(1.5,-1.5)$) {$H'$};
                    % \draw[>->]  (rb) -- (h') node [midway,above] {};
                    % \draw[>->]  (c) -- (h') node [midway,above] {};
                    % \draw[>->] (k'') -- (k) node[midway,above] {};
                    % \draw[>->] (k') -- (k'') node[midway,above] {};
                    \draw[>->] (k'') -- (c'') node[midway,above] {};
                    \draw[>->] (k'') -- (l'') node[midway,above] {};
                    \draw[<-<]  (r) -- (k) node [midway,above] {};
                    \draw[>->] (c) -- (h) node [midway, below] {$r'$};
                    \draw[>->] (r) -- (h) node[midway, right] {$m'$};
                    \draw[>->] (k) -- (c) node[midway, left] {};
                    \draw[>->] (k') -- (k'') node[midway, left] {};
                    \draw[>->] (k'') -- (k) node[midway, left] {};
                    % \draw[->] (rb) to (l);
                    % \draw[<-<] (rb) to (k);
                    \draw[>->] (c'') -- (c);
                    \draw[>->] (c'') --(f);
                    \draw[>->] (l'') --(f);
                    \draw[>->] (l'') --(r);

                    % \draw [->] (x) -- (h') node[midway] {!};
                    \draw [>->] (c') -- (x);
                    \draw [>->] (r') -- (x);
                    \draw [>->] (k') -- (r');
                    \draw [>->] (k') -- (c');
                    \draw [>->] (c') -- (c'');
                    % \draw [->] (h') -- (h);
                    \draw[>->] (r') -- (l'') node[midway,below] {};
                    \draw[>->] (x) -- (f) node[midway,right] {};
                    \draw[>->] (f) -- (h) node[midway,right] {};
                    % \node (rb) at ($\scl*(\sclx*1.5,-0.5)$) {$R_X$};
                    % \node (h') at ($\scl*(\sclx*1.5,-1.2)$) {$H'$};
                    % \draw[>->]  (rb) -- (h') node [midway,above] {};
                    % \draw[>->]  (c) -- (h') node [midway,above] {};
                    % \draw[>->]  (rb) -- (h') node [midway,above] {};
                    % \draw[>->] (rb) to (r);
                    % \draw[<-<] (rb) to (k);
                \end{tikzpicture}
            }
            \end{center} 
    We have $h_{XH} \mathop{=} h_{XF} \mathop{\star} h_{FH}$ because $K'C'XL'$ is a pushout and $h_{XH}$ 
        is the unique morphism such that 
        \begin{flalign*}
           H_{C'X} \mathop{\star} h_{XH} &= h_{C'C} \mathop{\star} h_{CH} \\
           &=  h_{C'C''} \mathop{\star} h_{C''C} \mathop{\star} h_{CH} \hspace{2cm} \text{by Equation~\eqref{antipattern:cpcppc}}
        \end{flalign*} and 
        \begin{flalign*}
        h_{L'X} \mathop{\star} h_{XH} &= h_{L'R} \mathop{\star} h_{RH} \\
        &= h_{L'L''} \mathop{\star} h_{L''R} \mathop{\star} h_{RH}, \hspace{2cm} \text{by Equation~\eqref{antipattern:lplpplppllpl}}
        \end{flalign*}
         but we have 
        \begin{flalign*}
            &h_{C'X} \mathop{\star} (h_{XF} \mathop{\star} h_{FH}) 
            \\= & (h_{C'C''} \mathop{\star} h_{C''F}) \mathop{\star} h_{FH} 
            \\= & h_{C'C''} \mathop{\star} (h_{C''C} \mathop{\star} h_{CH})
        \end{flalign*}
        and 
        \begin{flalign*}
            &h_{L'X} \mathop{\star} (h_{XF} \mathop{\star} h_{FH})
            \\ \mathop{=} & (h_{L'L''} \mathop{\star} h_{L''F}) \mathop{\star} h_{FH}
            \\ \mathop{=} & h_{L'L''} \mathop{\star} (h_{L''R} \mathop{\star} h_{RH}).
        \end{flalign*}
        This concludes our proof.

%     \color{red}

%         \begin{center}
%             \resizebox{0.3\textwidth}{!}{
%                 \begin{tikzpicture}
%                     \node (k) at (0,0) {K};
%                     \node (r) at (4,0) {R};
%                     \node (c) at (0,-3) {C};
%                     \node (x) at (2,-2) {X};
%                     \node (h) at (4,-3) {H};
%                     \node (f) at (3,-2.5) {F};
%                     % \node (rb) at ($\scl*(1.5,-0.5)$) {$R_X$};
%                     % \node (h') at ($\scl*(1.5,-1.5)$) {$H'$};
%                     % \draw[>->]  (rb) -- (h') node [midway,above] {};
%                     % \draw[>->]  (c) -- (h') node [midway,above] {};
%                     \draw[<-<]  (r) -- (k) node [midway,above] {};
%                     \draw[>->] (c) -- (h) node [midway, below] {$r'$};
%                     \draw[>->] (r) -- (h) node[midway, left] {$m'$};
%                     \draw[>->] (k) -- (c) node[midway, left] {};
%                     % \draw[->] (rb) to (l);
%                     % \draw[<-<] (rb) to (k);
%                     \node (k') at (1,-1) {K'};
%                     \node (r') at (2,-1) {L'};
%                     \node (c') at (1,-2) {C'};
%                     \node () at (1.5,-1.5) {$PO$};
%                     \node () at (3,-1.5) {$PB$};
%                     \node () at (1.5,-2.5) {$PB$};
%                     \node () at (1.5,-0.5) {$PB$};
%                     \node () at (0.5,-1.5) {$PB$};
%                     % \draw [->] (x) -- (h') node[midway] {!};
%                     \draw [>->] (c') -- (x);
%                     \draw [>->] (r') -- (x);
%                     \draw [>->] (k') -- (r');
%                     \draw [>->] (k') -- (c');
%                     \draw [>->] (c') -- (c);
%                     % \draw [->] (h') -- (h);
%                     \draw[>->] (r') -- (r) node[midway,below] {};
%                     \draw[>->] (x) -- (f) node[midway,right] {};
%                     \draw[>->] (f) -- (h) node[midway,right] {};
%                     % \node (rb) at ($\scl*(\sclx*1.5,-0.5)$) {$R_X$};
%                     % \node (h') at ($\scl*(\sclx*1.5,-1.2)$) {$H'$};
%                     % \draw[>->]  (rb) -- (h') node [midway,above] {};
%                     % \draw[>->]  (c) -- (h') node [midway,above] {};
%                     % \draw[>->]  (rb) -- (h') node [midway,above] {};
%                     % \draw[>->] (rb) to (r);
%                     % \draw[<-<] (rb) to (k);
%                     \draw[>->] (k') -- (k) ;
%                 \end{tikzpicture}
%             }
%             \end{center} 

%     We can now construct the following commutative diagram where $L''RHF$ and $C''CHF$ are pullbacks: \todo{do we need more details?}
%      \begin{center}
%             \resizebox{0.3\textwidth}{!}{
%                 \begin{tikzpicture}
%                     \node (k) at (0,0) {K};
%                     \node (r) at (4,0) {R};
%                     \node (c) at (0,-3) {C};
%                     \node (x) at (2,-2) {X};
%                     \node (h) at (4,-3) {H};
%                     \node (f) at (3,-2.5) {F};
%                     % \node (rb) at ($\scl*(1.5,-0.5)$) {$R_X$};
%                     % \node (h') at ($\scl*(1.5,-1.5)$) {$H'$};
%                     % \draw[>->]  (rb) -- (h') node [midway,above] {};
%                     % \draw[>->]  (c) -- (h') node [midway,above] {};
%                     \draw[<-<]  (r) -- (k) node [midway,above] {};
%                     \draw[>->] (c) -- (h) node [midway, below] {$r'$};
%                     \draw[>->] (r) -- (h) node[midway, left] {$m'$};
%                     \draw[>->] (k) -- (c) node[midway, left] {};
%                     % \draw[->] (rb) to (l);
%                     % \draw[<-<] (rb) to (k);
%                     \node (k') at (1,-1) {K'};
%                     \node (r') at (2,-1) {L'};
%                     \node (c') at (1,-2) {C'};
%                     \node () at (1.5,-1.5) {$PO$};
%                     \node (l'') at (3,-1.5) {$L''$};
%                     \node (c'') at (1.5,-2.5) {$C''$};
%                     \draw[>->] (c'') -- (c);
%                     \draw[>->] (c'') --(f);
%                     \draw[>->] (l'') --(f);
%                     \draw[>->] (l'') --(r);
%                     \node () at (1.5,-0.5) {$PB$};
%                     \node () at (0.5,-1.5) {$PB$};
%                     % \draw [->] (x) -- (h') node[midway] {!};
%                     \draw [>->] (c') -- (x);
%                     \draw [>->] (r') -- (x);
%                     \draw [>->] (k') -- (r');
%                     \draw [>->] (k') -- (c');
%                     \draw [>->] (c') -- (c);
%                     % \draw [->] (h') -- (h);
%                     \draw[>->] (r') -- (r) node[midway,below] {};
%                     \draw[>->] (x) -- (f) node[midway,right] {};
%                     \draw[>->] (f) -- (h) node[midway,right] {};
%                     % \node (rb) at ($\scl*(\sclx*1.5,-0.5)$) {$R_X$};
%                     % \node (h') at ($\scl*(\sclx*1.5,-1.2)$) {$H'$};
%                     % \draw[>->]  (rb) -- (h') node [midway,above] {};
%                     % \draw[>->]  (c) -- (h') node [midway,above] {};
%                     % \draw[>->]  (rb) -- (h') node [midway,above] {};
%                     % \draw[>->] (rb) to (r);
%                     % \draw[<-<] (rb) to (k);
%                     \draw[>->] (k') -- (k) ;
%                 \end{tikzpicture}
%             }
%             \end{center} 

%     Since $L''RHF$ and $C''CHF$ are pullbacks, and $L'XFHR$ and $C'CHFX$ are commutative diagrams, by the universal property of pullbacks, there are unique morphisms $h_{L'L''}:L' \rightarrowtail L''$ and $h_{C'C''}:C' \rightarrowtail C''$ such that $L'L''R$, $L'L''FX$, $C'C''FX$, and $C'CC''$ are commutative diagrams. Thus, the following diagram holds: 
%     \begin{center}
%             \resizebox{0.5\textwidth}{!}{
%                 \begin{tikzpicture}
%                     \node (k) at (-2,0) {K};
%                     \node (r) at (4,0) {R};
%                     \node (x) at (2,-2) {X};
%                     \node (h) at (4,-3) {H};
%                     \node (f) at (3,-2.5) {F};
%                     \node (k') at (1,-1) {K'};
%                     \node (r') at (2,-1) {L'};
%                     \node (c) at (-2,-3) {C};
%                     \node (c') at (1,-2) {C'};
%                     \node (c'') at (-0.5,-2.5) {$C''$};
%                     \node () at (1.5,-1.5) {$PO$};
%                     \node () at (2.5,-1.5) {$PB$};
%                     \node () at (3.5,-1.5) {$PB$};
%                     \node () at (1.5,-2.3) {$PB$};
%                     \node () at (1.5,-2.8) {$PB$};
%                     \node (l'') at (3,-0.5) {$L''$};
%                     % \node (rb) at ($\scl*(1.5,-0.5)$) {$R_X$};
%                     % \node (h') at ($\scl*(1.5,-1.5)$) {$H'$};
%                     % \draw[>->]  (rb) -- (h') node [midway,above] {};
%                     % \draw[>->]  (c) -- (h') node [midway,above] {};
%                     \draw[<-<]  (r) -- (k) node [midway,above] {};
%                     \draw[>->] (c) -- (h) node [midway, below] {$r'$};
%                     \draw[>->] (r) -- (h) node[midway, right] {$m'$};
%                     \draw[>->] (k) -- (c) node[midway, left] {};
%                     % \draw[->] (rb) to (l);
%                     % \draw[<-<] (rb) to (k);
%                     \draw[>->] (c'') -- (c);
%                     \draw[>->] (c'') --(f);
%                     \draw[>->] (l'') --(f);
%                     \draw[>->] (l'') --(r);
%                     \node () at (1.5,-0.5) {$PB$};
%                     \node () at (-0.5,-1.5) {$PB$};
%                     % \draw [->] (x) -- (h') node[midway] {!};
%                     \draw [>->] (c') -- (x);
%                     \draw [>->] (r') -- (x);
%                     \draw [>->] (k') -- (r');
%                     \draw [>->] (k') -- (c');
%                     \draw [>->] (c') -- (c'');
%                     % \draw [->] (h') -- (h);
%                     \draw[>->] (r') -- (l'') node[midway,below] {};
%                     \draw[>->] (x) -- (f) node[midway,right] {};
%                     \draw[>->] (f) -- (h) node[midway,right] {};
%                     % \node (rb) at ($\scl*(\sclx*1.5,-0.5)$) {$R_X$};
%                     % \node (h') at ($\scl*(\sclx*1.5,-1.2)$) {$H'$};
%                     % \draw[>->]  (rb) -- (h') node [midway,above] {};
%                     % \draw[>->]  (c) -- (h') node [midway,above] {};
%                     % \draw[>->]  (rb) -- (h') node [midway,above] {};
%                     % \draw[>->] (rb) to (r);
%                     % \draw[<-<] (rb) to (k);
%                     \draw[>->] (k') -- (k) ;
%                 \end{tikzpicture}
%             }
%             \end{center} 

%     Similarly, we can construct the following commutative diagram where $C''CKK''$ is pullback:
% \begin{center}
%             \resizebox{0.5\textwidth}{!}{
%                 \begin{tikzpicture}
%                     \node (k) at (-2,0) {K};
%                     \node (r) at (4,0) {R};
%                     \node (x) at (2,-2) {X};
%                     \node (h) at (4,-3) {H};
%                     \node (f) at (3,-2.5) {F};
%                     \node (k') at (1,-1) {K'};
%                     \node (r') at (2,-1) {L'};
%                     \node (c) at (-2,-3) {C};
%                     \node (c') at (1,-2) {C'};
%                     \node (c'') at (-0.5,-2.5) {$C''$};
%                     \node () at (1.5,-1.5) {$PO$};
%                                      \node () at (2.5,-1.5) {$PB$};
%                     \node () at (3.5,-1.5) {$PB$};
%                     \node () at (1.5,-2.3) {$PB$};
%                     \node () at (1.5,-2.8) {$PB$};
%                     \node (l'') at (3,-0.5) {$L''$};
%                     \node (k'') at (-0.5,-0.5) {$K''$};
%                     % \node (k3) at (0,-0.66) {$K^3$};
%                     \draw[>->]  (k') -- (k'') node [midway,above] {};
%                     % \draw[>->]  (k3) -- (k2) node [midway,above] {};
%                     \draw[>->]  (k'') -- (k) node [midway,above] {};
%                     \draw[>->]  (k'') -- (c'') node [midway,above] {};
%                     % \draw[>->]  (k'') -- (l'') node [midway,above] {};
%                     % \node (rb) at ($\scl*(1.5,-0.5)$) {$R_X$};
%                     % \node (h') at ($\scl*(1.5,-1.5)$) {$H'$};
%                     % \draw[>->]  (rb) -- (h') node [midway,above] {};
%                     % \draw[>->]  (c) -- (h') node [midway,above] {};
%                     \draw[<-<]  (r) -- (k) node [midway,above] {};
%                     \draw[>->] (c) -- (h) node [midway, below] {$r'$};
%                     \draw[>->] (r) -- (h) node[midway, right] {$m'$};
%                     \draw[>->] (k) -- (c) node[midway, left] {};
%                     % \draw[->] (rb) to (l);
%                     % \draw[<-<] (rb) to (k);
%                     \draw[>->] (c'') -- (c);
%                     \draw[>->] (c'') --(f);
%                     \draw[>->] (l'') --(f);
%                     \draw[>->] (l'') --(r);
%                     % \node () at (1.5,-0.5) {$PB$};
%                     % \node () at (-0.5,-1.5) {$PB$};
%                     % \draw [->] (x) -- (h') node[midway] {!};
%                     \draw [>->] (c') -- (x);
%                     \draw [>->] (r') -- (x);
%                     \draw [>->] (k') -- (r');
%                     \draw [>->] (k') -- (c');
%                     \draw [>->] (c') -- (c'');
%                     % \draw [->] (h') -- (h);
%                     \draw[>->] (r') -- (l'') node[midway,below] {};
%                     \draw[>->] (x) -- (f) node[midway,right] {};
%                     \draw[>->] (f) -- (h) node[midway,right] {};
%                     % \node (rb) at ($\scl*(\sclx*1.5,-0.5)$) {$R_X$};
%                     % \node (h') at ($\scl*(\sclx*1.5,-1.2)$) {$H'$};
%                     % \draw[>->]  (rb) -- (h') node [midway,above] {};
%                     % \draw[>->]  (c) -- (h') node [midway,above] {};
%                     % \draw[>->]  (rb) -- (h') node [midway,above] {};
%                     % \draw[>->] (rb) to (r);
%                     % \draw[<-<] (rb) to (k);
%                     % \draw[>->] (k') -- (k) ;
%                 \end{tikzpicture}
%             }
%             \end{center} 

   


%     % There is $h_{L'L''}:L' \rightarrowtail L''$ such that $h_{L'L} \mathop{=} h_{L'L''} \mathop{\star} h_{L''L}$ 
%     % because $L''LGF$ is a pullback, and the following diagram holds:



%     % \begin{center}
%     %     \resizebox{0.3\textwidth}{!}{
%     %         \begin{tikzpicture}
%     %             \node (k) at (0,0) {K};
%     %             \node (r) at (4,0) {L};
%     %             \node (c) at (0,-3) {C};
%     %             \node (h) at (4,-3) {G};
%     %             \draw[<-<]  (r) -- (k) node [midway,above] {};
%     %             \draw[>->] (c) -- (h) node [midway, below] {};
%     %             \draw[>->] (r) -- (h) node[midway, left] {$m$};
%     %             \draw[>->] (k) -- (c) node[midway, left] {};
%     %             \node (k') at (1,-1) {K'};
%     %             \node (r') at (2,-1) {L'};
%     %             \node (c') at (1,-2) {C'};
%     %             \node () at (1.5,-1.5) {$PO$};
%     %             \node () at (3,-1.5) {$PB$};
%     %             \node () at (1.5,-2.5) {$PB$};
%     %             \node () at (1.5,-0.5) {$PB$};
%     %             \node () at (0.5,-1.5) {$PB$};
%     %             \node (x) at (2,-2) {X};
%     %             \draw [>->] (c') -- (x);
%     %             \draw [>->] (r') -- (x);
%     %             \draw [>->] (k') -- (r');
%     %             \draw [>->] (k') -- (c');
%     %             \draw [>->] (c') -- (c);
%     %             \draw[>->] (r') -- (r);
%     %             \draw[>->] (x) -- (h) node[midway,right] {};
%     %             \draw[>->] (k') -- (k) ;
%     %         \end{tikzpicture}
%     %     }
%     %     \end{center} 
%     % where
%     % \begin{flalign*}
%     %     h_{L'L} \mathop{\star} m &= h_{L'X} \mathop{\star} h_{XG} \\
%     %                     &= h_{L'X} \mathop{\star} h_{XF} \mathop{\star} h_{FG}
%     %                     &\text{by Equation~\eqref{antipattern:eq:xxxhxg}}
%     % \end{flalign*} 
%     % Analogously, there exist $h_{C'C''}:C' \rightarrowtail C''$ such that $h_{C'C} \mathop{=} h_{C'C''} \mathop{\star} h_{C''C}$ 
%     % and $h_{K'K''}:K' \rightarrowtail K''$ such that $h_{K'K} \mathop{=} h_{K'K''} \mathop{\star} h_{K''K}$.
    
%      \todo{details}We can show that the following diagram holds:

%     \begin{center}
%         \resizebox{0.3\textwidth}{!}{
%             \begin{tikzpicture}
%                 \node (kk) at (-1,1) {K};
%                 \node (rr) at (5,1) {R};
%                 \node (cc) at (-1,-4) {C};
%                 \node (hh) at (5,-4) {H};
%                 \node (k) at (0,0) {K''};
%                 \node (r) at (4,0) {L''};
%                 \node (c) at (0,-3) {C''};
%                 \node (h) at (4,-3) {F};
%                 \draw[>->]  (k) -- (kk) node [midway,above] {};
%                 \draw[>->]  (c) -- (cc) node [midway,above] {};
%                 \draw[>->]  (h) -- (hh) node [midway,above] {};
%                 \draw[>->]  (r) -- (rr) node [midway,above] {};
%                 \draw[>->]  (cc) -- (hh) node [midway,below] {$r'$};
%                 \draw[>->]  (kk) -- (cc) node [midway,left] {$u$};
%                 \draw[>->]  (kk) -- (rr) node [midway,above] {$r$};
%                 \draw[>->]  (rr) -- (hh) node [midway,right] {$m'$};
%                 % \node (rb) at ($\scl*(1.5,-0.5)$) {$R_X$};
%                 % \node (h') at ($\scl*(1.5,-1.5)$) {$H'$};
%                 % \draw[>->]  (rb) -- (h') node [midway,above] {};
%                 % \draw[>->]  (c) -- (h') node [midway,above] {};
%                 \draw[<-<]  (r) -- (k) node [midway,above] {};
%                 \draw[>->] (c) -- (h) node [midway, below] {};
%                 \draw[>->] (r) -- (h) node[midway, left] {};
%                 \draw[>->] (k) -- (c) node[midway, left] {};
%                 % \draw[->] (rb) to (l);
%                 % \draw[<-<] (rb) to (k);
%                 \node (k') at (1,-1) {K'};
%                 \node (r') at (2,-1) {L'};
%                 \node (c') at (1,-2) {C'};
%                 \node () at (1.5,-1.5) {$PO$};
%                 \node () at (3,-1.5) {$PB$};
%                 \node () at (1.5,-2.5) {$PB$};
%                 \node () at (1.5,-0.5) {$PB$};
%                 \node () at (0.5,-1.5) {$PB$};
%                 \node () [at=($(cc)!0.5!(h)$)] {$PB$};
%                 \node () [at=($(cc)!0.5!(k)$)] {$PB$};
%                 \node () [at=($(rr)!0.5!(k)$)] {$PB$};
%                 \node () [at=($(rr)!0.5!(h)$)] {$PB$};
%                 \node (x) at (2,-2) {X};
%                 % \draw [->] (x) -- (h') node[midway] {!};
%                 \draw [>->] (c') -- (x);
%                 \draw [>->] (r') -- (x);
%                 \draw [>->] (k') -- (r');
%                 \draw [>->] (k') -- (c');
%                 \draw [>->] (c') -- (c);
%                 % \draw [->] (r') -- (rb);
%                 % \draw [->] (r') -- (rb);
%                 % \draw [->] (h') -- (h);
%                 \draw[>->] (r') -- (r) node[midway,below] {};
%                 \draw[>->] (x) -- (h) node[midway,right] {};
%                 % \node (rb) at ($\scl*(\sclx*1.5,-0.5)$) {$R_X$};
%                 % \node (h') at ($\scl*(\sclx*1.5,-1.2)$) {$H'$};
%                 % \draw[>->]  (rb) -- (h') node [midway,above] {};
%                 % \draw[>->]  (c) -- (h') node [midway,above] {};
%                 % \draw[>->]  (rb) -- (h') node [midway,above] {};
%                 % \draw[>->] (rb) to (r);
%                 % \draw[<-<] (rb) to (k);
%                 \draw[>->] (k') -- (k) ;
%             \end{tikzpicture}
%         }
%         \end{center} 
%         % Hyp ?? $h_{L'R} \mathop{=} Id_{L'L''} \mathop{\star} h_{L''R}$ !!\\



%         We have $h_{XH} \mathop{=} h_{XF} \mathop{\star} h_{FH}$ because $K'C'XL'$ is a pushout and $h_{XH}$ 
%         is the unique morphism such that $H_{C'X} \mathop{\star} h_{XH} \mathop{=} h_{C'C} \mathop{\star} h_{CH} \mathop{=}  h_{C'C''} \mathop{\star} h_{C''C} \mathop{\star} h_{CH}$ and 
%         $h_{L'X} \mathop{\star} h_{XH} \mathop{=} h_{L'R} \mathop{\star} h_{RH} \mathop{=} h_{L'L''} \mathop{\star} h_{L''R} \mathop{\star} h_{RH}$, but we have 
%         \begin{flalign*}
%             &h_{C'X} \mathop{\star} (h_{XF} \mathop{\star} h_{FH}) 
%             \\= & (h_{C'C''} \mathop{\star} h_{C''F}) \mathop{\star} h_{FH} 
%             \\= & h_{C'C''} \mathop{\star} (h_{C''C} \mathop{\star} h_{CH})
%         \end{flalign*}
%         and 
%         \begin{flalign*}
%             &h_{L'X} \mathop{\star} (h_{XF} \mathop{\star} h_{FH})
%             \\ \mathop{=} & (h_{L'L''} \mathop{\star} h_{L''F}) \mathop{\star} h_{FH}
%             \\ \mathop{=} & h_{L'L''} \mathop{\star} (h_{L''R} \mathop{\star} h_{RH})
%         \end{flalign*}
%         This concludes our proof.
%                 \color{black}
\end{proof}


% \begin{lemma}
%     \label{lem:w_g_geq_w_h_leq}
%     Let $\rho \mathop{=} (L \overset{l}{\leftarrowtail} K \overset{r}{\rightarrowtail} R)$ be an injective DPO rewriting rule,
%     \( (X, \emptyset) \) a ruler-graph,
%     and \( G \mathop{\Rightarrow}_{\rho,\mathfrak{M}} H \) a rewriting step. 
%     If $\rho$ is ?? then:\[
%         |Mono(\mathcal{X},G)_{NF}| - |Mono(\mathcal{X},H)_{NF}|
%         \mathop{\geq} 
%         |Mono(X,L)| - |Mono(X,R)|
%     \]
% \end{lemma}
% \begin{proof}
%     \label{proof:lem:w_g_geq_w_h_leq}
%     Let the following diagram be the witness diagram of the rewriting step
%     \begin{center}
%         \begin{tikzpicture}
%             \node (k) at (0,1) {K};
%             \node (l) at (-2,1) {L};
%             \node (r) at (2,1) {R};
%             \node (c) at (0,-1) {C};
%             \node (g) at (-2,-1) {G};
%             \node (h) at (2,-1) {H};
%             \draw[<-<]  (l) -- (k) node [midway,below] {$l$};
%             \draw[>->]  (k) -- (r) node [midway,below] {$r$};
%             \draw[>->] (c) -- (g) node [midway, above] {$l'$};
%             \draw[>->] (c) -- (h) node [midway,above] {$r'$};
%             \draw[>->] (l) -- (g) node[midway, right] {$m$};
%             \draw[>->] (r) -- (h) node[midway, left] {$m'$};
%             \draw[>->] (k) -- (c) node[midway, left] {};
%             \node () [at=($(l)!0.5!(c)$)] {$PO$};
%             \node () [at=($(r)!0.5!(c)$)] {$PO$};
%         \end{tikzpicture}
%     \end{center}
%     We have 
%     \begin{flalign*}
%         &\card{\operatorname{Mono}(X,G)_{NF}} 
%         \mathop{=} & \card{\operatorname{Mono}(X,G)}
%             % \overset{\operatorname{def}}{=} 
%             % w_{T_\Sigma}(t_G) 
%         \\
%         % =& \Psi_{G}\\
%         =& 
%         \card{\operatorname{Mono}(X,G,l')
%         \uplus
%          \operatorname{Mono}(X,G,\lnot l',m)
%         \uplus
%          \operatorname{Mono}(X,G,\lnot l',\lnot m)}\\
%         =& 
%         \card{\operatorname{Mono}(X,G,l')}
%         +
%         \textcolor{red}{\card{\operatorname{Mono}(X,G,\lnot l',m)}}
%         +
%         \card{\operatorname{Mono}(X,G,\lnot l',\lnot m)}
%         \\
%     =& \card{\operatorname{Mono}(X,C)}
%             +
%             \card{\operatorname{Mono}(X, L, \lnot l)} 
%             +
%             \card{\operatorname{Mono}(X, G,\lnot l', \lnot m)} 
%         &  \text{by Lemma~\ref{lem:decomp_w_u}}
%     \end{flalign*}
%     and
%     \begin{flalign*}
%         % ligne 1
%         &\card{\operatorname{Mono}(X,H)_{NF}}
%         \mathop{=} &\card{\operatorname{Mono}(X,H)}
%         % \overset{\operatorname{def}}{=}  
%         %     w_{T_\Sigma}(t_H )
%         \\
%         =& 
%         \card{\operatorname{Mono}(X,H,r')
%         \uplus
%             \operatorname{Mono}(X,H,\lnot r',m')
%         \uplus
%             \operatorname{Mono}(X,H,\lnot r',\lnot m')}
%         \\
%         =& 
%         \card{\operatorname{Mono}(X,H,r')}
%         +
%         \textcolor{red}{\card{\operatorname{Mono}(X,H,\lnot r',m')}}
%         +
%         \card{\operatorname{Mono}(X,H,\lnot r',\lnot m')}
%         \\
%         % \mathop{=} & \Psi_{H} \\
%         % ligne 2
%         \mathop{=} &
%             \card{\operatorname{Mono}(X,C)}
%             +
%             \card{\operatorname{Mono}(X, R, \lnot r)} 
%             +
%             \card{\operatorname{Mono}(X, H, \lnot r', \lnot m')}
%         % ( 
%         %             w_{T_\Sigma} (h_{RH} \mathop{\star} t_H )\mathop{+}w_{T_\Sigma} (h_{CH} \mathop{\star} t_H  -h_{KC}) 
%         % )
%         %        \mathop{+} w_{T_\Sigma}(t_H  - \{h_{RH},h_{CH} \})  
%         &  \text{by Lemma~\ref{lem:decomp_w_u}}
%     \end{flalign*}
%     Therefore, the following inequality holds
%      \begin{flalign*}
%            &\card{\operatorname{Mono}(X,G)} - \card{\operatorname{Mono}(X,H)} &\\
%            =&(\card{\operatorname{Mono}(X, L, \lnot l)} - \card{\operatorname{Mono}(X, R, \lnot r)})\mathop{+}
%            \textcolor{red}{(\card{\operatorname{Mono}(X, G, \lnot l', \lnot m)} - \card{\operatorname{Mono}(X, H, \lnot r', \lnot m')})}  \\
%        \geq& \textcolor{red}{\card{\operatorname{Mono}(X, L, \lnot l)} - \card{\operatorname{Mono}(X, R, \lnot r)}} \hspace{3cm} \text{by Lemma~\ref{lem:w_u_l_not_geq_r_not}} \\
%           =& \card{\operatorname{Mono}(X, L)} - \card{\operatorname{Mono}(X, R)} \hspace{3cm} \text{by Lemma~\ref{lem:xlnlmxrnr}}
%      \end{flalign*}
% \end{proof}

 
% \begin{lemma}
%     \label{lem:w_g_geq_w_h_leq_with_forbidden_contexts}
%     Let $\rho \mathop{=} (L \overset{l}{\leftarrowtail} K \overset{r}{\rightarrowtail} R)$ be an injective DPO rewriting rule,
%     \( (X, F_X) \) a ruler-graph,
%     and \( G \mathop{\Rightarrow}_{\rho,\mathfrak{M}} H \) a rewriting step. 
%     If $\rho$ is ?? then:\[
%         |Mono(\mathcal{X},G)_{NF}| - |Mono(\mathcal{X},H)_{NF}|
%         \mathop{\geq} 
%         |Mono(X,L)| - |Mono(X,R)|
%     \]
% \end{lemma}
% \begin{proof}
%     \label{proof:lem:w_g_geq_w_h_leq_with_forbidden_contexts}
%     Let the following diagram be the witness diagram of the rewriting step
%     \begin{center}
%         \begin{tikzpicture}
%             \node (k) at (0,1) {K};
%             \node (l) at (-2,1) {L};
%             \node (r) at (2,1) {R};
%             \node (c) at (0,-1) {C};
%             \node (g) at (-2,-1) {G};
%             \node (h) at (2,-1) {H};
%             \draw[<-<]  (l) -- (k) node [midway,below] {$l$};
%             \draw[>->]  (k) -- (r) node [midway,below] {$r$};
%             \draw[>->] (c) -- (g) node [midway, above] {$l'$};
%             \draw[>->] (c) -- (h) node [midway,above] {$r'$};
%             \draw[>->] (l) -- (g) node[midway, right] {$m$};
%             \draw[>->] (r) -- (h) node[midway, left] {$m'$};
%             \draw[>->] (k) -- (c) node[midway, left] {};
%             \node () [at=($(l)!0.5!(c)$)] {$PO$};
%             \node () [at=($(r)!0.5!(c)$)] {$PO$};
%         \end{tikzpicture}
%     \end{center}
%     We have 
%     \begin{flalign*}
%         &\card{\operatorname{Mono}(X,G)_{NF}} 
%         \mathop{=} & \card{\operatorname{Mono}(X,G)}
%             % \overset{\operatorname{def}}{=} 
%             % w_{T_\Sigma}(t_G) 
%         \\
%         % =& \Psi_{G}\\
%         =& 
%         \card{\operatorname{Mono}(X,G,l')
%         \uplus
%          \operatorname{Mono}(X,G,\lnot l',m)
%         \uplus
%          \operatorname{Mono}(X,G,\lnot l',\lnot m)}\\
%         =& 
%         \card{\operatorname{Mono}(X,G,l')}
%         +
%         \textcolor{red}{\card{\operatorname{Mono}(X,G,\lnot l',m)}}
%         +
%         \card{\operatorname{Mono}(X,G,\lnot l',\lnot m)}
%         \\
%     =& \card{\operatorname{Mono}(X,C)}
%             +
%             \card{\operatorname{Mono}(X, L, \lnot l)} 
%             +
%             \card{\operatorname{Mono}(X, G,\lnot l', \lnot m)} 
%         &  \text{by Lemma~\ref{lem:decomp_w_u}}
%     \end{flalign*}
%     and
%     \begin{flalign*}
%         % ligne 1
%         &\card{\operatorname{Mono}(X,H)_{NF}}
%         \mathop{=} &\card{\operatorname{Mono}(X,H)}
%         % \overset{\operatorname{def}}{=}  
%         %     w_{T_\Sigma}(t_H )
%         \\
%         =& 
%         \card{\operatorname{Mono}(X,H,r')
%         \uplus
%             \operatorname{Mono}(X,H,\lnot r',m')
%         \uplus
%             \operatorname{Mono}(X,H,\lnot r',\lnot m')}
%         \\
%         =& 
%         \card{\operatorname{Mono}(X,H,r')}
%         +
%         \textcolor{red}{\card{\operatorname{Mono}(X,H,\lnot r',m')}}
%         +
%         \card{\operatorname{Mono}(X,H,\lnot r',\lnot m')}
%         \\
%         % \mathop{=} & \Psi_{H} \\
%         % ligne 2
%         \mathop{=} &
%             \card{\operatorname{Mono}(X,C)}
%             +
%             \card{\operatorname{Mono}(X, R, \lnot r)} 
%             +
%             \card{\operatorname{Mono}(X, H, \lnot r', \lnot m')}
%         % ( 
%         %             w_{T_\Sigma} (h_{RH} \mathop{\star} t_H )\mathop{+}w_{T_\Sigma} (h_{CH} \mathop{\star} t_H  -h_{KC}) 
%         % )
%         %        \mathop{+} w_{T_\Sigma}(t_H  - \{h_{RH},h_{CH} \})  
%         &  \text{by Lemma~\ref{lem:decomp_w_u}}
%     \end{flalign*}
%     Therefore, the following inequality holds
%      \begin{flalign*}
%            &\card{\operatorname{Mono}(X,G)} - \card{\operatorname{Mono}(X,H)} &\\
%            =&(\card{\operatorname{Mono}(X, L, \lnot l)} - \card{\operatorname{Mono}(X, R, \lnot r)})\mathop{+}
%            \textcolor{red}{(\card{\operatorname{Mono}(X, G, \lnot l', \lnot m)} - \card{\operatorname{Mono}(X, H, \lnot r', \lnot m')})}  \\
%        \geq& \textcolor{red}{\card{\operatorname{Mono}(X, L, \lnot l)} - \card{\operatorname{Mono}(X, R, \lnot r)}} \hspace{3cm} \text{by Lemma~\ref{lem:w_u_l_not_geq_r_not}} \\
%           =& \card{\operatorname{Mono}(X, L)} - \card{\operatorname{Mono}(X, R)} \hspace{3cm} \text{by Lemma~\ref{lem:xlnlmxrnr}}
%      \end{flalign*}
%     \qed
% \end{proof}



\noindent
\begin*{Theorem~\textbf{\ref{antipattern:thm:termination_grs}}}(Termination)

 Let \(\mathcal{A}\) and \(\mathcal{B}\) be sets of injective DPO rewriting rules, $\mathbb{X}$ a set of ruler-graphs, and $s_\mathbb{X}$ a weight function. If the following hold:
    \begin{enumerate}
        \item  for every $\rho \mathop{\in} \mathcal{A} \mathop{\cup} \mathcal{B}$ and for every \( \mathcal{X} \mathop{=} (X,\mathcal{C})\mathop{\in} \mathbb{X} \), 
        % the number of $X$-occurrences that are not included in any occurrences of the forbidden context $F \mathop{\in} F_x$ in a $\rho$-rewriting step is predictable,
        $\rho$ is $X$-non-increasing, and if $\mathcal{C}= \set{f:X \rightarrowtail F}$ then $\rho^{-1}$ is $X$-non-increasing and $F$-non-increasing, and
        % $\rho^{-1}$ is $F$-non-increasing if $\mathcal{X}= (X,f:X \rightarrowtail F)$ and, $\rho$ and $\rho^{-1}$ are $X$-non-increasing
        \item for every \(\rho \mathop{\in} \mathcal{A}\), we have
        % \( w_{s_\mathbb{X}}(lhs(\rho)) \mathop{>} w_{s_\mathbb{X}}(rhs(\rho)) \),
        $ \sum_{\mathcal{X} \mathop{\in} \mathbb{X}}^{}s_\mathbb{X}(\mathcal{X}) \mathop{*} 
            \Lambda(\mathcal{X},\rho) \mathop{>} 0 $, and
        \item for every \(\rho \mathop{\in} \mathcal{B}\), we have   
        % \( w_{s_\mathbb{X}}(lhs(\rho)) \mathop{\geq} w_{s_\mathbb{X}}(rhs(\rho)) \).
        $ 
            \sum_{\mathcal{X} \mathop{\in} \mathbb{X}}^{}s_\mathbb{X}(\mathcal{X}) \mathop{*} \Lambda(\mathcal{X},\rho) \mathop{\geq} 0 
        $.
    \end{enumerate}
    Then \(\mathop{\Rightarrow}_{\mathcal{A},\mathcal{M}}\) terminates relative to \(\mathop{\Rightarrow}_{\mathcal{B},\mathcal{M}}\).
\end*{}


% Let \(\mathcal{A}\) and \(\mathcal{B}\) be sets of injective DPO rewriting rules, $\mathbb{X}$ a set of ruler-graphs and $s_\mathbb{X}$ a weight function. If the following hold:
%     \begin{enumerate}
%         \item  for every $\rho \mathop{\in} \mathcal{A} \mathop{\cup} \mathcal{B}$ and for every \( \mathcal{X} \mathop{\in} \mathbb{X} \) with underlying graph $X$, 
%         % the number of $X$-occurrences that are not included in any occurrences of the forbidden context $F \mathop{\in} F_x$ in a $\rho$-rewriting step is predictable,
%         $\rho$ is $X$-non-increasing, and if $\mathcal{X}= (X,f \mathop{\colon} X \rightarrowtail F)$ then $\rho^{-1}$ is $X$-non-increasing and $F$-non-increasing,
%         % $\rho^{-1}$ is $F$-non-increasing if $\mathcal{X}= (X,f \mathop{\colon} X \rightarrowtail F)$ and, $\rho$ and $\rho^{-1}$ are $X$-non-increasing
%         \item for every \(\rho \mathop{\in} \mathcal{A}\), we have
%         % \( w_{s_\mathbb{X}}(lhs(\rho)) \mathop{>} w_{s_\mathbb{X}}(rhs(\rho)) \),
%         $ \sum_{\mathcal{X} \mathop{\in} \mathbb{X}}^{}s_\mathbb{X}(X) * 
%             \Lambda(\mathcal{X},\rho) \mathop{>} 0 $
%         \item for every \(\rho \mathop{\in} \mathcal{B}\), we have   
%         % \( w_{s_\mathbb{X}}(lhs(\rho)) \mathop{\geq} w_{s_\mathbb{X}}(rhs(\rho)) \).
%         $ 
%             \sum_{\mathcal{X} \mathop{\in} \mathbb{X}}^{}s_\mathbb{X}(X) * \Lambda(\mathcal{X},\rho) \mathop{\geq} 0 
%         $
%     \end{enumerate}
%     Then \(\mathop{\Rightarrow}_{\mathcal{A},\mathcal{M}}\) is relatively terminating to \(\mathop{\Rightarrow}_{\mathcal{B},\mathcal{M}}\).

\begin{proof} 
    \label{antipattern:proof:thm:termination_grs}
    
    By the assumption 1 and Lemma~\ref{antipattern:lem:w_g_geq_w_h_leq}, the following inequality holds for all rewriting steps $G \mathop{\Rightarrow}_{\rho, \mathcal{M}} H$  with $\rho \mathop{\in} \mathcal{A} \mathop{\cup} \mathcal{B}$:
      \[
        w_{s_\mathbb{X}}(G) - w_{s_\mathbb{X}}(H) 
        \mathop{\geq} 
        \sum_{\mathcal{X} \mathop{\in} \mathbb{X}}^{}s_\mathbb{X}(\mathcal{X}) * \left( 
            \card{\operatorname{Mono}(X,L)} -
            \card{\operatorname{Mono}(X,R)}
            \right).
    \]
    \noindent From the assumptions 2 and 3, we deduce 
    \begin{itemize}
        \item \( w_{s_\mathbb{X}}(G) \mathop{>} w_{s_\mathbb{X}}(H) \) for every rule \(\rho \mathop{\in} \mathcal{A}\);
        \item  \( w_{s_\mathbb{X}}(G) \mathop{\geq} w_{s_\mathbb{X}}(H) \) for every rule \(\rho \mathop{\in} \mathcal{B}\).
    \end{itemize}
    Since $w_{s_\mathbb{X}}(G) \mathop{\in} \mathbb{N}$ for all graph $G$. Every rewriting chain can only have a finite number of rewriting steps using rules from $\mathcal{A}$.
\end{proof} 

       


 
% %%%%%%%%%%%%%%%%%%%%%%%%%%%%%%%%%%%%%%%%%%%%%%%%%%%%%%%%%%%%
% %%%%   Chapter : LyonParallel — A Termination Tool for DPO Graph Rewriting
% %%%%%%%%%%%%%%%%%%%%%%%%%%%%%%%%%%%%%%%%%%%%%%%%%%%%%%%%%%%%
\chapter{LyonParallel\textemdash A Tool for Relative Termination DPO Graph Rewriting} 
\label{chap:lyonparallel}
LyonParallel is a prototype tool, written in OCaml, for the analysis of DPO graph rewriting systems.
It implements the type graph method presented by Endrullis and Overbeek~\cite{endrullis2024generalized_icgt} and its extension in Chapter~\ref{chap:nwf}, and it also implements the morphism counting method in Chapter~\ref{chap:subgraph_counting} and its extension in Chapter~\ref{chap:antipattern}. 
For DPO graph rewriting systems with monic matches,
the type graph method presented by Endrullis and Overbeek~\cite{endrullis2024generalized_icgt} is more powerful than the versions presented by Zantema et al. and Bruggink et al.~\cite{zantema2014termination,bruggink2014termination,bruggink2015proving}, which are implemented in TORPAcyc~\cite{TORPAcyc} and Grez~\cite{grez}. At the time of writing and to the best of our knowledge, the method presented by Overbeek and Endrullis has no publicly available implementation.

LyonParallel implements the notions of DPO rewriting systems on finite directed, edge-labeled multigraphs introduced in Chapter~\ref{chap:preliminaries}. Using the \texttt{REPL} (Read-Eval-Print Loop)\textemdash{}an interactive environment that reads user commands, evaluates them, prints results, and waits for the next input\textemdash{}users can select a set of rules encoded in the tool, then choose a method to partition the rules of the system into two sets $\mathcal{A}$ and $\mathcal{B}$ such that $\mathcal{A}$ terminates relative to $\mathcal{B}$. If such a partition is found, the tool replaces the current rule set with \(\mathcal{B}\) and a new iteration begins; otherwise, users may try a different method. When the rule set becomes empty, the termination of the original system is proved.
\section{DPO graph rewriting systems}
\label{lyonparallel:sec:implementation_of_dpo_graph_rewriting_systems}
\subsection{Labeled graphs}
\label{lyonparallel:sec:labeled_graphs}
A labeled graph is encoded in the tool as a structure with five fields: a list of nodes, a list of edges, a mapping that associates to each edge a source node, a mapping that associates to each edge a target node, and a mapping that associates to each edge a label.
The representation of a labeled graph is constructed using the function \colorbox{Ivory2}{$MGraph.fromList$} given 2 parameters. The first parameter is a list of integers representing the nodes of the graph. The second parameter is a list of four-tuples representing the labeled edges of the graph, where each tuple \colorbox{Ivory2}{$(s,lab,t,id)$} of type \colorbox{Ivory2}{$\textit{int}\times \textit{string}\times \textit{int}\times \textit{int}$} represents an edge with source node $s$, label $lab$, target node $t$, and identifier $id$.

For example, consider the graph rewriting rule:
% \begin{figure}[H]
%       \centering 

      \begin{center}
      \begin{tikzpicture}
          \graphbox{$L$}{0mm}{0mm}{34mm}{15mm}{2mm}{-5mm}{
              \coordinate (o) at (0mm,-3mm); 
              \node[draw,circle] (l1) at ($(o)+(-10mm,0mm)$) {1};
              \node[draw,circle] (l2) at ($(l1)+(2,0)$) {2};
              \node[draw,circle] (l3) at ($(l1)+(1,0)$) {3};
              \draw[->] (l1) -- (l3) node[midway,above] {$a$};
              \draw[->] (l3) -- (l2) node[midway,above] {$a$};
          }     
          \graphbox{$K$}{40mm}{0mm}{24mm}{15mm}{2mm}{-5mm}{
              \coordinate (o) at (5mm,-3mm); 
              \node[draw,circle] (l1) at ($(o)+(-10mm,0mm)$) {1};
              \node[draw,circle] (l2) at ($(l1)+(1,0)$) {2};
              % \node[draw,circle] (l3) at ($(l1)+(1,0)$) {$\ $};
              % \draw[->] (l1) -- (l3) node[midway,above] {$a$};
              % \draw[->] (l3) -- (l2) node[midway,above] {$a$};
          }    
          \graphbox{$R$}{70mm}{0mm}{45mm}{15mm}{2mm}{-5mm}{
              \coordinate (o) at (-5mm,-3mm); 
              \node[draw,circle] (l1) at ($(o)+(-10mm,0mm)$) {1};
              \node[draw,circle] (l2) at ($(l1)+(3,0)$) {2};
              \node[draw,circle] (l3) at ($(l1)+(1,0)$) {4};
              \node[draw,circle] (l4) at ($(l1)+(2,0)$) {5};
              \draw[->] (l1) -- (l3) node[midway,above] {$a$};
              \draw[->] (l3) -- (l4) node[midway,above] {$b$};
              \draw[->] (l4) -- (l2) node[midway,above] {$a$};
          }    
          \node () at (37mm,-8mm) {$\overset{l}{\leftarrowtail}$};
          \node () at (67mm,-8mm) {$\overset{r}{\rightarrowtail}$};
          % \draw[>->] (51mm,2mm) -- (52mm,3mm);
      \end{tikzpicture}
    \end{center}
%       \caption{}
%       \label{fig:implemented:grsaa_0}
%   \end{figure}
\noindent The representation of graphs $K$,$L$, and $R$ in OCaml can be constructed by the following code:
\begin{minted}{ocaml}
    let graph_K = MGraph.fromList [1;2] []
    let graph_L = Mgraph.fromList [1;2;3] 
                    [(1,"a",3,1);(3,"a",2,2)]
    let graph_R = Mgraph.fromList [1;2;4;5] 
                    [(1,"a",4,1);(4,"b",5,2);(5,"a",2,3)]
\end{minted}

\subsection{Homomorphisms}
\label{lyonparallel:sec:homomorphisms}
A homomorphism between two labeled graphs is encoded as a mapping from the nodes of the source graph to the nodes of the target graph, which preserves the edge labels and the source and target nodes of each edge. Its representation can be constructed using the function
\colorbox{Ivory2}{$GraphHomomorphism.fromList$}, given 4 arguments:
    \begin{enumerate}
        \item the first argument is the source graph,
        \item the second argument is the target graph,
        \item the third argument is a list of pairs, where each pair $(s,t)$ means that node $s$ of the source graph is mapped to node $t$ of the target graph,
        \item the fourth argument is a list of pairs similar to the third argument, but for the edges.
    \end{enumerate}
For example, consider morphisms $l$ and $r$ of the rewriting rule:
  \begin{center}
      \begin{tikzpicture}
          \graphbox{$L$}{0mm}{0mm}{34mm}{15mm}{2mm}{-5mm}{
              \coordinate (o) at (0mm,-3mm); 
              \node[draw,circle] (l1) at ($(o)+(-10mm,0mm)$) {1};
              \node[draw,circle] (l2) at ($(l1)+(2,0)$) {2};
              \node[draw,circle] (l3) at ($(l1)+(1,0)$) {3};
              \draw[->] (l1) -- (l3) node[midway,above] {$a$};
              \draw[->] (l3) -- (l2) node[midway,above] {$a$};
          }     
          \graphbox{$K$}{40mm}{0mm}{24mm}{15mm}{2mm}{-5mm}{
              \coordinate (o) at (5mm,-3mm); 
              \node[draw,circle] (l1) at ($(o)+(-10mm,0mm)$) {1};
              \node[draw,circle] (l2) at ($(l1)+(1,0)$) {2};
              % \node[draw,circle] (l3) at ($(l1)+(1,0)$) {$\ $};
              % \draw[->] (l1) -- (l3) node[midway,above] {$a$};
              % \draw[->] (l3) -- (l2) node[midway,above] {$a$};
          }    
          \graphbox{$R$}{70mm}{0mm}{45mm}{15mm}{2mm}{-5mm}{
              \coordinate (o) at (-5mm,-3mm); 
              \node[draw,circle] (l1) at ($(o)+(-10mm,0mm)$) {1};
              \node[draw,circle] (l2) at ($(l1)+(3,0)$) {2};
              \node[draw,circle] (l3) at ($(l1)+(1,0)$) {4};
              \node[draw,circle] (l4) at ($(l1)+(2,0)$) {5};
              \draw[->] (l1) -- (l3) node[midway,above] {$a$};
              \draw[->] (l3) -- (l4) node[midway,above] {$b$};
              \draw[->] (l4) -- (l2) node[midway,above] {$a$};
          }    
          \node () at (37mm,-8mm) {$\overset{l}{\leftarrowtail}$};
          \node () at (67mm,-8mm) {$\overset{r}{\rightarrowtail}$};
          % \draw[>->] (51mm,2mm) -- (52mm,3mm);
      \end{tikzpicture}
    \end{center}
\noindent Their representation in OCaml can be constructed using the following code, where \colorbox{Ivory2}{$morphism\_l$} represents $l$ and \colorbox{Ivory2}{$morphism\_r$} represents $r$, respectively, and \colorbox{Ivory2}{$graph\_K$}, \colorbox{Ivory2}{$graph\_L$}, and \colorbox{Ivory2}{$graph\_R$} represent graphs $K$, $L$, and $R$, respectively (see Section~\ref{lyonparallel:sec:labeled_graphs}):
\begin{minted}{ocaml}
    let morphism_l = GraphHomomorphism.fromList 
                        graph_K graph_L [(1,1);(2,2)] []
    let morphism_r = GraphHomomorphism.fromList
                        graph_K graph_R [(1,1);(2,2)] []
\end{minted}

\subsection{DPO rewriting rules and systems}
A DPO rewriting rule is represented by a structure with two fields: a left-hand-side homomorphism and a right-hand side homomorphism. The representation of a rewriting rule can be constructed using the function
 $$\colorbox{Ivory2}{$GraphRewritingSystem.DPOrule.fromHomos$}$$
given the representations of the left- and right-hand side homomorphisms.
For example, 
consider the rewriting rule:
  \begin{center}
      \begin{tikzpicture}
          \graphbox{$L$}{0mm}{0mm}{34mm}{15mm}{2mm}{-5mm}{
              \coordinate (o) at (0mm,-3mm); 
              \node[draw,circle] (l1) at ($(o)+(-10mm,0mm)$) {1};
              \node[draw,circle] (l2) at ($(l1)+(2,0)$) {2};
              \node[draw,circle] (l3) at ($(l1)+(1,0)$) {3};
              \draw[->] (l1) -- (l3) node[midway,above] {$a$};
              \draw[->] (l3) -- (l2) node[midway,above] {$a$};
          }     
          \graphbox{$K$}{40mm}{0mm}{24mm}{15mm}{2mm}{-5mm}{
              \coordinate (o) at (5mm,-3mm); 
              \node[draw,circle] (l1) at ($(o)+(-10mm,0mm)$) {1};
              \node[draw,circle] (l2) at ($(l1)+(1,0)$) {2};
              % \node[draw,circle] (l3) at ($(l1)+(1,0)$) {$\ $};
              % \draw[->] (l1) -- (l3) node[midway,above] {$a$};
              % \draw[->] (l3) -- (l2) node[midway,above] {$a$};
          }    
          \graphbox{$R$}{70mm}{0mm}{45mm}{15mm}{2mm}{-5mm}{
              \coordinate (o) at (-5mm,-3mm); 
              \node[draw,circle] (l1) at ($(o)+(-10mm,0mm)$) {1};
              \node[draw,circle] (l2) at ($(l1)+(3,0)$) {2};
              \node[draw,circle] (l3) at ($(l1)+(1,0)$) {4};
              \node[draw,circle] (l4) at ($(l1)+(2,0)$) {5};
              \draw[->] (l1) -- (l3) node[midway,above] {$a$};
              \draw[->] (l3) -- (l4) node[midway,above] {$b$};
              \draw[->] (l4) -- (l2) node[midway,above] {$a$};
          }    
          \node () at (37mm,-8mm) {$\overset{l}{\leftarrowtail}$};
          \node () at (67mm,-8mm) {$\overset{r}{\rightarrowtail}$};
          % \draw[>->] (51mm,2mm) -- (52mm,3mm);
      \end{tikzpicture}
    \end{center}
\noindent The representation of the rewriting rule is constructed using the following OCaml expression, which creates the representation of a DPO rule from the representations of the left and right morphisms. Here, \colorbox{Ivory2}{$morphism\_l$} and \colorbox{Ivory2}{$morphism\_r$} denote the representations of $l$ and $r$, respectively (see Section~\ref{lyonparallel:sec:homomorphisms}):
\begin{minted}{ocaml}
    let rl = GraphRewritingSystem.DPOrule.fromHomos 
                morphism_l morphism_r
\end{minted}

A DPO graph rewriting system is represented by a structure with three fields: the set of rules, the name of the system, and a Boolean indicating whether the system is restricted to monic matches. A representation of the DPO graph rewriting system
can be constructed using the function 
$$\colorbox{Ivory2}{$GraphRewritingSystem.fromRulesListAndName$}$$ given three arguments: a list of rewriting rules and the name of the system; the third parameter is \colorbox{Ivory2}{$false$} by default.

For example, consider the DPO graph rewriting system that consists of the rule:
  \begin{center}
      \begin{tikzpicture}
          \graphbox{$L$}{0mm}{0mm}{34mm}{15mm}{2mm}{-5mm}{
              \coordinate (o) at (0mm,-3mm); 
              \node[draw,circle] (l1) at ($(o)+(-10mm,0mm)$) {1};
              \node[draw,circle] (l2) at ($(l1)+(2,0)$) {2};
              \node[draw,circle] (l3) at ($(l1)+(1,0)$) {3};
              \draw[->] (l1) -- (l3) node[midway,above] {$a$};
              \draw[->] (l3) -- (l2) node[midway,above] {$a$};
          }     
          \graphbox{$K$}{40mm}{0mm}{24mm}{15mm}{2mm}{-5mm}{
              \coordinate (o) at (5mm,-3mm); 
              \node[draw,circle] (l1) at ($(o)+(-10mm,0mm)$) {1};
              \node[draw,circle] (l2) at ($(l1)+(1,0)$) {2};
              % \node[draw,circle] (l3) at ($(l1)+(1,0)$) {$\ $};
              % \draw[->] (l1) -- (l3) node[midway,above] {$a$};
              % \draw[->] (l3) -- (l2) node[midway,above] {$a$};
          }    
          \graphbox{$R$}{70mm}{0mm}{45mm}{15mm}{2mm}{-5mm}{
              \coordinate (o) at (-5mm,-3mm); 
              \node[draw,circle] (l1) at ($(o)+(-10mm,0mm)$) {1};
              \node[draw,circle] (l2) at ($(l1)+(3,0)$) {2};
              \node[draw,circle] (l3) at ($(l1)+(1,0)$) {4};
              \node[draw,circle] (l4) at ($(l1)+(2,0)$) {5};
              \draw[->] (l1) -- (l3) node[midway,above] {$a$};
              \draw[->] (l3) -- (l4) node[midway,above] {$b$};
              \draw[->] (l4) -- (l2) node[midway,above] {$a$};
          }    
          \node () at (37mm,-8mm) {$\overset{l}{\leftarrowtail}$};
          \node () at (67mm,-8mm) {$\overset{r}{\rightarrowtail}$};
          % \draw[>->] (51mm,2mm) -- (52mm,3mm);
      \end{tikzpicture}
    \end{center}
It is represented by the following OCaml expression that creates a system from the singleton list containing the representation of the rule. Here, \colorbox{Ivory2}{$rl$} denotes the representation of the rewriting rule.
\begin{minted}{ocaml}
    let bruggink14_ex_4_and_6 =
        GraphRewritingSystem.fromRulesListAndName 
            [rl] "bruggink14_ex_4_and_6"
\end{minted}
The function \colorbox{Ivory2}{$GraphRewritingSystem.fromRulesListAndName$} produces a system named \colorbox{Ivory2}{$bruggink14\_ex\_4\_and\_6$} from the singleton list \colorbox{Ivory2}{$[rl]$}.

\subsection{User-defined rewriting systems}
User-defined rewriting systems should be placed in the file 
    $$\colorbox{Ivory2}{$lib/concretGraphRewritingSystems.ml$}$$
 To expose them, add each instance to the list 
 $$\colorbox{Ivory2}{$available\_graph\_rewriting\_systems$}$$ in the same file.

For example, for the rewriting rule named \colorbox{Ivory2}{$bruggink14\_ex\_4\_and\_6$}, the following code should be added to the end of the file:
\begin{minted}{ocaml}
    let available_graph_rewriting_systems = 
        bruggink14_ex_4_and_6 :: available_graph_rewriting_systems 
\end{minted}

\section{Ruler-Graphs}
A ruler-graph with one forbidden context as defined in Chapter~\ref{chap:antipattern} is represented by a structure with an underlying graph $X$ and an injective graph homomorphism from $X$. 
Consider a ruler-graph $(X, f)$ where $X$ is the graph \raisebox{2pt}{
            \scalebox{0.7}{\tikz[baseline=-0.5ex]{
            \node [draw,circle] (z) at (-1,0) {};
            \node [draw,circle] (x) at (0,0) {};
            \node[draw,circle] (y) at (1,0) {};
            \draw[->] (z)--(x) node[midway, above] {$a$};
            \draw[->] (x)--(y) node[midway, above] {$a$};
        }}} 
and the injective graph homomorphism is shown below:
%  in Figure~\ref{fig:intro:graph_transformation_rule_anti_patternsfs}.
%  \begin{figure}[H]
%     \centering
\begin{center}
\begin{tikzpicture}
      \graphbox{$L$}{30mm}{0mm}{34mm}{20mm}{2mm}{-5mm}{
          \coordinate (o) at (0mm,-3mm); 
          \node[draw,circle] (l1) at ($(o)+(-10mm,0mm)$) {1};
          \node[draw,circle] (l2) at ($(l1)+(2,0)$) {2};
          \node[draw,circle] (l3) at ($(l1)+(1,0)$) {3};
          \draw[->] (l1) -- (l3) node[midway,above] {$a$};
          \draw[->] (l3) -- (l2) node[midway,above] {$a$};
      }        
      \graphbox{$F$}{70mm}{0mm}{45mm}{20mm}{2mm}{-5mm}{
        \coordinate (o) at (0mm,-3mm); 
        \node[draw,circle] (l1) at ($(o)+(-10mm,0mm)$) {1};
        \node[draw,circle] (l2) at ($(l1)+(2,0)$) {2};
        \node[draw,circle] (l3) at ($(l1)+(1,0)$) {3};
        \draw[->] (l1) -- (l3) node[midway,above] {$a$};
        \draw[->] (l3) -- (l2) node[midway,above] {$a$};
        \draw[->] (l3) edge [loop below] node {$c$} (l3);
      }    
    %   \node () at (37mm,-10mm) {$\leftarrowtail$};
      \node () at (67mm,-10mm) {$\rightarrowtail$};
  \end{tikzpicture}
%   \caption{}
%   \label{fig:intro:graph_transformation_rule_anti_patternsfs}
%  \end{figure} 
\end{center}
Its representation in OCaml
        can be constructed using the following code:
\begin{figure}[H]
\begin{minted}{ocaml} 
    let graph_X = MGraph.fromList 
                    [1;2;3] [(1,"a",3,1);(3,"a",2,2)]
    let ruler_graph_aa_notin_aca = 
        let graph_F = MGraph.fromList 
                [1;2;3] [(1,"a",3,1);(3,"c",3,3);(3,"a",2,2)] in
        let f = Homo.fromList graph_X graph_F 
                    [(1,1);(2,2);(3,3)] [(1,1);(2,2)] in
        {x: graph_X; fx = Some f}
\end{minted}
    \caption{}
    \label{fig:implemented:ruler_graph_representation}
\end{figure}

Instances of ruler-graphs can be defined in the file \colorbox{Ivory2}{$lib/ruler\_graph.ml$} of \textbf{LyonParallel}. To expose these instances to the REPL, they should be added to the list \colorbox{Ivory2}{$ruler\_graphs$} in the same file. For example, for the instance of ruler-graph \colorbox{Ivory2}{$ruler\_graph\_aa\_notin\_aca$} shown in Figure~\ref{fig:implemented:ruler_graph_representation}, the following code should be added to the end of the file:
\begin{minted}{ocaml}
    let ruler_graphs = ruler_graph_aa_notin_aca :: ruler_graphs 
\end{minted}

% As another example, the representation of graph \raisebox{2pt}{
%             \scalebox{0.7}{\tikz[baseline=-0.5ex]{
%             \node [draw,circle] (z) at (-1,0) {};
%             \node [draw,circle] (x) at (0,0) {};
%             \node[draw,circle] (y) at (1,0) {};
%             \draw[->] (z)--(x) node[midway, above] {$a$};
%             \draw[->] (x) edge[loop above] node {$c$} (x);
%             \draw[->] (x)--(y) node[midway, above] {$a$};
%         }}}
%    is constructed using the following code:
% \begin{minted}{ocaml}
%     let y = MGraph.fromList [1;2;3] [(1,"a",3,1);(3,"c",3,3);(3,"a",2,2)]
% \end{minted}

% For example, the ruler-graph with underlying graph \raisebox{2pt}{
%             \scalebox{0.7}{\tikz[baseline=-0.5ex]{
%             \node [draw,circle] (z) at (-1,0) {};
%             \node [draw,circle] (x) at (0,0) {};
%             \node[draw,circle] (y) at (1,0) {};
%             \draw[->] (z)--(x) node[midway, above] {$a$};
%             \draw[->] (x)--(y) node[midway, above] {$a$};
%         }}} and forbidden context \raisebox{2pt}{
%             \scalebox{0.7}{\tikz[baseline=-0.5ex]{
%             \node [draw,circle] (z) at (-1,0) {};
%             \node [draw,circle] (x) at (0,0) {};
%             \node[draw,circle] (y) at (1,0) {};
%             \draw[->] (z)--(x) node[midway, above] {$a$};
%             \draw[->] (x) edge[loop above] node {$c$} (x);
%             \draw[->] (x)--(y) node[midway, above] {$a$};
%         }}} is encoded as the following ocaml structure in the file \mintinline{latex}{Parallel/lib/ruler_graph.ml} of \textbf{LyonParallel}:
% \begin{figure}[H]
%     \begin{minted}{ocaml}
%     let aa_not_in_aca =   
%         let x = MGraph.fromList [1;2;3] [(1,"a",3,1);(3,"a",2,2)] in
%         let f = Homo.fromList 
%             [1;2;3] [(1,"a",3,1);(3,"a",2,2)]
%             [1;2;3] [(1,"a",3,1);(3,"c",3,3);(3,"a",2,2)]
%             [(1,1);(2,2);(3,3)] [(1,1);(2,2)]  in
%         {x; fx = Some f}
%     \end{minted}
%     \caption{}
%     \label{fig:implemented:aa_not_in_aca}
% \end{figure}
 
\section{Installation and Execution of LyonParallel}
\label{lyonparallel:sec:installation}
This section shows the installation process using \colorbox{Ivory2}{OPAM}\textemdash{}a package manager for OCaml. To install LyonParallel, one needs to have \colorbox{Ivory2}{OPAM} installed on one's system, then 
run the following command
\begin{minted}[style=bw]{bash}
    >> opam switch create lyonParallel 5.2.1
\end{minted}
    % \colorbox{Ivory2}{opam switch create lyonParallel 5.2.1}
 and follow the instructions to create and switch to a new \colorbox{Ivory2}{OPAM} switch. 
 The dependencies of LyonParallel are listed in the file \colorbox{Ivory2}{$LyonParallel.opam$} in the root directory of the project. To install these dependencies, run the following command:
\begin{minted}{bash}
    >> opam install .
\end{minted}
The project can be built and installed by running the following command:
\begin{minted}{bash}
    >> dune build && dune install
\end{minted}
The following command uninstalls the tool:
\begin{minted}{bash}
    >> opam remove LyonParallel
\end{minted}
Finally, the following command launches the interactive REPL of the tool:
\begin{minted}{bash}
    >> LyonParallel
\end{minted}
Upon launching, one should see:
\begin{minted}[style=bw]{bash}
    >> Type "help" for a list of commands.
\end{minted}
\section{Relative Termination Analysis using LyonParallel}
\label{lyonparallel:sec:termination}
The LyonParallel REPL lets the user select a set of rewriting rules and run the tool's relative termination analysis techniques to test whether the set can be partitioned into two subsets \(\mathcal{A}\) and \(\mathcal{B}\) such that \(\mathcal{A}\) terminates relative to \(\mathcal{B}\). If such a partition is found, the REPL replaces the analyzed rule set with \(\mathcal{B}\) for further analysis; otherwise, the user can try alternative methods or different partitions.

The REPL command \colorbox{Ivory2}{$systems$} lists the available DPO graph rewriting systems; user-defined systems can be added as described in~\textsection~\ref{lyonparallel:sec:implementation_of_dpo_graph_rewriting_systems}. The predefined system consisting of the rule shown below is named \colorbox{Ivory2}{$bruggink14\_ex\_4\_and\_6$}:
\begin{center}
      \begin{tikzpicture}
          \graphbox{$L$}{0mm}{0mm}{34mm}{15mm}{2mm}{-5mm}{
              \coordinate (o) at (0mm,-3mm); 
              \node[draw,circle] (l1) at ($(o)+(-10mm,0mm)$) {1};
              \node[draw,circle] (l2) at ($(l1)+(2,0)$) {2};
              \node[draw,circle] (l3) at ($(l1)+(1,0)$) {3};
              \draw[->] (l1) -- (l3) node[midway,above] {$a$};
              \draw[->] (l3) -- (l2) node[midway,above] {$a$};
          }     
          \graphbox{$K$}{40mm}{0mm}{24mm}{15mm}{2mm}{-5mm}{
              \coordinate (o) at (5mm,-3mm); 
              \node[draw,circle] (l1) at ($(o)+(-10mm,0mm)$) {1};
              \node[draw,circle] (l2) at ($(l1)+(1,0)$) {2};
              % \node[draw,circle] (l3) at ($(l1)+(1,0)$) {$\ $};
              % \draw[->] (l1) -- (l3) node[midway,above] {$a$};
              % \draw[->] (l3) -- (l2) node[midway,above] {$a$};
          }    
          \graphbox{$R$}{70mm}{0mm}{45mm}{15mm}{2mm}{-5mm}{
              \coordinate (o) at (-5mm,-3mm); 
              \node[draw,circle] (l1) at ($(o)+(-10mm,0mm)$) {1};
              \node[draw,circle] (l2) at ($(l1)+(3,0)$) {2};
              \node[draw,circle] (l3) at ($(l1)+(1,0)$) {4};
              \node[draw,circle] (l4) at ($(l1)+(2,0)$) {5};
              \draw[->] (l1) -- (l3) node[midway,above] {$a$};
              \draw[->] (l3) -- (l4) node[midway,above] {$b$};
              \draw[->] (l4) -- (l2) node[midway,above] {$a$};
          }    
          \node () at (37mm,-8mm) {$\overset{l}{\leftarrowtail}$};
          \node () at (67mm,-8mm) {$\overset{r}{\rightarrowtail}$};
          % \draw[>->] (51mm,2mm) -- (52mm,3mm);
      \end{tikzpicture}
    \end{center}
It appears in the list. Select it at the REPL by name:
\begin{minted}[style=bw]{bash}
select_system_by_name bruggink14_ex_4_and_6
\end{minted}


\subsection{Termination Analysis using the Morphism Counting Method}
Consider the graph rewriting system presented in Example~\ref{subgraph_counting:ex_contrib_variant} with a single rule:
    \begin{center}
            \begin{tikzpicture}[scale=0.9]
                \graphbox{$L$}{0mm}{0mm}{35mm}{35mm}{2mm}{-5mm}{
                    \coordinate (delta) at (0,-18mm);
                    \node[draw,circle] (l1) at ($(delta)+(-1,1.5)$) {1};
                    \node[draw,circle] (l2) at ($(delta)+(1,1.5)$) {2};
                    \node[draw,circle] (l3) at ($(delta)+(0,0)$) {3};
                    \draw[->] (l1) -- (l3) node[midway,left] {$s$};
                    \draw[->] (l2) -- (l3) node[midway,right] {$s$};
                    \draw[->] (l3) edge [loop below] node {0} (l3);
                }
                \graphbox{$K$}{40mm}{0mm}{35mm}{35mm}{2mm}{-5mm}{
                    \coordinate (delta) at (0,-18mm);
                    \coordinate (interfaceorigin) at ($(delta) +(5,0)$);
                    \node[draw,circle] (r1) at ($(delta) +(-1,1.5)$) {1};
                    \node[draw,circle] (r2) at ($(delta) +(0.5,1.5)$) {2};
                    \node[draw,circle] (r3) at ($(delta)+(0,0)$) {3};
                    % \draw[->] (r1) -- (r3) node[midway,left] {$s$};
                    % \draw[->] (r3) edge [loop below] node {0} (r3);
                }
                \graphbox{$R$}{80mm}{0mm}{50mm}{35mm}{2mm}{-5mm}{
                    \coordinate (delta) at (-10mm,-18mm);
                    \node[draw,circle] (r1) at ($(delta)+(-1,1.5)$) {1};
                    \node[draw,circle] (r2) at ($(delta)+(0.5,1.5)$) {2};
                    \node[draw,circle] (r3) at ($(delta)+(0,0)$) {3};
                    \node[draw,circle] (r4) at ($(delta)+(1,0)$) {4};
                    \draw[->] (r1) -- (r3) node[midway,left] {$s$};
                    \draw[->] (r2) -- (r4) node[midway,right] {$s$};
                    \draw[->] (r4) edge [loop below] node {0} (r4);
                    \draw[->] (r3) edge [loop below] node {0} (r3);
                    \node[draw,circle] (r5) at ($(r2)+(1.5,0)$) {};
                    \draw[->] (r5) edge [loop below] node {0} (r5);
                    \draw[->] (r5) edge [loop right] node {0} (r5);
                    \draw[->] (r5) edge [loop left] node {0} (r5);
                }
                \node () at (38mm,-18mm) {$\leftarrowtail$};
                \node () at (77mm,-18mm) {$\rightarrowtail$};
            \end{tikzpicture}
        \end{center}
Its termination cannot be proved using the weighted type graph method presented in~\cite{zantema2014termination,bruggink2014termination,bruggink2015proving,endrullis2024generalized_icgt} and in Chapter~\ref{chap:weighted_type_graphs}. However, its relative termination with respect to the empty set can be proved using the morphism counting method.

The following sequence of commands selects the system, then uses the morphism counting method without forbidden patterns presented in Chapter~\ref{chap:subgraph_counting} to analyse the relative termination of the system.
\begin{minted}[style=bw]{bash}
    >> select_system_by_name plump18_ex6_rule_copy_variant
    >> subgraph_counting_no_forbidden_context
\end{minted}
The following command summarizes the result of the termination analysis.
 \begin{minted}[style=bw]{bash}
    >> recap
\end{minted}

The message displayed contains \colorbox{Ivory2}{$Terminating:\ yes$} when the method has proved the rule's relative termination with respect to the empty set; otherwise, the message contains \colorbox{Ivory2}{$Terminating:\ unknown$} and \colorbox{Ivory2}{$Remaining\ Rules:$}, followed by a list of rules to be analyzed.

Our implementation automatically tries every subgraph of the left-hand-side graphs of the system's rules as a candidate rule-graph until one succeeds. The chosen rule graph is then displayed in the analysis summary. For the system considered, the summary contains the information, shown below, which describes the employed rule graph.
% \begin{figure}[H]
\begin{center}
    \begin{minted}[style=bw]{bash}
        rule graph : nodes : [ 1;2;3 ]
        arrows : [ (1,s,3,1);(2,s,3,3) ]
    \end{minted}
%     \caption{}
%     \label{fig:kdsljsssmkkggnn}
% \end{figure}
\end{center}
% in Figure~\ref{fig:kdsljsssmkkggnn},  
The ruler-graph is illustrated below.
%  in Figure~\ref{fig:dkgnnnaqqmmamnv}.
% \begin{figure}[H]
%     \centering
\begin{center}
        \begin{tikzpicture}
                \graphbox{$L$}{0mm}{0mm}{35mm}{35mm}{2mm}{-5mm}{
                        \coordinate (delta) at (0,-18mm);
                        \node[draw,circle] (l1) at ($(delta)+(-1,1.5)$) {1};
                        \node[draw,circle] (l2) at ($(delta)+(1,1.5)$) {2};
                        \node[draw,circle] (l3) at ($(delta)+(0,0)$) {3};
                        \draw[->] (l1) -- (l3) node[midway,left] {$s$};
                        \draw[->] (l2) -- (l3) node[midway,right] {$s$};
                }
        \end{tikzpicture}
%     \caption{}
%     \label{fig:dkgnnnaqqmmamnv}
% \end{figure}
    \end{center}

\subsection{Termination Analysis using the Morphism Counting with Antipattern}
Consider the graph rewriting system shown in Example~\ref{antipattern:ex:grs_aca}, whose termination cannot be proven either by the morphism counting method without forbidden patterns introduced in Chapter~\ref{chap:subgraph_counting} or by the subgraph counting method presented in~\cite{overbeek2024termination_lmcs}. Its unique rule is illustrated below.
\begin{center}
    \begin{tikzpicture}
        \graphbox{$L$}{0mm}{0mm}{34mm}{20mm}{2mm}{-5mm}{
            \coordinate (o) at (0mm,-3mm); 
            \node[draw,circle] (l1) at ($(o)+(-10mm,0mm)$) {1};
            \node[draw,circle] (l2) at ($(l1)+(2,0)$) {2};
            \node[draw,circle] (l3) at ($(l1)+(1,0)$) {3};
            \draw[->] (l1) -- (l3) node[midway,above] {$a$};
            \draw[->] (l3) -- (l2) node[midway,above] {$a$};
        }     
        \graphbox{$K$}{40mm}{0mm}{24mm}{20mm}{2mm}{-5mm}{
            \coordinate (o) at (5mm,-3mm); 
            \node[draw,circle] (l1) at ($(o)+(-10mm,0mm)$) {1}; 
            \node[draw,circle] (l2) at ($(l1)+(1,0)$) {2};
        }    
        \graphbox{$R$}{70mm}{0mm}{45mm}{20mm}{2mm}{-5mm}{
            \coordinate (o) at (0mm,-3mm); 
            \node[draw,circle] (l1) at ($(o)+(-10mm,0mm)$) {1};
            \node[draw,circle] (l2) at ($(l1)+(2,0)$) {2};
            \node[draw,circle] (l3) at ($(l1)+(1,0)$) {3};
            \draw[->] (l1) -- (l3) node[midway,above] {$a$};
            \draw[->] (l3) -- (l2) node[midway,above] {$a$};
            \draw[->] (l3) edge [loop below] node {$c$} (l3);
        }    

        \node () at (37mm,-10mm) {$\leftarrowtail$};
        \node () at (67mm,-10mm) {$\rightarrowtail$};
    \end{tikzpicture}
\end{center}
The rule is predefined in the tool and can be selected with the REPL command:
 \begin{minted}[style=bw]{bash}
   >> select_system_by_name bruggink14_ex1
 \end{minted}
The REPL provides the command \colorbox{Ivory2}{$rulergraphs$} to list available ruler-graphs. User-defined ruler-graphs can be added to this list using the method presented in the previous section. The ruler-graph shown in Figure~\ref{fig:implemented:ruler_graph_representation} is predefined in the tool and appears in the list of available ruler-graphs.

The ruler-graph shown in Figure~\ref{fig:implemented:ruler_graph_representation} is predefined and can be selected with the REPL command:
 \begin{minted}[style=bw]{bash}
   >> select_ruler_graph aa_not_in_aca
 \end{minted}
To analyse the system's relative termination using the morphism counting method with a forbidden pattern (see Chapter~\ref{chap:antipattern}) with the ruler-graph, run:
 \begin{minted}[style=bw]{bash}
   >> subgraph_counting_one_forbidden_context
\end{minted}
The following command displays a summary of the termination analysis.
 \begin{minted}[style=bw]{bash}
   >> recap
\end{minted} 
The message displayed contains \colorbox{Ivory2}{$Terminating:\ yes$}, which shows that the method applied successfully proved the relative termination of the rule with respect to the empty set of rules.

\subsection{Termination Analysis using the Type Graph Method}
Consider the rewriting rule shown below:
% in Figure~\ref{fig:implemented:grsaa_0ssssdfss}.
% \begin{figure}[H]
%       \centering 
\begin{center}
      \begin{tikzpicture}
          \graphbox{$L$}{0mm}{0mm}{34mm}{15mm}{2mm}{-5mm}{
              \coordinate (o) at (0mm,-3mm); 
              \node[draw,circle] (l1) at ($(o)+(-10mm,0mm)$) {1};
              \node[draw,circle] (l2) at ($(l1)+(2,0)$) {2};
              \node[draw,circle] (l3) at ($(l1)+(1,0)$) {3};
              \draw[->] (l1) -- (l3) node[midway,above] {$a$};
              \draw[->] (l3) -- (l2) node[midway,above] {$a$};
          }     
          \graphbox{$K$}{40mm}{0mm}{24mm}{15mm}{2mm}{-5mm}{
              \coordinate (o) at (5mm,-3mm); 
              \node[draw,circle] (l1) at ($(o)+(-10mm,0mm)$) {1};
              \node[draw,circle] (l2) at ($(l1)+(1,0)$) {2};
          }    
          \graphbox{$R$}{70mm}{0mm}{45mm}{15mm}{2mm}{-5mm}{
              \coordinate (o) at (-5mm,-3mm); 
              \node[draw,circle] (l1) at ($(o)+(-10mm,0mm)$) {1};
              \node[draw,circle] (l2) at ($(l1)+(3,0)$) {2};
              \node[draw,circle] (l3) at ($(l1)+(1,0)$) {4};
              \node[draw,circle] (l4) at ($(l1)+(2,0)$) {5};
              \draw[->] (l1) -- (l3) node[midway,above] {$a$};
              \draw[->] (l3) -- (l4) node[midway,above] {$b$};
              \draw[->] (l4) -- (l2) node[midway,above] {$a$};
          }    
          \node () at (37mm,-8mm) {$\overset{l}{\leftarrowtail}$};
          \node () at (67mm,-8mm) {$\overset{r}{\rightarrowtail}$};
      \end{tikzpicture}
    \end{center}
The rewriting system consisting of the rule is predefined in the tool and can be selected with the REPL command:
 \begin{minted}[style=bw]{bash}
   >> select_system_by_name bruggink14_ex_4_and_6
\end{minted}
For a selected system, the following command analyses relative termination of some non-empty subset of rules of the system with respect to the set of the remaining rules using the type graph method with the natural arithmetic semiring. The timeout is set to 30 seconds.
 \begin{minted}[style=bw]{bash}
   >> type_graph 30 n
\end{minted}
This command is an instance of the following command pattern
\begin{minted}[style=bw]{bash}
    type_graph [timeout] [semiring_1 semiring_2 ...]
\end{minted}
From the command pattern, we can see that more semirings can be selected. In fact, there are 6 semirings available in the tool, which can be selected by adding flags in the set $\set{a,t,n,A,T,N}$ where $A$, $T$, $N$ stand for the arctic, tropical, and arithmetic semirings over the (extended) natural numbers;  $a$, $t$, $n$ stand for the arctic, tropical, and arithmetic semirings over the (extended) real numbers. If multiple semirings are selected, the tool runs the type graph method in parallel for each chosen semiring. When some rules are removed the tool automatically starts a new iteration on the remaining rules using the same selected semirings. 

For example, we can have the following command which launches the type graph method with the timeout set to 30 seconds and all 6 available semirings.
\begin{minted}[style=bw]{bash} 
    >> type_graph 30 a n t A N T
\end{minted}
Type graph methods with different semirings will be launched in parallel and cooperate with each other. When a non-empty subset $A$ of rules is proved to be terminating relative to the set of remaining rules $B$, a new iteration starts with the set of rules $B$.

The command \colorbox{Ivory2}{$recap$} displays a summary of the termination analysis. The summary contains \colorbox{Ivory2}{$Terminating:\ yes$} and \colorbox{Ivory2}{$Remaining\ Rules:\ none$}, showing that the method successfully proved the rule's relative termination with respect to the empty rule set.
The summary contains also the information shown below.
%  in Figure~\ref{fig:implemented:type_graph_summarydfsaldkfjsdl}. 
% \begin{figure}[H]
\begin{center}
    \begin{minted}[style=bw]{bash} 
        1 rule eliminated : rule 0 
        0 rule remaind : None

        0--a->0 : 0.0
        1--a->1 : 1.0/4.0 
        0--a->1 : 0.0
        0--b->0 : 1.0/16.0 
        1--a->0 : 1.0/8.0 

        0--a->0 : 0
        1--a->1 : 4
        0--a->1 : 0
        0--b->0 : 1
        1--a->0 : 2
    \end{minted}
%     \caption{}
%     \label{fig:implemented:type_graph_summarydfsaldkfjsdl}
% \end{figure}
\end{center}
The first block indicates that the rule whose representation is indexed by \colorbox{Ivory2}{0} (here the system's unique rule) can be eliminated using the type graph method over the real arctic semiring; no rules remain to be analyzed. The second block gives the information for constructing the type graph $\mathcal{T}$: two nodes and 5 edges with rational weights. The third block shows an equivalent type graph $\mathcal{T}'$ with integer weights obtained by multiplying the rational weights by their least common multiple. These type graphs are shown below:
% in Figure~\ref{fig:implemented:type_gdfdgggggg}.
%  \begin{figure}[H]
%     \centering
\begin{center}
    \begin{tikzpicture}
      \graphbox{$\mathcal{T}$}{20mm}{0mm}{45mm}{25mm}{2mm}{-5mm}{
          \coordinate (o) at (3mm,-4mm); 
          \node[draw,circle] (l0) at ($(o)+(-10mm,0mm)$) {$0$};
          \node[draw,circle] (l1) at ($(l0)+(2,0)$) {$1$};
            \draw[->] (l0) edge [loop below] node {$a^{0.0}$} (l0);
            \draw[->] (l0) edge [loop left] node {$b^{\frac{1.0}{16.0}}$} (l0);
            \draw[->] (l1) edge [loop below] node {$a^{\frac{1.0}{4.0}}$} (l1);
          \draw[->] (l0) to[bend left] node[midway,above] {$a^{0.0}$} (l1) ;
          \draw[->] (l1) to[bend left] node[midway,below] {$a^{\frac{1.0}{8.0}}$} (l0) ;
      }   
    \end{tikzpicture}
    \begin{tikzpicture}     
      \graphbox{$\mathcal{T}'$}{70mm}{0mm}{45mm}{25mm}{2mm}{-5mm}{
          \coordinate (o) at (3mm,-4mm); 
          \node[draw,circle] (l0) at ($(o)+(-10mm,0mm)$) {$0$};
          \node[draw,circle] (l1) at ($(l0)+(2,0)$) {$1$};
            \draw[->] (l0) edge [loop below] node {$a^{0}$} (l0);
            \draw[->] (l0) edge [loop left] node {$b^{1}$} (l0);
            \draw[->] (l1) edge [loop below] node {$a^{4}$} (l1);
          \draw[->] (l0) to[bend left] node[midway,above] {$a^{0}$} (l1) ;
          \draw[->] (l1) to[bend left] node[midway,below] {$a^{2}$} (l0);
      }    
  \end{tikzpicture}
\end{center}


% The system can be selected by using the command \colorbox{Ivory2}{select 3} which select the system with number \textit{3}. The REPL will display the following message to confirm the selection:
% \begin{center}
%     \colorbox{Ivory2}{$bruggink\_2014\_ex\_4\_and\_6$ selected}
% \end{center}

\section{Availability and license}

LyonParallel is available at \url{https://github.com/Qi-tchi/LyonParallel/tree/thesis}. The project is distributed under the LGPL-2.1 License.
\chapter{Conclusion and Future Work}
\paragraph{Conclusion}
This dissertation focused on automated techniques for proving the relative termination of DPO graph rewriting systems on edge-labeled directed multigraphs.

We first extended an existing technique, the type graph method, to non-well-founded semirings, enabling more efficient implementation of the method in some cases.
We then developed a new machine-checkable sufficient condition, called morphism counting, for proving relative termination of injective DPO graph rewriting systems on edge-labeled directed multigraphs, based on counting monomorphisms from specific subgraphs to the graph before and after transformation. This technique can handle cases, such as~\autoref{subgraph_counting:ex_contrib_variant}, that existing interpretation-based approaches cannot.
Finally, we extended our morphism counting method to count specific subgraphs that are not included in forbidden contexts. This extension successfully proves termination for systems like the ones presented
in~\autoref{antipattern:ex:grs_aca} and~\autoref{antipattern:ex:endrullis:d3:termination} which prior interpretation-based approaches~\cite{zantema2014termination,bruggink2014termination,bruggink2015proving,endrullis2024generalized_arxiv_v2,overbeek2024termination_lmcs} and the morphism counting method presented in Chapter~\ref{chap:subgraph_counting} fail to handle.

The implementation of the type graph method and our extension, along with our morphism counting technique and its extension, was integrated into a unified tool, called LyonParallel~\footnote{LyonParallel is available at \url{https://github.com/Qi-tchi/LyonParallel/tree/icgt2025}}.
 
\paragraph{Future Work}

Extending the type graph method to other algebraic graph rewriting formalisms, such as PBPO+, could be a promising direction.

Another interesting direction is to extend the morphism counting technique to graph rewriting systems with negative application conditions. Negative application conditions are conditions that must not be satisfied in the context of a rule application. This could lead to a more practical termination tool, as negative application conditions are often used in graph rewriting systems to express constraints on the application of rules.

It is also worth investigating the difference between our morphism counting technique and
 the morphism counting technique developed by Overbeek and Endrullis for PBPO+. 
 
The implementation of the modular termination technique~\cite{plump2018modular} which complements the type graph method and the morphism counting technique, could be another interesting but more practical direction. This could lead to a more powerful termination tool.
  
Finally, confluence is another important property of graph rewriting systems, which ensures that the final result of a sequence of rewriting steps is independent of the order in which the rules are applied. Future work could explore automated techniques for proving confluence of graph rewriting systems.

\printbibliography
\end{document}   